\documentclass[11pt]{article}

\usepackage[margin=1in]{geometry}
\usepackage{amsmath,amssymb,amsthm}
\usepackage{booktabs}
\usepackage{longtable}
\usepackage{enumitem}
\usepackage{hyperref}
\hypersetup{colorlinks=true,linkcolor=blue,urlcolor=blue,citecolor=blue}

\newtheorem{theorem}{Theorem}[section]
\newtheorem{lemma}[theorem]{Lemma}
\newtheorem{corollary}[theorem]{Corollary}
\newtheorem{proposition}[theorem]{Proposition}
\theoremstyle{definition}
\newtheorem{definition}[theorem]{Definition}
\newtheorem{conjecture}[theorem]{Conjecture}
\theoremstyle{remark}
\newtheorem{remark}[theorem]{Remark}

\newcommand{\phigr}{\varphi}
\newcommand{\ZmodQ}{\mathbb{Z}/q\mathbb{Z}}
\newcommand{\lean}[1]{\texttt{#1}}

\newcommand{\nProved}{254}
\newcommand{\nConditional}{4}
\newcommand{\nConjecture}{2}
\newcommand{\nDefinition}{75}
\newcommand{\nTheorems}{258}
\newcommand{\nModules}{32}


\title{\bf A Machine-Verified Core for the Pentagonal Structure of Prime Residues,\\
with a Reproducible Discipline for AI-Assisted Mathematics}
\author{Christopher Brock\\ \small with proofs generated by Aristotle (Harmonic) and independently verified by AXLE (Axiom)}
\date{\today}

\begin{document}
\maketitle

\begin{abstract}
We report a machine-verified core of the \emph{Brockian} program on the pentagonal,
golden-ratio structure of prime residues (the ``Curved Number Line''). The core comprises
\nTheorems{} Lean~4 theorems (\nProved{} unconditional, \nConditional{} conditional) across
\nModules{} modules, checked against Mathlib and, crucially, \emph{independently re-verified}
by a second, cloud-hosted Lean engine. The mathematical highlights are: an exact
admissibility law counting principal residues of prime $k$-tuples ($q-2$ admissible starts
modulo a prime $q$, specialising to the twin constraint mod~$3$ and the Brockian case mod~$5$);
a rigidity theorem showing that the golden value $\phigr-1$ lies in the adjacency spectrum of
the cycle graph $C_p$ for \emph{exactly one} prime, $p=5$; the identity that the algebraic
connectivity of the pentagon is golden, $\lambda_2(C_5)=2-1/\phigr$; a faithful realisation of
the dihedral group $D_5$ inside $\mathrm{Aut}(C_5)$; an exact local-covariance kernel for the
Goldbach residual; and a ``silver'' spectral alphabet governing a sieve on prime constellations.
Equally, we contribute the \emph{method}: an audit ledger of named failure modes, a four-tier
register (\textsc{proved}/\textsc{computation}/\textsc{conditional}/\textsc{conjecture}), a
\emph{triple-verification} gate (local build $+$ axiom check $+$ independent re-check), and a
machine-generated registry that makes it structurally impossible to label a result verified
that the build does not earn. We are deliberate about what is \emph{not} proved: the
Riemann-Hypothesis attack is formalised only as an open schema, and its status is recorded as
such. All artefacts, statements, axioms, and verdicts are public and reproducible.
\end{abstract}

\section{Introduction}

The ``Curved Number Line'' is a family of observations organising the integers on a golden
spiral whose fifth-power symmetry sends prime residues into a small number of rays and
constrains how prime constellations may sit on them. Over several years these observations
accreted into a large body of exploratory mathematics and, more recently, into hundreds of
machine-generated Lean~4 proof attempts. Exploratory AI-assisted mathematics of this kind
faces a specific credibility problem: a proof assistant guarantees a proof, never its
\emph{statement}, and a generative prover will happily return a plausible-looking theorem whose
hypotheses are vacuous, whose operator is secretly the zero map, or whose ``proof'' is a
disguised \texttt{sorry}. Left unaudited, such a corpus is indistinguishable from numerology.

This paper does two things. First (\S\ref{sec:method}) it describes the discipline we used to
separate the real from the decorative: an audit ledger that names each failure mode as it is
encountered, a small register vocabulary applied to every declaration, and a
\emph{triple-verification} gate in which a result is called \textsc{proved} only if it (i)
compiles, (ii) depends on no axioms beyond the three Mathlib admits by default
(\lean{propext}, \lean{Classical.choice}, \lean{Quot.sound}) and uses no compiler-trusting
\lean{native\_decide}, and (iii) is independently re-verified by a \emph{second}, unrelated
Lean engine. Second (\S\S\ref{sec:core}--\ref{sec:analysis}) it presents the mathematical
core that survived this discipline: \nTheorems{} theorems whose exact statements, axiom
footprints, and independent verdicts are published as a machine-generated registry
(Table~\ref{tab:registry}).

We are equally explicit about the program's horizon (\S\ref{sec:open}). The Brockian attack on
the Riemann Hypothesis, via a putative Hilbert--P\'olya operator on the prime-Gaussian
potential, is present here only as a \emph{formalised schema}: its gates are stated, its one
genuinely analytic obligation is reduced to a named classical criterion, and its central
hypothesis is shown to be exactly of open-problem strength. Nothing in this paper should be
read as progress \emph{on} the Riemann Hypothesis; the contribution is a verified substrate
and an honest accounting.

\paragraph{Contributions.}
\begin{enumerate}[leftmargin=*,itemsep=2pt]
  \item A verified core of \nProved{} unconditional and \nConditional{} conditional Lean~4
  theorems on the pentagonal/golden structure of prime residues, each independently re-checked
  (Table~\ref{tab:registry}).
  \item Two rigidity results singling out $p=5$: $\phigr-1\in\mathrm{spec}(C_p)\iff p=5$, and
  $\lambda_2(C_5)=2-1/\phigr$.
  \item An exact finite admissibility law $q-2$ for prime constellations, with its analytic
  companion (positive, summable singular series).
  \item A reproducible discipline --- ledger, registers, triple verification, generated
  registry --- offered as a template for honest AI-assisted mathematics, together with a
  catalogue of the failure modes it is designed to catch.
\end{enumerate}

\section{The verification discipline}\label{sec:method}

\subsection{Registers}
Every declaration in the core carries exactly one \emph{register}, and the register is
\emph{derived} from the compiled artefact, never hand-asserted:
\begin{description}[leftmargin=*,itemsep=1pt]
  \item[\textsc{proved}] sorry-free; \lean{\#print axioms} $\subseteq\{$\lean{propext},
  \lean{Classical.choice}, \lean{Quot.sound}$\}$; no \lean{native\_decide} (which trusts the
  compiler and adds \lean{Lean.ofReduceBool}); and independently re-verified (\S\ref{sec:triple}).
  Kernel-checked \lean{decide} on finite goals is admitted here, as it adds no axioms.
  \item[\textsc{computation}] finite checks that rely on \lean{native\_decide}; recorded as
  computations, never promoted to \textsc{proved}.
  \item[\textsc{conditional}] depends on a named hypothesis, tagged with its \emph{rung}:
  \emph{classical} (a true theorem not yet in Mathlib), \emph{literature} (a published result
  used as a hypothesis with citation), or \emph{open} (an open-problem-strength assumption ---
  admissible only as a schema, never as evidence).
  \item[\textsc{conjecture}] a named \lean{Prop} container, never a theorem.
\end{description}

\subsection{Triple verification}\label{sec:triple}
A \textsc{proved} label requires three independent checks to agree:
(1) a local \lean{lake build} on the pinned toolchain; (2) a local axiom audit via
\lean{\#print axioms}; and (3) an independent re-check by \emph{AXLE}~\cite{axle}, a
cloud-hosted Lean~4 engine that recompiles each proof against Mathlib at a named environment
(\lean{lean-4.32.0}) and, via its \lean{verify\_proof} tool, checks that the proof discharges
the intended statement rather than a drifted one. Because generative provers can self-report
success while smuggling a custom axiom, no engine is trusted on its own word: the three gates
run on every proof regardless of its provenance. A subtle but important point: a \lean{sorry}
raises a \emph{warning}, not an error, so a naive ``did it compile'' check passes it; our gate
treats any sorry/admit warning, and any \lean{sorry}-derived axiom, as a hard failure.

\subsection{The registry as single source of truth}
The three checks feed a machine-generated registry (\lean{registry/theorems.json}), one entry
per declaration recording its statement, module, axiom set, independent verdict, register, and
provenance. The register is \emph{computed} from these facts, so the registry cannot claim more
than the build earns; the table in this paper (Table~\ref{tab:registry}) and the project's
public results page are both generated from it. Declarations rejected during the audit are
recorded, with their named failure mode, in a companion exclusion list --- dropping decorative
theorems silently would defeat the purpose.

\subsection{Named failure modes}
The audit accumulated a catalogue of concrete ways a machine-generated ``theorem'' can be
hollow; each was met at least once in the corpus and is now screened for. Representative
entries: \emph{vacuous bounded-operator hypotheses} (a bounded operator cannot have the
unbounded spectrum a Hilbert--P\'olya operator needs, so the conditional holds vacuously);
\emph{$\mathbb{R}$-mod collapse} (a phase written \texttt{x \% (2*$\pi$)} is identically $0$ in
a field, trivialising every angular claim); \emph{Nat-division exponents} (\texttt{x\^{}(1/2)}
with a natural-number literal is $x^0=1$); \emph{degenerate witnesses} (an existential identity
satisfied by the zero object); \emph{definitional theater} (a term ``proved'' equal to the
definition it unfolds); and \emph{overtitling} (a three-line lemma wearing a famous theorem's
name). These are not hypothetical: the discipline exists because each occurred.

\section{Golden algebra and the ray ring}\label{sec:core}

Write $\phigr=(1+\sqrt5)/2$ for the golden ratio (Mathlib's \lean{Real.goldenRatio}). The
foundational module verifies the golden identities and the arithmetic of \emph{rays}. Residues
modulo $5$ partition the nonzero classes; the reduction $\mathbb{N}\to\mathbb{Z}/5\mathbb{Z}$ is
a ring homomorphism (\lean{ray\_add}, \lean{ray\_mul}), and a value lands on the zero ray iff it
is a multiple of $5$ (\lean{ray\_zero\_iff\_dvd}). The key non-elementary fact is imported from
Mathlib's Dirichlet theorem:

\begin{theorem}[\lean{each\_ray\_has\_infinitely\_many\_primes}, \textsc{proved}]\label{thm:dirichlet}
For each unit $i\in\mathbb{Z}/5\mathbb{Z}$, the set $\{p\ \text{prime}:p\equiv i \pmod 5\}$ is
infinite.
\end{theorem}

\noindent Alongside sit $\phigr^2=\phigr+1$ (\lean{phi\_sq}), $\cos(\pi/5)=\phigr/2$
(\lean{cos\_pi\_div\_five\_eq\_phi\_div\_two}), Binet's formula (\lean{binet\_formula}), and
$5\mid F_{5k}$ (\lean{fib\_five\_dvd}). These are the substrate the later modules build on.

\section{The admissibility law and the transition kernel}\label{sec:admissibility}

For a prime modulus $q$ and a gap $g$, call a residue $a\in\ZmodQ$ an \emph{admissible start}
if neither $a$ nor $a+g$ is $\equiv 0$, i.e.\ $a\notin\{0,-g\}$.

\begin{theorem}[\lean{universal\_admissibility\_count}, \textsc{proved}]\label{thm:qnu}
For any modulus $q\ge 1$ and gap $g\neq 0$ in $\ZmodQ$, the number of admissible starts is
exactly $q-2$.
\end{theorem}

\begin{corollary}[\lean{admissibility\_count\_three}, \lean{admissibility\_count\_five}]
Modulo $3$ there is exactly $1$ admissible start (the twin-prime constraint); modulo $5$ there
are exactly $3$ (the Brockian case).
\end{corollary}

This finite law is repackaged as a transition kernel on the active residues. Defining
\lean{kernel q g} to count admissible starts of each label stepping to the next, the double sum
over the kernel equals the admissible-start count:

\begin{theorem}[\lean{totalSum\_eq}, \textsc{proved}]
For $q\ge1$ and any gap $g$, $\ \sum_i\sum_j \lean{kernel}\,q\,g\,i\,j
=\lvert\lean{admissibleStarts}\,q\,g\rvert$, and this equals $q-2$ for $g\neq0$.
\end{theorem}

\noindent The kernel specialises to the constellation classification: the twin, cousin, sexy,
and quadruplet configurations pin (or free) residues exactly as recorded by
\lean{twin\_pins\_mod\_three}, \lean{cousin\_pins\_mod\_three}, \lean{sexy\_free\_mod\_three},
\lean{quadruplet\_pins\_mod\_five}, and the forbidden twin transition
(\lean{forbidden\_transition}). On the analytic side, the Hardy--Littlewood local factor built
from the same invariant $\nu_p$ is positive and its finite product is positive under
admissibility (\lean{local\_factor\_pos}, \lean{singular\_series\_finite\_pos},
\lean{local\_factor\_asymptotic}); positivity of the full singular series is recorded as
\textsc{conditional} on convergence to a positive limit --- a classical analytic fact stated as
an explicit hypothesis rather than assumed as an axiom:

\begin{theorem}[\lean{singular\_series\_pos}, \textsc{conditional} (rung: classical)]
If the finite singular-series products converge to some $S>0$, then the singular series of $G$
is positive.
\end{theorem}

\section{Spectral rigidity of the pentagon}\label{sec:spectral}

Why five? Two verified rigidity theorems answer this in the language of the cycle graph. Model
the adjacency spectrum of $C_n$ concretely as
$\lean{cycleSpectrum}\,n=\{\,\mu : \exists k,\ \mu=2\cos(2\pi k/n)\,\}$ (its circulant
eigenvalues).

\begin{theorem}[\lean{golden\_unique\_to\_five}, \textsc{proved}]\label{thm:unique}
For every prime $p$, $\ \phigr-1\in\lean{cycleSpectrum}\,p \iff p=5$.
\end{theorem}

\noindent The forward direction is the content: from $\phigr-1=2\cos(2\pi k/p)=2\cos(2\pi/5)$
one deduces $5k=p(5m\pm1)$ for some integer $m$, hence $5\mid p$, hence $p=5$. Thus among all
primes, exactly one pentagon ``sings'' the golden ratio. The companion
\lean{neg\_golden\_in\_C5\_spectrum} records $-\phigr=2\cos(4\pi/5)\in\mathrm{spec}(C_5)$.

The second rigidity is metric. The Laplacian of the $2$-regular graph $C_5$ has eigenvalues
$2-2\cos(2\pi k/5)$; its second-smallest (the algebraic connectivity) is golden:

\begin{theorem}[\lean{pentagon\_lambda2\_phi}, \textsc{proved}]\label{thm:lambda2}
$2-1/\phigr$ is a Laplacian eigenvalue of $C_5$, is positive, and is at most the other nonzero
eigenvalue $2-2\cos(4\pi/5)$; hence $\lambda_2(C_5)=2-1/\phigr=3-\phigr$, with ratio
$\lambda_2/2=1-1/(2\phigr)$.
\end{theorem}

Finally, the symmetry group. Mathlib~4.32 provides $C_n$ but no graph-automorphism-group API,
so we build the action by hand. The rotation $x\mapsto x+a$ and reflection $x\mapsto -b-x$ are
verified to preserve adjacency (\lean{rot\_map\_adj}, \lean{refl\_map\_adj}), assembling into a
group homomorphism $D_5\to\mathrm{Aut}(C_5)$ that is injective:

\begin{theorem}[\lean{dihedral\_action\_faithful}, \textsc{proved}]
The dihedral group $D_5$ acts faithfully on $C_5$ by graph automorphisms; consequently
$10\le\lvert\mathrm{Aut}(C_5)\rvert$ (\lean{ten\_le\_card\_aut}).
\end{theorem}

\begin{remark}
The reverse bound $\lvert\mathrm{Aut}(C_5)\rvert\le 10$ --- which would upgrade this to the full
isomorphism $\mathrm{Aut}(C_5)\cong D_5$ --- is \emph{not} proved: Mathlib~4.32 offers no
automorphism enumeration for cycle graphs, and we decline to assert what we cannot verify. The
open direction is recorded as such.
\end{remark}

\section{A Goldbach comb and a silver sieve}\label{sec:analysis}

Two further modules verify exact structure without claiming the hard theorems behind them.

\paragraph{Goldbach local covariance.} For a prime $p$, the number of ordered pairs of nonzero
residues summing to $c$ is $g_p(c)=p-2+[c=0]$ (\lean{gCount\_eq}), which centres to the spike
$[c=0]-1/p$ (\lean{gCount\_centered}). The autocorrelation of the resulting weight is exact:

\begin{theorem}[\lean{local\_covariance}, \textsc{proved}]
For prime $p$ and shift $h$,
$\ \frac1p\sum_{c}g_p(c)\,g_p(c+h)=\bigl((p-1)^2/p\bigr)^2\,K_p(h)$, with $K_p$ lifting when
$p\mid h$ and dipping otherwise.
\end{theorem}

\noindent The passage from this local kernel to the global Goldbach residual is stated only as
a named conjecture (\S\ref{sec:open}); it is not claimed.

\paragraph{The silver alphabet.} A sieve on prime triples yields a $3\times3$ Hamiltonian whose
spectrum is the ``silver'' set: eigenvalues $2-\sqrt2$, $2$, $2+\sqrt2$, with trace $6$ and
determinant $4$ (\lean{H3\_ground}, \lean{H3\_middle}, \lean{H3\_top}, \lean{H3\_trace},
\lean{H3\_det}), and a finite rigidity of the gap set (\lean{silver\_gap\_rigidity\_finite}).
Companion facts fix the twin admissibility count and a phase--depth torus of minimal period $25$
(\lean{tau\_minimal\_period}, \lean{tau\_injective\_on\_period}).

\section{What is \emph{not} proved: the open program}\label{sec:open}

The value of the discipline is as much in its refusals as its theorems.
\begin{itemize}[leftmargin=*,itemsep=2pt]
  \item \textbf{Riemann Hypothesis.} The Hilbert--P\'olya attack is a \emph{schema}. The
  intended operator is unbounded and densely defined (a bounded operator has compact spectrum
  and cannot carry the zeros' unbounded ordinates --- a constraint we learned from a failed
  formalisation). The prime-Gaussian potential is a bounded self-adjoint multiplication
  operator; essential self-adjointness reduces to Weyl's limit-point criterion, a classical
  theorem absent from Mathlib. No claim of progress on RH is made or implied.
  \item \textbf{Goldbach.} The local covariance kernel (\S\ref{sec:analysis}) is exact; its
  transfer to the global residual is a preregistered conjecture with a stated falsifier, never
  a theorem.
  \item \textbf{$\mathrm{Aut}(C_5)\cong D_5$.} The faithful action is proved; the reverse
  cardinality bound is open (\S\ref{sec:spectral}).
\end{itemize}

\section{Reproducibility}\label{sec:repro}

The entire core is public. Each theorem's exact statement, axiom footprint, and independent
AXLE verdict are in \lean{registry/theorems.json}; Table~\ref{tab:registry} is generated from
it. The pinned environment is \lean{leanprover/lean4:v4.32.0} with Mathlib \lean{v4.32.0}.
Independent verification is performed by AXLE at \lean{lean-4.32.0}; local reproduction is
\lean{lake exe cache get \&\& lake build} followed by regenerating the registry.

\begin{remark}[Honest limitation]
At the time of writing, independent verification is carried by AXLE (a third-party cloud
re-check at the exact toolchain); a from-source local \lean{lake build} was not completed in our
environment because the Mathlib cache was unreachable there, and continuous integration runs the
local build where the cache is available. The registry marks the local-build leg accordingly
rather than overstating it.
\end{remark}

\section{Related work}
The formalisation rests on Lean~4 and Mathlib~\cite{mathlib}. Proof generation and porting used
Aristotle~\cite{aristotle}; independent verification used AXLE~\cite{axle}. The methodological
contribution --- registers, triple verification, and a generated registry --- is in the spirit
of reproducible and adversarially-audited formal mathematics, adapted to a workflow in which
proofs are machine-generated and must be treated as untrusted until independently checked.

\section{Conclusion}
A large exploratory corpus, disciplined by a named-failure-mode ledger and a
triple-verification gate, yields a compact, independently-verified core: \nTheorems{}
theorems on the pentagonal structure of prime residues, two of which single out $p=5$ by
spectral rigidity. The same discipline draws a bright line around what remains open. We offer
both the mathematics and the method, and invite others to check every line.

\appendix
\section{The verified registry}
Table~\ref{tab:registry} is generated verbatim from \lean{registry/theorems.json}.
\begin{longtable}{@{}p{0.44\textwidth}llcc@{}}
\caption{The verified core: every PROVED/CONDITIONAL declaration in the registry, its module, register, whether its \texttt{\#print axioms} is clean (only \texttt{propext, Classical.choice, Quot.sound}), and its independent AXLE verdict at \texttt{lean-4.32.0}.}\label{tab:registry}\\
\toprule Theorem & Module & Register & Axioms & AXLE \\ \midrule
\endfirsthead
\toprule Theorem & Module & Register & Axioms & AXLE \\ \midrule \endhead
\texttt{admissibility\_count\_five} & Admissibility & PROVED & \checkmark & verified \\
\texttt{admissibility\_count\_three} & Admissibility & PROVED & \checkmark & verified \\
\texttt{universal\_admissibility\_count} & Admissibility & PROVED & \checkmark & verified \\
\texttt{admissibleResidues\_crt\_card} & CRT & PROVED & \checkmark & verified \\
\texttt{admissibleResidues\_crt\_card\_two\_primes} & CRT & PROVED & \checkmark & verified \\
\texttt{admissible\_count\_three\_five} & CRT & PROVED & \checkmark & verified \\
\texttt{localIntegerTupleAdmissible\_iff\_localNu\_lt} & CriterionScaffold & PROVED & \checkmark & verified \\
\texttt{localNu\_eq\_card\_localResidueSet} & CriterionScaffold & PROVED & \checkmark & verified \\
\texttt{localTupleAdmissible\_iff\_exists\_avoids} & CriterionScaffold & PROVED & \checkmark & verified \\
\texttt{localTupleAdmissible\_iff\_obstruction\_lt} & CriterionScaffold & PROVED & \checkmark & verified \\
\texttt{not\_localTupleAdmissible\_iff\_modulus\_le\_obstruction} & CriterionScaffold & PROVED & \checkmark & verified \\
\texttt{not\_localTupleAdmissible\_iff\_obstruction\_eq\_modulus} & CriterionScaffold & PROVED & \checkmark & verified \\
\texttt{primeLocalAdmissible\_iff\_every\_prime\_has\_local\_start} & CriterionScaffold & PROVED & \checkmark & verified \\
\texttt{residueSet\_card\_le\_modulus} & CriterionScaffold & PROVED & \checkmark & verified \\
\texttt{admissibleTupleResidues\_prodCRT\_card} & AdmissibilityCRTGeneral & PROVED & \checkmark & verified \\
\texttt{admissibleTupleResidues\_prodCRT\_primes\_card} & AdmissibilityCRTGeneral & PROVED & \checkmark & verified \\
\texttt{admissibleTuple\_pi\_card} & AdmissibilityCRTGeneral & PROVED & \checkmark & verified \\
\texttt{admissible\_crt\_count\_fifteen} & AdmissibilityCRTGeneral & PROVED & \checkmark & verified \\
\texttt{admissible\_ktuple\_count\_fifteen\_factors} & AdmissibilityCRTGeneral & PROVED & \checkmark & verified \\
\texttt{neZero\_prod} & AdmissibilityCRTGeneral & PROVED & \checkmark & verified \\
\texttt{pairwise\_coprime\_of\_primes} & AdmissibilityCRTGeneral & PROVED & \checkmark & verified \\
\texttt{admissibility\_count\_dichotomy} & AdmissibilityDiagonal & PROVED & \checkmark & verified \\
\texttt{admissibleResidues\_zero\_eq} & AdmissibilityDiagonal & PROVED & \checkmark & verified \\
\texttt{diagonal\_admissibility\_count} & AdmissibilityDiagonal & PROVED & \checkmark & verified \\
\texttt{diagonal\_admissibility\_count\_of\_eq\_zero} & AdmissibilityDiagonal & PROVED & \checkmark & verified \\
\texttt{diagonal\_count\_five} & AdmissibilityDiagonal & PROVED & \checkmark & verified \\
\texttt{diagonal\_count\_three} & AdmissibilityDiagonal & PROVED & \checkmark & verified \\
\texttt{admissible\_iff\_card\_image\_lt} & AdmissibilityHLCriterion & PROVED & \checkmark & verified \\
\texttt{admissible\_iff\_count\_pos} & AdmissibilityHLCriterion & PROVED & \checkmark & verified \\
\texttt{admissible\_iff\_exists\_avoiding\_start} & AdmissibilityHLCriterion & PROVED & \checkmark & verified \\
\texttt{admissible\_iff\_nu\_lt} & AdmissibilityHLCriterion & PROVED & \checkmark & verified \\
\texttt{admissible\_iff\_nu\_lt\_of\_le\_card} & AdmissibilityHLCriterion & PROVED & \checkmark & verified \\
\texttt{admissible\_zero\_two} & AdmissibilityHLCriterion & PROVED & \checkmark & verified \\
\texttt{not\_admissible\_zero\_two\_four} & AdmissibilityHLCriterion & PROVED & \checkmark & verified \\
\texttt{omitsResidue\_iff\_ne\_univ} & AdmissibilityHLCriterion & PROVED & \checkmark & verified \\
\texttt{omitsResidue\_iff\_nu\_lt} & AdmissibilityHLCriterion & PROVED & \checkmark & verified \\
\texttt{admissibleTupleResidues\_card} & AdmissibilityKTuple & PROVED & \checkmark & verified \\
\texttt{admissibleTupleResidues\_card\_pair} & AdmissibilityKTuple & PROVED & \checkmark & verified \\
\texttt{admissibleTupleResidues\_card\_triple} & AdmissibilityKTuple & PROVED & \checkmark & verified \\
\texttt{admissibleTupleResidues\_crt\_card} & AdmissibilityKTuple & PROVED & \checkmark & verified \\
\texttt{admissibleTupleResidues\_crt\_card\_pair} & AdmissibilityKTuple & PROVED & \checkmark & verified \\
\texttt{admissibleTupleResidues\_pair\_eq} & AdmissibilityKTuple & PROVED & \checkmark & verified \\
\texttt{admissible\_ktuple\_count\_five} & AdmissibilityKTuple & PROVED & \checkmark & verified \\
\texttt{admissible\_ktuple\_count\_three} & AdmissibilityKTuple & PROVED & \checkmark & verified \\
\texttt{crt\_filter\_card} & AdmissibilityKTuple & PROVED & \checkmark & verified \\
\texttt{mem\_admissibleTupleResidues} & AdmissibilityKTuple & PROVED & \checkmark & verified \\
\texttt{additiveAut\_card} & AffineSymmetry & PROVED & \checkmark & verified \\
\texttt{additiveAut\_card\_five} & AffineSymmetry & PROVED & \checkmark & verified \\
\texttt{dihedralToPerm\_card} & AffineSymmetry & PROVED & \checkmark & verified \\
\texttt{dihedralToPerm\_injective} & AffineSymmetry & PROVED & \checkmark & verified \\
\texttt{dihedralToPerm\_range\_le\_affineGroup} & AffineSymmetry & PROVED & \checkmark & verified \\
\texttt{symmetry\_separation} & AffineSymmetry & PROVED & \checkmark & verified \\
\texttt{units\_isCyclic} & AffineSymmetry & PROVED & \checkmark & verified \\
\texttt{amicable\_1184\_1210} & AmicableNumbers & PROVED & \checkmark & verified \\
\texttt{amicable\_220\_284} & AmicableNumbers & PROVED & \checkmark & verified \\
\texttt{amicable\_2620\_2924} & AmicableNumbers & PROVED & \checkmark & verified \\
\texttt{amicable\_not\_perfect} & AmicableNumbers & PROVED & \checkmark & verified \\
\texttt{amicable\_symm} & AmicableNumbers & PROVED & \checkmark & verified \\
\texttt{perfect\_iff\_aliquot\_fixed} & AmicableNumbers & PROVED & \checkmark & verified \\
\texttt{andricaInt\_iff\_sqrt} & AndricaConjecture & PROVED & \checkmark & verified \\
\texttt{andrica\_113\_127} & AndricaConjecture & PROVED & \checkmark & verified \\
\texttt{andrica\_23\_29} & AndricaConjecture & PROVED & \checkmark & verified \\
\texttt{andrica\_2\_3} & AndricaConjecture & PROVED & \checkmark & verified \\
\texttt{andrica\_7\_11} & AndricaConjecture & PROVED & \checkmark & verified \\
\texttt{andrica\_89\_97} & AndricaConjecture & PROVED & \checkmark & verified \\
\texttt{dihedral\_action\_faithful} & Automorphism & PROVED & \checkmark & verified \\
\texttt{mul\_rr} & Automorphism & PROVED & \checkmark & verified \\
\texttt{mul\_rsr} & Automorphism & PROVED & \checkmark & verified \\
\texttt{mul\_srr} & Automorphism & PROVED & \checkmark & verified \\
\texttt{mul\_srsr} & Automorphism & PROVED & \checkmark & verified \\
\texttt{refl\_map\_adj} & Automorphism & PROVED & \checkmark & verified \\
\texttt{rot\_map\_adj} & Automorphism & PROVED & \checkmark & verified \\
\texttt{ten\_le\_card\_aut} & Automorphism & PROVED & \checkmark & verified \\
\texttt{aut\_card\_eq\_ten} & Full & PROVED & \checkmark & verified \\
\texttt{aut\_equiv\_dihedral} & Full & PROVED & \checkmark & verified \\
\texttt{card\_aut\_le\_ten} & Full & PROVED & \checkmark & verified \\
\texttt{dihedralHom\_bijective} & Full & PROVED & \checkmark & verified \\
\texttt{dihedralHom\_surjective} & Full & PROVED & \checkmark & verified \\
\texttt{betrothed\_1050\_1925} & BetrothedNumbers & PROVED & \checkmark & verified \\
\texttt{betrothed\_140\_195} & BetrothedNumbers & PROVED & \checkmark & verified \\
\texttt{betrothed\_140\_195\_opposite\_parity} & BetrothedNumbers & PROVED & \checkmark & verified \\
\texttt{betrothed\_48\_75} & BetrothedNumbers & PROVED & \checkmark & verified \\
\texttt{betrothed\_48\_75\_opposite\_parity} & BetrothedNumbers & PROVED & \checkmark & verified \\
\texttt{betrothed\_ne} & BetrothedNumbers & PROVED & \checkmark & verified \\
\texttt{betrothed\_symm} & BetrothedNumbers & PROVED & \checkmark & verified \\
\texttt{brocard\_11\_13} & BrocardGap & PROVED & \checkmark & verified \\
\texttt{brocard\_13\_17} & BrocardGap & PROVED & \checkmark & verified \\
\texttt{brocard\_3\_5} & BrocardGap & PROVED & \checkmark & verified \\
\texttt{brocard\_5\_7} & BrocardGap & PROVED & \checkmark & verified \\
\texttt{brocard\_7\_11} & BrocardGap & PROVED & \checkmark & verified \\
\texttt{brown\_4\_5} & BrocardProblem & PROVED & \checkmark & verified \\
\texttt{brown\_5\_11} & BrocardProblem & PROVED & \checkmark & verified \\
\texttt{brown\_7\_71} & BrocardProblem & PROVED & \checkmark & verified \\
\texttt{brown\_factorization} & BrocardProblem & PROVED & \checkmark & verified \\
\texttt{brown\_m\_odd} & BrocardProblem & PROVED & \checkmark & verified \\
\texttt{no\_brown\_8\_9\_10} & BrocardProblem & PROVED & \checkmark & verified \\
\texttt{not\_square\_between} & BrocardProblem & PROVED & \checkmark & verified \\
\texttt{c5DistinctEigs\_card} & C5SpectralMultiplicities & PROVED & \checkmark & verified \\
\texttt{c5LapMode\_four} & C5SpectralMultiplicities & PROVED & \checkmark & verified \\
\texttt{c5LapMode\_one} & C5SpectralMultiplicities & PROVED & \checkmark & verified \\
\texttt{c5LapMode\_three} & C5SpectralMultiplicities & PROVED & \checkmark & verified \\
\texttt{c5LapMode\_two} & C5SpectralMultiplicities & PROVED & \checkmark & verified \\
\texttt{c5LapMode\_zero} & C5SpectralMultiplicities & PROVED & \checkmark & verified \\
\texttt{c5LaplacianMultiset\_eq} & C5SpectralMultiplicities & PROVED & \checkmark & verified \\
\texttt{c5Mode\_four} & C5SpectralMultiplicities & PROVED & \checkmark & verified \\
\texttt{c5Mode\_mem\_cycleSpectrum} & C5SpectralMultiplicities & PROVED & \checkmark & verified \\
\texttt{c5Mode\_one} & C5SpectralMultiplicities & PROVED & \checkmark & verified \\
\texttt{c5Mode\_one\_eq\_four} & C5SpectralMultiplicities & PROVED & \checkmark & verified \\
\texttt{c5Mode\_three} & C5SpectralMultiplicities & PROVED & \checkmark & verified \\
\texttt{c5Mode\_two} & C5SpectralMultiplicities & PROVED & \checkmark & verified \\
\texttt{c5Mode\_two\_eq\_three} & C5SpectralMultiplicities & PROVED & \checkmark & verified \\
\texttt{c5Mode\_zero} & C5SpectralMultiplicities & PROVED & \checkmark & verified \\
\texttt{c5SpectrumMultiset\_card} & C5SpectralMultiplicities & PROVED & \checkmark & verified \\
\texttt{c5SpectrumMultiset\_eq} & C5SpectralMultiplicities & PROVED & \checkmark & verified \\
\texttt{c5SpectrumMultiset\_toFinset} & C5SpectralMultiplicities & PROVED & \checkmark & verified \\
\texttt{c5\_arg\_four} & C5SpectralMultiplicities & PROVED & \checkmark & verified \\
\texttt{c5\_arg\_one} & C5SpectralMultiplicities & PROVED & \checkmark & verified \\
\texttt{c5\_arg\_three} & C5SpectralMultiplicities & PROVED & \checkmark & verified \\
\texttt{c5\_arg\_two} & C5SpectralMultiplicities & PROVED & \checkmark & verified \\
\texttt{c5\_eigs\_pairwise\_distinct} & C5SpectralMultiplicities & PROVED & \checkmark & verified \\
\texttt{goldenRatio\_lt\_two} & C5SpectralMultiplicities & PROVED & \checkmark & verified \\
\texttt{golden\_sub\_one\_mem\_C5} & C5SpectralMultiplicities & PROVED & \checkmark & verified \\
\texttt{golden\_unique\_to\_five\_setlevel} & C5SpectralMultiplicities & PROVED & \checkmark & verified \\
\texttt{lap\_gap\_eq\_connectivity} & C5SpectralMultiplicities & PROVED & \checkmark & verified \\
\texttt{lap\_large\_eq} & C5SpectralMultiplicities & PROVED & \checkmark & verified \\
\texttt{mem\_c5SpectrumMultiset\_iff} & C5SpectralMultiplicities & PROVED & \checkmark & verified \\
\texttt{multiplicity\_connectivity\_gap} & C5SpectralMultiplicities & PROVED & \checkmark & verified \\
\texttt{multiplicity\_golden\_sub\_one} & C5SpectralMultiplicities & PROVED & \checkmark & verified \\
\texttt{multiplicity\_lap\_gap} & C5SpectralMultiplicities & PROVED & \checkmark & verified \\
\texttt{multiplicity\_lap\_large} & C5SpectralMultiplicities & PROVED & \checkmark & verified \\
\texttt{multiplicity\_lap\_two\_plus\_phi} & C5SpectralMultiplicities & PROVED & \checkmark & verified \\
\texttt{multiplicity\_lap\_zero} & C5SpectralMultiplicities & PROVED & \checkmark & verified \\
\texttt{multiplicity\_neg\_golden} & C5SpectralMultiplicities & PROVED & \checkmark & verified \\
\texttt{multiplicity\_two} & C5SpectralMultiplicities & PROVED & \checkmark & verified \\
\texttt{neg\_golden\_mem\_C5} & C5SpectralMultiplicities & PROVED & \checkmark & verified \\
\texttt{two\_mem\_C5} & C5SpectralMultiplicities & PROVED & \checkmark & verified \\
\texttt{korselt\_1105} & CarmichaelKorselt & PROVED & \checkmark & verified \\
\texttt{korselt\_1729} & CarmichaelKorselt & PROVED & \checkmark & verified \\
\texttt{korselt\_561} & CarmichaelKorselt & PROVED & \checkmark & verified \\
\texttt{korselt\_odd} & CarmichaelKorselt & PROVED & \checkmark & verified \\
\texttt{sqfree\_prod3} & CarmichaelKorselt & PROVED & \checkmark & verified \\
\texttt{collatz\_four} & CollatzPartial & PROVED & \checkmark & verified \\
\texttt{collatz\_one} & CollatzPartial & PROVED & \checkmark & verified \\
\texttt{collatz\_two} & CollatzPartial & PROVED & \checkmark & verified \\
\texttt{collatz\_two\_mul} & CollatzPartial & PROVED & \checkmark & verified \\
\texttt{descent\_mod4\_one} & CollatzPartial & PROVED & \checkmark & verified \\
\texttt{reaches1\_mul\_pow\_two} & CollatzPartial & PROVED & \checkmark & verified \\
\texttt{reaches1\_one} & CollatzPartial & PROVED & \checkmark & verified \\
\texttt{reaches1\_pow\_two} & CollatzPartial & PROVED & \checkmark & verified \\
\texttt{reaches1\_two\_mul} & CollatzPartial & PROVED & \checkmark & verified \\
\texttt{trivial\_cycle} & CollatzPartial & PROVED & \checkmark & verified \\
\texttt{cos\_2pi\_5} & Connectivity & PROVED & \checkmark & verified \\
\texttt{lambda2\_eq} & Connectivity & PROVED & \checkmark & verified \\
\texttt{one\_div\_phi} & Connectivity & PROVED & \checkmark & verified \\
\texttt{pentagon\_lambda2\_phi} & Connectivity & PROVED & \checkmark & verified \\
\texttt{pentagon\_ratio} & Connectivity & PROVED & \checkmark & verified \\
\texttt{two\_cos\_4pi\_5} & Connectivity & PROVED & \checkmark & verified \\
\texttt{algebraic\_connectivity\_C5} & ConnectivityGoldenBridge & PROVED & \checkmark & verified \\
\texttt{algebraic\_connectivity\_C5\_schema} & ConnectivityGoldenBridge & PROVED & \checkmark & verified \\
\texttt{algebraic\_connectivity\_schema} & ConnectivityGoldenBridge & PROVED & \checkmark & verified \\
\texttt{cos\_two\_pi\_div\_five\_eq} & ConnectivityGoldenBridge & PROVED & \checkmark & verified \\
\texttt{golden\_sub\_one\_eq\_two\_cos} & ConnectivityGoldenBridge & PROVED & \checkmark & verified \\
\texttt{golden\_sub\_one\_mem\_adjacency\_C5} & ConnectivityGoldenBridge & PROVED & \checkmark & verified \\
\texttt{inv\_phi\_eq\_phi\_sub\_one} & ConnectivityGoldenBridge & PROVED & \checkmark & verified \\
\texttt{lambda2\_eq\_cos\_form} & ConnectivityGoldenBridge & PROVED & \checkmark & verified \\
\texttt{lambda2\_eq\_three\_minus\_phi} & ConnectivityGoldenBridge & PROVED & \checkmark & verified \\
\texttt{lambda2\_eq\_two\_minus\_phi\_sub\_one} & ConnectivityGoldenBridge & PROVED & \checkmark & verified \\
\texttt{lambda2\_eq\_two\_minus\_two\_cos} & ConnectivityGoldenBridge & PROVED & \checkmark & verified \\
\texttt{lambda2\_le\_large\_eig} & ConnectivityGoldenBridge & PROVED & \checkmark & verified \\
\texttt{lambda2\_le\_two\_plus\_phi} & ConnectivityGoldenBridge & PROVED & \checkmark & verified \\
\texttt{lambda2\_mem\_laplacianEigs5} & ConnectivityGoldenBridge & PROVED & \checkmark & verified \\
\texttt{lambda2\_over\_degree} & ConnectivityGoldenBridge & PROVED & \checkmark & verified \\
\texttt{lambda2\_pos} & ConnectivityGoldenBridge & PROVED & \checkmark & verified \\
\texttt{lambda2\_triple\_identity} & ConnectivityGoldenBridge & PROVED & \checkmark & verified \\
\texttt{laplacian\_gap\_from\_adjacency\_mode} & ConnectivityGoldenBridge & PROVED & \checkmark & verified \\
\texttt{large\_eig\_eq\_two\_plus\_phi} & ConnectivityGoldenBridge & PROVED & \checkmark & verified \\
\texttt{neg\_phi\_mem\_adjacency\_C5} & ConnectivityGoldenBridge & PROVED & \checkmark & verified \\
\texttt{two\_cos\_four\_pi\_div\_five\_eq\_neg\_phi} & ConnectivityGoldenBridge & PROVED & \checkmark & verified \\
\texttt{two\_cos\_four\_pi\_div\_five\_geometry} & ConnectivityGoldenBridge & PROVED & \checkmark & verified \\
\texttt{binet\_formula} & Core & PROVED & \checkmark & verified \\
\texttt{cos\_2pi\_5} & Core & PROVED & \checkmark & verified \\
\texttt{cos\_pi\_div\_five\_eq\_phi\_div\_two} & Core & PROVED & \checkmark & verified \\
\texttt{each\_ray\_has\_infinitely\_many\_primes} & Core & PROVED & \checkmark & verified \\
\texttt{fib\_five\_dvd} & Core & PROVED & \checkmark & verified \\
\texttt{fifth\_root\_of\_unity} & Core & PROVED & \checkmark & verified \\
\texttt{one\_lt\_phi} & Core & PROVED & \checkmark & verified \\
\texttt{phi\_pos} & Core & PROVED & \checkmark & verified \\
\texttt{phi\_sq} & Core & PROVED & \checkmark & verified \\
\texttt{ray\_add} & Core & PROVED & \checkmark & verified \\
\texttt{ray\_mul} & Core & PROVED & \checkmark & verified \\
\texttt{ray\_ne\_zero\_infinite} & Core & PROVED & \checkmark & verified \\
\texttt{ray\_zero\_iff\_dvd} & Core & PROVED & \checkmark & verified \\
\texttt{aeval\_spectralGen\_five\_X\_sq\_add\_X\_sub\_one} & CosAlgebraicInteger & PROVED & \checkmark & verified \\
\texttt{aeval\_spectralGen\_seven\_cubic7} & CosAlgebraicInteger & PROVED & \checkmark & verified \\
\texttt{aeval\_spectralGen\_three\_X\_add\_one} & CosAlgebraicInteger & PROVED & \checkmark & verified \\
\texttt{degree\_five\_pack} & CosAlgebraicInteger & PROVED & \checkmark & verified \\
\texttt{degree\_seven\_pack} & CosAlgebraicInteger & PROVED & \checkmark & verified \\
\texttt{degree\_three\_pack} & CosAlgebraicInteger & PROVED & \checkmark & verified \\
\texttt{isIntegral\_and\_degree} & CosAlgebraicInteger & PROVED & \checkmark & verified \\
\texttt{isIntegral\_spectralGen} & CosAlgebraicInteger & PROVED & \checkmark & verified \\
\texttt{isIntegral\_spectralGen\_ℚ} & CosAlgebraicInteger & PROVED & \checkmark & verified \\
\texttt{isIntegral\_two\_cos\_two\_pi\_div} & CosAlgebraicInteger & PROVED & \checkmark & verified \\
\texttt{isIntegral\_two\_cos\_two\_pi\_div\_ℚ} & CosAlgebraicInteger & PROVED & \checkmark & verified \\
\texttt{quadratic\_iff\_five\_pack} & CosAlgebraicInteger & PROVED & \checkmark & verified \\
\texttt{real\_subfield\_degree\_pack} & CosAlgebraicInteger & PROVED & \checkmark & verified \\
\texttt{two\_pi\_div\_eq\_rat\_mul\_pi} & CosAlgebraicInteger & PROVED & \checkmark & verified \\
\texttt{coeff\_zero\_minpoly\_five} & CosTraceNorm & PROVED & \checkmark & verified \\
\texttt{coeff\_zero\_minpoly\_seven} & CosTraceNorm & PROVED & \checkmark & verified \\
\texttt{coeff\_zero\_minpoly\_three} & CosTraceNorm & PROVED & \checkmark & verified \\
\texttt{isIntegral\_spectralGen} & CosTraceNorm & PROVED & \checkmark & verified \\
\texttt{isIntegral\_spectralGen\_five\_ℚ} & CosTraceNorm & PROVED & \checkmark & verified \\
\texttt{isIntegral\_spectralGen\_seven\_ℚ} & CosTraceNorm & PROVED & \checkmark & verified \\
\texttt{isIntegral\_spectralGen\_three\_ℚ} & CosTraceNorm & PROVED & \checkmark & verified \\
\texttt{isIntegral\_spectralGen\_ℚ} & CosTraceNorm & PROVED & \checkmark & verified \\
\texttt{isIntegral\_two\_cos\_two\_pi\_div} & CosTraceNorm & PROVED & \checkmark & verified \\
\texttt{isIntegral\_two\_cos\_two\_pi\_div\_ℚ} & CosTraceNorm & PROVED & \checkmark & verified \\
\texttt{minpoly\_five} & CosTraceNorm & PROVED & \checkmark & verified \\
\texttt{minpoly\_seven} & CosTraceNorm & PROVED & \checkmark & verified \\
\texttt{minpoly\_three} & CosTraceNorm & PROVED & \checkmark & verified \\
\texttt{nextCoeff\_minpoly\_five} & CosTraceNorm & PROVED & \checkmark & verified \\
\texttt{nextCoeff\_minpoly\_seven} & CosTraceNorm & PROVED & \checkmark & verified \\
\texttt{nextCoeff\_minpoly\_three} & CosTraceNorm & PROVED & \checkmark & verified \\
\texttt{norm\_adjoin\_gen\_eq\_coeff\_zero} & CosTraceNorm & PROVED & \checkmark & verified \\
\texttt{norm\_spectralGen\_five} & CosTraceNorm & PROVED & \checkmark & verified \\
\texttt{norm\_spectralGen\_seven} & CosTraceNorm & PROVED & \checkmark & verified \\
\texttt{norm\_spectralGen\_three} & CosTraceNorm & PROVED & \checkmark & verified \\
\texttt{trace\_adjoin\_gen\_eq\_neg\_nextCoeff} & CosTraceNorm & PROVED & \checkmark & verified \\
\texttt{trace\_norm\_pack} & CosTraceNorm & PROVED & \checkmark & verified \\
\texttt{trace\_spectralGen\_five} & CosTraceNorm & PROVED & \checkmark & verified \\
\texttt{trace\_spectralGen\_seven} & CosTraceNorm & PROVED & \checkmark & verified \\
\texttt{trace\_spectralGen\_three} & CosTraceNorm & PROVED & \checkmark & verified \\
\texttt{degree\_eightHundredEightyOne} & CosTraceNormEightHundredEightyOne & PROVED & \checkmark & verified \\
\texttt{eightHundredEightyOne\_ne\_two} & CosTraceNormEightHundredEightyOne & PROVED & \checkmark & verified \\
\texttt{eightHundredEightyOne\_pack} & CosTraceNormEightHundredEightyOne & PROVED & \checkmark & verified \\
\texttt{isIntegral\_and\_degree\_eightHundredEightyOne} & CosTraceNormEightHundredEightyOne & PROVED & \checkmark & verified \\
\texttt{isIntegral\_spectralGen\_eightHundredEightyOne} & CosTraceNormEightHundredEightyOne & PROVED & \checkmark & verified \\
\texttt{isIntegral\_spectralGen\_eightHundredEightyOne\_Q} & CosTraceNormEightHundredEightyOne & PROVED & \checkmark & verified \\
\texttt{prime\_eightHundredEightyOne} & CosTraceNormEightHundredEightyOne & PROVED & \checkmark & verified \\
\texttt{degree\_eightHundredEightySeven} & CosTraceNormEightHundredEightySeven & PROVED & \checkmark & verified \\
\texttt{eightHundredEightySeven\_ne\_two} & CosTraceNormEightHundredEightySeven & PROVED & \checkmark & verified \\
\texttt{eightHundredEightySeven\_pack} & CosTraceNormEightHundredEightySeven & PROVED & \checkmark & verified \\
\texttt{isIntegral\_and\_degree\_eightHundredEightySeven} & CosTraceNormEightHundredEightySeven & PROVED & \checkmark & verified \\
\texttt{isIntegral\_spectralGen\_eightHundredEightySeven} & CosTraceNormEightHundredEightySeven & PROVED & \checkmark & verified \\
\texttt{isIntegral\_spectralGen\_eightHundredEightySeven\_Q} & CosTraceNormEightHundredEightySeven & PROVED & \checkmark & verified \\
\texttt{prime\_eightHundredEightySeven} & CosTraceNormEightHundredEightySeven & PROVED & \checkmark & verified \\
\texttt{degree\_eightHundredEightyThree} & CosTraceNormEightHundredEightyThree & PROVED & \checkmark & verified \\
\texttt{eightHundredEightyThree\_ne\_two} & CosTraceNormEightHundredEightyThree & PROVED & \checkmark & verified \\
\texttt{eightHundredEightyThree\_pack} & CosTraceNormEightHundredEightyThree & PROVED & \checkmark & verified \\
\texttt{isIntegral\_and\_degree\_eightHundredEightyThree} & CosTraceNormEightHundredEightyThree & PROVED & \checkmark & verified \\
\texttt{isIntegral\_spectralGen\_eightHundredEightyThree} & CosTraceNormEightHundredEightyThree & PROVED & \checkmark & verified \\
\texttt{isIntegral\_spectralGen\_eightHundredEightyThree\_Q} & CosTraceNormEightHundredEightyThree & PROVED & \checkmark & verified \\
\texttt{prime\_eightHundredEightyThree} & CosTraceNormEightHundredEightyThree & PROVED & \checkmark & verified \\
\texttt{degree\_eightHundredEleven} & CosTraceNormEightHundredEleven & PROVED & \checkmark & verified \\
\texttt{eightHundredEleven\_ne\_two} & CosTraceNormEightHundredEleven & PROVED & \checkmark & verified \\
\texttt{eightHundredEleven\_pack} & CosTraceNormEightHundredEleven & PROVED & \checkmark & verified \\
\texttt{isIntegral\_and\_degree\_eightHundredEleven} & CosTraceNormEightHundredEleven & PROVED & \checkmark & verified \\
\texttt{isIntegral\_spectralGen\_eightHundredEleven} & CosTraceNormEightHundredEleven & PROVED & \checkmark & verified \\
\texttt{isIntegral\_spectralGen\_eightHundredEleven\_Q} & CosTraceNormEightHundredEleven & PROVED & \checkmark & verified \\
\texttt{prime\_eightHundredEleven} & CosTraceNormEightHundredEleven & PROVED & \checkmark & verified \\
\texttt{degree\_eightHundredFiftyNine} & CosTraceNormEightHundredFiftyNine & PROVED & \checkmark & verified \\
\texttt{eightHundredFiftyNine\_ne\_two} & CosTraceNormEightHundredFiftyNine & PROVED & \checkmark & verified \\
\texttt{eightHundredFiftyNine\_pack} & CosTraceNormEightHundredFiftyNine & PROVED & \checkmark & verified \\
\texttt{isIntegral\_and\_degree\_eightHundredFiftyNine} & CosTraceNormEightHundredFiftyNine & PROVED & \checkmark & verified \\
\texttt{isIntegral\_spectralGen\_eightHundredFiftyNine} & CosTraceNormEightHundredFiftyNine & PROVED & \checkmark & verified \\
\texttt{isIntegral\_spectralGen\_eightHundredFiftyNine\_Q} & CosTraceNormEightHundredFiftyNine & PROVED & \checkmark & verified \\
\texttt{prime\_eightHundredFiftyNine} & CosTraceNormEightHundredFiftyNine & PROVED & \checkmark & verified \\
\texttt{degree\_eightHundredFiftySeven} & CosTraceNormEightHundredFiftySeven & PROVED & \checkmark & verified \\
\texttt{eightHundredFiftySeven\_ne\_two} & CosTraceNormEightHundredFiftySeven & PROVED & \checkmark & verified \\
\texttt{eightHundredFiftySeven\_pack} & CosTraceNormEightHundredFiftySeven & PROVED & \checkmark & verified \\
\texttt{isIntegral\_and\_degree\_eightHundredFiftySeven} & CosTraceNormEightHundredFiftySeven & PROVED & \checkmark & verified \\
\texttt{isIntegral\_spectralGen\_eightHundredFiftySeven} & CosTraceNormEightHundredFiftySeven & PROVED & \checkmark & verified \\
\texttt{isIntegral\_spectralGen\_eightHundredFiftySeven\_Q} & CosTraceNormEightHundredFiftySeven & PROVED & \checkmark & verified \\
\texttt{prime\_eightHundredFiftySeven} & CosTraceNormEightHundredFiftySeven & PROVED & \checkmark & verified \\
\texttt{degree\_eightHundredFiftyThree} & CosTraceNormEightHundredFiftyThree & PROVED & \checkmark & verified \\
\texttt{eightHundredFiftyThree\_ne\_two} & CosTraceNormEightHundredFiftyThree & PROVED & \checkmark & verified \\
\texttt{eightHundredFiftyThree\_pack} & CosTraceNormEightHundredFiftyThree & PROVED & \checkmark & verified \\
\texttt{isIntegral\_and\_degree\_eightHundredFiftyThree} & CosTraceNormEightHundredFiftyThree & PROVED & \checkmark & verified \\
\texttt{isIntegral\_spectralGen\_eightHundredFiftyThree} & CosTraceNormEightHundredFiftyThree & PROVED & \checkmark & verified \\
\texttt{isIntegral\_spectralGen\_eightHundredFiftyThree\_Q} & CosTraceNormEightHundredFiftyThree & PROVED & \checkmark & verified \\
\texttt{prime\_eightHundredFiftyThree} & CosTraceNormEightHundredFiftyThree & PROVED & \checkmark & verified \\
\texttt{degree\_eightHundredNine} & CosTraceNormEightHundredNine & PROVED & \checkmark & verified \\
\texttt{eightHundredNine\_ne\_two} & CosTraceNormEightHundredNine & PROVED & \checkmark & verified \\
\texttt{eightHundredNine\_pack} & CosTraceNormEightHundredNine & PROVED & \checkmark & verified \\
\texttt{isIntegral\_and\_degree\_eightHundredNine} & CosTraceNormEightHundredNine & PROVED & \checkmark & verified \\
\texttt{isIntegral\_spectralGen\_eightHundredNine} & CosTraceNormEightHundredNine & PROVED & \checkmark & verified \\
\texttt{isIntegral\_spectralGen\_eightHundredNine\_Q} & CosTraceNormEightHundredNine & PROVED & \checkmark & verified \\
\texttt{prime\_eightHundredNine} & CosTraceNormEightHundredNine & PROVED & \checkmark & verified \\
\texttt{degree\_eightHundredSeventySeven} & CosTraceNormEightHundredSeventySeven & PROVED & \checkmark & verified \\
\texttt{eightHundredSeventySeven\_ne\_two} & CosTraceNormEightHundredSeventySeven & PROVED & \checkmark & verified \\
\texttt{eightHundredSeventySeven\_pack} & CosTraceNormEightHundredSeventySeven & PROVED & \checkmark & verified \\
\texttt{isIntegral\_and\_degree\_eightHundredSeventySeven} & CosTraceNormEightHundredSeventySeven & PROVED & \checkmark & verified \\
\texttt{isIntegral\_spectralGen\_eightHundredSeventySeven} & CosTraceNormEightHundredSeventySeven & PROVED & \checkmark & verified \\
\texttt{isIntegral\_spectralGen\_eightHundredSeventySeven\_Q} & CosTraceNormEightHundredSeventySeven & PROVED & \checkmark & verified \\
\texttt{prime\_eightHundredSeventySeven} & CosTraceNormEightHundredSeventySeven & PROVED & \checkmark & verified \\
\texttt{degree\_eightHundredSixtyThree} & CosTraceNormEightHundredSixtyThree & PROVED & \checkmark & verified \\
\texttt{eightHundredSixtyThree\_ne\_two} & CosTraceNormEightHundredSixtyThree & PROVED & \checkmark & verified \\
\texttt{eightHundredSixtyThree\_pack} & CosTraceNormEightHundredSixtyThree & PROVED & \checkmark & verified \\
\texttt{isIntegral\_and\_degree\_eightHundredSixtyThree} & CosTraceNormEightHundredSixtyThree & PROVED & \checkmark & verified \\
\texttt{isIntegral\_spectralGen\_eightHundredSixtyThree} & CosTraceNormEightHundredSixtyThree & PROVED & \checkmark & verified \\
\texttt{isIntegral\_spectralGen\_eightHundredSixtyThree\_Q} & CosTraceNormEightHundredSixtyThree & PROVED & \checkmark & verified \\
\texttt{prime\_eightHundredSixtyThree} & CosTraceNormEightHundredSixtyThree & PROVED & \checkmark & verified \\
\texttt{degree\_eightHundredThirtyNine} & CosTraceNormEightHundredThirtyNine & PROVED & \checkmark & verified \\
\texttt{eightHundredThirtyNine\_ne\_two} & CosTraceNormEightHundredThirtyNine & PROVED & \checkmark & verified \\
\texttt{eightHundredThirtyNine\_pack} & CosTraceNormEightHundredThirtyNine & PROVED & \checkmark & verified \\
\texttt{isIntegral\_and\_degree\_eightHundredThirtyNine} & CosTraceNormEightHundredThirtyNine & PROVED & \checkmark & verified \\
\texttt{isIntegral\_spectralGen\_eightHundredThirtyNine} & CosTraceNormEightHundredThirtyNine & PROVED & \checkmark & verified \\
\texttt{isIntegral\_spectralGen\_eightHundredThirtyNine\_Q} & CosTraceNormEightHundredThirtyNine & PROVED & \checkmark & verified \\
\texttt{prime\_eightHundredThirtyNine} & CosTraceNormEightHundredThirtyNine & PROVED & \checkmark & verified \\
\texttt{degree\_eightHundredTwentyNine} & CosTraceNormEightHundredTwentyNine & PROVED & \checkmark & verified \\
\texttt{eightHundredTwentyNine\_ne\_two} & CosTraceNormEightHundredTwentyNine & PROVED & \checkmark & verified \\
\texttt{eightHundredTwentyNine\_pack} & CosTraceNormEightHundredTwentyNine & PROVED & \checkmark & verified \\
\texttt{isIntegral\_and\_degree\_eightHundredTwentyNine} & CosTraceNormEightHundredTwentyNine & PROVED & \checkmark & verified \\
\texttt{isIntegral\_spectralGen\_eightHundredTwentyNine} & CosTraceNormEightHundredTwentyNine & PROVED & \checkmark & verified \\
\texttt{isIntegral\_spectralGen\_eightHundredTwentyNine\_Q} & CosTraceNormEightHundredTwentyNine & PROVED & \checkmark & verified \\
\texttt{prime\_eightHundredTwentyNine} & CosTraceNormEightHundredTwentyNine & PROVED & \checkmark & verified \\
\texttt{degree\_eightHundredTwentyOne} & CosTraceNormEightHundredTwentyOne & PROVED & \checkmark & verified \\
\texttt{eightHundredTwentyOne\_ne\_two} & CosTraceNormEightHundredTwentyOne & PROVED & \checkmark & verified \\
\texttt{eightHundredTwentyOne\_pack} & CosTraceNormEightHundredTwentyOne & PROVED & \checkmark & verified \\
\texttt{isIntegral\_and\_degree\_eightHundredTwentyOne} & CosTraceNormEightHundredTwentyOne & PROVED & \checkmark & verified \\
\texttt{isIntegral\_spectralGen\_eightHundredTwentyOne} & CosTraceNormEightHundredTwentyOne & PROVED & \checkmark & verified \\
\texttt{isIntegral\_spectralGen\_eightHundredTwentyOne\_Q} & CosTraceNormEightHundredTwentyOne & PROVED & \checkmark & verified \\
\texttt{prime\_eightHundredTwentyOne} & CosTraceNormEightHundredTwentyOne & PROVED & \checkmark & verified \\
\texttt{degree\_eightHundredTwentySeven} & CosTraceNormEightHundredTwentySeven & PROVED & \checkmark & verified \\
\texttt{eightHundredTwentySeven\_ne\_two} & CosTraceNormEightHundredTwentySeven & PROVED & \checkmark & verified \\
\texttt{eightHundredTwentySeven\_pack} & CosTraceNormEightHundredTwentySeven & PROVED & \checkmark & verified \\
\texttt{isIntegral\_and\_degree\_eightHundredTwentySeven} & CosTraceNormEightHundredTwentySeven & PROVED & \checkmark & verified \\
\texttt{isIntegral\_spectralGen\_eightHundredTwentySeven} & CosTraceNormEightHundredTwentySeven & PROVED & \checkmark & verified \\
\texttt{isIntegral\_spectralGen\_eightHundredTwentySeven\_Q} & CosTraceNormEightHundredTwentySeven & PROVED & \checkmark & verified \\
\texttt{prime\_eightHundredTwentySeven} & CosTraceNormEightHundredTwentySeven & PROVED & \checkmark & verified \\
\texttt{degree\_eightHundredTwentyThree} & CosTraceNormEightHundredTwentyThree & PROVED & \checkmark & verified \\
\texttt{eightHundredTwentyThree\_ne\_two} & CosTraceNormEightHundredTwentyThree & PROVED & \checkmark & verified \\
\texttt{eightHundredTwentyThree\_pack} & CosTraceNormEightHundredTwentyThree & PROVED & \checkmark & verified \\
\texttt{isIntegral\_and\_degree\_eightHundredTwentyThree} & CosTraceNormEightHundredTwentyThree & PROVED & \checkmark & verified \\
\texttt{isIntegral\_spectralGen\_eightHundredTwentyThree} & CosTraceNormEightHundredTwentyThree & PROVED & \checkmark & verified \\
\texttt{isIntegral\_spectralGen\_eightHundredTwentyThree\_Q} & CosTraceNormEightHundredTwentyThree & PROVED & \checkmark & verified \\
\texttt{prime\_eightHundredTwentyThree} & CosTraceNormEightHundredTwentyThree & PROVED & \checkmark & verified \\
\texttt{degree\_eightyNine} & CosTraceNormEightyNine & PROVED & \checkmark & verified \\
\texttt{eightyNine\_ne\_two} & CosTraceNormEightyNine & PROVED & \checkmark & verified \\
\texttt{eightyNine\_pack} & CosTraceNormEightyNine & PROVED & \checkmark & verified \\
\texttt{isIntegral\_and\_degree\_eightyNine} & CosTraceNormEightyNine & PROVED & \checkmark & verified \\
\texttt{isIntegral\_spectralGen\_eightyNine} & CosTraceNormEightyNine & PROVED & \checkmark & verified \\
\texttt{isIntegral\_spectralGen\_eightyNine\_Q} & CosTraceNormEightyNine & PROVED & \checkmark & verified \\
\texttt{prime\_eightyNine} & CosTraceNormEightyNine & PROVED & \checkmark & verified \\
\texttt{degree\_eightyThree} & CosTraceNormEightyThree & PROVED & \checkmark & verified \\
\texttt{eightyThree\_ne\_two} & CosTraceNormEightyThree & PROVED & \checkmark & verified \\
\texttt{eightyThree\_pack} & CosTraceNormEightyThree & PROVED & \checkmark & verified \\
\texttt{isIntegral\_and\_degree\_eightyThree} & CosTraceNormEightyThree & PROVED & \checkmark & verified \\
\texttt{isIntegral\_spectralGen\_eightyThree} & CosTraceNormEightyThree & PROVED & \checkmark & verified \\
\texttt{isIntegral\_spectralGen\_eightyThree\_Q} & CosTraceNormEightyThree & PROVED & \checkmark & verified \\
\texttt{prime\_eightyThree} & CosTraceNormEightyThree & PROVED & \checkmark & verified \\
\texttt{P11\_monic} & CosTraceNormEleven & PROVED & \checkmark & verified \\
\texttt{P11\_natDegree} & CosTraceNormEleven & PROVED & \checkmark & verified \\
\texttt{Psi\_eleven} & CosTraceNormEleven & PROVED & \checkmark & verified \\
\texttt{coeff\_zero\_minpoly\_eleven} & CosTraceNormEleven & PROVED & \checkmark & verified \\
\texttt{degree\_eleven} & CosTraceNormEleven & PROVED & \checkmark & verified \\
\texttt{degree\_eleven\_pack} & CosTraceNormEleven & PROVED & \checkmark & verified \\
\texttt{degree\_odd\_prime} & CosTraceNormEleven & PROVED & \checkmark & verified \\
\texttt{eleven\_ne\_two} & CosTraceNormEleven & PROVED & \checkmark & verified \\
\texttt{eleven\_pack} & CosTraceNormEleven & PROVED & \checkmark & verified \\
\texttt{isIntegral\_and\_degree\_eleven} & CosTraceNormEleven & PROVED & \checkmark & verified \\
\texttt{isIntegral\_and\_degree\_odd} & CosTraceNormEleven & PROVED & \checkmark & verified \\
\texttt{isIntegral\_spectralGen\_eleven} & CosTraceNormEleven & PROVED & \checkmark & verified \\
\texttt{isIntegral\_spectralGen\_eleven\_ℚ} & CosTraceNormEleven & PROVED & \checkmark & verified \\
\texttt{isIntegral\_spectralGen\_odd} & CosTraceNormEleven & PROVED & \checkmark & verified \\
\texttt{isIntegral\_spectralGen\_odd\_ℚ} & CosTraceNormEleven & PROVED & \checkmark & verified \\
\texttt{minpoly\_eleven} & CosTraceNormEleven & PROVED & \checkmark & verified \\
\texttt{minpoly\_eq\_Psi} & CosTraceNormEleven & PROVED & \checkmark & verified \\
\texttt{nextCoeff\_minpoly\_eleven} & CosTraceNormEleven & PROVED & \checkmark & verified \\
\texttt{norm\_adjoin\_gen\_eq\_coeff\_zero} & CosTraceNormEleven & PROVED & \checkmark & verified \\
\texttt{norm\_spectralGen\_eleven} & CosTraceNormEleven & PROVED & \checkmark & verified \\
\texttt{prime\_eleven} & CosTraceNormEleven & PROVED & \checkmark & verified \\
\texttt{trace\_adjoin\_gen\_eq\_neg\_nextCoeff} & CosTraceNormEleven & PROVED & \checkmark & verified \\
\texttt{trace\_norm\_eleven\_pack} & CosTraceNormEleven & PROVED & \checkmark & verified \\
\texttt{trace\_spectralGen\_eleven} & CosTraceNormEleven & PROVED & \checkmark & verified \\
\texttt{degree\_fiftyNine} & CosTraceNormFiftyNine & PROVED & \checkmark & verified \\
\texttt{fiftyNine\_ne\_two} & CosTraceNormFiftyNine & PROVED & \checkmark & verified \\
\texttt{fiftyNine\_pack} & CosTraceNormFiftyNine & PROVED & \checkmark & verified \\
\texttt{isIntegral\_and\_degree\_fiftyNine} & CosTraceNormFiftyNine & PROVED & \checkmark & verified \\
\texttt{isIntegral\_spectralGen\_fiftyNine} & CosTraceNormFiftyNine & PROVED & \checkmark & verified \\
\texttt{isIntegral\_spectralGen\_fiftyNine\_Q} & CosTraceNormFiftyNine & PROVED & \checkmark & verified \\
\texttt{prime\_fiftyNine} & CosTraceNormFiftyNine & PROVED & \checkmark & verified \\
\texttt{degree\_fiftyThree} & CosTraceNormFiftyThree & PROVED & \checkmark & verified \\
\texttt{fiftyThree\_ne\_two} & CosTraceNormFiftyThree & PROVED & \checkmark & verified \\
\texttt{fiftyThree\_pack} & CosTraceNormFiftyThree & PROVED & \checkmark & verified \\
\texttt{isIntegral\_and\_degree\_fiftyThree} & CosTraceNormFiftyThree & PROVED & \checkmark & verified \\
\texttt{isIntegral\_spectralGen\_fiftyThree} & CosTraceNormFiftyThree & PROVED & \checkmark & verified \\
\texttt{isIntegral\_spectralGen\_fiftyThree\_Q} & CosTraceNormFiftyThree & PROVED & \checkmark & verified \\
\texttt{prime\_fiftyThree} & CosTraceNormFiftyThree & PROVED & \checkmark & verified \\
\texttt{degree\_fiveHundredEightySeven} & CosTraceNormFiveHundredEightySeven & PROVED & \checkmark & verified \\
\texttt{fiveHundredEightySeven\_ne\_two} & CosTraceNormFiveHundredEightySeven & PROVED & \checkmark & verified \\
\texttt{fiveHundredEightySeven\_pack} & CosTraceNormFiveHundredEightySeven & PROVED & \checkmark & verified \\
\texttt{isIntegral\_and\_degree\_fiveHundredEightySeven} & CosTraceNormFiveHundredEightySeven & PROVED & \checkmark & verified \\
\texttt{isIntegral\_spectralGen\_fiveHundredEightySeven} & CosTraceNormFiveHundredEightySeven & PROVED & \checkmark & verified \\
\texttt{isIntegral\_spectralGen\_fiveHundredEightySeven\_Q} & CosTraceNormFiveHundredEightySeven & PROVED & \checkmark & verified \\
\texttt{prime\_fiveHundredEightySeven} & CosTraceNormFiveHundredEightySeven & PROVED & \checkmark & verified \\
\texttt{degree\_fiveHundredFiftySeven} & CosTraceNormFiveHundredFiftySeven & PROVED & \checkmark & verified \\
\texttt{fiveHundredFiftySeven\_ne\_two} & CosTraceNormFiveHundredFiftySeven & PROVED & \checkmark & verified \\
\texttt{fiveHundredFiftySeven\_pack} & CosTraceNormFiveHundredFiftySeven & PROVED & \checkmark & verified \\
\texttt{isIntegral\_and\_degree\_fiveHundredFiftySeven} & CosTraceNormFiveHundredFiftySeven & PROVED & \checkmark & verified \\
\texttt{isIntegral\_spectralGen\_fiveHundredFiftySeven} & CosTraceNormFiveHundredFiftySeven & PROVED & \checkmark & verified \\
\texttt{isIntegral\_spectralGen\_fiveHundredFiftySeven\_Q} & CosTraceNormFiveHundredFiftySeven & PROVED & \checkmark & verified \\
\texttt{prime\_fiveHundredFiftySeven} & CosTraceNormFiveHundredFiftySeven & PROVED & \checkmark & verified \\
\texttt{degree\_fiveHundredFortyOne} & CosTraceNormFiveHundredFortyOne & PROVED & \checkmark & verified \\
\texttt{fiveHundredFortyOne\_ne\_two} & CosTraceNormFiveHundredFortyOne & PROVED & \checkmark & verified \\
\texttt{fiveHundredFortyOne\_pack} & CosTraceNormFiveHundredFortyOne & PROVED & \checkmark & verified \\
\texttt{isIntegral\_and\_degree\_fiveHundredFortyOne} & CosTraceNormFiveHundredFortyOne & PROVED & \checkmark & verified \\
\texttt{isIntegral\_spectralGen\_fiveHundredFortyOne} & CosTraceNormFiveHundredFortyOne & PROVED & \checkmark & verified \\
\texttt{isIntegral\_spectralGen\_fiveHundredFortyOne\_Q} & CosTraceNormFiveHundredFortyOne & PROVED & \checkmark & verified \\
\texttt{prime\_fiveHundredFortyOne} & CosTraceNormFiveHundredFortyOne & PROVED & \checkmark & verified \\
\texttt{degree\_fiveHundredFortySeven} & CosTraceNormFiveHundredFortySeven & PROVED & \checkmark & verified \\
\texttt{fiveHundredFortySeven\_ne\_two} & CosTraceNormFiveHundredFortySeven & PROVED & \checkmark & verified \\
\texttt{fiveHundredFortySeven\_pack} & CosTraceNormFiveHundredFortySeven & PROVED & \checkmark & verified \\
\texttt{isIntegral\_and\_degree\_fiveHundredFortySeven} & CosTraceNormFiveHundredFortySeven & PROVED & \checkmark & verified \\
\texttt{isIntegral\_spectralGen\_fiveHundredFortySeven} & CosTraceNormFiveHundredFortySeven & PROVED & \checkmark & verified \\
\texttt{isIntegral\_spectralGen\_fiveHundredFortySeven\_Q} & CosTraceNormFiveHundredFortySeven & PROVED & \checkmark & verified \\
\texttt{prime\_fiveHundredFortySeven} & CosTraceNormFiveHundredFortySeven & PROVED & \checkmark & verified \\
\texttt{degree\_fiveHundredNine} & CosTraceNormFiveHundredNine & PROVED & \checkmark & verified \\
\texttt{fiveHundredNine\_ne\_two} & CosTraceNormFiveHundredNine & PROVED & \checkmark & verified \\
\texttt{fiveHundredNine\_pack} & CosTraceNormFiveHundredNine & PROVED & \checkmark & verified \\
\texttt{isIntegral\_and\_degree\_fiveHundredNine} & CosTraceNormFiveHundredNine & PROVED & \checkmark & verified \\
\texttt{isIntegral\_spectralGen\_fiveHundredNine} & CosTraceNormFiveHundredNine & PROVED & \checkmark & verified \\
\texttt{isIntegral\_spectralGen\_fiveHundredNine\_Q} & CosTraceNormFiveHundredNine & PROVED & \checkmark & verified \\
\texttt{prime\_fiveHundredNine} & CosTraceNormFiveHundredNine & PROVED & \checkmark & verified \\
\texttt{degree\_fiveHundredNinetyNine} & CosTraceNormFiveHundredNinetyNine & PROVED & \checkmark & verified \\
\texttt{fiveHundredNinetyNine\_ne\_two} & CosTraceNormFiveHundredNinetyNine & PROVED & \checkmark & verified \\
\texttt{fiveHundredNinetyNine\_pack} & CosTraceNormFiveHundredNinetyNine & PROVED & \checkmark & verified \\
\texttt{isIntegral\_and\_degree\_fiveHundredNinetyNine} & CosTraceNormFiveHundredNinetyNine & PROVED & \checkmark & verified \\
\texttt{isIntegral\_spectralGen\_fiveHundredNinetyNine} & CosTraceNormFiveHundredNinetyNine & PROVED & \checkmark & verified \\
\texttt{isIntegral\_spectralGen\_fiveHundredNinetyNine\_Q} & CosTraceNormFiveHundredNinetyNine & PROVED & \checkmark & verified \\
\texttt{prime\_fiveHundredNinetyNine} & CosTraceNormFiveHundredNinetyNine & PROVED & \checkmark & verified \\
\texttt{degree\_fiveHundredNinetyThree} & CosTraceNormFiveHundredNinetyThree & PROVED & \checkmark & verified \\
\texttt{fiveHundredNinetyThree\_ne\_two} & CosTraceNormFiveHundredNinetyThree & PROVED & \checkmark & verified \\
\texttt{fiveHundredNinetyThree\_pack} & CosTraceNormFiveHundredNinetyThree & PROVED & \checkmark & verified \\
\texttt{isIntegral\_and\_degree\_fiveHundredNinetyThree} & CosTraceNormFiveHundredNinetyThree & PROVED & \checkmark & verified \\
\texttt{isIntegral\_spectralGen\_fiveHundredNinetyThree} & CosTraceNormFiveHundredNinetyThree & PROVED & \checkmark & verified \\
\texttt{isIntegral\_spectralGen\_fiveHundredNinetyThree\_Q} & CosTraceNormFiveHundredNinetyThree & PROVED & \checkmark & verified \\
\texttt{prime\_fiveHundredNinetyThree} & CosTraceNormFiveHundredNinetyThree & PROVED & \checkmark & verified \\
\texttt{degree\_fiveHundredSeventyOne} & CosTraceNormFiveHundredSeventyOne & PROVED & \checkmark & verified \\
\texttt{fiveHundredSeventyOne\_ne\_two} & CosTraceNormFiveHundredSeventyOne & PROVED & \checkmark & verified \\
\texttt{fiveHundredSeventyOne\_pack} & CosTraceNormFiveHundredSeventyOne & PROVED & \checkmark & verified \\
\texttt{isIntegral\_and\_degree\_fiveHundredSeventyOne} & CosTraceNormFiveHundredSeventyOne & PROVED & \checkmark & verified \\
\texttt{isIntegral\_spectralGen\_fiveHundredSeventyOne} & CosTraceNormFiveHundredSeventyOne & PROVED & \checkmark & verified \\
\texttt{isIntegral\_spectralGen\_fiveHundredSeventyOne\_Q} & CosTraceNormFiveHundredSeventyOne & PROVED & \checkmark & verified \\
\texttt{prime\_fiveHundredSeventyOne} & CosTraceNormFiveHundredSeventyOne & PROVED & \checkmark & verified \\
\texttt{degree\_fiveHundredSeventySeven} & CosTraceNormFiveHundredSeventySeven & PROVED & \checkmark & verified \\
\texttt{fiveHundredSeventySeven\_ne\_two} & CosTraceNormFiveHundredSeventySeven & PROVED & \checkmark & verified \\
\texttt{fiveHundredSeventySeven\_pack} & CosTraceNormFiveHundredSeventySeven & PROVED & \checkmark & verified \\
\texttt{isIntegral\_and\_degree\_fiveHundredSeventySeven} & CosTraceNormFiveHundredSeventySeven & PROVED & \checkmark & verified \\
\texttt{isIntegral\_spectralGen\_fiveHundredSeventySeven} & CosTraceNormFiveHundredSeventySeven & PROVED & \checkmark & verified \\
\texttt{isIntegral\_spectralGen\_fiveHundredSeventySeven\_Q} & CosTraceNormFiveHundredSeventySeven & PROVED & \checkmark & verified \\
\texttt{prime\_fiveHundredSeventySeven} & CosTraceNormFiveHundredSeventySeven & PROVED & \checkmark & verified \\
\texttt{degree\_fiveHundredSixtyNine} & CosTraceNormFiveHundredSixtyNine & PROVED & \checkmark & verified \\
\texttt{fiveHundredSixtyNine\_ne\_two} & CosTraceNormFiveHundredSixtyNine & PROVED & \checkmark & verified \\
\texttt{fiveHundredSixtyNine\_pack} & CosTraceNormFiveHundredSixtyNine & PROVED & \checkmark & verified \\
\texttt{isIntegral\_and\_degree\_fiveHundredSixtyNine} & CosTraceNormFiveHundredSixtyNine & PROVED & \checkmark & verified \\
\texttt{isIntegral\_spectralGen\_fiveHundredSixtyNine} & CosTraceNormFiveHundredSixtyNine & PROVED & \checkmark & verified \\
\texttt{isIntegral\_spectralGen\_fiveHundredSixtyNine\_Q} & CosTraceNormFiveHundredSixtyNine & PROVED & \checkmark & verified \\
\texttt{prime\_fiveHundredSixtyNine} & CosTraceNormFiveHundredSixtyNine & PROVED & \checkmark & verified \\
\texttt{degree\_fiveHundredSixtyThree} & CosTraceNormFiveHundredSixtyThree & PROVED & \checkmark & verified \\
\texttt{fiveHundredSixtyThree\_ne\_two} & CosTraceNormFiveHundredSixtyThree & PROVED & \checkmark & verified \\
\texttt{fiveHundredSixtyThree\_pack} & CosTraceNormFiveHundredSixtyThree & PROVED & \checkmark & verified \\
\texttt{isIntegral\_and\_degree\_fiveHundredSixtyThree} & CosTraceNormFiveHundredSixtyThree & PROVED & \checkmark & verified \\
\texttt{isIntegral\_spectralGen\_fiveHundredSixtyThree} & CosTraceNormFiveHundredSixtyThree & PROVED & \checkmark & verified \\
\texttt{isIntegral\_spectralGen\_fiveHundredSixtyThree\_Q} & CosTraceNormFiveHundredSixtyThree & PROVED & \checkmark & verified \\
\texttt{prime\_fiveHundredSixtyThree} & CosTraceNormFiveHundredSixtyThree & PROVED & \checkmark & verified \\
\texttt{degree\_fiveHundredThree} & CosTraceNormFiveHundredThree & PROVED & \checkmark & verified \\
\texttt{fiveHundredThree\_ne\_two} & CosTraceNormFiveHundredThree & PROVED & \checkmark & verified \\
\texttt{fiveHundredThree\_pack} & CosTraceNormFiveHundredThree & PROVED & \checkmark & verified \\
\texttt{isIntegral\_and\_degree\_fiveHundredThree} & CosTraceNormFiveHundredThree & PROVED & \checkmark & verified \\
\texttt{isIntegral\_spectralGen\_fiveHundredThree} & CosTraceNormFiveHundredThree & PROVED & \checkmark & verified \\
\texttt{isIntegral\_spectralGen\_fiveHundredThree\_Q} & CosTraceNormFiveHundredThree & PROVED & \checkmark & verified \\
\texttt{prime\_fiveHundredThree} & CosTraceNormFiveHundredThree & PROVED & \checkmark & verified \\
\texttt{degree\_fiveHundredTwentyOne} & CosTraceNormFiveHundredTwentyOne & PROVED & \checkmark & verified \\
\texttt{fiveHundredTwentyOne\_ne\_two} & CosTraceNormFiveHundredTwentyOne & PROVED & \checkmark & verified \\
\texttt{fiveHundredTwentyOne\_pack} & CosTraceNormFiveHundredTwentyOne & PROVED & \checkmark & verified \\
\texttt{isIntegral\_and\_degree\_fiveHundredTwentyOne} & CosTraceNormFiveHundredTwentyOne & PROVED & \checkmark & verified \\
\texttt{isIntegral\_spectralGen\_fiveHundredTwentyOne} & CosTraceNormFiveHundredTwentyOne & PROVED & \checkmark & verified \\
\texttt{isIntegral\_spectralGen\_fiveHundredTwentyOne\_Q} & CosTraceNormFiveHundredTwentyOne & PROVED & \checkmark & verified \\
\texttt{prime\_fiveHundredTwentyOne} & CosTraceNormFiveHundredTwentyOne & PROVED & \checkmark & verified \\
\texttt{degree\_fiveHundredTwentyThree} & CosTraceNormFiveHundredTwentyThree & PROVED & \checkmark & verified \\
\texttt{fiveHundredTwentyThree\_ne\_two} & CosTraceNormFiveHundredTwentyThree & PROVED & \checkmark & verified \\
\texttt{fiveHundredTwentyThree\_pack} & CosTraceNormFiveHundredTwentyThree & PROVED & \checkmark & verified \\
\texttt{isIntegral\_and\_degree\_fiveHundredTwentyThree} & CosTraceNormFiveHundredTwentyThree & PROVED & \checkmark & verified \\
\texttt{isIntegral\_spectralGen\_fiveHundredTwentyThree} & CosTraceNormFiveHundredTwentyThree & PROVED & \checkmark & verified \\
\texttt{isIntegral\_spectralGen\_fiveHundredTwentyThree\_Q} & CosTraceNormFiveHundredTwentyThree & PROVED & \checkmark & verified \\
\texttt{prime\_fiveHundredTwentyThree} & CosTraceNormFiveHundredTwentyThree & PROVED & \checkmark & verified \\
\texttt{degree\_fortyOne} & CosTraceNormFortyOne & PROVED & \checkmark & verified \\
\texttt{fortyOne\_ne\_two} & CosTraceNormFortyOne & PROVED & \checkmark & verified \\
\texttt{fortyOne\_pack} & CosTraceNormFortyOne & PROVED & \checkmark & verified \\
\texttt{isIntegral\_and\_degree\_fortyOne} & CosTraceNormFortyOne & PROVED & \checkmark & verified \\
\texttt{isIntegral\_spectralGen\_fortyOne} & CosTraceNormFortyOne & PROVED & \checkmark & verified \\
\texttt{isIntegral\_spectralGen\_fortyOne\_Q} & CosTraceNormFortyOne & PROVED & \checkmark & verified \\
\texttt{prime\_fortyOne} & CosTraceNormFortyOne & PROVED & \checkmark & verified \\
\texttt{degree\_fortySeven} & CosTraceNormFortySeven & PROVED & \checkmark & verified \\
\texttt{fortySeven\_ne\_two} & CosTraceNormFortySeven & PROVED & \checkmark & verified \\
\texttt{fortySeven\_pack} & CosTraceNormFortySeven & PROVED & \checkmark & verified \\
\texttt{isIntegral\_and\_degree\_fortySeven} & CosTraceNormFortySeven & PROVED & \checkmark & verified \\
\texttt{isIntegral\_spectralGen\_fortySeven} & CosTraceNormFortySeven & PROVED & \checkmark & verified \\
\texttt{isIntegral\_spectralGen\_fortySeven\_Q} & CosTraceNormFortySeven & PROVED & \checkmark & verified \\
\texttt{prime\_fortySeven} & CosTraceNormFortySeven & PROVED & \checkmark & verified \\
\texttt{degree\_fortyThree} & CosTraceNormFortyThree & PROVED & \checkmark & verified \\
\texttt{fortyThree\_ne\_two} & CosTraceNormFortyThree & PROVED & \checkmark & verified \\
\texttt{fortyThree\_pack} & CosTraceNormFortyThree & PROVED & \checkmark & verified \\
\texttt{isIntegral\_and\_degree\_fortyThree} & CosTraceNormFortyThree & PROVED & \checkmark & verified \\
\texttt{isIntegral\_spectralGen\_fortyThree} & CosTraceNormFortyThree & PROVED & \checkmark & verified \\
\texttt{isIntegral\_spectralGen\_fortyThree\_Q} & CosTraceNormFortyThree & PROVED & \checkmark & verified \\
\texttt{prime\_fortyThree} & CosTraceNormFortyThree & PROVED & \checkmark & verified \\
\texttt{degree\_fourHundredEightySeven} & CosTraceNormFourHundredEightySeven & PROVED & \checkmark & verified \\
\texttt{fourHundredEightySeven\_ne\_two} & CosTraceNormFourHundredEightySeven & PROVED & \checkmark & verified \\
\texttt{fourHundredEightySeven\_pack} & CosTraceNormFourHundredEightySeven & PROVED & \checkmark & verified \\
\texttt{isIntegral\_and\_degree\_fourHundredEightySeven} & CosTraceNormFourHundredEightySeven & PROVED & \checkmark & verified \\
\texttt{isIntegral\_spectralGen\_fourHundredEightySeven} & CosTraceNormFourHundredEightySeven & PROVED & \checkmark & verified \\
\texttt{isIntegral\_spectralGen\_fourHundredEightySeven\_Q} & CosTraceNormFourHundredEightySeven & PROVED & \checkmark & verified \\
\texttt{prime\_fourHundredEightySeven} & CosTraceNormFourHundredEightySeven & PROVED & \checkmark & verified \\
\texttt{degree\_fourHundredFiftySeven} & CosTraceNormFourHundredFiftySeven & PROVED & \checkmark & verified \\
\texttt{fourHundredFiftySeven\_ne\_two} & CosTraceNormFourHundredFiftySeven & PROVED & \checkmark & verified \\
\texttt{fourHundredFiftySeven\_pack} & CosTraceNormFourHundredFiftySeven & PROVED & \checkmark & verified \\
\texttt{isIntegral\_and\_degree\_fourHundredFiftySeven} & CosTraceNormFourHundredFiftySeven & PROVED & \checkmark & verified \\
\texttt{isIntegral\_spectralGen\_fourHundredFiftySeven} & CosTraceNormFourHundredFiftySeven & PROVED & \checkmark & verified \\
\texttt{isIntegral\_spectralGen\_fourHundredFiftySeven\_Q} & CosTraceNormFourHundredFiftySeven & PROVED & \checkmark & verified \\
\texttt{prime\_fourHundredFiftySeven} & CosTraceNormFourHundredFiftySeven & PROVED & \checkmark & verified \\
\texttt{degree\_fourHundredFortyNine} & CosTraceNormFourHundredFortyNine & PROVED & \checkmark & verified \\
\texttt{fourHundredFortyNine\_ne\_two} & CosTraceNormFourHundredFortyNine & PROVED & \checkmark & verified \\
\texttt{fourHundredFortyNine\_pack} & CosTraceNormFourHundredFortyNine & PROVED & \checkmark & verified \\
\texttt{isIntegral\_and\_degree\_fourHundredFortyNine} & CosTraceNormFourHundredFortyNine & PROVED & \checkmark & verified \\
\texttt{isIntegral\_spectralGen\_fourHundredFortyNine} & CosTraceNormFourHundredFortyNine & PROVED & \checkmark & verified \\
\texttt{isIntegral\_spectralGen\_fourHundredFortyNine\_Q} & CosTraceNormFourHundredFortyNine & PROVED & \checkmark & verified \\
\texttt{prime\_fourHundredFortyNine} & CosTraceNormFourHundredFortyNine & PROVED & \checkmark & verified \\
\texttt{degree\_fourHundredFortyThree} & CosTraceNormFourHundredFortyThree & PROVED & \checkmark & verified \\
\texttt{fourHundredFortyThree\_ne\_two} & CosTraceNormFourHundredFortyThree & PROVED & \checkmark & verified \\
\texttt{fourHundredFortyThree\_pack} & CosTraceNormFourHundredFortyThree & PROVED & \checkmark & verified \\
\texttt{isIntegral\_and\_degree\_fourHundredFortyThree} & CosTraceNormFourHundredFortyThree & PROVED & \checkmark & verified \\
\texttt{isIntegral\_spectralGen\_fourHundredFortyThree} & CosTraceNormFourHundredFortyThree & PROVED & \checkmark & verified \\
\texttt{isIntegral\_spectralGen\_fourHundredFortyThree\_Q} & CosTraceNormFourHundredFortyThree & PROVED & \checkmark & verified \\
\texttt{prime\_fourHundredFortyThree} & CosTraceNormFourHundredFortyThree & PROVED & \checkmark & verified \\
\texttt{degree\_fourHundredNine} & CosTraceNormFourHundredNine & PROVED & \checkmark & verified \\
\texttt{fourHundredNine\_ne\_two} & CosTraceNormFourHundredNine & PROVED & \checkmark & verified \\
\texttt{fourHundredNine\_pack} & CosTraceNormFourHundredNine & PROVED & \checkmark & verified \\
\texttt{isIntegral\_and\_degree\_fourHundredNine} & CosTraceNormFourHundredNine & PROVED & \checkmark & verified \\
\texttt{isIntegral\_spectralGen\_fourHundredNine} & CosTraceNormFourHundredNine & PROVED & \checkmark & verified \\
\texttt{isIntegral\_spectralGen\_fourHundredNine\_Q} & CosTraceNormFourHundredNine & PROVED & \checkmark & verified \\
\texttt{prime\_fourHundredNine} & CosTraceNormFourHundredNine & PROVED & \checkmark & verified \\
\texttt{degree\_fourHundredNineteen} & CosTraceNormFourHundredNineteen & PROVED & \checkmark & verified \\
\texttt{fourHundredNineteen\_ne\_two} & CosTraceNormFourHundredNineteen & PROVED & \checkmark & verified \\
\texttt{fourHundredNineteen\_pack} & CosTraceNormFourHundredNineteen & PROVED & \checkmark & verified \\
\texttt{isIntegral\_and\_degree\_fourHundredNineteen} & CosTraceNormFourHundredNineteen & PROVED & \checkmark & verified \\
\texttt{isIntegral\_spectralGen\_fourHundredNineteen} & CosTraceNormFourHundredNineteen & PROVED & \checkmark & verified \\
\texttt{isIntegral\_spectralGen\_fourHundredNineteen\_Q} & CosTraceNormFourHundredNineteen & PROVED & \checkmark & verified \\
\texttt{prime\_fourHundredNineteen} & CosTraceNormFourHundredNineteen & PROVED & \checkmark & verified \\
\texttt{degree\_fourHundredNinetyNine} & CosTraceNormFourHundredNinetyNine & PROVED & \checkmark & verified \\
\texttt{fourHundredNinetyNine\_ne\_two} & CosTraceNormFourHundredNinetyNine & PROVED & \checkmark & verified \\
\texttt{fourHundredNinetyNine\_pack} & CosTraceNormFourHundredNinetyNine & PROVED & \checkmark & verified \\
\texttt{isIntegral\_and\_degree\_fourHundredNinetyNine} & CosTraceNormFourHundredNinetyNine & PROVED & \checkmark & verified \\
\texttt{isIntegral\_spectralGen\_fourHundredNinetyNine} & CosTraceNormFourHundredNinetyNine & PROVED & \checkmark & verified \\
\texttt{isIntegral\_spectralGen\_fourHundredNinetyNine\_Q} & CosTraceNormFourHundredNinetyNine & PROVED & \checkmark & verified \\
\texttt{prime\_fourHundredNinetyNine} & CosTraceNormFourHundredNinetyNine & PROVED & \checkmark & verified \\
\texttt{degree\_fourHundredNinetyOne} & CosTraceNormFourHundredNinetyOne & PROVED & \checkmark & verified \\
\texttt{fourHundredNinetyOne\_ne\_two} & CosTraceNormFourHundredNinetyOne & PROVED & \checkmark & verified \\
\texttt{fourHundredNinetyOne\_pack} & CosTraceNormFourHundredNinetyOne & PROVED & \checkmark & verified \\
\texttt{isIntegral\_and\_degree\_fourHundredNinetyOne} & CosTraceNormFourHundredNinetyOne & PROVED & \checkmark & verified \\
\texttt{isIntegral\_spectralGen\_fourHundredNinetyOne} & CosTraceNormFourHundredNinetyOne & PROVED & \checkmark & verified \\
\texttt{isIntegral\_spectralGen\_fourHundredNinetyOne\_Q} & CosTraceNormFourHundredNinetyOne & PROVED & \checkmark & verified \\
\texttt{prime\_fourHundredNinetyOne} & CosTraceNormFourHundredNinetyOne & PROVED & \checkmark & verified \\
\texttt{degree\_fourHundredOne} & CosTraceNormFourHundredOne & PROVED & \checkmark & verified \\
\texttt{fourHundredOne\_ne\_two} & CosTraceNormFourHundredOne & PROVED & \checkmark & verified \\
\texttt{fourHundredOne\_pack} & CosTraceNormFourHundredOne & PROVED & \checkmark & verified \\
\texttt{isIntegral\_and\_degree\_fourHundredOne} & CosTraceNormFourHundredOne & PROVED & \checkmark & verified \\
\texttt{isIntegral\_spectralGen\_fourHundredOne} & CosTraceNormFourHundredOne & PROVED & \checkmark & verified \\
\texttt{isIntegral\_spectralGen\_fourHundredOne\_Q} & CosTraceNormFourHundredOne & PROVED & \checkmark & verified \\
\texttt{prime\_fourHundredOne} & CosTraceNormFourHundredOne & PROVED & \checkmark & verified \\
\texttt{degree\_fourHundredSeventyNine} & CosTraceNormFourHundredSeventyNine & PROVED & \checkmark & verified \\
\texttt{fourHundredSeventyNine\_ne\_two} & CosTraceNormFourHundredSeventyNine & PROVED & \checkmark & verified \\
\texttt{fourHundredSeventyNine\_pack} & CosTraceNormFourHundredSeventyNine & PROVED & \checkmark & verified \\
\texttt{isIntegral\_and\_degree\_fourHundredSeventyNine} & CosTraceNormFourHundredSeventyNine & PROVED & \checkmark & verified \\
\texttt{isIntegral\_spectralGen\_fourHundredSeventyNine} & CosTraceNormFourHundredSeventyNine & PROVED & \checkmark & verified \\
\texttt{isIntegral\_spectralGen\_fourHundredSeventyNine\_Q} & CosTraceNormFourHundredSeventyNine & PROVED & \checkmark & verified \\
\texttt{prime\_fourHundredSeventyNine} & CosTraceNormFourHundredSeventyNine & PROVED & \checkmark & verified \\
\texttt{degree\_fourHundredSixtyOne} & CosTraceNormFourHundredSixtyOne & PROVED & \checkmark & verified \\
\texttt{fourHundredSixtyOne\_ne\_two} & CosTraceNormFourHundredSixtyOne & PROVED & \checkmark & verified \\
\texttt{fourHundredSixtyOne\_pack} & CosTraceNormFourHundredSixtyOne & PROVED & \checkmark & verified \\
\texttt{isIntegral\_and\_degree\_fourHundredSixtyOne} & CosTraceNormFourHundredSixtyOne & PROVED & \checkmark & verified \\
\texttt{isIntegral\_spectralGen\_fourHundredSixtyOne} & CosTraceNormFourHundredSixtyOne & PROVED & \checkmark & verified \\
\texttt{isIntegral\_spectralGen\_fourHundredSixtyOne\_Q} & CosTraceNormFourHundredSixtyOne & PROVED & \checkmark & verified \\
\texttt{prime\_fourHundredSixtyOne} & CosTraceNormFourHundredSixtyOne & PROVED & \checkmark & verified \\
\texttt{degree\_fourHundredSixtySeven} & CosTraceNormFourHundredSixtySeven & PROVED & \checkmark & verified \\
\texttt{fourHundredSixtySeven\_ne\_two} & CosTraceNormFourHundredSixtySeven & PROVED & \checkmark & verified \\
\texttt{fourHundredSixtySeven\_pack} & CosTraceNormFourHundredSixtySeven & PROVED & \checkmark & verified \\
\texttt{isIntegral\_and\_degree\_fourHundredSixtySeven} & CosTraceNormFourHundredSixtySeven & PROVED & \checkmark & verified \\
\texttt{isIntegral\_spectralGen\_fourHundredSixtySeven} & CosTraceNormFourHundredSixtySeven & PROVED & \checkmark & verified \\
\texttt{isIntegral\_spectralGen\_fourHundredSixtySeven\_Q} & CosTraceNormFourHundredSixtySeven & PROVED & \checkmark & verified \\
\texttt{prime\_fourHundredSixtySeven} & CosTraceNormFourHundredSixtySeven & PROVED & \checkmark & verified \\
\texttt{degree\_fourHundredSixtyThree} & CosTraceNormFourHundredSixtyThree & PROVED & \checkmark & verified \\
\texttt{fourHundredSixtyThree\_ne\_two} & CosTraceNormFourHundredSixtyThree & PROVED & \checkmark & verified \\
\texttt{fourHundredSixtyThree\_pack} & CosTraceNormFourHundredSixtyThree & PROVED & \checkmark & verified \\
\texttt{isIntegral\_and\_degree\_fourHundredSixtyThree} & CosTraceNormFourHundredSixtyThree & PROVED & \checkmark & verified \\
\texttt{isIntegral\_spectralGen\_fourHundredSixtyThree} & CosTraceNormFourHundredSixtyThree & PROVED & \checkmark & verified \\
\texttt{isIntegral\_spectralGen\_fourHundredSixtyThree\_Q} & CosTraceNormFourHundredSixtyThree & PROVED & \checkmark & verified \\
\texttt{prime\_fourHundredSixtyThree} & CosTraceNormFourHundredSixtyThree & PROVED & \checkmark & verified \\
\texttt{degree\_fourHundredThirtyNine} & CosTraceNormFourHundredThirtyNine & PROVED & \checkmark & verified \\
\texttt{fourHundredThirtyNine\_ne\_two} & CosTraceNormFourHundredThirtyNine & PROVED & \checkmark & verified \\
\texttt{fourHundredThirtyNine\_pack} & CosTraceNormFourHundredThirtyNine & PROVED & \checkmark & verified \\
\texttt{isIntegral\_and\_degree\_fourHundredThirtyNine} & CosTraceNormFourHundredThirtyNine & PROVED & \checkmark & verified \\
\texttt{isIntegral\_spectralGen\_fourHundredThirtyNine} & CosTraceNormFourHundredThirtyNine & PROVED & \checkmark & verified \\
\texttt{isIntegral\_spectralGen\_fourHundredThirtyNine\_Q} & CosTraceNormFourHundredThirtyNine & PROVED & \checkmark & verified \\
\texttt{prime\_fourHundredThirtyNine} & CosTraceNormFourHundredThirtyNine & PROVED & \checkmark & verified \\
\texttt{degree\_fourHundredThirtyOne} & CosTraceNormFourHundredThirtyOne & PROVED & \checkmark & verified \\
\texttt{fourHundredThirtyOne\_ne\_two} & CosTraceNormFourHundredThirtyOne & PROVED & \checkmark & verified \\
\texttt{fourHundredThirtyOne\_pack} & CosTraceNormFourHundredThirtyOne & PROVED & \checkmark & verified \\
\texttt{isIntegral\_and\_degree\_fourHundredThirtyOne} & CosTraceNormFourHundredThirtyOne & PROVED & \checkmark & verified \\
\texttt{isIntegral\_spectralGen\_fourHundredThirtyOne} & CosTraceNormFourHundredThirtyOne & PROVED & \checkmark & verified \\
\texttt{isIntegral\_spectralGen\_fourHundredThirtyOne\_Q} & CosTraceNormFourHundredThirtyOne & PROVED & \checkmark & verified \\
\texttt{prime\_fourHundredThirtyOne} & CosTraceNormFourHundredThirtyOne & PROVED & \checkmark & verified \\
\texttt{degree\_fourHundredThirtyThree} & CosTraceNormFourHundredThirtyThree & PROVED & \checkmark & verified \\
\texttt{fourHundredThirtyThree\_ne\_two} & CosTraceNormFourHundredThirtyThree & PROVED & \checkmark & verified \\
\texttt{fourHundredThirtyThree\_pack} & CosTraceNormFourHundredThirtyThree & PROVED & \checkmark & verified \\
\texttt{isIntegral\_and\_degree\_fourHundredThirtyThree} & CosTraceNormFourHundredThirtyThree & PROVED & \checkmark & verified \\
\texttt{isIntegral\_spectralGen\_fourHundredThirtyThree} & CosTraceNormFourHundredThirtyThree & PROVED & \checkmark & verified \\
\texttt{isIntegral\_spectralGen\_fourHundredThirtyThree\_Q} & CosTraceNormFourHundredThirtyThree & PROVED & \checkmark & verified \\
\texttt{prime\_fourHundredThirtyThree} & CosTraceNormFourHundredThirtyThree & PROVED & \checkmark & verified \\
\texttt{degree\_fourHundredTwentyOne} & CosTraceNormFourHundredTwentyOne & PROVED & \checkmark & verified \\
\texttt{fourHundredTwentyOne\_ne\_two} & CosTraceNormFourHundredTwentyOne & PROVED & \checkmark & verified \\
\texttt{fourHundredTwentyOne\_pack} & CosTraceNormFourHundredTwentyOne & PROVED & \checkmark & verified \\
\texttt{isIntegral\_and\_degree\_fourHundredTwentyOne} & CosTraceNormFourHundredTwentyOne & PROVED & \checkmark & verified \\
\texttt{isIntegral\_spectralGen\_fourHundredTwentyOne} & CosTraceNormFourHundredTwentyOne & PROVED & \checkmark & verified \\
\texttt{isIntegral\_spectralGen\_fourHundredTwentyOne\_Q} & CosTraceNormFourHundredTwentyOne & PROVED & \checkmark & verified \\
\texttt{prime\_fourHundredTwentyOne} & CosTraceNormFourHundredTwentyOne & PROVED & \checkmark & verified \\
\texttt{degree\_nineHundredEightyThree} & CosTraceNormNineHundredEightyThree & PROVED & \checkmark & verified \\
\texttt{isIntegral\_and\_degree\_nineHundredEightyThree} & CosTraceNormNineHundredEightyThree & PROVED & \checkmark & verified \\
\texttt{isIntegral\_spectralGen\_nineHundredEightyThree} & CosTraceNormNineHundredEightyThree & PROVED & \checkmark & verified \\
\texttt{isIntegral\_spectralGen\_nineHundredEightyThree\_Q} & CosTraceNormNineHundredEightyThree & PROVED & \checkmark & verified \\
\texttt{nineHundredEightyThree\_ne\_two} & CosTraceNormNineHundredEightyThree & PROVED & \checkmark & verified \\
\texttt{nineHundredEightyThree\_pack} & CosTraceNormNineHundredEightyThree & PROVED & \checkmark & verified \\
\texttt{prime\_nineHundredEightyThree} & CosTraceNormNineHundredEightyThree & PROVED & \checkmark & verified \\
\texttt{degree\_nineHundredEleven} & CosTraceNormNineHundredEleven & PROVED & \checkmark & verified \\
\texttt{isIntegral\_and\_degree\_nineHundredEleven} & CosTraceNormNineHundredEleven & PROVED & \checkmark & verified \\
\texttt{isIntegral\_spectralGen\_nineHundredEleven} & CosTraceNormNineHundredEleven & PROVED & \checkmark & verified \\
\texttt{isIntegral\_spectralGen\_nineHundredEleven\_Q} & CosTraceNormNineHundredEleven & PROVED & \checkmark & verified \\
\texttt{nineHundredEleven\_ne\_two} & CosTraceNormNineHundredEleven & PROVED & \checkmark & verified \\
\texttt{nineHundredEleven\_pack} & CosTraceNormNineHundredEleven & PROVED & \checkmark & verified \\
\texttt{prime\_nineHundredEleven} & CosTraceNormNineHundredEleven & PROVED & \checkmark & verified \\
\texttt{degree\_nineHundredFiftyThree} & CosTraceNormNineHundredFiftyThree & PROVED & \checkmark & verified \\
\texttt{isIntegral\_and\_degree\_nineHundredFiftyThree} & CosTraceNormNineHundredFiftyThree & PROVED & \checkmark & verified \\
\texttt{isIntegral\_spectralGen\_nineHundredFiftyThree} & CosTraceNormNineHundredFiftyThree & PROVED & \checkmark & verified \\
\texttt{isIntegral\_spectralGen\_nineHundredFiftyThree\_Q} & CosTraceNormNineHundredFiftyThree & PROVED & \checkmark & verified \\
\texttt{nineHundredFiftyThree\_ne\_two} & CosTraceNormNineHundredFiftyThree & PROVED & \checkmark & verified \\
\texttt{nineHundredFiftyThree\_pack} & CosTraceNormNineHundredFiftyThree & PROVED & \checkmark & verified \\
\texttt{prime\_nineHundredFiftyThree} & CosTraceNormNineHundredFiftyThree & PROVED & \checkmark & verified \\
\texttt{degree\_nineHundredFortyOne} & CosTraceNormNineHundredFortyOne & PROVED & \checkmark & verified \\
\texttt{isIntegral\_and\_degree\_nineHundredFortyOne} & CosTraceNormNineHundredFortyOne & PROVED & \checkmark & verified \\
\texttt{isIntegral\_spectralGen\_nineHundredFortyOne} & CosTraceNormNineHundredFortyOne & PROVED & \checkmark & verified \\
\texttt{isIntegral\_spectralGen\_nineHundredFortyOne\_Q} & CosTraceNormNineHundredFortyOne & PROVED & \checkmark & verified \\
\texttt{nineHundredFortyOne\_ne\_two} & CosTraceNormNineHundredFortyOne & PROVED & \checkmark & verified \\
\texttt{nineHundredFortyOne\_pack} & CosTraceNormNineHundredFortyOne & PROVED & \checkmark & verified \\
\texttt{prime\_nineHundredFortyOne} & CosTraceNormNineHundredFortyOne & PROVED & \checkmark & verified \\
\texttt{degree\_nineHundredFortySeven} & CosTraceNormNineHundredFortySeven & PROVED & \checkmark & verified \\
\texttt{isIntegral\_and\_degree\_nineHundredFortySeven} & CosTraceNormNineHundredFortySeven & PROVED & \checkmark & verified \\
\texttt{isIntegral\_spectralGen\_nineHundredFortySeven} & CosTraceNormNineHundredFortySeven & PROVED & \checkmark & verified \\
\texttt{isIntegral\_spectralGen\_nineHundredFortySeven\_Q} & CosTraceNormNineHundredFortySeven & PROVED & \checkmark & verified \\
\texttt{nineHundredFortySeven\_ne\_two} & CosTraceNormNineHundredFortySeven & PROVED & \checkmark & verified \\
\texttt{nineHundredFortySeven\_pack} & CosTraceNormNineHundredFortySeven & PROVED & \checkmark & verified \\
\texttt{prime\_nineHundredFortySeven} & CosTraceNormNineHundredFortySeven & PROVED & \checkmark & verified \\
\texttt{degree\_nineHundredNineteen} & CosTraceNormNineHundredNineteen & PROVED & \checkmark & verified \\
\texttt{isIntegral\_and\_degree\_nineHundredNineteen} & CosTraceNormNineHundredNineteen & PROVED & \checkmark & verified \\
\texttt{isIntegral\_spectralGen\_nineHundredNineteen} & CosTraceNormNineHundredNineteen & PROVED & \checkmark & verified \\
\texttt{isIntegral\_spectralGen\_nineHundredNineteen\_Q} & CosTraceNormNineHundredNineteen & PROVED & \checkmark & verified \\
\texttt{nineHundredNineteen\_ne\_two} & CosTraceNormNineHundredNineteen & PROVED & \checkmark & verified \\
\texttt{nineHundredNineteen\_pack} & CosTraceNormNineHundredNineteen & PROVED & \checkmark & verified \\
\texttt{prime\_nineHundredNineteen} & CosTraceNormNineHundredNineteen & PROVED & \checkmark & verified \\
\texttt{degree\_nineHundredNinetyOne} & CosTraceNormNineHundredNinetyOne & PROVED & \checkmark & verified \\
\texttt{isIntegral\_and\_degree\_nineHundredNinetyOne} & CosTraceNormNineHundredNinetyOne & PROVED & \checkmark & verified \\
\texttt{isIntegral\_spectralGen\_nineHundredNinetyOne} & CosTraceNormNineHundredNinetyOne & PROVED & \checkmark & verified \\
\texttt{isIntegral\_spectralGen\_nineHundredNinetyOne\_Q} & CosTraceNormNineHundredNinetyOne & PROVED & \checkmark & verified \\
\texttt{nineHundredNinetyOne\_ne\_two} & CosTraceNormNineHundredNinetyOne & PROVED & \checkmark & verified \\
\texttt{nineHundredNinetyOne\_pack} & CosTraceNormNineHundredNinetyOne & PROVED & \checkmark & verified \\
\texttt{prime\_nineHundredNinetyOne} & CosTraceNormNineHundredNinetyOne & PROVED & \checkmark & verified \\
\texttt{degree\_nineHundredNinetySeven} & CosTraceNormNineHundredNinetySeven & PROVED & \checkmark & verified \\
\texttt{isIntegral\_and\_degree\_nineHundredNinetySeven} & CosTraceNormNineHundredNinetySeven & PROVED & \checkmark & verified \\
\texttt{isIntegral\_spectralGen\_nineHundredNinetySeven} & CosTraceNormNineHundredNinetySeven & PROVED & \checkmark & verified \\
\texttt{isIntegral\_spectralGen\_nineHundredNinetySeven\_Q} & CosTraceNormNineHundredNinetySeven & PROVED & \checkmark & verified \\
\texttt{nineHundredNinetySeven\_ne\_two} & CosTraceNormNineHundredNinetySeven & PROVED & \checkmark & verified \\
\texttt{nineHundredNinetySeven\_pack} & CosTraceNormNineHundredNinetySeven & PROVED & \checkmark & verified \\
\texttt{prime\_nineHundredNinetySeven} & CosTraceNormNineHundredNinetySeven & PROVED & \checkmark & verified \\
\texttt{degree\_nineHundredSeven} & CosTraceNormNineHundredSeven & PROVED & \checkmark & verified \\
\texttt{isIntegral\_and\_degree\_nineHundredSeven} & CosTraceNormNineHundredSeven & PROVED & \checkmark & verified \\
\texttt{isIntegral\_spectralGen\_nineHundredSeven} & CosTraceNormNineHundredSeven & PROVED & \checkmark & verified \\
\texttt{isIntegral\_spectralGen\_nineHundredSeven\_Q} & CosTraceNormNineHundredSeven & PROVED & \checkmark & verified \\
\texttt{nineHundredSeven\_ne\_two} & CosTraceNormNineHundredSeven & PROVED & \checkmark & verified \\
\texttt{nineHundredSeven\_pack} & CosTraceNormNineHundredSeven & PROVED & \checkmark & verified \\
\texttt{prime\_nineHundredSeven} & CosTraceNormNineHundredSeven & PROVED & \checkmark & verified \\
\texttt{degree\_nineHundredSeventyOne} & CosTraceNormNineHundredSeventyOne & PROVED & \checkmark & verified \\
\texttt{isIntegral\_and\_degree\_nineHundredSeventyOne} & CosTraceNormNineHundredSeventyOne & PROVED & \checkmark & verified \\
\texttt{isIntegral\_spectralGen\_nineHundredSeventyOne} & CosTraceNormNineHundredSeventyOne & PROVED & \checkmark & verified \\
\texttt{isIntegral\_spectralGen\_nineHundredSeventyOne\_Q} & CosTraceNormNineHundredSeventyOne & PROVED & \checkmark & verified \\
\texttt{nineHundredSeventyOne\_ne\_two} & CosTraceNormNineHundredSeventyOne & PROVED & \checkmark & verified \\
\texttt{nineHundredSeventyOne\_pack} & CosTraceNormNineHundredSeventyOne & PROVED & \checkmark & verified \\
\texttt{prime\_nineHundredSeventyOne} & CosTraceNormNineHundredSeventyOne & PROVED & \checkmark & verified \\
\texttt{degree\_nineHundredSeventySeven} & CosTraceNormNineHundredSeventySeven & PROVED & \checkmark & verified \\
\texttt{isIntegral\_and\_degree\_nineHundredSeventySeven} & CosTraceNormNineHundredSeventySeven & PROVED & \checkmark & verified \\
\texttt{isIntegral\_spectralGen\_nineHundredSeventySeven} & CosTraceNormNineHundredSeventySeven & PROVED & \checkmark & verified \\
\texttt{isIntegral\_spectralGen\_nineHundredSeventySeven\_Q} & CosTraceNormNineHundredSeventySeven & PROVED & \checkmark & verified \\
\texttt{nineHundredSeventySeven\_ne\_two} & CosTraceNormNineHundredSeventySeven & PROVED & \checkmark & verified \\
\texttt{nineHundredSeventySeven\_pack} & CosTraceNormNineHundredSeventySeven & PROVED & \checkmark & verified \\
\texttt{prime\_nineHundredSeventySeven} & CosTraceNormNineHundredSeventySeven & PROVED & \checkmark & verified \\
\texttt{degree\_nineHundredSixtySeven} & CosTraceNormNineHundredSixtySeven & PROVED & \checkmark & verified \\
\texttt{isIntegral\_and\_degree\_nineHundredSixtySeven} & CosTraceNormNineHundredSixtySeven & PROVED & \checkmark & verified \\
\texttt{isIntegral\_spectralGen\_nineHundredSixtySeven} & CosTraceNormNineHundredSixtySeven & PROVED & \checkmark & verified \\
\texttt{isIntegral\_spectralGen\_nineHundredSixtySeven\_Q} & CosTraceNormNineHundredSixtySeven & PROVED & \checkmark & verified \\
\texttt{nineHundredSixtySeven\_ne\_two} & CosTraceNormNineHundredSixtySeven & PROVED & \checkmark & verified \\
\texttt{nineHundredSixtySeven\_pack} & CosTraceNormNineHundredSixtySeven & PROVED & \checkmark & verified \\
\texttt{prime\_nineHundredSixtySeven} & CosTraceNormNineHundredSixtySeven & PROVED & \checkmark & verified \\
\texttt{degree\_nineHundredThirtySeven} & CosTraceNormNineHundredThirtySeven & PROVED & \checkmark & verified \\
\texttt{isIntegral\_and\_degree\_nineHundredThirtySeven} & CosTraceNormNineHundredThirtySeven & PROVED & \checkmark & verified \\
\texttt{isIntegral\_spectralGen\_nineHundredThirtySeven} & CosTraceNormNineHundredThirtySeven & PROVED & \checkmark & verified \\
\texttt{isIntegral\_spectralGen\_nineHundredThirtySeven\_Q} & CosTraceNormNineHundredThirtySeven & PROVED & \checkmark & verified \\
\texttt{nineHundredThirtySeven\_ne\_two} & CosTraceNormNineHundredThirtySeven & PROVED & \checkmark & verified \\
\texttt{nineHundredThirtySeven\_pack} & CosTraceNormNineHundredThirtySeven & PROVED & \checkmark & verified \\
\texttt{prime\_nineHundredThirtySeven} & CosTraceNormNineHundredThirtySeven & PROVED & \checkmark & verified \\
\texttt{degree\_nineHundredTwentyNine} & CosTraceNormNineHundredTwentyNine & PROVED & \checkmark & verified \\
\texttt{isIntegral\_and\_degree\_nineHundredTwentyNine} & CosTraceNormNineHundredTwentyNine & PROVED & \checkmark & verified \\
\texttt{isIntegral\_spectralGen\_nineHundredTwentyNine} & CosTraceNormNineHundredTwentyNine & PROVED & \checkmark & verified \\
\texttt{isIntegral\_spectralGen\_nineHundredTwentyNine\_Q} & CosTraceNormNineHundredTwentyNine & PROVED & \checkmark & verified \\
\texttt{nineHundredTwentyNine\_ne\_two} & CosTraceNormNineHundredTwentyNine & PROVED & \checkmark & verified \\
\texttt{nineHundredTwentyNine\_pack} & CosTraceNormNineHundredTwentyNine & PROVED & \checkmark & verified \\
\texttt{prime\_nineHundredTwentyNine} & CosTraceNormNineHundredTwentyNine & PROVED & \checkmark & verified \\
\texttt{degree\_nineteen} & CosTraceNormNineteen & PROVED & \checkmark & verified \\
\texttt{isIntegral\_and\_degree\_nineteen} & CosTraceNormNineteen & PROVED & \checkmark & verified \\
\texttt{isIntegral\_spectralGen\_nineteen} & CosTraceNormNineteen & PROVED & \checkmark & verified \\
\texttt{isIntegral\_spectralGen\_nineteen\_Q} & CosTraceNormNineteen & PROVED & \checkmark & verified \\
\texttt{nineteen\_ne\_two} & CosTraceNormNineteen & PROVED & \checkmark & verified \\
\texttt{nineteen\_pack} & CosTraceNormNineteen & PROVED & \checkmark & verified \\
\texttt{prime\_nineteen} & CosTraceNormNineteen & PROVED & \checkmark & verified \\
\texttt{degree\_ninetySeven} & CosTraceNormNinetySeven & PROVED & \checkmark & verified \\
\texttt{isIntegral\_and\_degree\_ninetySeven} & CosTraceNormNinetySeven & PROVED & \checkmark & verified \\
\texttt{isIntegral\_spectralGen\_ninetySeven} & CosTraceNormNinetySeven & PROVED & \checkmark & verified \\
\texttt{isIntegral\_spectralGen\_ninetySeven\_Q} & CosTraceNormNinetySeven & PROVED & \checkmark & verified \\
\texttt{ninetySeven\_ne\_two} & CosTraceNormNinetySeven & PROVED & \checkmark & verified \\
\texttt{ninetySeven\_pack} & CosTraceNormNinetySeven & PROVED & \checkmark & verified \\
\texttt{prime\_ninetySeven} & CosTraceNormNinetySeven & PROVED & \checkmark & verified \\
\texttt{degree\_oneHundredEightyOne} & CosTraceNormOneHundredEightyOne & PROVED & \checkmark & verified \\
\texttt{isIntegral\_and\_degree\_oneHundredEightyOne} & CosTraceNormOneHundredEightyOne & PROVED & \checkmark & verified \\
\texttt{isIntegral\_spectralGen\_oneHundredEightyOne} & CosTraceNormOneHundredEightyOne & PROVED & \checkmark & verified \\
\texttt{isIntegral\_spectralGen\_oneHundredEightyOne\_Q} & CosTraceNormOneHundredEightyOne & PROVED & \checkmark & verified \\
\texttt{oneHundredEightyOne\_ne\_two} & CosTraceNormOneHundredEightyOne & PROVED & \checkmark & verified \\
\texttt{oneHundredEightyOne\_pack} & CosTraceNormOneHundredEightyOne & PROVED & \checkmark & verified \\
\texttt{prime\_oneHundredEightyOne} & CosTraceNormOneHundredEightyOne & PROVED & \checkmark & verified \\
\texttt{degree\_oneHundredFiftyOne} & CosTraceNormOneHundredFiftyOne & PROVED & \checkmark & verified \\
\texttt{isIntegral\_and\_degree\_oneHundredFiftyOne} & CosTraceNormOneHundredFiftyOne & PROVED & \checkmark & verified \\
\texttt{isIntegral\_spectralGen\_oneHundredFiftyOne} & CosTraceNormOneHundredFiftyOne & PROVED & \checkmark & verified \\
\texttt{isIntegral\_spectralGen\_oneHundredFiftyOne\_Q} & CosTraceNormOneHundredFiftyOne & PROVED & \checkmark & verified \\
\texttt{oneHundredFiftyOne\_ne\_two} & CosTraceNormOneHundredFiftyOne & PROVED & \checkmark & verified \\
\texttt{oneHundredFiftyOne\_pack} & CosTraceNormOneHundredFiftyOne & PROVED & \checkmark & verified \\
\texttt{prime\_oneHundredFiftyOne} & CosTraceNormOneHundredFiftyOne & PROVED & \checkmark & verified \\
\texttt{degree\_oneHundredFiftySeven} & CosTraceNormOneHundredFiftySeven & PROVED & \checkmark & verified \\
\texttt{isIntegral\_and\_degree\_oneHundredFiftySeven} & CosTraceNormOneHundredFiftySeven & PROVED & \checkmark & verified \\
\texttt{isIntegral\_spectralGen\_oneHundredFiftySeven} & CosTraceNormOneHundredFiftySeven & PROVED & \checkmark & verified \\
\texttt{isIntegral\_spectralGen\_oneHundredFiftySeven\_Q} & CosTraceNormOneHundredFiftySeven & PROVED & \checkmark & verified \\
\texttt{oneHundredFiftySeven\_ne\_two} & CosTraceNormOneHundredFiftySeven & PROVED & \checkmark & verified \\
\texttt{oneHundredFiftySeven\_pack} & CosTraceNormOneHundredFiftySeven & PROVED & \checkmark & verified \\
\texttt{prime\_oneHundredFiftySeven} & CosTraceNormOneHundredFiftySeven & PROVED & \checkmark & verified \\
\texttt{degree\_oneHundredFortyNine} & CosTraceNormOneHundredFortyNine & PROVED & \checkmark & verified \\
\texttt{isIntegral\_and\_degree\_oneHundredFortyNine} & CosTraceNormOneHundredFortyNine & PROVED & \checkmark & verified \\
\texttt{isIntegral\_spectralGen\_oneHundredFortyNine} & CosTraceNormOneHundredFortyNine & PROVED & \checkmark & verified \\
\texttt{isIntegral\_spectralGen\_oneHundredFortyNine\_Q} & CosTraceNormOneHundredFortyNine & PROVED & \checkmark & verified \\
\texttt{oneHundredFortyNine\_ne\_two} & CosTraceNormOneHundredFortyNine & PROVED & \checkmark & verified \\
\texttt{oneHundredFortyNine\_pack} & CosTraceNormOneHundredFortyNine & PROVED & \checkmark & verified \\
\texttt{prime\_oneHundredFortyNine} & CosTraceNormOneHundredFortyNine & PROVED & \checkmark & verified \\
\texttt{degree\_oneHundredNine} & CosTraceNormOneHundredNine & PROVED & \checkmark & verified \\
\texttt{isIntegral\_and\_degree\_oneHundredNine} & CosTraceNormOneHundredNine & PROVED & \checkmark & verified \\
\texttt{isIntegral\_spectralGen\_oneHundredNine} & CosTraceNormOneHundredNine & PROVED & \checkmark & verified \\
\texttt{isIntegral\_spectralGen\_oneHundredNine\_Q} & CosTraceNormOneHundredNine & PROVED & \checkmark & verified \\
\texttt{oneHundredNine\_ne\_two} & CosTraceNormOneHundredNine & PROVED & \checkmark & verified \\
\texttt{oneHundredNine\_pack} & CosTraceNormOneHundredNine & PROVED & \checkmark & verified \\
\texttt{prime\_oneHundredNine} & CosTraceNormOneHundredNine & PROVED & \checkmark & verified \\
\texttt{degree\_oneHundredNinetyNine} & CosTraceNormOneHundredNinetyNine & PROVED & \checkmark & verified \\
\texttt{isIntegral\_and\_degree\_oneHundredNinetyNine} & CosTraceNormOneHundredNinetyNine & PROVED & \checkmark & verified \\
\texttt{isIntegral\_spectralGen\_oneHundredNinetyNine} & CosTraceNormOneHundredNinetyNine & PROVED & \checkmark & verified \\
\texttt{isIntegral\_spectralGen\_oneHundredNinetyNine\_Q} & CosTraceNormOneHundredNinetyNine & PROVED & \checkmark & verified \\
\texttt{oneHundredNinetyNine\_ne\_two} & CosTraceNormOneHundredNinetyNine & PROVED & \checkmark & verified \\
\texttt{oneHundredNinetyNine\_pack} & CosTraceNormOneHundredNinetyNine & PROVED & \checkmark & verified \\
\texttt{prime\_oneHundredNinetyNine} & CosTraceNormOneHundredNinetyNine & PROVED & \checkmark & verified \\
\texttt{degree\_oneHundredNinetyOne} & CosTraceNormOneHundredNinetyOne & PROVED & \checkmark & verified \\
\texttt{isIntegral\_and\_degree\_oneHundredNinetyOne} & CosTraceNormOneHundredNinetyOne & PROVED & \checkmark & verified \\
\texttt{isIntegral\_spectralGen\_oneHundredNinetyOne} & CosTraceNormOneHundredNinetyOne & PROVED & \checkmark & verified \\
\texttt{isIntegral\_spectralGen\_oneHundredNinetyOne\_Q} & CosTraceNormOneHundredNinetyOne & PROVED & \checkmark & verified \\
\texttt{oneHundredNinetyOne\_ne\_two} & CosTraceNormOneHundredNinetyOne & PROVED & \checkmark & verified \\
\texttt{oneHundredNinetyOne\_pack} & CosTraceNormOneHundredNinetyOne & PROVED & \checkmark & verified \\
\texttt{prime\_oneHundredNinetyOne} & CosTraceNormOneHundredNinetyOne & PROVED & \checkmark & verified \\
\texttt{degree\_oneHundredNinetySeven} & CosTraceNormOneHundredNinetySeven & PROVED & \checkmark & verified \\
\texttt{isIntegral\_and\_degree\_oneHundredNinetySeven} & CosTraceNormOneHundredNinetySeven & PROVED & \checkmark & verified \\
\texttt{isIntegral\_spectralGen\_oneHundredNinetySeven} & CosTraceNormOneHundredNinetySeven & PROVED & \checkmark & verified \\
\texttt{isIntegral\_spectralGen\_oneHundredNinetySeven\_Q} & CosTraceNormOneHundredNinetySeven & PROVED & \checkmark & verified \\
\texttt{oneHundredNinetySeven\_ne\_two} & CosTraceNormOneHundredNinetySeven & PROVED & \checkmark & verified \\
\texttt{oneHundredNinetySeven\_pack} & CosTraceNormOneHundredNinetySeven & PROVED & \checkmark & verified \\
\texttt{prime\_oneHundredNinetySeven} & CosTraceNormOneHundredNinetySeven & PROVED & \checkmark & verified \\
\texttt{degree\_oneHundredNinetyThree} & CosTraceNormOneHundredNinetyThree & PROVED & \checkmark & verified \\
\texttt{isIntegral\_and\_degree\_oneHundredNinetyThree} & CosTraceNormOneHundredNinetyThree & PROVED & \checkmark & verified \\
\texttt{isIntegral\_spectralGen\_oneHundredNinetyThree} & CosTraceNormOneHundredNinetyThree & PROVED & \checkmark & verified \\
\texttt{isIntegral\_spectralGen\_oneHundredNinetyThree\_Q} & CosTraceNormOneHundredNinetyThree & PROVED & \checkmark & verified \\
\texttt{oneHundredNinetyThree\_ne\_two} & CosTraceNormOneHundredNinetyThree & PROVED & \checkmark & verified \\
\texttt{oneHundredNinetyThree\_pack} & CosTraceNormOneHundredNinetyThree & PROVED & \checkmark & verified \\
\texttt{prime\_oneHundredNinetyThree} & CosTraceNormOneHundredNinetyThree & PROVED & \checkmark & verified \\
\texttt{degree\_oneHundredOne} & CosTraceNormOneHundredOne & PROVED & \checkmark & verified \\
\texttt{isIntegral\_and\_degree\_oneHundredOne} & CosTraceNormOneHundredOne & PROVED & \checkmark & verified \\
\texttt{isIntegral\_spectralGen\_oneHundredOne} & CosTraceNormOneHundredOne & PROVED & \checkmark & verified \\
\texttt{isIntegral\_spectralGen\_oneHundredOne\_Q} & CosTraceNormOneHundredOne & PROVED & \checkmark & verified \\
\texttt{oneHundredOne\_ne\_two} & CosTraceNormOneHundredOne & PROVED & \checkmark & verified \\
\texttt{oneHundredOne\_pack} & CosTraceNormOneHundredOne & PROVED & \checkmark & verified \\
\texttt{prime\_oneHundredOne} & CosTraceNormOneHundredOne & PROVED & \checkmark & verified \\
\texttt{degree\_oneHundredSeven} & CosTraceNormOneHundredSeven & PROVED & \checkmark & verified \\
\texttt{isIntegral\_and\_degree\_oneHundredSeven} & CosTraceNormOneHundredSeven & PROVED & \checkmark & verified \\
\texttt{isIntegral\_spectralGen\_oneHundredSeven} & CosTraceNormOneHundredSeven & PROVED & \checkmark & verified \\
\texttt{isIntegral\_spectralGen\_oneHundredSeven\_Q} & CosTraceNormOneHundredSeven & PROVED & \checkmark & verified \\
\texttt{oneHundredSeven\_ne\_two} & CosTraceNormOneHundredSeven & PROVED & \checkmark & verified \\
\texttt{oneHundredSeven\_pack} & CosTraceNormOneHundredSeven & PROVED & \checkmark & verified \\
\texttt{prime\_oneHundredSeven} & CosTraceNormOneHundredSeven & PROVED & \checkmark & verified \\
\texttt{degree\_oneHundredSeventyNine} & CosTraceNormOneHundredSeventyNine & PROVED & \checkmark & verified \\
\texttt{isIntegral\_and\_degree\_oneHundredSeventyNine} & CosTraceNormOneHundredSeventyNine & PROVED & \checkmark & verified \\
\texttt{isIntegral\_spectralGen\_oneHundredSeventyNine} & CosTraceNormOneHundredSeventyNine & PROVED & \checkmark & verified \\
\texttt{isIntegral\_spectralGen\_oneHundredSeventyNine\_Q} & CosTraceNormOneHundredSeventyNine & PROVED & \checkmark & verified \\
\texttt{oneHundredSeventyNine\_ne\_two} & CosTraceNormOneHundredSeventyNine & PROVED & \checkmark & verified \\
\texttt{oneHundredSeventyNine\_pack} & CosTraceNormOneHundredSeventyNine & PROVED & \checkmark & verified \\
\texttt{prime\_oneHundredSeventyNine} & CosTraceNormOneHundredSeventyNine & PROVED & \checkmark & verified \\
\texttt{degree\_oneHundredSeventyThree} & CosTraceNormOneHundredSeventyThree & PROVED & \checkmark & verified \\
\texttt{isIntegral\_and\_degree\_oneHundredSeventyThree} & CosTraceNormOneHundredSeventyThree & PROVED & \checkmark & verified \\
\texttt{isIntegral\_spectralGen\_oneHundredSeventyThree} & CosTraceNormOneHundredSeventyThree & PROVED & \checkmark & verified \\
\texttt{isIntegral\_spectralGen\_oneHundredSeventyThree\_Q} & CosTraceNormOneHundredSeventyThree & PROVED & \checkmark & verified \\
\texttt{oneHundredSeventyThree\_ne\_two} & CosTraceNormOneHundredSeventyThree & PROVED & \checkmark & verified \\
\texttt{oneHundredSeventyThree\_pack} & CosTraceNormOneHundredSeventyThree & PROVED & \checkmark & verified \\
\texttt{prime\_oneHundredSeventyThree} & CosTraceNormOneHundredSeventyThree & PROVED & \checkmark & verified \\
\texttt{degree\_oneHundredSixtySeven} & CosTraceNormOneHundredSixtySeven & PROVED & \checkmark & verified \\
\texttt{isIntegral\_and\_degree\_oneHundredSixtySeven} & CosTraceNormOneHundredSixtySeven & PROVED & \checkmark & verified \\
\texttt{isIntegral\_spectralGen\_oneHundredSixtySeven} & CosTraceNormOneHundredSixtySeven & PROVED & \checkmark & verified \\
\texttt{isIntegral\_spectralGen\_oneHundredSixtySeven\_Q} & CosTraceNormOneHundredSixtySeven & PROVED & \checkmark & verified \\
\texttt{oneHundredSixtySeven\_ne\_two} & CosTraceNormOneHundredSixtySeven & PROVED & \checkmark & verified \\
\texttt{oneHundredSixtySeven\_pack} & CosTraceNormOneHundredSixtySeven & PROVED & \checkmark & verified \\
\texttt{prime\_oneHundredSixtySeven} & CosTraceNormOneHundredSixtySeven & PROVED & \checkmark & verified \\
\texttt{degree\_oneHundredSixtyThree} & CosTraceNormOneHundredSixtyThree & PROVED & \checkmark & verified \\
\texttt{isIntegral\_and\_degree\_oneHundredSixtyThree} & CosTraceNormOneHundredSixtyThree & PROVED & \checkmark & verified \\
\texttt{isIntegral\_spectralGen\_oneHundredSixtyThree} & CosTraceNormOneHundredSixtyThree & PROVED & \checkmark & verified \\
\texttt{isIntegral\_spectralGen\_oneHundredSixtyThree\_Q} & CosTraceNormOneHundredSixtyThree & PROVED & \checkmark & verified \\
\texttt{oneHundredSixtyThree\_ne\_two} & CosTraceNormOneHundredSixtyThree & PROVED & \checkmark & verified \\
\texttt{oneHundredSixtyThree\_pack} & CosTraceNormOneHundredSixtyThree & PROVED & \checkmark & verified \\
\texttt{prime\_oneHundredSixtyThree} & CosTraceNormOneHundredSixtyThree & PROVED & \checkmark & verified \\
\texttt{degree\_oneHundredThirteen} & CosTraceNormOneHundredThirteen & PROVED & \checkmark & verified \\
\texttt{isIntegral\_and\_degree\_oneHundredThirteen} & CosTraceNormOneHundredThirteen & PROVED & \checkmark & verified \\
\texttt{isIntegral\_spectralGen\_oneHundredThirteen} & CosTraceNormOneHundredThirteen & PROVED & \checkmark & verified \\
\texttt{isIntegral\_spectralGen\_oneHundredThirteen\_Q} & CosTraceNormOneHundredThirteen & PROVED & \checkmark & verified \\
\texttt{oneHundredThirteen\_ne\_two} & CosTraceNormOneHundredThirteen & PROVED & \checkmark & verified \\
\texttt{oneHundredThirteen\_pack} & CosTraceNormOneHundredThirteen & PROVED & \checkmark & verified \\
\texttt{prime\_oneHundredThirteen} & CosTraceNormOneHundredThirteen & PROVED & \checkmark & verified \\
\texttt{degree\_oneHundredThirtyNine} & CosTraceNormOneHundredThirtyNine & PROVED & \checkmark & verified \\
\texttt{isIntegral\_and\_degree\_oneHundredThirtyNine} & CosTraceNormOneHundredThirtyNine & PROVED & \checkmark & verified \\
\texttt{isIntegral\_spectralGen\_oneHundredThirtyNine} & CosTraceNormOneHundredThirtyNine & PROVED & \checkmark & verified \\
\texttt{isIntegral\_spectralGen\_oneHundredThirtyNine\_Q} & CosTraceNormOneHundredThirtyNine & PROVED & \checkmark & verified \\
\texttt{oneHundredThirtyNine\_ne\_two} & CosTraceNormOneHundredThirtyNine & PROVED & \checkmark & verified \\
\texttt{oneHundredThirtyNine\_pack} & CosTraceNormOneHundredThirtyNine & PROVED & \checkmark & verified \\
\texttt{prime\_oneHundredThirtyNine} & CosTraceNormOneHundredThirtyNine & PROVED & \checkmark & verified \\
\texttt{degree\_oneHundredThirtyOne} & CosTraceNormOneHundredThirtyOne & PROVED & \checkmark & verified \\
\texttt{isIntegral\_and\_degree\_oneHundredThirtyOne} & CosTraceNormOneHundredThirtyOne & PROVED & \checkmark & verified \\
\texttt{isIntegral\_spectralGen\_oneHundredThirtyOne} & CosTraceNormOneHundredThirtyOne & PROVED & \checkmark & verified \\
\texttt{isIntegral\_spectralGen\_oneHundredThirtyOne\_Q} & CosTraceNormOneHundredThirtyOne & PROVED & \checkmark & verified \\
\texttt{oneHundredThirtyOne\_ne\_two} & CosTraceNormOneHundredThirtyOne & PROVED & \checkmark & verified \\
\texttt{oneHundredThirtyOne\_pack} & CosTraceNormOneHundredThirtyOne & PROVED & \checkmark & verified \\
\texttt{prime\_oneHundredThirtyOne} & CosTraceNormOneHundredThirtyOne & PROVED & \checkmark & verified \\
\texttt{degree\_oneHundredThirtySeven} & CosTraceNormOneHundredThirtySeven & PROVED & \checkmark & verified \\
\texttt{isIntegral\_and\_degree\_oneHundredThirtySeven} & CosTraceNormOneHundredThirtySeven & PROVED & \checkmark & verified \\
\texttt{isIntegral\_spectralGen\_oneHundredThirtySeven} & CosTraceNormOneHundredThirtySeven & PROVED & \checkmark & verified \\
\texttt{isIntegral\_spectralGen\_oneHundredThirtySeven\_Q} & CosTraceNormOneHundredThirtySeven & PROVED & \checkmark & verified \\
\texttt{oneHundredThirtySeven\_ne\_two} & CosTraceNormOneHundredThirtySeven & PROVED & \checkmark & verified \\
\texttt{oneHundredThirtySeven\_pack} & CosTraceNormOneHundredThirtySeven & PROVED & \checkmark & verified \\
\texttt{prime\_oneHundredThirtySeven} & CosTraceNormOneHundredThirtySeven & PROVED & \checkmark & verified \\
\texttt{degree\_oneHundredThree} & CosTraceNormOneHundredThree & PROVED & \checkmark & verified \\
\texttt{isIntegral\_and\_degree\_oneHundredThree} & CosTraceNormOneHundredThree & PROVED & \checkmark & verified \\
\texttt{isIntegral\_spectralGen\_oneHundredThree} & CosTraceNormOneHundredThree & PROVED & \checkmark & verified \\
\texttt{isIntegral\_spectralGen\_oneHundredThree\_Q} & CosTraceNormOneHundredThree & PROVED & \checkmark & verified \\
\texttt{oneHundredThree\_ne\_two} & CosTraceNormOneHundredThree & PROVED & \checkmark & verified \\
\texttt{oneHundredThree\_pack} & CosTraceNormOneHundredThree & PROVED & \checkmark & verified \\
\texttt{prime\_oneHundredThree} & CosTraceNormOneHundredThree & PROVED & \checkmark & verified \\
\texttt{degree\_oneHundredTwentySeven} & CosTraceNormOneHundredTwentySeven & PROVED & \checkmark & verified \\
\texttt{isIntegral\_and\_degree\_oneHundredTwentySeven} & CosTraceNormOneHundredTwentySeven & PROVED & \checkmark & verified \\
\texttt{isIntegral\_spectralGen\_oneHundredTwentySeven} & CosTraceNormOneHundredTwentySeven & PROVED & \checkmark & verified \\
\texttt{isIntegral\_spectralGen\_oneHundredTwentySeven\_Q} & CosTraceNormOneHundredTwentySeven & PROVED & \checkmark & verified \\
\texttt{oneHundredTwentySeven\_ne\_two} & CosTraceNormOneHundredTwentySeven & PROVED & \checkmark & verified \\
\texttt{oneHundredTwentySeven\_pack} & CosTraceNormOneHundredTwentySeven & PROVED & \checkmark & verified \\
\texttt{prime\_oneHundredTwentySeven} & CosTraceNormOneHundredTwentySeven & PROVED & \checkmark & verified \\
\texttt{degree\_oneThousandEightySeven} & CosTraceNormOneThousandEightySeven & PROVED & \checkmark & verified \\
\texttt{isIntegral\_and\_degree\_oneThousandEightySeven} & CosTraceNormOneThousandEightySeven & PROVED & \checkmark & verified \\
\texttt{isIntegral\_spectralGen\_oneThousandEightySeven} & CosTraceNormOneThousandEightySeven & PROVED & \checkmark & verified \\
\texttt{isIntegral\_spectralGen\_oneThousandEightySeven\_Q} & CosTraceNormOneThousandEightySeven & PROVED & \checkmark & verified \\
\texttt{oneThousandEightySeven\_ne\_two} & CosTraceNormOneThousandEightySeven & PROVED & \checkmark & verified \\
\texttt{oneThousandEightySeven\_pack} & CosTraceNormOneThousandEightySeven & PROVED & \checkmark & verified \\
\texttt{prime\_oneThousandEightySeven} & CosTraceNormOneThousandEightySeven & PROVED & \checkmark & verified \\
\texttt{degree\_oneThousandFiftyOne} & CosTraceNormOneThousandFiftyOne & PROVED & \checkmark & verified \\
\texttt{isIntegral\_and\_degree\_oneThousandFiftyOne} & CosTraceNormOneThousandFiftyOne & PROVED & \checkmark & verified \\
\texttt{isIntegral\_spectralGen\_oneThousandFiftyOne} & CosTraceNormOneThousandFiftyOne & PROVED & \checkmark & verified \\
\texttt{isIntegral\_spectralGen\_oneThousandFiftyOne\_Q} & CosTraceNormOneThousandFiftyOne & PROVED & \checkmark & verified \\
\texttt{oneThousandFiftyOne\_ne\_two} & CosTraceNormOneThousandFiftyOne & PROVED & \checkmark & verified \\
\texttt{oneThousandFiftyOne\_pack} & CosTraceNormOneThousandFiftyOne & PROVED & \checkmark & verified \\
\texttt{prime\_oneThousandFiftyOne} & CosTraceNormOneThousandFiftyOne & PROVED & \checkmark & verified \\
\texttt{degree\_oneThousandFortyNine} & CosTraceNormOneThousandFortyNine & PROVED & \checkmark & verified \\
\texttt{isIntegral\_and\_degree\_oneThousandFortyNine} & CosTraceNormOneThousandFortyNine & PROVED & \checkmark & verified \\
\texttt{isIntegral\_spectralGen\_oneThousandFortyNine} & CosTraceNormOneThousandFortyNine & PROVED & \checkmark & verified \\
\texttt{isIntegral\_spectralGen\_oneThousandFortyNine\_Q} & CosTraceNormOneThousandFortyNine & PROVED & \checkmark & verified \\
\texttt{oneThousandFortyNine\_ne\_two} & CosTraceNormOneThousandFortyNine & PROVED & \checkmark & verified \\
\texttt{oneThousandFortyNine\_pack} & CosTraceNormOneThousandFortyNine & PROVED & \checkmark & verified \\
\texttt{prime\_oneThousandFortyNine} & CosTraceNormOneThousandFortyNine & PROVED & \checkmark & verified \\
\texttt{degree\_oneThousandNine} & CosTraceNormOneThousandNine & PROVED & \checkmark & verified \\
\texttt{isIntegral\_and\_degree\_oneThousandNine} & CosTraceNormOneThousandNine & PROVED & \checkmark & verified \\
\texttt{isIntegral\_spectralGen\_oneThousandNine} & CosTraceNormOneThousandNine & PROVED & \checkmark & verified \\
\texttt{isIntegral\_spectralGen\_oneThousandNine\_Q} & CosTraceNormOneThousandNine & PROVED & \checkmark & verified \\
\texttt{oneThousandNine\_ne\_two} & CosTraceNormOneThousandNine & PROVED & \checkmark & verified \\
\texttt{oneThousandNine\_pack} & CosTraceNormOneThousandNine & PROVED & \checkmark & verified \\
\texttt{prime\_oneThousandNine} & CosTraceNormOneThousandNine & PROVED & \checkmark & verified \\
\texttt{degree\_oneThousandNineteen} & CosTraceNormOneThousandNineteen & PROVED & \checkmark & verified \\
\texttt{isIntegral\_and\_degree\_oneThousandNineteen} & CosTraceNormOneThousandNineteen & PROVED & \checkmark & verified \\
\texttt{isIntegral\_spectralGen\_oneThousandNineteen} & CosTraceNormOneThousandNineteen & PROVED & \checkmark & verified \\
\texttt{isIntegral\_spectralGen\_oneThousandNineteen\_Q} & CosTraceNormOneThousandNineteen & PROVED & \checkmark & verified \\
\texttt{oneThousandNineteen\_ne\_two} & CosTraceNormOneThousandNineteen & PROVED & \checkmark & verified \\
\texttt{oneThousandNineteen\_pack} & CosTraceNormOneThousandNineteen & PROVED & \checkmark & verified \\
\texttt{prime\_oneThousandNineteen} & CosTraceNormOneThousandNineteen & PROVED & \checkmark & verified \\
\texttt{degree\_oneThousandNinetyOne} & CosTraceNormOneThousandNinetyOne & PROVED & \checkmark & verified \\
\texttt{isIntegral\_and\_degree\_oneThousandNinetyOne} & CosTraceNormOneThousandNinetyOne & PROVED & \checkmark & verified \\
\texttt{isIntegral\_spectralGen\_oneThousandNinetyOne} & CosTraceNormOneThousandNinetyOne & PROVED & \checkmark & verified \\
\texttt{isIntegral\_spectralGen\_oneThousandNinetyOne\_Q} & CosTraceNormOneThousandNinetyOne & PROVED & \checkmark & verified \\
\texttt{oneThousandNinetyOne\_ne\_two} & CosTraceNormOneThousandNinetyOne & PROVED & \checkmark & verified \\
\texttt{oneThousandNinetyOne\_pack} & CosTraceNormOneThousandNinetyOne & PROVED & \checkmark & verified \\
\texttt{prime\_oneThousandNinetyOne} & CosTraceNormOneThousandNinetyOne & PROVED & \checkmark & verified \\
\texttt{degree\_oneThousandNinetyThree} & CosTraceNormOneThousandNinetyThree & PROVED & \checkmark & verified \\
\texttt{isIntegral\_and\_degree\_oneThousandNinetyThree} & CosTraceNormOneThousandNinetyThree & PROVED & \checkmark & verified \\
\texttt{isIntegral\_spectralGen\_oneThousandNinetyThree} & CosTraceNormOneThousandNinetyThree & PROVED & \checkmark & verified \\
\texttt{isIntegral\_spectralGen\_oneThousandNinetyThree\_Q} & CosTraceNormOneThousandNinetyThree & PROVED & \checkmark & verified \\
\texttt{oneThousandNinetyThree\_ne\_two} & CosTraceNormOneThousandNinetyThree & PROVED & \checkmark & verified \\
\texttt{oneThousandNinetyThree\_pack} & CosTraceNormOneThousandNinetyThree & PROVED & \checkmark & verified \\
\texttt{prime\_oneThousandNinetyThree} & CosTraceNormOneThousandNinetyThree & PROVED & \checkmark & verified \\
\texttt{degree\_oneThousandSixtyNine} & CosTraceNormOneThousandSixtyNine & PROVED & \checkmark & verified \\
\texttt{isIntegral\_and\_degree\_oneThousandSixtyNine} & CosTraceNormOneThousandSixtyNine & PROVED & \checkmark & verified \\
\texttt{isIntegral\_spectralGen\_oneThousandSixtyNine} & CosTraceNormOneThousandSixtyNine & PROVED & \checkmark & verified \\
\texttt{isIntegral\_spectralGen\_oneThousandSixtyNine\_Q} & CosTraceNormOneThousandSixtyNine & PROVED & \checkmark & verified \\
\texttt{oneThousandSixtyNine\_ne\_two} & CosTraceNormOneThousandSixtyNine & PROVED & \checkmark & verified \\
\texttt{oneThousandSixtyNine\_pack} & CosTraceNormOneThousandSixtyNine & PROVED & \checkmark & verified \\
\texttt{prime\_oneThousandSixtyNine} & CosTraceNormOneThousandSixtyNine & PROVED & \checkmark & verified \\
\texttt{degree\_oneThousandSixtyOne} & CosTraceNormOneThousandSixtyOne & PROVED & \checkmark & verified \\
\texttt{isIntegral\_and\_degree\_oneThousandSixtyOne} & CosTraceNormOneThousandSixtyOne & PROVED & \checkmark & verified \\
\texttt{isIntegral\_spectralGen\_oneThousandSixtyOne} & CosTraceNormOneThousandSixtyOne & PROVED & \checkmark & verified \\
\texttt{isIntegral\_spectralGen\_oneThousandSixtyOne\_Q} & CosTraceNormOneThousandSixtyOne & PROVED & \checkmark & verified \\
\texttt{oneThousandSixtyOne\_ne\_two} & CosTraceNormOneThousandSixtyOne & PROVED & \checkmark & verified \\
\texttt{oneThousandSixtyOne\_pack} & CosTraceNormOneThousandSixtyOne & PROVED & \checkmark & verified \\
\texttt{prime\_oneThousandSixtyOne} & CosTraceNormOneThousandSixtyOne & PROVED & \checkmark & verified \\
\texttt{degree\_oneThousandSixtyThree} & CosTraceNormOneThousandSixtyThree & PROVED & \checkmark & verified \\
\texttt{isIntegral\_and\_degree\_oneThousandSixtyThree} & CosTraceNormOneThousandSixtyThree & PROVED & \checkmark & verified \\
\texttt{isIntegral\_spectralGen\_oneThousandSixtyThree} & CosTraceNormOneThousandSixtyThree & PROVED & \checkmark & verified \\
\texttt{isIntegral\_spectralGen\_oneThousandSixtyThree\_Q} & CosTraceNormOneThousandSixtyThree & PROVED & \checkmark & verified \\
\texttt{oneThousandSixtyThree\_ne\_two} & CosTraceNormOneThousandSixtyThree & PROVED & \checkmark & verified \\
\texttt{oneThousandSixtyThree\_pack} & CosTraceNormOneThousandSixtyThree & PROVED & \checkmark & verified \\
\texttt{prime\_oneThousandSixtyThree} & CosTraceNormOneThousandSixtyThree & PROVED & \checkmark & verified \\
\texttt{degree\_oneThousandThirteen} & CosTraceNormOneThousandThirteen & PROVED & \checkmark & verified \\
\texttt{isIntegral\_and\_degree\_oneThousandThirteen} & CosTraceNormOneThousandThirteen & PROVED & \checkmark & verified \\
\texttt{isIntegral\_spectralGen\_oneThousandThirteen} & CosTraceNormOneThousandThirteen & PROVED & \checkmark & verified \\
\texttt{isIntegral\_spectralGen\_oneThousandThirteen\_Q} & CosTraceNormOneThousandThirteen & PROVED & \checkmark & verified \\
\texttt{oneThousandThirteen\_ne\_two} & CosTraceNormOneThousandThirteen & PROVED & \checkmark & verified \\
\texttt{oneThousandThirteen\_pack} & CosTraceNormOneThousandThirteen & PROVED & \checkmark & verified \\
\texttt{prime\_oneThousandThirteen} & CosTraceNormOneThousandThirteen & PROVED & \checkmark & verified \\
\texttt{degree\_oneThousandThirtyNine} & CosTraceNormOneThousandThirtyNine & PROVED & \checkmark & verified \\
\texttt{isIntegral\_and\_degree\_oneThousandThirtyNine} & CosTraceNormOneThousandThirtyNine & PROVED & \checkmark & verified \\
\texttt{isIntegral\_spectralGen\_oneThousandThirtyNine} & CosTraceNormOneThousandThirtyNine & PROVED & \checkmark & verified \\
\texttt{isIntegral\_spectralGen\_oneThousandThirtyNine\_Q} & CosTraceNormOneThousandThirtyNine & PROVED & \checkmark & verified \\
\texttt{oneThousandThirtyNine\_ne\_two} & CosTraceNormOneThousandThirtyNine & PROVED & \checkmark & verified \\
\texttt{oneThousandThirtyNine\_pack} & CosTraceNormOneThousandThirtyNine & PROVED & \checkmark & verified \\
\texttt{prime\_oneThousandThirtyNine} & CosTraceNormOneThousandThirtyNine & PROVED & \checkmark & verified \\
\texttt{degree\_oneThousandThirtyOne} & CosTraceNormOneThousandThirtyOne & PROVED & \checkmark & verified \\
\texttt{isIntegral\_and\_degree\_oneThousandThirtyOne} & CosTraceNormOneThousandThirtyOne & PROVED & \checkmark & verified \\
\texttt{isIntegral\_spectralGen\_oneThousandThirtyOne} & CosTraceNormOneThousandThirtyOne & PROVED & \checkmark & verified \\
\texttt{isIntegral\_spectralGen\_oneThousandThirtyOne\_Q} & CosTraceNormOneThousandThirtyOne & PROVED & \checkmark & verified \\
\texttt{oneThousandThirtyOne\_ne\_two} & CosTraceNormOneThousandThirtyOne & PROVED & \checkmark & verified \\
\texttt{oneThousandThirtyOne\_pack} & CosTraceNormOneThousandThirtyOne & PROVED & \checkmark & verified \\
\texttt{prime\_oneThousandThirtyOne} & CosTraceNormOneThousandThirtyOne & PROVED & \checkmark & verified \\
\texttt{degree\_oneThousandThirtyThree} & CosTraceNormOneThousandThirtyThree & PROVED & \checkmark & verified \\
\texttt{isIntegral\_and\_degree\_oneThousandThirtyThree} & CosTraceNormOneThousandThirtyThree & PROVED & \checkmark & verified \\
\texttt{isIntegral\_spectralGen\_oneThousandThirtyThree} & CosTraceNormOneThousandThirtyThree & PROVED & \checkmark & verified \\
\texttt{isIntegral\_spectralGen\_oneThousandThirtyThree\_Q} & CosTraceNormOneThousandThirtyThree & PROVED & \checkmark & verified \\
\texttt{oneThousandThirtyThree\_ne\_two} & CosTraceNormOneThousandThirtyThree & PROVED & \checkmark & verified \\
\texttt{oneThousandThirtyThree\_pack} & CosTraceNormOneThousandThirtyThree & PROVED & \checkmark & verified \\
\texttt{prime\_oneThousandThirtyThree} & CosTraceNormOneThousandThirtyThree & PROVED & \checkmark & verified \\
\texttt{degree\_oneThousandTwentyOne} & CosTraceNormOneThousandTwentyOne & PROVED & \checkmark & verified \\
\texttt{isIntegral\_and\_degree\_oneThousandTwentyOne} & CosTraceNormOneThousandTwentyOne & PROVED & \checkmark & verified \\
\texttt{isIntegral\_spectralGen\_oneThousandTwentyOne} & CosTraceNormOneThousandTwentyOne & PROVED & \checkmark & verified \\
\texttt{isIntegral\_spectralGen\_oneThousandTwentyOne\_Q} & CosTraceNormOneThousandTwentyOne & PROVED & \checkmark & verified \\
\texttt{oneThousandTwentyOne\_ne\_two} & CosTraceNormOneThousandTwentyOne & PROVED & \checkmark & verified \\
\texttt{oneThousandTwentyOne\_pack} & CosTraceNormOneThousandTwentyOne & PROVED & \checkmark & verified \\
\texttt{prime\_oneThousandTwentyOne} & CosTraceNormOneThousandTwentyOne & PROVED & \checkmark & verified \\
\texttt{degree\_sevenHundredEightySeven} & CosTraceNormSevenHundredEightySeven & PROVED & \checkmark & verified \\
\texttt{isIntegral\_and\_degree\_sevenHundredEightySeven} & CosTraceNormSevenHundredEightySeven & PROVED & \checkmark & verified \\
\texttt{isIntegral\_spectralGen\_sevenHundredEightySeven} & CosTraceNormSevenHundredEightySeven & PROVED & \checkmark & verified \\
\texttt{isIntegral\_spectralGen\_sevenHundredEightySeven\_Q} & CosTraceNormSevenHundredEightySeven & PROVED & \checkmark & verified \\
\texttt{prime\_sevenHundredEightySeven} & CosTraceNormSevenHundredEightySeven & PROVED & \checkmark & verified \\
\texttt{sevenHundredEightySeven\_ne\_two} & CosTraceNormSevenHundredEightySeven & PROVED & \checkmark & verified \\
\texttt{sevenHundredEightySeven\_pack} & CosTraceNormSevenHundredEightySeven & PROVED & \checkmark & verified \\
\texttt{degree\_sevenHundredFiftyOne} & CosTraceNormSevenHundredFiftyOne & PROVED & \checkmark & verified \\
\texttt{isIntegral\_and\_degree\_sevenHundredFiftyOne} & CosTraceNormSevenHundredFiftyOne & PROVED & \checkmark & verified \\
\texttt{isIntegral\_spectralGen\_sevenHundredFiftyOne} & CosTraceNormSevenHundredFiftyOne & PROVED & \checkmark & verified \\
\texttt{isIntegral\_spectralGen\_sevenHundredFiftyOne\_Q} & CosTraceNormSevenHundredFiftyOne & PROVED & \checkmark & verified \\
\texttt{prime\_sevenHundredFiftyOne} & CosTraceNormSevenHundredFiftyOne & PROVED & \checkmark & verified \\
\texttt{sevenHundredFiftyOne\_ne\_two} & CosTraceNormSevenHundredFiftyOne & PROVED & \checkmark & verified \\
\texttt{sevenHundredFiftyOne\_pack} & CosTraceNormSevenHundredFiftyOne & PROVED & \checkmark & verified \\
\texttt{degree\_sevenHundredFiftySeven} & CosTraceNormSevenHundredFiftySeven & PROVED & \checkmark & verified \\
\texttt{isIntegral\_and\_degree\_sevenHundredFiftySeven} & CosTraceNormSevenHundredFiftySeven & PROVED & \checkmark & verified \\
\texttt{isIntegral\_spectralGen\_sevenHundredFiftySeven} & CosTraceNormSevenHundredFiftySeven & PROVED & \checkmark & verified \\
\texttt{isIntegral\_spectralGen\_sevenHundredFiftySeven\_Q} & CosTraceNormSevenHundredFiftySeven & PROVED & \checkmark & verified \\
\texttt{prime\_sevenHundredFiftySeven} & CosTraceNormSevenHundredFiftySeven & PROVED & \checkmark & verified \\
\texttt{sevenHundredFiftySeven\_ne\_two} & CosTraceNormSevenHundredFiftySeven & PROVED & \checkmark & verified \\
\texttt{sevenHundredFiftySeven\_pack} & CosTraceNormSevenHundredFiftySeven & PROVED & \checkmark & verified \\
\texttt{degree\_sevenHundredFortyThree} & CosTraceNormSevenHundredFortyThree & PROVED & \checkmark & verified \\
\texttt{isIntegral\_and\_degree\_sevenHundredFortyThree} & CosTraceNormSevenHundredFortyThree & PROVED & \checkmark & verified \\
\texttt{isIntegral\_spectralGen\_sevenHundredFortyThree} & CosTraceNormSevenHundredFortyThree & PROVED & \checkmark & verified \\
\texttt{isIntegral\_spectralGen\_sevenHundredFortyThree\_Q} & CosTraceNormSevenHundredFortyThree & PROVED & \checkmark & verified \\
\texttt{prime\_sevenHundredFortyThree} & CosTraceNormSevenHundredFortyThree & PROVED & \checkmark & verified \\
\texttt{sevenHundredFortyThree\_ne\_two} & CosTraceNormSevenHundredFortyThree & PROVED & \checkmark & verified \\
\texttt{sevenHundredFortyThree\_pack} & CosTraceNormSevenHundredFortyThree & PROVED & \checkmark & verified \\
\texttt{degree\_sevenHundredNine} & CosTraceNormSevenHundredNine & PROVED & \checkmark & verified \\
\texttt{isIntegral\_and\_degree\_sevenHundredNine} & CosTraceNormSevenHundredNine & PROVED & \checkmark & verified \\
\texttt{isIntegral\_spectralGen\_sevenHundredNine} & CosTraceNormSevenHundredNine & PROVED & \checkmark & verified \\
\texttt{isIntegral\_spectralGen\_sevenHundredNine\_Q} & CosTraceNormSevenHundredNine & PROVED & \checkmark & verified \\
\texttt{prime\_sevenHundredNine} & CosTraceNormSevenHundredNine & PROVED & \checkmark & verified \\
\texttt{sevenHundredNine\_ne\_two} & CosTraceNormSevenHundredNine & PROVED & \checkmark & verified \\
\texttt{sevenHundredNine\_pack} & CosTraceNormSevenHundredNine & PROVED & \checkmark & verified \\
\texttt{degree\_sevenHundredNineteen} & CosTraceNormSevenHundredNineteen & PROVED & \checkmark & verified \\
\texttt{isIntegral\_and\_degree\_sevenHundredNineteen} & CosTraceNormSevenHundredNineteen & PROVED & \checkmark & verified \\
\texttt{isIntegral\_spectralGen\_sevenHundredNineteen} & CosTraceNormSevenHundredNineteen & PROVED & \checkmark & verified \\
\texttt{isIntegral\_spectralGen\_sevenHundredNineteen\_Q} & CosTraceNormSevenHundredNineteen & PROVED & \checkmark & verified \\
\texttt{prime\_sevenHundredNineteen} & CosTraceNormSevenHundredNineteen & PROVED & \checkmark & verified \\
\texttt{sevenHundredNineteen\_ne\_two} & CosTraceNormSevenHundredNineteen & PROVED & \checkmark & verified \\
\texttt{sevenHundredNineteen\_pack} & CosTraceNormSevenHundredNineteen & PROVED & \checkmark & verified \\
\texttt{degree\_sevenHundredNinetySeven} & CosTraceNormSevenHundredNinetySeven & PROVED & \checkmark & verified \\
\texttt{isIntegral\_and\_degree\_sevenHundredNinetySeven} & CosTraceNormSevenHundredNinetySeven & PROVED & \checkmark & verified \\
\texttt{isIntegral\_spectralGen\_sevenHundredNinetySeven} & CosTraceNormSevenHundredNinetySeven & PROVED & \checkmark & verified \\
\texttt{isIntegral\_spectralGen\_sevenHundredNinetySeven\_Q} & CosTraceNormSevenHundredNinetySeven & PROVED & \checkmark & verified \\
\texttt{prime\_sevenHundredNinetySeven} & CosTraceNormSevenHundredNinetySeven & PROVED & \checkmark & verified \\
\texttt{sevenHundredNinetySeven\_ne\_two} & CosTraceNormSevenHundredNinetySeven & PROVED & \checkmark & verified \\
\texttt{sevenHundredNinetySeven\_pack} & CosTraceNormSevenHundredNinetySeven & PROVED & \checkmark & verified \\
\texttt{degree\_sevenHundredOne} & CosTraceNormSevenHundredOne & PROVED & \checkmark & verified \\
\texttt{isIntegral\_and\_degree\_sevenHundredOne} & CosTraceNormSevenHundredOne & PROVED & \checkmark & verified \\
\texttt{isIntegral\_spectralGen\_sevenHundredOne} & CosTraceNormSevenHundredOne & PROVED & \checkmark & verified \\
\texttt{isIntegral\_spectralGen\_sevenHundredOne\_Q} & CosTraceNormSevenHundredOne & PROVED & \checkmark & verified \\
\texttt{prime\_sevenHundredOne} & CosTraceNormSevenHundredOne & PROVED & \checkmark & verified \\
\texttt{sevenHundredOne\_ne\_two} & CosTraceNormSevenHundredOne & PROVED & \checkmark & verified \\
\texttt{sevenHundredOne\_pack} & CosTraceNormSevenHundredOne & PROVED & \checkmark & verified \\
\texttt{degree\_sevenHundredSeventyThree} & CosTraceNormSevenHundredSeventyThree & PROVED & \checkmark & verified \\
\texttt{isIntegral\_and\_degree\_sevenHundredSeventyThree} & CosTraceNormSevenHundredSeventyThree & PROVED & \checkmark & verified \\
\texttt{isIntegral\_spectralGen\_sevenHundredSeventyThree} & CosTraceNormSevenHundredSeventyThree & PROVED & \checkmark & verified \\
\texttt{isIntegral\_spectralGen\_sevenHundredSeventyThree\_Q} & CosTraceNormSevenHundredSeventyThree & PROVED & \checkmark & verified \\
\texttt{prime\_sevenHundredSeventyThree} & CosTraceNormSevenHundredSeventyThree & PROVED & \checkmark & verified \\
\texttt{sevenHundredSeventyThree\_ne\_two} & CosTraceNormSevenHundredSeventyThree & PROVED & \checkmark & verified \\
\texttt{sevenHundredSeventyThree\_pack} & CosTraceNormSevenHundredSeventyThree & PROVED & \checkmark & verified \\
\texttt{degree\_sevenHundredSixtyNine} & CosTraceNormSevenHundredSixtyNine & PROVED & \checkmark & verified \\
\texttt{isIntegral\_and\_degree\_sevenHundredSixtyNine} & CosTraceNormSevenHundredSixtyNine & PROVED & \checkmark & verified \\
\texttt{isIntegral\_spectralGen\_sevenHundredSixtyNine} & CosTraceNormSevenHundredSixtyNine & PROVED & \checkmark & verified \\
\texttt{isIntegral\_spectralGen\_sevenHundredSixtyNine\_Q} & CosTraceNormSevenHundredSixtyNine & PROVED & \checkmark & verified \\
\texttt{prime\_sevenHundredSixtyNine} & CosTraceNormSevenHundredSixtyNine & PROVED & \checkmark & verified \\
\texttt{sevenHundredSixtyNine\_ne\_two} & CosTraceNormSevenHundredSixtyNine & PROVED & \checkmark & verified \\
\texttt{sevenHundredSixtyNine\_pack} & CosTraceNormSevenHundredSixtyNine & PROVED & \checkmark & verified \\
\texttt{degree\_sevenHundredSixtyOne} & CosTraceNormSevenHundredSixtyOne & PROVED & \checkmark & verified \\
\texttt{isIntegral\_and\_degree\_sevenHundredSixtyOne} & CosTraceNormSevenHundredSixtyOne & PROVED & \checkmark & verified \\
\texttt{isIntegral\_spectralGen\_sevenHundredSixtyOne} & CosTraceNormSevenHundredSixtyOne & PROVED & \checkmark & verified \\
\texttt{isIntegral\_spectralGen\_sevenHundredSixtyOne\_Q} & CosTraceNormSevenHundredSixtyOne & PROVED & \checkmark & verified \\
\texttt{prime\_sevenHundredSixtyOne} & CosTraceNormSevenHundredSixtyOne & PROVED & \checkmark & verified \\
\texttt{sevenHundredSixtyOne\_ne\_two} & CosTraceNormSevenHundredSixtyOne & PROVED & \checkmark & verified \\
\texttt{sevenHundredSixtyOne\_pack} & CosTraceNormSevenHundredSixtyOne & PROVED & \checkmark & verified \\
\texttt{degree\_sevenHundredThirtyNine} & CosTraceNormSevenHundredThirtyNine & PROVED & \checkmark & verified \\
\texttt{isIntegral\_and\_degree\_sevenHundredThirtyNine} & CosTraceNormSevenHundredThirtyNine & PROVED & \checkmark & verified \\
\texttt{isIntegral\_spectralGen\_sevenHundredThirtyNine} & CosTraceNormSevenHundredThirtyNine & PROVED & \checkmark & verified \\
\texttt{isIntegral\_spectralGen\_sevenHundredThirtyNine\_Q} & CosTraceNormSevenHundredThirtyNine & PROVED & \checkmark & verified \\
\texttt{prime\_sevenHundredThirtyNine} & CosTraceNormSevenHundredThirtyNine & PROVED & \checkmark & verified \\
\texttt{sevenHundredThirtyNine\_ne\_two} & CosTraceNormSevenHundredThirtyNine & PROVED & \checkmark & verified \\
\texttt{sevenHundredThirtyNine\_pack} & CosTraceNormSevenHundredThirtyNine & PROVED & \checkmark & verified \\
\texttt{degree\_sevenHundredThirtyThree} & CosTraceNormSevenHundredThirtyThree & PROVED & \checkmark & verified \\
\texttt{isIntegral\_and\_degree\_sevenHundredThirtyThree} & CosTraceNormSevenHundredThirtyThree & PROVED & \checkmark & verified \\
\texttt{isIntegral\_spectralGen\_sevenHundredThirtyThree} & CosTraceNormSevenHundredThirtyThree & PROVED & \checkmark & verified \\
\texttt{isIntegral\_spectralGen\_sevenHundredThirtyThree\_Q} & CosTraceNormSevenHundredThirtyThree & PROVED & \checkmark & verified \\
\texttt{prime\_sevenHundredThirtyThree} & CosTraceNormSevenHundredThirtyThree & PROVED & \checkmark & verified \\
\texttt{sevenHundredThirtyThree\_ne\_two} & CosTraceNormSevenHundredThirtyThree & PROVED & \checkmark & verified \\
\texttt{sevenHundredThirtyThree\_pack} & CosTraceNormSevenHundredThirtyThree & PROVED & \checkmark & verified \\
\texttt{degree\_sevenHundredTwentySeven} & CosTraceNormSevenHundredTwentySeven & PROVED & \checkmark & verified \\
\texttt{isIntegral\_and\_degree\_sevenHundredTwentySeven} & CosTraceNormSevenHundredTwentySeven & PROVED & \checkmark & verified \\
\texttt{isIntegral\_spectralGen\_sevenHundredTwentySeven} & CosTraceNormSevenHundredTwentySeven & PROVED & \checkmark & verified \\
\texttt{isIntegral\_spectralGen\_sevenHundredTwentySeven\_Q} & CosTraceNormSevenHundredTwentySeven & PROVED & \checkmark & verified \\
\texttt{prime\_sevenHundredTwentySeven} & CosTraceNormSevenHundredTwentySeven & PROVED & \checkmark & verified \\
\texttt{sevenHundredTwentySeven\_ne\_two} & CosTraceNormSevenHundredTwentySeven & PROVED & \checkmark & verified \\
\texttt{sevenHundredTwentySeven\_pack} & CosTraceNormSevenHundredTwentySeven & PROVED & \checkmark & verified \\
\texttt{degree\_seventeen} & CosTraceNormSeventeen & PROVED & \checkmark & verified \\
\texttt{isIntegral\_and\_degree\_seventeen} & CosTraceNormSeventeen & PROVED & \checkmark & verified \\
\texttt{isIntegral\_spectralGen\_seventeen} & CosTraceNormSeventeen & PROVED & \checkmark & verified \\
\texttt{isIntegral\_spectralGen\_seventeen\_Q} & CosTraceNormSeventeen & PROVED & \checkmark & verified \\
\texttt{prime\_seventeen} & CosTraceNormSeventeen & PROVED & \checkmark & verified \\
\texttt{seventeen\_ne\_two} & CosTraceNormSeventeen & PROVED & \checkmark & verified \\
\texttt{seventeen\_pack} & CosTraceNormSeventeen & PROVED & \checkmark & verified \\
\texttt{degree\_seventyNine} & CosTraceNormSeventyNine & PROVED & \checkmark & verified \\
\texttt{isIntegral\_and\_degree\_seventyNine} & CosTraceNormSeventyNine & PROVED & \checkmark & verified \\
\texttt{isIntegral\_spectralGen\_seventyNine} & CosTraceNormSeventyNine & PROVED & \checkmark & verified \\
\texttt{isIntegral\_spectralGen\_seventyNine\_Q} & CosTraceNormSeventyNine & PROVED & \checkmark & verified \\
\texttt{prime\_seventyNine} & CosTraceNormSeventyNine & PROVED & \checkmark & verified \\
\texttt{seventyNine\_ne\_two} & CosTraceNormSeventyNine & PROVED & \checkmark & verified \\
\texttt{seventyNine\_pack} & CosTraceNormSeventyNine & PROVED & \checkmark & verified \\
\texttt{degree\_seventyOne} & CosTraceNormSeventyOne & PROVED & \checkmark & verified \\
\texttt{isIntegral\_and\_degree\_seventyOne} & CosTraceNormSeventyOne & PROVED & \checkmark & verified \\
\texttt{isIntegral\_spectralGen\_seventyOne} & CosTraceNormSeventyOne & PROVED & \checkmark & verified \\
\texttt{isIntegral\_spectralGen\_seventyOne\_Q} & CosTraceNormSeventyOne & PROVED & \checkmark & verified \\
\texttt{prime\_seventyOne} & CosTraceNormSeventyOne & PROVED & \checkmark & verified \\
\texttt{seventyOne\_ne\_two} & CosTraceNormSeventyOne & PROVED & \checkmark & verified \\
\texttt{seventyOne\_pack} & CosTraceNormSeventyOne & PROVED & \checkmark & verified \\
\texttt{degree\_seventyThree} & CosTraceNormSeventyThree & PROVED & \checkmark & verified \\
\texttt{isIntegral\_and\_degree\_seventyThree} & CosTraceNormSeventyThree & PROVED & \checkmark & verified \\
\texttt{isIntegral\_spectralGen\_seventyThree} & CosTraceNormSeventyThree & PROVED & \checkmark & verified \\
\texttt{isIntegral\_spectralGen\_seventyThree\_Q} & CosTraceNormSeventyThree & PROVED & \checkmark & verified \\
\texttt{prime\_seventyThree} & CosTraceNormSeventyThree & PROVED & \checkmark & verified \\
\texttt{seventyThree\_ne\_two} & CosTraceNormSeventyThree & PROVED & \checkmark & verified \\
\texttt{seventyThree\_pack} & CosTraceNormSeventyThree & PROVED & \checkmark & verified \\
\texttt{degree\_sixHundredEightyThree} & CosTraceNormSixHundredEightyThree & PROVED & \checkmark & verified \\
\texttt{isIntegral\_and\_degree\_sixHundredEightyThree} & CosTraceNormSixHundredEightyThree & PROVED & \checkmark & verified \\
\texttt{isIntegral\_spectralGen\_sixHundredEightyThree} & CosTraceNormSixHundredEightyThree & PROVED & \checkmark & verified \\
\texttt{isIntegral\_spectralGen\_sixHundredEightyThree\_Q} & CosTraceNormSixHundredEightyThree & PROVED & \checkmark & verified \\
\texttt{prime\_sixHundredEightyThree} & CosTraceNormSixHundredEightyThree & PROVED & \checkmark & verified \\
\texttt{sixHundredEightyThree\_ne\_two} & CosTraceNormSixHundredEightyThree & PROVED & \checkmark & verified \\
\texttt{sixHundredEightyThree\_pack} & CosTraceNormSixHundredEightyThree & PROVED & \checkmark & verified \\
\texttt{degree\_sixHundredFiftyNine} & CosTraceNormSixHundredFiftyNine & PROVED & \checkmark & verified \\
\texttt{isIntegral\_and\_degree\_sixHundredFiftyNine} & CosTraceNormSixHundredFiftyNine & PROVED & \checkmark & verified \\
\texttt{isIntegral\_spectralGen\_sixHundredFiftyNine} & CosTraceNormSixHundredFiftyNine & PROVED & \checkmark & verified \\
\texttt{isIntegral\_spectralGen\_sixHundredFiftyNine\_Q} & CosTraceNormSixHundredFiftyNine & PROVED & \checkmark & verified \\
\texttt{prime\_sixHundredFiftyNine} & CosTraceNormSixHundredFiftyNine & PROVED & \checkmark & verified \\
\texttt{sixHundredFiftyNine\_ne\_two} & CosTraceNormSixHundredFiftyNine & PROVED & \checkmark & verified \\
\texttt{sixHundredFiftyNine\_pack} & CosTraceNormSixHundredFiftyNine & PROVED & \checkmark & verified \\
\texttt{degree\_sixHundredFiftyThree} & CosTraceNormSixHundredFiftyThree & PROVED & \checkmark & verified \\
\texttt{isIntegral\_and\_degree\_sixHundredFiftyThree} & CosTraceNormSixHundredFiftyThree & PROVED & \checkmark & verified \\
\texttt{isIntegral\_spectralGen\_sixHundredFiftyThree} & CosTraceNormSixHundredFiftyThree & PROVED & \checkmark & verified \\
\texttt{isIntegral\_spectralGen\_sixHundredFiftyThree\_Q} & CosTraceNormSixHundredFiftyThree & PROVED & \checkmark & verified \\
\texttt{prime\_sixHundredFiftyThree} & CosTraceNormSixHundredFiftyThree & PROVED & \checkmark & verified \\
\texttt{sixHundredFiftyThree\_ne\_two} & CosTraceNormSixHundredFiftyThree & PROVED & \checkmark & verified \\
\texttt{sixHundredFiftyThree\_pack} & CosTraceNormSixHundredFiftyThree & PROVED & \checkmark & verified \\
\texttt{degree\_sixHundredFortyOne} & CosTraceNormSixHundredFortyOne & PROVED & \checkmark & verified \\
\texttt{isIntegral\_and\_degree\_sixHundredFortyOne} & CosTraceNormSixHundredFortyOne & PROVED & \checkmark & verified \\
\texttt{isIntegral\_spectralGen\_sixHundredFortyOne} & CosTraceNormSixHundredFortyOne & PROVED & \checkmark & verified \\
\texttt{isIntegral\_spectralGen\_sixHundredFortyOne\_Q} & CosTraceNormSixHundredFortyOne & PROVED & \checkmark & verified \\
\texttt{prime\_sixHundredFortyOne} & CosTraceNormSixHundredFortyOne & PROVED & \checkmark & verified \\
\texttt{sixHundredFortyOne\_ne\_two} & CosTraceNormSixHundredFortyOne & PROVED & \checkmark & verified \\
\texttt{sixHundredFortyOne\_pack} & CosTraceNormSixHundredFortyOne & PROVED & \checkmark & verified \\
\texttt{degree\_sixHundredFortySeven} & CosTraceNormSixHundredFortySeven & PROVED & \checkmark & verified \\
\texttt{isIntegral\_and\_degree\_sixHundredFortySeven} & CosTraceNormSixHundredFortySeven & PROVED & \checkmark & verified \\
\texttt{isIntegral\_spectralGen\_sixHundredFortySeven} & CosTraceNormSixHundredFortySeven & PROVED & \checkmark & verified \\
\texttt{isIntegral\_spectralGen\_sixHundredFortySeven\_Q} & CosTraceNormSixHundredFortySeven & PROVED & \checkmark & verified \\
\texttt{prime\_sixHundredFortySeven} & CosTraceNormSixHundredFortySeven & PROVED & \checkmark & verified \\
\texttt{sixHundredFortySeven\_ne\_two} & CosTraceNormSixHundredFortySeven & PROVED & \checkmark & verified \\
\texttt{sixHundredFortySeven\_pack} & CosTraceNormSixHundredFortySeven & PROVED & \checkmark & verified \\
\texttt{degree\_sixHundredFortyThree} & CosTraceNormSixHundredFortyThree & PROVED & \checkmark & verified \\
\texttt{isIntegral\_and\_degree\_sixHundredFortyThree} & CosTraceNormSixHundredFortyThree & PROVED & \checkmark & verified \\
\texttt{isIntegral\_spectralGen\_sixHundredFortyThree} & CosTraceNormSixHundredFortyThree & PROVED & \checkmark & verified \\
\texttt{isIntegral\_spectralGen\_sixHundredFortyThree\_Q} & CosTraceNormSixHundredFortyThree & PROVED & \checkmark & verified \\
\texttt{prime\_sixHundredFortyThree} & CosTraceNormSixHundredFortyThree & PROVED & \checkmark & verified \\
\texttt{sixHundredFortyThree\_ne\_two} & CosTraceNormSixHundredFortyThree & PROVED & \checkmark & verified \\
\texttt{sixHundredFortyThree\_pack} & CosTraceNormSixHundredFortyThree & PROVED & \checkmark & verified \\
\texttt{degree\_sixHundredNineteen} & CosTraceNormSixHundredNineteen & PROVED & \checkmark & verified \\
\texttt{isIntegral\_and\_degree\_sixHundredNineteen} & CosTraceNormSixHundredNineteen & PROVED & \checkmark & verified \\
\texttt{isIntegral\_spectralGen\_sixHundredNineteen} & CosTraceNormSixHundredNineteen & PROVED & \checkmark & verified \\
\texttt{isIntegral\_spectralGen\_sixHundredNineteen\_Q} & CosTraceNormSixHundredNineteen & PROVED & \checkmark & verified \\
\texttt{prime\_sixHundredNineteen} & CosTraceNormSixHundredNineteen & PROVED & \checkmark & verified \\
\texttt{sixHundredNineteen\_ne\_two} & CosTraceNormSixHundredNineteen & PROVED & \checkmark & verified \\
\texttt{sixHundredNineteen\_pack} & CosTraceNormSixHundredNineteen & PROVED & \checkmark & verified \\
\texttt{degree\_sixHundredNinetyOne} & CosTraceNormSixHundredNinetyOne & PROVED & \checkmark & verified \\
\texttt{isIntegral\_and\_degree\_sixHundredNinetyOne} & CosTraceNormSixHundredNinetyOne & PROVED & \checkmark & verified \\
\texttt{isIntegral\_spectralGen\_sixHundredNinetyOne} & CosTraceNormSixHundredNinetyOne & PROVED & \checkmark & verified \\
\texttt{isIntegral\_spectralGen\_sixHundredNinetyOne\_Q} & CosTraceNormSixHundredNinetyOne & PROVED & \checkmark & verified \\
\texttt{prime\_sixHundredNinetyOne} & CosTraceNormSixHundredNinetyOne & PROVED & \checkmark & verified \\
\texttt{sixHundredNinetyOne\_ne\_two} & CosTraceNormSixHundredNinetyOne & PROVED & \checkmark & verified \\
\texttt{sixHundredNinetyOne\_pack} & CosTraceNormSixHundredNinetyOne & PROVED & \checkmark & verified \\
\texttt{degree\_sixHundredOne} & CosTraceNormSixHundredOne & PROVED & \checkmark & verified \\
\texttt{isIntegral\_and\_degree\_sixHundredOne} & CosTraceNormSixHundredOne & PROVED & \checkmark & verified \\
\texttt{isIntegral\_spectralGen\_sixHundredOne} & CosTraceNormSixHundredOne & PROVED & \checkmark & verified \\
\texttt{isIntegral\_spectralGen\_sixHundredOne\_Q} & CosTraceNormSixHundredOne & PROVED & \checkmark & verified \\
\texttt{prime\_sixHundredOne} & CosTraceNormSixHundredOne & PROVED & \checkmark & verified \\
\texttt{sixHundredOne\_ne\_two} & CosTraceNormSixHundredOne & PROVED & \checkmark & verified \\
\texttt{sixHundredOne\_pack} & CosTraceNormSixHundredOne & PROVED & \checkmark & verified \\
\texttt{degree\_sixHundredSeven} & CosTraceNormSixHundredSeven & PROVED & \checkmark & verified \\
\texttt{isIntegral\_and\_degree\_sixHundredSeven} & CosTraceNormSixHundredSeven & PROVED & \checkmark & verified \\
\texttt{isIntegral\_spectralGen\_sixHundredSeven} & CosTraceNormSixHundredSeven & PROVED & \checkmark & verified \\
\texttt{isIntegral\_spectralGen\_sixHundredSeven\_Q} & CosTraceNormSixHundredSeven & PROVED & \checkmark & verified \\
\texttt{prime\_sixHundredSeven} & CosTraceNormSixHundredSeven & PROVED & \checkmark & verified \\
\texttt{sixHundredSeven\_ne\_two} & CosTraceNormSixHundredSeven & PROVED & \checkmark & verified \\
\texttt{sixHundredSeven\_pack} & CosTraceNormSixHundredSeven & PROVED & \checkmark & verified \\
\texttt{degree\_sixHundredSeventeen} & CosTraceNormSixHundredSeventeen & PROVED & \checkmark & verified \\
\texttt{isIntegral\_and\_degree\_sixHundredSeventeen} & CosTraceNormSixHundredSeventeen & PROVED & \checkmark & verified \\
\texttt{isIntegral\_spectralGen\_sixHundredSeventeen} & CosTraceNormSixHundredSeventeen & PROVED & \checkmark & verified \\
\texttt{isIntegral\_spectralGen\_sixHundredSeventeen\_Q} & CosTraceNormSixHundredSeventeen & PROVED & \checkmark & verified \\
\texttt{prime\_sixHundredSeventeen} & CosTraceNormSixHundredSeventeen & PROVED & \checkmark & verified \\
\texttt{sixHundredSeventeen\_ne\_two} & CosTraceNormSixHundredSeventeen & PROVED & \checkmark & verified \\
\texttt{sixHundredSeventeen\_pack} & CosTraceNormSixHundredSeventeen & PROVED & \checkmark & verified \\
\texttt{degree\_sixHundredSeventySeven} & CosTraceNormSixHundredSeventySeven & PROVED & \checkmark & verified \\
\texttt{isIntegral\_and\_degree\_sixHundredSeventySeven} & CosTraceNormSixHundredSeventySeven & PROVED & \checkmark & verified \\
\texttt{isIntegral\_spectralGen\_sixHundredSeventySeven} & CosTraceNormSixHundredSeventySeven & PROVED & \checkmark & verified \\
\texttt{isIntegral\_spectralGen\_sixHundredSeventySeven\_Q} & CosTraceNormSixHundredSeventySeven & PROVED & \checkmark & verified \\
\texttt{prime\_sixHundredSeventySeven} & CosTraceNormSixHundredSeventySeven & PROVED & \checkmark & verified \\
\texttt{sixHundredSeventySeven\_ne\_two} & CosTraceNormSixHundredSeventySeven & PROVED & \checkmark & verified \\
\texttt{sixHundredSeventySeven\_pack} & CosTraceNormSixHundredSeventySeven & PROVED & \checkmark & verified \\
\texttt{degree\_sixHundredSeventyThree} & CosTraceNormSixHundredSeventyThree & PROVED & \checkmark & verified \\
\texttt{isIntegral\_and\_degree\_sixHundredSeventyThree} & CosTraceNormSixHundredSeventyThree & PROVED & \checkmark & verified \\
\texttt{isIntegral\_spectralGen\_sixHundredSeventyThree} & CosTraceNormSixHundredSeventyThree & PROVED & \checkmark & verified \\
\texttt{isIntegral\_spectralGen\_sixHundredSeventyThree\_Q} & CosTraceNormSixHundredSeventyThree & PROVED & \checkmark & verified \\
\texttt{prime\_sixHundredSeventyThree} & CosTraceNormSixHundredSeventyThree & PROVED & \checkmark & verified \\
\texttt{sixHundredSeventyThree\_ne\_two} & CosTraceNormSixHundredSeventyThree & PROVED & \checkmark & verified \\
\texttt{sixHundredSeventyThree\_pack} & CosTraceNormSixHundredSeventyThree & PROVED & \checkmark & verified \\
\texttt{degree\_sixHundredSixtyOne} & CosTraceNormSixHundredSixtyOne & PROVED & \checkmark & verified \\
\texttt{isIntegral\_and\_degree\_sixHundredSixtyOne} & CosTraceNormSixHundredSixtyOne & PROVED & \checkmark & verified \\
\texttt{isIntegral\_spectralGen\_sixHundredSixtyOne} & CosTraceNormSixHundredSixtyOne & PROVED & \checkmark & verified \\
\texttt{isIntegral\_spectralGen\_sixHundredSixtyOne\_Q} & CosTraceNormSixHundredSixtyOne & PROVED & \checkmark & verified \\
\texttt{prime\_sixHundredSixtyOne} & CosTraceNormSixHundredSixtyOne & PROVED & \checkmark & verified \\
\texttt{sixHundredSixtyOne\_ne\_two} & CosTraceNormSixHundredSixtyOne & PROVED & \checkmark & verified \\
\texttt{sixHundredSixtyOne\_pack} & CosTraceNormSixHundredSixtyOne & PROVED & \checkmark & verified \\
\texttt{degree\_sixHundredThirteen} & CosTraceNormSixHundredThirteen & PROVED & \checkmark & verified \\
\texttt{isIntegral\_and\_degree\_sixHundredThirteen} & CosTraceNormSixHundredThirteen & PROVED & \checkmark & verified \\
\texttt{isIntegral\_spectralGen\_sixHundredThirteen} & CosTraceNormSixHundredThirteen & PROVED & \checkmark & verified \\
\texttt{isIntegral\_spectralGen\_sixHundredThirteen\_Q} & CosTraceNormSixHundredThirteen & PROVED & \checkmark & verified \\
\texttt{prime\_sixHundredThirteen} & CosTraceNormSixHundredThirteen & PROVED & \checkmark & verified \\
\texttt{sixHundredThirteen\_ne\_two} & CosTraceNormSixHundredThirteen & PROVED & \checkmark & verified \\
\texttt{sixHundredThirteen\_pack} & CosTraceNormSixHundredThirteen & PROVED & \checkmark & verified \\
\texttt{degree\_sixHundredThirtyOne} & CosTraceNormSixHundredThirtyOne & PROVED & \checkmark & verified \\
\texttt{isIntegral\_and\_degree\_sixHundredThirtyOne} & CosTraceNormSixHundredThirtyOne & PROVED & \checkmark & verified \\
\texttt{isIntegral\_spectralGen\_sixHundredThirtyOne} & CosTraceNormSixHundredThirtyOne & PROVED & \checkmark & verified \\
\texttt{isIntegral\_spectralGen\_sixHundredThirtyOne\_Q} & CosTraceNormSixHundredThirtyOne & PROVED & \checkmark & verified \\
\texttt{prime\_sixHundredThirtyOne} & CosTraceNormSixHundredThirtyOne & PROVED & \checkmark & verified \\
\texttt{sixHundredThirtyOne\_ne\_two} & CosTraceNormSixHundredThirtyOne & PROVED & \checkmark & verified \\
\texttt{sixHundredThirtyOne\_pack} & CosTraceNormSixHundredThirtyOne & PROVED & \checkmark & verified \\
\texttt{degree\_sixtyOne} & CosTraceNormSixtyOne & PROVED & \checkmark & verified \\
\texttt{isIntegral\_and\_degree\_sixtyOne} & CosTraceNormSixtyOne & PROVED & \checkmark & verified \\
\texttt{isIntegral\_spectralGen\_sixtyOne} & CosTraceNormSixtyOne & PROVED & \checkmark & verified \\
\texttt{isIntegral\_spectralGen\_sixtyOne\_Q} & CosTraceNormSixtyOne & PROVED & \checkmark & verified \\
\texttt{prime\_sixtyOne} & CosTraceNormSixtyOne & PROVED & \checkmark & verified \\
\texttt{sixtyOne\_ne\_two} & CosTraceNormSixtyOne & PROVED & \checkmark & verified \\
\texttt{sixtyOne\_pack} & CosTraceNormSixtyOne & PROVED & \checkmark & verified \\
\texttt{degree\_sixtySeven} & CosTraceNormSixtySeven & PROVED & \checkmark & verified \\
\texttt{isIntegral\_and\_degree\_sixtySeven} & CosTraceNormSixtySeven & PROVED & \checkmark & verified \\
\texttt{isIntegral\_spectralGen\_sixtySeven} & CosTraceNormSixtySeven & PROVED & \checkmark & verified \\
\texttt{isIntegral\_spectralGen\_sixtySeven\_Q} & CosTraceNormSixtySeven & PROVED & \checkmark & verified \\
\texttt{prime\_sixtySeven} & CosTraceNormSixtySeven & PROVED & \checkmark & verified \\
\texttt{sixtySeven\_ne\_two} & CosTraceNormSixtySeven & PROVED & \checkmark & verified \\
\texttt{sixtySeven\_pack} & CosTraceNormSixtySeven & PROVED & \checkmark & verified \\
\texttt{degree\_thirteen} & CosTraceNormThirteen & PROVED & \checkmark & verified \\
\texttt{isIntegral\_and\_degree\_thirteen} & CosTraceNormThirteen & PROVED & \checkmark & verified \\
\texttt{isIntegral\_spectralGen\_thirteen} & CosTraceNormThirteen & PROVED & \checkmark & verified \\
\texttt{isIntegral\_spectralGen\_thirteen\_Q} & CosTraceNormThirteen & PROVED & \checkmark & verified \\
\texttt{prime\_thirteen} & CosTraceNormThirteen & PROVED & \checkmark & verified \\
\texttt{thirteen\_ne\_two} & CosTraceNormThirteen & PROVED & \checkmark & verified \\
\texttt{thirteen\_pack} & CosTraceNormThirteen & PROVED & \checkmark & verified \\
\texttt{degree\_thirtyOne} & CosTraceNormThirtyOne & PROVED & \checkmark & verified \\
\texttt{isIntegral\_and\_degree\_thirtyOne} & CosTraceNormThirtyOne & PROVED & \checkmark & verified \\
\texttt{isIntegral\_spectralGen\_thirtyOne} & CosTraceNormThirtyOne & PROVED & \checkmark & verified \\
\texttt{isIntegral\_spectralGen\_thirtyOne\_Q} & CosTraceNormThirtyOne & PROVED & \checkmark & verified \\
\texttt{prime\_thirtyOne} & CosTraceNormThirtyOne & PROVED & \checkmark & verified \\
\texttt{thirtyOne\_ne\_two} & CosTraceNormThirtyOne & PROVED & \checkmark & verified \\
\texttt{thirtyOne\_pack} & CosTraceNormThirtyOne & PROVED & \checkmark & verified \\
\texttt{degree\_thirtySeven} & CosTraceNormThirtySeven & PROVED & \checkmark & verified \\
\texttt{isIntegral\_and\_degree\_thirtySeven} & CosTraceNormThirtySeven & PROVED & \checkmark & verified \\
\texttt{isIntegral\_spectralGen\_thirtySeven} & CosTraceNormThirtySeven & PROVED & \checkmark & verified \\
\texttt{isIntegral\_spectralGen\_thirtySeven\_Q} & CosTraceNormThirtySeven & PROVED & \checkmark & verified \\
\texttt{prime\_thirtySeven} & CosTraceNormThirtySeven & PROVED & \checkmark & verified \\
\texttt{thirtySeven\_ne\_two} & CosTraceNormThirtySeven & PROVED & \checkmark & verified \\
\texttt{thirtySeven\_pack} & CosTraceNormThirtySeven & PROVED & \checkmark & verified \\
\texttt{degree\_threeHundredEightyNine} & CosTraceNormThreeHundredEightyNine & PROVED & \checkmark & verified \\
\texttt{isIntegral\_and\_degree\_threeHundredEightyNine} & CosTraceNormThreeHundredEightyNine & PROVED & \checkmark & verified \\
\texttt{isIntegral\_spectralGen\_threeHundredEightyNine} & CosTraceNormThreeHundredEightyNine & PROVED & \checkmark & verified \\
\texttt{isIntegral\_spectralGen\_threeHundredEightyNine\_Q} & CosTraceNormThreeHundredEightyNine & PROVED & \checkmark & verified \\
\texttt{prime\_threeHundredEightyNine} & CosTraceNormThreeHundredEightyNine & PROVED & \checkmark & verified \\
\texttt{threeHundredEightyNine\_ne\_two} & CosTraceNormThreeHundredEightyNine & PROVED & \checkmark & verified \\
\texttt{threeHundredEightyNine\_pack} & CosTraceNormThreeHundredEightyNine & PROVED & \checkmark & verified \\
\texttt{degree\_threeHundredEightyThree} & CosTraceNormThreeHundredEightyThree & PROVED & \checkmark & verified \\
\texttt{isIntegral\_and\_degree\_threeHundredEightyThree} & CosTraceNormThreeHundredEightyThree & PROVED & \checkmark & verified \\
\texttt{isIntegral\_spectralGen\_threeHundredEightyThree} & CosTraceNormThreeHundredEightyThree & PROVED & \checkmark & verified \\
\texttt{isIntegral\_spectralGen\_threeHundredEightyThree\_Q} & CosTraceNormThreeHundredEightyThree & PROVED & \checkmark & verified \\
\texttt{prime\_threeHundredEightyThree} & CosTraceNormThreeHundredEightyThree & PROVED & \checkmark & verified \\
\texttt{threeHundredEightyThree\_ne\_two} & CosTraceNormThreeHundredEightyThree & PROVED & \checkmark & verified \\
\texttt{threeHundredEightyThree\_pack} & CosTraceNormThreeHundredEightyThree & PROVED & \checkmark & verified \\
\texttt{degree\_threeHundredEleven} & CosTraceNormThreeHundredEleven & PROVED & \checkmark & verified \\
\texttt{isIntegral\_and\_degree\_threeHundredEleven} & CosTraceNormThreeHundredEleven & PROVED & \checkmark & verified \\
\texttt{isIntegral\_spectralGen\_threeHundredEleven} & CosTraceNormThreeHundredEleven & PROVED & \checkmark & verified \\
\texttt{isIntegral\_spectralGen\_threeHundredEleven\_Q} & CosTraceNormThreeHundredEleven & PROVED & \checkmark & verified \\
\texttt{prime\_threeHundredEleven} & CosTraceNormThreeHundredEleven & PROVED & \checkmark & verified \\
\texttt{threeHundredEleven\_ne\_two} & CosTraceNormThreeHundredEleven & PROVED & \checkmark & verified \\
\texttt{threeHundredEleven\_pack} & CosTraceNormThreeHundredEleven & PROVED & \checkmark & verified \\
\texttt{degree\_threeHundredFiftyNine} & CosTraceNormThreeHundredFiftyNine & PROVED & \checkmark & verified \\
\texttt{isIntegral\_and\_degree\_threeHundredFiftyNine} & CosTraceNormThreeHundredFiftyNine & PROVED & \checkmark & verified \\
\texttt{isIntegral\_spectralGen\_threeHundredFiftyNine} & CosTraceNormThreeHundredFiftyNine & PROVED & \checkmark & verified \\
\texttt{isIntegral\_spectralGen\_threeHundredFiftyNine\_Q} & CosTraceNormThreeHundredFiftyNine & PROVED & \checkmark & verified \\
\texttt{prime\_threeHundredFiftyNine} & CosTraceNormThreeHundredFiftyNine & PROVED & \checkmark & verified \\
\texttt{threeHundredFiftyNine\_ne\_two} & CosTraceNormThreeHundredFiftyNine & PROVED & \checkmark & verified \\
\texttt{threeHundredFiftyNine\_pack} & CosTraceNormThreeHundredFiftyNine & PROVED & \checkmark & verified \\
\texttt{degree\_threeHundredFiftyThree} & CosTraceNormThreeHundredFiftyThree & PROVED & \checkmark & verified \\
\texttt{isIntegral\_and\_degree\_threeHundredFiftyThree} & CosTraceNormThreeHundredFiftyThree & PROVED & \checkmark & verified \\
\texttt{isIntegral\_spectralGen\_threeHundredFiftyThree} & CosTraceNormThreeHundredFiftyThree & PROVED & \checkmark & verified \\
\texttt{isIntegral\_spectralGen\_threeHundredFiftyThree\_Q} & CosTraceNormThreeHundredFiftyThree & PROVED & \checkmark & verified \\
\texttt{prime\_threeHundredFiftyThree} & CosTraceNormThreeHundredFiftyThree & PROVED & \checkmark & verified \\
\texttt{threeHundredFiftyThree\_ne\_two} & CosTraceNormThreeHundredFiftyThree & PROVED & \checkmark & verified \\
\texttt{threeHundredFiftyThree\_pack} & CosTraceNormThreeHundredFiftyThree & PROVED & \checkmark & verified \\
\texttt{degree\_threeHundredFortyNine} & CosTraceNormThreeHundredFortyNine & PROVED & \checkmark & verified \\
\texttt{isIntegral\_and\_degree\_threeHundredFortyNine} & CosTraceNormThreeHundredFortyNine & PROVED & \checkmark & verified \\
\texttt{isIntegral\_spectralGen\_threeHundredFortyNine} & CosTraceNormThreeHundredFortyNine & PROVED & \checkmark & verified \\
\texttt{isIntegral\_spectralGen\_threeHundredFortyNine\_Q} & CosTraceNormThreeHundredFortyNine & PROVED & \checkmark & verified \\
\texttt{prime\_threeHundredFortyNine} & CosTraceNormThreeHundredFortyNine & PROVED & \checkmark & verified \\
\texttt{threeHundredFortyNine\_ne\_two} & CosTraceNormThreeHundredFortyNine & PROVED & \checkmark & verified \\
\texttt{threeHundredFortyNine\_pack} & CosTraceNormThreeHundredFortyNine & PROVED & \checkmark & verified \\
\texttt{degree\_threeHundredFortySeven} & CosTraceNormThreeHundredFortySeven & PROVED & \checkmark & verified \\
\texttt{isIntegral\_and\_degree\_threeHundredFortySeven} & CosTraceNormThreeHundredFortySeven & PROVED & \checkmark & verified \\
\texttt{isIntegral\_spectralGen\_threeHundredFortySeven} & CosTraceNormThreeHundredFortySeven & PROVED & \checkmark & verified \\
\texttt{isIntegral\_spectralGen\_threeHundredFortySeven\_Q} & CosTraceNormThreeHundredFortySeven & PROVED & \checkmark & verified \\
\texttt{prime\_threeHundredFortySeven} & CosTraceNormThreeHundredFortySeven & PROVED & \checkmark & verified \\
\texttt{threeHundredFortySeven\_ne\_two} & CosTraceNormThreeHundredFortySeven & PROVED & \checkmark & verified \\
\texttt{threeHundredFortySeven\_pack} & CosTraceNormThreeHundredFortySeven & PROVED & \checkmark & verified \\
\texttt{degree\_threeHundredNinetySeven} & CosTraceNormThreeHundredNinetySeven & PROVED & \checkmark & verified \\
\texttt{isIntegral\_and\_degree\_threeHundredNinetySeven} & CosTraceNormThreeHundredNinetySeven & PROVED & \checkmark & verified \\
\texttt{isIntegral\_spectralGen\_threeHundredNinetySeven} & CosTraceNormThreeHundredNinetySeven & PROVED & \checkmark & verified \\
\texttt{isIntegral\_spectralGen\_threeHundredNinetySeven\_Q} & CosTraceNormThreeHundredNinetySeven & PROVED & \checkmark & verified \\
\texttt{prime\_threeHundredNinetySeven} & CosTraceNormThreeHundredNinetySeven & PROVED & \checkmark & verified \\
\texttt{threeHundredNinetySeven\_ne\_two} & CosTraceNormThreeHundredNinetySeven & PROVED & \checkmark & verified \\
\texttt{threeHundredNinetySeven\_pack} & CosTraceNormThreeHundredNinetySeven & PROVED & \checkmark & verified \\
\texttt{degree\_threeHundredSeven} & CosTraceNormThreeHundredSeven & PROVED & \checkmark & verified \\
\texttt{isIntegral\_and\_degree\_threeHundredSeven} & CosTraceNormThreeHundredSeven & PROVED & \checkmark & verified \\
\texttt{isIntegral\_spectralGen\_threeHundredSeven} & CosTraceNormThreeHundredSeven & PROVED & \checkmark & verified \\
\texttt{isIntegral\_spectralGen\_threeHundredSeven\_Q} & CosTraceNormThreeHundredSeven & PROVED & \checkmark & verified \\
\texttt{prime\_threeHundredSeven} & CosTraceNormThreeHundredSeven & PROVED & \checkmark & verified \\
\texttt{threeHundredSeven\_ne\_two} & CosTraceNormThreeHundredSeven & PROVED & \checkmark & verified \\
\texttt{threeHundredSeven\_pack} & CosTraceNormThreeHundredSeven & PROVED & \checkmark & verified \\
\texttt{degree\_threeHundredSeventeen} & CosTraceNormThreeHundredSeventeen & PROVED & \checkmark & verified \\
\texttt{isIntegral\_and\_degree\_threeHundredSeventeen} & CosTraceNormThreeHundredSeventeen & PROVED & \checkmark & verified \\
\texttt{isIntegral\_spectralGen\_threeHundredSeventeen} & CosTraceNormThreeHundredSeventeen & PROVED & \checkmark & verified \\
\texttt{isIntegral\_spectralGen\_threeHundredSeventeen\_Q} & CosTraceNormThreeHundredSeventeen & PROVED & \checkmark & verified \\
\texttt{prime\_threeHundredSeventeen} & CosTraceNormThreeHundredSeventeen & PROVED & \checkmark & verified \\
\texttt{threeHundredSeventeen\_ne\_two} & CosTraceNormThreeHundredSeventeen & PROVED & \checkmark & verified \\
\texttt{threeHundredSeventeen\_pack} & CosTraceNormThreeHundredSeventeen & PROVED & \checkmark & verified \\
\texttt{degree\_threeHundredSeventyNine} & CosTraceNormThreeHundredSeventyNine & PROVED & \checkmark & verified \\
\texttt{isIntegral\_and\_degree\_threeHundredSeventyNine} & CosTraceNormThreeHundredSeventyNine & PROVED & \checkmark & verified \\
\texttt{isIntegral\_spectralGen\_threeHundredSeventyNine} & CosTraceNormThreeHundredSeventyNine & PROVED & \checkmark & verified \\
\texttt{isIntegral\_spectralGen\_threeHundredSeventyNine\_Q} & CosTraceNormThreeHundredSeventyNine & PROVED & \checkmark & verified \\
\texttt{prime\_threeHundredSeventyNine} & CosTraceNormThreeHundredSeventyNine & PROVED & \checkmark & verified \\
\texttt{threeHundredSeventyNine\_ne\_two} & CosTraceNormThreeHundredSeventyNine & PROVED & \checkmark & verified \\
\texttt{threeHundredSeventyNine\_pack} & CosTraceNormThreeHundredSeventyNine & PROVED & \checkmark & verified \\
\texttt{degree\_threeHundredSeventyThree} & CosTraceNormThreeHundredSeventyThree & PROVED & \checkmark & verified \\
\texttt{isIntegral\_and\_degree\_threeHundredSeventyThree} & CosTraceNormThreeHundredSeventyThree & PROVED & \checkmark & verified \\
\texttt{isIntegral\_spectralGen\_threeHundredSeventyThree} & CosTraceNormThreeHundredSeventyThree & PROVED & \checkmark & verified \\
\texttt{isIntegral\_spectralGen\_threeHundredSeventyThree\_Q} & CosTraceNormThreeHundredSeventyThree & PROVED & \checkmark & verified \\
\texttt{prime\_threeHundredSeventyThree} & CosTraceNormThreeHundredSeventyThree & PROVED & \checkmark & verified \\
\texttt{threeHundredSeventyThree\_ne\_two} & CosTraceNormThreeHundredSeventyThree & PROVED & \checkmark & verified \\
\texttt{threeHundredSeventyThree\_pack} & CosTraceNormThreeHundredSeventyThree & PROVED & \checkmark & verified \\
\texttt{degree\_threeHundredSixtySeven} & CosTraceNormThreeHundredSixtySeven & PROVED & \checkmark & verified \\
\texttt{isIntegral\_and\_degree\_threeHundredSixtySeven} & CosTraceNormThreeHundredSixtySeven & PROVED & \checkmark & verified \\
\texttt{isIntegral\_spectralGen\_threeHundredSixtySeven} & CosTraceNormThreeHundredSixtySeven & PROVED & \checkmark & verified \\
\texttt{isIntegral\_spectralGen\_threeHundredSixtySeven\_Q} & CosTraceNormThreeHundredSixtySeven & PROVED & \checkmark & verified \\
\texttt{prime\_threeHundredSixtySeven} & CosTraceNormThreeHundredSixtySeven & PROVED & \checkmark & verified \\
\texttt{threeHundredSixtySeven\_ne\_two} & CosTraceNormThreeHundredSixtySeven & PROVED & \checkmark & verified \\
\texttt{threeHundredSixtySeven\_pack} & CosTraceNormThreeHundredSixtySeven & PROVED & \checkmark & verified \\
\texttt{degree\_threeHundredThirteen} & CosTraceNormThreeHundredThirteen & PROVED & \checkmark & verified \\
\texttt{isIntegral\_and\_degree\_threeHundredThirteen} & CosTraceNormThreeHundredThirteen & PROVED & \checkmark & verified \\
\texttt{isIntegral\_spectralGen\_threeHundredThirteen} & CosTraceNormThreeHundredThirteen & PROVED & \checkmark & verified \\
\texttt{isIntegral\_spectralGen\_threeHundredThirteen\_Q} & CosTraceNormThreeHundredThirteen & PROVED & \checkmark & verified \\
\texttt{prime\_threeHundredThirteen} & CosTraceNormThreeHundredThirteen & PROVED & \checkmark & verified \\
\texttt{threeHundredThirteen\_ne\_two} & CosTraceNormThreeHundredThirteen & PROVED & \checkmark & verified \\
\texttt{threeHundredThirteen\_pack} & CosTraceNormThreeHundredThirteen & PROVED & \checkmark & verified \\
\texttt{degree\_threeHundredThirtyOne} & CosTraceNormThreeHundredThirtyOne & PROVED & \checkmark & verified \\
\texttt{isIntegral\_and\_degree\_threeHundredThirtyOne} & CosTraceNormThreeHundredThirtyOne & PROVED & \checkmark & verified \\
\texttt{isIntegral\_spectralGen\_threeHundredThirtyOne} & CosTraceNormThreeHundredThirtyOne & PROVED & \checkmark & verified \\
\texttt{isIntegral\_spectralGen\_threeHundredThirtyOne\_Q} & CosTraceNormThreeHundredThirtyOne & PROVED & \checkmark & verified \\
\texttt{prime\_threeHundredThirtyOne} & CosTraceNormThreeHundredThirtyOne & PROVED & \checkmark & verified \\
\texttt{threeHundredThirtyOne\_ne\_two} & CosTraceNormThreeHundredThirtyOne & PROVED & \checkmark & verified \\
\texttt{threeHundredThirtyOne\_pack} & CosTraceNormThreeHundredThirtyOne & PROVED & \checkmark & verified \\
\texttt{degree\_threeHundredThirtySeven} & CosTraceNormThreeHundredThirtySeven & PROVED & \checkmark & verified \\
\texttt{isIntegral\_and\_degree\_threeHundredThirtySeven} & CosTraceNormThreeHundredThirtySeven & PROVED & \checkmark & verified \\
\texttt{isIntegral\_spectralGen\_threeHundredThirtySeven} & CosTraceNormThreeHundredThirtySeven & PROVED & \checkmark & verified \\
\texttt{isIntegral\_spectralGen\_threeHundredThirtySeven\_Q} & CosTraceNormThreeHundredThirtySeven & PROVED & \checkmark & verified \\
\texttt{prime\_threeHundredThirtySeven} & CosTraceNormThreeHundredThirtySeven & PROVED & \checkmark & verified \\
\texttt{threeHundredThirtySeven\_ne\_two} & CosTraceNormThreeHundredThirtySeven & PROVED & \checkmark & verified \\
\texttt{threeHundredThirtySeven\_pack} & CosTraceNormThreeHundredThirtySeven & PROVED & \checkmark & verified \\
\texttt{degree\_twentyNine} & CosTraceNormTwentyNine & PROVED & \checkmark & verified \\
\texttt{isIntegral\_and\_degree\_twentyNine} & CosTraceNormTwentyNine & PROVED & \checkmark & verified \\
\texttt{isIntegral\_spectralGen\_twentyNine} & CosTraceNormTwentyNine & PROVED & \checkmark & verified \\
\texttt{isIntegral\_spectralGen\_twentyNine\_Q} & CosTraceNormTwentyNine & PROVED & \checkmark & verified \\
\texttt{prime\_twentyNine} & CosTraceNormTwentyNine & PROVED & \checkmark & verified \\
\texttt{twentyNine\_ne\_two} & CosTraceNormTwentyNine & PROVED & \checkmark & verified \\
\texttt{twentyNine\_pack} & CosTraceNormTwentyNine & PROVED & \checkmark & verified \\
\texttt{degree\_twentyThree} & CosTraceNormTwentyThree & PROVED & \checkmark & verified \\
\texttt{isIntegral\_and\_degree\_twentyThree} & CosTraceNormTwentyThree & PROVED & \checkmark & verified \\
\texttt{isIntegral\_spectralGen\_twentyThree} & CosTraceNormTwentyThree & PROVED & \checkmark & verified \\
\texttt{isIntegral\_spectralGen\_twentyThree\_Q} & CosTraceNormTwentyThree & PROVED & \checkmark & verified \\
\texttt{prime\_twentyThree} & CosTraceNormTwentyThree & PROVED & \checkmark & verified \\
\texttt{twentyThree\_ne\_two} & CosTraceNormTwentyThree & PROVED & \checkmark & verified \\
\texttt{twentyThree\_pack} & CosTraceNormTwentyThree & PROVED & \checkmark & verified \\
\texttt{degree\_twoHundredEightyOne} & CosTraceNormTwoHundredEightyOne & PROVED & \checkmark & verified \\
\texttt{isIntegral\_and\_degree\_twoHundredEightyOne} & CosTraceNormTwoHundredEightyOne & PROVED & \checkmark & verified \\
\texttt{isIntegral\_spectralGen\_twoHundredEightyOne} & CosTraceNormTwoHundredEightyOne & PROVED & \checkmark & verified \\
\texttt{isIntegral\_spectralGen\_twoHundredEightyOne\_Q} & CosTraceNormTwoHundredEightyOne & PROVED & \checkmark & verified \\
\texttt{prime\_twoHundredEightyOne} & CosTraceNormTwoHundredEightyOne & PROVED & \checkmark & verified \\
\texttt{twoHundredEightyOne\_ne\_two} & CosTraceNormTwoHundredEightyOne & PROVED & \checkmark & verified \\
\texttt{twoHundredEightyOne\_pack} & CosTraceNormTwoHundredEightyOne & PROVED & \checkmark & verified \\
\texttt{degree\_twoHundredEightyThree} & CosTraceNormTwoHundredEightyThree & PROVED & \checkmark & verified \\
\texttt{isIntegral\_and\_degree\_twoHundredEightyThree} & CosTraceNormTwoHundredEightyThree & PROVED & \checkmark & verified \\
\texttt{isIntegral\_spectralGen\_twoHundredEightyThree} & CosTraceNormTwoHundredEightyThree & PROVED & \checkmark & verified \\
\texttt{isIntegral\_spectralGen\_twoHundredEightyThree\_Q} & CosTraceNormTwoHundredEightyThree & PROVED & \checkmark & verified \\
\texttt{prime\_twoHundredEightyThree} & CosTraceNormTwoHundredEightyThree & PROVED & \checkmark & verified \\
\texttt{twoHundredEightyThree\_ne\_two} & CosTraceNormTwoHundredEightyThree & PROVED & \checkmark & verified \\
\texttt{twoHundredEightyThree\_pack} & CosTraceNormTwoHundredEightyThree & PROVED & \checkmark & verified \\
\texttt{degree\_twoHundredEleven} & CosTraceNormTwoHundredEleven & PROVED & \checkmark & verified \\
\texttt{isIntegral\_and\_degree\_twoHundredEleven} & CosTraceNormTwoHundredEleven & PROVED & \checkmark & verified \\
\texttt{isIntegral\_spectralGen\_twoHundredEleven} & CosTraceNormTwoHundredEleven & PROVED & \checkmark & verified \\
\texttt{isIntegral\_spectralGen\_twoHundredEleven\_Q} & CosTraceNormTwoHundredEleven & PROVED & \checkmark & verified \\
\texttt{prime\_twoHundredEleven} & CosTraceNormTwoHundredEleven & PROVED & \checkmark & verified \\
\texttt{twoHundredEleven\_ne\_two} & CosTraceNormTwoHundredEleven & PROVED & \checkmark & verified \\
\texttt{twoHundredEleven\_pack} & CosTraceNormTwoHundredEleven & PROVED & \checkmark & verified \\
\texttt{degree\_twoHundredFiftyOne} & CosTraceNormTwoHundredFiftyOne & PROVED & \checkmark & verified \\
\texttt{isIntegral\_and\_degree\_twoHundredFiftyOne} & CosTraceNormTwoHundredFiftyOne & PROVED & \checkmark & verified \\
\texttt{isIntegral\_spectralGen\_twoHundredFiftyOne} & CosTraceNormTwoHundredFiftyOne & PROVED & \checkmark & verified \\
\texttt{isIntegral\_spectralGen\_twoHundredFiftyOne\_Q} & CosTraceNormTwoHundredFiftyOne & PROVED & \checkmark & verified \\
\texttt{prime\_twoHundredFiftyOne} & CosTraceNormTwoHundredFiftyOne & PROVED & \checkmark & verified \\
\texttt{twoHundredFiftyOne\_ne\_two} & CosTraceNormTwoHundredFiftyOne & PROVED & \checkmark & verified \\
\texttt{twoHundredFiftyOne\_pack} & CosTraceNormTwoHundredFiftyOne & PROVED & \checkmark & verified \\
\texttt{degree\_twoHundredFiftySeven} & CosTraceNormTwoHundredFiftySeven & PROVED & \checkmark & verified \\
\texttt{isIntegral\_and\_degree\_twoHundredFiftySeven} & CosTraceNormTwoHundredFiftySeven & PROVED & \checkmark & verified \\
\texttt{isIntegral\_spectralGen\_twoHundredFiftySeven} & CosTraceNormTwoHundredFiftySeven & PROVED & \checkmark & verified \\
\texttt{isIntegral\_spectralGen\_twoHundredFiftySeven\_Q} & CosTraceNormTwoHundredFiftySeven & PROVED & \checkmark & verified \\
\texttt{prime\_twoHundredFiftySeven} & CosTraceNormTwoHundredFiftySeven & PROVED & \checkmark & verified \\
\texttt{twoHundredFiftySeven\_ne\_two} & CosTraceNormTwoHundredFiftySeven & PROVED & \checkmark & verified \\
\texttt{twoHundredFiftySeven\_pack} & CosTraceNormTwoHundredFiftySeven & PROVED & \checkmark & verified \\
\texttt{degree\_twoHundredFortyOne} & CosTraceNormTwoHundredFortyOne & PROVED & \checkmark & verified \\
\texttt{isIntegral\_and\_degree\_twoHundredFortyOne} & CosTraceNormTwoHundredFortyOne & PROVED & \checkmark & verified \\
\texttt{isIntegral\_spectralGen\_twoHundredFortyOne} & CosTraceNormTwoHundredFortyOne & PROVED & \checkmark & verified \\
\texttt{isIntegral\_spectralGen\_twoHundredFortyOne\_Q} & CosTraceNormTwoHundredFortyOne & PROVED & \checkmark & verified \\
\texttt{prime\_twoHundredFortyOne} & CosTraceNormTwoHundredFortyOne & PROVED & \checkmark & verified \\
\texttt{twoHundredFortyOne\_ne\_two} & CosTraceNormTwoHundredFortyOne & PROVED & \checkmark & verified \\
\texttt{twoHundredFortyOne\_pack} & CosTraceNormTwoHundredFortyOne & PROVED & \checkmark & verified \\
\texttt{degree\_twoHundredNinetyThree} & CosTraceNormTwoHundredNinetyThree & PROVED & \checkmark & verified \\
\texttt{isIntegral\_and\_degree\_twoHundredNinetyThree} & CosTraceNormTwoHundredNinetyThree & PROVED & \checkmark & verified \\
\texttt{isIntegral\_spectralGen\_twoHundredNinetyThree} & CosTraceNormTwoHundredNinetyThree & PROVED & \checkmark & verified \\
\texttt{isIntegral\_spectralGen\_twoHundredNinetyThree\_Q} & CosTraceNormTwoHundredNinetyThree & PROVED & \checkmark & verified \\
\texttt{prime\_twoHundredNinetyThree} & CosTraceNormTwoHundredNinetyThree & PROVED & \checkmark & verified \\
\texttt{twoHundredNinetyThree\_ne\_two} & CosTraceNormTwoHundredNinetyThree & PROVED & \checkmark & verified \\
\texttt{twoHundredNinetyThree\_pack} & CosTraceNormTwoHundredNinetyThree & PROVED & \checkmark & verified \\
\texttt{degree\_twoHundredSeventyOne} & CosTraceNormTwoHundredSeventyOne & PROVED & \checkmark & verified \\
\texttt{isIntegral\_and\_degree\_twoHundredSeventyOne} & CosTraceNormTwoHundredSeventyOne & PROVED & \checkmark & verified \\
\texttt{isIntegral\_spectralGen\_twoHundredSeventyOne} & CosTraceNormTwoHundredSeventyOne & PROVED & \checkmark & verified \\
\texttt{isIntegral\_spectralGen\_twoHundredSeventyOne\_Q} & CosTraceNormTwoHundredSeventyOne & PROVED & \checkmark & verified \\
\texttt{prime\_twoHundredSeventyOne} & CosTraceNormTwoHundredSeventyOne & PROVED & \checkmark & verified \\
\texttt{twoHundredSeventyOne\_ne\_two} & CosTraceNormTwoHundredSeventyOne & PROVED & \checkmark & verified \\
\texttt{twoHundredSeventyOne\_pack} & CosTraceNormTwoHundredSeventyOne & PROVED & \checkmark & verified \\
\texttt{degree\_twoHundredSeventySeven} & CosTraceNormTwoHundredSeventySeven & PROVED & \checkmark & verified \\
\texttt{isIntegral\_and\_degree\_twoHundredSeventySeven} & CosTraceNormTwoHundredSeventySeven & PROVED & \checkmark & verified \\
\texttt{isIntegral\_spectralGen\_twoHundredSeventySeven} & CosTraceNormTwoHundredSeventySeven & PROVED & \checkmark & verified \\
\texttt{isIntegral\_spectralGen\_twoHundredSeventySeven\_Q} & CosTraceNormTwoHundredSeventySeven & PROVED & \checkmark & verified \\
\texttt{prime\_twoHundredSeventySeven} & CosTraceNormTwoHundredSeventySeven & PROVED & \checkmark & verified \\
\texttt{twoHundredSeventySeven\_ne\_two} & CosTraceNormTwoHundredSeventySeven & PROVED & \checkmark & verified \\
\texttt{twoHundredSeventySeven\_pack} & CosTraceNormTwoHundredSeventySeven & PROVED & \checkmark & verified \\
\texttt{degree\_twoHundredSixtyNine} & CosTraceNormTwoHundredSixtyNine & PROVED & \checkmark & verified \\
\texttt{isIntegral\_and\_degree\_twoHundredSixtyNine} & CosTraceNormTwoHundredSixtyNine & PROVED & \checkmark & verified \\
\texttt{isIntegral\_spectralGen\_twoHundredSixtyNine} & CosTraceNormTwoHundredSixtyNine & PROVED & \checkmark & verified \\
\texttt{isIntegral\_spectralGen\_twoHundredSixtyNine\_Q} & CosTraceNormTwoHundredSixtyNine & PROVED & \checkmark & verified \\
\texttt{prime\_twoHundredSixtyNine} & CosTraceNormTwoHundredSixtyNine & PROVED & \checkmark & verified \\
\texttt{twoHundredSixtyNine\_ne\_two} & CosTraceNormTwoHundredSixtyNine & PROVED & \checkmark & verified \\
\texttt{twoHundredSixtyNine\_pack} & CosTraceNormTwoHundredSixtyNine & PROVED & \checkmark & verified \\
\texttt{degree\_twoHundredSixtyThree} & CosTraceNormTwoHundredSixtyThree & PROVED & \checkmark & verified \\
\texttt{isIntegral\_and\_degree\_twoHundredSixtyThree} & CosTraceNormTwoHundredSixtyThree & PROVED & \checkmark & verified \\
\texttt{isIntegral\_spectralGen\_twoHundredSixtyThree} & CosTraceNormTwoHundredSixtyThree & PROVED & \checkmark & verified \\
\texttt{isIntegral\_spectralGen\_twoHundredSixtyThree\_Q} & CosTraceNormTwoHundredSixtyThree & PROVED & \checkmark & verified \\
\texttt{prime\_twoHundredSixtyThree} & CosTraceNormTwoHundredSixtyThree & PROVED & \checkmark & verified \\
\texttt{twoHundredSixtyThree\_ne\_two} & CosTraceNormTwoHundredSixtyThree & PROVED & \checkmark & verified \\
\texttt{twoHundredSixtyThree\_pack} & CosTraceNormTwoHundredSixtyThree & PROVED & \checkmark & verified \\
\texttt{degree\_twoHundredThirtyNine} & CosTraceNormTwoHundredThirtyNine & PROVED & \checkmark & verified \\
\texttt{isIntegral\_and\_degree\_twoHundredThirtyNine} & CosTraceNormTwoHundredThirtyNine & PROVED & \checkmark & verified \\
\texttt{isIntegral\_spectralGen\_twoHundredThirtyNine} & CosTraceNormTwoHundredThirtyNine & PROVED & \checkmark & verified \\
\texttt{isIntegral\_spectralGen\_twoHundredThirtyNine\_Q} & CosTraceNormTwoHundredThirtyNine & PROVED & \checkmark & verified \\
\texttt{prime\_twoHundredThirtyNine} & CosTraceNormTwoHundredThirtyNine & PROVED & \checkmark & verified \\
\texttt{twoHundredThirtyNine\_ne\_two} & CosTraceNormTwoHundredThirtyNine & PROVED & \checkmark & verified \\
\texttt{twoHundredThirtyNine\_pack} & CosTraceNormTwoHundredThirtyNine & PROVED & \checkmark & verified \\
\texttt{degree\_twoHundredThirtyThree} & CosTraceNormTwoHundredThirtyThree & PROVED & \checkmark & verified \\
\texttt{isIntegral\_and\_degree\_twoHundredThirtyThree} & CosTraceNormTwoHundredThirtyThree & PROVED & \checkmark & verified \\
\texttt{isIntegral\_spectralGen\_twoHundredThirtyThree} & CosTraceNormTwoHundredThirtyThree & PROVED & \checkmark & verified \\
\texttt{isIntegral\_spectralGen\_twoHundredThirtyThree\_Q} & CosTraceNormTwoHundredThirtyThree & PROVED & \checkmark & verified \\
\texttt{prime\_twoHundredThirtyThree} & CosTraceNormTwoHundredThirtyThree & PROVED & \checkmark & verified \\
\texttt{twoHundredThirtyThree\_ne\_two} & CosTraceNormTwoHundredThirtyThree & PROVED & \checkmark & verified \\
\texttt{twoHundredThirtyThree\_pack} & CosTraceNormTwoHundredThirtyThree & PROVED & \checkmark & verified \\
\texttt{degree\_twoHundredTwentyNine} & CosTraceNormTwoHundredTwentyNine & PROVED & \checkmark & verified \\
\texttt{isIntegral\_and\_degree\_twoHundredTwentyNine} & CosTraceNormTwoHundredTwentyNine & PROVED & \checkmark & verified \\
\texttt{isIntegral\_spectralGen\_twoHundredTwentyNine} & CosTraceNormTwoHundredTwentyNine & PROVED & \checkmark & verified \\
\texttt{isIntegral\_spectralGen\_twoHundredTwentyNine\_Q} & CosTraceNormTwoHundredTwentyNine & PROVED & \checkmark & verified \\
\texttt{prime\_twoHundredTwentyNine} & CosTraceNormTwoHundredTwentyNine & PROVED & \checkmark & verified \\
\texttt{twoHundredTwentyNine\_ne\_two} & CosTraceNormTwoHundredTwentyNine & PROVED & \checkmark & verified \\
\texttt{twoHundredTwentyNine\_pack} & CosTraceNormTwoHundredTwentyNine & PROVED & \checkmark & verified \\
\texttt{degree\_twoHundredTwentySeven} & CosTraceNormTwoHundredTwentySeven & PROVED & \checkmark & verified \\
\texttt{isIntegral\_and\_degree\_twoHundredTwentySeven} & CosTraceNormTwoHundredTwentySeven & PROVED & \checkmark & verified \\
\texttt{isIntegral\_spectralGen\_twoHundredTwentySeven} & CosTraceNormTwoHundredTwentySeven & PROVED & \checkmark & verified \\
\texttt{isIntegral\_spectralGen\_twoHundredTwentySeven\_Q} & CosTraceNormTwoHundredTwentySeven & PROVED & \checkmark & verified \\
\texttt{prime\_twoHundredTwentySeven} & CosTraceNormTwoHundredTwentySeven & PROVED & \checkmark & verified \\
\texttt{twoHundredTwentySeven\_ne\_two} & CosTraceNormTwoHundredTwentySeven & PROVED & \checkmark & verified \\
\texttt{twoHundredTwentySeven\_pack} & CosTraceNormTwoHundredTwentySeven & PROVED & \checkmark & verified \\
\texttt{degree\_twoHundredTwentyThree} & CosTraceNormTwoHundredTwentyThree & PROVED & \checkmark & verified \\
\texttt{isIntegral\_and\_degree\_twoHundredTwentyThree} & CosTraceNormTwoHundredTwentyThree & PROVED & \checkmark & verified \\
\texttt{isIntegral\_spectralGen\_twoHundredTwentyThree} & CosTraceNormTwoHundredTwentyThree & PROVED & \checkmark & verified \\
\texttt{isIntegral\_spectralGen\_twoHundredTwentyThree\_Q} & CosTraceNormTwoHundredTwentyThree & PROVED & \checkmark & verified \\
\texttt{prime\_twoHundredTwentyThree} & CosTraceNormTwoHundredTwentyThree & PROVED & \checkmark & verified \\
\texttt{twoHundredTwentyThree\_ne\_two} & CosTraceNormTwoHundredTwentyThree & PROVED & \checkmark & verified \\
\texttt{twoHundredTwentyThree\_pack} & CosTraceNormTwoHundredTwentyThree & PROVED & \checkmark & verified \\
\texttt{cullen\_1\_prime} & CullenWoodall & PROVED & \checkmark & verified \\
\texttt{cullen\_2\_not\_prime} & CullenWoodall & PROVED & \checkmark & verified \\
\texttt{cullen\_3\_not\_prime} & CullenWoodall & PROVED & \checkmark & verified \\
\texttt{cullen\_odd} & CullenWoodall & PROVED & \checkmark & verified \\
\texttt{cullen\_sub\_woodall} & CullenWoodall & PROVED & \checkmark & verified \\
\texttt{woodall\_2\_prime} & CullenWoodall & PROVED & \checkmark & verified \\
\texttt{woodall\_3\_prime} & CullenWoodall & PROVED & \checkmark & verified \\
\texttt{woodall\_4\_not\_prime} & CullenWoodall & PROVED & \checkmark & verified \\
\texttt{woodall\_5\_not\_prime} & CullenWoodall & PROVED & \checkmark & verified \\
\texttt{woodall\_6\_prime} & CullenWoodall & PROVED & \checkmark & verified \\
\texttt{woodall\_odd} & CullenWoodall & PROVED & \checkmark & verified \\
\texttt{algebraic\_connectivity\_five} & CycleSpectrumFamily & PROVED & \checkmark & verified \\
\texttt{algebraic\_connectivity\_five\_props} & CycleSpectrumFamily & PROVED & \checkmark & verified \\
\texttt{algebraic\_connectivity\_le\_four} & CycleSpectrumFamily & PROVED & \checkmark & verified \\
\texttt{algebraic\_connectivity\_mem} & CycleSpectrumFamily & PROVED & \checkmark & verified \\
\texttt{algebraic\_connectivity\_pos} & CycleSpectrumFamily & PROVED & \checkmark & verified \\
\texttt{cos\_four\_pi\_div\_three} & CycleSpectrumFamily & PROVED & \checkmark & verified \\
\texttt{cos\_two\_pi\_div\_three} & CycleSpectrumFamily & PROVED & \checkmark & verified \\
\texttt{cycle3\_eig\_one} & CycleSpectrumFamily & PROVED & \checkmark & verified \\
\texttt{cycle3\_eig\_two} & CycleSpectrumFamily & PROVED & \checkmark & verified \\
\texttt{cycle3\_eig\_zero} & CycleSpectrumFamily & PROVED & \checkmark & verified \\
\texttt{cycle4\_eig\_one} & CycleSpectrumFamily & PROVED & \checkmark & verified \\
\texttt{cycle4\_eig\_three} & CycleSpectrumFamily & PROVED & \checkmark & verified \\
\texttt{cycle4\_eig\_two} & CycleSpectrumFamily & PROVED & \checkmark & verified \\
\texttt{cycle4\_eig\_zero} & CycleSpectrumFamily & PROVED & \checkmark & verified \\
\texttt{cycle6\_eig\_five} & CycleSpectrumFamily & PROVED & \checkmark & verified \\
\texttt{cycle6\_eig\_four} & CycleSpectrumFamily & PROVED & \checkmark & verified \\
\texttt{cycle6\_eig\_one} & CycleSpectrumFamily & PROVED & \checkmark & verified \\
\texttt{cycle6\_eig\_three} & CycleSpectrumFamily & PROVED & \checkmark & verified \\
\texttt{cycle6\_eig\_two} & CycleSpectrumFamily & PROVED & \checkmark & verified \\
\texttt{cycle6\_eig\_zero} & CycleSpectrumFamily & PROVED & \checkmark & verified \\
\texttt{cycleSpectrum\_ge\_neg\_two} & CycleSpectrumFamily & PROVED & \checkmark & verified \\
\texttt{cycleSpectrum\_le\_two} & CycleSpectrumFamily & PROVED & \checkmark & verified \\
\texttt{cycleSpectrum\_subset\_Icc} & CycleSpectrumFamily & PROVED & \checkmark & verified \\
\texttt{cycle\_eig\_one} & CycleSpectrumFamily & PROVED & \checkmark & verified \\
\texttt{golden\_in\_C5} & CycleSpectrumFamily & PROVED & \checkmark & verified \\
\texttt{golden\_unique\_among\_prime\_cycles} & CycleSpectrumFamily & PROVED & \checkmark & verified \\
\texttt{lambda\_max\_cycle} & CycleSpectrumFamily & PROVED & \checkmark & verified \\
\texttt{laplacianCycleSpectrum\_le\_four} & CycleSpectrumFamily & PROVED & \checkmark & verified \\
\texttt{laplacianCycleSpectrum\_nonneg} & CycleSpectrumFamily & PROVED & \checkmark & verified \\
\texttt{mem\_cycleSpectrum} & CycleSpectrumFamily & PROVED & \checkmark & verified \\
\texttt{mem\_laplacianCycleSpectrum} & CycleSpectrumFamily & PROVED & \checkmark & verified \\
\texttt{neg\_golden\_in\_C5} & CycleSpectrumFamily & PROVED & \checkmark & verified \\
\texttt{two\_cos\_pi\_div\_three} & CycleSpectrumFamily & PROVED & \checkmark & verified \\
\texttt{two\_cos\_two\_pi\_div\_five\_eq\_golden\_sub\_one} & CycleSpectrumFamily & PROVED & \checkmark & verified \\
\texttt{two\_cos\_two\_pi\_div\_three} & CycleSpectrumFamily & PROVED & \checkmark & verified \\
\texttt{two\_mem\_cycleSpectrum} & CycleSpectrumFamily & PROVED & \checkmark & verified \\
\texttt{zero\_mem\_laplacianCycleSpectrum} & CycleSpectrumFamily & PROVED & \checkmark & verified \\
\texttt{realSubfield\_facts\_general} & CyclotomicGaloisGroup & PROVED & \checkmark & verified \\
\texttt{realSubfield\_gal\_card} & CyclotomicGaloisGroup & PROVED & \checkmark & verified \\
\texttt{realSubfield\_gal\_equivUnitsQuotient} & CyclotomicGaloisGroup & PROVED & \checkmark & verified \\
\texttt{realSubfield\_gal\_isCyclic\_of\_prime} & CyclotomicGaloisGroup & PROVED & \checkmark & verified \\
\texttt{realSubfield\_gal\_isQuotientOfUnits} & CyclotomicGaloisGroup & PROVED & \checkmark & verified \\
\texttt{realSubfield\_gal\_units\_presentation} & CyclotomicGaloisGroup & PROVED & \checkmark & verified \\
\texttt{realSubfield\_isAbelianGalois} & CyclotomicGaloisGroup & PROVED & \checkmark & verified \\
\texttt{realSubfield\_isGalois} & CyclotomicGaloisGroup & PROVED & \checkmark & verified \\
\texttt{pentagon\_quadratic} & CyclotomicRealDegree & PROVED & \checkmark & verified \\
\texttt{quadratic\_iff\_mem} & CyclotomicRealDegree & PROVED & \checkmark & verified \\
\texttt{quadratic\_iff\_totient\_four} & CyclotomicRealDegree & PROVED & \checkmark & verified \\
\texttt{spectral\_degree\_general} & CyclotomicRealDegree & PROVED & \checkmark & verified \\
\texttt{spectral\_natDegree\_two\_mul} & CyclotomicRealDegree & PROVED & \checkmark & verified \\
\texttt{totient\_eq\_four\_iff} & CyclotomicRealDegree & PROVED & \checkmark & verified \\
\texttt{charInner\_eq\_pairSum} & D5CharacterComplete & PROVED & \checkmark & verified \\
\texttt{chiConjugate\_one} & D5CharacterComplete & PROVED & \checkmark & verified \\
\texttt{chiConjugate\_real} & D5CharacterComplete & PROVED & \checkmark & verified \\
\texttt{chiGolden\_one} & D5CharacterComplete & PROVED & \checkmark & verified \\
\texttt{chiGolden\_real} & D5CharacterComplete & PROVED & \checkmark & verified \\
\texttt{chiSign\_one} & D5CharacterComplete & PROVED & \checkmark & verified \\
\texttt{chiSign\_real} & D5CharacterComplete & PROVED & \checkmark & verified \\
\texttt{chiTrivial\_one} & D5CharacterComplete & PROVED & \checkmark & verified \\
\texttt{chiTrivial\_real} & D5CharacterComplete & PROVED & \checkmark & verified \\
\texttt{colInner\_one\_one} & D5CharacterComplete & PROVED & \checkmark & verified \\
\texttt{colInner\_one\_r1} & D5CharacterComplete & PROVED & \checkmark & verified \\
\texttt{colInner\_one\_r2} & D5CharacterComplete & PROVED & \checkmark & verified \\
\texttt{colInner\_one\_sr0} & D5CharacterComplete & PROVED & \checkmark & verified \\
\texttt{colInner\_r1\_r1} & D5CharacterComplete & PROVED & \checkmark & verified \\
\texttt{colInner\_r1\_r2} & D5CharacterComplete & PROVED & \checkmark & verified \\
\texttt{colInner\_r1\_sr0} & D5CharacterComplete & PROVED & \checkmark & verified \\
\texttt{colInner\_r2\_r2} & D5CharacterComplete & PROVED & \checkmark & verified \\
\texttt{colInner\_r2\_sr0} & D5CharacterComplete & PROVED & \checkmark & verified \\
\texttt{colInner\_sr0\_sr0} & D5CharacterComplete & PROVED & \checkmark & verified \\
\texttt{col\_orthogonal} & D5CharacterComplete & PROVED & \checkmark & verified \\
\texttt{conjRot\_col} & D5CharacterComplete & PROVED & \checkmark & verified \\
\texttt{conjRot\_one} & D5CharacterComplete & PROVED & \checkmark & verified \\
\texttt{conjRot\_real} & D5CharacterComplete & PROVED & \checkmark & verified \\
\texttt{conjRot\_two} & D5CharacterComplete & PROVED & \checkmark & verified \\
\texttt{conjRot\_zero} & D5CharacterComplete & PROVED & \checkmark & verified \\
\texttt{dimension\_identity} & D5CharacterComplete & PROVED & \checkmark & verified \\
\texttt{dimension\_identity\_card} & D5CharacterComplete & PROVED & \checkmark & verified \\
\texttt{goldenRatio\_sq\_complex} & D5CharacterComplete & PROVED & \checkmark & verified \\
\texttt{goldenRot\_col} & D5CharacterComplete & PROVED & \checkmark & verified \\
\texttt{goldenRot\_eq\_adjEigenvalue} & D5CharacterComplete & PROVED & \checkmark & verified \\
\texttt{goldenRot\_one} & D5CharacterComplete & PROVED & \checkmark & verified \\
\texttt{goldenRot\_real} & D5CharacterComplete & PROVED & \checkmark & verified \\
\texttt{goldenRot\_two} & D5CharacterComplete & PROVED & \checkmark & verified \\
\texttt{goldenRot\_zero} & D5CharacterComplete & PROVED & \checkmark & verified \\
\texttt{golden\_char\_eq\_two\_cos} & D5CharacterComplete & PROVED & \checkmark & verified \\
\texttt{golden\_char\_rotation\_class} & D5CharacterComplete & PROVED & \checkmark & verified \\
\texttt{omegaPow\_col\_mul} & D5CharacterComplete & PROVED & \checkmark & verified \\
\texttt{one\_eq\_r0} & D5CharacterComplete & PROVED & \checkmark & verified \\
\texttt{pairSum\_split} & D5CharacterComplete & PROVED & \checkmark & verified \\
\texttt{rot\_CC} & D5CharacterComplete & PROVED & \checkmark & verified \\
\texttt{rot\_C\_sum} & D5CharacterComplete & PROVED & \checkmark & verified \\
\texttt{rot\_GC} & D5CharacterComplete & PROVED & \checkmark & verified \\
\texttt{rot\_GG} & D5CharacterComplete & PROVED & \checkmark & verified \\
\texttt{rot\_G\_sum} & D5CharacterComplete & PROVED & \checkmark & verified \\
\texttt{row\_CC} & D5CharacterComplete & PROVED & \checkmark & verified \\
\texttt{row\_GC} & D5CharacterComplete & PROVED & \checkmark & verified \\
\texttt{row\_GG} & D5CharacterComplete & PROVED & \checkmark & verified \\
\texttt{row\_SC} & D5CharacterComplete & PROVED & \checkmark & verified \\
\texttt{row\_SG} & D5CharacterComplete & PROVED & \checkmark & verified \\
\texttt{row\_SS} & D5CharacterComplete & PROVED & \checkmark & verified \\
\texttt{row\_TC} & D5CharacterComplete & PROVED & \checkmark & verified \\
\texttt{row\_TG} & D5CharacterComplete & PROVED & \checkmark & verified \\
\texttt{row\_TS} & D5CharacterComplete & PROVED & \checkmark & verified \\
\texttt{row\_TT} & D5CharacterComplete & PROVED & \checkmark & verified \\
\texttt{row\_orthonormal} & D5CharacterComplete & PROVED & \checkmark & verified \\
\texttt{star\_omega} & D5CharacterComplete & PROVED & \checkmark & verified \\
\texttt{star\_omegaPow} & D5CharacterComplete & PROVED & \checkmark & verified \\
\texttt{sum\_dihedral5} & D5CharacterComplete & PROVED & \checkmark & verified \\
\texttt{sum\_omega\_binom} & D5CharacterComplete & PROVED & \checkmark & verified \\
\texttt{sum\_omega\_binom\_prod} & D5CharacterComplete & PROVED & \checkmark & verified \\
\texttt{d5Character\_eq\_sum\_fixed} & D5CharacterTable & PROVED & \checkmark & verified \\
\texttt{d5Character\_one} & D5CharacterTable & PROVED & \checkmark & verified \\
\texttt{d5Character\_reflection} & D5CharacterTable & PROVED & \checkmark & verified \\
\texttt{d5Character\_rotation} & D5CharacterTable & PROVED & \checkmark & verified \\
\texttt{d5Character\_rotation\_ne\_zero} & D5CharacterTable & PROVED & \checkmark & verified \\
\texttt{d5PermutationMatrix\_apply} & D5CharacterTable & PROVED & \checkmark & verified \\
\texttt{d5PermutationMatrix\_mulVec} & D5CharacterTable & PROVED & \checkmark & verified \\
\texttt{fourierCoeff\_add} & D5FourierInversion & PROVED & \checkmark & verified \\
\texttt{fourierCoeff\_eigenmode} & D5FourierInversion & PROVED & \checkmark & verified \\
\texttt{fourierCoeff\_smul} & D5FourierInversion & PROVED & \checkmark & verified \\
\texttt{fourier\_inversion} & D5FourierInversion & PROVED & \checkmark & verified \\
\texttt{isotypicProjector\_add} & D5FourierInversion & PROVED & \checkmark & verified \\
\texttt{isotypicProjector\_eq\_fourierCoeff\_smul} & D5FourierInversion & PROVED & \checkmark & verified \\
\texttt{isotypicProjector\_idempotent} & D5FourierInversion & PROVED & \checkmark & verified \\
\texttt{isotypicProjector\_orthogonal} & D5FourierInversion & PROVED & \checkmark & verified \\
\texttt{sum\_isotypicProjectors} & D5FourierInversion & PROVED & \checkmark & verified \\
\texttt{character\_orthogonality} & D5Isotypic & PROVED & \checkmark & verified \\
\texttt{coordSum\_eigenmode} & D5Isotypic & PROVED & \checkmark & verified \\
\texttt{d5Pull\_r\_apply} & D5Isotypic & PROVED & \checkmark & verified \\
\texttt{d5Pull\_r\_eigenmode} & D5Isotypic & PROVED & \checkmark & verified \\
\texttt{d5Pull\_r\_mul} & D5Isotypic & PROVED & \checkmark & verified \\
\texttt{d5Pull\_r\_one\_apply} & D5Isotypic & PROVED & \checkmark & verified \\
\texttt{d5Pull\_r\_one\_eigenmode} & D5Isotypic & PROVED & \checkmark & verified \\
\texttt{eigenmode\_mem\_zeroSumSubmodule} & D5Isotypic & PROVED & \checkmark & verified \\
\texttt{eigenmode\_zero} & D5Isotypic & PROVED & \checkmark & verified \\
\texttt{isotypicProjector\_apply} & D5Isotypic & PROVED & \checkmark & verified \\
\texttt{isotypicProjector\_eigenmode} & D5Isotypic & PROVED & \checkmark & verified \\
\texttt{isotypicProjector\_eigenmode\_of\_ne} & D5Isotypic & PROVED & \checkmark & verified \\
\texttt{isotypicProjector\_eigenmode\_self} & D5Isotypic & PROVED & \checkmark & verified \\
\texttt{isotypicProjector\_idempotent\_eigenmode} & D5Isotypic & PROVED & \checkmark & verified \\
\texttt{isotypicProjector\_idempotent\_self} & D5Isotypic & PROVED & \checkmark & verified \\
\texttt{isotypicProjector\_orthogonal\_eigenmode} & D5Isotypic & PROVED & \checkmark & verified \\
\texttt{isotypicProjector\_smul} & D5Isotypic & PROVED & \checkmark & verified \\
\texttt{omegaPow\_add} & D5Isotypic & PROVED & \checkmark & verified \\
\texttt{omegaPow\_neg} & D5Isotypic & PROVED & \checkmark & verified \\
\texttt{omega\_isPrimitiveRoot} & D5Isotypic & PROVED & \checkmark & verified \\
\texttt{omega\_pow\_eq\_one\_iff} & D5Isotypic & PROVED & \checkmark & verified \\
\texttt{omega\_pow\_five} & D5Isotypic & PROVED & \checkmark & verified \\
\texttt{omega\_pow\_modEq} & D5Isotypic & PROVED & \checkmark & verified \\
\texttt{orderOf\_omega} & D5Isotypic & PROVED & \checkmark & verified \\
\texttt{sum\_omegaPow} & D5Isotypic & PROVED & \checkmark & verified \\
\texttt{sum\_omegaPow\_ne\_zero} & D5Isotypic & PROVED & \checkmark & verified \\
\texttt{adjacencyLinear\_apply} & D5LaplacianModes & PROVED & \checkmark & verified \\
\texttt{adjacency\_add} & D5LaplacianModes & PROVED & \checkmark & verified \\
\texttt{adjacency\_apply} & D5LaplacianModes & PROVED & \checkmark & verified \\
\texttt{adjacency\_commute\_isotypicProjector\_eigenmode} & D5LaplacianModes & PROVED & \checkmark & verified \\
\texttt{adjacency\_eigenmode} & D5LaplacianModes & PROVED & \checkmark & verified \\
\texttt{adjacency\_eigenmode\_cos} & D5LaplacianModes & PROVED & \checkmark & verified \\
\texttt{adjacency\_eigenmode\_zero} & D5LaplacianModes & PROVED & \checkmark & verified \\
\texttt{adjacency\_eq\_pullbacks} & D5LaplacianModes & PROVED & \checkmark & verified \\
\texttt{adjacency\_isotypicProjector\_eigenmode} & D5LaplacianModes & PROVED & \checkmark & verified \\
\texttt{adjacency\_smul} & D5LaplacianModes & PROVED & \checkmark & verified \\
\texttt{laplacianLinear\_apply} & D5LaplacianModes & PROVED & \checkmark & verified \\
\texttt{laplacian\_add} & D5LaplacianModes & PROVED & \checkmark & verified \\
\texttt{laplacian\_apply} & D5LaplacianModes & PROVED & \checkmark & verified \\
\texttt{laplacian\_commute\_isotypicProjector\_eigenmode} & D5LaplacianModes & PROVED & \checkmark & verified \\
\texttt{laplacian\_eigenmode} & D5LaplacianModes & PROVED & \checkmark & verified \\
\texttt{laplacian\_eigenmode\_zero} & D5LaplacianModes & PROVED & \checkmark & verified \\
\texttt{laplacian\_eq} & D5LaplacianModes & PROVED & \checkmark & verified \\
\texttt{laplacian\_isotypicProjector\_eigenmode} & D5LaplacianModes & PROVED & \checkmark & verified \\
\texttt{laplacian\_smul} & D5LaplacianModes & PROVED & \checkmark & verified \\
\texttt{omegaPow\_add\_inv\_eq\_two\_cos} & D5LaplacianModes & PROVED & \checkmark & verified \\
\texttt{omegaPow\_eq\_exp} & D5LaplacianModes & PROVED & \checkmark & verified \\
\texttt{omegaPow\_eq\_exp\_mul\_I} & D5LaplacianModes & PROVED & \checkmark & verified \\
\texttt{autPull\_apply} & D5Representation & PROVED & \checkmark & verified \\
\texttt{autPull\_constant} & D5Representation & PROVED & \checkmark & verified \\
\texttt{autPull\_mem\_constantLine} & D5Representation & PROVED & \checkmark & verified \\
\texttt{autPull\_mem\_zeroSumSubmodule} & D5Representation & PROVED & \checkmark & verified \\
\texttt{constantLinear\_apply} & D5Representation & PROVED & \checkmark & verified \\
\texttt{constantVector\_mem\_zeroSumSubmodule\_iff} & D5Representation & PROVED & \checkmark & verified \\
\texttt{coordSumLinear\_apply} & D5Representation & PROVED & \checkmark & verified \\
\texttt{coordSum\_autPull} & D5Representation & PROVED & \checkmark & verified \\
\texttt{coordSum\_constantVector} & D5Representation & PROVED & \checkmark & verified \\
\texttt{coordSum\_d5Pull} & D5Representation & PROVED & \checkmark & verified \\
\texttt{d5Pull\_apply} & D5Representation & PROVED & \checkmark & verified \\
\texttt{d5Pull\_constant} & D5Representation & PROVED & \checkmark & verified \\
\texttt{d5Pull\_mem\_constantLine} & D5Representation & PROVED & \checkmark & verified \\
\texttt{d5Pull\_mem\_zeroSumSubmodule} & D5Representation & PROVED & \checkmark & verified \\
\texttt{asymptotic\_shape\_consistent} & Equidistribution & PROVED & \checkmark & verified \\
\texttt{configCount\_twelve\_five\_two\_one} & Equidistribution & PROVED & \checkmark & verified \\
\texttt{configCount\_twenty\_five\_two\_two} & Equidistribution & PROVED & \checkmark & verified \\
\texttt{equidistribution\_of\_asymptotic} & Equidistribution & CONDITIONAL & \checkmark & verified \\
\texttt{equidistribution\_of\_asymptotic\_exists} & Equidistribution & CONDITIONAL & \checkmark & verified \\
\texttt{prime\_pair\_config\_admissible} & Equidistribution & PROVED & \checkmark & verified \\
\texttt{pairCount\_deviation\_le\_scaled\_err\_of\_large\_pairs} & DeviationBound & PROVED & \checkmark & verified \\
\texttt{pairCount\_normalized\_deviation\_le\_scaled\_epsilon\_of\_large\_pairs} & DeviationBound & PROVED & \checkmark & verified \\
\texttt{perConfig\_deviation\_le\_err} & DeviationBound & PROVED & \checkmark & verified \\
\texttt{perConfig\_normalized\_deviation\_le\_epsilon} & DeviationBound & PROVED & \checkmark & verified \\
\texttt{perConfig\_normalized\_deviation\_le\_one\_third} & DeviationBound & PROVED & \checkmark & verified \\
\texttt{totalConfig\_deviation\_le\_scaled\_err} & DeviationBound & PROVED & \checkmark & verified \\
\texttt{totalConfig\_normalized\_deviation\_le\_one\_third\_scaled} & DeviationBound & PROVED & \checkmark & verified \\
\texttt{totalConfig\_normalized\_deviation\_le\_scaled\_epsilon} & DeviationBound & PROVED & \checkmark & verified \\
\texttt{configCount\_deviation\_bound} & FiniteScaffold & PROVED & \checkmark & verified \\
\texttt{configCount\_eq\_zero\_of\_not\_admissible\_of\_large\_pairs} & FiniteScaffold & PROVED & \checkmark & verified \\
\texttt{configCount\_le\_window} & FiniteScaffold & PROVED & \checkmark & verified \\
\texttt{mem\_pairStarts} & FiniteScaffold & PROVED & \checkmark & verified \\
\texttt{sum\_configCount\_univ\_eq\_pairCount} & FiniteScaffold & PROVED & \checkmark & verified \\
\texttt{totalConfigCount\_deviation\_bound} & FiniteScaffold & PROVED & \checkmark & verified \\
\texttt{totalConfigCount\_eq\_pairCount\_of\_large\_pairs} & FiniteScaffold & PROVED & \checkmark & verified \\
\texttt{totalConfigCount\_le\_admissible\_card\_mul\_window} & FiniteScaffold & PROVED & \checkmark & verified \\
\texttt{totalConfigCount\_le\_q\_sub\_two\_mul\_window} & FiniteScaffold & PROVED & \checkmark & verified \\
\texttt{admissible\_reflection\_symmetry} & EquidistributionBVReduction & PROVED & \checkmark & verified \\
\texttt{bv\_shape\_consistent} & EquidistributionBVReduction & PROVED & \checkmark & verified \\
\texttt{configCount\_density\_of\_BV} & EquidistributionBVReduction & CONDITIONAL & \checkmark & verified \\
\texttt{configCount\_over\_main\_tendsto} & EquidistributionBVReduction & CONDITIONAL & \checkmark & verified \\
\texttt{equidistribution\_of\_BV\_uniform} & EquidistributionBVReduction & CONDITIONAL & \checkmark & verified \\
\texttt{total\_over\_main\_tendsto} & EquidistributionBVReduction & CONDITIONAL & \checkmark & verified \\
\texttt{equidistribution\_of\_transitive\_symmetry} & EquidistributionUniformity & CONDITIONAL & \checkmark & verified \\
\texttt{equidistribution\_three} & EquidistributionUniformity & PROVED & \checkmark & verified \\
\texttt{iterate\_mem\_admissible} & EquidistributionUniformity & PROVED & \checkmark & verified \\
\texttt{reflect\_affine} & EquidistributionUniformity & PROVED & \checkmark & verified \\
\texttt{reflect\_five\_fixes\_four} & EquidistributionUniformity & PROVED & \checkmark & verified \\
\texttt{reflect\_five\_four\_orbit} & EquidistributionUniformity & PROVED & \checkmark & verified \\
\texttt{reflect\_five\_swaps\_one} & EquidistributionUniformity & PROVED & \checkmark & verified \\
\texttt{reflect\_five\_swaps\_two} & EquidistributionUniformity & PROVED & \checkmark & verified \\
\texttt{reflect\_involutive} & EquidistributionUniformity & PROVED & \checkmark & verified \\
\texttt{reflect\_preservesAdmissible} & EquidistributionUniformity & PROVED & \checkmark & verified \\
\texttt{reflect\_preserves\_admissible} & EquidistributionUniformity & PROVED & \checkmark & verified \\
\texttt{reflection\_not\_transitive\_five} & EquidistributionUniformity & PROVED & \checkmark & verified \\
\texttt{sing\_iterate} & EquidistributionUniformity & PROVED & \checkmark & verified \\
\texttt{sing\_uniform\_of\_transitive} & EquidistributionUniformity & CONDITIONAL & \checkmark & verified \\
\texttt{sing\_uniform\_three} & EquidistributionUniformity & PROVED & \checkmark & verified \\
\texttt{affine\_stabilizer\_five\_classification} & EquidistributionUniformityClosure & PROVED & \checkmark & verified \\
\texttt{affine\_stabilizer\_five\_fixes\_four} & EquidistributionUniformityClosure & PROVED & \checkmark & verified \\
\texttt{affine\_stabilizer\_five\_four\_orbit} & EquidistributionUniformityClosure & PROVED & \checkmark & verified \\
\texttt{affine\_stabilizer\_five\_not\_transitive} & EquidistributionUniformityClosure & PROVED & \checkmark & verified \\
\texttt{f\_le\_log} & Erdos236 & PROVED & \checkmark & verified \\
\texttt{entropy\_le\_log\_card} & ErdosPinned & PROVED & \checkmark & verified \\
\texttt{exp\_neg\_entropy\_le\_sum\_sq} & ErdosPinned & PROVED & \checkmark & verified \\
\texttt{erdosStraus\_dvd\_three} & ErdosStraus & PROVED & \checkmark & verified \\
\texttt{erdosStraus\_even} & ErdosStraus & PROVED & \checkmark & verified \\
\texttt{erdosStraus\_mod3\_two} & ErdosStraus & PROVED & \checkmark & verified \\
\texttt{erdosStraus\_mod4\_three} & ErdosStraus & PROVED & \checkmark & verified \\
\texttt{erdosStraus\_of\_covered} & ErdosStraus & PROVED & \checkmark & verified \\
\texttt{erdosStraus\_of\_dvd} & ErdosStraus & PROVED & \checkmark & verified \\
\texttt{erdosStraus\_of\_prime\_case} & ErdosStraus & PROVED & \checkmark & verified \\
\texttt{erdosStraus\_covered} & ErdosStrausResidues & PROVED & \checkmark & verified \\
\texttt{erdosStraus\_covered\_ext} & ErdosStrausResidues & PROVED & \checkmark & verified \\
\texttt{erdosStraus\_dvd\_five} & ErdosStrausResidues & PROVED & \checkmark & verified \\
\texttt{erdosStraus\_dvd\_seven} & ErdosStrausResidues & PROVED & \checkmark & verified \\
\texttt{erdosStraus\_open\_frontier\_mod24} & ErdosStrausResidues & PROVED & \checkmark & verified \\
\texttt{erdosStraus\_open\_reduces} & ErdosStrausResidues & PROVED & \checkmark & verified \\
\texttt{erdosStraus\_open\_reduces\_mod12} & ErdosStrausResidues & PROVED & \checkmark & verified \\
\texttt{fermat\_0\_prime} & FermatNumbers & PROVED & \checkmark & verified \\
\texttt{fermat\_1\_prime} & FermatNumbers & PROVED & \checkmark & verified \\
\texttt{fermat\_2\_prime} & FermatNumbers & PROVED & \checkmark & verified \\
\texttt{fermat\_3\_prime} & FermatNumbers & PROVED & \checkmark & verified \\
\texttt{fermat\_4\_prime} & FermatNumbers & PROVED & \checkmark & verified \\
\texttt{fermat\_5\_eq} & FermatNumbers & PROVED & \checkmark & verified \\
\texttt{fermat\_5\_not\_prime} & FermatNumbers & PROVED & \checkmark & verified \\
\texttt{fermat\_6\_eq} & FermatNumbers & PROVED & \checkmark & verified \\
\texttt{fermat\_6\_not\_prime} & FermatNumbers & PROVED & \checkmark & verified \\
\texttt{fermat\_coprime} & FermatNumbers & PROVED & \checkmark & verified \\
\texttt{fermat\_eq\_fermatNumber} & FermatNumbers & PROVED & \checkmark & verified \\
\texttt{conj\_omega} & Fin5InnerProduct & PROVED & \checkmark & verified \\
\texttt{conj\_omegaPow} & Fin5InnerProduct & PROVED & \checkmark & verified \\
\texttt{conj\_omega\_pow} & Fin5InnerProduct & PROVED & \checkmark & verified \\
\texttt{eigenmode\_orthogonal} & Fin5InnerProduct & PROVED & \checkmark & verified \\
\texttt{hermInner\_add\_left} & Fin5InnerProduct & PROVED & \checkmark & verified \\
\texttt{hermInner\_add\_right} & Fin5InnerProduct & PROVED & \checkmark & verified \\
\texttt{hermInner\_conj\_symm} & Fin5InnerProduct & PROVED & \checkmark & verified \\
\texttt{hermInner\_eigenmode} & Fin5InnerProduct & PROVED & \checkmark & verified \\
\texttt{hermInner\_eigenmode\_self} & Fin5InnerProduct & PROVED & \checkmark & verified \\
\texttt{hermInner\_eigenmode\_zero} & Fin5InnerProduct & PROVED & \checkmark & verified \\
\texttt{hermInner\_self} & Fin5InnerProduct & PROVED & \checkmark & verified \\
\texttt{hermInner\_self\_eq\_zero\_iff} & Fin5InnerProduct & PROVED & \checkmark & verified \\
\texttt{hermInner\_smul\_left} & Fin5InnerProduct & PROVED & \checkmark & verified \\
\texttt{hermInner\_smul\_right} & Fin5InnerProduct & PROVED & \checkmark & verified \\
\texttt{norm\_omega} & Fin5InnerProduct & PROVED & \checkmark & verified \\
\texttt{fortunate\_2} & FortunateNumbers & PROVED & \checkmark & verified \\
\texttt{fortunate\_210} & FortunateNumbers & PROVED & \checkmark & verified \\
\texttt{fortunate\_30} & FortunateNumbers & PROVED & \checkmark & verified \\
\texttt{fortunate\_6} & FortunateNumbers & PROVED & \checkmark & verified \\
\texttt{fortunate\_values\_prime} & FortunateNumbers & PROVED & \checkmark & verified \\
\texttt{primorial\_2} & FortunateNumbers & PROVED & \checkmark & verified \\
\texttt{primorial\_3} & FortunateNumbers & PROVED & \checkmark & verified \\
\texttt{primorial\_5} & FortunateNumbers & PROVED & \checkmark & verified \\
\texttt{primorial\_7} & FortunateNumbers & PROVED & \checkmark & verified \\
\texttt{downMs\_upPart} & FranklinFixedPoint & PROVED & \checkmark & verified \\
\texttt{downOverlap\_stair} & FranklinFixedPoint & PROVED & \checkmark & verified \\
\texttt{downPart\_largest} & FranklinFixedPoint & PROVED & \checkmark & verified \\
\texttt{downPart\_notFixed} & FranklinFixedPoint & PROVED & \checkmark & verified \\
\texttt{downPart\_sPart\_gt} & FranklinFixedPoint & PROVED & \checkmark & verified \\
\texttt{downPart\_tDiag} & FranklinFixedPoint & PROVED & \checkmark & verified \\
\texttt{down\_L\_ge} & FranklinFixedPoint & PROVED & \checkmark & verified \\
\texttt{fixedPart\_mem\_iff} & FranklinFixedPoint & PROVED & \checkmark & verified \\
\texttt{fixedPart\_overlap} & FranklinFixedPoint & PROVED & \checkmark & verified \\
\texttt{hd\_of\_mem} & FranklinFixedPoint & PROVED & \checkmark & verified \\
\texttt{hu\_of\_mem} & FranklinFixedPoint & PROVED & \checkmark & verified \\
\texttt{largestPart\_eq\_of} & FranklinFixedPoint & PROVED & \checkmark & verified \\
\texttt{largestPart\_of\_stair\_neg} & FranklinFixedPoint & PROVED & \checkmark & verified \\
\texttt{largestPart\_of\_stair\_pos} & FranklinFixedPoint & PROVED & \checkmark & verified \\
\texttt{mem\_stair\_natCast} & FranklinFixedPoint & PROVED & \checkmark & verified \\
\texttt{mem\_stair\_neg\_natCast} & FranklinFixedPoint & PROVED & \checkmark & verified \\
\texttt{parts\_ne\_zero\_of\_mem} & FranklinFixedPoint & PROVED & \checkmark & verified \\
\texttt{pentagonalNumberTheorem} & FranklinFixedPoint & PROVED & \checkmark & verified \\
\texttt{phi\_notFixed} & FranklinFixedPoint & PROVED & \checkmark & verified \\
\texttt{phi\_parts} & FranklinFixedPoint & PROVED & \checkmark & verified \\
\texttt{phi\_phi\_eq} & FranklinFixedPoint & PROVED & \checkmark & verified \\
\texttt{sPart\_eq\_of} & FranklinFixedPoint & PROVED & \checkmark & verified \\
\texttt{sPart\_of\_stair\_neg} & FranklinFixedPoint & PROVED & \checkmark & verified \\
\texttt{sPart\_of\_stair\_pos} & FranklinFixedPoint & PROVED & \checkmark & verified \\
\texttt{stair\_mem\_iff} & FranklinFixedPoint & PROVED & \checkmark & verified \\
\texttt{tDiag\_of\_stair\_neg} & FranklinFixedPoint & PROVED & \checkmark & verified \\
\texttt{tDiag\_of\_stair\_pos} & FranklinFixedPoint & PROVED & \checkmark & verified \\
\texttt{upMs\_downPart} & FranklinFixedPoint & PROVED & \checkmark & verified \\
\texttt{upOverlap\_stair} & FranklinFixedPoint & PROVED & \checkmark & verified \\
\texttt{upPart\_largest} & FranklinFixedPoint & PROVED & \checkmark & verified \\
\texttt{upPart\_le\_branch} & FranklinFixedPoint & PROVED & \checkmark & verified \\
\texttt{upPart\_notFixed} & FranklinFixedPoint & PROVED & \checkmark & verified \\
\texttt{upPart\_sPart} & FranklinFixedPoint & PROVED & \checkmark & verified \\
\texttt{up\_L\_gt} & FranklinFixedPoint & PROVED & \checkmark & verified \\
\texttt{largestPart\_mem} & FranklinInvolution & PROVED & \checkmark & verified \\
\texttt{le\_largestPart} & FranklinInvolution & PROVED & \checkmark & verified \\
\texttt{mem\_of\_lt\_tDiag} & FranklinInvolution & PROVED & \checkmark & verified \\
\texttt{one\_le\_tDiag} & FranklinInvolution & PROVED & \checkmark & verified \\
\texttt{sPart\_le} & FranklinInvolution & PROVED & \checkmark & verified \\
\texttt{sPart\_mem} & FranklinInvolution & PROVED & \checkmark & verified \\
\texttt{signOf\_ne\_zero} & FranklinInvolution & PROVED & \checkmark & verified \\
\texttt{signedSum\_eq\_fixed\_of\_involution} & FranklinInvolution & PROVED & \checkmark & verified \\
\texttt{tDiag\_gap\_exists} & FranklinInvolution & PROVED & \checkmark & verified \\
\texttt{tDiag\_notMem} & FranklinInvolution & PROVED & \checkmark & verified \\
\texttt{fixedPart\_subset} & FranklinInvolutionProof & PROVED & \checkmark & verified \\
\texttt{fixedPart\_sum} & FranklinInvolutionProof & PROVED & \checkmark & verified \\
\texttt{franklin\_sum\_invariant\_down} & FranklinInvolutionProof & PROVED & \checkmark & verified \\
\texttt{franklin\_sum\_invariant\_up} & FranklinInvolutionProof & PROVED & \checkmark & verified \\
\texttt{gauss\_int} & FranklinInvolutionProof & PROVED & \checkmark & verified \\
\texttt{neg\_one\_pow\_natAbs} & FranklinInvolutionProof & PROVED & \checkmark & verified \\
\texttt{stair\_card} & FranklinInvolutionProof & PROVED & \checkmark & verified \\
\texttt{stair\_nodup} & FranklinInvolutionProof & PROVED & \checkmark & verified \\
\texttt{stair\_pos} & FranklinInvolutionProof & PROVED & \checkmark & verified \\
\texttt{stair\_sum} & FranklinInvolutionProof & PROVED & \checkmark & verified \\
\texttt{stair\_sum\_eq} & FranklinInvolutionProof & PROVED & \checkmark & verified \\
\texttt{d\_mem} & FranklinMapConstruction & PROVED & \checkmark & verified \\
\texttt{downMs\_card} & FranklinMapConstruction & PROVED & \checkmark & verified \\
\texttt{downMs\_nodup} & FranklinMapConstruction & PROVED & \checkmark & verified \\
\texttt{downPart\_ne} & FranklinMapConstruction & PROVED & \checkmark & verified \\
\texttt{downPart\_sign} & FranklinMapConstruction & PROVED & \checkmark & verified \\
\texttt{nodup\_of\_sdiff} & FranklinMapConstruction & PROVED & \checkmark & verified \\
\texttt{phiMem} & FranklinMapConstruction & PROVED & \checkmark & verified \\
\texttt{phi\_ne} & FranklinMapConstruction & PROVED & \checkmark & verified \\
\texttt{phi\_parts\_nodup} & FranklinMapConstruction & PROVED & \checkmark & verified \\
\texttt{phi\_sign} & FranklinMapConstruction & PROVED & \checkmark & verified \\
\texttt{sPart\_le\_largest} & FranklinMapConstruction & PROVED & \checkmark & verified \\
\texttt{sPart\_pos} & FranklinMapConstruction & PROVED & \checkmark & verified \\
\texttt{upMs\_card} & FranklinMapConstruction & PROVED & \checkmark & verified \\
\texttt{upMs\_nodup} & FranklinMapConstruction & PROVED & \checkmark & verified \\
\texttt{upPart\_ne} & FranklinMapConstruction & PROVED & \checkmark & verified \\
\texttt{upPart\_sign} & FranklinMapConstruction & PROVED & \checkmark & verified \\
\texttt{essentiallySelfAdjoint\_fourierConj} & FreeLaplacianPlancherel & PROVED & \checkmark & verified \\
\texttt{fourierL2\_inner\_map} & FreeLaplacianPlancherel & PROVED & \checkmark & verified \\
\texttt{fourierL2\_norm\_map} & FreeLaplacianPlancherel & PROVED & \checkmark & verified \\
\texttt{freeLaplacian\_essentiallySelfAdjoint\_via\_plancherel} & FreeLaplacianPlancherel & CONDITIONAL & \checkmark & verified \\
\texttt{cycExt} & GaloisCyclicGroup & PROVED & \checkmark & verified \\
\texttt{primRoot} & GaloisCyclicGroup & PROVED & \checkmark & verified \\
\texttt{realSubfield\_facts} & GaloisCyclicGroup & PROVED & \checkmark & verified \\
\texttt{realSubfield\_gal\_card} & GaloisCyclicGroup & PROVED & \checkmark & verified \\
\texttt{realSubfield\_gal\_isCyclic} & GaloisCyclicGroup & PROVED & \checkmark & verified \\
\texttt{realSubfield\_isGalois} & GaloisCyclicGroup & PROVED & \checkmark & verified \\
\texttt{quadratic\_iff\_five\_general} & GaloisGeneralDegree & PROVED & \checkmark & verified \\
\texttt{real\_subfield\_degree} & GaloisGeneralDegree & PROVED & \checkmark & verified \\
\texttt{C\_facts} & GaloisMinPolyFamily & PROVED & \checkmark & verified \\
\texttt{Psi\_eq\_minpoly} & GaloisMinPolyFamily & PROVED & \checkmark & verified \\
\texttt{Psi\_five} & GaloisMinPolyFamily & PROVED & \checkmark & verified \\
\texttt{Psi\_monic} & GaloisMinPolyFamily & PROVED & \checkmark & verified \\
\texttt{Psi\_natDegree} & GaloisMinPolyFamily & PROVED & \checkmark & verified \\
\texttt{Psi\_seven} & GaloisMinPolyFamily & PROVED & \checkmark & verified \\
\texttt{aeval\_spectralGen\_Psi} & GaloisMinPolyFamily & PROVED & \checkmark & verified \\
\texttt{minpoly\_five} & GaloisMinPolyFamily & PROVED & \checkmark & verified \\
\texttt{minpoly\_seven} & GaloisMinPolyFamily & PROVED & \checkmark & verified \\
\texttt{psiAux} & GaloisMinPolyFamily & PROVED & \checkmark & verified \\
\texttt{aeval\_spectralGen\_eight} & GaloisNgonClassification & PROVED & \checkmark & verified \\
\texttt{aeval\_spectralGen\_ten} & GaloisNgonClassification & PROVED & \checkmark & verified \\
\texttt{aeval\_spectralGen\_twelve} & GaloisNgonClassification & PROVED & \checkmark & verified \\
\texttt{disc\_eight} & GaloisNgonClassification & PROVED & \checkmark & verified \\
\texttt{disc\_five} & GaloisNgonClassification & PROVED & \checkmark & verified \\
\texttt{disc\_ten} & GaloisNgonClassification & PROVED & \checkmark & verified \\
\texttt{disc\_twelve} & GaloisNgonClassification & PROVED & \checkmark & verified \\
\texttt{golden\_ngons\_are\_five\_and\_ten} & GaloisNgonClassification & PROVED & \checkmark & verified \\
\texttt{quadratic\_ngon\_tfae} & GaloisNgonClassification & PROVED & \checkmark & verified \\
\texttt{spectralGen\_eight} & GaloisNgonClassification & PROVED & \checkmark & verified \\
\texttt{spectralGen\_ten} & GaloisNgonClassification & PROVED & \checkmark & verified \\
\texttt{spectralGen\_twelve} & GaloisNgonClassification & PROVED & \checkmark & verified \\
\texttt{P7\_irreducible} & GaloisWhyFive & PROVED & \checkmark & verified \\
\texttt{P7\_monic} & GaloisWhyFive & PROVED & \checkmark & verified \\
\texttt{P7\_natDegree} & GaloisWhyFive & PROVED & \checkmark & verified \\
\texttt{Q5\_monic} & GaloisWhyFive & PROVED & \checkmark & verified \\
\texttt{Q5\_natDegree} & GaloisWhyFive & PROVED & \checkmark & verified \\
\texttt{aeval\_spectralGen\_seven} & GaloisWhyFive & PROVED & \checkmark & verified \\
\texttt{cubic7\_monic} & GaloisWhyFive & PROVED & \checkmark & verified \\
\texttt{cubic\_identity\_seven} & GaloisWhyFive & PROVED & \checkmark & verified \\
\texttt{degree\_five} & GaloisWhyFive & PROVED & \checkmark & verified \\
\texttt{degree\_seven} & GaloisWhyFive & PROVED & \checkmark & verified \\
\texttt{degree\_three} & GaloisWhyFive & PROVED & \checkmark & verified \\
\texttt{goldenSubOne\_irrational} & GaloisWhyFive & PROVED & \checkmark & verified \\
\texttt{quadratic\_iff\_five} & GaloisWhyFive & PROVED & \checkmark & verified \\
\texttt{spectralGen\_five} & GaloisWhyFive & PROVED & \checkmark & verified \\
\texttt{spectralGen\_seven} & GaloisWhyFive & PROVED & \checkmark & verified \\
\texttt{spectralGen\_three} & GaloisWhyFive & PROVED & \checkmark & verified \\
\texttt{why\_five} & GaloisWhyFive & PROVED & \checkmark & verified \\
\texttt{d5\_card} & Geometry & PROVED & \checkmark & verified \\
\texttt{golden\_ratio\_in\_C5\_spectrum} & Geometry & PROVED & \checkmark & verified \\
\texttt{pentagon\_golden\_diagonal} & Geometry & PROVED & \checkmark & verified \\
\texttt{pentagon\_two\_distances} & Geometry & PROVED & \checkmark & verified \\
\texttt{gilbreath\_row10\_head} & GilbreathConjecture & PROVED & \checkmark & verified \\
\texttt{gilbreath\_row1\_eq} & GilbreathConjecture & PROVED & \checkmark & verified \\
\texttt{gilbreath\_row1\_head} & GilbreathConjecture & PROVED & \checkmark & verified \\
\texttt{gilbreath\_row1\_head?} & GilbreathConjecture & PROVED & \checkmark & verified \\
\texttt{gilbreath\_row2\_head} & GilbreathConjecture & PROVED & \checkmark & verified \\
\texttt{gilbreath\_row3\_head} & GilbreathConjecture & PROVED & \checkmark & verified \\
\texttt{gilbreath\_row4\_head} & GilbreathConjecture & PROVED & \checkmark & verified \\
\texttt{gilbreath\_row5\_head} & GilbreathConjecture & PROVED & \checkmark & verified \\
\texttt{gilbreath\_row6\_head} & GilbreathConjecture & PROVED & \checkmark & verified \\
\texttt{gilbreath\_row7\_head} & GilbreathConjecture & PROVED & \checkmark & verified \\
\texttt{gilbreath\_row8\_head} & GilbreathConjecture & PROVED & \checkmark & verified \\
\texttt{gilbreath\_row9\_head} & GilbreathConjecture & PROVED & \checkmark & verified \\
\texttt{giugaNumber\_squarefree} & GiugaNumbers & PROVED & \checkmark & verified \\
\texttt{giuga\_30} & GiugaNumbers & PROVED & \checkmark & verified \\
\texttt{giuga\_858} & GiugaNumbers & PROVED & \checkmark & verified \\
\texttt{K23\_above\_even\_nonthree\_baseline\_iff} & CovarianceScaffold & PROVED & \checkmark & verified \\
\texttt{K23\_of\_not\_two\_dvd} & CovarianceScaffold & PROVED & \checkmark & verified \\
\texttt{K23\_of\_two\_dvd\_not\_three\_dvd} & CovarianceScaffold & PROVED & \checkmark & verified \\
\texttt{K23\_of\_two\_dvd\_three\_dvd} & CovarianceScaffold & PROVED & \checkmark & verified \\
\texttt{K23\_pos\_iff\_two\_dvd} & CovarianceScaffold & PROVED & \checkmark & verified \\
\texttt{Kp\_three\_excess\_neg\_iff} & CovarianceScaffold & PROVED & \checkmark & verified \\
\texttt{Kp\_three\_excess\_pos\_iff} & CovarianceScaffold & PROVED & \checkmark & verified \\
\texttt{Kp\_three\_gt\_one\_iff} & CovarianceScaffold & PROVED & \checkmark & verified \\
\texttt{Kp\_three\_of\_dvd} & CovarianceScaffold & PROVED & \checkmark & verified \\
\texttt{Kp\_three\_of\_not\_dvd} & CovarianceScaffold & PROVED & \checkmark & verified \\
\texttt{Kp\_two\_eq\_zero\_iff} & CovarianceScaffold & PROVED & \checkmark & verified \\
\texttt{Kp\_two\_pos\_iff} & CovarianceScaffold & PROVED & \checkmark & verified \\
\texttt{goldbachPairTuple\_card\_le\_two} & CovarianceScaffold & PROVED & \checkmark & verified \\
\texttt{goldbachPairTuple\_raw\_admissible\_of\_even} & CovarianceScaffold & PROVED & \checkmark & verified \\
\texttt{scaledK23Excess\_pos\_iff} & CovarianceScaffold & PROVED & \checkmark & verified \\
\texttt{singular\_series\_finite\_goldbachPairTuple\_pos\_of\_even} & CovarianceScaffold & PROVED & \checkmark & verified \\
\texttt{K23\_above\_even\_nonthree\_baseline\_iff} & LocalWheel & PROVED & \checkmark & verified \\
\texttt{K23\_aligned\_gt\_baseline} & LocalWheel & PROVED & \checkmark & verified \\
\texttt{K23\_eq} & LocalWheel & PROVED & \checkmark & verified \\
\texttt{K23\_eq\_nine\_quarters\_iff} & LocalWheel & PROVED & \checkmark & verified \\
\texttt{K23\_excess\_nonneg\_of\_even} & LocalWheel & PROVED & \checkmark & verified \\
\texttt{K23\_excess\_pos\_iff} & LocalWheel & PROVED & \checkmark & verified \\
\texttt{K23\_le\_nine\_quarters} & LocalWheel & PROVED & \checkmark & verified \\
\texttt{K23\_nonneg} & LocalWheel & PROVED & \checkmark & verified \\
\texttt{K23\_of\_not\_two\_dvd} & LocalWheel & PROVED & \checkmark & verified \\
\texttt{K23\_of\_two\_dvd\_not\_three\_dvd} & LocalWheel & PROVED & \checkmark & verified \\
\texttt{K23\_of\_two\_dvd\_three\_dvd} & LocalWheel & PROVED & \checkmark & verified \\
\texttt{K23\_pos\_iff\_two\_dvd} & LocalWheel & PROVED & \checkmark & verified \\
\texttt{Kp\_five} & LocalWheel & PROVED & \checkmark & verified \\
\texttt{Kp\_five\_of\_dvd} & LocalWheel & PROVED & \checkmark & verified \\
\texttt{Kp\_five\_of\_not\_dvd} & LocalWheel & PROVED & \checkmark & verified \\
\texttt{Kp\_seven} & LocalWheel & PROVED & \checkmark & verified \\
\texttt{Kp\_seven\_of\_dvd} & LocalWheel & PROVED & \checkmark & verified \\
\texttt{Kp\_seven\_of\_not\_dvd} & LocalWheel & PROVED & \checkmark & verified \\
\texttt{Kp\_three} & LocalWheel & PROVED & \checkmark & verified \\
\texttt{Kp\_three\_aligned\_gt\_misaligned} & LocalWheel & PROVED & \checkmark & verified \\
\texttt{Kp\_three\_gt\_one\_iff} & LocalWheel & PROVED & \checkmark & verified \\
\texttt{Kp\_three\_of\_dvd} & LocalWheel & PROVED & \checkmark & verified \\
\texttt{Kp\_three\_of\_not\_dvd} & LocalWheel & PROVED & \checkmark & verified \\
\texttt{gCount\_seven} & LocalWheel & PROVED & \checkmark & verified \\
\texttt{gCount\_seven\_of\_ne\_zero} & LocalWheel & PROVED & \checkmark & verified \\
\texttt{gCount\_seven\_zero} & LocalWheel & PROVED & \checkmark & verified \\
\texttt{gCount\_three} & LocalWheel & PROVED & \checkmark & verified \\
\texttt{gCount\_three\_eq\_gResidues\_card} & LocalWheel & PROVED & \checkmark & verified \\
\texttt{gCount\_three\_of\_ne\_zero} & LocalWheel & PROVED & \checkmark & verified \\
\texttt{gCount\_three\_zero} & LocalWheel & PROVED & \checkmark & verified \\
\texttt{gResidues\_three\_card} & LocalWheel & PROVED & \checkmark & verified \\
\texttt{local\_covariance\_three\_ne\_zero} & LocalWheel & PROVED & \checkmark & verified \\
\texttt{local\_covariance\_three\_zero} & LocalWheel & PROVED & \checkmark & verified \\
\texttt{local\_covariance\_three\_zero\_gt\_ne\_zero} & LocalWheel & PROVED & \checkmark & verified \\
\texttt{local\_covariance\_two\_ne\_zero} & LocalWheel & PROVED & \checkmark & verified \\
\texttt{local\_covariance\_two\_zero} & LocalWheel & PROVED & \checkmark & verified \\
\texttt{Kp\_two} & Parity & PROVED & \checkmark & verified \\
\texttt{Kp\_two\_of\_dvd} & Parity & PROVED & \checkmark & verified \\
\texttt{Kp\_two\_of\_not\_dvd} & Parity & PROVED & \checkmark & verified \\
\texttt{even\_ge\_four\_eq\_two\_plus\_even} & Parity & PROVED & \checkmark & verified \\
\texttt{even\_of\_odd\_prime\_add\_odd\_prime} & Parity & PROVED & \checkmark & verified \\
\texttt{gCount\_eq\_gResidues\_card} & Parity & PROVED & \checkmark & verified \\
\texttt{gCount\_five} & Parity & PROVED & \checkmark & verified \\
\texttt{gCount\_two} & Parity & PROVED & \checkmark & verified \\
\texttt{gCount\_two\_of\_ne\_zero} & Parity & PROVED & \checkmark & verified \\
\texttt{gCount\_two\_one} & Parity & PROVED & \checkmark & verified \\
\texttt{gCount\_two\_zero} & Parity & PROVED & \checkmark & verified \\
\texttt{gResidues\_five\_card} & Parity & PROVED & \checkmark & verified \\
\texttt{hasGoldbachRep\_odd\_iff} & Parity & PROVED & \checkmark & verified \\
\texttt{hasGoldbachRep\_odd\_imp\_two} & Parity & PROVED & \checkmark & verified \\
\texttt{hasGoldbachRep\_two\_plus\_prime} & Parity & PROVED & \checkmark & verified \\
\texttt{odd\_sub\_of\_even\_sub\_odd\_prime} & Parity & PROVED & \checkmark & verified \\
\texttt{K235\_above\_even\_baseline\_iff} & WheelExtended & PROVED & \checkmark & verified \\
\texttt{K235\_aligned\_gt\_K23\_aligned} & WheelExtended & PROVED & \checkmark & verified \\
\texttt{K235\_aligned\_gt\_baseline} & WheelExtended & PROVED & \checkmark & verified \\
\texttt{K235\_cases} & WheelExtended & PROVED & \checkmark & verified \\
\texttt{K235\_eq} & WheelExtended & PROVED & \checkmark & verified \\
\texttt{K235\_eq\_K23\_mul\_Kp\_five} & WheelExtended & PROVED & \checkmark & verified \\
\texttt{K235\_eq\_aligned\_iff} & WheelExtended & PROVED & \checkmark & verified \\
\texttt{K235\_excess\_nonneg\_of\_even} & WheelExtended & PROVED & \checkmark & verified \\
\texttt{K235\_le\_aligned} & WheelExtended & PROVED & \checkmark & verified \\
\texttt{K235\_nonneg} & WheelExtended & PROVED & \checkmark & verified \\
\texttt{K235\_of\_not\_two\_dvd} & WheelExtended & PROVED & \checkmark & verified \\
\texttt{K235\_of\_two\_dvd\_not\_three\_not\_five} & WheelExtended & PROVED & \checkmark & verified \\
\texttt{K235\_of\_two\_five\_dvd\_not\_three} & WheelExtended & PROVED & \checkmark & verified \\
\texttt{K235\_of\_two\_three\_dvd\_not\_five} & WheelExtended & PROVED & \checkmark & verified \\
\texttt{K235\_of\_two\_three\_five\_dvd} & WheelExtended & PROVED & \checkmark & verified \\
\texttt{K235\_pos\_iff\_two\_dvd} & WheelExtended & PROVED & \checkmark & verified \\
\texttt{Kp\_eleven} & WheelExtended & PROVED & \checkmark & verified \\
\texttt{Kp\_eleven\_aligned\_gt\_misaligned} & WheelExtended & PROVED & \checkmark & verified \\
\texttt{Kp\_eleven\_gt\_one\_iff} & WheelExtended & PROVED & \checkmark & verified \\
\texttt{Kp\_eleven\_le\_aligned} & WheelExtended & PROVED & \checkmark & verified \\
\texttt{Kp\_eleven\_of\_dvd} & WheelExtended & PROVED & \checkmark & verified \\
\texttt{Kp\_eleven\_of\_not\_dvd} & WheelExtended & PROVED & \checkmark & verified \\
\texttt{Kp\_eleven\_pos} & WheelExtended & PROVED & \checkmark & verified \\
\texttt{Kp\_thirteen} & WheelExtended & PROVED & \checkmark & verified \\
\texttt{Kp\_thirteen\_aligned\_gt\_misaligned} & WheelExtended & PROVED & \checkmark & verified \\
\texttt{Kp\_thirteen\_gt\_one\_iff} & WheelExtended & PROVED & \checkmark & verified \\
\texttt{Kp\_thirteen\_le\_aligned} & WheelExtended & PROVED & \checkmark & verified \\
\texttt{Kp\_thirteen\_of\_dvd} & WheelExtended & PROVED & \checkmark & verified \\
\texttt{Kp\_thirteen\_of\_not\_dvd} & WheelExtended & PROVED & \checkmark & verified \\
\texttt{Kp\_thirteen\_pos} & WheelExtended & PROVED & \checkmark & verified \\
\texttt{gCount\_eleven} & WheelExtended & PROVED & \checkmark & verified \\
\texttt{gCount\_eleven\_eq\_gResidues\_card} & WheelExtended & PROVED & \checkmark & verified \\
\texttt{gCount\_eleven\_of\_ne\_zero} & WheelExtended & PROVED & \checkmark & verified \\
\texttt{gCount\_eleven\_zero} & WheelExtended & PROVED & \checkmark & verified \\
\texttt{gCount\_thirteen} & WheelExtended & PROVED & \checkmark & verified \\
\texttt{gCount\_thirteen\_eq\_gResidues\_card} & WheelExtended & PROVED & \checkmark & verified \\
\texttt{gCount\_thirteen\_of\_ne\_zero} & WheelExtended & PROVED & \checkmark & verified \\
\texttt{gCount\_thirteen\_zero} & WheelExtended & PROVED & \checkmark & verified \\
\texttt{gResidues\_eleven\_card} & WheelExtended & PROVED & \checkmark & verified \\
\texttt{gResidues\_thirteen\_card} & WheelExtended & PROVED & \checkmark & verified \\
\texttt{K2357\_aligned\_gt\_K235\_aligned} & WheelK2357 & PROVED & \checkmark & verified \\
\texttt{K2357\_aligned\_gt\_K23\_aligned} & WheelK2357 & PROVED & \checkmark & verified \\
\texttt{K2357\_aligned\_gt\_baseline} & WheelK2357 & PROVED & \checkmark & verified \\
\texttt{K2357\_cases} & WheelK2357 & PROVED & \checkmark & verified \\
\texttt{K2357\_eq} & WheelK2357 & PROVED & \checkmark & verified \\
\texttt{K2357\_eq\_K235\_mul\_Kp\_seven} & WheelK2357 & PROVED & \checkmark & verified \\
\texttt{K2357\_eq\_K23\_mul\_Kp\_five\_mul\_Kp\_seven} & WheelK2357 & PROVED & \checkmark & verified \\
\texttt{K2357\_eq\_aligned\_iff} & WheelK2357 & PROVED & \checkmark & verified \\
\texttt{K2357\_eq\_zero\_iff\_not\_two\_dvd} & WheelK2357 & PROVED & \checkmark & verified \\
\texttt{K2357\_excess\_nonneg\_of\_even} & WheelK2357 & PROVED & \checkmark & verified \\
\texttt{K2357\_le\_aligned} & WheelK2357 & PROVED & \checkmark & verified \\
\texttt{K2357\_nonneg} & WheelK2357 & PROVED & \checkmark & verified \\
\texttt{K2357\_of\_not\_two\_dvd} & WheelK2357 & PROVED & \checkmark & verified \\
\texttt{K2357\_of\_two\_dvd\_not\_three\_not\_five\_not\_seven} & WheelK2357 & PROVED & \checkmark & verified \\
\texttt{K2357\_of\_two\_five\_dvd\_not\_three\_not\_seven} & WheelK2357 & PROVED & \checkmark & verified \\
\texttt{K2357\_of\_two\_five\_seven\_dvd\_not\_three} & WheelK2357 & PROVED & \checkmark & verified \\
\texttt{K2357\_of\_two\_seven\_dvd\_not\_three\_not\_five} & WheelK2357 & PROVED & \checkmark & verified \\
\texttt{K2357\_of\_two\_three\_dvd\_not\_five\_not\_seven} & WheelK2357 & PROVED & \checkmark & verified \\
\texttt{K2357\_of\_two\_three\_five\_dvd\_not\_seven} & WheelK2357 & PROVED & \checkmark & verified \\
\texttt{K2357\_of\_two\_three\_five\_seven\_dvd} & WheelK2357 & PROVED & \checkmark & verified \\
\texttt{K2357\_of\_two\_three\_seven\_dvd\_not\_five} & WheelK2357 & PROVED & \checkmark & verified \\
\texttt{K2357\_pos\_iff\_two\_dvd} & WheelK2357 & PROVED & \checkmark & verified \\
\texttt{K2\_11\_eq} & WheelK235711 & PROVED & \checkmark & verified \\
\texttt{K2\_11\_of\_not\_two\_dvd} & WheelK235711 & PROVED & \checkmark & verified \\
\texttt{K2\_11\_of\_two\_and\_eleven\_dvd} & WheelK235711 & PROVED & \checkmark & verified \\
\texttt{K2\_101\_eq} & WheelK2\_101 & PROVED & \checkmark & verified \\
\texttt{K2\_101\_of\_not\_two\_dvd} & WheelK2\_101 & PROVED & \checkmark & verified \\
\texttt{K2\_101\_of\_two\_and\_oneHundredOne\_dvd} & WheelK2\_101 & PROVED & \checkmark & verified \\
\texttt{Kp\_oneHundredOne} & WheelK2\_101 & PROVED & \checkmark & verified \\
\texttt{Kp\_oneHundredOne\_of\_dvd} & WheelK2\_101 & PROVED & \checkmark & verified \\
\texttt{Kp\_oneHundredOne\_of\_not\_dvd} & WheelK2\_101 & PROVED & \checkmark & verified \\
\texttt{K2\_103\_eq} & WheelK2\_103 & PROVED & \checkmark & verified \\
\texttt{K2\_103\_of\_not\_two\_dvd} & WheelK2\_103 & PROVED & \checkmark & verified \\
\texttt{K2\_103\_of\_two\_and\_oneHundredThree\_dvd} & WheelK2\_103 & PROVED & \checkmark & verified \\
\texttt{Kp\_oneHundredThree} & WheelK2\_103 & PROVED & \checkmark & verified \\
\texttt{Kp\_oneHundredThree\_of\_dvd} & WheelK2\_103 & PROVED & \checkmark & verified \\
\texttt{Kp\_oneHundredThree\_of\_not\_dvd} & WheelK2\_103 & PROVED & \checkmark & verified \\
\texttt{K2\_107\_eq} & WheelK2\_107 & PROVED & \checkmark & verified \\
\texttt{K2\_107\_of\_not\_two\_dvd} & WheelK2\_107 & PROVED & \checkmark & verified \\
\texttt{K2\_107\_of\_two\_and\_oneHundredSeven\_dvd} & WheelK2\_107 & PROVED & \checkmark & verified \\
\texttt{Kp\_oneHundredSeven} & WheelK2\_107 & PROVED & \checkmark & verified \\
\texttt{Kp\_oneHundredSeven\_of\_dvd} & WheelK2\_107 & PROVED & \checkmark & verified \\
\texttt{Kp\_oneHundredSeven\_of\_not\_dvd} & WheelK2\_107 & PROVED & \checkmark & verified \\
\texttt{K2\_109\_eq} & WheelK2\_109 & PROVED & \checkmark & verified \\
\texttt{K2\_109\_of\_not\_two\_dvd} & WheelK2\_109 & PROVED & \checkmark & verified \\
\texttt{K2\_109\_of\_two\_and\_oneHundredNine\_dvd} & WheelK2\_109 & PROVED & \checkmark & verified \\
\texttt{Kp\_oneHundredNine} & WheelK2\_109 & PROVED & \checkmark & verified \\
\texttt{Kp\_oneHundredNine\_of\_dvd} & WheelK2\_109 & PROVED & \checkmark & verified \\
\texttt{Kp\_oneHundredNine\_of\_not\_dvd} & WheelK2\_109 & PROVED & \checkmark & verified \\
\texttt{K2\_113\_eq} & WheelK2\_113 & PROVED & \checkmark & verified \\
\texttt{K2\_113\_of\_not\_two\_dvd} & WheelK2\_113 & PROVED & \checkmark & verified \\
\texttt{K2\_113\_of\_two\_and\_oneHundredThirteen\_dvd} & WheelK2\_113 & PROVED & \checkmark & verified \\
\texttt{Kp\_oneHundredThirteen} & WheelK2\_113 & PROVED & \checkmark & verified \\
\texttt{Kp\_oneHundredThirteen\_of\_dvd} & WheelK2\_113 & PROVED & \checkmark & verified \\
\texttt{Kp\_oneHundredThirteen\_of\_not\_dvd} & WheelK2\_113 & PROVED & \checkmark & verified \\
\texttt{K2\_127\_eq} & WheelK2\_127 & PROVED & \checkmark & verified \\
\texttt{K2\_127\_of\_not\_two\_dvd} & WheelK2\_127 & PROVED & \checkmark & verified \\
\texttt{K2\_127\_of\_two\_and\_oneHundredTwentySeven\_dvd} & WheelK2\_127 & PROVED & \checkmark & verified \\
\texttt{Kp\_oneHundredTwentySeven} & WheelK2\_127 & PROVED & \checkmark & verified \\
\texttt{Kp\_oneHundredTwentySeven\_of\_dvd} & WheelK2\_127 & PROVED & \checkmark & verified \\
\texttt{Kp\_oneHundredTwentySeven\_of\_not\_dvd} & WheelK2\_127 & PROVED & \checkmark & verified \\
\texttt{K2\_13\_eq} & WheelK2\_13 & PROVED & \checkmark & verified \\
\texttt{K2\_13\_of\_not\_two\_dvd} & WheelK2\_13 & PROVED & \checkmark & verified \\
\texttt{K2\_13\_of\_two\_and\_thirteen\_dvd} & WheelK2\_13 & PROVED & \checkmark & verified \\
\texttt{K2\_131\_eq} & WheelK2\_131 & PROVED & \checkmark & verified \\
\texttt{K2\_131\_of\_not\_two\_dvd} & WheelK2\_131 & PROVED & \checkmark & verified \\
\texttt{K2\_131\_of\_two\_and\_oneHundredThirtyOne\_dvd} & WheelK2\_131 & PROVED & \checkmark & verified \\
\texttt{Kp\_oneHundredThirtyOne} & WheelK2\_131 & PROVED & \checkmark & verified \\
\texttt{Kp\_oneHundredThirtyOne\_of\_dvd} & WheelK2\_131 & PROVED & \checkmark & verified \\
\texttt{Kp\_oneHundredThirtyOne\_of\_not\_dvd} & WheelK2\_131 & PROVED & \checkmark & verified \\
\texttt{K2\_137\_eq} & WheelK2\_137 & PROVED & \checkmark & verified \\
\texttt{K2\_137\_of\_not\_two\_dvd} & WheelK2\_137 & PROVED & \checkmark & verified \\
\texttt{K2\_137\_of\_two\_and\_oneHundredThirtySeven\_dvd} & WheelK2\_137 & PROVED & \checkmark & verified \\
\texttt{Kp\_oneHundredThirtySeven} & WheelK2\_137 & PROVED & \checkmark & verified \\
\texttt{Kp\_oneHundredThirtySeven\_of\_dvd} & WheelK2\_137 & PROVED & \checkmark & verified \\
\texttt{Kp\_oneHundredThirtySeven\_of\_not\_dvd} & WheelK2\_137 & PROVED & \checkmark & verified \\
\texttt{K2\_139\_eq} & WheelK2\_139 & PROVED & \checkmark & verified \\
\texttt{K2\_139\_of\_not\_two\_dvd} & WheelK2\_139 & PROVED & \checkmark & verified \\
\texttt{K2\_139\_of\_two\_and\_oneHundredThirtyNine\_dvd} & WheelK2\_139 & PROVED & \checkmark & verified \\
\texttt{Kp\_oneHundredThirtyNine} & WheelK2\_139 & PROVED & \checkmark & verified \\
\texttt{Kp\_oneHundredThirtyNine\_of\_dvd} & WheelK2\_139 & PROVED & \checkmark & verified \\
\texttt{Kp\_oneHundredThirtyNine\_of\_not\_dvd} & WheelK2\_139 & PROVED & \checkmark & verified \\
\texttt{K2\_149\_eq} & WheelK2\_149 & PROVED & \checkmark & verified \\
\texttt{K2\_149\_of\_not\_two\_dvd} & WheelK2\_149 & PROVED & \checkmark & verified \\
\texttt{K2\_149\_of\_two\_and\_oneHundredFortyNine\_dvd} & WheelK2\_149 & PROVED & \checkmark & verified \\
\texttt{Kp\_oneHundredFortyNine} & WheelK2\_149 & PROVED & \checkmark & verified \\
\texttt{Kp\_oneHundredFortyNine\_of\_dvd} & WheelK2\_149 & PROVED & \checkmark & verified \\
\texttt{Kp\_oneHundredFortyNine\_of\_not\_dvd} & WheelK2\_149 & PROVED & \checkmark & verified \\
\texttt{K2\_151\_eq} & WheelK2\_151 & PROVED & \checkmark & verified \\
\texttt{K2\_151\_of\_not\_two\_dvd} & WheelK2\_151 & PROVED & \checkmark & verified \\
\texttt{K2\_151\_of\_two\_and\_oneHundredFiftyOne\_dvd} & WheelK2\_151 & PROVED & \checkmark & verified \\
\texttt{Kp\_oneHundredFiftyOne} & WheelK2\_151 & PROVED & \checkmark & verified \\
\texttt{Kp\_oneHundredFiftyOne\_of\_dvd} & WheelK2\_151 & PROVED & \checkmark & verified \\
\texttt{Kp\_oneHundredFiftyOne\_of\_not\_dvd} & WheelK2\_151 & PROVED & \checkmark & verified \\
\texttt{K2\_157\_eq} & WheelK2\_157 & PROVED & \checkmark & verified \\
\texttt{K2\_157\_of\_not\_two\_dvd} & WheelK2\_157 & PROVED & \checkmark & verified \\
\texttt{K2\_157\_of\_two\_and\_oneHundredFiftySeven\_dvd} & WheelK2\_157 & PROVED & \checkmark & verified \\
\texttt{Kp\_oneHundredFiftySeven} & WheelK2\_157 & PROVED & \checkmark & verified \\
\texttt{Kp\_oneHundredFiftySeven\_of\_dvd} & WheelK2\_157 & PROVED & \checkmark & verified \\
\texttt{Kp\_oneHundredFiftySeven\_of\_not\_dvd} & WheelK2\_157 & PROVED & \checkmark & verified \\
\texttt{K2\_163\_eq} & WheelK2\_163 & PROVED & \checkmark & verified \\
\texttt{K2\_163\_of\_not\_two\_dvd} & WheelK2\_163 & PROVED & \checkmark & verified \\
\texttt{K2\_163\_of\_two\_and\_oneHundredSixtyThree\_dvd} & WheelK2\_163 & PROVED & \checkmark & verified \\
\texttt{Kp\_oneHundredSixtyThree} & WheelK2\_163 & PROVED & \checkmark & verified \\
\texttt{Kp\_oneHundredSixtyThree\_of\_dvd} & WheelK2\_163 & PROVED & \checkmark & verified \\
\texttt{Kp\_oneHundredSixtyThree\_of\_not\_dvd} & WheelK2\_163 & PROVED & \checkmark & verified \\
\texttt{K2\_167\_eq} & WheelK2\_167 & PROVED & \checkmark & verified \\
\texttt{K2\_167\_of\_not\_two\_dvd} & WheelK2\_167 & PROVED & \checkmark & verified \\
\texttt{K2\_167\_of\_two\_and\_oneHundredSixtySeven\_dvd} & WheelK2\_167 & PROVED & \checkmark & verified \\
\texttt{Kp\_oneHundredSixtySeven} & WheelK2\_167 & PROVED & \checkmark & verified \\
\texttt{Kp\_oneHundredSixtySeven\_of\_dvd} & WheelK2\_167 & PROVED & \checkmark & verified \\
\texttt{Kp\_oneHundredSixtySeven\_of\_not\_dvd} & WheelK2\_167 & PROVED & \checkmark & verified \\
\texttt{K2\_17\_eq} & WheelK2\_17 & PROVED & \checkmark & verified \\
\texttt{K2\_17\_of\_not\_two\_dvd} & WheelK2\_17 & PROVED & \checkmark & verified \\
\texttt{K2\_17\_of\_two\_and\_seventeen\_dvd} & WheelK2\_17 & PROVED & \checkmark & verified \\
\texttt{Kp\_seventeen} & WheelK2\_17 & PROVED & \checkmark & verified \\
\texttt{Kp\_seventeen\_of\_dvd} & WheelK2\_17 & PROVED & \checkmark & verified \\
\texttt{Kp\_seventeen\_of\_not\_dvd} & WheelK2\_17 & PROVED & \checkmark & verified \\
\texttt{K2\_173\_eq} & WheelK2\_173 & PROVED & \checkmark & verified \\
\texttt{K2\_173\_of\_not\_two\_dvd} & WheelK2\_173 & PROVED & \checkmark & verified \\
\texttt{K2\_173\_of\_two\_and\_oneHundredSeventyThree\_dvd} & WheelK2\_173 & PROVED & \checkmark & verified \\
\texttt{Kp\_oneHundredSeventyThree} & WheelK2\_173 & PROVED & \checkmark & verified \\
\texttt{Kp\_oneHundredSeventyThree\_of\_dvd} & WheelK2\_173 & PROVED & \checkmark & verified \\
\texttt{Kp\_oneHundredSeventyThree\_of\_not\_dvd} & WheelK2\_173 & PROVED & \checkmark & verified \\
\texttt{K2\_179\_eq} & WheelK2\_179 & PROVED & \checkmark & verified \\
\texttt{K2\_179\_of\_not\_two\_dvd} & WheelK2\_179 & PROVED & \checkmark & verified \\
\texttt{K2\_179\_of\_two\_and\_oneHundredSeventyNine\_dvd} & WheelK2\_179 & PROVED & \checkmark & verified \\
\texttt{Kp\_oneHundredSeventyNine} & WheelK2\_179 & PROVED & \checkmark & verified \\
\texttt{Kp\_oneHundredSeventyNine\_of\_dvd} & WheelK2\_179 & PROVED & \checkmark & verified \\
\texttt{Kp\_oneHundredSeventyNine\_of\_not\_dvd} & WheelK2\_179 & PROVED & \checkmark & verified \\
\texttt{K2\_181\_eq} & WheelK2\_181 & PROVED & \checkmark & verified \\
\texttt{K2\_181\_of\_not\_two\_dvd} & WheelK2\_181 & PROVED & \checkmark & verified \\
\texttt{K2\_181\_of\_two\_and\_oneHundredEightyOne\_dvd} & WheelK2\_181 & PROVED & \checkmark & verified \\
\texttt{Kp\_oneHundredEightyOne} & WheelK2\_181 & PROVED & \checkmark & verified \\
\texttt{Kp\_oneHundredEightyOne\_of\_dvd} & WheelK2\_181 & PROVED & \checkmark & verified \\
\texttt{Kp\_oneHundredEightyOne\_of\_not\_dvd} & WheelK2\_181 & PROVED & \checkmark & verified \\
\texttt{K2\_19\_eq} & WheelK2\_19 & PROVED & \checkmark & verified \\
\texttt{K2\_19\_of\_not\_two\_dvd} & WheelK2\_19 & PROVED & \checkmark & verified \\
\texttt{K2\_19\_of\_two\_and\_nineteen\_dvd} & WheelK2\_19 & PROVED & \checkmark & verified \\
\texttt{Kp\_nineteen} & WheelK2\_19 & PROVED & \checkmark & verified \\
\texttt{Kp\_nineteen\_of\_dvd} & WheelK2\_19 & PROVED & \checkmark & verified \\
\texttt{Kp\_nineteen\_of\_not\_dvd} & WheelK2\_19 & PROVED & \checkmark & verified \\
\texttt{K2\_191\_eq} & WheelK2\_191 & PROVED & \checkmark & verified \\
\texttt{K2\_191\_of\_not\_two\_dvd} & WheelK2\_191 & PROVED & \checkmark & verified \\
\texttt{K2\_191\_of\_two\_and\_oneHundredNinetyOne\_dvd} & WheelK2\_191 & PROVED & \checkmark & verified \\
\texttt{Kp\_oneHundredNinetyOne} & WheelK2\_191 & PROVED & \checkmark & verified \\
\texttt{Kp\_oneHundredNinetyOne\_of\_dvd} & WheelK2\_191 & PROVED & \checkmark & verified \\
\texttt{Kp\_oneHundredNinetyOne\_of\_not\_dvd} & WheelK2\_191 & PROVED & \checkmark & verified \\
\texttt{K2\_193\_eq} & WheelK2\_193 & PROVED & \checkmark & verified \\
\texttt{K2\_193\_of\_not\_two\_dvd} & WheelK2\_193 & PROVED & \checkmark & verified \\
\texttt{K2\_193\_of\_two\_and\_oneHundredNinetyThree\_dvd} & WheelK2\_193 & PROVED & \checkmark & verified \\
\texttt{Kp\_oneHundredNinetyThree} & WheelK2\_193 & PROVED & \checkmark & verified \\
\texttt{Kp\_oneHundredNinetyThree\_of\_dvd} & WheelK2\_193 & PROVED & \checkmark & verified \\
\texttt{Kp\_oneHundredNinetyThree\_of\_not\_dvd} & WheelK2\_193 & PROVED & \checkmark & verified \\
\texttt{K2\_197\_eq} & WheelK2\_197 & PROVED & \checkmark & verified \\
\texttt{K2\_197\_of\_not\_two\_dvd} & WheelK2\_197 & PROVED & \checkmark & verified \\
\texttt{K2\_197\_of\_two\_and\_oneHundredNinetySeven\_dvd} & WheelK2\_197 & PROVED & \checkmark & verified \\
\texttt{Kp\_oneHundredNinetySeven} & WheelK2\_197 & PROVED & \checkmark & verified \\
\texttt{Kp\_oneHundredNinetySeven\_of\_dvd} & WheelK2\_197 & PROVED & \checkmark & verified \\
\texttt{Kp\_oneHundredNinetySeven\_of\_not\_dvd} & WheelK2\_197 & PROVED & \checkmark & verified \\
\texttt{K2\_199\_eq} & WheelK2\_199 & PROVED & \checkmark & verified \\
\texttt{K2\_199\_of\_not\_two\_dvd} & WheelK2\_199 & PROVED & \checkmark & verified \\
\texttt{K2\_199\_of\_two\_and\_oneHundredNinetyNine\_dvd} & WheelK2\_199 & PROVED & \checkmark & verified \\
\texttt{Kp\_oneHundredNinetyNine} & WheelK2\_199 & PROVED & \checkmark & verified \\
\texttt{Kp\_oneHundredNinetyNine\_of\_dvd} & WheelK2\_199 & PROVED & \checkmark & verified \\
\texttt{Kp\_oneHundredNinetyNine\_of\_not\_dvd} & WheelK2\_199 & PROVED & \checkmark & verified \\
\texttt{K2\_211\_eq} & WheelK2\_211 & PROVED & \checkmark & verified \\
\texttt{K2\_211\_of\_not\_two\_dvd} & WheelK2\_211 & PROVED & \checkmark & verified \\
\texttt{K2\_211\_of\_two\_and\_twoHundredEleven\_dvd} & WheelK2\_211 & PROVED & \checkmark & verified \\
\texttt{Kp\_twoHundredEleven} & WheelK2\_211 & PROVED & \checkmark & verified \\
\texttt{Kp\_twoHundredEleven\_of\_dvd} & WheelK2\_211 & PROVED & \checkmark & verified \\
\texttt{Kp\_twoHundredEleven\_of\_not\_dvd} & WheelK2\_211 & PROVED & \checkmark & verified \\
\texttt{K2\_223\_eq} & WheelK2\_223 & PROVED & \checkmark & verified \\
\texttt{K2\_223\_of\_not\_two\_dvd} & WheelK2\_223 & PROVED & \checkmark & verified \\
\texttt{K2\_223\_of\_two\_and\_twoHundredTwentyThree\_dvd} & WheelK2\_223 & PROVED & \checkmark & verified \\
\texttt{Kp\_twoHundredTwentyThree} & WheelK2\_223 & PROVED & \checkmark & verified \\
\texttt{Kp\_twoHundredTwentyThree\_of\_dvd} & WheelK2\_223 & PROVED & \checkmark & verified \\
\texttt{Kp\_twoHundredTwentyThree\_of\_not\_dvd} & WheelK2\_223 & PROVED & \checkmark & verified \\
\texttt{K2\_227\_eq} & WheelK2\_227 & PROVED & \checkmark & verified \\
\texttt{K2\_227\_of\_not\_two\_dvd} & WheelK2\_227 & PROVED & \checkmark & verified \\
\texttt{K2\_227\_of\_two\_and\_twoHundredTwentySeven\_dvd} & WheelK2\_227 & PROVED & \checkmark & verified \\
\texttt{Kp\_twoHundredTwentySeven} & WheelK2\_227 & PROVED & \checkmark & verified \\
\texttt{Kp\_twoHundredTwentySeven\_of\_dvd} & WheelK2\_227 & PROVED & \checkmark & verified \\
\texttt{Kp\_twoHundredTwentySeven\_of\_not\_dvd} & WheelK2\_227 & PROVED & \checkmark & verified \\
\texttt{K2\_229\_eq} & WheelK2\_229 & PROVED & \checkmark & verified \\
\texttt{K2\_229\_of\_not\_two\_dvd} & WheelK2\_229 & PROVED & \checkmark & verified \\
\texttt{K2\_229\_of\_two\_and\_twoHundredTwentyNine\_dvd} & WheelK2\_229 & PROVED & \checkmark & verified \\
\texttt{Kp\_twoHundredTwentyNine} & WheelK2\_229 & PROVED & \checkmark & verified \\
\texttt{Kp\_twoHundredTwentyNine\_of\_dvd} & WheelK2\_229 & PROVED & \checkmark & verified \\
\texttt{Kp\_twoHundredTwentyNine\_of\_not\_dvd} & WheelK2\_229 & PROVED & \checkmark & verified \\
\texttt{K2\_23\_eq} & WheelK2\_23 & PROVED & \checkmark & verified \\
\texttt{K2\_23\_of\_not\_two\_dvd} & WheelK2\_23 & PROVED & \checkmark & verified \\
\texttt{K2\_23\_of\_two\_and\_twentyThree\_dvd} & WheelK2\_23 & PROVED & \checkmark & verified \\
\texttt{Kp\_twentyThree} & WheelK2\_23 & PROVED & \checkmark & verified \\
\texttt{Kp\_twentyThree\_of\_dvd} & WheelK2\_23 & PROVED & \checkmark & verified \\
\texttt{Kp\_twentyThree\_of\_not\_dvd} & WheelK2\_23 & PROVED & \checkmark & verified \\
\texttt{K2\_233\_eq} & WheelK2\_233 & PROVED & \checkmark & verified \\
\texttt{K2\_233\_of\_not\_two\_dvd} & WheelK2\_233 & PROVED & \checkmark & verified \\
\texttt{K2\_233\_of\_two\_and\_twoHundredThirtyThree\_dvd} & WheelK2\_233 & PROVED & \checkmark & verified \\
\texttt{Kp\_twoHundredThirtyThree} & WheelK2\_233 & PROVED & \checkmark & verified \\
\texttt{Kp\_twoHundredThirtyThree\_of\_dvd} & WheelK2\_233 & PROVED & \checkmark & verified \\
\texttt{Kp\_twoHundredThirtyThree\_of\_not\_dvd} & WheelK2\_233 & PROVED & \checkmark & verified \\
\texttt{K2\_239\_eq} & WheelK2\_239 & PROVED & \checkmark & verified \\
\texttt{K2\_239\_of\_not\_two\_dvd} & WheelK2\_239 & PROVED & \checkmark & verified \\
\texttt{K2\_239\_of\_two\_and\_twoHundredThirtyNine\_dvd} & WheelK2\_239 & PROVED & \checkmark & verified \\
\texttt{Kp\_twoHundredThirtyNine} & WheelK2\_239 & PROVED & \checkmark & verified \\
\texttt{Kp\_twoHundredThirtyNine\_of\_dvd} & WheelK2\_239 & PROVED & \checkmark & verified \\
\texttt{Kp\_twoHundredThirtyNine\_of\_not\_dvd} & WheelK2\_239 & PROVED & \checkmark & verified \\
\texttt{K2\_241\_eq} & WheelK2\_241 & PROVED & \checkmark & verified \\
\texttt{K2\_241\_of\_not\_two\_dvd} & WheelK2\_241 & PROVED & \checkmark & verified \\
\texttt{K2\_241\_of\_two\_and\_twoHundredFortyOne\_dvd} & WheelK2\_241 & PROVED & \checkmark & verified \\
\texttt{Kp\_twoHundredFortyOne} & WheelK2\_241 & PROVED & \checkmark & verified \\
\texttt{Kp\_twoHundredFortyOne\_of\_dvd} & WheelK2\_241 & PROVED & \checkmark & verified \\
\texttt{Kp\_twoHundredFortyOne\_of\_not\_dvd} & WheelK2\_241 & PROVED & \checkmark & verified \\
\texttt{K2\_251\_eq} & WheelK2\_251 & PROVED & \checkmark & verified \\
\texttt{K2\_251\_of\_not\_two\_dvd} & WheelK2\_251 & PROVED & \checkmark & verified \\
\texttt{K2\_251\_of\_two\_and\_twoHundredFiftyOne\_dvd} & WheelK2\_251 & PROVED & \checkmark & verified \\
\texttt{Kp\_twoHundredFiftyOne} & WheelK2\_251 & PROVED & \checkmark & verified \\
\texttt{Kp\_twoHundredFiftyOne\_of\_dvd} & WheelK2\_251 & PROVED & \checkmark & verified \\
\texttt{Kp\_twoHundredFiftyOne\_of\_not\_dvd} & WheelK2\_251 & PROVED & \checkmark & verified \\
\texttt{K2\_257\_eq} & WheelK2\_257 & PROVED & \checkmark & verified \\
\texttt{K2\_257\_of\_not\_two\_dvd} & WheelK2\_257 & PROVED & \checkmark & verified \\
\texttt{K2\_257\_of\_two\_and\_twoHundredFiftySeven\_dvd} & WheelK2\_257 & PROVED & \checkmark & verified \\
\texttt{Kp\_twoHundredFiftySeven} & WheelK2\_257 & PROVED & \checkmark & verified \\
\texttt{Kp\_twoHundredFiftySeven\_of\_dvd} & WheelK2\_257 & PROVED & \checkmark & verified \\
\texttt{Kp\_twoHundredFiftySeven\_of\_not\_dvd} & WheelK2\_257 & PROVED & \checkmark & verified \\
\texttt{K2\_263\_eq} & WheelK2\_263 & PROVED & \checkmark & verified \\
\texttt{K2\_263\_of\_not\_two\_dvd} & WheelK2\_263 & PROVED & \checkmark & verified \\
\texttt{K2\_263\_of\_two\_and\_twoHundredSixtyThree\_dvd} & WheelK2\_263 & PROVED & \checkmark & verified \\
\texttt{Kp\_twoHundredSixtyThree} & WheelK2\_263 & PROVED & \checkmark & verified \\
\texttt{Kp\_twoHundredSixtyThree\_of\_dvd} & WheelK2\_263 & PROVED & \checkmark & verified \\
\texttt{Kp\_twoHundredSixtyThree\_of\_not\_dvd} & WheelK2\_263 & PROVED & \checkmark & verified \\
\texttt{K2\_269\_eq} & WheelK2\_269 & PROVED & \checkmark & verified \\
\texttt{K2\_269\_of\_not\_two\_dvd} & WheelK2\_269 & PROVED & \checkmark & verified \\
\texttt{K2\_269\_of\_two\_and\_twoHundredSixtyNine\_dvd} & WheelK2\_269 & PROVED & \checkmark & verified \\
\texttt{Kp\_twoHundredSixtyNine} & WheelK2\_269 & PROVED & \checkmark & verified \\
\texttt{Kp\_twoHundredSixtyNine\_of\_dvd} & WheelK2\_269 & PROVED & \checkmark & verified \\
\texttt{Kp\_twoHundredSixtyNine\_of\_not\_dvd} & WheelK2\_269 & PROVED & \checkmark & verified \\
\texttt{K2\_271\_eq} & WheelK2\_271 & PROVED & \checkmark & verified \\
\texttt{K2\_271\_of\_not\_two\_dvd} & WheelK2\_271 & PROVED & \checkmark & verified \\
\texttt{K2\_271\_of\_two\_and\_twoHundredSeventyOne\_dvd} & WheelK2\_271 & PROVED & \checkmark & verified \\
\texttt{Kp\_twoHundredSeventyOne} & WheelK2\_271 & PROVED & \checkmark & verified \\
\texttt{Kp\_twoHundredSeventyOne\_of\_dvd} & WheelK2\_271 & PROVED & \checkmark & verified \\
\texttt{Kp\_twoHundredSeventyOne\_of\_not\_dvd} & WheelK2\_271 & PROVED & \checkmark & verified \\
\texttt{K2\_277\_eq} & WheelK2\_277 & PROVED & \checkmark & verified \\
\texttt{K2\_277\_of\_not\_two\_dvd} & WheelK2\_277 & PROVED & \checkmark & verified \\
\texttt{K2\_277\_of\_two\_and\_twoHundredSeventySeven\_dvd} & WheelK2\_277 & PROVED & \checkmark & verified \\
\texttt{Kp\_twoHundredSeventySeven} & WheelK2\_277 & PROVED & \checkmark & verified \\
\texttt{Kp\_twoHundredSeventySeven\_of\_dvd} & WheelK2\_277 & PROVED & \checkmark & verified \\
\texttt{Kp\_twoHundredSeventySeven\_of\_not\_dvd} & WheelK2\_277 & PROVED & \checkmark & verified \\
\texttt{K2\_281\_eq} & WheelK2\_281 & PROVED & \checkmark & verified \\
\texttt{K2\_281\_of\_not\_two\_dvd} & WheelK2\_281 & PROVED & \checkmark & verified \\
\texttt{K2\_281\_of\_two\_and\_twoHundredEightyOne\_dvd} & WheelK2\_281 & PROVED & \checkmark & verified \\
\texttt{Kp\_twoHundredEightyOne} & WheelK2\_281 & PROVED & \checkmark & verified \\
\texttt{Kp\_twoHundredEightyOne\_of\_dvd} & WheelK2\_281 & PROVED & \checkmark & verified \\
\texttt{Kp\_twoHundredEightyOne\_of\_not\_dvd} & WheelK2\_281 & PROVED & \checkmark & verified \\
\texttt{K2\_283\_eq} & WheelK2\_283 & PROVED & \checkmark & verified \\
\texttt{K2\_283\_of\_not\_two\_dvd} & WheelK2\_283 & PROVED & \checkmark & verified \\
\texttt{K2\_283\_of\_two\_and\_twoHundredEightyThree\_dvd} & WheelK2\_283 & PROVED & \checkmark & verified \\
\texttt{Kp\_twoHundredEightyThree} & WheelK2\_283 & PROVED & \checkmark & verified \\
\texttt{Kp\_twoHundredEightyThree\_of\_dvd} & WheelK2\_283 & PROVED & \checkmark & verified \\
\texttt{Kp\_twoHundredEightyThree\_of\_not\_dvd} & WheelK2\_283 & PROVED & \checkmark & verified \\
\texttt{K2\_29\_eq} & WheelK2\_29 & PROVED & \checkmark & verified \\
\texttt{K2\_29\_of\_not\_two\_dvd} & WheelK2\_29 & PROVED & \checkmark & verified \\
\texttt{K2\_29\_of\_two\_and\_twentyNine\_dvd} & WheelK2\_29 & PROVED & \checkmark & verified \\
\texttt{Kp\_twentyNine} & WheelK2\_29 & PROVED & \checkmark & verified \\
\texttt{Kp\_twentyNine\_of\_dvd} & WheelK2\_29 & PROVED & \checkmark & verified \\
\texttt{Kp\_twentyNine\_of\_not\_dvd} & WheelK2\_29 & PROVED & \checkmark & verified \\
\texttt{K2\_293\_eq} & WheelK2\_293 & PROVED & \checkmark & verified \\
\texttt{K2\_293\_of\_not\_two\_dvd} & WheelK2\_293 & PROVED & \checkmark & verified \\
\texttt{K2\_293\_of\_two\_and\_twoHundredNinetyThree\_dvd} & WheelK2\_293 & PROVED & \checkmark & verified \\
\texttt{Kp\_twoHundredNinetyThree} & WheelK2\_293 & PROVED & \checkmark & verified \\
\texttt{Kp\_twoHundredNinetyThree\_of\_dvd} & WheelK2\_293 & PROVED & \checkmark & verified \\
\texttt{Kp\_twoHundredNinetyThree\_of\_not\_dvd} & WheelK2\_293 & PROVED & \checkmark & verified \\
\texttt{K2\_307\_eq} & WheelK2\_307 & PROVED & \checkmark & verified \\
\texttt{K2\_307\_of\_not\_two\_dvd} & WheelK2\_307 & PROVED & \checkmark & verified \\
\texttt{K2\_307\_of\_two\_and\_threeHundredSeven\_dvd} & WheelK2\_307 & PROVED & \checkmark & verified \\
\texttt{Kp\_threeHundredSeven} & WheelK2\_307 & PROVED & \checkmark & verified \\
\texttt{Kp\_threeHundredSeven\_of\_dvd} & WheelK2\_307 & PROVED & \checkmark & verified \\
\texttt{Kp\_threeHundredSeven\_of\_not\_dvd} & WheelK2\_307 & PROVED & \checkmark & verified \\
\texttt{K2\_31\_eq} & WheelK2\_31 & PROVED & \checkmark & verified \\
\texttt{K2\_31\_of\_not\_two\_dvd} & WheelK2\_31 & PROVED & \checkmark & verified \\
\texttt{K2\_31\_of\_two\_and\_thirtyOne\_dvd} & WheelK2\_31 & PROVED & \checkmark & verified \\
\texttt{Kp\_thirtyOne} & WheelK2\_31 & PROVED & \checkmark & verified \\
\texttt{Kp\_thirtyOne\_of\_dvd} & WheelK2\_31 & PROVED & \checkmark & verified \\
\texttt{Kp\_thirtyOne\_of\_not\_dvd} & WheelK2\_31 & PROVED & \checkmark & verified \\
\texttt{K2\_311\_eq} & WheelK2\_311 & PROVED & \checkmark & verified \\
\texttt{K2\_311\_of\_not\_two\_dvd} & WheelK2\_311 & PROVED & \checkmark & verified \\
\texttt{K2\_311\_of\_two\_and\_threeHundredEleven\_dvd} & WheelK2\_311 & PROVED & \checkmark & verified \\
\texttt{Kp\_threeHundredEleven} & WheelK2\_311 & PROVED & \checkmark & verified \\
\texttt{Kp\_threeHundredEleven\_of\_dvd} & WheelK2\_311 & PROVED & \checkmark & verified \\
\texttt{Kp\_threeHundredEleven\_of\_not\_dvd} & WheelK2\_311 & PROVED & \checkmark & verified \\
\texttt{K2\_313\_eq} & WheelK2\_313 & PROVED & \checkmark & verified \\
\texttt{K2\_313\_of\_not\_two\_dvd} & WheelK2\_313 & PROVED & \checkmark & verified \\
\texttt{K2\_313\_of\_two\_and\_threeHundredThirteen\_dvd} & WheelK2\_313 & PROVED & \checkmark & verified \\
\texttt{Kp\_threeHundredThirteen} & WheelK2\_313 & PROVED & \checkmark & verified \\
\texttt{Kp\_threeHundredThirteen\_of\_dvd} & WheelK2\_313 & PROVED & \checkmark & verified \\
\texttt{Kp\_threeHundredThirteen\_of\_not\_dvd} & WheelK2\_313 & PROVED & \checkmark & verified \\
\texttt{K2\_317\_eq} & WheelK2\_317 & PROVED & \checkmark & verified \\
\texttt{K2\_317\_of\_not\_two\_dvd} & WheelK2\_317 & PROVED & \checkmark & verified \\
\texttt{K2\_317\_of\_two\_and\_threeHundredSeventeen\_dvd} & WheelK2\_317 & PROVED & \checkmark & verified \\
\texttt{Kp\_threeHundredSeventeen} & WheelK2\_317 & PROVED & \checkmark & verified \\
\texttt{Kp\_threeHundredSeventeen\_of\_dvd} & WheelK2\_317 & PROVED & \checkmark & verified \\
\texttt{Kp\_threeHundredSeventeen\_of\_not\_dvd} & WheelK2\_317 & PROVED & \checkmark & verified \\
\texttt{K2\_331\_eq} & WheelK2\_331 & PROVED & \checkmark & verified \\
\texttt{K2\_331\_of\_not\_two\_dvd} & WheelK2\_331 & PROVED & \checkmark & verified \\
\texttt{K2\_331\_of\_two\_and\_threeHundredThirtyOne\_dvd} & WheelK2\_331 & PROVED & \checkmark & verified \\
\texttt{Kp\_threeHundredThirtyOne} & WheelK2\_331 & PROVED & \checkmark & verified \\
\texttt{Kp\_threeHundredThirtyOne\_of\_dvd} & WheelK2\_331 & PROVED & \checkmark & verified \\
\texttt{Kp\_threeHundredThirtyOne\_of\_not\_dvd} & WheelK2\_331 & PROVED & \checkmark & verified \\
\texttt{K2\_337\_eq} & WheelK2\_337 & PROVED & \checkmark & verified \\
\texttt{K2\_337\_of\_not\_two\_dvd} & WheelK2\_337 & PROVED & \checkmark & verified \\
\texttt{K2\_337\_of\_two\_and\_threeHundredThirtySeven\_dvd} & WheelK2\_337 & PROVED & \checkmark & verified \\
\texttt{Kp\_threeHundredThirtySeven} & WheelK2\_337 & PROVED & \checkmark & verified \\
\texttt{Kp\_threeHundredThirtySeven\_of\_dvd} & WheelK2\_337 & PROVED & \checkmark & verified \\
\texttt{Kp\_threeHundredThirtySeven\_of\_not\_dvd} & WheelK2\_337 & PROVED & \checkmark & verified \\
\texttt{K2\_347\_eq} & WheelK2\_347 & PROVED & \checkmark & verified \\
\texttt{K2\_347\_of\_not\_two\_dvd} & WheelK2\_347 & PROVED & \checkmark & verified \\
\texttt{K2\_347\_of\_two\_and\_threeHundredFortySeven\_dvd} & WheelK2\_347 & PROVED & \checkmark & verified \\
\texttt{Kp\_threeHundredFortySeven} & WheelK2\_347 & PROVED & \checkmark & verified \\
\texttt{Kp\_threeHundredFortySeven\_of\_dvd} & WheelK2\_347 & PROVED & \checkmark & verified \\
\texttt{Kp\_threeHundredFortySeven\_of\_not\_dvd} & WheelK2\_347 & PROVED & \checkmark & verified \\
\texttt{K2\_349\_eq} & WheelK2\_349 & PROVED & \checkmark & verified \\
\texttt{K2\_349\_of\_not\_two\_dvd} & WheelK2\_349 & PROVED & \checkmark & verified \\
\texttt{K2\_349\_of\_two\_and\_threeHundredFortyNine\_dvd} & WheelK2\_349 & PROVED & \checkmark & verified \\
\texttt{Kp\_threeHundredFortyNine} & WheelK2\_349 & PROVED & \checkmark & verified \\
\texttt{Kp\_threeHundredFortyNine\_of\_dvd} & WheelK2\_349 & PROVED & \checkmark & verified \\
\texttt{Kp\_threeHundredFortyNine\_of\_not\_dvd} & WheelK2\_349 & PROVED & \checkmark & verified \\
\texttt{K2\_353\_eq} & WheelK2\_353 & PROVED & \checkmark & verified \\
\texttt{K2\_353\_of\_not\_two\_dvd} & WheelK2\_353 & PROVED & \checkmark & verified \\
\texttt{K2\_353\_of\_two\_and\_threeHundredFiftyThree\_dvd} & WheelK2\_353 & PROVED & \checkmark & verified \\
\texttt{Kp\_threeHundredFiftyThree} & WheelK2\_353 & PROVED & \checkmark & verified \\
\texttt{Kp\_threeHundredFiftyThree\_of\_dvd} & WheelK2\_353 & PROVED & \checkmark & verified \\
\texttt{Kp\_threeHundredFiftyThree\_of\_not\_dvd} & WheelK2\_353 & PROVED & \checkmark & verified \\
\texttt{K2\_359\_eq} & WheelK2\_359 & PROVED & \checkmark & verified \\
\texttt{K2\_359\_of\_not\_two\_dvd} & WheelK2\_359 & PROVED & \checkmark & verified \\
\texttt{K2\_359\_of\_two\_and\_threeHundredFiftyNine\_dvd} & WheelK2\_359 & PROVED & \checkmark & verified \\
\texttt{Kp\_threeHundredFiftyNine} & WheelK2\_359 & PROVED & \checkmark & verified \\
\texttt{Kp\_threeHundredFiftyNine\_of\_dvd} & WheelK2\_359 & PROVED & \checkmark & verified \\
\texttt{Kp\_threeHundredFiftyNine\_of\_not\_dvd} & WheelK2\_359 & PROVED & \checkmark & verified \\
\texttt{K2\_367\_eq} & WheelK2\_367 & PROVED & \checkmark & verified \\
\texttt{K2\_367\_of\_not\_two\_dvd} & WheelK2\_367 & PROVED & \checkmark & verified \\
\texttt{K2\_367\_of\_two\_and\_threeHundredSixtySeven\_dvd} & WheelK2\_367 & PROVED & \checkmark & verified \\
\texttt{Kp\_threeHundredSixtySeven} & WheelK2\_367 & PROVED & \checkmark & verified \\
\texttt{Kp\_threeHundredSixtySeven\_of\_dvd} & WheelK2\_367 & PROVED & \checkmark & verified \\
\texttt{Kp\_threeHundredSixtySeven\_of\_not\_dvd} & WheelK2\_367 & PROVED & \checkmark & verified \\
\texttt{K2\_37\_eq} & WheelK2\_37 & PROVED & \checkmark & verified \\
\texttt{K2\_37\_of\_not\_two\_dvd} & WheelK2\_37 & PROVED & \checkmark & verified \\
\texttt{K2\_37\_of\_two\_and\_thirtySeven\_dvd} & WheelK2\_37 & PROVED & \checkmark & verified \\
\texttt{Kp\_thirtySeven} & WheelK2\_37 & PROVED & \checkmark & verified \\
\texttt{Kp\_thirtySeven\_of\_dvd} & WheelK2\_37 & PROVED & \checkmark & verified \\
\texttt{Kp\_thirtySeven\_of\_not\_dvd} & WheelK2\_37 & PROVED & \checkmark & verified \\
\texttt{K2\_373\_eq} & WheelK2\_373 & PROVED & \checkmark & verified \\
\texttt{K2\_373\_of\_not\_two\_dvd} & WheelK2\_373 & PROVED & \checkmark & verified \\
\texttt{K2\_373\_of\_two\_and\_threeHundredSeventyThree\_dvd} & WheelK2\_373 & PROVED & \checkmark & verified \\
\texttt{Kp\_threeHundredSeventyThree} & WheelK2\_373 & PROVED & \checkmark & verified \\
\texttt{Kp\_threeHundredSeventyThree\_of\_dvd} & WheelK2\_373 & PROVED & \checkmark & verified \\
\texttt{Kp\_threeHundredSeventyThree\_of\_not\_dvd} & WheelK2\_373 & PROVED & \checkmark & verified \\
\texttt{K2\_379\_eq} & WheelK2\_379 & PROVED & \checkmark & verified \\
\texttt{K2\_379\_of\_not\_two\_dvd} & WheelK2\_379 & PROVED & \checkmark & verified \\
\texttt{K2\_379\_of\_two\_and\_threeHundredSeventyNine\_dvd} & WheelK2\_379 & PROVED & \checkmark & verified \\
\texttt{Kp\_threeHundredSeventyNine} & WheelK2\_379 & PROVED & \checkmark & verified \\
\texttt{Kp\_threeHundredSeventyNine\_of\_dvd} & WheelK2\_379 & PROVED & \checkmark & verified \\
\texttt{Kp\_threeHundredSeventyNine\_of\_not\_dvd} & WheelK2\_379 & PROVED & \checkmark & verified \\
\texttt{K2\_383\_eq} & WheelK2\_383 & PROVED & \checkmark & verified \\
\texttt{K2\_383\_of\_not\_two\_dvd} & WheelK2\_383 & PROVED & \checkmark & verified \\
\texttt{K2\_383\_of\_two\_and\_threeHundredEightyThree\_dvd} & WheelK2\_383 & PROVED & \checkmark & verified \\
\texttt{Kp\_threeHundredEightyThree} & WheelK2\_383 & PROVED & \checkmark & verified \\
\texttt{Kp\_threeHundredEightyThree\_of\_dvd} & WheelK2\_383 & PROVED & \checkmark & verified \\
\texttt{Kp\_threeHundredEightyThree\_of\_not\_dvd} & WheelK2\_383 & PROVED & \checkmark & verified \\
\texttt{K2\_389\_eq} & WheelK2\_389 & PROVED & \checkmark & verified \\
\texttt{K2\_389\_of\_not\_two\_dvd} & WheelK2\_389 & PROVED & \checkmark & verified \\
\texttt{K2\_389\_of\_two\_and\_threeHundredEightyNine\_dvd} & WheelK2\_389 & PROVED & \checkmark & verified \\
\texttt{Kp\_threeHundredEightyNine} & WheelK2\_389 & PROVED & \checkmark & verified \\
\texttt{Kp\_threeHundredEightyNine\_of\_dvd} & WheelK2\_389 & PROVED & \checkmark & verified \\
\texttt{Kp\_threeHundredEightyNine\_of\_not\_dvd} & WheelK2\_389 & PROVED & \checkmark & verified \\
\texttt{K2\_397\_eq} & WheelK2\_397 & PROVED & \checkmark & verified \\
\texttt{K2\_397\_of\_not\_two\_dvd} & WheelK2\_397 & PROVED & \checkmark & verified \\
\texttt{K2\_397\_of\_two\_and\_threeHundredNinetySeven\_dvd} & WheelK2\_397 & PROVED & \checkmark & verified \\
\texttt{Kp\_threeHundredNinetySeven} & WheelK2\_397 & PROVED & \checkmark & verified \\
\texttt{Kp\_threeHundredNinetySeven\_of\_dvd} & WheelK2\_397 & PROVED & \checkmark & verified \\
\texttt{Kp\_threeHundredNinetySeven\_of\_not\_dvd} & WheelK2\_397 & PROVED & \checkmark & verified \\
\texttt{K2\_401\_eq} & WheelK2\_401 & PROVED & \checkmark & verified \\
\texttt{K2\_401\_of\_not\_two\_dvd} & WheelK2\_401 & PROVED & \checkmark & verified \\
\texttt{K2\_401\_of\_two\_and\_fourHundredOne\_dvd} & WheelK2\_401 & PROVED & \checkmark & verified \\
\texttt{Kp\_fourHundredOne} & WheelK2\_401 & PROVED & \checkmark & verified \\
\texttt{Kp\_fourHundredOne\_of\_dvd} & WheelK2\_401 & PROVED & \checkmark & verified \\
\texttt{Kp\_fourHundredOne\_of\_not\_dvd} & WheelK2\_401 & PROVED & \checkmark & verified \\
\texttt{K2\_409\_eq} & WheelK2\_409 & PROVED & \checkmark & verified \\
\texttt{K2\_409\_of\_not\_two\_dvd} & WheelK2\_409 & PROVED & \checkmark & verified \\
\texttt{K2\_409\_of\_two\_and\_fourHundredNine\_dvd} & WheelK2\_409 & PROVED & \checkmark & verified \\
\texttt{Kp\_fourHundredNine} & WheelK2\_409 & PROVED & \checkmark & verified \\
\texttt{Kp\_fourHundredNine\_of\_dvd} & WheelK2\_409 & PROVED & \checkmark & verified \\
\texttt{Kp\_fourHundredNine\_of\_not\_dvd} & WheelK2\_409 & PROVED & \checkmark & verified \\
\texttt{K2\_41\_eq} & WheelK2\_41 & PROVED & \checkmark & verified \\
\texttt{K2\_41\_of\_not\_two\_dvd} & WheelK2\_41 & PROVED & \checkmark & verified \\
\texttt{K2\_41\_of\_two\_and\_fortyOne\_dvd} & WheelK2\_41 & PROVED & \checkmark & verified \\
\texttt{Kp\_fortyOne} & WheelK2\_41 & PROVED & \checkmark & verified \\
\texttt{Kp\_fortyOne\_of\_dvd} & WheelK2\_41 & PROVED & \checkmark & verified \\
\texttt{Kp\_fortyOne\_of\_not\_dvd} & WheelK2\_41 & PROVED & \checkmark & verified \\
\texttt{K2\_419\_eq} & WheelK2\_419 & PROVED & \checkmark & verified \\
\texttt{K2\_419\_of\_not\_two\_dvd} & WheelK2\_419 & PROVED & \checkmark & verified \\
\texttt{K2\_419\_of\_two\_and\_fourHundredNineteen\_dvd} & WheelK2\_419 & PROVED & \checkmark & verified \\
\texttt{Kp\_fourHundredNineteen} & WheelK2\_419 & PROVED & \checkmark & verified \\
\texttt{Kp\_fourHundredNineteen\_of\_dvd} & WheelK2\_419 & PROVED & \checkmark & verified \\
\texttt{Kp\_fourHundredNineteen\_of\_not\_dvd} & WheelK2\_419 & PROVED & \checkmark & verified \\
\texttt{K2\_421\_eq} & WheelK2\_421 & PROVED & \checkmark & verified \\
\texttt{K2\_421\_of\_not\_two\_dvd} & WheelK2\_421 & PROVED & \checkmark & verified \\
\texttt{K2\_421\_of\_two\_and\_fourHundredTwentyOne\_dvd} & WheelK2\_421 & PROVED & \checkmark & verified \\
\texttt{Kp\_fourHundredTwentyOne} & WheelK2\_421 & PROVED & \checkmark & verified \\
\texttt{Kp\_fourHundredTwentyOne\_of\_dvd} & WheelK2\_421 & PROVED & \checkmark & verified \\
\texttt{Kp\_fourHundredTwentyOne\_of\_not\_dvd} & WheelK2\_421 & PROVED & \checkmark & verified \\
\texttt{K2\_43\_eq} & WheelK2\_43 & PROVED & \checkmark & verified \\
\texttt{K2\_43\_of\_not\_two\_dvd} & WheelK2\_43 & PROVED & \checkmark & verified \\
\texttt{K2\_43\_of\_two\_and\_fortyThree\_dvd} & WheelK2\_43 & PROVED & \checkmark & verified \\
\texttt{Kp\_fortyThree} & WheelK2\_43 & PROVED & \checkmark & verified \\
\texttt{Kp\_fortyThree\_of\_dvd} & WheelK2\_43 & PROVED & \checkmark & verified \\
\texttt{Kp\_fortyThree\_of\_not\_dvd} & WheelK2\_43 & PROVED & \checkmark & verified \\
\texttt{K2\_431\_eq} & WheelK2\_431 & PROVED & \checkmark & verified \\
\texttt{K2\_431\_of\_not\_two\_dvd} & WheelK2\_431 & PROVED & \checkmark & verified \\
\texttt{K2\_431\_of\_two\_and\_fourHundredThirtyOne\_dvd} & WheelK2\_431 & PROVED & \checkmark & verified \\
\texttt{Kp\_fourHundredThirtyOne} & WheelK2\_431 & PROVED & \checkmark & verified \\
\texttt{Kp\_fourHundredThirtyOne\_of\_dvd} & WheelK2\_431 & PROVED & \checkmark & verified \\
\texttt{Kp\_fourHundredThirtyOne\_of\_not\_dvd} & WheelK2\_431 & PROVED & \checkmark & verified \\
\texttt{K2\_433\_eq} & WheelK2\_433 & PROVED & \checkmark & verified \\
\texttt{K2\_433\_of\_not\_two\_dvd} & WheelK2\_433 & PROVED & \checkmark & verified \\
\texttt{K2\_433\_of\_two\_and\_fourHundredThirtyThree\_dvd} & WheelK2\_433 & PROVED & \checkmark & verified \\
\texttt{Kp\_fourHundredThirtyThree} & WheelK2\_433 & PROVED & \checkmark & verified \\
\texttt{Kp\_fourHundredThirtyThree\_of\_dvd} & WheelK2\_433 & PROVED & \checkmark & verified \\
\texttt{Kp\_fourHundredThirtyThree\_of\_not\_dvd} & WheelK2\_433 & PROVED & \checkmark & verified \\
\texttt{K2\_439\_eq} & WheelK2\_439 & PROVED & \checkmark & verified \\
\texttt{K2\_439\_of\_not\_two\_dvd} & WheelK2\_439 & PROVED & \checkmark & verified \\
\texttt{K2\_439\_of\_two\_and\_fourHundredThirtyNine\_dvd} & WheelK2\_439 & PROVED & \checkmark & verified \\
\texttt{Kp\_fourHundredThirtyNine} & WheelK2\_439 & PROVED & \checkmark & verified \\
\texttt{Kp\_fourHundredThirtyNine\_of\_dvd} & WheelK2\_439 & PROVED & \checkmark & verified \\
\texttt{Kp\_fourHundredThirtyNine\_of\_not\_dvd} & WheelK2\_439 & PROVED & \checkmark & verified \\
\texttt{K2\_443\_eq} & WheelK2\_443 & PROVED & \checkmark & verified \\
\texttt{K2\_443\_of\_not\_two\_dvd} & WheelK2\_443 & PROVED & \checkmark & verified \\
\texttt{K2\_443\_of\_two\_and\_fourHundredFortyThree\_dvd} & WheelK2\_443 & PROVED & \checkmark & verified \\
\texttt{Kp\_fourHundredFortyThree} & WheelK2\_443 & PROVED & \checkmark & verified \\
\texttt{Kp\_fourHundredFortyThree\_of\_dvd} & WheelK2\_443 & PROVED & \checkmark & verified \\
\texttt{Kp\_fourHundredFortyThree\_of\_not\_dvd} & WheelK2\_443 & PROVED & \checkmark & verified \\
\texttt{K2\_449\_eq} & WheelK2\_449 & PROVED & \checkmark & verified \\
\texttt{K2\_449\_of\_not\_two\_dvd} & WheelK2\_449 & PROVED & \checkmark & verified \\
\texttt{K2\_449\_of\_two\_and\_fourHundredFortyNine\_dvd} & WheelK2\_449 & PROVED & \checkmark & verified \\
\texttt{Kp\_fourHundredFortyNine} & WheelK2\_449 & PROVED & \checkmark & verified \\
\texttt{Kp\_fourHundredFortyNine\_of\_dvd} & WheelK2\_449 & PROVED & \checkmark & verified \\
\texttt{Kp\_fourHundredFortyNine\_of\_not\_dvd} & WheelK2\_449 & PROVED & \checkmark & verified \\
\texttt{K2\_457\_eq} & WheelK2\_457 & PROVED & \checkmark & verified \\
\texttt{K2\_457\_of\_not\_two\_dvd} & WheelK2\_457 & PROVED & \checkmark & verified \\
\texttt{K2\_457\_of\_two\_and\_fourHundredFiftySeven\_dvd} & WheelK2\_457 & PROVED & \checkmark & verified \\
\texttt{Kp\_fourHundredFiftySeven} & WheelK2\_457 & PROVED & \checkmark & verified \\
\texttt{Kp\_fourHundredFiftySeven\_of\_dvd} & WheelK2\_457 & PROVED & \checkmark & verified \\
\texttt{Kp\_fourHundredFiftySeven\_of\_not\_dvd} & WheelK2\_457 & PROVED & \checkmark & verified \\
\texttt{K2\_461\_eq} & WheelK2\_461 & PROVED & \checkmark & verified \\
\texttt{K2\_461\_of\_not\_two\_dvd} & WheelK2\_461 & PROVED & \checkmark & verified \\
\texttt{K2\_461\_of\_two\_and\_fourHundredSixtyOne\_dvd} & WheelK2\_461 & PROVED & \checkmark & verified \\
\texttt{Kp\_fourHundredSixtyOne} & WheelK2\_461 & PROVED & \checkmark & verified \\
\texttt{Kp\_fourHundredSixtyOne\_of\_dvd} & WheelK2\_461 & PROVED & \checkmark & verified \\
\texttt{Kp\_fourHundredSixtyOne\_of\_not\_dvd} & WheelK2\_461 & PROVED & \checkmark & verified \\
\texttt{K2\_463\_eq} & WheelK2\_463 & PROVED & \checkmark & verified \\
\texttt{K2\_463\_of\_not\_two\_dvd} & WheelK2\_463 & PROVED & \checkmark & verified \\
\texttt{K2\_463\_of\_two\_and\_fourHundredSixtyThree\_dvd} & WheelK2\_463 & PROVED & \checkmark & verified \\
\texttt{Kp\_fourHundredSixtyThree} & WheelK2\_463 & PROVED & \checkmark & verified \\
\texttt{Kp\_fourHundredSixtyThree\_of\_dvd} & WheelK2\_463 & PROVED & \checkmark & verified \\
\texttt{Kp\_fourHundredSixtyThree\_of\_not\_dvd} & WheelK2\_463 & PROVED & \checkmark & verified \\
\texttt{K2\_467\_eq} & WheelK2\_467 & PROVED & \checkmark & verified \\
\texttt{K2\_467\_of\_not\_two\_dvd} & WheelK2\_467 & PROVED & \checkmark & verified \\
\texttt{K2\_467\_of\_two\_and\_fourHundredSixtySeven\_dvd} & WheelK2\_467 & PROVED & \checkmark & verified \\
\texttt{Kp\_fourHundredSixtySeven} & WheelK2\_467 & PROVED & \checkmark & verified \\
\texttt{Kp\_fourHundredSixtySeven\_of\_dvd} & WheelK2\_467 & PROVED & \checkmark & verified \\
\texttt{Kp\_fourHundredSixtySeven\_of\_not\_dvd} & WheelK2\_467 & PROVED & \checkmark & verified \\
\texttt{K2\_47\_eq} & WheelK2\_47 & PROVED & \checkmark & verified \\
\texttt{K2\_47\_of\_not\_two\_dvd} & WheelK2\_47 & PROVED & \checkmark & verified \\
\texttt{K2\_47\_of\_two\_and\_fortySeven\_dvd} & WheelK2\_47 & PROVED & \checkmark & verified \\
\texttt{Kp\_fortySeven} & WheelK2\_47 & PROVED & \checkmark & verified \\
\texttt{Kp\_fortySeven\_of\_dvd} & WheelK2\_47 & PROVED & \checkmark & verified \\
\texttt{Kp\_fortySeven\_of\_not\_dvd} & WheelK2\_47 & PROVED & \checkmark & verified \\
\texttt{K2\_479\_eq} & WheelK2\_479 & PROVED & \checkmark & verified \\
\texttt{K2\_479\_of\_not\_two\_dvd} & WheelK2\_479 & PROVED & \checkmark & verified \\
\texttt{K2\_479\_of\_two\_and\_fourHundredSeventyNine\_dvd} & WheelK2\_479 & PROVED & \checkmark & verified \\
\texttt{Kp\_fourHundredSeventyNine} & WheelK2\_479 & PROVED & \checkmark & verified \\
\texttt{Kp\_fourHundredSeventyNine\_of\_dvd} & WheelK2\_479 & PROVED & \checkmark & verified \\
\texttt{Kp\_fourHundredSeventyNine\_of\_not\_dvd} & WheelK2\_479 & PROVED & \checkmark & verified \\
\texttt{K2\_487\_eq} & WheelK2\_487 & PROVED & \checkmark & verified \\
\texttt{K2\_487\_of\_not\_two\_dvd} & WheelK2\_487 & PROVED & \checkmark & verified \\
\texttt{K2\_487\_of\_two\_and\_fourHundredEightySeven\_dvd} & WheelK2\_487 & PROVED & \checkmark & verified \\
\texttt{Kp\_fourHundredEightySeven} & WheelK2\_487 & PROVED & \checkmark & verified \\
\texttt{Kp\_fourHundredEightySeven\_of\_dvd} & WheelK2\_487 & PROVED & \checkmark & verified \\
\texttt{Kp\_fourHundredEightySeven\_of\_not\_dvd} & WheelK2\_487 & PROVED & \checkmark & verified \\
\texttt{K2\_491\_eq} & WheelK2\_491 & PROVED & \checkmark & verified \\
\texttt{K2\_491\_of\_not\_two\_dvd} & WheelK2\_491 & PROVED & \checkmark & verified \\
\texttt{K2\_491\_of\_two\_and\_fourHundredNinetyOne\_dvd} & WheelK2\_491 & PROVED & \checkmark & verified \\
\texttt{Kp\_fourHundredNinetyOne} & WheelK2\_491 & PROVED & \checkmark & verified \\
\texttt{Kp\_fourHundredNinetyOne\_of\_dvd} & WheelK2\_491 & PROVED & \checkmark & verified \\
\texttt{Kp\_fourHundredNinetyOne\_of\_not\_dvd} & WheelK2\_491 & PROVED & \checkmark & verified \\
\texttt{K2\_499\_eq} & WheelK2\_499 & PROVED & \checkmark & verified \\
\texttt{K2\_499\_of\_not\_two\_dvd} & WheelK2\_499 & PROVED & \checkmark & verified \\
\texttt{K2\_499\_of\_two\_and\_fourHundredNinetyNine\_dvd} & WheelK2\_499 & PROVED & \checkmark & verified \\
\texttt{Kp\_fourHundredNinetyNine} & WheelK2\_499 & PROVED & \checkmark & verified \\
\texttt{Kp\_fourHundredNinetyNine\_of\_dvd} & WheelK2\_499 & PROVED & \checkmark & verified \\
\texttt{Kp\_fourHundredNinetyNine\_of\_not\_dvd} & WheelK2\_499 & PROVED & \checkmark & verified \\
\texttt{K2\_503\_eq} & WheelK2\_503 & PROVED & \checkmark & verified \\
\texttt{K2\_503\_of\_not\_two\_dvd} & WheelK2\_503 & PROVED & \checkmark & verified \\
\texttt{K2\_503\_of\_two\_and\_fiveHundredThree\_dvd} & WheelK2\_503 & PROVED & \checkmark & verified \\
\texttt{Kp\_fiveHundredThree} & WheelK2\_503 & PROVED & \checkmark & verified \\
\texttt{Kp\_fiveHundredThree\_of\_dvd} & WheelK2\_503 & PROVED & \checkmark & verified \\
\texttt{Kp\_fiveHundredThree\_of\_not\_dvd} & WheelK2\_503 & PROVED & \checkmark & verified \\
\texttt{K2\_509\_eq} & WheelK2\_509 & PROVED & \checkmark & verified \\
\texttt{K2\_509\_of\_not\_two\_dvd} & WheelK2\_509 & PROVED & \checkmark & verified \\
\texttt{K2\_509\_of\_two\_and\_fiveHundredNine\_dvd} & WheelK2\_509 & PROVED & \checkmark & verified \\
\texttt{Kp\_fiveHundredNine} & WheelK2\_509 & PROVED & \checkmark & verified \\
\texttt{Kp\_fiveHundredNine\_of\_dvd} & WheelK2\_509 & PROVED & \checkmark & verified \\
\texttt{Kp\_fiveHundredNine\_of\_not\_dvd} & WheelK2\_509 & PROVED & \checkmark & verified \\
\texttt{K2\_521\_eq} & WheelK2\_521 & PROVED & \checkmark & verified \\
\texttt{K2\_521\_of\_not\_two\_dvd} & WheelK2\_521 & PROVED & \checkmark & verified \\
\texttt{K2\_521\_of\_two\_and\_fiveHundredTwentyOne\_dvd} & WheelK2\_521 & PROVED & \checkmark & verified \\
\texttt{Kp\_fiveHundredTwentyOne} & WheelK2\_521 & PROVED & \checkmark & verified \\
\texttt{Kp\_fiveHundredTwentyOne\_of\_dvd} & WheelK2\_521 & PROVED & \checkmark & verified \\
\texttt{Kp\_fiveHundredTwentyOne\_of\_not\_dvd} & WheelK2\_521 & PROVED & \checkmark & verified \\
\texttt{K2\_523\_eq} & WheelK2\_523 & PROVED & \checkmark & verified \\
\texttt{K2\_523\_of\_not\_two\_dvd} & WheelK2\_523 & PROVED & \checkmark & verified \\
\texttt{K2\_523\_of\_two\_and\_fiveHundredTwentyThree\_dvd} & WheelK2\_523 & PROVED & \checkmark & verified \\
\texttt{Kp\_fiveHundredTwentyThree} & WheelK2\_523 & PROVED & \checkmark & verified \\
\texttt{Kp\_fiveHundredTwentyThree\_of\_dvd} & WheelK2\_523 & PROVED & \checkmark & verified \\
\texttt{Kp\_fiveHundredTwentyThree\_of\_not\_dvd} & WheelK2\_523 & PROVED & \checkmark & verified \\
\texttt{K2\_53\_eq} & WheelK2\_53 & PROVED & \checkmark & verified \\
\texttt{K2\_53\_of\_not\_two\_dvd} & WheelK2\_53 & PROVED & \checkmark & verified \\
\texttt{K2\_53\_of\_two\_and\_fiftyThree\_dvd} & WheelK2\_53 & PROVED & \checkmark & verified \\
\texttt{Kp\_fiftyThree} & WheelK2\_53 & PROVED & \checkmark & verified \\
\texttt{Kp\_fiftyThree\_of\_dvd} & WheelK2\_53 & PROVED & \checkmark & verified \\
\texttt{Kp\_fiftyThree\_of\_not\_dvd} & WheelK2\_53 & PROVED & \checkmark & verified \\
\texttt{K2\_541\_eq} & WheelK2\_541 & PROVED & \checkmark & verified \\
\texttt{K2\_541\_of\_not\_two\_dvd} & WheelK2\_541 & PROVED & \checkmark & verified \\
\texttt{K2\_541\_of\_two\_and\_fiveHundredFortyOne\_dvd} & WheelK2\_541 & PROVED & \checkmark & verified \\
\texttt{Kp\_fiveHundredFortyOne} & WheelK2\_541 & PROVED & \checkmark & verified \\
\texttt{Kp\_fiveHundredFortyOne\_of\_dvd} & WheelK2\_541 & PROVED & \checkmark & verified \\
\texttt{Kp\_fiveHundredFortyOne\_of\_not\_dvd} & WheelK2\_541 & PROVED & \checkmark & verified \\
\texttt{K2\_547\_eq} & WheelK2\_547 & PROVED & \checkmark & verified \\
\texttt{K2\_547\_of\_not\_two\_dvd} & WheelK2\_547 & PROVED & \checkmark & verified \\
\texttt{K2\_547\_of\_two\_and\_fiveHundredFortySeven\_dvd} & WheelK2\_547 & PROVED & \checkmark & verified \\
\texttt{Kp\_fiveHundredFortySeven} & WheelK2\_547 & PROVED & \checkmark & verified \\
\texttt{Kp\_fiveHundredFortySeven\_of\_dvd} & WheelK2\_547 & PROVED & \checkmark & verified \\
\texttt{Kp\_fiveHundredFortySeven\_of\_not\_dvd} & WheelK2\_547 & PROVED & \checkmark & verified \\
\texttt{K2\_557\_eq} & WheelK2\_557 & PROVED & \checkmark & verified \\
\texttt{K2\_557\_of\_not\_two\_dvd} & WheelK2\_557 & PROVED & \checkmark & verified \\
\texttt{K2\_557\_of\_two\_and\_fiveHundredFiftySeven\_dvd} & WheelK2\_557 & PROVED & \checkmark & verified \\
\texttt{Kp\_fiveHundredFiftySeven} & WheelK2\_557 & PROVED & \checkmark & verified \\
\texttt{Kp\_fiveHundredFiftySeven\_of\_dvd} & WheelK2\_557 & PROVED & \checkmark & verified \\
\texttt{Kp\_fiveHundredFiftySeven\_of\_not\_dvd} & WheelK2\_557 & PROVED & \checkmark & verified \\
\texttt{K2\_563\_eq} & WheelK2\_563 & PROVED & \checkmark & verified \\
\texttt{K2\_563\_of\_not\_two\_dvd} & WheelK2\_563 & PROVED & \checkmark & verified \\
\texttt{K2\_563\_of\_two\_and\_fiveHundredSixtyThree\_dvd} & WheelK2\_563 & PROVED & \checkmark & verified \\
\texttt{Kp\_fiveHundredSixtyThree} & WheelK2\_563 & PROVED & \checkmark & verified \\
\texttt{Kp\_fiveHundredSixtyThree\_of\_dvd} & WheelK2\_563 & PROVED & \checkmark & verified \\
\texttt{Kp\_fiveHundredSixtyThree\_of\_not\_dvd} & WheelK2\_563 & PROVED & \checkmark & verified \\
\texttt{K2\_569\_eq} & WheelK2\_569 & PROVED & \checkmark & verified \\
\texttt{K2\_569\_of\_not\_two\_dvd} & WheelK2\_569 & PROVED & \checkmark & verified \\
\texttt{K2\_569\_of\_two\_and\_fiveHundredSixtyNine\_dvd} & WheelK2\_569 & PROVED & \checkmark & verified \\
\texttt{Kp\_fiveHundredSixtyNine} & WheelK2\_569 & PROVED & \checkmark & verified \\
\texttt{Kp\_fiveHundredSixtyNine\_of\_dvd} & WheelK2\_569 & PROVED & \checkmark & verified \\
\texttt{Kp\_fiveHundredSixtyNine\_of\_not\_dvd} & WheelK2\_569 & PROVED & \checkmark & verified \\
\texttt{K2\_571\_eq} & WheelK2\_571 & PROVED & \checkmark & verified \\
\texttt{K2\_571\_of\_not\_two\_dvd} & WheelK2\_571 & PROVED & \checkmark & verified \\
\texttt{K2\_571\_of\_two\_and\_fiveHundredSeventyOne\_dvd} & WheelK2\_571 & PROVED & \checkmark & verified \\
\texttt{Kp\_fiveHundredSeventyOne} & WheelK2\_571 & PROVED & \checkmark & verified \\
\texttt{Kp\_fiveHundredSeventyOne\_of\_dvd} & WheelK2\_571 & PROVED & \checkmark & verified \\
\texttt{Kp\_fiveHundredSeventyOne\_of\_not\_dvd} & WheelK2\_571 & PROVED & \checkmark & verified \\
\texttt{K2\_577\_eq} & WheelK2\_577 & PROVED & \checkmark & verified \\
\texttt{K2\_577\_of\_not\_two\_dvd} & WheelK2\_577 & PROVED & \checkmark & verified \\
\texttt{K2\_577\_of\_two\_and\_fiveHundredSeventySeven\_dvd} & WheelK2\_577 & PROVED & \checkmark & verified \\
\texttt{Kp\_fiveHundredSeventySeven} & WheelK2\_577 & PROVED & \checkmark & verified \\
\texttt{Kp\_fiveHundredSeventySeven\_of\_dvd} & WheelK2\_577 & PROVED & \checkmark & verified \\
\texttt{Kp\_fiveHundredSeventySeven\_of\_not\_dvd} & WheelK2\_577 & PROVED & \checkmark & verified \\
\texttt{K2\_587\_eq} & WheelK2\_587 & PROVED & \checkmark & verified \\
\texttt{K2\_587\_of\_not\_two\_dvd} & WheelK2\_587 & PROVED & \checkmark & verified \\
\texttt{K2\_587\_of\_two\_and\_fiveHundredEightySeven\_dvd} & WheelK2\_587 & PROVED & \checkmark & verified \\
\texttt{Kp\_fiveHundredEightySeven} & WheelK2\_587 & PROVED & \checkmark & verified \\
\texttt{Kp\_fiveHundredEightySeven\_of\_dvd} & WheelK2\_587 & PROVED & \checkmark & verified \\
\texttt{Kp\_fiveHundredEightySeven\_of\_not\_dvd} & WheelK2\_587 & PROVED & \checkmark & verified \\
\texttt{K2\_59\_eq} & WheelK2\_59 & PROVED & \checkmark & verified \\
\texttt{K2\_59\_of\_not\_two\_dvd} & WheelK2\_59 & PROVED & \checkmark & verified \\
\texttt{K2\_59\_of\_two\_and\_fiftyNine\_dvd} & WheelK2\_59 & PROVED & \checkmark & verified \\
\texttt{Kp\_fiftyNine} & WheelK2\_59 & PROVED & \checkmark & verified \\
\texttt{Kp\_fiftyNine\_of\_dvd} & WheelK2\_59 & PROVED & \checkmark & verified \\
\texttt{Kp\_fiftyNine\_of\_not\_dvd} & WheelK2\_59 & PROVED & \checkmark & verified \\
\texttt{K2\_593\_eq} & WheelK2\_593 & PROVED & \checkmark & verified \\
\texttt{K2\_593\_of\_not\_two\_dvd} & WheelK2\_593 & PROVED & \checkmark & verified \\
\texttt{K2\_593\_of\_two\_and\_fiveHundredNinetyThree\_dvd} & WheelK2\_593 & PROVED & \checkmark & verified \\
\texttt{Kp\_fiveHundredNinetyThree} & WheelK2\_593 & PROVED & \checkmark & verified \\
\texttt{Kp\_fiveHundredNinetyThree\_of\_dvd} & WheelK2\_593 & PROVED & \checkmark & verified \\
\texttt{Kp\_fiveHundredNinetyThree\_of\_not\_dvd} & WheelK2\_593 & PROVED & \checkmark & verified \\
\texttt{K2\_599\_eq} & WheelK2\_599 & PROVED & \checkmark & verified \\
\texttt{K2\_599\_of\_not\_two\_dvd} & WheelK2\_599 & PROVED & \checkmark & verified \\
\texttt{K2\_599\_of\_two\_and\_fiveHundredNinetyNine\_dvd} & WheelK2\_599 & PROVED & \checkmark & verified \\
\texttt{Kp\_fiveHundredNinetyNine} & WheelK2\_599 & PROVED & \checkmark & verified \\
\texttt{Kp\_fiveHundredNinetyNine\_of\_dvd} & WheelK2\_599 & PROVED & \checkmark & verified \\
\texttt{Kp\_fiveHundredNinetyNine\_of\_not\_dvd} & WheelK2\_599 & PROVED & \checkmark & verified \\
\texttt{K2\_601\_eq} & WheelK2\_601 & PROVED & \checkmark & verified \\
\texttt{K2\_601\_of\_not\_two\_dvd} & WheelK2\_601 & PROVED & \checkmark & verified \\
\texttt{K2\_601\_of\_two\_and\_sixHundredOne\_dvd} & WheelK2\_601 & PROVED & \checkmark & verified \\
\texttt{Kp\_sixHundredOne} & WheelK2\_601 & PROVED & \checkmark & verified \\
\texttt{Kp\_sixHundredOne\_of\_dvd} & WheelK2\_601 & PROVED & \checkmark & verified \\
\texttt{Kp\_sixHundredOne\_of\_not\_dvd} & WheelK2\_601 & PROVED & \checkmark & verified \\
\texttt{K2\_607\_eq} & WheelK2\_607 & PROVED & \checkmark & verified \\
\texttt{K2\_607\_of\_not\_two\_dvd} & WheelK2\_607 & PROVED & \checkmark & verified \\
\texttt{K2\_607\_of\_two\_and\_sixHundredSeven\_dvd} & WheelK2\_607 & PROVED & \checkmark & verified \\
\texttt{Kp\_sixHundredSeven} & WheelK2\_607 & PROVED & \checkmark & verified \\
\texttt{Kp\_sixHundredSeven\_of\_dvd} & WheelK2\_607 & PROVED & \checkmark & verified \\
\texttt{Kp\_sixHundredSeven\_of\_not\_dvd} & WheelK2\_607 & PROVED & \checkmark & verified \\
\texttt{K2\_61\_eq} & WheelK2\_61 & PROVED & \checkmark & verified \\
\texttt{K2\_61\_of\_not\_two\_dvd} & WheelK2\_61 & PROVED & \checkmark & verified \\
\texttt{K2\_61\_of\_two\_and\_sixtyOne\_dvd} & WheelK2\_61 & PROVED & \checkmark & verified \\
\texttt{Kp\_sixtyOne} & WheelK2\_61 & PROVED & \checkmark & verified \\
\texttt{Kp\_sixtyOne\_of\_dvd} & WheelK2\_61 & PROVED & \checkmark & verified \\
\texttt{Kp\_sixtyOne\_of\_not\_dvd} & WheelK2\_61 & PROVED & \checkmark & verified \\
\texttt{K2\_613\_eq} & WheelK2\_613 & PROVED & \checkmark & verified \\
\texttt{K2\_613\_of\_not\_two\_dvd} & WheelK2\_613 & PROVED & \checkmark & verified \\
\texttt{K2\_613\_of\_two\_and\_sixHundredThirteen\_dvd} & WheelK2\_613 & PROVED & \checkmark & verified \\
\texttt{Kp\_sixHundredThirteen} & WheelK2\_613 & PROVED & \checkmark & verified \\
\texttt{Kp\_sixHundredThirteen\_of\_dvd} & WheelK2\_613 & PROVED & \checkmark & verified \\
\texttt{Kp\_sixHundredThirteen\_of\_not\_dvd} & WheelK2\_613 & PROVED & \checkmark & verified \\
\texttt{K2\_617\_eq} & WheelK2\_617 & PROVED & \checkmark & verified \\
\texttt{K2\_617\_of\_not\_two\_dvd} & WheelK2\_617 & PROVED & \checkmark & verified \\
\texttt{K2\_617\_of\_two\_and\_sixHundredSeventeen\_dvd} & WheelK2\_617 & PROVED & \checkmark & verified \\
\texttt{Kp\_sixHundredSeventeen} & WheelK2\_617 & PROVED & \checkmark & verified \\
\texttt{Kp\_sixHundredSeventeen\_of\_dvd} & WheelK2\_617 & PROVED & \checkmark & verified \\
\texttt{Kp\_sixHundredSeventeen\_of\_not\_dvd} & WheelK2\_617 & PROVED & \checkmark & verified \\
\texttt{K2\_619\_eq} & WheelK2\_619 & PROVED & \checkmark & verified \\
\texttt{K2\_619\_of\_not\_two\_dvd} & WheelK2\_619 & PROVED & \checkmark & verified \\
\texttt{K2\_619\_of\_two\_and\_sixHundredNineteen\_dvd} & WheelK2\_619 & PROVED & \checkmark & verified \\
\texttt{Kp\_sixHundredNineteen} & WheelK2\_619 & PROVED & \checkmark & verified \\
\texttt{Kp\_sixHundredNineteen\_of\_dvd} & WheelK2\_619 & PROVED & \checkmark & verified \\
\texttt{Kp\_sixHundredNineteen\_of\_not\_dvd} & WheelK2\_619 & PROVED & \checkmark & verified \\
\texttt{K2\_631\_eq} & WheelK2\_631 & PROVED & \checkmark & verified \\
\texttt{K2\_631\_of\_not\_two\_dvd} & WheelK2\_631 & PROVED & \checkmark & verified \\
\texttt{K2\_631\_of\_two\_and\_sixHundredThirtyOne\_dvd} & WheelK2\_631 & PROVED & \checkmark & verified \\
\texttt{Kp\_sixHundredThirtyOne} & WheelK2\_631 & PROVED & \checkmark & verified \\
\texttt{Kp\_sixHundredThirtyOne\_of\_dvd} & WheelK2\_631 & PROVED & \checkmark & verified \\
\texttt{Kp\_sixHundredThirtyOne\_of\_not\_dvd} & WheelK2\_631 & PROVED & \checkmark & verified \\
\texttt{K2\_641\_eq} & WheelK2\_641 & PROVED & \checkmark & verified \\
\texttt{K2\_641\_of\_not\_two\_dvd} & WheelK2\_641 & PROVED & \checkmark & verified \\
\texttt{K2\_641\_of\_two\_and\_sixHundredFortyOne\_dvd} & WheelK2\_641 & PROVED & \checkmark & verified \\
\texttt{Kp\_sixHundredFortyOne} & WheelK2\_641 & PROVED & \checkmark & verified \\
\texttt{Kp\_sixHundredFortyOne\_of\_dvd} & WheelK2\_641 & PROVED & \checkmark & verified \\
\texttt{Kp\_sixHundredFortyOne\_of\_not\_dvd} & WheelK2\_641 & PROVED & \checkmark & verified \\
\texttt{K2\_643\_eq} & WheelK2\_643 & PROVED & \checkmark & verified \\
\texttt{K2\_643\_of\_not\_two\_dvd} & WheelK2\_643 & PROVED & \checkmark & verified \\
\texttt{K2\_643\_of\_two\_and\_sixHundredFortyThree\_dvd} & WheelK2\_643 & PROVED & \checkmark & verified \\
\texttt{Kp\_sixHundredFortyThree} & WheelK2\_643 & PROVED & \checkmark & verified \\
\texttt{Kp\_sixHundredFortyThree\_of\_dvd} & WheelK2\_643 & PROVED & \checkmark & verified \\
\texttt{Kp\_sixHundredFortyThree\_of\_not\_dvd} & WheelK2\_643 & PROVED & \checkmark & verified \\
\texttt{K2\_647\_eq} & WheelK2\_647 & PROVED & \checkmark & verified \\
\texttt{K2\_647\_of\_not\_two\_dvd} & WheelK2\_647 & PROVED & \checkmark & verified \\
\texttt{K2\_647\_of\_two\_and\_sixHundredFortySeven\_dvd} & WheelK2\_647 & PROVED & \checkmark & verified \\
\texttt{Kp\_sixHundredFortySeven} & WheelK2\_647 & PROVED & \checkmark & verified \\
\texttt{Kp\_sixHundredFortySeven\_of\_dvd} & WheelK2\_647 & PROVED & \checkmark & verified \\
\texttt{Kp\_sixHundredFortySeven\_of\_not\_dvd} & WheelK2\_647 & PROVED & \checkmark & verified \\
\texttt{K2\_653\_eq} & WheelK2\_653 & PROVED & \checkmark & verified \\
\texttt{K2\_653\_of\_not\_two\_dvd} & WheelK2\_653 & PROVED & \checkmark & verified \\
\texttt{K2\_653\_of\_two\_and\_sixHundredFiftyThree\_dvd} & WheelK2\_653 & PROVED & \checkmark & verified \\
\texttt{Kp\_sixHundredFiftyThree} & WheelK2\_653 & PROVED & \checkmark & verified \\
\texttt{Kp\_sixHundredFiftyThree\_of\_dvd} & WheelK2\_653 & PROVED & \checkmark & verified \\
\texttt{Kp\_sixHundredFiftyThree\_of\_not\_dvd} & WheelK2\_653 & PROVED & \checkmark & verified \\
\texttt{K2\_659\_eq} & WheelK2\_659 & PROVED & \checkmark & verified \\
\texttt{K2\_659\_of\_not\_two\_dvd} & WheelK2\_659 & PROVED & \checkmark & verified \\
\texttt{K2\_659\_of\_two\_and\_sixHundredFiftyNine\_dvd} & WheelK2\_659 & PROVED & \checkmark & verified \\
\texttt{Kp\_sixHundredFiftyNine} & WheelK2\_659 & PROVED & \checkmark & verified \\
\texttt{Kp\_sixHundredFiftyNine\_of\_dvd} & WheelK2\_659 & PROVED & \checkmark & verified \\
\texttt{Kp\_sixHundredFiftyNine\_of\_not\_dvd} & WheelK2\_659 & PROVED & \checkmark & verified \\
\texttt{K2\_661\_eq} & WheelK2\_661 & PROVED & \checkmark & verified \\
\texttt{K2\_661\_of\_not\_two\_dvd} & WheelK2\_661 & PROVED & \checkmark & verified \\
\texttt{K2\_661\_of\_two\_and\_sixHundredSixtyOne\_dvd} & WheelK2\_661 & PROVED & \checkmark & verified \\
\texttt{Kp\_sixHundredSixtyOne} & WheelK2\_661 & PROVED & \checkmark & verified \\
\texttt{Kp\_sixHundredSixtyOne\_of\_dvd} & WheelK2\_661 & PROVED & \checkmark & verified \\
\texttt{Kp\_sixHundredSixtyOne\_of\_not\_dvd} & WheelK2\_661 & PROVED & \checkmark & verified \\
\texttt{K2\_67\_eq} & WheelK2\_67 & PROVED & \checkmark & verified \\
\texttt{K2\_67\_of\_not\_two\_dvd} & WheelK2\_67 & PROVED & \checkmark & verified \\
\texttt{K2\_67\_of\_two\_and\_sixtySeven\_dvd} & WheelK2\_67 & PROVED & \checkmark & verified \\
\texttt{Kp\_sixtySeven} & WheelK2\_67 & PROVED & \checkmark & verified \\
\texttt{Kp\_sixtySeven\_of\_dvd} & WheelK2\_67 & PROVED & \checkmark & verified \\
\texttt{Kp\_sixtySeven\_of\_not\_dvd} & WheelK2\_67 & PROVED & \checkmark & verified \\
\texttt{K2\_673\_eq} & WheelK2\_673 & PROVED & \checkmark & verified \\
\texttt{K2\_673\_of\_not\_two\_dvd} & WheelK2\_673 & PROVED & \checkmark & verified \\
\texttt{K2\_673\_of\_two\_and\_sixHundredSeventyThree\_dvd} & WheelK2\_673 & PROVED & \checkmark & verified \\
\texttt{Kp\_sixHundredSeventyThree} & WheelK2\_673 & PROVED & \checkmark & verified \\
\texttt{Kp\_sixHundredSeventyThree\_of\_dvd} & WheelK2\_673 & PROVED & \checkmark & verified \\
\texttt{Kp\_sixHundredSeventyThree\_of\_not\_dvd} & WheelK2\_673 & PROVED & \checkmark & verified \\
\texttt{K2\_677\_eq} & WheelK2\_677 & PROVED & \checkmark & verified \\
\texttt{K2\_677\_of\_not\_two\_dvd} & WheelK2\_677 & PROVED & \checkmark & verified \\
\texttt{K2\_677\_of\_two\_and\_sixHundredSeventySeven\_dvd} & WheelK2\_677 & PROVED & \checkmark & verified \\
\texttt{Kp\_sixHundredSeventySeven} & WheelK2\_677 & PROVED & \checkmark & verified \\
\texttt{Kp\_sixHundredSeventySeven\_of\_dvd} & WheelK2\_677 & PROVED & \checkmark & verified \\
\texttt{Kp\_sixHundredSeventySeven\_of\_not\_dvd} & WheelK2\_677 & PROVED & \checkmark & verified \\
\texttt{K2\_683\_eq} & WheelK2\_683 & PROVED & \checkmark & verified \\
\texttt{K2\_683\_of\_not\_two\_dvd} & WheelK2\_683 & PROVED & \checkmark & verified \\
\texttt{K2\_683\_of\_two\_and\_sixHundredEightyThree\_dvd} & WheelK2\_683 & PROVED & \checkmark & verified \\
\texttt{Kp\_sixHundredEightyThree} & WheelK2\_683 & PROVED & \checkmark & verified \\
\texttt{Kp\_sixHundredEightyThree\_of\_dvd} & WheelK2\_683 & PROVED & \checkmark & verified \\
\texttt{Kp\_sixHundredEightyThree\_of\_not\_dvd} & WheelK2\_683 & PROVED & \checkmark & verified \\
\texttt{K2\_691\_eq} & WheelK2\_691 & PROVED & \checkmark & verified \\
\texttt{K2\_691\_of\_not\_two\_dvd} & WheelK2\_691 & PROVED & \checkmark & verified \\
\texttt{K2\_691\_of\_two\_and\_sixHundredNinetyOne\_dvd} & WheelK2\_691 & PROVED & \checkmark & verified \\
\texttt{Kp\_sixHundredNinetyOne} & WheelK2\_691 & PROVED & \checkmark & verified \\
\texttt{Kp\_sixHundredNinetyOne\_of\_dvd} & WheelK2\_691 & PROVED & \checkmark & verified \\
\texttt{Kp\_sixHundredNinetyOne\_of\_not\_dvd} & WheelK2\_691 & PROVED & \checkmark & verified \\
\texttt{K2\_701\_eq} & WheelK2\_701 & PROVED & \checkmark & verified \\
\texttt{K2\_701\_of\_not\_two\_dvd} & WheelK2\_701 & PROVED & \checkmark & verified \\
\texttt{K2\_701\_of\_two\_and\_sevenHundredOne\_dvd} & WheelK2\_701 & PROVED & \checkmark & verified \\
\texttt{Kp\_sevenHundredOne} & WheelK2\_701 & PROVED & \checkmark & verified \\
\texttt{Kp\_sevenHundredOne\_of\_dvd} & WheelK2\_701 & PROVED & \checkmark & verified \\
\texttt{Kp\_sevenHundredOne\_of\_not\_dvd} & WheelK2\_701 & PROVED & \checkmark & verified \\
\texttt{K2\_709\_eq} & WheelK2\_709 & PROVED & \checkmark & verified \\
\texttt{K2\_709\_of\_not\_two\_dvd} & WheelK2\_709 & PROVED & \checkmark & verified \\
\texttt{K2\_709\_of\_two\_and\_sevenHundredNine\_dvd} & WheelK2\_709 & PROVED & \checkmark & verified \\
\texttt{Kp\_sevenHundredNine} & WheelK2\_709 & PROVED & \checkmark & verified \\
\texttt{Kp\_sevenHundredNine\_of\_dvd} & WheelK2\_709 & PROVED & \checkmark & verified \\
\texttt{Kp\_sevenHundredNine\_of\_not\_dvd} & WheelK2\_709 & PROVED & \checkmark & verified \\
\texttt{K2\_71\_eq} & WheelK2\_71 & PROVED & \checkmark & verified \\
\texttt{K2\_71\_of\_not\_two\_dvd} & WheelK2\_71 & PROVED & \checkmark & verified \\
\texttt{K2\_71\_of\_two\_and\_seventyOne\_dvd} & WheelK2\_71 & PROVED & \checkmark & verified \\
\texttt{Kp\_seventyOne} & WheelK2\_71 & PROVED & \checkmark & verified \\
\texttt{Kp\_seventyOne\_of\_dvd} & WheelK2\_71 & PROVED & \checkmark & verified \\
\texttt{Kp\_seventyOne\_of\_not\_dvd} & WheelK2\_71 & PROVED & \checkmark & verified \\
\texttt{K2\_719\_eq} & WheelK2\_719 & PROVED & \checkmark & verified \\
\texttt{K2\_719\_of\_not\_two\_dvd} & WheelK2\_719 & PROVED & \checkmark & verified \\
\texttt{K2\_719\_of\_two\_and\_sevenHundredNineteen\_dvd} & WheelK2\_719 & PROVED & \checkmark & verified \\
\texttt{Kp\_sevenHundredNineteen} & WheelK2\_719 & PROVED & \checkmark & verified \\
\texttt{Kp\_sevenHundredNineteen\_of\_dvd} & WheelK2\_719 & PROVED & \checkmark & verified \\
\texttt{Kp\_sevenHundredNineteen\_of\_not\_dvd} & WheelK2\_719 & PROVED & \checkmark & verified \\
\texttt{K2\_727\_eq} & WheelK2\_727 & PROVED & \checkmark & verified \\
\texttt{K2\_727\_of\_not\_two\_dvd} & WheelK2\_727 & PROVED & \checkmark & verified \\
\texttt{K2\_727\_of\_two\_and\_sevenHundredTwentySeven\_dvd} & WheelK2\_727 & PROVED & \checkmark & verified \\
\texttt{Kp\_sevenHundredTwentySeven} & WheelK2\_727 & PROVED & \checkmark & verified \\
\texttt{Kp\_sevenHundredTwentySeven\_of\_dvd} & WheelK2\_727 & PROVED & \checkmark & verified \\
\texttt{Kp\_sevenHundredTwentySeven\_of\_not\_dvd} & WheelK2\_727 & PROVED & \checkmark & verified \\
\texttt{K2\_73\_eq} & WheelK2\_73 & PROVED & \checkmark & verified \\
\texttt{K2\_73\_of\_not\_two\_dvd} & WheelK2\_73 & PROVED & \checkmark & verified \\
\texttt{K2\_73\_of\_two\_and\_seventyThree\_dvd} & WheelK2\_73 & PROVED & \checkmark & verified \\
\texttt{Kp\_seventyThree} & WheelK2\_73 & PROVED & \checkmark & verified \\
\texttt{Kp\_seventyThree\_of\_dvd} & WheelK2\_73 & PROVED & \checkmark & verified \\
\texttt{Kp\_seventyThree\_of\_not\_dvd} & WheelK2\_73 & PROVED & \checkmark & verified \\
\texttt{K2\_733\_eq} & WheelK2\_733 & PROVED & \checkmark & verified \\
\texttt{K2\_733\_of\_not\_two\_dvd} & WheelK2\_733 & PROVED & \checkmark & verified \\
\texttt{K2\_733\_of\_two\_and\_sevenHundredThirtyThree\_dvd} & WheelK2\_733 & PROVED & \checkmark & verified \\
\texttt{Kp\_sevenHundredThirtyThree} & WheelK2\_733 & PROVED & \checkmark & verified \\
\texttt{Kp\_sevenHundredThirtyThree\_of\_dvd} & WheelK2\_733 & PROVED & \checkmark & verified \\
\texttt{Kp\_sevenHundredThirtyThree\_of\_not\_dvd} & WheelK2\_733 & PROVED & \checkmark & verified \\
\texttt{K2\_739\_eq} & WheelK2\_739 & PROVED & \checkmark & verified \\
\texttt{K2\_739\_of\_not\_two\_dvd} & WheelK2\_739 & PROVED & \checkmark & verified \\
\texttt{K2\_739\_of\_two\_and\_sevenHundredThirtyNine\_dvd} & WheelK2\_739 & PROVED & \checkmark & verified \\
\texttt{Kp\_sevenHundredThirtyNine} & WheelK2\_739 & PROVED & \checkmark & verified \\
\texttt{Kp\_sevenHundredThirtyNine\_of\_dvd} & WheelK2\_739 & PROVED & \checkmark & verified \\
\texttt{Kp\_sevenHundredThirtyNine\_of\_not\_dvd} & WheelK2\_739 & PROVED & \checkmark & verified \\
\texttt{K2\_743\_eq} & WheelK2\_743 & PROVED & \checkmark & verified \\
\texttt{K2\_743\_of\_not\_two\_dvd} & WheelK2\_743 & PROVED & \checkmark & verified \\
\texttt{K2\_743\_of\_two\_and\_sevenHundredFortyThree\_dvd} & WheelK2\_743 & PROVED & \checkmark & verified \\
\texttt{Kp\_sevenHundredFortyThree} & WheelK2\_743 & PROVED & \checkmark & verified \\
\texttt{Kp\_sevenHundredFortyThree\_of\_dvd} & WheelK2\_743 & PROVED & \checkmark & verified \\
\texttt{Kp\_sevenHundredFortyThree\_of\_not\_dvd} & WheelK2\_743 & PROVED & \checkmark & verified \\
\texttt{K2\_751\_eq} & WheelK2\_751 & PROVED & \checkmark & verified \\
\texttt{K2\_751\_of\_not\_two\_dvd} & WheelK2\_751 & PROVED & \checkmark & verified \\
\texttt{K2\_751\_of\_two\_and\_sevenHundredFiftyOne\_dvd} & WheelK2\_751 & PROVED & \checkmark & verified \\
\texttt{Kp\_sevenHundredFiftyOne} & WheelK2\_751 & PROVED & \checkmark & verified \\
\texttt{Kp\_sevenHundredFiftyOne\_of\_dvd} & WheelK2\_751 & PROVED & \checkmark & verified \\
\texttt{Kp\_sevenHundredFiftyOne\_of\_not\_dvd} & WheelK2\_751 & PROVED & \checkmark & verified \\
\texttt{K2\_757\_eq} & WheelK2\_757 & PROVED & \checkmark & verified \\
\texttt{K2\_757\_of\_not\_two\_dvd} & WheelK2\_757 & PROVED & \checkmark & verified \\
\texttt{K2\_757\_of\_two\_and\_sevenHundredFiftySeven\_dvd} & WheelK2\_757 & PROVED & \checkmark & verified \\
\texttt{Kp\_sevenHundredFiftySeven} & WheelK2\_757 & PROVED & \checkmark & verified \\
\texttt{Kp\_sevenHundredFiftySeven\_of\_dvd} & WheelK2\_757 & PROVED & \checkmark & verified \\
\texttt{Kp\_sevenHundredFiftySeven\_of\_not\_dvd} & WheelK2\_757 & PROVED & \checkmark & verified \\
\texttt{K2\_761\_eq} & WheelK2\_761 & PROVED & \checkmark & verified \\
\texttt{K2\_761\_of\_not\_two\_dvd} & WheelK2\_761 & PROVED & \checkmark & verified \\
\texttt{K2\_761\_of\_two\_and\_sevenHundredSixtyOne\_dvd} & WheelK2\_761 & PROVED & \checkmark & verified \\
\texttt{Kp\_sevenHundredSixtyOne} & WheelK2\_761 & PROVED & \checkmark & verified \\
\texttt{Kp\_sevenHundredSixtyOne\_of\_dvd} & WheelK2\_761 & PROVED & \checkmark & verified \\
\texttt{Kp\_sevenHundredSixtyOne\_of\_not\_dvd} & WheelK2\_761 & PROVED & \checkmark & verified \\
\texttt{K2\_769\_eq} & WheelK2\_769 & PROVED & \checkmark & verified \\
\texttt{K2\_769\_of\_not\_two\_dvd} & WheelK2\_769 & PROVED & \checkmark & verified \\
\texttt{K2\_769\_of\_two\_and\_sevenHundredSixtyNine\_dvd} & WheelK2\_769 & PROVED & \checkmark & verified \\
\texttt{Kp\_sevenHundredSixtyNine} & WheelK2\_769 & PROVED & \checkmark & verified \\
\texttt{Kp\_sevenHundredSixtyNine\_of\_dvd} & WheelK2\_769 & PROVED & \checkmark & verified \\
\texttt{Kp\_sevenHundredSixtyNine\_of\_not\_dvd} & WheelK2\_769 & PROVED & \checkmark & verified \\
\texttt{K2\_773\_eq} & WheelK2\_773 & PROVED & \checkmark & verified \\
\texttt{K2\_773\_of\_not\_two\_dvd} & WheelK2\_773 & PROVED & \checkmark & verified \\
\texttt{K2\_773\_of\_two\_and\_sevenHundredSeventyThree\_dvd} & WheelK2\_773 & PROVED & \checkmark & verified \\
\texttt{Kp\_sevenHundredSeventyThree} & WheelK2\_773 & PROVED & \checkmark & verified \\
\texttt{Kp\_sevenHundredSeventyThree\_of\_dvd} & WheelK2\_773 & PROVED & \checkmark & verified \\
\texttt{Kp\_sevenHundredSeventyThree\_of\_not\_dvd} & WheelK2\_773 & PROVED & \checkmark & verified \\
\texttt{K2\_787\_eq} & WheelK2\_787 & PROVED & \checkmark & verified \\
\texttt{K2\_787\_of\_not\_two\_dvd} & WheelK2\_787 & PROVED & \checkmark & verified \\
\texttt{K2\_787\_of\_two\_and\_sevenHundredEightySeven\_dvd} & WheelK2\_787 & PROVED & \checkmark & verified \\
\texttt{Kp\_sevenHundredEightySeven} & WheelK2\_787 & PROVED & \checkmark & verified \\
\texttt{Kp\_sevenHundredEightySeven\_of\_dvd} & WheelK2\_787 & PROVED & \checkmark & verified \\
\texttt{Kp\_sevenHundredEightySeven\_of\_not\_dvd} & WheelK2\_787 & PROVED & \checkmark & verified \\
\texttt{K2\_79\_eq} & WheelK2\_79 & PROVED & \checkmark & verified \\
\texttt{K2\_79\_of\_not\_two\_dvd} & WheelK2\_79 & PROVED & \checkmark & verified \\
\texttt{K2\_79\_of\_two\_and\_seventyNine\_dvd} & WheelK2\_79 & PROVED & \checkmark & verified \\
\texttt{Kp\_seventyNine} & WheelK2\_79 & PROVED & \checkmark & verified \\
\texttt{Kp\_seventyNine\_of\_dvd} & WheelK2\_79 & PROVED & \checkmark & verified \\
\texttt{Kp\_seventyNine\_of\_not\_dvd} & WheelK2\_79 & PROVED & \checkmark & verified \\
\texttt{K2\_797\_eq} & WheelK2\_797 & PROVED & \checkmark & verified \\
\texttt{K2\_797\_of\_not\_two\_dvd} & WheelK2\_797 & PROVED & \checkmark & verified \\
\texttt{K2\_797\_of\_two\_and\_sevenHundredNinetySeven\_dvd} & WheelK2\_797 & PROVED & \checkmark & verified \\
\texttt{Kp\_sevenHundredNinetySeven} & WheelK2\_797 & PROVED & \checkmark & verified \\
\texttt{Kp\_sevenHundredNinetySeven\_of\_dvd} & WheelK2\_797 & PROVED & \checkmark & verified \\
\texttt{Kp\_sevenHundredNinetySeven\_of\_not\_dvd} & WheelK2\_797 & PROVED & \checkmark & verified \\
\texttt{K2\_83\_eq} & WheelK2\_83 & PROVED & \checkmark & verified \\
\texttt{K2\_83\_of\_not\_two\_dvd} & WheelK2\_83 & PROVED & \checkmark & verified \\
\texttt{K2\_83\_of\_two\_and\_eightyThree\_dvd} & WheelK2\_83 & PROVED & \checkmark & verified \\
\texttt{Kp\_eightyThree} & WheelK2\_83 & PROVED & \checkmark & verified \\
\texttt{Kp\_eightyThree\_of\_dvd} & WheelK2\_83 & PROVED & \checkmark & verified \\
\texttt{Kp\_eightyThree\_of\_not\_dvd} & WheelK2\_83 & PROVED & \checkmark & verified \\
\texttt{K2\_89\_eq} & WheelK2\_89 & PROVED & \checkmark & verified \\
\texttt{K2\_89\_of\_not\_two\_dvd} & WheelK2\_89 & PROVED & \checkmark & verified \\
\texttt{K2\_89\_of\_two\_and\_eightyNine\_dvd} & WheelK2\_89 & PROVED & \checkmark & verified \\
\texttt{Kp\_eightyNine} & WheelK2\_89 & PROVED & \checkmark & verified \\
\texttt{Kp\_eightyNine\_of\_dvd} & WheelK2\_89 & PROVED & \checkmark & verified \\
\texttt{Kp\_eightyNine\_of\_not\_dvd} & WheelK2\_89 & PROVED & \checkmark & verified \\
\texttt{K2\_97\_eq} & WheelK2\_97 & PROVED & \checkmark & verified \\
\texttt{K2\_97\_of\_not\_two\_dvd} & WheelK2\_97 & PROVED & \checkmark & verified \\
\texttt{K2\_97\_of\_two\_and\_ninetySeven\_dvd} & WheelK2\_97 & PROVED & \checkmark & verified \\
\texttt{Kp\_ninetySeven} & WheelK2\_97 & PROVED & \checkmark & verified \\
\texttt{Kp\_ninetySeven\_of\_dvd} & WheelK2\_97 & PROVED & \checkmark & verified \\
\texttt{Kp\_ninetySeven\_of\_not\_dvd} & WheelK2\_97 & PROVED & \checkmark & verified \\
\texttt{gCount\_centered} & GoldbachComb & PROVED & \checkmark & verified \\
\texttt{gCount\_eq} & GoldbachComb & PROVED & \checkmark & verified \\
\texttt{local\_covariance} & GoldbachComb & PROVED & \checkmark & verified \\
\texttt{factor\_ratio} & GoldbachLemmas & PROVED & \checkmark & verified \\
\texttt{gFactor\_eq\_one\_add} & GoldbachLemmas & PROVED & \checkmark & verified \\
\texttt{gFactor\_pos} & GoldbachLemmas & PROVED & \checkmark & verified \\
\texttt{gFactor\_pos\_prime} & GoldbachLemmas & PROVED & \checkmark & verified \\
\texttt{gProduct\_le\_of\_subset} & GoldbachLemmas & PROVED & \checkmark & verified \\
\texttt{gProduct\_pos} & GoldbachLemmas & PROVED & \checkmark & verified \\
\texttt{gResidues\_card\_ne\_zero} & GoldbachLemmas & PROVED & \checkmark & verified \\
\texttt{gResidues\_card\_zero} & GoldbachLemmas & PROVED & \checkmark & verified \\
\texttt{localDensity\_ne} & GoldbachLemmas & PROVED & \checkmark & verified \\
\texttt{localDensity\_zero} & GoldbachLemmas & PROVED & \checkmark & verified \\
\texttt{one\_le\_gProduct} & GoldbachLemmas & PROVED & \checkmark & verified \\
\texttt{one\_lt\_gFactor} & GoldbachLemmas & PROVED & \checkmark & verified \\
\texttt{one\_lt\_gFactor\_prime} & GoldbachLemmas & PROVED & \checkmark & verified \\
\texttt{tFactor\_eq} & GoldbachLemmas & PROVED & \checkmark & verified \\
\texttt{tFactor\_lt\_one} & GoldbachLemmas & PROVED & \checkmark & verified \\
\texttt{tFactor\_pos} & GoldbachLemmas & PROVED & \checkmark & verified \\
\texttt{three\_le\_of\_prime\_ne\_two} & GoldbachLemmas & PROVED & \checkmark & verified \\
\texttt{goldbachCount\_four} & GoldbachSchema & PROVED & \checkmark & verified \\
\texttt{goldbachCount\_ten} & GoldbachSchema & PROVED & \checkmark & verified \\
\texttt{goldbach\_beyond\_of\_model} & GoldbachSchema & CONDITIONAL & \checkmark & verified \\
\texttt{goldbach\_from\_spectral\_model} & GoldbachSchema & CONDITIONAL & \checkmark & verified \\
\texttt{hasGoldbachRep\_eight} & GoldbachSchema & PROVED & \checkmark & verified \\
\texttt{hasGoldbachRep\_four} & GoldbachSchema & PROVED & \checkmark & verified \\
\texttt{hasGoldbachRep\_of\_count\_pos} & GoldbachSchema & PROVED & \checkmark & verified \\
\texttt{hasGoldbachRep\_six} & GoldbachSchema & PROVED & \checkmark & verified \\
\texttt{admissibleUnits\_card} & GoldbachSelectionRule & PROVED & \checkmark & verified \\
\texttt{admissibleUnits\_card\_totient} & GoldbachSelectionRule & PROVED & \checkmark & verified \\
\texttt{admissibleUnits\_eq\_sdiff} & GoldbachSelectionRule & PROVED & \checkmark & verified \\
\texttt{admissibleUnits\_reflection\_eq\_residues} & GoldbachSelectionRule & PROVED & \checkmark & verified \\
\texttt{admissibleUnits\_translation\_eq\_residues} & GoldbachSelectionRule & PROVED & \checkmark & verified \\
\texttt{admissible\_card\_dihedral} & GoldbachSelectionRule & PROVED & \checkmark & verified \\
\texttt{admissible\_reflection\_card} & GoldbachSelectionRule & PROVED & \checkmark & verified \\
\texttt{admissible\_translation\_card} & GoldbachSelectionRule & PROVED & \checkmark & verified \\
\texttt{admissible\_translation\_matches\_dichotomy} & GoldbachSelectionRule & PROVED & \checkmark & verified \\
\texttt{gapPairs\_card} & GoldbachSelectionRule & PROVED & \checkmark & verified \\
\texttt{gap\_pairs\_eq\_translation\_admissible} & GoldbachSelectionRule & PROVED & \checkmark & verified \\
\texttt{goldbachPairs\_card} & GoldbachSelectionRule & PROVED & \checkmark & verified \\
\texttt{goldbach\_pairs\_eq\_reflection\_admissible} & GoldbachSelectionRule & PROVED & \checkmark & verified \\
\texttt{mem\_admissibleUnits\_iff} & GoldbachSelectionRule & PROVED & \checkmark & verified \\
\texttt{reflection\_eq\_dihedral} & GoldbachSelectionRule & PROVED & \checkmark & verified \\
\texttt{sigma\_card} & GoldbachSelectionRule & PROVED & \checkmark & verified \\
\texttt{translation\_eq\_dihedral} & GoldbachSelectionRule & PROVED & \checkmark & verified \\
\texttt{units\_card} & GoldbachSelectionRule & PROVED & \checkmark & verified \\
\texttt{golden\_in\_cycleSpectrum\_iff\_five\_dvd} & GoldenDivisibility & PROVED & \checkmark & verified \\
\texttt{golden\_unique\_to\_five\_recovered} & GoldenDivisibility & PROVED & \checkmark & verified \\
\texttt{golden\_quadratic\_roots} & GoldenSpectralCharacterization & PROVED & \checkmark & verified \\
\texttt{golden\_ratio\_spectral\_characterization} & GoldenSpectralCharacterization & PROVED & \checkmark & verified \\
\texttt{golden\_roots\_mem\_pentagon\_spectrum} & GoldenSpectralCharacterization & PROVED & \checkmark & verified \\
\texttt{golden\_roots\_sign\_split} & GoldenSpectralCharacterization & PROVED & \checkmark & verified \\
\texttt{C5\_membership\_layer} & GoldenUniqueness & PROVED & \checkmark & verified \\
\texttt{algebraic\_connectivity\_C5\_eq} & GoldenUniqueness & PROVED & \checkmark & verified \\
\texttt{algebraic\_connectivity\_C5\_props} & GoldenUniqueness & PROVED & \checkmark & verified \\
\texttt{algebraic\_layer} & GoldenUniqueness & PROVED & \checkmark & verified \\
\texttt{cos\_two\_pi\_div\_five\_eq\_phi\_sub\_one\_div\_two} & GoldenUniqueness & PROVED & \checkmark & verified \\
\texttt{five\_carries\_golden} & GoldenUniqueness & PROVED & \checkmark & verified \\
\texttt{golden\_algebraic\_identity} & GoldenUniqueness & PROVED & \checkmark & verified \\
\texttt{golden\_in\_C5\_family} & GoldenUniqueness & PROVED & \checkmark & verified \\
\texttt{golden\_mem\_C5} & GoldenUniqueness & PROVED & \checkmark & verified \\
\texttt{golden\_not\_in\_prime\_cycle\_ne\_five} & GoldenUniqueness & PROVED & \checkmark & verified \\
\texttt{golden\_ratio\_in\_C5\_geometry} & GoldenUniqueness & PROVED & \checkmark & verified \\
\texttt{golden\_sub\_one\_eq\_two\_cos} & GoldenUniqueness & PROVED & \checkmark & verified \\
\texttt{golden\_unique\_among\_prime\_cycles} & GoldenUniqueness & PROVED & \checkmark & verified \\
\texttt{golden\_unique\_to\_five} & GoldenUniqueness & PROVED & \checkmark & verified \\
\texttt{inv\_golden\_eq\_sub\_one} & GoldenUniqueness & PROVED & \checkmark & verified \\
\texttt{neg\_golden\_mem\_C5} & GoldenUniqueness & PROVED & \checkmark & verified \\
\texttt{prime\_cycle\_golden\_forces\_five} & GoldenUniqueness & PROVED & \checkmark & verified \\
\texttt{prime\_rigidity\_layer} & GoldenUniqueness & PROVED & \checkmark & verified \\
\texttt{two\_cos\_four\_pi\_div\_five\_eq\_neg\_golden} & GoldenUniqueness & PROVED & \checkmark & verified \\
\texttt{two\_cos\_fundamental\_mode\_C5} & GoldenUniqueness & PROVED & \checkmark & verified \\
\texttt{hyperperfect\_1\_28} & HyperperfectNumbers & PROVED & \checkmark & verified \\
\texttt{hyperperfect\_1\_6} & HyperperfectNumbers & PROVED & \checkmark & verified \\
\texttt{hyperperfect\_2\_21} & HyperperfectNumbers & PROVED & \checkmark & verified \\
\texttt{hyperperfect\_2\_2133} & HyperperfectNumbers & PROVED & \checkmark & verified \\
\texttt{hyperperfect\_6\_301} & HyperperfectNumbers & PROVED & \checkmark & verified \\
\texttt{hyperperfect\_one\_iff\_sigma\_two\_mul} & HyperperfectNumbers & PROVED & \checkmark & verified \\
\texttt{not\_hyperperfect\_2\_6} & HyperperfectNumbers & PROVED & \checkmark & verified \\
\texttt{nsq\_1} & LandauNSquaredPlusOne & PROVED & \checkmark & verified \\
\texttt{nsq\_10} & LandauNSquaredPlusOne & PROVED & \checkmark & verified \\
\texttt{nsq\_14} & LandauNSquaredPlusOne & PROVED & \checkmark & verified \\
\texttt{nsq\_16} & LandauNSquaredPlusOne & PROVED & \checkmark & verified \\
\texttt{nsq\_2} & LandauNSquaredPlusOne & PROVED & \checkmark & verified \\
\texttt{nsq\_20} & LandauNSquaredPlusOne & PROVED & \checkmark & verified \\
\texttt{nsq\_24} & LandauNSquaredPlusOne & PROVED & \checkmark & verified \\
\texttt{nsq\_26} & LandauNSquaredPlusOne & PROVED & \checkmark & verified \\
\texttt{nsq\_4} & LandauNSquaredPlusOne & PROVED & \checkmark & verified \\
\texttt{nsq\_6} & LandauNSquaredPlusOne & PROVED & \checkmark & verified \\
\texttt{nsq\_even\_of\_prime} & LandauNSquaredPlusOne & PROVED & \checkmark & verified \\
\texttt{nsq\_mod\_two\_of\_prime} & LandauNSquaredPlusOne & PROVED & \checkmark & verified \\
\texttt{bertrand\_holds} & LegendreConjecture & PROVED & \checkmark & verified \\
\texttt{legendre\_1} & LegendreConjecture & PROVED & \checkmark & verified \\
\texttt{legendre\_10} & LegendreConjecture & PROVED & \checkmark & verified \\
\texttt{legendre\_11} & LegendreConjecture & PROVED & \checkmark & verified \\
\texttt{legendre\_12} & LegendreConjecture & PROVED & \checkmark & verified \\
\texttt{legendre\_2} & LegendreConjecture & PROVED & \checkmark & verified \\
\texttt{legendre\_3} & LegendreConjecture & PROVED & \checkmark & verified \\
\texttt{legendre\_4} & LegendreConjecture & PROVED & \checkmark & verified \\
\texttt{legendre\_5} & LegendreConjecture & PROVED & \checkmark & verified \\
\texttt{legendre\_6} & LegendreConjecture & PROVED & \checkmark & verified \\
\texttt{legendre\_7} & LegendreConjecture & PROVED & \checkmark & verified \\
\texttt{legendre\_8} & LegendreConjecture & PROVED & \checkmark & verified \\
\texttt{legendre\_9} & LegendreConjecture & PROVED & \checkmark & verified \\
\texttt{lehmer\_odd} & LehmerTotient & PROVED & \checkmark & verified \\
\texttt{lehmer\_squarefree} & LehmerTotient & PROVED & \checkmark & verified \\
\texttt{lehmer\_three\_primes} & LehmerTotient & PROVED & \checkmark & verified \\
\texttt{countermodel\_assoc} & MagmaLawRefutations & PROVED & \checkmark & verified \\
\texttt{countermodel\_comm} & MagmaLawRefutations & PROVED & \checkmark & verified \\
\texttt{countermodel\_idem} & MagmaLawRefutations & PROVED & \checkmark & verified \\
\texttt{countermodel\_mid} & MagmaLawRefutations & PROVED & \checkmark & verified \\
\texttt{not\_entails\_assoc} & MagmaLawRefutations & PROVED & \checkmark & verified \\
\texttt{not\_entails\_comm} & MagmaLawRefutations & PROVED & \checkmark & verified \\
\texttt{not\_entails\_idem} & MagmaLawRefutations & PROVED & \checkmark & verified \\
\texttt{not\_entails\_mid} & MagmaLawRefutations & PROVED & \checkmark & verified \\
\texttt{eq\_two\_pow\_mul\_odd} & MersennePerfect & PROVED & \checkmark & verified \\
\texttt{eq\_two\_pow\_mul\_prime\_mersenne\_of\_even\_perfect} & MersennePerfect & PROVED & \checkmark & verified \\
\texttt{even\_two\_pow\_mul\_mersenne\_of\_prime} & MersennePerfect & PROVED & \checkmark & verified \\
\texttt{infinitude\_equiv} & MersennePerfect & PROVED & \checkmark & verified \\
\texttt{mersenne\_2\_prime} & MersennePerfect & PROVED & \checkmark & verified \\
\texttt{mersenne\_3\_prime} & MersennePerfect & PROVED & \checkmark & verified \\
\texttt{mersenne\_5\_prime} & MersennePerfect & PROVED & \checkmark & verified \\
\texttt{mersenne\_7\_prime} & MersennePerfect & PROVED & \checkmark & verified \\
\texttt{ne\_zero\_of\_prime\_mersenne} & MersennePerfect & PROVED & \checkmark & verified \\
\texttt{perfect\_28} & MersennePerfect & PROVED & \checkmark & verified \\
\texttt{perfect\_496} & MersennePerfect & PROVED & \checkmark & verified \\
\texttt{perfect\_6} & MersennePerfect & PROVED & \checkmark & verified \\
\texttt{perfect\_8128} & MersennePerfect & PROVED & \checkmark & verified \\
\texttt{perfect\_two\_pow\_mul\_mersenne\_of\_prime} & MersennePerfect & PROVED & \checkmark & verified \\
\texttt{sigma\_two\_pow\_eq\_mersenne\_succ} & MersennePerfect & PROVED & \checkmark & verified \\
\texttt{H3\_ground\_metallic} & MetallicFamily & PROVED & \checkmark & verified \\
\texttt{inv\_metallicMean\_eq\_sub} & MetallicFamily & PROVED & \checkmark & verified \\
\texttt{metallicConj\_sq} & MetallicFamily & PROVED & \checkmark & verified \\
\texttt{metallicMean\_add\_conj} & MetallicFamily & PROVED & \checkmark & verified \\
\texttt{metallicMean\_mul\_conj} & MetallicFamily & PROVED & \checkmark & verified \\
\texttt{metallicMean\_one} & MetallicFamily & PROVED & \checkmark & verified \\
\texttt{metallicMean\_pos\_of\_nonneg} & MetallicFamily & PROVED & \checkmark & verified \\
\texttt{metallicMean\_sq} & MetallicFamily & PROVED & \checkmark & verified \\
\texttt{metallicMean\_sub\_conj} & MetallicFamily & PROVED & \checkmark & verified \\
\texttt{metallicMean\_two} & MetallicFamily & PROVED & \checkmark & verified \\
\texttt{metallic\_one\_unique\_to\_five} & MetallicFamily & PROVED & \checkmark & verified \\
\texttt{metallic\_radicand\_nonneg} & MetallicFamily & PROVED & \checkmark & verified \\
\texttt{silverGap\_eq\_three\_sub\_metallicMean\_two} & MetallicFamily & PROVED & \checkmark & verified \\
\texttt{det\_eq\_prod\_roots} & MetallicRealization & PROVED & \checkmark & verified \\
\texttt{eigvec\_conj\_ne\_zero} & MetallicRealization & PROVED & \checkmark & verified \\
\texttt{eigvec\_mean\_ne\_zero} & MetallicRealization & PROVED & \checkmark & verified \\
\texttt{fibQ\_charpoly} & MetallicRealization & PROVED & \checkmark & verified \\
\texttt{fibQ\_mulVec\_golden} & MetallicRealization & PROVED & \checkmark & verified \\
\texttt{goldenRatio\_irrational} & MetallicRealization & PROVED & \checkmark & verified \\
\texttt{golden\_hasEigenvalue} & MetallicRealization & PROVED & \checkmark & verified \\
\texttt{metallicConj\_isRoot} & MetallicRealization & PROVED & \checkmark & verified \\
\texttt{metallicMatrix\_charpoly} & MetallicRealization & PROVED & \checkmark & verified \\
\texttt{metallicMatrix\_charpoly\_natDegree} & MetallicRealization & PROVED & \checkmark & verified \\
\texttt{metallicMatrix\_det} & MetallicRealization & PROVED & \checkmark & verified \\
\texttt{metallicMatrix\_eigenvalue\_iff} & MetallicRealization & PROVED & \checkmark & verified \\
\texttt{metallicMatrix\_hasEigenvalue\_conj} & MetallicRealization & PROVED & \checkmark & verified \\
\texttt{metallicMatrix\_hasEigenvalue\_mean} & MetallicRealization & PROVED & \checkmark & verified \\
\texttt{metallicMatrix\_isSymm} & MetallicRealization & PROVED & \checkmark & verified \\
\texttt{metallicMatrix\_mulVec\_conj} & MetallicRealization & PROVED & \checkmark & verified \\
\texttt{metallicMatrix\_mulVec\_mean} & MetallicRealization & PROVED & \checkmark & verified \\
\texttt{metallicMatrix\_one} & MetallicRealization & PROVED & \checkmark & verified \\
\texttt{metallicMatrix\_trace} & MetallicRealization & PROVED & \checkmark & verified \\
\texttt{metallicMean\_irrational} & MetallicRealization & PROVED & \checkmark & verified \\
\texttt{metallicMean\_isRoot} & MetallicRealization & PROVED & \checkmark & verified \\
\texttt{metallicMean\_notMem\_cycleSpectrum} & MetallicRealization & PROVED & \checkmark & verified \\
\texttt{metallicMean\_two\_notMem\_cycleSpectrum} & MetallicRealization & PROVED & \checkmark & verified \\
\texttt{metallicPoly\_factor} & MetallicRealization & PROVED & \checkmark & verified \\
\texttt{metallicPoly\_natDegree} & MetallicRealization & PROVED & \checkmark & verified \\
\texttt{trace\_eq\_sum\_roots} & MetallicRealization & PROVED & \checkmark & verified \\
\texttt{two\_lt\_metallicMean} & MetallicRealization & PROVED & \checkmark & verified \\
\texttt{brockian\_admissible\_count} & NewEra & PROVED & \checkmark & verified \\
\texttt{decaying\_potential\_cannot\_realize\_large\_zeros} & NewEra & PROVED & \checkmark & verified \\
\texttt{golden\_lives\_on\_the\_pentagon} & NewEra & PROVED & \checkmark & verified \\
\texttt{neg\_golden\_lives\_on\_the\_pentagon} & NewEra & PROVED & \checkmark & verified \\
\texttt{pentagon\_cosine\_is\_golden} & NewEra & PROVED & \checkmark & verified \\
\texttt{primeGaussian\_not\_confining} & NewEra & PROVED & \checkmark & verified \\
\texttt{quadratic\_is\_confining} & NewEra & PROVED & \checkmark & verified \\
\texttt{twin\_admissible\_count} & NewEra & PROVED & \checkmark & verified \\
\texttt{why\_five} & NewEra & PROVED & \checkmark & verified \\
\texttt{card\_distincts\_le\_partition} & OddDistinctPartition & PROVED & \checkmark & verified \\
\texttt{card\_distincts\_le\_powerset} & OddDistinctPartition & PROVED & \checkmark & verified \\
\texttt{card\_oddDistincts\_le\_distincts} & OddDistinctPartition & PROVED & \checkmark & verified \\
\texttt{card\_oddDistincts\_le\_odds} & OddDistinctPartition & PROVED & \checkmark & verified \\
\texttt{card\_odds\_le\_partition} & OddDistinctPartition & PROVED & \checkmark & verified \\
\texttt{countRestricted\_two\_eq\_distincts} & OddDistinctPartition & PROVED & \checkmark & verified \\
\texttt{distincts\_one\_card} & OddDistinctPartition & PROVED & \checkmark & verified \\
\texttt{distincts\_zero\_card} & OddDistinctPartition & PROVED & \checkmark & verified \\
\texttt{euler\_odd\_eq\_distinct} & OddDistinctPartition & PROVED & \checkmark & verified \\
\texttt{euler\_one} & OddDistinctPartition & PROVED & \checkmark & verified \\
\texttt{euler\_via\_glaisher} & OddDistinctPartition & PROVED & \checkmark & verified \\
\texttt{euler\_zero} & OddDistinctPartition & PROVED & \checkmark & verified \\
\texttt{glaisher} & OddDistinctPartition & PROVED & \checkmark & verified \\
\texttt{mem\_distincts\_iff} & OddDistinctPartition & PROVED & \checkmark & verified \\
\texttt{mem\_odds\_iff} & OddDistinctPartition & PROVED & \checkmark & verified \\
\texttt{oddDistincts\_subset\_distincts} & OddDistinctPartition & PROVED & \checkmark & verified \\
\texttt{oddDistincts\_subset\_odds} & OddDistinctPartition & PROVED & \checkmark & verified \\
\texttt{odds\_one\_card} & OddDistinctPartition & PROVED & \checkmark & verified \\
\texttt{odds\_zero\_card} & OddDistinctPartition & PROVED & \checkmark & verified \\
\texttt{parts\_nodup\_of\_mem\_distincts} & OddDistinctPartition & PROVED & \checkmark & verified \\
\texttt{parts\_odd\_of\_mem\_odds} & OddDistinctPartition & PROVED & \checkmark & verified \\
\texttt{parts\_subset\_Icc} & OddDistinctPartition & PROVED & \checkmark & verified \\
\texttt{powerSeries\_odds\_eq\_distincts} & OddDistinctPartition & PROVED & \checkmark & verified \\
\texttt{toFinset\_inj\_on\_distincts} & OddDistinctPartition & PROVED & \checkmark & verified \\
\texttt{even\_factorization\_of\_isSquare} & OddPerfectConstraints & PROVED & \checkmark & verified \\
\texttt{oddPerfect\_not\_prime\_pow} & OddPerfectConstraints & PROVED & \checkmark & verified \\
\texttt{oddPerfect\_not\_square} & OddPerfectConstraints & PROVED & \checkmark & verified \\
\texttt{oddPerfect\_not\_two\_dvd} & OddPerfectConstraints & PROVED & \checkmark & verified \\
\texttt{oddPerfect\_one\_lt} & OddPerfectConstraints & PROVED & \checkmark & verified \\
\texttt{oddPerfect\_pos} & OddPerfectConstraints & PROVED & \checkmark & verified \\
\texttt{odd\_of\_prime\_dvd\_odd} & OddPerfectConstraints & PROVED & \checkmark & verified \\
\texttt{oppermannUpper\_imp\_legendre} & OppermannConjecture & PROVED & \checkmark & verified \\
\texttt{oppermann\_10} & OppermannConjecture & PROVED & \checkmark & verified \\
\texttt{oppermann\_2} & OppermannConjecture & PROVED & \checkmark & verified \\
\texttt{oppermann\_3} & OppermannConjecture & PROVED & \checkmark & verified \\
\texttt{oppermann\_4} & OppermannConjecture & PROVED & \checkmark & verified \\
\texttt{oppermann\_5} & OppermannConjecture & PROVED & \checkmark & verified \\
\texttt{oppermann\_6} & OppermannConjecture & PROVED & \checkmark & verified \\
\texttt{oppermann\_7} & OppermannConjecture & PROVED & \checkmark & verified \\
\texttt{oppermann\_8} & OppermannConjecture & PROVED & \checkmark & verified \\
\texttt{oppermann\_9} & OppermannConjecture & PROVED & \checkmark & verified \\
\texttt{harmonic\_140} & OreHarmonicNumbers & PROVED & \checkmark & verified \\
\texttt{harmonic\_270} & OreHarmonicNumbers & PROVED & \checkmark & verified \\
\texttt{harmonic\_28} & OreHarmonicNumbers & PROVED & \checkmark & verified \\
\texttt{harmonic\_496} & OreHarmonicNumbers & PROVED & \checkmark & verified \\
\texttt{harmonic\_6} & OreHarmonicNumbers & PROVED & \checkmark & verified \\
\texttt{not\_harmonic\_4} & OreHarmonicNumbers & PROVED & \checkmark & verified \\
\texttt{eleven\_dvd\_one\_add\_odd\_pow} & PalindromicPrimes & PROVED & \checkmark & verified \\
\texttt{eleven\_unique\_even\_palindromic\_prime} & PalindromicPrimes & PROVED & \checkmark & verified \\
\texttt{even\_digit\_palindrome\_dvd\_11} & PalindromicPrimes & PROVED & \checkmark & verified \\
\texttt{even\_len\_palindrome\_dvd\_11\_examples} & PalindromicPrimes & PROVED & \checkmark & verified \\
\texttt{even\_palindrome\_dvd\_11} & PalindromicPrimes & PROVED & \checkmark & verified \\
\texttt{not\_palindromic\_13} & PalindromicPrimes & PROVED & \checkmark & verified \\
\texttt{palindromic\_101} & PalindromicPrimes & PROVED & \checkmark & verified \\
\texttt{palindromic\_11} & PalindromicPrimes & PROVED & \checkmark & verified \\
\texttt{palindromic\_131} & PalindromicPrimes & PROVED & \checkmark & verified \\
\texttt{palindromic\_151} & PalindromicPrimes & PROVED & \checkmark & verified \\
\texttt{palindromic\_2} & PalindromicPrimes & PROVED & \checkmark & verified \\
\texttt{palindromic\_313} & PalindromicPrimes & PROVED & \checkmark & verified \\
\texttt{palindromic\_353} & PalindromicPrimes & PROVED & \checkmark & verified \\
\texttt{palindromic\_727} & PalindromicPrimes & PROVED & \checkmark & verified \\
\texttt{palindromic\_757} & PalindromicPrimes & PROVED & \checkmark & verified \\
\texttt{palindromic\_919} & PalindromicPrimes & PROVED & \checkmark & verified \\
\texttt{two\_digit\_palindrome\_dvd\_11} & PalindromicPrimes & PROVED & \checkmark & verified \\
\texttt{factor\_eq\_geo} & PartitionRecurrence & PROVED & \checkmark & verified \\
\texttt{geo\_mul} & PartitionRecurrence & PROVED & \checkmark & verified \\
\texttt{partitionGF\_eq\_tprod\_geo} & PartitionRecurrence & PROVED & \checkmark & verified \\
\texttt{partitionGF\_mul\_pentagonalProduct} & PartitionRecurrence & PROVED & \checkmark & verified \\
\texttt{partition\_pentagonal\_convolution} & PartitionRecurrence & PROVED & \checkmark & verified \\
\texttt{partition\_pentagonal\_convolution\_range} & PartitionRecurrence & PROVED & \checkmark & verified \\
\texttt{partition\_recurrence} & PartitionRecurrence & PROVED & \checkmark & verified \\
\texttt{pentCoeff\_zero} & PartitionRecurrence & PROVED & \checkmark & verified \\
\texttt{A\_ae\_add} & Penrose & PROVED & \checkmark & verified \\
\texttt{A\_ae\_coeFn} & Penrose & PROVED & \checkmark & verified \\
\texttt{A\_ae\_memLp} & Penrose & PROVED & \checkmark & verified \\
\texttt{A\_ae\_smul} & Penrose & PROVED & \checkmark & verified \\
\texttt{A\_raw\_bound} & Penrose & PROVED & \checkmark & verified \\
\texttt{A\_raw\_ptwise\_bound} & Penrose & PROVED & \checkmark & verified \\
\texttt{D\_ae\_add} & Penrose & PROVED & \checkmark & verified \\
\texttt{D\_ae\_coeFn} & Penrose & PROVED & \checkmark & verified \\
\texttt{D\_ae\_memLp} & Penrose & PROVED & \checkmark & verified \\
\texttt{D\_ae\_smul} & Penrose & PROVED & \checkmark & verified \\
\texttt{D\_raw\_bound\_ptwise} & Penrose & PROVED & \checkmark & verified \\
\texttt{D\_raw\_norm\_bound} & Penrose & PROVED & \checkmark & verified \\
\texttt{Delta\_bounded} & Penrose & PROVED & \checkmark & verified \\
\texttt{adjacent\_loopless} & Penrose & PROVED & \checkmark & verified \\
\texttt{adjacent\_symm} & Penrose & PROVED & \checkmark & verified \\
\texttt{ae\_eq\_of\_count} & Penrose & PROVED & \checkmark & verified \\
\texttt{deg\_fn\_bound} & Penrose & PROVED & \checkmark & verified \\
\texttt{degree\_bound} & Penrose & PROVED & \checkmark & verified \\
\texttt{lintegral\_count\_eq\_tsum} & Penrose & PROVED & \checkmark & verified \\
\texttt{lintegral\_sum\_neighbors\_le} & Penrose & PROVED & \checkmark & verified \\
\texttt{memLp\_A\_raw} & Penrose & PROVED & \checkmark & verified \\
\texttt{memLp\_D\_raw} & Penrose & PROVED & \checkmark & verified \\
\texttt{neighbor\_subset} & Penrose & PROVED & \checkmark & verified \\
\texttt{norm\_A\_le} & Penrose & PROVED & \checkmark & verified \\
\texttt{norm\_D\_le} & Penrose & PROVED & \checkmark & verified \\
\texttt{pentagonVertex\_norm} & Penrose & PROVED & \checkmark & verified \\
\texttt{phi\_equation} & Penrose & PROVED & \checkmark & verified \\
\texttt{phi\_gt\_one} & Penrose & PROVED & \checkmark & verified \\
\texttt{phi\_ne\_zero} & Penrose & PROVED & \checkmark & verified \\
\texttt{phi\_pos} & Penrose & PROVED & \checkmark & verified \\
\texttt{phi\_product\_conjugate} & Penrose & PROVED & \checkmark & verified \\
\texttt{phi\_reciprocal} & Penrose & PROVED & \checkmark & verified \\
\texttt{phi\_squared} & Penrose & PROVED & \checkmark & verified \\
\texttt{phi\_sum\_conjugate} & Penrose & PROVED & \checkmark & verified \\
\texttt{potentialNeighbors\_finite} & Penrose & PROVED & \checkmark & verified \\
\texttt{rotation\_is\_multiplication} & Penrose & PROVED & \checkmark & verified \\
\texttt{sqrt5\_gt\_one} & Penrose & PROVED & \checkmark & verified \\
\texttt{sqrt5\_gt\_two} & Penrose & PROVED & \checkmark & verified \\
\texttt{sum\_sq\_le\_card\_mul\_sum\_sq} & Penrose & PROVED & \checkmark & verified \\
\texttt{zeta5\_norm} & Penrose & PROVED & \checkmark & verified \\
\texttt{zeta5\_pow\_five} & Penrose & PROVED & \checkmark & verified \\
\texttt{chiConjugate\_one} & PentagonCharacterMultiplicity & PROVED & \checkmark & verified \\
\texttt{chiGolden\_one} & PentagonCharacterMultiplicity & PROVED & \checkmark & verified \\
\texttt{golden\_isotypic\_multiplicity} & PentagonCharacterMultiplicity & PROVED & \checkmark & verified \\
\texttt{golden\_multiplicity\_eq\_irrep\_dim} & PentagonCharacterMultiplicity & PROVED & \checkmark & verified \\
\texttt{golden\_self\_inner} & PentagonCharacterMultiplicity & PROVED & \checkmark & verified \\
\texttt{multiplicity\_table} & PentagonCharacterMultiplicity & PROVED & \checkmark & verified \\
\texttt{neg\_golden\_multiplicity\_eq\_irrep\_dim} & PentagonCharacterMultiplicity & PROVED & \checkmark & verified \\
\texttt{permInner\_golden} & PentagonCharacterMultiplicity & PROVED & \checkmark & verified \\
\texttt{adjacency\_comm\_d5} & PentagonEquivariance & PROVED & \checkmark & verified \\
\texttt{adjacency\_comm\_rot} & PentagonEquivariance & PROVED & \checkmark & verified \\
\texttt{d5Pull\_sr\_apply} & PentagonEquivariance & PROVED & \checkmark & verified \\
\texttt{golden\_eigenspace\_invariant\_d5} & PentagonEquivariance & PROVED & \checkmark & verified \\
\texttt{golden\_eigenspace\_invariant\_rot} & PentagonEquivariance & PROVED & \checkmark & verified \\
\texttt{golden\_eigenspace\_is\_subrep} & PentagonEquivariance & PROVED & \checkmark & verified \\
\texttt{golden\_eigenspace\_is\_subrep\_d5} & PentagonEquivariance & PROVED & \checkmark & verified \\
\texttt{pentagon\_grand\_equivalence} & PentagonGrandEquivalence & PROVED & \checkmark & verified \\
\texttt{adjEigenvalue\_eq\_two\_cos} & PentagonIsotypic & PROVED & \checkmark & verified \\
\texttt{adjEigenvalue\_four} & PentagonIsotypic & PROVED & \checkmark & verified \\
\texttt{adjEigenvalue\_neg} & PentagonIsotypic & PROVED & \checkmark & verified \\
\texttt{adjEigenvalue\_one} & PentagonIsotypic & PROVED & \checkmark & verified \\
\texttt{adjEigenvalue\_three} & PentagonIsotypic & PROVED & \checkmark & verified \\
\texttt{adjEigenvalue\_two} & PentagonIsotypic & PROVED & \checkmark & verified \\
\texttt{adjEigenvalue\_zero} & PentagonIsotypic & PROVED & \checkmark & verified \\
\texttt{adjacency\_apply} & PentagonIsotypic & PROVED & \checkmark & verified \\
\texttt{adjacency\_eigenmode} & PentagonIsotypic & PROVED & \checkmark & verified \\
\texttt{adjacency\_isotypicProjector} & PentagonIsotypic & PROVED & \checkmark & verified \\
\texttt{eigenmode\_linearIndependent} & PentagonIsotypic & PROVED & \checkmark & verified \\
\texttt{eigenmode\_ne\_zero} & PentagonIsotypic & PROVED & \checkmark & verified \\
\texttt{eigenvalues\_distinct} & PentagonIsotypic & PROVED & \checkmark & verified \\
\texttt{golden\_eigenfrequencies} & PentagonIsotypic & PROVED & \checkmark & verified \\
\texttt{golden\_eigenvector} & PentagonIsotypic & PROVED & \checkmark & verified \\
\texttt{isotypicProjector\_add} & PentagonIsotypic & PROVED & \checkmark & verified \\
\texttt{isotypicProjector\_comp} & PentagonIsotypic & PROVED & \checkmark & verified \\
\texttt{isotypicProjector\_completeness} & PentagonIsotypic & PROVED & \checkmark & verified \\
\texttt{isotypicProjector\_idempotent} & PentagonIsotypic & PROVED & \checkmark & verified \\
\texttt{isotypicProjector\_orthogonal} & PentagonIsotypic & PROVED & \checkmark & verified \\
\texttt{multiplicities} & PentagonIsotypic & PROVED & \checkmark & verified \\
\texttt{neg\_golden\_eigenfrequencies} & PentagonIsotypic & PROVED & \checkmark & verified \\
\texttt{rot\_isotypic} & PentagonIsotypic & PROVED & \checkmark & verified \\
\texttt{two\_eigenfrequency} & PentagonIsotypic & PROVED & \checkmark & verified \\
\texttt{adjL\_apply} & PentagonMultiplicities & PROVED & \checkmark & verified \\
\texttt{adjL\_eigenmode} & PentagonMultiplicities & PROVED & \checkmark & verified \\
\texttt{eigenBasis\_apply} & PentagonMultiplicities & PROVED & \checkmark & verified \\
\texttt{eigenIndices\_golden} & PentagonMultiplicities & PROVED & \checkmark & verified \\
\texttt{eigenIndices\_neg\_golden} & PentagonMultiplicities & PROVED & \checkmark & verified \\
\texttt{eigenIndices\_two} & PentagonMultiplicities & PROVED & \checkmark & verified \\
\texttt{eigenspace\_eq\_span\_group} & PentagonMultiplicities & PROVED & \checkmark & verified \\
\texttt{eigenspace\_golden\_eq} & PentagonMultiplicities & PROVED & \checkmark & verified \\
\texttt{eigenspace\_neg\_golden\_eq} & PentagonMultiplicities & PROVED & \checkmark & verified \\
\texttt{eigenspace\_two\_eq} & PentagonMultiplicities & PROVED & \checkmark & verified \\
\texttt{eigenspaces\_span\_top} & PentagonMultiplicities & PROVED & \checkmark & verified \\
\texttt{finrank\_eigenspace\_golden} & PentagonMultiplicities & PROVED & \checkmark & verified \\
\texttt{finrank\_eigenspace\_neg\_golden} & PentagonMultiplicities & PROVED & \checkmark & verified \\
\texttt{finrank\_eigenspace\_pair} & PentagonMultiplicities & PROVED & \checkmark & verified \\
\texttt{finrank\_eigenspace\_two} & PentagonMultiplicities & PROVED & \checkmark & verified \\
\texttt{finrank\_vertexSpace} & PentagonMultiplicities & PROVED & \checkmark & verified \\
\texttt{hasEigenvalue\_golden} & PentagonMultiplicities & PROVED & \checkmark & verified \\
\texttt{hasEigenvalue\_neg\_golden} & PentagonMultiplicities & PROVED & \checkmark & verified \\
\texttt{hasEigenvalue\_two} & PentagonMultiplicities & PROVED & \checkmark & verified \\
\texttt{hasEigenvector\_eigenmode} & PentagonMultiplicities & PROVED & \checkmark & verified \\
\texttt{multiplicities\_sum\_eq\_finrank} & PentagonMultiplicities & PROVED & \checkmark & verified \\
\texttt{range\_matrix\_two} & PentagonMultiplicities & PROVED & \checkmark & verified \\
\texttt{repr\_eq\_zero\_of\_ne} & PentagonMultiplicities & PROVED & \checkmark & verified \\
\texttt{concrete\_permInner\_golden} & PentagonTraceBridge & PROVED & \checkmark & verified \\
\texttt{d5Character\_eq\_permCharacter} & PentagonTraceBridge & PROVED & \checkmark & verified \\
\texttt{permCharacter\_eq\_trace\_d5PermutationMatrix} & PentagonTraceBridge & PROVED & \checkmark & verified \\
\texttt{partition\_zero\_card} & PentagonalPartition & PROVED & \checkmark & verified \\
\texttt{pent\_injective} & PentagonalPartition & PROVED & \checkmark & verified \\
\texttt{pent\_lt\_succ} & PentagonalPartition & PROVED & \checkmark & verified \\
\texttt{pent\_nonneg} & PentagonalPartition & PROVED & \checkmark & verified \\
\texttt{pent\_reflect} & PentagonalPartition & PROVED & \checkmark & verified \\
\texttt{pent\_succ} & PentagonalPartition & PROVED & \checkmark & verified \\
\texttt{pent\_values} & PentagonalPartition & PROVED & \checkmark & verified \\
\texttt{two\_dvd\_pentNum} & PentagonalPartition & PROVED & \checkmark & verified \\
\texttt{two\_mul\_pent} & PentagonalPartition & PROVED & \checkmark & verified \\
\texttt{two\_mul\_pent\_expand} & PentagonalPartition & PROVED & \checkmark & verified \\
\texttt{coeff\_genFun\_pstChar} & PentagonalTheoremFranklin & PROVED & \checkmark & verified \\
\texttt{genFun\_pstChar\_eq\_prod} & PentagonalTheoremFranklin & PROVED & \checkmark & verified \\
\texttt{natCast\_pentagonal\_eq\_pent} & PentagonalTheoremFranklin & PROVED & \checkmark & verified \\
\texttt{prod\_pstChar\_eq} & PentagonalTheoremFranklin & PROVED & \checkmark & verified \\
\texttt{not\_pt\_4} & PerfectTotient & PROVED & \checkmark & verified \\
\texttt{pt\_111} & PerfectTotient & PROVED & \checkmark & verified \\
\texttt{pt\_15} & PerfectTotient & PROVED & \checkmark & verified \\
\texttt{pt\_183} & PerfectTotient & PROVED & \checkmark & verified \\
\texttt{pt\_243} & PerfectTotient & PROVED & \checkmark & verified \\
\texttt{pt\_27} & PerfectTotient & PROVED & \checkmark & verified \\
\texttt{pt\_3} & PerfectTotient & PROVED & \checkmark & verified \\
\texttt{pt\_39} & PerfectTotient & PROVED & \checkmark & verified \\
\texttt{pt\_81} & PerfectTotient & PROVED & \checkmark & verified \\
\texttt{pt\_9} & PerfectTotient & PROVED & \checkmark & verified \\
\texttt{cousin\_13} & PolignacPrimes & PROVED & \checkmark & verified \\
\texttt{cousin\_19} & PolignacPrimes & PROVED & \checkmark & verified \\
\texttt{cousin\_3} & PolignacPrimes & PROVED & \checkmark & verified \\
\texttt{cousin\_37} & PolignacPrimes & PROVED & \checkmark & verified \\
\texttt{cousin\_43} & PolignacPrimes & PROVED & \checkmark & verified \\
\texttt{cousin\_67} & PolignacPrimes & PROVED & \checkmark & verified \\
\texttt{cousin\_7} & PolignacPrimes & PROVED & \checkmark & verified \\
\texttt{cousin\_79} & PolignacPrimes & PROVED & \checkmark & verified \\
\texttt{cousin\_form\_6k} & PolignacPrimes & PROVED & \checkmark & verified \\
\texttt{cousin\_mod\_six} & PolignacPrimes & PROVED & \checkmark & verified \\
\texttt{sexy\_11} & PolignacPrimes & PROVED & \checkmark & verified \\
\texttt{sexy\_13} & PolignacPrimes & PROVED & \checkmark & verified \\
\texttt{sexy\_17} & PolignacPrimes & PROVED & \checkmark & verified \\
\texttt{sexy\_23} & PolignacPrimes & PROVED & \checkmark & verified \\
\texttt{sexy\_31} & PolignacPrimes & PROVED & \checkmark & verified \\
\texttt{sexy\_47} & PolignacPrimes & PROVED & \checkmark & verified \\
\texttt{sexy\_5} & PolignacPrimes & PROVED & \checkmark & verified \\
\texttt{sexy\_7} & PolignacPrimes & PROVED & \checkmark & verified \\
\texttt{sexy\_same\_mod\_six} & PolignacPrimes & PROVED & \checkmark & verified \\
\texttt{almostPerfect\_1} & QuasiperfectNumbers & PROVED & \checkmark & verified \\
\texttt{almostPerfect\_16} & QuasiperfectNumbers & PROVED & \checkmark & verified \\
\texttt{almostPerfect\_2} & QuasiperfectNumbers & PROVED & \checkmark & verified \\
\texttt{almostPerfect\_pow\_two} & QuasiperfectNumbers & PROVED & \checkmark & verified \\
\texttt{not\_quasiperfect\_of\_almostPerfect} & QuasiperfectNumbers & PROVED & \checkmark & verified \\
\texttt{not\_quasiperfect\_of\_perfectσ} & QuasiperfectNumbers & PROVED & \checkmark & verified \\
\texttt{perfectσ\_28} & QuasiperfectNumbers & PROVED & \checkmark & verified \\
\texttt{perfectσ\_6} & QuasiperfectNumbers & PROVED & \checkmark & verified \\
\texttt{sigma\_two\_pow} & QuasiperfectNumbers & PROVED & \checkmark & verified \\
\texttt{coeff\_partitionGF} & RamanujanCongruence & PROVED & \checkmark & verified \\
\texttt{coeff\_partitionGF\_zero} & RamanujanCongruence & PROVED & \checkmark & verified \\
\texttt{five\_dvd\_of\_dissection} & RamanujanCongruence & PROVED & \checkmark & verified \\
\texttt{prime\_of\_repunit\_prime} & RepunitPrimes & PROVED & \checkmark & verified \\
\texttt{repunit\_1\_not\_prime} & RepunitPrimes & PROVED & \checkmark & verified \\
\texttt{repunit\_2\_prime} & RepunitPrimes & PROVED & \checkmark & verified \\
\texttt{repunit\_3\_not\_prime} & RepunitPrimes & PROVED & \checkmark & verified \\
\texttt{repunit\_4\_not\_prime} & RepunitPrimes & PROVED & \checkmark & verified \\
\texttt{repunit\_5\_not\_prime} & RepunitPrimes & PROVED & \checkmark & verified \\
\texttt{repunit\_6\_not\_prime} & RepunitPrimes & PROVED & \checkmark & verified \\
\texttt{repunit\_dvd\_of\_dvd} & RepunitPrimes & PROVED & \checkmark & verified \\
\texttt{repunit\_mul\_nine\_add\_one} & RepunitPrimes & PROVED & \checkmark & verified \\
\texttt{repunit\_strictMono} & RepunitPrimes & PROVED & \checkmark & verified \\
\texttt{Gammaℝ\_ne\_zero\_of\_nontrivial} & RiemannScaffold & PROVED & \checkmark & verified \\
\texttt{RH\_of\_BrockianSystem} & RiemannScaffold & CONDITIONAL & \checkmark & verified \\
\texttt{RiemannHypothesis\_of\_forall\_xi\_zero} & RiemannScaffold & PROVED & \checkmark & verified \\
\texttt{riemannXi\_eq\_zero\_of\_nontrivial\_zeta\_zero} & RiemannScaffold & PROVED & \checkmark & verified \\
\texttt{symmetric\_eigenvalue\_im\_zero} & RiemannScaffold & PROVED & \checkmark & verified \\
\texttt{completedRiemannZeta\_functional\_equation} & RiemannXiFunctionalEquation & PROVED & \checkmark & verified \\
\texttt{riemannXi\_one\_sub} & RiemannXiFunctionalEquation & PROVED & \checkmark & verified \\
\texttt{riemannXi\_one\_sub\_eq\_zero} & RiemannXiFunctionalEquation & PROVED & \checkmark & verified \\
\texttt{riemannXi\_one\_sub\_eq\_zero\_iff} & RiemannXiFunctionalEquation & PROVED & \checkmark & verified \\
\texttt{criticalLine\_re\_iff\_reflect\_re\_eq} & RiemannXiSymmetry & PROVED & \checkmark & verified \\
\texttt{reflect\_fixed\_iff} & RiemannXiSymmetry & PROVED & \checkmark & verified \\
\texttt{reflect\_re} & RiemannXiSymmetry & PROVED & \checkmark & verified \\
\texttt{reflect\_re\_eq\_half\_iff} & RiemannXiSymmetry & PROVED & \checkmark & verified \\
\texttt{reflect\_re\_eq\_self\_iff} & RiemannXiSymmetry & PROVED & \checkmark & verified \\
\texttt{reflect\_reflect} & RiemannXiSymmetry & PROVED & \checkmark & verified \\
\texttt{riemannXi\_eq\_zero\_iff\_reflect} & RiemannXiSymmetry & PROVED & \checkmark & verified \\
\texttt{riemannXi\_eq\_zero\_reflect} & RiemannXiSymmetry & PROVED & \checkmark & verified \\
\texttt{riemannXi\_reflect} & RiemannXiSymmetry & PROVED & \checkmark & verified \\
\texttt{riemannXi\_reflect\_eq\_zero\_iff} & RiemannXiSymmetry & PROVED & \checkmark & verified \\
\texttt{riemannXi\_reflect\_zero\_and\_criticalLine} & RiemannXiSymmetry & PROVED & \checkmark & verified \\
\texttt{riemannXi\_zeroSet\_image\_reflect} & RiemannXiSymmetry & PROVED & \checkmark & verified \\
\texttt{riemannXi\_zeroSet\_preimage\_reflect} & RiemannXiSymmetry & PROVED & \checkmark & verified \\
\texttt{riemannXi\_zero\_pair\_of\_reflect\_zero} & RiemannXiSymmetry & PROVED & \checkmark & verified \\
\texttt{riemannXi\_zero\_pair\_of\_zero} & RiemannXiSymmetry & PROVED & \checkmark & verified \\
\texttt{composite\_of} & RieselCovering & PROVED & \checkmark & verified \\
\texttt{covering\_table} & RieselCovering & PROVED & \checkmark & verified \\
\texttt{riesel\_509203} & RieselCovering & PROVED & \checkmark & verified \\
\texttt{two\_pow\_periodic} & RieselCovering & PROVED & \checkmark & verified \\
\texttt{not\_ruthAaron\_6} & RuthAaronPairs & PROVED & \checkmark & verified \\
\texttt{ruthAaron\_125} & RuthAaronPairs & PROVED & \checkmark & verified \\
\texttt{ruthAaron\_5} & RuthAaronPairs & PROVED & \checkmark & verified \\
\texttt{ruthAaron\_714} & RuthAaronPairs & PROVED & \checkmark & verified \\
\texttt{ruthAaron\_77} & RuthAaronPairs & PROVED & \checkmark & verified \\
\texttt{sopfr\_125} & RuthAaronPairs & PROVED & \checkmark & verified \\
\texttt{sopfr\_126} & RuthAaronPairs & PROVED & \checkmark & verified \\
\texttt{sopfr\_5} & RuthAaronPairs & PROVED & \checkmark & verified \\
\texttt{sopfr\_6} & RuthAaronPairs & PROVED & \checkmark & verified \\
\texttt{sopfr\_7} & RuthAaronPairs & PROVED & \checkmark & verified \\
\texttt{sopfr\_714} & RuthAaronPairs & PROVED & \checkmark & verified \\
\texttt{sopfr\_715} & RuthAaronPairs & PROVED & \checkmark & verified \\
\texttt{sopfr\_77} & RuthAaronPairs & PROVED & \checkmark & verified \\
\texttt{sopfr\_78} & RuthAaronPairs & PROVED & \checkmark & verified \\
\texttt{sopfr\_eq} & RuthAaronPairs & PROVED & \checkmark & verified \\
\texttt{brockian\_sanity} & Sanity & PROVED & \checkmark & verified \\
\texttt{composite\_of} & SierpinskiCovering & PROVED & \checkmark & verified \\
\texttt{covering\_table} & SierpinskiCovering & PROVED & \checkmark & verified \\
\texttt{sierpinski\_78557} & SierpinskiCovering & PROVED & \checkmark & verified \\
\texttt{two\_pow\_periodic} & SierpinskiCovering & PROVED & \checkmark & verified \\
\texttt{H3\_det} & Sieve & PROVED & \checkmark & verified \\
\texttt{H3\_ground} & Sieve & PROVED & \checkmark & verified \\
\texttt{H3\_middle} & Sieve & PROVED & \checkmark & verified \\
\texttt{H3\_top} & Sieve & PROVED & \checkmark & verified \\
\texttt{H3\_trace} & Sieve & PROVED & \checkmark & verified \\
\texttt{compatible\_closure} & Sieve & PROVED & \checkmark & verified \\
\texttt{no\_adjacent\_admissible} & Sieve & PROVED & \checkmark & verified \\
\texttt{run3\_signature} & Sieve & PROVED & \checkmark & verified \\
\texttt{run\_cap} & Sieve & PROVED & \checkmark & verified \\
\texttt{silver\_gap\_rigidity\_finite} & Sieve & PROVED & \checkmark & verified \\
\texttt{tau\_injective\_on\_period} & Sieve & PROVED & \checkmark & verified \\
\texttt{tau\_minimal\_period} & Sieve & PROVED & \checkmark & verified \\
\texttt{tau\_period} & Sieve & PROVED & \checkmark & verified \\
\texttt{twin\_admissible\_card} & Sieve & PROVED & \checkmark & verified \\
\texttt{twin\_pins\_mod\_three} & Sieve & PROVED & \checkmark & verified \\
\texttt{localFactorAt\_eq} & SingularSeries & PROVED & \checkmark & verified \\
\texttt{localFactorAt\_of\_not\_prime} & SingularSeries & PROVED & \checkmark & verified \\
\texttt{localFactorAt\_pos} & SingularSeries & PROVED & \checkmark & verified \\
\texttt{local\_factor\_asymptotic} & SingularSeries & PROVED & \checkmark & verified \\
\texttt{local\_factor\_denom\_ne\_zero} & SingularSeries & PROVED & \checkmark & verified \\
\texttt{local\_factor\_of\_nu\_p\_eq\_p} & SingularSeries & PROVED & \checkmark & verified \\
\texttt{local\_factor\_of\_nu\_p\_eq\_zero} & SingularSeries & PROVED & \checkmark & verified \\
\texttt{local\_factor\_pos} & SingularSeries & PROVED & \checkmark & verified \\
\texttt{nu\_p'} & SingularSeries & PROVED & \checkmark & verified \\
\texttt{nu\_p\_eq\_image\_card} & SingularSeries & PROVED & \checkmark & verified \\
\texttt{nu\_p\_lt\_p\_of\_admissible} & SingularSeries & PROVED & \checkmark & verified \\
\texttt{singular\_series\_finite\_pos} & SingularSeries & PROVED & \checkmark & verified \\
\texttt{singular\_series\_pos} & SingularSeries & PROVED & \checkmark & verified \\
\texttt{err\_bound} & Convergence & PROVED & \checkmark & verified \\
\texttt{localFactor\_sub\_one\_bound} & Convergence & PROVED & \checkmark & verified \\
\texttt{nu\_p\_eq\_card\_of\_lt} & Convergence & PROVED & \checkmark & verified \\
\texttt{singularSeriesFinite\_tendsto\_pos} & Convergence & PROVED & \checkmark & verified \\
\texttt{singular\_series\_pos'} & Convergence & PROVED & \checkmark & verified \\
\texttt{summable\_localFactorAt\_sub\_one} & Convergence & PROVED & \checkmark & verified \\
\texttt{evenPair\_card\_eighteen} & EvenMore & PROVED & \checkmark & verified \\
\texttt{evenPair\_card\_fourteen} & EvenMore & PROVED & \checkmark & verified \\
\texttt{evenPair\_card\_sixteen} & EvenMore & PROVED & \checkmark & verified \\
\texttt{evenPair\_card\_twelve} & EvenMore & PROVED & \checkmark & verified \\
\texttt{evenPair\_card\_twenty} & EvenMore & PROVED & \checkmark & verified \\
\texttt{isAdmissible\_evenPair\_eighteen} & EvenMore & PROVED & \checkmark & verified \\
\texttt{isAdmissible\_evenPair\_fourteen} & EvenMore & PROVED & \checkmark & verified \\
\texttt{isAdmissible\_evenPair\_sixteen} & EvenMore & PROVED & \checkmark & verified \\
\texttt{isAdmissible\_evenPair\_twelve} & EvenMore & PROVED & \checkmark & verified \\
\texttt{isAdmissible\_evenPair\_twenty} & EvenMore & PROVED & \checkmark & verified \\
\texttt{localFactorAt\_eighteen\_two} & EvenMore & PROVED & \checkmark & verified \\
\texttt{localFactorAt\_fourteen\_two} & EvenMore & PROVED & \checkmark & verified \\
\texttt{localFactorAt\_sixteen\_two} & EvenMore & PROVED & \checkmark & verified \\
\texttt{localFactorAt\_twelve\_five} & EvenMore & PROVED & \checkmark & verified \\
\texttt{localFactorAt\_twelve\_odd\_ne\_three} & EvenMore & PROVED & \checkmark & verified \\
\texttt{localFactorAt\_twelve\_three} & EvenMore & PROVED & \checkmark & verified \\
\texttt{localFactorAt\_twelve\_two} & EvenMore & PROVED & \checkmark & verified \\
\texttt{localFactorAt\_twenty\_two} & EvenMore & PROVED & \checkmark & verified \\
\texttt{localFactor\_eighteen\_two} & EvenMore & PROVED & \checkmark & verified \\
\texttt{localFactor\_fourteen\_two} & EvenMore & PROVED & \checkmark & verified \\
\texttt{localFactor\_sixteen\_two} & EvenMore & PROVED & \checkmark & verified \\
\texttt{localFactor\_twelve\_five} & EvenMore & PROVED & \checkmark & verified \\
\texttt{localFactor\_twelve\_odd\_ne\_three} & EvenMore & PROVED & \checkmark & verified \\
\texttt{localFactor\_twelve\_three} & EvenMore & PROVED & \checkmark & verified \\
\texttt{localFactor\_twelve\_two} & EvenMore & PROVED & \checkmark & verified \\
\texttt{localFactor\_twenty\_two} & EvenMore & PROVED & \checkmark & verified \\
\texttt{nu\_p\_eighteen} & EvenMore & PROVED & \checkmark & verified \\
\texttt{nu\_p\_eighteen\_two} & EvenMore & PROVED & \checkmark & verified \\
\texttt{nu\_p\_fourteen} & EvenMore & PROVED & \checkmark & verified \\
\texttt{nu\_p\_fourteen\_two} & EvenMore & PROVED & \checkmark & verified \\
\texttt{nu\_p\_sixteen} & EvenMore & PROVED & \checkmark & verified \\
\texttt{nu\_p\_sixteen\_two} & EvenMore & PROVED & \checkmark & verified \\
\texttt{nu\_p\_twelve} & EvenMore & PROVED & \checkmark & verified \\
\texttt{nu\_p\_twelve\_odd\_ne\_three} & EvenMore & PROVED & \checkmark & verified \\
\texttt{nu\_p\_twelve\_three} & EvenMore & PROVED & \checkmark & verified \\
\texttt{nu\_p\_twelve\_two} & EvenMore & PROVED & \checkmark & verified \\
\texttt{nu\_p\_twenty} & EvenMore & PROVED & \checkmark & verified \\
\texttt{nu\_p\_twenty\_two} & EvenMore & PROVED & \checkmark & verified \\
\texttt{singular\_series\_finite\_pos\_evenPair\_eighteen} & EvenMore & PROVED & \checkmark & verified \\
\texttt{singular\_series\_finite\_pos\_evenPair\_fourteen} & EvenMore & PROVED & \checkmark & verified \\
\texttt{singular\_series\_finite\_pos\_evenPair\_sixteen} & EvenMore & PROVED & \checkmark & verified \\
\texttt{singular\_series\_finite\_pos\_evenPair\_twelve} & EvenMore & PROVED & \checkmark & verified \\
\texttt{singular\_series\_finite\_pos\_evenPair\_twenty} & EvenMore & PROVED & \checkmark & verified \\
\texttt{singular\_series\_pos\_evenPair\_eighteen} & EvenMore & PROVED & \checkmark & verified \\
\texttt{singular\_series\_pos\_evenPair\_fourteen} & EvenMore & PROVED & \checkmark & verified \\
\texttt{singular\_series\_pos\_evenPair\_sixteen} & EvenMore & PROVED & \checkmark & verified \\
\texttt{singular\_series\_pos\_evenPair\_twelve} & EvenMore & PROVED & \checkmark & verified \\
\texttt{singular\_series\_pos\_evenPair\_twenty} & EvenMore & PROVED & \checkmark & verified \\
\texttt{evenPair\_card\_le\_two} & Examples & PROVED & \checkmark & verified \\
\texttt{isAdmissible\_evenPair} & Examples & PROVED & \checkmark & verified \\
\texttt{isAdmissible\_evenPair\_four} & Examples & PROVED & \checkmark & verified \\
\texttt{isAdmissible\_evenPair\_six} & Examples & PROVED & \checkmark & verified \\
\texttt{isAdmissible\_evenPair\_two} & Examples & PROVED & \checkmark & verified \\
\texttt{isAdmissible\_twinGap} & Examples & PROVED & \checkmark & verified \\
\texttt{singular\_series\_finite\_pos\_twinGap} & Examples & PROVED & \checkmark & verified \\
\texttt{singular\_series\_pos\_evenPair} & Examples & PROVED & \checkmark & verified \\
\texttt{singular\_series\_pos\_evenPair\_four} & Examples & PROVED & \checkmark & verified \\
\texttt{singular\_series\_pos\_evenPair\_six} & Examples & PROVED & \checkmark & verified \\
\texttt{singular\_series\_pos\_evenPair\_two} & Examples & PROVED & \checkmark & verified \\
\texttt{singular\_series\_pos\_twinGap} & Examples & PROVED & \checkmark & verified \\
\texttt{twinGap\_eq} & Examples & PROVED & \checkmark & verified \\
\texttt{evenPair\_card\_oneThousandEight} & Gaps10021010 & PROVED & \checkmark & verified \\
\texttt{evenPair\_card\_oneThousandFour} & Gaps10021010 & PROVED & \checkmark & verified \\
\texttt{evenPair\_card\_oneThousandSix} & Gaps10021010 & PROVED & \checkmark & verified \\
\texttt{evenPair\_card\_oneThousandTen} & Gaps10021010 & PROVED & \checkmark & verified \\
\texttt{evenPair\_card\_oneThousandTwo} & Gaps10021010 & PROVED & \checkmark & verified \\
\texttt{isAdmissible\_evenPair\_oneThousandEight} & Gaps10021010 & PROVED & \checkmark & verified \\
\texttt{isAdmissible\_evenPair\_oneThousandFour} & Gaps10021010 & PROVED & \checkmark & verified \\
\texttt{isAdmissible\_evenPair\_oneThousandSix} & Gaps10021010 & PROVED & \checkmark & verified \\
\texttt{isAdmissible\_evenPair\_oneThousandTen} & Gaps10021010 & PROVED & \checkmark & verified \\
\texttt{isAdmissible\_evenPair\_oneThousandTwo} & Gaps10021010 & PROVED & \checkmark & verified \\
\texttt{localFactor\_oneThousandTen\_two} & Gaps10021010 & PROVED & \checkmark & verified \\
\texttt{localFactor\_oneThousandTwo\_two} & Gaps10021010 & PROVED & \checkmark & verified \\
\texttt{nu\_p\_oneThousandEight} & Gaps10021010 & PROVED & \checkmark & verified \\
\texttt{nu\_p\_oneThousandFour} & Gaps10021010 & PROVED & \checkmark & verified \\
\texttt{nu\_p\_oneThousandSix} & Gaps10021010 & PROVED & \checkmark & verified \\
\texttt{nu\_p\_oneThousandTen} & Gaps10021010 & PROVED & \checkmark & verified \\
\texttt{nu\_p\_oneThousandTen\_two} & Gaps10021010 & PROVED & \checkmark & verified \\
\texttt{nu\_p\_oneThousandTwo} & Gaps10021010 & PROVED & \checkmark & verified \\
\texttt{nu\_p\_oneThousandTwo\_two} & Gaps10021010 & PROVED & \checkmark & verified \\
\texttt{singular\_series\_finite\_pos\_evenPair\_oneThousandEight} & Gaps10021010 & PROVED & \checkmark & verified \\
\texttt{singular\_series\_finite\_pos\_evenPair\_oneThousandFour} & Gaps10021010 & PROVED & \checkmark & verified \\
\texttt{singular\_series\_finite\_pos\_evenPair\_oneThousandSix} & Gaps10021010 & PROVED & \checkmark & verified \\
\texttt{singular\_series\_finite\_pos\_evenPair\_oneThousandTen} & Gaps10021010 & PROVED & \checkmark & verified \\
\texttt{singular\_series\_finite\_pos\_evenPair\_oneThousandTwo} & Gaps10021010 & PROVED & \checkmark & verified \\
\texttt{singular\_series\_pos\_evenPair\_oneThousandEight} & Gaps10021010 & PROVED & \checkmark & verified \\
\texttt{singular\_series\_pos\_evenPair\_oneThousandFour} & Gaps10021010 & PROVED & \checkmark & verified \\
\texttt{singular\_series\_pos\_evenPair\_oneThousandSix} & Gaps10021010 & PROVED & \checkmark & verified \\
\texttt{singular\_series\_pos\_evenPair\_oneThousandTen} & Gaps10021010 & PROVED & \checkmark & verified \\
\texttt{singular\_series\_pos\_evenPair\_oneThousandTwo} & Gaps10021010 & PROVED & \checkmark & verified \\
\texttt{evenPair\_card\_oneThousandEighteen} & Gaps10121020 & PROVED & \checkmark & verified \\
\texttt{evenPair\_card\_oneThousandFourteen} & Gaps10121020 & PROVED & \checkmark & verified \\
\texttt{evenPair\_card\_oneThousandSixteen} & Gaps10121020 & PROVED & \checkmark & verified \\
\texttt{evenPair\_card\_oneThousandTwelve} & Gaps10121020 & PROVED & \checkmark & verified \\
\texttt{evenPair\_card\_oneThousandTwenty} & Gaps10121020 & PROVED & \checkmark & verified \\
\texttt{isAdmissible\_evenPair\_oneThousandEighteen} & Gaps10121020 & PROVED & \checkmark & verified \\
\texttt{isAdmissible\_evenPair\_oneThousandFourteen} & Gaps10121020 & PROVED & \checkmark & verified \\
\texttt{isAdmissible\_evenPair\_oneThousandSixteen} & Gaps10121020 & PROVED & \checkmark & verified \\
\texttt{isAdmissible\_evenPair\_oneThousandTwelve} & Gaps10121020 & PROVED & \checkmark & verified \\
\texttt{isAdmissible\_evenPair\_oneThousandTwenty} & Gaps10121020 & PROVED & \checkmark & verified \\
\texttt{localFactor\_oneThousandTwelve\_two} & Gaps10121020 & PROVED & \checkmark & verified \\
\texttt{localFactor\_oneThousandTwenty\_two} & Gaps10121020 & PROVED & \checkmark & verified \\
\texttt{nu\_p\_oneThousandEighteen} & Gaps10121020 & PROVED & \checkmark & verified \\
\texttt{nu\_p\_oneThousandFourteen} & Gaps10121020 & PROVED & \checkmark & verified \\
\texttt{nu\_p\_oneThousandSixteen} & Gaps10121020 & PROVED & \checkmark & verified \\
\texttt{nu\_p\_oneThousandTwelve} & Gaps10121020 & PROVED & \checkmark & verified \\
\texttt{nu\_p\_oneThousandTwelve\_two} & Gaps10121020 & PROVED & \checkmark & verified \\
\texttt{nu\_p\_oneThousandTwenty} & Gaps10121020 & PROVED & \checkmark & verified \\
\texttt{nu\_p\_oneThousandTwenty\_two} & Gaps10121020 & PROVED & \checkmark & verified \\
\texttt{singular\_series\_finite\_pos\_evenPair\_oneThousandEighteen} & Gaps10121020 & PROVED & \checkmark & verified \\
\texttt{singular\_series\_finite\_pos\_evenPair\_oneThousandFourteen} & Gaps10121020 & PROVED & \checkmark & verified \\
\texttt{singular\_series\_finite\_pos\_evenPair\_oneThousandSixteen} & Gaps10121020 & PROVED & \checkmark & verified \\
\texttt{singular\_series\_finite\_pos\_evenPair\_oneThousandTwelve} & Gaps10121020 & PROVED & \checkmark & verified \\
\texttt{singular\_series\_finite\_pos\_evenPair\_oneThousandTwenty} & Gaps10121020 & PROVED & \checkmark & verified \\
\texttt{singular\_series\_pos\_evenPair\_oneThousandEighteen} & Gaps10121020 & PROVED & \checkmark & verified \\
\texttt{singular\_series\_pos\_evenPair\_oneThousandFourteen} & Gaps10121020 & PROVED & \checkmark & verified \\
\texttt{singular\_series\_pos\_evenPair\_oneThousandSixteen} & Gaps10121020 & PROVED & \checkmark & verified \\
\texttt{singular\_series\_pos\_evenPair\_oneThousandTwelve} & Gaps10121020 & PROVED & \checkmark & verified \\
\texttt{singular\_series\_pos\_evenPair\_oneThousandTwenty} & Gaps10121020 & PROVED & \checkmark & verified \\
\texttt{evenPair\_card\_oneHundredEight} & Gaps102110 & PROVED & \checkmark & verified \\
\texttt{evenPair\_card\_oneHundredFour} & Gaps102110 & PROVED & \checkmark & verified \\
\texttt{evenPair\_card\_oneHundredSix} & Gaps102110 & PROVED & \checkmark & verified \\
\texttt{evenPair\_card\_oneHundredTen} & Gaps102110 & PROVED & \checkmark & verified \\
\texttt{evenPair\_card\_oneHundredTwo} & Gaps102110 & PROVED & \checkmark & verified \\
\texttt{isAdmissible\_evenPair\_oneHundredEight} & Gaps102110 & PROVED & \checkmark & verified \\
\texttt{isAdmissible\_evenPair\_oneHundredFour} & Gaps102110 & PROVED & \checkmark & verified \\
\texttt{isAdmissible\_evenPair\_oneHundredSix} & Gaps102110 & PROVED & \checkmark & verified \\
\texttt{isAdmissible\_evenPair\_oneHundredTen} & Gaps102110 & PROVED & \checkmark & verified \\
\texttt{isAdmissible\_evenPair\_oneHundredTwo} & Gaps102110 & PROVED & \checkmark & verified \\
\texttt{localFactor\_oneHundredTen\_two} & Gaps102110 & PROVED & \checkmark & verified \\
\texttt{localFactor\_oneHundredTwo\_two} & Gaps102110 & PROVED & \checkmark & verified \\
\texttt{nu\_p\_oneHundredEight} & Gaps102110 & PROVED & \checkmark & verified \\
\texttt{nu\_p\_oneHundredFour} & Gaps102110 & PROVED & \checkmark & verified \\
\texttt{nu\_p\_oneHundredSix} & Gaps102110 & PROVED & \checkmark & verified \\
\texttt{nu\_p\_oneHundredTen} & Gaps102110 & PROVED & \checkmark & verified \\
\texttt{nu\_p\_oneHundredTen\_two} & Gaps102110 & PROVED & \checkmark & verified \\
\texttt{nu\_p\_oneHundredTwo} & Gaps102110 & PROVED & \checkmark & verified \\
\texttt{nu\_p\_oneHundredTwo\_two} & Gaps102110 & PROVED & \checkmark & verified \\
\texttt{singular\_series\_finite\_pos\_evenPair\_oneHundredEight} & Gaps102110 & PROVED & \checkmark & verified \\
\texttt{singular\_series\_finite\_pos\_evenPair\_oneHundredFour} & Gaps102110 & PROVED & \checkmark & verified \\
\texttt{singular\_series\_finite\_pos\_evenPair\_oneHundredSix} & Gaps102110 & PROVED & \checkmark & verified \\
\texttt{singular\_series\_finite\_pos\_evenPair\_oneHundredTen} & Gaps102110 & PROVED & \checkmark & verified \\
\texttt{singular\_series\_finite\_pos\_evenPair\_oneHundredTwo} & Gaps102110 & PROVED & \checkmark & verified \\
\texttt{singular\_series\_pos\_evenPair\_oneHundredEight} & Gaps102110 & PROVED & \checkmark & verified \\
\texttt{singular\_series\_pos\_evenPair\_oneHundredFour} & Gaps102110 & PROVED & \checkmark & verified \\
\texttt{singular\_series\_pos\_evenPair\_oneHundredSix} & Gaps102110 & PROVED & \checkmark & verified \\
\texttt{singular\_series\_pos\_evenPair\_oneHundredTen} & Gaps102110 & PROVED & \checkmark & verified \\
\texttt{singular\_series\_pos\_evenPair\_oneHundredTwo} & Gaps102110 & PROVED & \checkmark & verified \\
\texttt{evenPair\_card\_oneThousandThirty} & Gaps10221030 & PROVED & \checkmark & verified \\
\texttt{evenPair\_card\_oneThousandTwentyEight} & Gaps10221030 & PROVED & \checkmark & verified \\
\texttt{evenPair\_card\_oneThousandTwentyFour} & Gaps10221030 & PROVED & \checkmark & verified \\
\texttt{evenPair\_card\_oneThousandTwentySix} & Gaps10221030 & PROVED & \checkmark & verified \\
\texttt{evenPair\_card\_oneThousandTwentyTwo} & Gaps10221030 & PROVED & \checkmark & verified \\
\texttt{isAdmissible\_evenPair\_oneThousandThirty} & Gaps10221030 & PROVED & \checkmark & verified \\
\texttt{isAdmissible\_evenPair\_oneThousandTwentyEight} & Gaps10221030 & PROVED & \checkmark & verified \\
\texttt{isAdmissible\_evenPair\_oneThousandTwentyFour} & Gaps10221030 & PROVED & \checkmark & verified \\
\texttt{isAdmissible\_evenPair\_oneThousandTwentySix} & Gaps10221030 & PROVED & \checkmark & verified \\
\texttt{isAdmissible\_evenPair\_oneThousandTwentyTwo} & Gaps10221030 & PROVED & \checkmark & verified \\
\texttt{localFactor\_oneThousandThirty\_two} & Gaps10221030 & PROVED & \checkmark & verified \\
\texttt{localFactor\_oneThousandTwentyTwo\_two} & Gaps10221030 & PROVED & \checkmark & verified \\
\texttt{nu\_p\_oneThousandThirty} & Gaps10221030 & PROVED & \checkmark & verified \\
\texttt{nu\_p\_oneThousandThirty\_two} & Gaps10221030 & PROVED & \checkmark & verified \\
\texttt{nu\_p\_oneThousandTwentyEight} & Gaps10221030 & PROVED & \checkmark & verified \\
\texttt{nu\_p\_oneThousandTwentyFour} & Gaps10221030 & PROVED & \checkmark & verified \\
\texttt{nu\_p\_oneThousandTwentySix} & Gaps10221030 & PROVED & \checkmark & verified \\
\texttt{nu\_p\_oneThousandTwentyTwo} & Gaps10221030 & PROVED & \checkmark & verified \\
\texttt{nu\_p\_oneThousandTwentyTwo\_two} & Gaps10221030 & PROVED & \checkmark & verified \\
\texttt{singular\_series\_finite\_pos\_evenPair\_oneThousandThirty} & Gaps10221030 & PROVED & \checkmark & verified \\
\texttt{singular\_series\_finite\_pos\_evenPair\_oneThousandTwentyEight} & Gaps10221030 & PROVED & \checkmark & verified \\
\texttt{singular\_series\_finite\_pos\_evenPair\_oneThousandTwentyFour} & Gaps10221030 & PROVED & \checkmark & verified \\
\texttt{singular\_series\_finite\_pos\_evenPair\_oneThousandTwentySix} & Gaps10221030 & PROVED & \checkmark & verified \\
\texttt{singular\_series\_finite\_pos\_evenPair\_oneThousandTwentyTwo} & Gaps10221030 & PROVED & \checkmark & verified \\
\texttt{singular\_series\_pos\_evenPair\_oneThousandThirty} & Gaps10221030 & PROVED & \checkmark & verified \\
\texttt{singular\_series\_pos\_evenPair\_oneThousandTwentyEight} & Gaps10221030 & PROVED & \checkmark & verified \\
\texttt{singular\_series\_pos\_evenPair\_oneThousandTwentyFour} & Gaps10221030 & PROVED & \checkmark & verified \\
\texttt{singular\_series\_pos\_evenPair\_oneThousandTwentySix} & Gaps10221030 & PROVED & \checkmark & verified \\
\texttt{singular\_series\_pos\_evenPair\_oneThousandTwentyTwo} & Gaps10221030 & PROVED & \checkmark & verified \\
\texttt{evenPair\_card\_oneThousandForty} & Gaps10321040 & PROVED & \checkmark & verified \\
\texttt{evenPair\_card\_oneThousandThirtyEight} & Gaps10321040 & PROVED & \checkmark & verified \\
\texttt{evenPair\_card\_oneThousandThirtyFour} & Gaps10321040 & PROVED & \checkmark & verified \\
\texttt{evenPair\_card\_oneThousandThirtySix} & Gaps10321040 & PROVED & \checkmark & verified \\
\texttt{evenPair\_card\_oneThousandThirtyTwo} & Gaps10321040 & PROVED & \checkmark & verified \\
\texttt{isAdmissible\_evenPair\_oneThousandForty} & Gaps10321040 & PROVED & \checkmark & verified \\
\texttt{isAdmissible\_evenPair\_oneThousandThirtyEight} & Gaps10321040 & PROVED & \checkmark & verified \\
\texttt{isAdmissible\_evenPair\_oneThousandThirtyFour} & Gaps10321040 & PROVED & \checkmark & verified \\
\texttt{isAdmissible\_evenPair\_oneThousandThirtySix} & Gaps10321040 & PROVED & \checkmark & verified \\
\texttt{isAdmissible\_evenPair\_oneThousandThirtyTwo} & Gaps10321040 & PROVED & \checkmark & verified \\
\texttt{localFactor\_oneThousandForty\_two} & Gaps10321040 & PROVED & \checkmark & verified \\
\texttt{localFactor\_oneThousandThirtyTwo\_two} & Gaps10321040 & PROVED & \checkmark & verified \\
\texttt{nu\_p\_oneThousandForty} & Gaps10321040 & PROVED & \checkmark & verified \\
\texttt{nu\_p\_oneThousandForty\_two} & Gaps10321040 & PROVED & \checkmark & verified \\
\texttt{nu\_p\_oneThousandThirtyEight} & Gaps10321040 & PROVED & \checkmark & verified \\
\texttt{nu\_p\_oneThousandThirtyFour} & Gaps10321040 & PROVED & \checkmark & verified \\
\texttt{nu\_p\_oneThousandThirtySix} & Gaps10321040 & PROVED & \checkmark & verified \\
\texttt{nu\_p\_oneThousandThirtyTwo} & Gaps10321040 & PROVED & \checkmark & verified \\
\texttt{nu\_p\_oneThousandThirtyTwo\_two} & Gaps10321040 & PROVED & \checkmark & verified \\
\texttt{singular\_series\_finite\_pos\_evenPair\_oneThousandForty} & Gaps10321040 & PROVED & \checkmark & verified \\
\texttt{singular\_series\_finite\_pos\_evenPair\_oneThousandThirtyEight} & Gaps10321040 & PROVED & \checkmark & verified \\
\texttt{singular\_series\_finite\_pos\_evenPair\_oneThousandThirtyFour} & Gaps10321040 & PROVED & \checkmark & verified \\
\texttt{singular\_series\_finite\_pos\_evenPair\_oneThousandThirtySix} & Gaps10321040 & PROVED & \checkmark & verified \\
\texttt{singular\_series\_finite\_pos\_evenPair\_oneThousandThirtyTwo} & Gaps10321040 & PROVED & \checkmark & verified \\
\texttt{singular\_series\_pos\_evenPair\_oneThousandForty} & Gaps10321040 & PROVED & \checkmark & verified \\
\texttt{singular\_series\_pos\_evenPair\_oneThousandThirtyEight} & Gaps10321040 & PROVED & \checkmark & verified \\
\texttt{singular\_series\_pos\_evenPair\_oneThousandThirtyFour} & Gaps10321040 & PROVED & \checkmark & verified \\
\texttt{singular\_series\_pos\_evenPair\_oneThousandThirtySix} & Gaps10321040 & PROVED & \checkmark & verified \\
\texttt{singular\_series\_pos\_evenPair\_oneThousandThirtyTwo} & Gaps10321040 & PROVED & \checkmark & verified \\
\texttt{evenPair\_card\_oneThousandFifty} & Gaps10421050 & PROVED & \checkmark & verified \\
\texttt{evenPair\_card\_oneThousandFortyEight} & Gaps10421050 & PROVED & \checkmark & verified \\
\texttt{evenPair\_card\_oneThousandFortyFour} & Gaps10421050 & PROVED & \checkmark & verified \\
\texttt{evenPair\_card\_oneThousandFortySix} & Gaps10421050 & PROVED & \checkmark & verified \\
\texttt{evenPair\_card\_oneThousandFortyTwo} & Gaps10421050 & PROVED & \checkmark & verified \\
\texttt{isAdmissible\_evenPair\_oneThousandFifty} & Gaps10421050 & PROVED & \checkmark & verified \\
\texttt{isAdmissible\_evenPair\_oneThousandFortyEight} & Gaps10421050 & PROVED & \checkmark & verified \\
\texttt{isAdmissible\_evenPair\_oneThousandFortyFour} & Gaps10421050 & PROVED & \checkmark & verified \\
\texttt{isAdmissible\_evenPair\_oneThousandFortySix} & Gaps10421050 & PROVED & \checkmark & verified \\
\texttt{isAdmissible\_evenPair\_oneThousandFortyTwo} & Gaps10421050 & PROVED & \checkmark & verified \\
\texttt{localFactor\_oneThousandFifty\_two} & Gaps10421050 & PROVED & \checkmark & verified \\
\texttt{localFactor\_oneThousandFortyTwo\_two} & Gaps10421050 & PROVED & \checkmark & verified \\
\texttt{nu\_p\_oneThousandFifty} & Gaps10421050 & PROVED & \checkmark & verified \\
\texttt{nu\_p\_oneThousandFifty\_two} & Gaps10421050 & PROVED & \checkmark & verified \\
\texttt{nu\_p\_oneThousandFortyEight} & Gaps10421050 & PROVED & \checkmark & verified \\
\texttt{nu\_p\_oneThousandFortyFour} & Gaps10421050 & PROVED & \checkmark & verified \\
\texttt{nu\_p\_oneThousandFortySix} & Gaps10421050 & PROVED & \checkmark & verified \\
\texttt{nu\_p\_oneThousandFortyTwo} & Gaps10421050 & PROVED & \checkmark & verified \\
\texttt{nu\_p\_oneThousandFortyTwo\_two} & Gaps10421050 & PROVED & \checkmark & verified \\
\texttt{singular\_series\_finite\_pos\_evenPair\_oneThousandFifty} & Gaps10421050 & PROVED & \checkmark & verified \\
\texttt{singular\_series\_finite\_pos\_evenPair\_oneThousandFortyEight} & Gaps10421050 & PROVED & \checkmark & verified \\
\texttt{singular\_series\_finite\_pos\_evenPair\_oneThousandFortyFour} & Gaps10421050 & PROVED & \checkmark & verified \\
\texttt{singular\_series\_finite\_pos\_evenPair\_oneThousandFortySix} & Gaps10421050 & PROVED & \checkmark & verified \\
\texttt{singular\_series\_finite\_pos\_evenPair\_oneThousandFortyTwo} & Gaps10421050 & PROVED & \checkmark & verified \\
\texttt{singular\_series\_pos\_evenPair\_oneThousandFifty} & Gaps10421050 & PROVED & \checkmark & verified \\
\texttt{singular\_series\_pos\_evenPair\_oneThousandFortyEight} & Gaps10421050 & PROVED & \checkmark & verified \\
\texttt{singular\_series\_pos\_evenPair\_oneThousandFortyFour} & Gaps10421050 & PROVED & \checkmark & verified \\
\texttt{singular\_series\_pos\_evenPair\_oneThousandFortySix} & Gaps10421050 & PROVED & \checkmark & verified \\
\texttt{singular\_series\_pos\_evenPair\_oneThousandFortyTwo} & Gaps10421050 & PROVED & \checkmark & verified \\
\texttt{evenPair\_card\_oneThousandFiftyEight} & Gaps10521060 & PROVED & \checkmark & verified \\
\texttt{evenPair\_card\_oneThousandFiftyFour} & Gaps10521060 & PROVED & \checkmark & verified \\
\texttt{evenPair\_card\_oneThousandFiftySix} & Gaps10521060 & PROVED & \checkmark & verified \\
\texttt{evenPair\_card\_oneThousandFiftyTwo} & Gaps10521060 & PROVED & \checkmark & verified \\
\texttt{evenPair\_card\_oneThousandSixty} & Gaps10521060 & PROVED & \checkmark & verified \\
\texttt{isAdmissible\_evenPair\_oneThousandFiftyEight} & Gaps10521060 & PROVED & \checkmark & verified \\
\texttt{isAdmissible\_evenPair\_oneThousandFiftyFour} & Gaps10521060 & PROVED & \checkmark & verified \\
\texttt{isAdmissible\_evenPair\_oneThousandFiftySix} & Gaps10521060 & PROVED & \checkmark & verified \\
\texttt{isAdmissible\_evenPair\_oneThousandFiftyTwo} & Gaps10521060 & PROVED & \checkmark & verified \\
\texttt{isAdmissible\_evenPair\_oneThousandSixty} & Gaps10521060 & PROVED & \checkmark & verified \\
\texttt{localFactor\_oneThousandFiftyTwo\_two} & Gaps10521060 & PROVED & \checkmark & verified \\
\texttt{localFactor\_oneThousandSixty\_two} & Gaps10521060 & PROVED & \checkmark & verified \\
\texttt{nu\_p\_oneThousandFiftyEight} & Gaps10521060 & PROVED & \checkmark & verified \\
\texttt{nu\_p\_oneThousandFiftyFour} & Gaps10521060 & PROVED & \checkmark & verified \\
\texttt{nu\_p\_oneThousandFiftySix} & Gaps10521060 & PROVED & \checkmark & verified \\
\texttt{nu\_p\_oneThousandFiftyTwo} & Gaps10521060 & PROVED & \checkmark & verified \\
\texttt{nu\_p\_oneThousandFiftyTwo\_two} & Gaps10521060 & PROVED & \checkmark & verified \\
\texttt{nu\_p\_oneThousandSixty} & Gaps10521060 & PROVED & \checkmark & verified \\
\texttt{nu\_p\_oneThousandSixty\_two} & Gaps10521060 & PROVED & \checkmark & verified \\
\texttt{singular\_series\_finite\_pos\_evenPair\_oneThousandFiftyEight} & Gaps10521060 & PROVED & \checkmark & verified \\
\texttt{singular\_series\_finite\_pos\_evenPair\_oneThousandFiftyFour} & Gaps10521060 & PROVED & \checkmark & verified \\
\texttt{singular\_series\_finite\_pos\_evenPair\_oneThousandFiftySix} & Gaps10521060 & PROVED & \checkmark & verified \\
\texttt{singular\_series\_finite\_pos\_evenPair\_oneThousandFiftyTwo} & Gaps10521060 & PROVED & \checkmark & verified \\
\texttt{singular\_series\_finite\_pos\_evenPair\_oneThousandSixty} & Gaps10521060 & PROVED & \checkmark & verified \\
\texttt{singular\_series\_pos\_evenPair\_oneThousandFiftyEight} & Gaps10521060 & PROVED & \checkmark & verified \\
\texttt{singular\_series\_pos\_evenPair\_oneThousandFiftyFour} & Gaps10521060 & PROVED & \checkmark & verified \\
\texttt{singular\_series\_pos\_evenPair\_oneThousandFiftySix} & Gaps10521060 & PROVED & \checkmark & verified \\
\texttt{singular\_series\_pos\_evenPair\_oneThousandFiftyTwo} & Gaps10521060 & PROVED & \checkmark & verified \\
\texttt{singular\_series\_pos\_evenPair\_oneThousandSixty} & Gaps10521060 & PROVED & \checkmark & verified \\
\texttt{evenPair\_card\_oneThousandSeventy} & Gaps10621070 & PROVED & \checkmark & verified \\
\texttt{evenPair\_card\_oneThousandSixtyEight} & Gaps10621070 & PROVED & \checkmark & verified \\
\texttt{evenPair\_card\_oneThousandSixtyFour} & Gaps10621070 & PROVED & \checkmark & verified \\
\texttt{evenPair\_card\_oneThousandSixtySix} & Gaps10621070 & PROVED & \checkmark & verified \\
\texttt{evenPair\_card\_oneThousandSixtyTwo} & Gaps10621070 & PROVED & \checkmark & verified \\
\texttt{isAdmissible\_evenPair\_oneThousandSeventy} & Gaps10621070 & PROVED & \checkmark & verified \\
\texttt{isAdmissible\_evenPair\_oneThousandSixtyEight} & Gaps10621070 & PROVED & \checkmark & verified \\
\texttt{isAdmissible\_evenPair\_oneThousandSixtyFour} & Gaps10621070 & PROVED & \checkmark & verified \\
\texttt{isAdmissible\_evenPair\_oneThousandSixtySix} & Gaps10621070 & PROVED & \checkmark & verified \\
\texttt{isAdmissible\_evenPair\_oneThousandSixtyTwo} & Gaps10621070 & PROVED & \checkmark & verified \\
\texttt{localFactor\_oneThousandSeventy\_two} & Gaps10621070 & PROVED & \checkmark & verified \\
\texttt{localFactor\_oneThousandSixtyTwo\_two} & Gaps10621070 & PROVED & \checkmark & verified \\
\texttt{nu\_p\_oneThousandSeventy} & Gaps10621070 & PROVED & \checkmark & verified \\
\texttt{nu\_p\_oneThousandSeventy\_two} & Gaps10621070 & PROVED & \checkmark & verified \\
\texttt{nu\_p\_oneThousandSixtyEight} & Gaps10621070 & PROVED & \checkmark & verified \\
\texttt{nu\_p\_oneThousandSixtyFour} & Gaps10621070 & PROVED & \checkmark & verified \\
\texttt{nu\_p\_oneThousandSixtySix} & Gaps10621070 & PROVED & \checkmark & verified \\
\texttt{nu\_p\_oneThousandSixtyTwo} & Gaps10621070 & PROVED & \checkmark & verified \\
\texttt{nu\_p\_oneThousandSixtyTwo\_two} & Gaps10621070 & PROVED & \checkmark & verified \\
\texttt{singular\_series\_finite\_pos\_evenPair\_oneThousandSeventy} & Gaps10621070 & PROVED & \checkmark & verified \\
\texttt{singular\_series\_finite\_pos\_evenPair\_oneThousandSixtyEight} & Gaps10621070 & PROVED & \checkmark & verified \\
\texttt{singular\_series\_finite\_pos\_evenPair\_oneThousandSixtyFour} & Gaps10621070 & PROVED & \checkmark & verified \\
\texttt{singular\_series\_finite\_pos\_evenPair\_oneThousandSixtySix} & Gaps10621070 & PROVED & \checkmark & verified \\
\texttt{singular\_series\_finite\_pos\_evenPair\_oneThousandSixtyTwo} & Gaps10621070 & PROVED & \checkmark & verified \\
\texttt{singular\_series\_pos\_evenPair\_oneThousandSeventy} & Gaps10621070 & PROVED & \checkmark & verified \\
\texttt{singular\_series\_pos\_evenPair\_oneThousandSixtyEight} & Gaps10621070 & PROVED & \checkmark & verified \\
\texttt{singular\_series\_pos\_evenPair\_oneThousandSixtyFour} & Gaps10621070 & PROVED & \checkmark & verified \\
\texttt{singular\_series\_pos\_evenPair\_oneThousandSixtySix} & Gaps10621070 & PROVED & \checkmark & verified \\
\texttt{singular\_series\_pos\_evenPair\_oneThousandSixtyTwo} & Gaps10621070 & PROVED & \checkmark & verified \\
\texttt{evenPair\_card\_oneThousandEighty} & Gaps10721080 & PROVED & \checkmark & verified \\
\texttt{evenPair\_card\_oneThousandSeventyEight} & Gaps10721080 & PROVED & \checkmark & verified \\
\texttt{evenPair\_card\_oneThousandSeventyFour} & Gaps10721080 & PROVED & \checkmark & verified \\
\texttt{evenPair\_card\_oneThousandSeventySix} & Gaps10721080 & PROVED & \checkmark & verified \\
\texttt{evenPair\_card\_oneThousandSeventyTwo} & Gaps10721080 & PROVED & \checkmark & verified \\
\texttt{isAdmissible\_evenPair\_oneThousandEighty} & Gaps10721080 & PROVED & \checkmark & verified \\
\texttt{isAdmissible\_evenPair\_oneThousandSeventyEight} & Gaps10721080 & PROVED & \checkmark & verified \\
\texttt{isAdmissible\_evenPair\_oneThousandSeventyFour} & Gaps10721080 & PROVED & \checkmark & verified \\
\texttt{isAdmissible\_evenPair\_oneThousandSeventySix} & Gaps10721080 & PROVED & \checkmark & verified \\
\texttt{isAdmissible\_evenPair\_oneThousandSeventyTwo} & Gaps10721080 & PROVED & \checkmark & verified \\
\texttt{localFactor\_oneThousandEighty\_two} & Gaps10721080 & PROVED & \checkmark & verified \\
\texttt{localFactor\_oneThousandSeventyTwo\_two} & Gaps10721080 & PROVED & \checkmark & verified \\
\texttt{nu\_p\_oneThousandEighty} & Gaps10721080 & PROVED & \checkmark & verified \\
\texttt{nu\_p\_oneThousandEighty\_two} & Gaps10721080 & PROVED & \checkmark & verified \\
\texttt{nu\_p\_oneThousandSeventyEight} & Gaps10721080 & PROVED & \checkmark & verified \\
\texttt{nu\_p\_oneThousandSeventyFour} & Gaps10721080 & PROVED & \checkmark & verified \\
\texttt{nu\_p\_oneThousandSeventySix} & Gaps10721080 & PROVED & \checkmark & verified \\
\texttt{nu\_p\_oneThousandSeventyTwo} & Gaps10721080 & PROVED & \checkmark & verified \\
\texttt{nu\_p\_oneThousandSeventyTwo\_two} & Gaps10721080 & PROVED & \checkmark & verified \\
\texttt{singular\_series\_finite\_pos\_evenPair\_oneThousandEighty} & Gaps10721080 & PROVED & \checkmark & verified \\
\texttt{singular\_series\_finite\_pos\_evenPair\_oneThousandSeventyEight} & Gaps10721080 & PROVED & \checkmark & verified \\
\texttt{singular\_series\_finite\_pos\_evenPair\_oneThousandSeventyFour} & Gaps10721080 & PROVED & \checkmark & verified \\
\texttt{singular\_series\_finite\_pos\_evenPair\_oneThousandSeventySix} & Gaps10721080 & PROVED & \checkmark & verified \\
\texttt{singular\_series\_finite\_pos\_evenPair\_oneThousandSeventyTwo} & Gaps10721080 & PROVED & \checkmark & verified \\
\texttt{singular\_series\_pos\_evenPair\_oneThousandEighty} & Gaps10721080 & PROVED & \checkmark & verified \\
\texttt{singular\_series\_pos\_evenPair\_oneThousandSeventyEight} & Gaps10721080 & PROVED & \checkmark & verified \\
\texttt{singular\_series\_pos\_evenPair\_oneThousandSeventyFour} & Gaps10721080 & PROVED & \checkmark & verified \\
\texttt{singular\_series\_pos\_evenPair\_oneThousandSeventySix} & Gaps10721080 & PROVED & \checkmark & verified \\
\texttt{singular\_series\_pos\_evenPair\_oneThousandSeventyTwo} & Gaps10721080 & PROVED & \checkmark & verified \\
\texttt{evenPair\_card\_oneThousandEightyEight} & Gaps10821090 & PROVED & \checkmark & verified \\
\texttt{evenPair\_card\_oneThousandEightyFour} & Gaps10821090 & PROVED & \checkmark & verified \\
\texttt{evenPair\_card\_oneThousandEightySix} & Gaps10821090 & PROVED & \checkmark & verified \\
\texttt{evenPair\_card\_oneThousandEightyTwo} & Gaps10821090 & PROVED & \checkmark & verified \\
\texttt{evenPair\_card\_oneThousandNinety} & Gaps10821090 & PROVED & \checkmark & verified \\
\texttt{isAdmissible\_evenPair\_oneThousandEightyEight} & Gaps10821090 & PROVED & \checkmark & verified \\
\texttt{isAdmissible\_evenPair\_oneThousandEightyFour} & Gaps10821090 & PROVED & \checkmark & verified \\
\texttt{isAdmissible\_evenPair\_oneThousandEightySix} & Gaps10821090 & PROVED & \checkmark & verified \\
\texttt{isAdmissible\_evenPair\_oneThousandEightyTwo} & Gaps10821090 & PROVED & \checkmark & verified \\
\texttt{isAdmissible\_evenPair\_oneThousandNinety} & Gaps10821090 & PROVED & \checkmark & verified \\
\texttt{localFactor\_oneThousandEightyTwo\_two} & Gaps10821090 & PROVED & \checkmark & verified \\
\texttt{localFactor\_oneThousandNinety\_two} & Gaps10821090 & PROVED & \checkmark & verified \\
\texttt{nu\_p\_oneThousandEightyEight} & Gaps10821090 & PROVED & \checkmark & verified \\
\texttt{nu\_p\_oneThousandEightyFour} & Gaps10821090 & PROVED & \checkmark & verified \\
\texttt{nu\_p\_oneThousandEightySix} & Gaps10821090 & PROVED & \checkmark & verified \\
\texttt{nu\_p\_oneThousandEightyTwo} & Gaps10821090 & PROVED & \checkmark & verified \\
\texttt{nu\_p\_oneThousandEightyTwo\_two} & Gaps10821090 & PROVED & \checkmark & verified \\
\texttt{nu\_p\_oneThousandNinety} & Gaps10821090 & PROVED & \checkmark & verified \\
\texttt{nu\_p\_oneThousandNinety\_two} & Gaps10821090 & PROVED & \checkmark & verified \\
\texttt{singular\_series\_finite\_pos\_evenPair\_oneThousandEightyEight} & Gaps10821090 & PROVED & \checkmark & verified \\
\texttt{singular\_series\_finite\_pos\_evenPair\_oneThousandEightyFour} & Gaps10821090 & PROVED & \checkmark & verified \\
\texttt{singular\_series\_finite\_pos\_evenPair\_oneThousandEightySix} & Gaps10821090 & PROVED & \checkmark & verified \\
\texttt{singular\_series\_finite\_pos\_evenPair\_oneThousandEightyTwo} & Gaps10821090 & PROVED & \checkmark & verified \\
\texttt{singular\_series\_finite\_pos\_evenPair\_oneThousandNinety} & Gaps10821090 & PROVED & \checkmark & verified \\
\texttt{singular\_series\_pos\_evenPair\_oneThousandEightyEight} & Gaps10821090 & PROVED & \checkmark & verified \\
\texttt{singular\_series\_pos\_evenPair\_oneThousandEightyFour} & Gaps10821090 & PROVED & \checkmark & verified \\
\texttt{singular\_series\_pos\_evenPair\_oneThousandEightySix} & Gaps10821090 & PROVED & \checkmark & verified \\
\texttt{singular\_series\_pos\_evenPair\_oneThousandEightyTwo} & Gaps10821090 & PROVED & \checkmark & verified \\
\texttt{singular\_series\_pos\_evenPair\_oneThousandNinety} & Gaps10821090 & PROVED & \checkmark & verified \\
\texttt{evenPair\_card\_oneThousandNinetyEight} & Gaps10921100 & PROVED & \checkmark & verified \\
\texttt{evenPair\_card\_oneThousandNinetyFour} & Gaps10921100 & PROVED & \checkmark & verified \\
\texttt{evenPair\_card\_oneThousandNinetySix} & Gaps10921100 & PROVED & \checkmark & verified \\
\texttt{evenPair\_card\_oneThousandNinetyTwo} & Gaps10921100 & PROVED & \checkmark & verified \\
\texttt{evenPair\_card\_oneThousandOneHundred} & Gaps10921100 & PROVED & \checkmark & verified \\
\texttt{isAdmissible\_evenPair\_oneThousandNinetyEight} & Gaps10921100 & PROVED & \checkmark & verified \\
\texttt{isAdmissible\_evenPair\_oneThousandNinetyFour} & Gaps10921100 & PROVED & \checkmark & verified \\
\texttt{isAdmissible\_evenPair\_oneThousandNinetySix} & Gaps10921100 & PROVED & \checkmark & verified \\
\texttt{isAdmissible\_evenPair\_oneThousandNinetyTwo} & Gaps10921100 & PROVED & \checkmark & verified \\
\texttt{isAdmissible\_evenPair\_oneThousandOneHundred} & Gaps10921100 & PROVED & \checkmark & verified \\
\texttt{localFactor\_oneThousandNinetyTwo\_two} & Gaps10921100 & PROVED & \checkmark & verified \\
\texttt{localFactor\_oneThousandOneHundred\_two} & Gaps10921100 & PROVED & \checkmark & verified \\
\texttt{nu\_p\_oneThousandNinetyEight} & Gaps10921100 & PROVED & \checkmark & verified \\
\texttt{nu\_p\_oneThousandNinetyFour} & Gaps10921100 & PROVED & \checkmark & verified \\
\texttt{nu\_p\_oneThousandNinetySix} & Gaps10921100 & PROVED & \checkmark & verified \\
\texttt{nu\_p\_oneThousandNinetyTwo} & Gaps10921100 & PROVED & \checkmark & verified \\
\texttt{nu\_p\_oneThousandNinetyTwo\_two} & Gaps10921100 & PROVED & \checkmark & verified \\
\texttt{nu\_p\_oneThousandOneHundred} & Gaps10921100 & PROVED & \checkmark & verified \\
\texttt{nu\_p\_oneThousandOneHundred\_two} & Gaps10921100 & PROVED & \checkmark & verified \\
\texttt{singular\_series\_finite\_pos\_evenPair\_oneThousandNinetyEight} & Gaps10921100 & PROVED & \checkmark & verified \\
\texttt{singular\_series\_finite\_pos\_evenPair\_oneThousandNinetyFour} & Gaps10921100 & PROVED & \checkmark & verified \\
\texttt{singular\_series\_finite\_pos\_evenPair\_oneThousandNinetySix} & Gaps10921100 & PROVED & \checkmark & verified \\
\texttt{singular\_series\_finite\_pos\_evenPair\_oneThousandNinetyTwo} & Gaps10921100 & PROVED & \checkmark & verified \\
\texttt{singular\_series\_finite\_pos\_evenPair\_oneThousandOneHundred} & Gaps10921100 & PROVED & \checkmark & verified \\
\texttt{singular\_series\_pos\_evenPair\_oneThousandNinetyEight} & Gaps10921100 & PROVED & \checkmark & verified \\
\texttt{singular\_series\_pos\_evenPair\_oneThousandNinetyFour} & Gaps10921100 & PROVED & \checkmark & verified \\
\texttt{singular\_series\_pos\_evenPair\_oneThousandNinetySix} & Gaps10921100 & PROVED & \checkmark & verified \\
\texttt{singular\_series\_pos\_evenPair\_oneThousandNinetyTwo} & Gaps10921100 & PROVED & \checkmark & verified \\
\texttt{singular\_series\_pos\_evenPair\_oneThousandOneHundred} & Gaps10921100 & PROVED & \checkmark & verified \\
\texttt{evenPair\_card\_oneThousandOneHundredEight} & Gaps11021110 & PROVED & \checkmark & verified \\
\texttt{evenPair\_card\_oneThousandOneHundredFour} & Gaps11021110 & PROVED & \checkmark & verified \\
\texttt{evenPair\_card\_oneThousandOneHundredSix} & Gaps11021110 & PROVED & \checkmark & verified \\
\texttt{evenPair\_card\_oneThousandOneHundredTen} & Gaps11021110 & PROVED & \checkmark & verified \\
\texttt{evenPair\_card\_oneThousandOneHundredTwo} & Gaps11021110 & PROVED & \checkmark & verified \\
\texttt{isAdmissible\_evenPair\_oneThousandOneHundredEight} & Gaps11021110 & PROVED & \checkmark & verified \\
\texttt{isAdmissible\_evenPair\_oneThousandOneHundredFour} & Gaps11021110 & PROVED & \checkmark & verified \\
\texttt{isAdmissible\_evenPair\_oneThousandOneHundredSix} & Gaps11021110 & PROVED & \checkmark & verified \\
\texttt{isAdmissible\_evenPair\_oneThousandOneHundredTen} & Gaps11021110 & PROVED & \checkmark & verified \\
\texttt{isAdmissible\_evenPair\_oneThousandOneHundredTwo} & Gaps11021110 & PROVED & \checkmark & verified \\
\texttt{localFactor\_oneThousandOneHundredTen\_two} & Gaps11021110 & PROVED & \checkmark & verified \\
\texttt{localFactor\_oneThousandOneHundredTwo\_two} & Gaps11021110 & PROVED & \checkmark & verified \\
\texttt{nu\_p\_oneThousandOneHundredEight} & Gaps11021110 & PROVED & \checkmark & verified \\
\texttt{nu\_p\_oneThousandOneHundredFour} & Gaps11021110 & PROVED & \checkmark & verified \\
\texttt{nu\_p\_oneThousandOneHundredSix} & Gaps11021110 & PROVED & \checkmark & verified \\
\texttt{nu\_p\_oneThousandOneHundredTen} & Gaps11021110 & PROVED & \checkmark & verified \\
\texttt{nu\_p\_oneThousandOneHundredTen\_two} & Gaps11021110 & PROVED & \checkmark & verified \\
\texttt{nu\_p\_oneThousandOneHundredTwo} & Gaps11021110 & PROVED & \checkmark & verified \\
\texttt{nu\_p\_oneThousandOneHundredTwo\_two} & Gaps11021110 & PROVED & \checkmark & verified \\
\texttt{singular\_series\_finite\_pos\_evenPair\_oneThousandOneHundredEight} & Gaps11021110 & PROVED & \checkmark & verified \\
\texttt{singular\_series\_finite\_pos\_evenPair\_oneThousandOneHundredFour} & Gaps11021110 & PROVED & \checkmark & verified \\
\texttt{singular\_series\_finite\_pos\_evenPair\_oneThousandOneHundredSix} & Gaps11021110 & PROVED & \checkmark & verified \\
\texttt{singular\_series\_finite\_pos\_evenPair\_oneThousandOneHundredTen} & Gaps11021110 & PROVED & \checkmark & verified \\
\texttt{singular\_series\_finite\_pos\_evenPair\_oneThousandOneHundredTwo} & Gaps11021110 & PROVED & \checkmark & verified \\
\texttt{singular\_series\_pos\_evenPair\_oneThousandOneHundredEight} & Gaps11021110 & PROVED & \checkmark & verified \\
\texttt{singular\_series\_pos\_evenPair\_oneThousandOneHundredFour} & Gaps11021110 & PROVED & \checkmark & verified \\
\texttt{singular\_series\_pos\_evenPair\_oneThousandOneHundredSix} & Gaps11021110 & PROVED & \checkmark & verified \\
\texttt{singular\_series\_pos\_evenPair\_oneThousandOneHundredTen} & Gaps11021110 & PROVED & \checkmark & verified \\
\texttt{singular\_series\_pos\_evenPair\_oneThousandOneHundredTwo} & Gaps11021110 & PROVED & \checkmark & verified \\
\texttt{evenPair\_card\_oneThousandOneHundredEighteen} & Gaps11121120 & PROVED & \checkmark & verified \\
\texttt{evenPair\_card\_oneThousandOneHundredFourteen} & Gaps11121120 & PROVED & \checkmark & verified \\
\texttt{evenPair\_card\_oneThousandOneHundredSixteen} & Gaps11121120 & PROVED & \checkmark & verified \\
\texttt{evenPair\_card\_oneThousandOneHundredTwelve} & Gaps11121120 & PROVED & \checkmark & verified \\
\texttt{evenPair\_card\_oneThousandOneHundredTwenty} & Gaps11121120 & PROVED & \checkmark & verified \\
\texttt{isAdmissible\_evenPair\_oneThousandOneHundredEighteen} & Gaps11121120 & PROVED & \checkmark & verified \\
\texttt{isAdmissible\_evenPair\_oneThousandOneHundredFourteen} & Gaps11121120 & PROVED & \checkmark & verified \\
\texttt{isAdmissible\_evenPair\_oneThousandOneHundredSixteen} & Gaps11121120 & PROVED & \checkmark & verified \\
\texttt{isAdmissible\_evenPair\_oneThousandOneHundredTwelve} & Gaps11121120 & PROVED & \checkmark & verified \\
\texttt{isAdmissible\_evenPair\_oneThousandOneHundredTwenty} & Gaps11121120 & PROVED & \checkmark & verified \\
\texttt{localFactor\_oneThousandOneHundredTwelve\_two} & Gaps11121120 & PROVED & \checkmark & verified \\
\texttt{localFactor\_oneThousandOneHundredTwenty\_two} & Gaps11121120 & PROVED & \checkmark & verified \\
\texttt{nu\_p\_oneThousandOneHundredEighteen} & Gaps11121120 & PROVED & \checkmark & verified \\
\texttt{nu\_p\_oneThousandOneHundredFourteen} & Gaps11121120 & PROVED & \checkmark & verified \\
\texttt{nu\_p\_oneThousandOneHundredSixteen} & Gaps11121120 & PROVED & \checkmark & verified \\
\texttt{nu\_p\_oneThousandOneHundredTwelve} & Gaps11121120 & PROVED & \checkmark & verified \\
\texttt{nu\_p\_oneThousandOneHundredTwelve\_two} & Gaps11121120 & PROVED & \checkmark & verified \\
\texttt{nu\_p\_oneThousandOneHundredTwenty} & Gaps11121120 & PROVED & \checkmark & verified \\
\texttt{nu\_p\_oneThousandOneHundredTwenty\_two} & Gaps11121120 & PROVED & \checkmark & verified \\
\texttt{singular\_series\_finite\_pos\_evenPair\_oneThousandOneHundredEighteen} & Gaps11121120 & PROVED & \checkmark & verified \\
\texttt{singular\_series\_finite\_pos\_evenPair\_oneThousandOneHundredFourteen} & Gaps11121120 & PROVED & \checkmark & verified \\
\texttt{singular\_series\_finite\_pos\_evenPair\_oneThousandOneHundredSixteen} & Gaps11121120 & PROVED & \checkmark & verified \\
\texttt{singular\_series\_finite\_pos\_evenPair\_oneThousandOneHundredTwelve} & Gaps11121120 & PROVED & \checkmark & verified \\
\texttt{singular\_series\_finite\_pos\_evenPair\_oneThousandOneHundredTwenty} & Gaps11121120 & PROVED & \checkmark & verified \\
\texttt{singular\_series\_pos\_evenPair\_oneThousandOneHundredEighteen} & Gaps11121120 & PROVED & \checkmark & verified \\
\texttt{singular\_series\_pos\_evenPair\_oneThousandOneHundredFourteen} & Gaps11121120 & PROVED & \checkmark & verified \\
\texttt{singular\_series\_pos\_evenPair\_oneThousandOneHundredSixteen} & Gaps11121120 & PROVED & \checkmark & verified \\
\texttt{singular\_series\_pos\_evenPair\_oneThousandOneHundredTwelve} & Gaps11121120 & PROVED & \checkmark & verified \\
\texttt{singular\_series\_pos\_evenPair\_oneThousandOneHundredTwenty} & Gaps11121120 & PROVED & \checkmark & verified \\
\texttt{evenPair\_card\_oneHundredEighteen} & Gaps112120 & PROVED & \checkmark & verified \\
\texttt{evenPair\_card\_oneHundredFourteen} & Gaps112120 & PROVED & \checkmark & verified \\
\texttt{evenPair\_card\_oneHundredSixteen} & Gaps112120 & PROVED & \checkmark & verified \\
\texttt{evenPair\_card\_oneHundredTwelve} & Gaps112120 & PROVED & \checkmark & verified \\
\texttt{evenPair\_card\_oneHundredTwenty} & Gaps112120 & PROVED & \checkmark & verified \\
\texttt{isAdmissible\_evenPair\_oneHundredEighteen} & Gaps112120 & PROVED & \checkmark & verified \\
\texttt{isAdmissible\_evenPair\_oneHundredFourteen} & Gaps112120 & PROVED & \checkmark & verified \\
\texttt{isAdmissible\_evenPair\_oneHundredSixteen} & Gaps112120 & PROVED & \checkmark & verified \\
\texttt{isAdmissible\_evenPair\_oneHundredTwelve} & Gaps112120 & PROVED & \checkmark & verified \\
\texttt{isAdmissible\_evenPair\_oneHundredTwenty} & Gaps112120 & PROVED & \checkmark & verified \\
\texttt{localFactor\_oneHundredTwelve\_two} & Gaps112120 & PROVED & \checkmark & verified \\
\texttt{localFactor\_oneHundredTwenty\_two} & Gaps112120 & PROVED & \checkmark & verified \\
\texttt{nu\_p\_oneHundredEighteen} & Gaps112120 & PROVED & \checkmark & verified \\
\texttt{nu\_p\_oneHundredFourteen} & Gaps112120 & PROVED & \checkmark & verified \\
\texttt{nu\_p\_oneHundredSixteen} & Gaps112120 & PROVED & \checkmark & verified \\
\texttt{nu\_p\_oneHundredTwelve} & Gaps112120 & PROVED & \checkmark & verified \\
\texttt{nu\_p\_oneHundredTwelve\_two} & Gaps112120 & PROVED & \checkmark & verified \\
\texttt{nu\_p\_oneHundredTwenty} & Gaps112120 & PROVED & \checkmark & verified \\
\texttt{nu\_p\_oneHundredTwenty\_two} & Gaps112120 & PROVED & \checkmark & verified \\
\texttt{singular\_series\_finite\_pos\_evenPair\_oneHundredEighteen} & Gaps112120 & PROVED & \checkmark & verified \\
\texttt{singular\_series\_finite\_pos\_evenPair\_oneHundredFourteen} & Gaps112120 & PROVED & \checkmark & verified \\
\texttt{singular\_series\_finite\_pos\_evenPair\_oneHundredSixteen} & Gaps112120 & PROVED & \checkmark & verified \\
\texttt{singular\_series\_finite\_pos\_evenPair\_oneHundredTwelve} & Gaps112120 & PROVED & \checkmark & verified \\
\texttt{singular\_series\_finite\_pos\_evenPair\_oneHundredTwenty} & Gaps112120 & PROVED & \checkmark & verified \\
\texttt{singular\_series\_pos\_evenPair\_oneHundredEighteen} & Gaps112120 & PROVED & \checkmark & verified \\
\texttt{singular\_series\_pos\_evenPair\_oneHundredFourteen} & Gaps112120 & PROVED & \checkmark & verified \\
\texttt{singular\_series\_pos\_evenPair\_oneHundredSixteen} & Gaps112120 & PROVED & \checkmark & verified \\
\texttt{singular\_series\_pos\_evenPair\_oneHundredTwelve} & Gaps112120 & PROVED & \checkmark & verified \\
\texttt{singular\_series\_pos\_evenPair\_oneHundredTwenty} & Gaps112120 & PROVED & \checkmark & verified \\
\texttt{evenPair\_card\_oneThousandOneHundredThirty} & Gaps11221130 & PROVED & \checkmark & verified \\
\texttt{evenPair\_card\_oneThousandOneHundredTwentyEight} & Gaps11221130 & PROVED & \checkmark & verified \\
\texttt{evenPair\_card\_oneThousandOneHundredTwentyFour} & Gaps11221130 & PROVED & \checkmark & verified \\
\texttt{evenPair\_card\_oneThousandOneHundredTwentySix} & Gaps11221130 & PROVED & \checkmark & verified \\
\texttt{evenPair\_card\_oneThousandOneHundredTwentyTwo} & Gaps11221130 & PROVED & \checkmark & verified \\
\texttt{isAdmissible\_evenPair\_oneThousandOneHundredThirty} & Gaps11221130 & PROVED & \checkmark & verified \\
\texttt{isAdmissible\_evenPair\_oneThousandOneHundredTwentyEight} & Gaps11221130 & PROVED & \checkmark & verified \\
\texttt{isAdmissible\_evenPair\_oneThousandOneHundredTwentyFour} & Gaps11221130 & PROVED & \checkmark & verified \\
\texttt{isAdmissible\_evenPair\_oneThousandOneHundredTwentySix} & Gaps11221130 & PROVED & \checkmark & verified \\
\texttt{isAdmissible\_evenPair\_oneThousandOneHundredTwentyTwo} & Gaps11221130 & PROVED & \checkmark & verified \\
\texttt{localFactor\_oneThousandOneHundredThirty\_two} & Gaps11221130 & PROVED & \checkmark & verified \\
\texttt{localFactor\_oneThousandOneHundredTwentyTwo\_two} & Gaps11221130 & PROVED & \checkmark & verified \\
\texttt{nu\_p\_oneThousandOneHundredThirty} & Gaps11221130 & PROVED & \checkmark & verified \\
\texttt{nu\_p\_oneThousandOneHundredThirty\_two} & Gaps11221130 & PROVED & \checkmark & verified \\
\texttt{nu\_p\_oneThousandOneHundredTwentyEight} & Gaps11221130 & PROVED & \checkmark & verified \\
\texttt{nu\_p\_oneThousandOneHundredTwentyFour} & Gaps11221130 & PROVED & \checkmark & verified \\
\texttt{nu\_p\_oneThousandOneHundredTwentySix} & Gaps11221130 & PROVED & \checkmark & verified \\
\texttt{nu\_p\_oneThousandOneHundredTwentyTwo} & Gaps11221130 & PROVED & \checkmark & verified \\
\texttt{nu\_p\_oneThousandOneHundredTwentyTwo\_two} & Gaps11221130 & PROVED & \checkmark & verified \\
\texttt{singular\_series\_finite\_pos\_evenPair\_oneThousandOneHundredThirty} & Gaps11221130 & PROVED & \checkmark & verified \\
\texttt{singular\_series\_finite\_pos\_evenPair\_oneThousandOneHundredTwentyEight} & Gaps11221130 & PROVED & \checkmark & verified \\
\texttt{singular\_series\_finite\_pos\_evenPair\_oneThousandOneHundredTwentyFour} & Gaps11221130 & PROVED & \checkmark & verified \\
\texttt{singular\_series\_finite\_pos\_evenPair\_oneThousandOneHundredTwentySix} & Gaps11221130 & PROVED & \checkmark & verified \\
\texttt{singular\_series\_finite\_pos\_evenPair\_oneThousandOneHundredTwentyTwo} & Gaps11221130 & PROVED & \checkmark & verified \\
\texttt{singular\_series\_pos\_evenPair\_oneThousandOneHundredThirty} & Gaps11221130 & PROVED & \checkmark & verified \\
\texttt{singular\_series\_pos\_evenPair\_oneThousandOneHundredTwentyEight} & Gaps11221130 & PROVED & \checkmark & verified \\
\texttt{singular\_series\_pos\_evenPair\_oneThousandOneHundredTwentyFour} & Gaps11221130 & PROVED & \checkmark & verified \\
\texttt{singular\_series\_pos\_evenPair\_oneThousandOneHundredTwentySix} & Gaps11221130 & PROVED & \checkmark & verified \\
\texttt{singular\_series\_pos\_evenPair\_oneThousandOneHundredTwentyTwo} & Gaps11221130 & PROVED & \checkmark & verified \\
\texttt{evenPair\_card\_oneThousandOneHundredForty} & Gaps11321140 & PROVED & \checkmark & verified \\
\texttt{evenPair\_card\_oneThousandOneHundredThirtyEight} & Gaps11321140 & PROVED & \checkmark & verified \\
\texttt{evenPair\_card\_oneThousandOneHundredThirtyFour} & Gaps11321140 & PROVED & \checkmark & verified \\
\texttt{evenPair\_card\_oneThousandOneHundredThirtySix} & Gaps11321140 & PROVED & \checkmark & verified \\
\texttt{evenPair\_card\_oneThousandOneHundredThirtyTwo} & Gaps11321140 & PROVED & \checkmark & verified \\
\texttt{isAdmissible\_evenPair\_oneThousandOneHundredForty} & Gaps11321140 & PROVED & \checkmark & verified \\
\texttt{isAdmissible\_evenPair\_oneThousandOneHundredThirtyEight} & Gaps11321140 & PROVED & \checkmark & verified \\
\texttt{isAdmissible\_evenPair\_oneThousandOneHundredThirtyFour} & Gaps11321140 & PROVED & \checkmark & verified \\
\texttt{isAdmissible\_evenPair\_oneThousandOneHundredThirtySix} & Gaps11321140 & PROVED & \checkmark & verified \\
\texttt{isAdmissible\_evenPair\_oneThousandOneHundredThirtyTwo} & Gaps11321140 & PROVED & \checkmark & verified \\
\texttt{localFactor\_oneThousandOneHundredForty\_two} & Gaps11321140 & PROVED & \checkmark & verified \\
\texttt{localFactor\_oneThousandOneHundredThirtyTwo\_two} & Gaps11321140 & PROVED & \checkmark & verified \\
\texttt{nu\_p\_oneThousandOneHundredForty} & Gaps11321140 & PROVED & \checkmark & verified \\
\texttt{nu\_p\_oneThousandOneHundredForty\_two} & Gaps11321140 & PROVED & \checkmark & verified \\
\texttt{nu\_p\_oneThousandOneHundredThirtyEight} & Gaps11321140 & PROVED & \checkmark & verified \\
\texttt{nu\_p\_oneThousandOneHundredThirtyFour} & Gaps11321140 & PROVED & \checkmark & verified \\
\texttt{nu\_p\_oneThousandOneHundredThirtySix} & Gaps11321140 & PROVED & \checkmark & verified \\
\texttt{nu\_p\_oneThousandOneHundredThirtyTwo} & Gaps11321140 & PROVED & \checkmark & verified \\
\texttt{nu\_p\_oneThousandOneHundredThirtyTwo\_two} & Gaps11321140 & PROVED & \checkmark & verified \\
\texttt{singular\_series\_finite\_pos\_evenPair\_oneThousandOneHundredForty} & Gaps11321140 & PROVED & \checkmark & verified \\
\texttt{singular\_series\_finite\_pos\_evenPair\_oneThousandOneHundredThirtyEight} & Gaps11321140 & PROVED & \checkmark & verified \\
\texttt{singular\_series\_finite\_pos\_evenPair\_oneThousandOneHundredThirtyFour} & Gaps11321140 & PROVED & \checkmark & verified \\
\texttt{singular\_series\_finite\_pos\_evenPair\_oneThousandOneHundredThirtySix} & Gaps11321140 & PROVED & \checkmark & verified \\
\texttt{singular\_series\_finite\_pos\_evenPair\_oneThousandOneHundredThirtyTwo} & Gaps11321140 & PROVED & \checkmark & verified \\
\texttt{singular\_series\_pos\_evenPair\_oneThousandOneHundredForty} & Gaps11321140 & PROVED & \checkmark & verified \\
\texttt{singular\_series\_pos\_evenPair\_oneThousandOneHundredThirtyEight} & Gaps11321140 & PROVED & \checkmark & verified \\
\texttt{singular\_series\_pos\_evenPair\_oneThousandOneHundredThirtyFour} & Gaps11321140 & PROVED & \checkmark & verified \\
\texttt{singular\_series\_pos\_evenPair\_oneThousandOneHundredThirtySix} & Gaps11321140 & PROVED & \checkmark & verified \\
\texttt{singular\_series\_pos\_evenPair\_oneThousandOneHundredThirtyTwo} & Gaps11321140 & PROVED & \checkmark & verified \\
\texttt{evenPair\_card\_oneThousandOneHundredFifty} & Gaps11421150 & PROVED & \checkmark & verified \\
\texttt{evenPair\_card\_oneThousandOneHundredFortyEight} & Gaps11421150 & PROVED & \checkmark & verified \\
\texttt{evenPair\_card\_oneThousandOneHundredFortyFour} & Gaps11421150 & PROVED & \checkmark & verified \\
\texttt{evenPair\_card\_oneThousandOneHundredFortySix} & Gaps11421150 & PROVED & \checkmark & verified \\
\texttt{evenPair\_card\_oneThousandOneHundredFortyTwo} & Gaps11421150 & PROVED & \checkmark & verified \\
\texttt{isAdmissible\_evenPair\_oneThousandOneHundredFifty} & Gaps11421150 & PROVED & \checkmark & verified \\
\texttt{isAdmissible\_evenPair\_oneThousandOneHundredFortyEight} & Gaps11421150 & PROVED & \checkmark & verified \\
\texttt{isAdmissible\_evenPair\_oneThousandOneHundredFortyFour} & Gaps11421150 & PROVED & \checkmark & verified \\
\texttt{isAdmissible\_evenPair\_oneThousandOneHundredFortySix} & Gaps11421150 & PROVED & \checkmark & verified \\
\texttt{isAdmissible\_evenPair\_oneThousandOneHundredFortyTwo} & Gaps11421150 & PROVED & \checkmark & verified \\
\texttt{localFactor\_oneThousandOneHundredFifty\_two} & Gaps11421150 & PROVED & \checkmark & verified \\
\texttt{localFactor\_oneThousandOneHundredFortyTwo\_two} & Gaps11421150 & PROVED & \checkmark & verified \\
\texttt{nu\_p\_oneThousandOneHundredFifty} & Gaps11421150 & PROVED & \checkmark & verified \\
\texttt{nu\_p\_oneThousandOneHundredFifty\_two} & Gaps11421150 & PROVED & \checkmark & verified \\
\texttt{nu\_p\_oneThousandOneHundredFortyEight} & Gaps11421150 & PROVED & \checkmark & verified \\
\texttt{nu\_p\_oneThousandOneHundredFortyFour} & Gaps11421150 & PROVED & \checkmark & verified \\
\texttt{nu\_p\_oneThousandOneHundredFortySix} & Gaps11421150 & PROVED & \checkmark & verified \\
\texttt{nu\_p\_oneThousandOneHundredFortyTwo} & Gaps11421150 & PROVED & \checkmark & verified \\
\texttt{nu\_p\_oneThousandOneHundredFortyTwo\_two} & Gaps11421150 & PROVED & \checkmark & verified \\
\texttt{singular\_series\_finite\_pos\_evenPair\_oneThousandOneHundredFifty} & Gaps11421150 & PROVED & \checkmark & verified \\
\texttt{singular\_series\_finite\_pos\_evenPair\_oneThousandOneHundredFortyEight} & Gaps11421150 & PROVED & \checkmark & verified \\
\texttt{singular\_series\_finite\_pos\_evenPair\_oneThousandOneHundredFortyFour} & Gaps11421150 & PROVED & \checkmark & verified \\
\texttt{singular\_series\_finite\_pos\_evenPair\_oneThousandOneHundredFortySix} & Gaps11421150 & PROVED & \checkmark & verified \\
\texttt{singular\_series\_finite\_pos\_evenPair\_oneThousandOneHundredFortyTwo} & Gaps11421150 & PROVED & \checkmark & verified \\
\texttt{singular\_series\_pos\_evenPair\_oneThousandOneHundredFifty} & Gaps11421150 & PROVED & \checkmark & verified \\
\texttt{singular\_series\_pos\_evenPair\_oneThousandOneHundredFortyEight} & Gaps11421150 & PROVED & \checkmark & verified \\
\texttt{singular\_series\_pos\_evenPair\_oneThousandOneHundredFortyFour} & Gaps11421150 & PROVED & \checkmark & verified \\
\texttt{singular\_series\_pos\_evenPair\_oneThousandOneHundredFortySix} & Gaps11421150 & PROVED & \checkmark & verified \\
\texttt{singular\_series\_pos\_evenPair\_oneThousandOneHundredFortyTwo} & Gaps11421150 & PROVED & \checkmark & verified \\
\texttt{evenPair\_card\_oneThousandOneHundredFiftyEight} & Gaps11521160 & PROVED & \checkmark & verified \\
\texttt{evenPair\_card\_oneThousandOneHundredFiftyFour} & Gaps11521160 & PROVED & \checkmark & verified \\
\texttt{evenPair\_card\_oneThousandOneHundredFiftySix} & Gaps11521160 & PROVED & \checkmark & verified \\
\texttt{evenPair\_card\_oneThousandOneHundredFiftyTwo} & Gaps11521160 & PROVED & \checkmark & verified \\
\texttt{evenPair\_card\_oneThousandOneHundredSixty} & Gaps11521160 & PROVED & \checkmark & verified \\
\texttt{isAdmissible\_evenPair\_oneThousandOneHundredFiftyEight} & Gaps11521160 & PROVED & \checkmark & verified \\
\texttt{isAdmissible\_evenPair\_oneThousandOneHundredFiftyFour} & Gaps11521160 & PROVED & \checkmark & verified \\
\texttt{isAdmissible\_evenPair\_oneThousandOneHundredFiftySix} & Gaps11521160 & PROVED & \checkmark & verified \\
\texttt{isAdmissible\_evenPair\_oneThousandOneHundredFiftyTwo} & Gaps11521160 & PROVED & \checkmark & verified \\
\texttt{isAdmissible\_evenPair\_oneThousandOneHundredSixty} & Gaps11521160 & PROVED & \checkmark & verified \\
\texttt{localFactor\_oneThousandOneHundredFiftyTwo\_two} & Gaps11521160 & PROVED & \checkmark & verified \\
\texttt{localFactor\_oneThousandOneHundredSixty\_two} & Gaps11521160 & PROVED & \checkmark & verified \\
\texttt{nu\_p\_oneThousandOneHundredFiftyEight} & Gaps11521160 & PROVED & \checkmark & verified \\
\texttt{nu\_p\_oneThousandOneHundredFiftyFour} & Gaps11521160 & PROVED & \checkmark & verified \\
\texttt{nu\_p\_oneThousandOneHundredFiftySix} & Gaps11521160 & PROVED & \checkmark & verified \\
\texttt{nu\_p\_oneThousandOneHundredFiftyTwo} & Gaps11521160 & PROVED & \checkmark & verified \\
\texttt{nu\_p\_oneThousandOneHundredFiftyTwo\_two} & Gaps11521160 & PROVED & \checkmark & verified \\
\texttt{nu\_p\_oneThousandOneHundredSixty} & Gaps11521160 & PROVED & \checkmark & verified \\
\texttt{nu\_p\_oneThousandOneHundredSixty\_two} & Gaps11521160 & PROVED & \checkmark & verified \\
\texttt{singular\_series\_finite\_pos\_evenPair\_oneThousandOneHundredFiftyEight} & Gaps11521160 & PROVED & \checkmark & verified \\
\texttt{singular\_series\_finite\_pos\_evenPair\_oneThousandOneHundredFiftyFour} & Gaps11521160 & PROVED & \checkmark & verified \\
\texttt{singular\_series\_finite\_pos\_evenPair\_oneThousandOneHundredFiftySix} & Gaps11521160 & PROVED & \checkmark & verified \\
\texttt{singular\_series\_finite\_pos\_evenPair\_oneThousandOneHundredFiftyTwo} & Gaps11521160 & PROVED & \checkmark & verified \\
\texttt{singular\_series\_finite\_pos\_evenPair\_oneThousandOneHundredSixty} & Gaps11521160 & PROVED & \checkmark & verified \\
\texttt{singular\_series\_pos\_evenPair\_oneThousandOneHundredFiftyEight} & Gaps11521160 & PROVED & \checkmark & verified \\
\texttt{singular\_series\_pos\_evenPair\_oneThousandOneHundredFiftyFour} & Gaps11521160 & PROVED & \checkmark & verified \\
\texttt{singular\_series\_pos\_evenPair\_oneThousandOneHundredFiftySix} & Gaps11521160 & PROVED & \checkmark & verified \\
\texttt{singular\_series\_pos\_evenPair\_oneThousandOneHundredFiftyTwo} & Gaps11521160 & PROVED & \checkmark & verified \\
\texttt{singular\_series\_pos\_evenPair\_oneThousandOneHundredSixty} & Gaps11521160 & PROVED & \checkmark & verified \\
\texttt{evenPair\_card\_oneThousandOneHundredSeventy} & Gaps11621170 & PROVED & \checkmark & verified \\
\texttt{evenPair\_card\_oneThousandOneHundredSixtyEight} & Gaps11621170 & PROVED & \checkmark & verified \\
\texttt{evenPair\_card\_oneThousandOneHundredSixtyFour} & Gaps11621170 & PROVED & \checkmark & verified \\
\texttt{evenPair\_card\_oneThousandOneHundredSixtySix} & Gaps11621170 & PROVED & \checkmark & verified \\
\texttt{evenPair\_card\_oneThousandOneHundredSixtyTwo} & Gaps11621170 & PROVED & \checkmark & verified \\
\texttt{isAdmissible\_evenPair\_oneThousandOneHundredSeventy} & Gaps11621170 & PROVED & \checkmark & verified \\
\texttt{isAdmissible\_evenPair\_oneThousandOneHundredSixtyEight} & Gaps11621170 & PROVED & \checkmark & verified \\
\texttt{isAdmissible\_evenPair\_oneThousandOneHundredSixtyFour} & Gaps11621170 & PROVED & \checkmark & verified \\
\texttt{isAdmissible\_evenPair\_oneThousandOneHundredSixtySix} & Gaps11621170 & PROVED & \checkmark & verified \\
\texttt{isAdmissible\_evenPair\_oneThousandOneHundredSixtyTwo} & Gaps11621170 & PROVED & \checkmark & verified \\
\texttt{localFactor\_oneThousandOneHundredSeventy\_two} & Gaps11621170 & PROVED & \checkmark & verified \\
\texttt{localFactor\_oneThousandOneHundredSixtyTwo\_two} & Gaps11621170 & PROVED & \checkmark & verified \\
\texttt{nu\_p\_oneThousandOneHundredSeventy} & Gaps11621170 & PROVED & \checkmark & verified \\
\texttt{nu\_p\_oneThousandOneHundredSeventy\_two} & Gaps11621170 & PROVED & \checkmark & verified \\
\texttt{nu\_p\_oneThousandOneHundredSixtyEight} & Gaps11621170 & PROVED & \checkmark & verified \\
\texttt{nu\_p\_oneThousandOneHundredSixtyFour} & Gaps11621170 & PROVED & \checkmark & verified \\
\texttt{nu\_p\_oneThousandOneHundredSixtySix} & Gaps11621170 & PROVED & \checkmark & verified \\
\texttt{nu\_p\_oneThousandOneHundredSixtyTwo} & Gaps11621170 & PROVED & \checkmark & verified \\
\texttt{nu\_p\_oneThousandOneHundredSixtyTwo\_two} & Gaps11621170 & PROVED & \checkmark & verified \\
\texttt{singular\_series\_finite\_pos\_evenPair\_oneThousandOneHundredSeventy} & Gaps11621170 & PROVED & \checkmark & verified \\
\texttt{singular\_series\_finite\_pos\_evenPair\_oneThousandOneHundredSixtyEight} & Gaps11621170 & PROVED & \checkmark & verified \\
\texttt{singular\_series\_finite\_pos\_evenPair\_oneThousandOneHundredSixtyFour} & Gaps11621170 & PROVED & \checkmark & verified \\
\texttt{singular\_series\_finite\_pos\_evenPair\_oneThousandOneHundredSixtySix} & Gaps11621170 & PROVED & \checkmark & verified \\
\texttt{singular\_series\_finite\_pos\_evenPair\_oneThousandOneHundredSixtyTwo} & Gaps11621170 & PROVED & \checkmark & verified \\
\texttt{singular\_series\_pos\_evenPair\_oneThousandOneHundredSeventy} & Gaps11621170 & PROVED & \checkmark & verified \\
\texttt{singular\_series\_pos\_evenPair\_oneThousandOneHundredSixtyEight} & Gaps11621170 & PROVED & \checkmark & verified \\
\texttt{singular\_series\_pos\_evenPair\_oneThousandOneHundredSixtyFour} & Gaps11621170 & PROVED & \checkmark & verified \\
\texttt{singular\_series\_pos\_evenPair\_oneThousandOneHundredSixtySix} & Gaps11621170 & PROVED & \checkmark & verified \\
\texttt{singular\_series\_pos\_evenPair\_oneThousandOneHundredSixtyTwo} & Gaps11621170 & PROVED & \checkmark & verified \\
\texttt{evenPair\_card\_oneThousandOneHundredEighty} & Gaps11721180 & PROVED & \checkmark & verified \\
\texttt{evenPair\_card\_oneThousandOneHundredSeventyEight} & Gaps11721180 & PROVED & \checkmark & verified \\
\texttt{evenPair\_card\_oneThousandOneHundredSeventyFour} & Gaps11721180 & PROVED & \checkmark & verified \\
\texttt{evenPair\_card\_oneThousandOneHundredSeventySix} & Gaps11721180 & PROVED & \checkmark & verified \\
\texttt{evenPair\_card\_oneThousandOneHundredSeventyTwo} & Gaps11721180 & PROVED & \checkmark & verified \\
\texttt{isAdmissible\_evenPair\_oneThousandOneHundredEighty} & Gaps11721180 & PROVED & \checkmark & verified \\
\texttt{isAdmissible\_evenPair\_oneThousandOneHundredSeventyEight} & Gaps11721180 & PROVED & \checkmark & verified \\
\texttt{isAdmissible\_evenPair\_oneThousandOneHundredSeventyFour} & Gaps11721180 & PROVED & \checkmark & verified \\
\texttt{isAdmissible\_evenPair\_oneThousandOneHundredSeventySix} & Gaps11721180 & PROVED & \checkmark & verified \\
\texttt{isAdmissible\_evenPair\_oneThousandOneHundredSeventyTwo} & Gaps11721180 & PROVED & \checkmark & verified \\
\texttt{localFactor\_oneThousandOneHundredEighty\_two} & Gaps11721180 & PROVED & \checkmark & verified \\
\texttt{localFactor\_oneThousandOneHundredSeventyTwo\_two} & Gaps11721180 & PROVED & \checkmark & verified \\
\texttt{nu\_p\_oneThousandOneHundredEighty} & Gaps11721180 & PROVED & \checkmark & verified \\
\texttt{nu\_p\_oneThousandOneHundredEighty\_two} & Gaps11721180 & PROVED & \checkmark & verified \\
\texttt{nu\_p\_oneThousandOneHundredSeventyEight} & Gaps11721180 & PROVED & \checkmark & verified \\
\texttt{nu\_p\_oneThousandOneHundredSeventyFour} & Gaps11721180 & PROVED & \checkmark & verified \\
\texttt{nu\_p\_oneThousandOneHundredSeventySix} & Gaps11721180 & PROVED & \checkmark & verified \\
\texttt{nu\_p\_oneThousandOneHundredSeventyTwo} & Gaps11721180 & PROVED & \checkmark & verified \\
\texttt{nu\_p\_oneThousandOneHundredSeventyTwo\_two} & Gaps11721180 & PROVED & \checkmark & verified \\
\texttt{singular\_series\_finite\_pos\_evenPair\_oneThousandOneHundredEighty} & Gaps11721180 & PROVED & \checkmark & verified \\
\texttt{singular\_series\_finite\_pos\_evenPair\_oneThousandOneHundredSeventyEight} & Gaps11721180 & PROVED & \checkmark & verified \\
\texttt{singular\_series\_finite\_pos\_evenPair\_oneThousandOneHundredSeventyFour} & Gaps11721180 & PROVED & \checkmark & verified \\
\texttt{singular\_series\_finite\_pos\_evenPair\_oneThousandOneHundredSeventySix} & Gaps11721180 & PROVED & \checkmark & verified \\
\texttt{singular\_series\_finite\_pos\_evenPair\_oneThousandOneHundredSeventyTwo} & Gaps11721180 & PROVED & \checkmark & verified \\
\texttt{singular\_series\_pos\_evenPair\_oneThousandOneHundredEighty} & Gaps11721180 & PROVED & \checkmark & verified \\
\texttt{singular\_series\_pos\_evenPair\_oneThousandOneHundredSeventyEight} & Gaps11721180 & PROVED & \checkmark & verified \\
\texttt{singular\_series\_pos\_evenPair\_oneThousandOneHundredSeventyFour} & Gaps11721180 & PROVED & \checkmark & verified \\
\texttt{singular\_series\_pos\_evenPair\_oneThousandOneHundredSeventySix} & Gaps11721180 & PROVED & \checkmark & verified \\
\texttt{singular\_series\_pos\_evenPair\_oneThousandOneHundredSeventyTwo} & Gaps11721180 & PROVED & \checkmark & verified \\
\texttt{evenPair\_card\_oneThousandOneHundredEightyEight} & Gaps11821190 & PROVED & \checkmark & verified \\
\texttt{evenPair\_card\_oneThousandOneHundredEightyFour} & Gaps11821190 & PROVED & \checkmark & verified \\
\texttt{evenPair\_card\_oneThousandOneHundredEightySix} & Gaps11821190 & PROVED & \checkmark & verified \\
\texttt{evenPair\_card\_oneThousandOneHundredEightyTwo} & Gaps11821190 & PROVED & \checkmark & verified \\
\texttt{evenPair\_card\_oneThousandOneHundredNinety} & Gaps11821190 & PROVED & \checkmark & verified \\
\texttt{isAdmissible\_evenPair\_oneThousandOneHundredEightyEight} & Gaps11821190 & PROVED & \checkmark & verified \\
\texttt{isAdmissible\_evenPair\_oneThousandOneHundredEightyFour} & Gaps11821190 & PROVED & \checkmark & verified \\
\texttt{isAdmissible\_evenPair\_oneThousandOneHundredEightySix} & Gaps11821190 & PROVED & \checkmark & verified \\
\texttt{isAdmissible\_evenPair\_oneThousandOneHundredEightyTwo} & Gaps11821190 & PROVED & \checkmark & verified \\
\texttt{isAdmissible\_evenPair\_oneThousandOneHundredNinety} & Gaps11821190 & PROVED & \checkmark & verified \\
\texttt{localFactor\_oneThousandOneHundredEightyTwo\_two} & Gaps11821190 & PROVED & \checkmark & verified \\
\texttt{localFactor\_oneThousandOneHundredNinety\_two} & Gaps11821190 & PROVED & \checkmark & verified \\
\texttt{nu\_p\_oneThousandOneHundredEightyEight} & Gaps11821190 & PROVED & \checkmark & verified \\
\texttt{nu\_p\_oneThousandOneHundredEightyFour} & Gaps11821190 & PROVED & \checkmark & verified \\
\texttt{nu\_p\_oneThousandOneHundredEightySix} & Gaps11821190 & PROVED & \checkmark & verified \\
\texttt{nu\_p\_oneThousandOneHundredEightyTwo} & Gaps11821190 & PROVED & \checkmark & verified \\
\texttt{nu\_p\_oneThousandOneHundredEightyTwo\_two} & Gaps11821190 & PROVED & \checkmark & verified \\
\texttt{nu\_p\_oneThousandOneHundredNinety} & Gaps11821190 & PROVED & \checkmark & verified \\
\texttt{nu\_p\_oneThousandOneHundredNinety\_two} & Gaps11821190 & PROVED & \checkmark & verified \\
\texttt{singular\_series\_finite\_pos\_evenPair\_oneThousandOneHundredEightyEight} & Gaps11821190 & PROVED & \checkmark & verified \\
\texttt{singular\_series\_finite\_pos\_evenPair\_oneThousandOneHundredEightyFour} & Gaps11821190 & PROVED & \checkmark & verified \\
\texttt{singular\_series\_finite\_pos\_evenPair\_oneThousandOneHundredEightySix} & Gaps11821190 & PROVED & \checkmark & verified \\
\texttt{singular\_series\_finite\_pos\_evenPair\_oneThousandOneHundredEightyTwo} & Gaps11821190 & PROVED & \checkmark & verified \\
\texttt{singular\_series\_finite\_pos\_evenPair\_oneThousandOneHundredNinety} & Gaps11821190 & PROVED & \checkmark & verified \\
\texttt{singular\_series\_pos\_evenPair\_oneThousandOneHundredEightyEight} & Gaps11821190 & PROVED & \checkmark & verified \\
\texttt{singular\_series\_pos\_evenPair\_oneThousandOneHundredEightyFour} & Gaps11821190 & PROVED & \checkmark & verified \\
\texttt{singular\_series\_pos\_evenPair\_oneThousandOneHundredEightySix} & Gaps11821190 & PROVED & \checkmark & verified \\
\texttt{singular\_series\_pos\_evenPair\_oneThousandOneHundredEightyTwo} & Gaps11821190 & PROVED & \checkmark & verified \\
\texttt{singular\_series\_pos\_evenPair\_oneThousandOneHundredNinety} & Gaps11821190 & PROVED & \checkmark & verified \\
\texttt{evenPair\_card\_oneThousandOneHundredNinetyEight} & Gaps11921200 & PROVED & \checkmark & verified \\
\texttt{evenPair\_card\_oneThousandOneHundredNinetyFour} & Gaps11921200 & PROVED & \checkmark & verified \\
\texttt{evenPair\_card\_oneThousandOneHundredNinetySix} & Gaps11921200 & PROVED & \checkmark & verified \\
\texttt{evenPair\_card\_oneThousandOneHundredNinetyTwo} & Gaps11921200 & PROVED & \checkmark & verified \\
\texttt{evenPair\_card\_oneThousandTwoHundred} & Gaps11921200 & PROVED & \checkmark & verified \\
\texttt{isAdmissible\_evenPair\_oneThousandOneHundredNinetyEight} & Gaps11921200 & PROVED & \checkmark & verified \\
\texttt{isAdmissible\_evenPair\_oneThousandOneHundredNinetyFour} & Gaps11921200 & PROVED & \checkmark & verified \\
\texttt{isAdmissible\_evenPair\_oneThousandOneHundredNinetySix} & Gaps11921200 & PROVED & \checkmark & verified \\
\texttt{isAdmissible\_evenPair\_oneThousandOneHundredNinetyTwo} & Gaps11921200 & PROVED & \checkmark & verified \\
\texttt{isAdmissible\_evenPair\_oneThousandTwoHundred} & Gaps11921200 & PROVED & \checkmark & verified \\
\texttt{localFactor\_oneThousandOneHundredNinetyTwo\_two} & Gaps11921200 & PROVED & \checkmark & verified \\
\texttt{localFactor\_oneThousandTwoHundred\_two} & Gaps11921200 & PROVED & \checkmark & verified \\
\texttt{nu\_p\_oneThousandOneHundredNinetyEight} & Gaps11921200 & PROVED & \checkmark & verified \\
\texttt{nu\_p\_oneThousandOneHundredNinetyFour} & Gaps11921200 & PROVED & \checkmark & verified \\
\texttt{nu\_p\_oneThousandOneHundredNinetySix} & Gaps11921200 & PROVED & \checkmark & verified \\
\texttt{nu\_p\_oneThousandOneHundredNinetyTwo} & Gaps11921200 & PROVED & \checkmark & verified \\
\texttt{nu\_p\_oneThousandOneHundredNinetyTwo\_two} & Gaps11921200 & PROVED & \checkmark & verified \\
\texttt{nu\_p\_oneThousandTwoHundred} & Gaps11921200 & PROVED & \checkmark & verified \\
\texttt{nu\_p\_oneThousandTwoHundred\_two} & Gaps11921200 & PROVED & \checkmark & verified \\
\texttt{singular\_series\_finite\_pos\_evenPair\_oneThousandOneHundredNinetyEight} & Gaps11921200 & PROVED & \checkmark & verified \\
\texttt{singular\_series\_finite\_pos\_evenPair\_oneThousandOneHundredNinetyFour} & Gaps11921200 & PROVED & \checkmark & verified \\
\texttt{singular\_series\_finite\_pos\_evenPair\_oneThousandOneHundredNinetySix} & Gaps11921200 & PROVED & \checkmark & verified \\
\texttt{singular\_series\_finite\_pos\_evenPair\_oneThousandOneHundredNinetyTwo} & Gaps11921200 & PROVED & \checkmark & verified \\
\texttt{singular\_series\_finite\_pos\_evenPair\_oneThousandTwoHundred} & Gaps11921200 & PROVED & \checkmark & verified \\
\texttt{singular\_series\_pos\_evenPair\_oneThousandOneHundredNinetyEight} & Gaps11921200 & PROVED & \checkmark & verified \\
\texttt{singular\_series\_pos\_evenPair\_oneThousandOneHundredNinetyFour} & Gaps11921200 & PROVED & \checkmark & verified \\
\texttt{singular\_series\_pos\_evenPair\_oneThousandOneHundredNinetySix} & Gaps11921200 & PROVED & \checkmark & verified \\
\texttt{singular\_series\_pos\_evenPair\_oneThousandOneHundredNinetyTwo} & Gaps11921200 & PROVED & \checkmark & verified \\
\texttt{singular\_series\_pos\_evenPair\_oneThousandTwoHundred} & Gaps11921200 & PROVED & \checkmark & verified \\
\texttt{evenPair\_card\_oneThousandTwoHundredEight} & Gaps12021210 & PROVED & \checkmark & verified \\
\texttt{evenPair\_card\_oneThousandTwoHundredFour} & Gaps12021210 & PROVED & \checkmark & verified \\
\texttt{evenPair\_card\_oneThousandTwoHundredSix} & Gaps12021210 & PROVED & \checkmark & verified \\
\texttt{evenPair\_card\_oneThousandTwoHundredTen} & Gaps12021210 & PROVED & \checkmark & verified \\
\texttt{evenPair\_card\_oneThousandTwoHundredTwo} & Gaps12021210 & PROVED & \checkmark & verified \\
\texttt{isAdmissible\_evenPair\_oneThousandTwoHundredEight} & Gaps12021210 & PROVED & \checkmark & verified \\
\texttt{isAdmissible\_evenPair\_oneThousandTwoHundredFour} & Gaps12021210 & PROVED & \checkmark & verified \\
\texttt{isAdmissible\_evenPair\_oneThousandTwoHundredSix} & Gaps12021210 & PROVED & \checkmark & verified \\
\texttt{isAdmissible\_evenPair\_oneThousandTwoHundredTen} & Gaps12021210 & PROVED & \checkmark & verified \\
\texttt{isAdmissible\_evenPair\_oneThousandTwoHundredTwo} & Gaps12021210 & PROVED & \checkmark & verified \\
\texttt{localFactor\_oneThousandTwoHundredTen\_two} & Gaps12021210 & PROVED & \checkmark & verified \\
\texttt{localFactor\_oneThousandTwoHundredTwo\_two} & Gaps12021210 & PROVED & \checkmark & verified \\
\texttt{nu\_p\_oneThousandTwoHundredEight} & Gaps12021210 & PROVED & \checkmark & verified \\
\texttt{nu\_p\_oneThousandTwoHundredFour} & Gaps12021210 & PROVED & \checkmark & verified \\
\texttt{nu\_p\_oneThousandTwoHundredSix} & Gaps12021210 & PROVED & \checkmark & verified \\
\texttt{nu\_p\_oneThousandTwoHundredTen} & Gaps12021210 & PROVED & \checkmark & verified \\
\texttt{nu\_p\_oneThousandTwoHundredTen\_two} & Gaps12021210 & PROVED & \checkmark & verified \\
\texttt{nu\_p\_oneThousandTwoHundredTwo} & Gaps12021210 & PROVED & \checkmark & verified \\
\texttt{nu\_p\_oneThousandTwoHundredTwo\_two} & Gaps12021210 & PROVED & \checkmark & verified \\
\texttt{singular\_series\_finite\_pos\_evenPair\_oneThousandTwoHundredEight} & Gaps12021210 & PROVED & \checkmark & verified \\
\texttt{singular\_series\_finite\_pos\_evenPair\_oneThousandTwoHundredFour} & Gaps12021210 & PROVED & \checkmark & verified \\
\texttt{singular\_series\_finite\_pos\_evenPair\_oneThousandTwoHundredSix} & Gaps12021210 & PROVED & \checkmark & verified \\
\texttt{singular\_series\_finite\_pos\_evenPair\_oneThousandTwoHundredTen} & Gaps12021210 & PROVED & \checkmark & verified \\
\texttt{singular\_series\_finite\_pos\_evenPair\_oneThousandTwoHundredTwo} & Gaps12021210 & PROVED & \checkmark & verified \\
\texttt{singular\_series\_pos\_evenPair\_oneThousandTwoHundredEight} & Gaps12021210 & PROVED & \checkmark & verified \\
\texttt{singular\_series\_pos\_evenPair\_oneThousandTwoHundredFour} & Gaps12021210 & PROVED & \checkmark & verified \\
\texttt{singular\_series\_pos\_evenPair\_oneThousandTwoHundredSix} & Gaps12021210 & PROVED & \checkmark & verified \\
\texttt{singular\_series\_pos\_evenPair\_oneThousandTwoHundredTen} & Gaps12021210 & PROVED & \checkmark & verified \\
\texttt{singular\_series\_pos\_evenPair\_oneThousandTwoHundredTwo} & Gaps12021210 & PROVED & \checkmark & verified \\
\texttt{evenPair\_card\_oneThousandTwoHundredEighteen} & Gaps12121220 & PROVED & \checkmark & verified \\
\texttt{evenPair\_card\_oneThousandTwoHundredFourteen} & Gaps12121220 & PROVED & \checkmark & verified \\
\texttt{evenPair\_card\_oneThousandTwoHundredSixteen} & Gaps12121220 & PROVED & \checkmark & verified \\
\texttt{evenPair\_card\_oneThousandTwoHundredTwelve} & Gaps12121220 & PROVED & \checkmark & verified \\
\texttt{evenPair\_card\_oneThousandTwoHundredTwenty} & Gaps12121220 & PROVED & \checkmark & verified \\
\texttt{isAdmissible\_evenPair\_oneThousandTwoHundredEighteen} & Gaps12121220 & PROVED & \checkmark & verified \\
\texttt{isAdmissible\_evenPair\_oneThousandTwoHundredFourteen} & Gaps12121220 & PROVED & \checkmark & verified \\
\texttt{isAdmissible\_evenPair\_oneThousandTwoHundredSixteen} & Gaps12121220 & PROVED & \checkmark & verified \\
\texttt{isAdmissible\_evenPair\_oneThousandTwoHundredTwelve} & Gaps12121220 & PROVED & \checkmark & verified \\
\texttt{isAdmissible\_evenPair\_oneThousandTwoHundredTwenty} & Gaps12121220 & PROVED & \checkmark & verified \\
\texttt{localFactor\_oneThousandTwoHundredTwelve\_two} & Gaps12121220 & PROVED & \checkmark & verified \\
\texttt{localFactor\_oneThousandTwoHundredTwenty\_two} & Gaps12121220 & PROVED & \checkmark & verified \\
\texttt{nu\_p\_oneThousandTwoHundredEighteen} & Gaps12121220 & PROVED & \checkmark & verified \\
\texttt{nu\_p\_oneThousandTwoHundredFourteen} & Gaps12121220 & PROVED & \checkmark & verified \\
\texttt{nu\_p\_oneThousandTwoHundredSixteen} & Gaps12121220 & PROVED & \checkmark & verified \\
\texttt{nu\_p\_oneThousandTwoHundredTwelve} & Gaps12121220 & PROVED & \checkmark & verified \\
\texttt{nu\_p\_oneThousandTwoHundredTwelve\_two} & Gaps12121220 & PROVED & \checkmark & verified \\
\texttt{nu\_p\_oneThousandTwoHundredTwenty} & Gaps12121220 & PROVED & \checkmark & verified \\
\texttt{nu\_p\_oneThousandTwoHundredTwenty\_two} & Gaps12121220 & PROVED & \checkmark & verified \\
\texttt{singular\_series\_finite\_pos\_evenPair\_oneThousandTwoHundredEighteen} & Gaps12121220 & PROVED & \checkmark & verified \\
\texttt{singular\_series\_finite\_pos\_evenPair\_oneThousandTwoHundredFourteen} & Gaps12121220 & PROVED & \checkmark & verified \\
\texttt{singular\_series\_finite\_pos\_evenPair\_oneThousandTwoHundredSixteen} & Gaps12121220 & PROVED & \checkmark & verified \\
\texttt{singular\_series\_finite\_pos\_evenPair\_oneThousandTwoHundredTwelve} & Gaps12121220 & PROVED & \checkmark & verified \\
\texttt{singular\_series\_finite\_pos\_evenPair\_oneThousandTwoHundredTwenty} & Gaps12121220 & PROVED & \checkmark & verified \\
\texttt{singular\_series\_pos\_evenPair\_oneThousandTwoHundredEighteen} & Gaps12121220 & PROVED & \checkmark & verified \\
\texttt{singular\_series\_pos\_evenPair\_oneThousandTwoHundredFourteen} & Gaps12121220 & PROVED & \checkmark & verified \\
\texttt{singular\_series\_pos\_evenPair\_oneThousandTwoHundredSixteen} & Gaps12121220 & PROVED & \checkmark & verified \\
\texttt{singular\_series\_pos\_evenPair\_oneThousandTwoHundredTwelve} & Gaps12121220 & PROVED & \checkmark & verified \\
\texttt{singular\_series\_pos\_evenPair\_oneThousandTwoHundredTwenty} & Gaps12121220 & PROVED & \checkmark & verified \\
\texttt{evenPair\_card\_oneHundredThirty} & Gaps122130 & PROVED & \checkmark & verified \\
\texttt{evenPair\_card\_oneHundredTwentyEight} & Gaps122130 & PROVED & \checkmark & verified \\
\texttt{evenPair\_card\_oneHundredTwentyFour} & Gaps122130 & PROVED & \checkmark & verified \\
\texttt{evenPair\_card\_oneHundredTwentySix} & Gaps122130 & PROVED & \checkmark & verified \\
\texttt{evenPair\_card\_oneHundredTwentyTwo} & Gaps122130 & PROVED & \checkmark & verified \\
\texttt{isAdmissible\_evenPair\_oneHundredThirty} & Gaps122130 & PROVED & \checkmark & verified \\
\texttt{isAdmissible\_evenPair\_oneHundredTwentyEight} & Gaps122130 & PROVED & \checkmark & verified \\
\texttt{isAdmissible\_evenPair\_oneHundredTwentyFour} & Gaps122130 & PROVED & \checkmark & verified \\
\texttt{isAdmissible\_evenPair\_oneHundredTwentySix} & Gaps122130 & PROVED & \checkmark & verified \\
\texttt{isAdmissible\_evenPair\_oneHundredTwentyTwo} & Gaps122130 & PROVED & \checkmark & verified \\
\texttt{localFactor\_oneHundredThirty\_two} & Gaps122130 & PROVED & \checkmark & verified \\
\texttt{localFactor\_oneHundredTwentyTwo\_two} & Gaps122130 & PROVED & \checkmark & verified \\
\texttt{nu\_p\_oneHundredThirty} & Gaps122130 & PROVED & \checkmark & verified \\
\texttt{nu\_p\_oneHundredThirty\_two} & Gaps122130 & PROVED & \checkmark & verified \\
\texttt{nu\_p\_oneHundredTwentyEight} & Gaps122130 & PROVED & \checkmark & verified \\
\texttt{nu\_p\_oneHundredTwentyFour} & Gaps122130 & PROVED & \checkmark & verified \\
\texttt{nu\_p\_oneHundredTwentySix} & Gaps122130 & PROVED & \checkmark & verified \\
\texttt{nu\_p\_oneHundredTwentyTwo} & Gaps122130 & PROVED & \checkmark & verified \\
\texttt{nu\_p\_oneHundredTwentyTwo\_two} & Gaps122130 & PROVED & \checkmark & verified \\
\texttt{singular\_series\_finite\_pos\_evenPair\_oneHundredThirty} & Gaps122130 & PROVED & \checkmark & verified \\
\texttt{singular\_series\_finite\_pos\_evenPair\_oneHundredTwentyEight} & Gaps122130 & PROVED & \checkmark & verified \\
\texttt{singular\_series\_finite\_pos\_evenPair\_oneHundredTwentyFour} & Gaps122130 & PROVED & \checkmark & verified \\
\texttt{singular\_series\_finite\_pos\_evenPair\_oneHundredTwentySix} & Gaps122130 & PROVED & \checkmark & verified \\
\texttt{singular\_series\_finite\_pos\_evenPair\_oneHundredTwentyTwo} & Gaps122130 & PROVED & \checkmark & verified \\
\texttt{singular\_series\_pos\_evenPair\_oneHundredThirty} & Gaps122130 & PROVED & \checkmark & verified \\
\texttt{singular\_series\_pos\_evenPair\_oneHundredTwentyEight} & Gaps122130 & PROVED & \checkmark & verified \\
\texttt{singular\_series\_pos\_evenPair\_oneHundredTwentyFour} & Gaps122130 & PROVED & \checkmark & verified \\
\texttt{singular\_series\_pos\_evenPair\_oneHundredTwentySix} & Gaps122130 & PROVED & \checkmark & verified \\
\texttt{singular\_series\_pos\_evenPair\_oneHundredTwentyTwo} & Gaps122130 & PROVED & \checkmark & verified \\
\texttt{evenPair\_card\_oneThousandTwoHundredThirty} & Gaps12221230 & PROVED & \checkmark & verified \\
\texttt{evenPair\_card\_oneThousandTwoHundredTwentyEight} & Gaps12221230 & PROVED & \checkmark & verified \\
\texttt{evenPair\_card\_oneThousandTwoHundredTwentyFour} & Gaps12221230 & PROVED & \checkmark & verified \\
\texttt{evenPair\_card\_oneThousandTwoHundredTwentySix} & Gaps12221230 & PROVED & \checkmark & verified \\
\texttt{evenPair\_card\_oneThousandTwoHundredTwentyTwo} & Gaps12221230 & PROVED & \checkmark & verified \\
\texttt{isAdmissible\_evenPair\_oneThousandTwoHundredThirty} & Gaps12221230 & PROVED & \checkmark & verified \\
\texttt{isAdmissible\_evenPair\_oneThousandTwoHundredTwentyEight} & Gaps12221230 & PROVED & \checkmark & verified \\
\texttt{isAdmissible\_evenPair\_oneThousandTwoHundredTwentyFour} & Gaps12221230 & PROVED & \checkmark & verified \\
\texttt{isAdmissible\_evenPair\_oneThousandTwoHundredTwentySix} & Gaps12221230 & PROVED & \checkmark & verified \\
\texttt{isAdmissible\_evenPair\_oneThousandTwoHundredTwentyTwo} & Gaps12221230 & PROVED & \checkmark & verified \\
\texttt{localFactor\_oneThousandTwoHundredThirty\_two} & Gaps12221230 & PROVED & \checkmark & verified \\
\texttt{localFactor\_oneThousandTwoHundredTwentyTwo\_two} & Gaps12221230 & PROVED & \checkmark & verified \\
\texttt{nu\_p\_oneThousandTwoHundredThirty} & Gaps12221230 & PROVED & \checkmark & verified \\
\texttt{nu\_p\_oneThousandTwoHundredThirty\_two} & Gaps12221230 & PROVED & \checkmark & verified \\
\texttt{nu\_p\_oneThousandTwoHundredTwentyEight} & Gaps12221230 & PROVED & \checkmark & verified \\
\texttt{nu\_p\_oneThousandTwoHundredTwentyFour} & Gaps12221230 & PROVED & \checkmark & verified \\
\texttt{nu\_p\_oneThousandTwoHundredTwentySix} & Gaps12221230 & PROVED & \checkmark & verified \\
\texttt{nu\_p\_oneThousandTwoHundredTwentyTwo} & Gaps12221230 & PROVED & \checkmark & verified \\
\texttt{nu\_p\_oneThousandTwoHundredTwentyTwo\_two} & Gaps12221230 & PROVED & \checkmark & verified \\
\texttt{singular\_series\_finite\_pos\_evenPair\_oneThousandTwoHundredThirty} & Gaps12221230 & PROVED & \checkmark & verified \\
\texttt{singular\_series\_finite\_pos\_evenPair\_oneThousandTwoHundredTwentyEight} & Gaps12221230 & PROVED & \checkmark & verified \\
\texttt{singular\_series\_finite\_pos\_evenPair\_oneThousandTwoHundredTwentyFour} & Gaps12221230 & PROVED & \checkmark & verified \\
\texttt{singular\_series\_finite\_pos\_evenPair\_oneThousandTwoHundredTwentySix} & Gaps12221230 & PROVED & \checkmark & verified \\
\texttt{singular\_series\_finite\_pos\_evenPair\_oneThousandTwoHundredTwentyTwo} & Gaps12221230 & PROVED & \checkmark & verified \\
\texttt{singular\_series\_pos\_evenPair\_oneThousandTwoHundredThirty} & Gaps12221230 & PROVED & \checkmark & verified \\
\texttt{singular\_series\_pos\_evenPair\_oneThousandTwoHundredTwentyEight} & Gaps12221230 & PROVED & \checkmark & verified \\
\texttt{singular\_series\_pos\_evenPair\_oneThousandTwoHundredTwentyFour} & Gaps12221230 & PROVED & \checkmark & verified \\
\texttt{singular\_series\_pos\_evenPair\_oneThousandTwoHundredTwentySix} & Gaps12221230 & PROVED & \checkmark & verified \\
\texttt{singular\_series\_pos\_evenPair\_oneThousandTwoHundredTwentyTwo} & Gaps12221230 & PROVED & \checkmark & verified \\
\texttt{evenPair\_card\_oneThousandTwoHundredForty} & Gaps12321240 & PROVED & \checkmark & verified \\
\texttt{evenPair\_card\_oneThousandTwoHundredThirtyEight} & Gaps12321240 & PROVED & \checkmark & verified \\
\texttt{evenPair\_card\_oneThousandTwoHundredThirtyFour} & Gaps12321240 & PROVED & \checkmark & verified \\
\texttt{evenPair\_card\_oneThousandTwoHundredThirtySix} & Gaps12321240 & PROVED & \checkmark & verified \\
\texttt{evenPair\_card\_oneThousandTwoHundredThirtyTwo} & Gaps12321240 & PROVED & \checkmark & verified \\
\texttt{isAdmissible\_evenPair\_oneThousandTwoHundredForty} & Gaps12321240 & PROVED & \checkmark & verified \\
\texttt{isAdmissible\_evenPair\_oneThousandTwoHundredThirtyEight} & Gaps12321240 & PROVED & \checkmark & verified \\
\texttt{isAdmissible\_evenPair\_oneThousandTwoHundredThirtyFour} & Gaps12321240 & PROVED & \checkmark & verified \\
\texttt{isAdmissible\_evenPair\_oneThousandTwoHundredThirtySix} & Gaps12321240 & PROVED & \checkmark & verified \\
\texttt{isAdmissible\_evenPair\_oneThousandTwoHundredThirtyTwo} & Gaps12321240 & PROVED & \checkmark & verified \\
\texttt{localFactor\_oneThousandTwoHundredForty\_two} & Gaps12321240 & PROVED & \checkmark & verified \\
\texttt{localFactor\_oneThousandTwoHundredThirtyTwo\_two} & Gaps12321240 & PROVED & \checkmark & verified \\
\texttt{nu\_p\_oneThousandTwoHundredForty} & Gaps12321240 & PROVED & \checkmark & verified \\
\texttt{nu\_p\_oneThousandTwoHundredForty\_two} & Gaps12321240 & PROVED & \checkmark & verified \\
\texttt{nu\_p\_oneThousandTwoHundredThirtyEight} & Gaps12321240 & PROVED & \checkmark & verified \\
\texttt{nu\_p\_oneThousandTwoHundredThirtyFour} & Gaps12321240 & PROVED & \checkmark & verified \\
\texttt{nu\_p\_oneThousandTwoHundredThirtySix} & Gaps12321240 & PROVED & \checkmark & verified \\
\texttt{nu\_p\_oneThousandTwoHundredThirtyTwo} & Gaps12321240 & PROVED & \checkmark & verified \\
\texttt{nu\_p\_oneThousandTwoHundredThirtyTwo\_two} & Gaps12321240 & PROVED & \checkmark & verified \\
\texttt{singular\_series\_finite\_pos\_evenPair\_oneThousandTwoHundredForty} & Gaps12321240 & PROVED & \checkmark & verified \\
\texttt{singular\_series\_finite\_pos\_evenPair\_oneThousandTwoHundredThirtyEight} & Gaps12321240 & PROVED & \checkmark & verified \\
\texttt{singular\_series\_finite\_pos\_evenPair\_oneThousandTwoHundredThirtyFour} & Gaps12321240 & PROVED & \checkmark & verified \\
\texttt{singular\_series\_finite\_pos\_evenPair\_oneThousandTwoHundredThirtySix} & Gaps12321240 & PROVED & \checkmark & verified \\
\texttt{singular\_series\_finite\_pos\_evenPair\_oneThousandTwoHundredThirtyTwo} & Gaps12321240 & PROVED & \checkmark & verified \\
\texttt{singular\_series\_pos\_evenPair\_oneThousandTwoHundredForty} & Gaps12321240 & PROVED & \checkmark & verified \\
\texttt{singular\_series\_pos\_evenPair\_oneThousandTwoHundredThirtyEight} & Gaps12321240 & PROVED & \checkmark & verified \\
\texttt{singular\_series\_pos\_evenPair\_oneThousandTwoHundredThirtyFour} & Gaps12321240 & PROVED & \checkmark & verified \\
\texttt{singular\_series\_pos\_evenPair\_oneThousandTwoHundredThirtySix} & Gaps12321240 & PROVED & \checkmark & verified \\
\texttt{singular\_series\_pos\_evenPair\_oneThousandTwoHundredThirtyTwo} & Gaps12321240 & PROVED & \checkmark & verified \\
\texttt{evenPair\_card\_oneThousandTwoHundredFifty} & Gaps12421250 & PROVED & \checkmark & verified \\
\texttt{evenPair\_card\_oneThousandTwoHundredFortyEight} & Gaps12421250 & PROVED & \checkmark & verified \\
\texttt{evenPair\_card\_oneThousandTwoHundredFortyFour} & Gaps12421250 & PROVED & \checkmark & verified \\
\texttt{evenPair\_card\_oneThousandTwoHundredFortySix} & Gaps12421250 & PROVED & \checkmark & verified \\
\texttt{evenPair\_card\_oneThousandTwoHundredFortyTwo} & Gaps12421250 & PROVED & \checkmark & verified \\
\texttt{isAdmissible\_evenPair\_oneThousandTwoHundredFifty} & Gaps12421250 & PROVED & \checkmark & verified \\
\texttt{isAdmissible\_evenPair\_oneThousandTwoHundredFortyEight} & Gaps12421250 & PROVED & \checkmark & verified \\
\texttt{isAdmissible\_evenPair\_oneThousandTwoHundredFortyFour} & Gaps12421250 & PROVED & \checkmark & verified \\
\texttt{isAdmissible\_evenPair\_oneThousandTwoHundredFortySix} & Gaps12421250 & PROVED & \checkmark & verified \\
\texttt{isAdmissible\_evenPair\_oneThousandTwoHundredFortyTwo} & Gaps12421250 & PROVED & \checkmark & verified \\
\texttt{localFactor\_oneThousandTwoHundredFifty\_two} & Gaps12421250 & PROVED & \checkmark & verified \\
\texttt{localFactor\_oneThousandTwoHundredFortyTwo\_two} & Gaps12421250 & PROVED & \checkmark & verified \\
\texttt{nu\_p\_oneThousandTwoHundredFifty} & Gaps12421250 & PROVED & \checkmark & verified \\
\texttt{nu\_p\_oneThousandTwoHundredFifty\_two} & Gaps12421250 & PROVED & \checkmark & verified \\
\texttt{nu\_p\_oneThousandTwoHundredFortyEight} & Gaps12421250 & PROVED & \checkmark & verified \\
\texttt{nu\_p\_oneThousandTwoHundredFortyFour} & Gaps12421250 & PROVED & \checkmark & verified \\
\texttt{nu\_p\_oneThousandTwoHundredFortySix} & Gaps12421250 & PROVED & \checkmark & verified \\
\texttt{nu\_p\_oneThousandTwoHundredFortyTwo} & Gaps12421250 & PROVED & \checkmark & verified \\
\texttt{nu\_p\_oneThousandTwoHundredFortyTwo\_two} & Gaps12421250 & PROVED & \checkmark & verified \\
\texttt{singular\_series\_finite\_pos\_evenPair\_oneThousandTwoHundredFifty} & Gaps12421250 & PROVED & \checkmark & verified \\
\texttt{singular\_series\_finite\_pos\_evenPair\_oneThousandTwoHundredFortyEight} & Gaps12421250 & PROVED & \checkmark & verified \\
\texttt{singular\_series\_finite\_pos\_evenPair\_oneThousandTwoHundredFortyFour} & Gaps12421250 & PROVED & \checkmark & verified \\
\texttt{singular\_series\_finite\_pos\_evenPair\_oneThousandTwoHundredFortySix} & Gaps12421250 & PROVED & \checkmark & verified \\
\texttt{singular\_series\_finite\_pos\_evenPair\_oneThousandTwoHundredFortyTwo} & Gaps12421250 & PROVED & \checkmark & verified \\
\texttt{singular\_series\_pos\_evenPair\_oneThousandTwoHundredFifty} & Gaps12421250 & PROVED & \checkmark & verified \\
\texttt{singular\_series\_pos\_evenPair\_oneThousandTwoHundredFortyEight} & Gaps12421250 & PROVED & \checkmark & verified \\
\texttt{singular\_series\_pos\_evenPair\_oneThousandTwoHundredFortyFour} & Gaps12421250 & PROVED & \checkmark & verified \\
\texttt{singular\_series\_pos\_evenPair\_oneThousandTwoHundredFortySix} & Gaps12421250 & PROVED & \checkmark & verified \\
\texttt{singular\_series\_pos\_evenPair\_oneThousandTwoHundredFortyTwo} & Gaps12421250 & PROVED & \checkmark & verified \\
\texttt{evenPair\_card\_oneThousandTwoHundredFiftyEight} & Gaps12521260 & PROVED & \checkmark & verified \\
\texttt{evenPair\_card\_oneThousandTwoHundredFiftyFour} & Gaps12521260 & PROVED & \checkmark & verified \\
\texttt{evenPair\_card\_oneThousandTwoHundredFiftySix} & Gaps12521260 & PROVED & \checkmark & verified \\
\texttt{evenPair\_card\_oneThousandTwoHundredFiftyTwo} & Gaps12521260 & PROVED & \checkmark & verified \\
\texttt{evenPair\_card\_oneThousandTwoHundredSixty} & Gaps12521260 & PROVED & \checkmark & verified \\
\texttt{isAdmissible\_evenPair\_oneThousandTwoHundredFiftyEight} & Gaps12521260 & PROVED & \checkmark & verified \\
\texttt{isAdmissible\_evenPair\_oneThousandTwoHundredFiftyFour} & Gaps12521260 & PROVED & \checkmark & verified \\
\texttt{isAdmissible\_evenPair\_oneThousandTwoHundredFiftySix} & Gaps12521260 & PROVED & \checkmark & verified \\
\texttt{isAdmissible\_evenPair\_oneThousandTwoHundredFiftyTwo} & Gaps12521260 & PROVED & \checkmark & verified \\
\texttt{isAdmissible\_evenPair\_oneThousandTwoHundredSixty} & Gaps12521260 & PROVED & \checkmark & verified \\
\texttt{localFactor\_oneThousandTwoHundredFiftyTwo\_two} & Gaps12521260 & PROVED & \checkmark & verified \\
\texttt{localFactor\_oneThousandTwoHundredSixty\_two} & Gaps12521260 & PROVED & \checkmark & verified \\
\texttt{nu\_p\_oneThousandTwoHundredFiftyEight} & Gaps12521260 & PROVED & \checkmark & verified \\
\texttt{nu\_p\_oneThousandTwoHundredFiftyFour} & Gaps12521260 & PROVED & \checkmark & verified \\
\texttt{nu\_p\_oneThousandTwoHundredFiftySix} & Gaps12521260 & PROVED & \checkmark & verified \\
\texttt{nu\_p\_oneThousandTwoHundredFiftyTwo} & Gaps12521260 & PROVED & \checkmark & verified \\
\texttt{nu\_p\_oneThousandTwoHundredFiftyTwo\_two} & Gaps12521260 & PROVED & \checkmark & verified \\
\texttt{nu\_p\_oneThousandTwoHundredSixty} & Gaps12521260 & PROVED & \checkmark & verified \\
\texttt{nu\_p\_oneThousandTwoHundredSixty\_two} & Gaps12521260 & PROVED & \checkmark & verified \\
\texttt{singular\_series\_finite\_pos\_evenPair\_oneThousandTwoHundredFiftyEight} & Gaps12521260 & PROVED & \checkmark & verified \\
\texttt{singular\_series\_finite\_pos\_evenPair\_oneThousandTwoHundredFiftyFour} & Gaps12521260 & PROVED & \checkmark & verified \\
\texttt{singular\_series\_finite\_pos\_evenPair\_oneThousandTwoHundredFiftySix} & Gaps12521260 & PROVED & \checkmark & verified \\
\texttt{singular\_series\_finite\_pos\_evenPair\_oneThousandTwoHundredFiftyTwo} & Gaps12521260 & PROVED & \checkmark & verified \\
\texttt{singular\_series\_finite\_pos\_evenPair\_oneThousandTwoHundredSixty} & Gaps12521260 & PROVED & \checkmark & verified \\
\texttt{singular\_series\_pos\_evenPair\_oneThousandTwoHundredFiftyEight} & Gaps12521260 & PROVED & \checkmark & verified \\
\texttt{singular\_series\_pos\_evenPair\_oneThousandTwoHundredFiftyFour} & Gaps12521260 & PROVED & \checkmark & verified \\
\texttt{singular\_series\_pos\_evenPair\_oneThousandTwoHundredFiftySix} & Gaps12521260 & PROVED & \checkmark & verified \\
\texttt{singular\_series\_pos\_evenPair\_oneThousandTwoHundredFiftyTwo} & Gaps12521260 & PROVED & \checkmark & verified \\
\texttt{singular\_series\_pos\_evenPair\_oneThousandTwoHundredSixty} & Gaps12521260 & PROVED & \checkmark & verified \\
\texttt{evenPair\_card\_oneThousandTwoHundredSeventy} & Gaps12621270 & PROVED & \checkmark & verified \\
\texttt{evenPair\_card\_oneThousandTwoHundredSixtyEight} & Gaps12621270 & PROVED & \checkmark & verified \\
\texttt{evenPair\_card\_oneThousandTwoHundredSixtyFour} & Gaps12621270 & PROVED & \checkmark & verified \\
\texttt{evenPair\_card\_oneThousandTwoHundredSixtySix} & Gaps12621270 & PROVED & \checkmark & verified \\
\texttt{evenPair\_card\_oneThousandTwoHundredSixtyTwo} & Gaps12621270 & PROVED & \checkmark & verified \\
\texttt{isAdmissible\_evenPair\_oneThousandTwoHundredSeventy} & Gaps12621270 & PROVED & \checkmark & verified \\
\texttt{isAdmissible\_evenPair\_oneThousandTwoHundredSixtyEight} & Gaps12621270 & PROVED & \checkmark & verified \\
\texttt{isAdmissible\_evenPair\_oneThousandTwoHundredSixtyFour} & Gaps12621270 & PROVED & \checkmark & verified \\
\texttt{isAdmissible\_evenPair\_oneThousandTwoHundredSixtySix} & Gaps12621270 & PROVED & \checkmark & verified \\
\texttt{isAdmissible\_evenPair\_oneThousandTwoHundredSixtyTwo} & Gaps12621270 & PROVED & \checkmark & verified \\
\texttt{localFactor\_oneThousandTwoHundredSeventy\_two} & Gaps12621270 & PROVED & \checkmark & verified \\
\texttt{localFactor\_oneThousandTwoHundredSixtyTwo\_two} & Gaps12621270 & PROVED & \checkmark & verified \\
\texttt{nu\_p\_oneThousandTwoHundredSeventy} & Gaps12621270 & PROVED & \checkmark & verified \\
\texttt{nu\_p\_oneThousandTwoHundredSeventy\_two} & Gaps12621270 & PROVED & \checkmark & verified \\
\texttt{nu\_p\_oneThousandTwoHundredSixtyEight} & Gaps12621270 & PROVED & \checkmark & verified \\
\texttt{nu\_p\_oneThousandTwoHundredSixtyFour} & Gaps12621270 & PROVED & \checkmark & verified \\
\texttt{nu\_p\_oneThousandTwoHundredSixtySix} & Gaps12621270 & PROVED & \checkmark & verified \\
\texttt{nu\_p\_oneThousandTwoHundredSixtyTwo} & Gaps12621270 & PROVED & \checkmark & verified \\
\texttt{nu\_p\_oneThousandTwoHundredSixtyTwo\_two} & Gaps12621270 & PROVED & \checkmark & verified \\
\texttt{singular\_series\_finite\_pos\_evenPair\_oneThousandTwoHundredSeventy} & Gaps12621270 & PROVED & \checkmark & verified \\
\texttt{singular\_series\_finite\_pos\_evenPair\_oneThousandTwoHundredSixtyEight} & Gaps12621270 & PROVED & \checkmark & verified \\
\texttt{singular\_series\_finite\_pos\_evenPair\_oneThousandTwoHundredSixtyFour} & Gaps12621270 & PROVED & \checkmark & verified \\
\texttt{singular\_series\_finite\_pos\_evenPair\_oneThousandTwoHundredSixtySix} & Gaps12621270 & PROVED & \checkmark & verified \\
\texttt{singular\_series\_finite\_pos\_evenPair\_oneThousandTwoHundredSixtyTwo} & Gaps12621270 & PROVED & \checkmark & verified \\
\texttt{singular\_series\_pos\_evenPair\_oneThousandTwoHundredSeventy} & Gaps12621270 & PROVED & \checkmark & verified \\
\texttt{singular\_series\_pos\_evenPair\_oneThousandTwoHundredSixtyEight} & Gaps12621270 & PROVED & \checkmark & verified \\
\texttt{singular\_series\_pos\_evenPair\_oneThousandTwoHundredSixtyFour} & Gaps12621270 & PROVED & \checkmark & verified \\
\texttt{singular\_series\_pos\_evenPair\_oneThousandTwoHundredSixtySix} & Gaps12621270 & PROVED & \checkmark & verified \\
\texttt{singular\_series\_pos\_evenPair\_oneThousandTwoHundredSixtyTwo} & Gaps12621270 & PROVED & \checkmark & verified \\
\texttt{evenPair\_card\_oneThousandTwoHundredEighty} & Gaps12721280 & PROVED & \checkmark & verified \\
\texttt{evenPair\_card\_oneThousandTwoHundredSeventyEight} & Gaps12721280 & PROVED & \checkmark & verified \\
\texttt{evenPair\_card\_oneThousandTwoHundredSeventyFour} & Gaps12721280 & PROVED & \checkmark & verified \\
\texttt{evenPair\_card\_oneThousandTwoHundredSeventySix} & Gaps12721280 & PROVED & \checkmark & verified \\
\texttt{evenPair\_card\_oneThousandTwoHundredSeventyTwo} & Gaps12721280 & PROVED & \checkmark & verified \\
\texttt{isAdmissible\_evenPair\_oneThousandTwoHundredEighty} & Gaps12721280 & PROVED & \checkmark & verified \\
\texttt{isAdmissible\_evenPair\_oneThousandTwoHundredSeventyEight} & Gaps12721280 & PROVED & \checkmark & verified \\
\texttt{isAdmissible\_evenPair\_oneThousandTwoHundredSeventyFour} & Gaps12721280 & PROVED & \checkmark & verified \\
\texttt{isAdmissible\_evenPair\_oneThousandTwoHundredSeventySix} & Gaps12721280 & PROVED & \checkmark & verified \\
\texttt{isAdmissible\_evenPair\_oneThousandTwoHundredSeventyTwo} & Gaps12721280 & PROVED & \checkmark & verified \\
\texttt{localFactor\_oneThousandTwoHundredEighty\_two} & Gaps12721280 & PROVED & \checkmark & verified \\
\texttt{localFactor\_oneThousandTwoHundredSeventyTwo\_two} & Gaps12721280 & PROVED & \checkmark & verified \\
\texttt{nu\_p\_oneThousandTwoHundredEighty} & Gaps12721280 & PROVED & \checkmark & verified \\
\texttt{nu\_p\_oneThousandTwoHundredEighty\_two} & Gaps12721280 & PROVED & \checkmark & verified \\
\texttt{nu\_p\_oneThousandTwoHundredSeventyEight} & Gaps12721280 & PROVED & \checkmark & verified \\
\texttt{nu\_p\_oneThousandTwoHundredSeventyFour} & Gaps12721280 & PROVED & \checkmark & verified \\
\texttt{nu\_p\_oneThousandTwoHundredSeventySix} & Gaps12721280 & PROVED & \checkmark & verified \\
\texttt{nu\_p\_oneThousandTwoHundredSeventyTwo} & Gaps12721280 & PROVED & \checkmark & verified \\
\texttt{nu\_p\_oneThousandTwoHundredSeventyTwo\_two} & Gaps12721280 & PROVED & \checkmark & verified \\
\texttt{singular\_series\_finite\_pos\_evenPair\_oneThousandTwoHundredEighty} & Gaps12721280 & PROVED & \checkmark & verified \\
\texttt{singular\_series\_finite\_pos\_evenPair\_oneThousandTwoHundredSeventyEight} & Gaps12721280 & PROVED & \checkmark & verified \\
\texttt{singular\_series\_finite\_pos\_evenPair\_oneThousandTwoHundredSeventyFour} & Gaps12721280 & PROVED & \checkmark & verified \\
\texttt{singular\_series\_finite\_pos\_evenPair\_oneThousandTwoHundredSeventySix} & Gaps12721280 & PROVED & \checkmark & verified \\
\texttt{singular\_series\_finite\_pos\_evenPair\_oneThousandTwoHundredSeventyTwo} & Gaps12721280 & PROVED & \checkmark & verified \\
\texttt{singular\_series\_pos\_evenPair\_oneThousandTwoHundredEighty} & Gaps12721280 & PROVED & \checkmark & verified \\
\texttt{singular\_series\_pos\_evenPair\_oneThousandTwoHundredSeventyEight} & Gaps12721280 & PROVED & \checkmark & verified \\
\texttt{singular\_series\_pos\_evenPair\_oneThousandTwoHundredSeventyFour} & Gaps12721280 & PROVED & \checkmark & verified \\
\texttt{singular\_series\_pos\_evenPair\_oneThousandTwoHundredSeventySix} & Gaps12721280 & PROVED & \checkmark & verified \\
\texttt{singular\_series\_pos\_evenPair\_oneThousandTwoHundredSeventyTwo} & Gaps12721280 & PROVED & \checkmark & verified \\
\texttt{evenPair\_card\_oneThousandTwoHundredEightyEight} & Gaps12821290 & PROVED & \checkmark & verified \\
\texttt{evenPair\_card\_oneThousandTwoHundredEightyFour} & Gaps12821290 & PROVED & \checkmark & verified \\
\texttt{evenPair\_card\_oneThousandTwoHundredEightySix} & Gaps12821290 & PROVED & \checkmark & verified \\
\texttt{evenPair\_card\_oneThousandTwoHundredEightyTwo} & Gaps12821290 & PROVED & \checkmark & verified \\
\texttt{evenPair\_card\_oneThousandTwoHundredNinety} & Gaps12821290 & PROVED & \checkmark & verified \\
\texttt{isAdmissible\_evenPair\_oneThousandTwoHundredEightyEight} & Gaps12821290 & PROVED & \checkmark & verified \\
\texttt{isAdmissible\_evenPair\_oneThousandTwoHundredEightyFour} & Gaps12821290 & PROVED & \checkmark & verified \\
\texttt{isAdmissible\_evenPair\_oneThousandTwoHundredEightySix} & Gaps12821290 & PROVED & \checkmark & verified \\
\texttt{isAdmissible\_evenPair\_oneThousandTwoHundredEightyTwo} & Gaps12821290 & PROVED & \checkmark & verified \\
\texttt{isAdmissible\_evenPair\_oneThousandTwoHundredNinety} & Gaps12821290 & PROVED & \checkmark & verified \\
\texttt{localFactor\_oneThousandTwoHundredEightyTwo\_two} & Gaps12821290 & PROVED & \checkmark & verified \\
\texttt{localFactor\_oneThousandTwoHundredNinety\_two} & Gaps12821290 & PROVED & \checkmark & verified \\
\texttt{nu\_p\_oneThousandTwoHundredEightyEight} & Gaps12821290 & PROVED & \checkmark & verified \\
\texttt{nu\_p\_oneThousandTwoHundredEightyFour} & Gaps12821290 & PROVED & \checkmark & verified \\
\texttt{nu\_p\_oneThousandTwoHundredEightySix} & Gaps12821290 & PROVED & \checkmark & verified \\
\texttt{nu\_p\_oneThousandTwoHundredEightyTwo} & Gaps12821290 & PROVED & \checkmark & verified \\
\texttt{nu\_p\_oneThousandTwoHundredEightyTwo\_two} & Gaps12821290 & PROVED & \checkmark & verified \\
\texttt{nu\_p\_oneThousandTwoHundredNinety} & Gaps12821290 & PROVED & \checkmark & verified \\
\texttt{nu\_p\_oneThousandTwoHundredNinety\_two} & Gaps12821290 & PROVED & \checkmark & verified \\
\texttt{singular\_series\_finite\_pos\_evenPair\_oneThousandTwoHundredEightyEight} & Gaps12821290 & PROVED & \checkmark & verified \\
\texttt{singular\_series\_finite\_pos\_evenPair\_oneThousandTwoHundredEightyFour} & Gaps12821290 & PROVED & \checkmark & verified \\
\texttt{singular\_series\_finite\_pos\_evenPair\_oneThousandTwoHundredEightySix} & Gaps12821290 & PROVED & \checkmark & verified \\
\texttt{singular\_series\_finite\_pos\_evenPair\_oneThousandTwoHundredEightyTwo} & Gaps12821290 & PROVED & \checkmark & verified \\
\texttt{singular\_series\_finite\_pos\_evenPair\_oneThousandTwoHundredNinety} & Gaps12821290 & PROVED & \checkmark & verified \\
\texttt{singular\_series\_pos\_evenPair\_oneThousandTwoHundredEightyEight} & Gaps12821290 & PROVED & \checkmark & verified \\
\texttt{singular\_series\_pos\_evenPair\_oneThousandTwoHundredEightyFour} & Gaps12821290 & PROVED & \checkmark & verified \\
\texttt{singular\_series\_pos\_evenPair\_oneThousandTwoHundredEightySix} & Gaps12821290 & PROVED & \checkmark & verified \\
\texttt{singular\_series\_pos\_evenPair\_oneThousandTwoHundredEightyTwo} & Gaps12821290 & PROVED & \checkmark & verified \\
\texttt{singular\_series\_pos\_evenPair\_oneThousandTwoHundredNinety} & Gaps12821290 & PROVED & \checkmark & verified \\
\texttt{evenPair\_card\_oneThousandThreeHundred} & Gaps12921300 & PROVED & \checkmark & verified \\
\texttt{evenPair\_card\_oneThousandTwoHundredNinetyEight} & Gaps12921300 & PROVED & \checkmark & verified \\
\texttt{evenPair\_card\_oneThousandTwoHundredNinetyFour} & Gaps12921300 & PROVED & \checkmark & verified \\
\texttt{evenPair\_card\_oneThousandTwoHundredNinetySix} & Gaps12921300 & PROVED & \checkmark & verified \\
\texttt{evenPair\_card\_oneThousandTwoHundredNinetyTwo} & Gaps12921300 & PROVED & \checkmark & verified \\
\texttt{isAdmissible\_evenPair\_oneThousandThreeHundred} & Gaps12921300 & PROVED & \checkmark & verified \\
\texttt{isAdmissible\_evenPair\_oneThousandTwoHundredNinetyEight} & Gaps12921300 & PROVED & \checkmark & verified \\
\texttt{isAdmissible\_evenPair\_oneThousandTwoHundredNinetyFour} & Gaps12921300 & PROVED & \checkmark & verified \\
\texttt{isAdmissible\_evenPair\_oneThousandTwoHundredNinetySix} & Gaps12921300 & PROVED & \checkmark & verified \\
\texttt{isAdmissible\_evenPair\_oneThousandTwoHundredNinetyTwo} & Gaps12921300 & PROVED & \checkmark & verified \\
\texttt{localFactor\_oneThousandThreeHundred\_two} & Gaps12921300 & PROVED & \checkmark & verified \\
\texttt{localFactor\_oneThousandTwoHundredNinetyTwo\_two} & Gaps12921300 & PROVED & \checkmark & verified \\
\texttt{nu\_p\_oneThousandThreeHundred} & Gaps12921300 & PROVED & \checkmark & verified \\
\texttt{nu\_p\_oneThousandThreeHundred\_two} & Gaps12921300 & PROVED & \checkmark & verified \\
\texttt{nu\_p\_oneThousandTwoHundredNinetyEight} & Gaps12921300 & PROVED & \checkmark & verified \\
\texttt{nu\_p\_oneThousandTwoHundredNinetyFour} & Gaps12921300 & PROVED & \checkmark & verified \\
\texttt{nu\_p\_oneThousandTwoHundredNinetySix} & Gaps12921300 & PROVED & \checkmark & verified \\
\texttt{nu\_p\_oneThousandTwoHundredNinetyTwo} & Gaps12921300 & PROVED & \checkmark & verified \\
\texttt{nu\_p\_oneThousandTwoHundredNinetyTwo\_two} & Gaps12921300 & PROVED & \checkmark & verified \\
\texttt{singular\_series\_finite\_pos\_evenPair\_oneThousandThreeHundred} & Gaps12921300 & PROVED & \checkmark & verified \\
\texttt{singular\_series\_finite\_pos\_evenPair\_oneThousandTwoHundredNinetyEight} & Gaps12921300 & PROVED & \checkmark & verified \\
\texttt{singular\_series\_finite\_pos\_evenPair\_oneThousandTwoHundredNinetyFour} & Gaps12921300 & PROVED & \checkmark & verified \\
\texttt{singular\_series\_finite\_pos\_evenPair\_oneThousandTwoHundredNinetySix} & Gaps12921300 & PROVED & \checkmark & verified \\
\texttt{singular\_series\_finite\_pos\_evenPair\_oneThousandTwoHundredNinetyTwo} & Gaps12921300 & PROVED & \checkmark & verified \\
\texttt{singular\_series\_pos\_evenPair\_oneThousandThreeHundred} & Gaps12921300 & PROVED & \checkmark & verified \\
\texttt{singular\_series\_pos\_evenPair\_oneThousandTwoHundredNinetyEight} & Gaps12921300 & PROVED & \checkmark & verified \\
\texttt{singular\_series\_pos\_evenPair\_oneThousandTwoHundredNinetyFour} & Gaps12921300 & PROVED & \checkmark & verified \\
\texttt{singular\_series\_pos\_evenPair\_oneThousandTwoHundredNinetySix} & Gaps12921300 & PROVED & \checkmark & verified \\
\texttt{singular\_series\_pos\_evenPair\_oneThousandTwoHundredNinetyTwo} & Gaps12921300 & PROVED & \checkmark & verified \\
\texttt{evenPair\_card\_oneThousandThreeHundredEight} & Gaps13021310 & PROVED & \checkmark & verified \\
\texttt{evenPair\_card\_oneThousandThreeHundredFour} & Gaps13021310 & PROVED & \checkmark & verified \\
\texttt{evenPair\_card\_oneThousandThreeHundredSix} & Gaps13021310 & PROVED & \checkmark & verified \\
\texttt{evenPair\_card\_oneThousandThreeHundredTen} & Gaps13021310 & PROVED & \checkmark & verified \\
\texttt{evenPair\_card\_oneThousandThreeHundredTwo} & Gaps13021310 & PROVED & \checkmark & verified \\
\texttt{isAdmissible\_evenPair\_oneThousandThreeHundredEight} & Gaps13021310 & PROVED & \checkmark & verified \\
\texttt{isAdmissible\_evenPair\_oneThousandThreeHundredFour} & Gaps13021310 & PROVED & \checkmark & verified \\
\texttt{isAdmissible\_evenPair\_oneThousandThreeHundredSix} & Gaps13021310 & PROVED & \checkmark & verified \\
\texttt{isAdmissible\_evenPair\_oneThousandThreeHundredTen} & Gaps13021310 & PROVED & \checkmark & verified \\
\texttt{isAdmissible\_evenPair\_oneThousandThreeHundredTwo} & Gaps13021310 & PROVED & \checkmark & verified \\
\texttt{localFactor\_oneThousandThreeHundredTen\_two} & Gaps13021310 & PROVED & \checkmark & verified \\
\texttt{localFactor\_oneThousandThreeHundredTwo\_two} & Gaps13021310 & PROVED & \checkmark & verified \\
\texttt{nu\_p\_oneThousandThreeHundredEight} & Gaps13021310 & PROVED & \checkmark & verified \\
\texttt{nu\_p\_oneThousandThreeHundredFour} & Gaps13021310 & PROVED & \checkmark & verified \\
\texttt{nu\_p\_oneThousandThreeHundredSix} & Gaps13021310 & PROVED & \checkmark & verified \\
\texttt{nu\_p\_oneThousandThreeHundredTen} & Gaps13021310 & PROVED & \checkmark & verified \\
\texttt{nu\_p\_oneThousandThreeHundredTen\_two} & Gaps13021310 & PROVED & \checkmark & verified \\
\texttt{nu\_p\_oneThousandThreeHundredTwo} & Gaps13021310 & PROVED & \checkmark & verified \\
\texttt{nu\_p\_oneThousandThreeHundredTwo\_two} & Gaps13021310 & PROVED & \checkmark & verified \\
\texttt{singular\_series\_finite\_pos\_evenPair\_oneThousandThreeHundredEight} & Gaps13021310 & PROVED & \checkmark & verified \\
\texttt{singular\_series\_finite\_pos\_evenPair\_oneThousandThreeHundredFour} & Gaps13021310 & PROVED & \checkmark & verified \\
\texttt{singular\_series\_finite\_pos\_evenPair\_oneThousandThreeHundredSix} & Gaps13021310 & PROVED & \checkmark & verified \\
\texttt{singular\_series\_finite\_pos\_evenPair\_oneThousandThreeHundredTen} & Gaps13021310 & PROVED & \checkmark & verified \\
\texttt{singular\_series\_finite\_pos\_evenPair\_oneThousandThreeHundredTwo} & Gaps13021310 & PROVED & \checkmark & verified \\
\texttt{singular\_series\_pos\_evenPair\_oneThousandThreeHundredEight} & Gaps13021310 & PROVED & \checkmark & verified \\
\texttt{singular\_series\_pos\_evenPair\_oneThousandThreeHundredFour} & Gaps13021310 & PROVED & \checkmark & verified \\
\texttt{singular\_series\_pos\_evenPair\_oneThousandThreeHundredSix} & Gaps13021310 & PROVED & \checkmark & verified \\
\texttt{singular\_series\_pos\_evenPair\_oneThousandThreeHundredTen} & Gaps13021310 & PROVED & \checkmark & verified \\
\texttt{singular\_series\_pos\_evenPair\_oneThousandThreeHundredTwo} & Gaps13021310 & PROVED & \checkmark & verified \\
\texttt{evenPair\_card\_oneThousandThreeHundredEighteen} & Gaps13121320 & PROVED & \checkmark & verified \\
\texttt{evenPair\_card\_oneThousandThreeHundredFourteen} & Gaps13121320 & PROVED & \checkmark & verified \\
\texttt{evenPair\_card\_oneThousandThreeHundredSixteen} & Gaps13121320 & PROVED & \checkmark & verified \\
\texttt{evenPair\_card\_oneThousandThreeHundredTwelve} & Gaps13121320 & PROVED & \checkmark & verified \\
\texttt{evenPair\_card\_oneThousandThreeHundredTwenty} & Gaps13121320 & PROVED & \checkmark & verified \\
\texttt{isAdmissible\_evenPair\_oneThousandThreeHundredEighteen} & Gaps13121320 & PROVED & \checkmark & verified \\
\texttt{isAdmissible\_evenPair\_oneThousandThreeHundredFourteen} & Gaps13121320 & PROVED & \checkmark & verified \\
\texttt{isAdmissible\_evenPair\_oneThousandThreeHundredSixteen} & Gaps13121320 & PROVED & \checkmark & verified \\
\texttt{isAdmissible\_evenPair\_oneThousandThreeHundredTwelve} & Gaps13121320 & PROVED & \checkmark & verified \\
\texttt{isAdmissible\_evenPair\_oneThousandThreeHundredTwenty} & Gaps13121320 & PROVED & \checkmark & verified \\
\texttt{localFactor\_oneThousandThreeHundredTwelve\_two} & Gaps13121320 & PROVED & \checkmark & verified \\
\texttt{localFactor\_oneThousandThreeHundredTwenty\_two} & Gaps13121320 & PROVED & \checkmark & verified \\
\texttt{nu\_p\_oneThousandThreeHundredEighteen} & Gaps13121320 & PROVED & \checkmark & verified \\
\texttt{nu\_p\_oneThousandThreeHundredFourteen} & Gaps13121320 & PROVED & \checkmark & verified \\
\texttt{nu\_p\_oneThousandThreeHundredSixteen} & Gaps13121320 & PROVED & \checkmark & verified \\
\texttt{nu\_p\_oneThousandThreeHundredTwelve} & Gaps13121320 & PROVED & \checkmark & verified \\
\texttt{nu\_p\_oneThousandThreeHundredTwelve\_two} & Gaps13121320 & PROVED & \checkmark & verified \\
\texttt{nu\_p\_oneThousandThreeHundredTwenty} & Gaps13121320 & PROVED & \checkmark & verified \\
\texttt{nu\_p\_oneThousandThreeHundredTwenty\_two} & Gaps13121320 & PROVED & \checkmark & verified \\
\texttt{singular\_series\_finite\_pos\_evenPair\_oneThousandThreeHundredEighteen} & Gaps13121320 & PROVED & \checkmark & verified \\
\texttt{singular\_series\_finite\_pos\_evenPair\_oneThousandThreeHundredFourteen} & Gaps13121320 & PROVED & \checkmark & verified \\
\texttt{singular\_series\_finite\_pos\_evenPair\_oneThousandThreeHundredSixteen} & Gaps13121320 & PROVED & \checkmark & verified \\
\texttt{singular\_series\_finite\_pos\_evenPair\_oneThousandThreeHundredTwelve} & Gaps13121320 & PROVED & \checkmark & verified \\
\texttt{singular\_series\_finite\_pos\_evenPair\_oneThousandThreeHundredTwenty} & Gaps13121320 & PROVED & \checkmark & verified \\
\texttt{singular\_series\_pos\_evenPair\_oneThousandThreeHundredEighteen} & Gaps13121320 & PROVED & \checkmark & verified \\
\texttt{singular\_series\_pos\_evenPair\_oneThousandThreeHundredFourteen} & Gaps13121320 & PROVED & \checkmark & verified \\
\texttt{singular\_series\_pos\_evenPair\_oneThousandThreeHundredSixteen} & Gaps13121320 & PROVED & \checkmark & verified \\
\texttt{singular\_series\_pos\_evenPair\_oneThousandThreeHundredTwelve} & Gaps13121320 & PROVED & \checkmark & verified \\
\texttt{singular\_series\_pos\_evenPair\_oneThousandThreeHundredTwenty} & Gaps13121320 & PROVED & \checkmark & verified \\
\texttt{evenPair\_card\_oneHundredForty} & Gaps132140 & PROVED & \checkmark & verified \\
\texttt{evenPair\_card\_oneHundredThirtyEight} & Gaps132140 & PROVED & \checkmark & verified \\
\texttt{evenPair\_card\_oneHundredThirtyFour} & Gaps132140 & PROVED & \checkmark & verified \\
\texttt{evenPair\_card\_oneHundredThirtySix} & Gaps132140 & PROVED & \checkmark & verified \\
\texttt{evenPair\_card\_oneHundredThirtyTwo} & Gaps132140 & PROVED & \checkmark & verified \\
\texttt{isAdmissible\_evenPair\_oneHundredForty} & Gaps132140 & PROVED & \checkmark & verified \\
\texttt{isAdmissible\_evenPair\_oneHundredThirtyEight} & Gaps132140 & PROVED & \checkmark & verified \\
\texttt{isAdmissible\_evenPair\_oneHundredThirtyFour} & Gaps132140 & PROVED & \checkmark & verified \\
\texttt{isAdmissible\_evenPair\_oneHundredThirtySix} & Gaps132140 & PROVED & \checkmark & verified \\
\texttt{isAdmissible\_evenPair\_oneHundredThirtyTwo} & Gaps132140 & PROVED & \checkmark & verified \\
\texttt{localFactor\_oneHundredForty\_two} & Gaps132140 & PROVED & \checkmark & verified \\
\texttt{localFactor\_oneHundredThirtyTwo\_two} & Gaps132140 & PROVED & \checkmark & verified \\
\texttt{nu\_p\_oneHundredForty} & Gaps132140 & PROVED & \checkmark & verified \\
\texttt{nu\_p\_oneHundredForty\_two} & Gaps132140 & PROVED & \checkmark & verified \\
\texttt{nu\_p\_oneHundredThirtyEight} & Gaps132140 & PROVED & \checkmark & verified \\
\texttt{nu\_p\_oneHundredThirtyFour} & Gaps132140 & PROVED & \checkmark & verified \\
\texttt{nu\_p\_oneHundredThirtySix} & Gaps132140 & PROVED & \checkmark & verified \\
\texttt{nu\_p\_oneHundredThirtyTwo} & Gaps132140 & PROVED & \checkmark & verified \\
\texttt{nu\_p\_oneHundredThirtyTwo\_two} & Gaps132140 & PROVED & \checkmark & verified \\
\texttt{singular\_series\_finite\_pos\_evenPair\_oneHundredForty} & Gaps132140 & PROVED & \checkmark & verified \\
\texttt{singular\_series\_finite\_pos\_evenPair\_oneHundredThirtyEight} & Gaps132140 & PROVED & \checkmark & verified \\
\texttt{singular\_series\_finite\_pos\_evenPair\_oneHundredThirtyFour} & Gaps132140 & PROVED & \checkmark & verified \\
\texttt{singular\_series\_finite\_pos\_evenPair\_oneHundredThirtySix} & Gaps132140 & PROVED & \checkmark & verified \\
\texttt{singular\_series\_finite\_pos\_evenPair\_oneHundredThirtyTwo} & Gaps132140 & PROVED & \checkmark & verified \\
\texttt{singular\_series\_pos\_evenPair\_oneHundredForty} & Gaps132140 & PROVED & \checkmark & verified \\
\texttt{singular\_series\_pos\_evenPair\_oneHundredThirtyEight} & Gaps132140 & PROVED & \checkmark & verified \\
\texttt{singular\_series\_pos\_evenPair\_oneHundredThirtyFour} & Gaps132140 & PROVED & \checkmark & verified \\
\texttt{singular\_series\_pos\_evenPair\_oneHundredThirtySix} & Gaps132140 & PROVED & \checkmark & verified \\
\texttt{singular\_series\_pos\_evenPair\_oneHundredThirtyTwo} & Gaps132140 & PROVED & \checkmark & verified \\
\texttt{evenPair\_card\_oneThousandThreeHundredThirty} & Gaps13221330 & PROVED & \checkmark & verified \\
\texttt{evenPair\_card\_oneThousandThreeHundredTwentyEight} & Gaps13221330 & PROVED & \checkmark & verified \\
\texttt{evenPair\_card\_oneThousandThreeHundredTwentyFour} & Gaps13221330 & PROVED & \checkmark & verified \\
\texttt{evenPair\_card\_oneThousandThreeHundredTwentySix} & Gaps13221330 & PROVED & \checkmark & verified \\
\texttt{evenPair\_card\_oneThousandThreeHundredTwentyTwo} & Gaps13221330 & PROVED & \checkmark & verified \\
\texttt{isAdmissible\_evenPair\_oneThousandThreeHundredThirty} & Gaps13221330 & PROVED & \checkmark & verified \\
\texttt{isAdmissible\_evenPair\_oneThousandThreeHundredTwentyEight} & Gaps13221330 & PROVED & \checkmark & verified \\
\texttt{isAdmissible\_evenPair\_oneThousandThreeHundredTwentyFour} & Gaps13221330 & PROVED & \checkmark & verified \\
\texttt{isAdmissible\_evenPair\_oneThousandThreeHundredTwentySix} & Gaps13221330 & PROVED & \checkmark & verified \\
\texttt{isAdmissible\_evenPair\_oneThousandThreeHundredTwentyTwo} & Gaps13221330 & PROVED & \checkmark & verified \\
\texttt{localFactor\_oneThousandThreeHundredThirty\_two} & Gaps13221330 & PROVED & \checkmark & verified \\
\texttt{localFactor\_oneThousandThreeHundredTwentyTwo\_two} & Gaps13221330 & PROVED & \checkmark & verified \\
\texttt{nu\_p\_oneThousandThreeHundredThirty} & Gaps13221330 & PROVED & \checkmark & verified \\
\texttt{nu\_p\_oneThousandThreeHundredThirty\_two} & Gaps13221330 & PROVED & \checkmark & verified \\
\texttt{nu\_p\_oneThousandThreeHundredTwentyEight} & Gaps13221330 & PROVED & \checkmark & verified \\
\texttt{nu\_p\_oneThousandThreeHundredTwentyFour} & Gaps13221330 & PROVED & \checkmark & verified \\
\texttt{nu\_p\_oneThousandThreeHundredTwentySix} & Gaps13221330 & PROVED & \checkmark & verified \\
\texttt{nu\_p\_oneThousandThreeHundredTwentyTwo} & Gaps13221330 & PROVED & \checkmark & verified \\
\texttt{nu\_p\_oneThousandThreeHundredTwentyTwo\_two} & Gaps13221330 & PROVED & \checkmark & verified \\
\texttt{singular\_series\_finite\_pos\_evenPair\_oneThousandThreeHundredThirty} & Gaps13221330 & PROVED & \checkmark & verified \\
\texttt{singular\_series\_finite\_pos\_evenPair\_oneThousandThreeHundredTwentyEight} & Gaps13221330 & PROVED & \checkmark & verified \\
\texttt{singular\_series\_finite\_pos\_evenPair\_oneThousandThreeHundredTwentyFour} & Gaps13221330 & PROVED & \checkmark & verified \\
\texttt{singular\_series\_finite\_pos\_evenPair\_oneThousandThreeHundredTwentySix} & Gaps13221330 & PROVED & \checkmark & verified \\
\texttt{singular\_series\_finite\_pos\_evenPair\_oneThousandThreeHundredTwentyTwo} & Gaps13221330 & PROVED & \checkmark & verified \\
\texttt{singular\_series\_pos\_evenPair\_oneThousandThreeHundredThirty} & Gaps13221330 & PROVED & \checkmark & verified \\
\texttt{singular\_series\_pos\_evenPair\_oneThousandThreeHundredTwentyEight} & Gaps13221330 & PROVED & \checkmark & verified \\
\texttt{singular\_series\_pos\_evenPair\_oneThousandThreeHundredTwentyFour} & Gaps13221330 & PROVED & \checkmark & verified \\
\texttt{singular\_series\_pos\_evenPair\_oneThousandThreeHundredTwentySix} & Gaps13221330 & PROVED & \checkmark & verified \\
\texttt{singular\_series\_pos\_evenPair\_oneThousandThreeHundredTwentyTwo} & Gaps13221330 & PROVED & \checkmark & verified \\
\texttt{evenPair\_card\_oneThousandThreeHundredForty} & Gaps13321340 & PROVED & \checkmark & verified \\
\texttt{evenPair\_card\_oneThousandThreeHundredThirtyEight} & Gaps13321340 & PROVED & \checkmark & verified \\
\texttt{evenPair\_card\_oneThousandThreeHundredThirtyFour} & Gaps13321340 & PROVED & \checkmark & verified \\
\texttt{evenPair\_card\_oneThousandThreeHundredThirtySix} & Gaps13321340 & PROVED & \checkmark & verified \\
\texttt{evenPair\_card\_oneThousandThreeHundredThirtyTwo} & Gaps13321340 & PROVED & \checkmark & verified \\
\texttt{isAdmissible\_evenPair\_oneThousandThreeHundredForty} & Gaps13321340 & PROVED & \checkmark & verified \\
\texttt{isAdmissible\_evenPair\_oneThousandThreeHundredThirtyEight} & Gaps13321340 & PROVED & \checkmark & verified \\
\texttt{isAdmissible\_evenPair\_oneThousandThreeHundredThirtyFour} & Gaps13321340 & PROVED & \checkmark & verified \\
\texttt{isAdmissible\_evenPair\_oneThousandThreeHundredThirtySix} & Gaps13321340 & PROVED & \checkmark & verified \\
\texttt{isAdmissible\_evenPair\_oneThousandThreeHundredThirtyTwo} & Gaps13321340 & PROVED & \checkmark & verified \\
\texttt{localFactor\_oneThousandThreeHundredForty\_two} & Gaps13321340 & PROVED & \checkmark & verified \\
\texttt{localFactor\_oneThousandThreeHundredThirtyTwo\_two} & Gaps13321340 & PROVED & \checkmark & verified \\
\texttt{nu\_p\_oneThousandThreeHundredForty} & Gaps13321340 & PROVED & \checkmark & verified \\
\texttt{nu\_p\_oneThousandThreeHundredForty\_two} & Gaps13321340 & PROVED & \checkmark & verified \\
\texttt{nu\_p\_oneThousandThreeHundredThirtyEight} & Gaps13321340 & PROVED & \checkmark & verified \\
\texttt{nu\_p\_oneThousandThreeHundredThirtyFour} & Gaps13321340 & PROVED & \checkmark & verified \\
\texttt{nu\_p\_oneThousandThreeHundredThirtySix} & Gaps13321340 & PROVED & \checkmark & verified \\
\texttt{nu\_p\_oneThousandThreeHundredThirtyTwo} & Gaps13321340 & PROVED & \checkmark & verified \\
\texttt{nu\_p\_oneThousandThreeHundredThirtyTwo\_two} & Gaps13321340 & PROVED & \checkmark & verified \\
\texttt{singular\_series\_finite\_pos\_evenPair\_oneThousandThreeHundredForty} & Gaps13321340 & PROVED & \checkmark & verified \\
\texttt{singular\_series\_finite\_pos\_evenPair\_oneThousandThreeHundredThirtyEight} & Gaps13321340 & PROVED & \checkmark & verified \\
\texttt{singular\_series\_finite\_pos\_evenPair\_oneThousandThreeHundredThirtyFour} & Gaps13321340 & PROVED & \checkmark & verified \\
\texttt{singular\_series\_finite\_pos\_evenPair\_oneThousandThreeHundredThirtySix} & Gaps13321340 & PROVED & \checkmark & verified \\
\texttt{singular\_series\_finite\_pos\_evenPair\_oneThousandThreeHundredThirtyTwo} & Gaps13321340 & PROVED & \checkmark & verified \\
\texttt{singular\_series\_pos\_evenPair\_oneThousandThreeHundredForty} & Gaps13321340 & PROVED & \checkmark & verified \\
\texttt{singular\_series\_pos\_evenPair\_oneThousandThreeHundredThirtyEight} & Gaps13321340 & PROVED & \checkmark & verified \\
\texttt{singular\_series\_pos\_evenPair\_oneThousandThreeHundredThirtyFour} & Gaps13321340 & PROVED & \checkmark & verified \\
\texttt{singular\_series\_pos\_evenPair\_oneThousandThreeHundredThirtySix} & Gaps13321340 & PROVED & \checkmark & verified \\
\texttt{singular\_series\_pos\_evenPair\_oneThousandThreeHundredThirtyTwo} & Gaps13321340 & PROVED & \checkmark & verified \\
\texttt{evenPair\_card\_oneThousandThreeHundredFifty} & Gaps13421350 & PROVED & \checkmark & verified \\
\texttt{evenPair\_card\_oneThousandThreeHundredFortyEight} & Gaps13421350 & PROVED & \checkmark & verified \\
\texttt{evenPair\_card\_oneThousandThreeHundredFortyFour} & Gaps13421350 & PROVED & \checkmark & verified \\
\texttt{evenPair\_card\_oneThousandThreeHundredFortySix} & Gaps13421350 & PROVED & \checkmark & verified \\
\texttt{evenPair\_card\_oneThousandThreeHundredFortyTwo} & Gaps13421350 & PROVED & \checkmark & verified \\
\texttt{isAdmissible\_evenPair\_oneThousandThreeHundredFifty} & Gaps13421350 & PROVED & \checkmark & verified \\
\texttt{isAdmissible\_evenPair\_oneThousandThreeHundredFortyEight} & Gaps13421350 & PROVED & \checkmark & verified \\
\texttt{isAdmissible\_evenPair\_oneThousandThreeHundredFortyFour} & Gaps13421350 & PROVED & \checkmark & verified \\
\texttt{isAdmissible\_evenPair\_oneThousandThreeHundredFortySix} & Gaps13421350 & PROVED & \checkmark & verified \\
\texttt{isAdmissible\_evenPair\_oneThousandThreeHundredFortyTwo} & Gaps13421350 & PROVED & \checkmark & verified \\
\texttt{localFactor\_oneThousandThreeHundredFifty\_two} & Gaps13421350 & PROVED & \checkmark & verified \\
\texttt{localFactor\_oneThousandThreeHundredFortyTwo\_two} & Gaps13421350 & PROVED & \checkmark & verified \\
\texttt{nu\_p\_oneThousandThreeHundredFifty} & Gaps13421350 & PROVED & \checkmark & verified \\
\texttt{nu\_p\_oneThousandThreeHundredFifty\_two} & Gaps13421350 & PROVED & \checkmark & verified \\
\texttt{nu\_p\_oneThousandThreeHundredFortyEight} & Gaps13421350 & PROVED & \checkmark & verified \\
\texttt{nu\_p\_oneThousandThreeHundredFortyFour} & Gaps13421350 & PROVED & \checkmark & verified \\
\texttt{nu\_p\_oneThousandThreeHundredFortySix} & Gaps13421350 & PROVED & \checkmark & verified \\
\texttt{nu\_p\_oneThousandThreeHundredFortyTwo} & Gaps13421350 & PROVED & \checkmark & verified \\
\texttt{nu\_p\_oneThousandThreeHundredFortyTwo\_two} & Gaps13421350 & PROVED & \checkmark & verified \\
\texttt{singular\_series\_finite\_pos\_evenPair\_oneThousandThreeHundredFifty} & Gaps13421350 & PROVED & \checkmark & verified \\
\texttt{singular\_series\_finite\_pos\_evenPair\_oneThousandThreeHundredFortyEight} & Gaps13421350 & PROVED & \checkmark & verified \\
\texttt{singular\_series\_finite\_pos\_evenPair\_oneThousandThreeHundredFortyFour} & Gaps13421350 & PROVED & \checkmark & verified \\
\texttt{singular\_series\_finite\_pos\_evenPair\_oneThousandThreeHundredFortySix} & Gaps13421350 & PROVED & \checkmark & verified \\
\texttt{singular\_series\_finite\_pos\_evenPair\_oneThousandThreeHundredFortyTwo} & Gaps13421350 & PROVED & \checkmark & verified \\
\texttt{singular\_series\_pos\_evenPair\_oneThousandThreeHundredFifty} & Gaps13421350 & PROVED & \checkmark & verified \\
\texttt{singular\_series\_pos\_evenPair\_oneThousandThreeHundredFortyEight} & Gaps13421350 & PROVED & \checkmark & verified \\
\texttt{singular\_series\_pos\_evenPair\_oneThousandThreeHundredFortyFour} & Gaps13421350 & PROVED & \checkmark & verified \\
\texttt{singular\_series\_pos\_evenPair\_oneThousandThreeHundredFortySix} & Gaps13421350 & PROVED & \checkmark & verified \\
\texttt{singular\_series\_pos\_evenPair\_oneThousandThreeHundredFortyTwo} & Gaps13421350 & PROVED & \checkmark & verified \\
\texttt{evenPair\_card\_oneThousandThreeHundredFiftyEight} & Gaps13521360 & PROVED & \checkmark & verified \\
\texttt{evenPair\_card\_oneThousandThreeHundredFiftyFour} & Gaps13521360 & PROVED & \checkmark & verified \\
\texttt{evenPair\_card\_oneThousandThreeHundredFiftySix} & Gaps13521360 & PROVED & \checkmark & verified \\
\texttt{evenPair\_card\_oneThousandThreeHundredFiftyTwo} & Gaps13521360 & PROVED & \checkmark & verified \\
\texttt{evenPair\_card\_oneThousandThreeHundredSixty} & Gaps13521360 & PROVED & \checkmark & verified \\
\texttt{isAdmissible\_evenPair\_oneThousandThreeHundredFiftyEight} & Gaps13521360 & PROVED & \checkmark & verified \\
\texttt{isAdmissible\_evenPair\_oneThousandThreeHundredFiftyFour} & Gaps13521360 & PROVED & \checkmark & verified \\
\texttt{isAdmissible\_evenPair\_oneThousandThreeHundredFiftySix} & Gaps13521360 & PROVED & \checkmark & verified \\
\texttt{isAdmissible\_evenPair\_oneThousandThreeHundredFiftyTwo} & Gaps13521360 & PROVED & \checkmark & verified \\
\texttt{isAdmissible\_evenPair\_oneThousandThreeHundredSixty} & Gaps13521360 & PROVED & \checkmark & verified \\
\texttt{localFactor\_oneThousandThreeHundredFiftyTwo\_two} & Gaps13521360 & PROVED & \checkmark & verified \\
\texttt{localFactor\_oneThousandThreeHundredSixty\_two} & Gaps13521360 & PROVED & \checkmark & verified \\
\texttt{nu\_p\_oneThousandThreeHundredFiftyEight} & Gaps13521360 & PROVED & \checkmark & verified \\
\texttt{nu\_p\_oneThousandThreeHundredFiftyFour} & Gaps13521360 & PROVED & \checkmark & verified \\
\texttt{nu\_p\_oneThousandThreeHundredFiftySix} & Gaps13521360 & PROVED & \checkmark & verified \\
\texttt{nu\_p\_oneThousandThreeHundredFiftyTwo} & Gaps13521360 & PROVED & \checkmark & verified \\
\texttt{nu\_p\_oneThousandThreeHundredFiftyTwo\_two} & Gaps13521360 & PROVED & \checkmark & verified \\
\texttt{nu\_p\_oneThousandThreeHundredSixty} & Gaps13521360 & PROVED & \checkmark & verified \\
\texttt{nu\_p\_oneThousandThreeHundredSixty\_two} & Gaps13521360 & PROVED & \checkmark & verified \\
\texttt{singular\_series\_finite\_pos\_evenPair\_oneThousandThreeHundredFiftyEight} & Gaps13521360 & PROVED & \checkmark & verified \\
\texttt{singular\_series\_finite\_pos\_evenPair\_oneThousandThreeHundredFiftyFour} & Gaps13521360 & PROVED & \checkmark & verified \\
\texttt{singular\_series\_finite\_pos\_evenPair\_oneThousandThreeHundredFiftySix} & Gaps13521360 & PROVED & \checkmark & verified \\
\texttt{singular\_series\_finite\_pos\_evenPair\_oneThousandThreeHundredFiftyTwo} & Gaps13521360 & PROVED & \checkmark & verified \\
\texttt{singular\_series\_finite\_pos\_evenPair\_oneThousandThreeHundredSixty} & Gaps13521360 & PROVED & \checkmark & verified \\
\texttt{singular\_series\_pos\_evenPair\_oneThousandThreeHundredFiftyEight} & Gaps13521360 & PROVED & \checkmark & verified \\
\texttt{singular\_series\_pos\_evenPair\_oneThousandThreeHundredFiftyFour} & Gaps13521360 & PROVED & \checkmark & verified \\
\texttt{singular\_series\_pos\_evenPair\_oneThousandThreeHundredFiftySix} & Gaps13521360 & PROVED & \checkmark & verified \\
\texttt{singular\_series\_pos\_evenPair\_oneThousandThreeHundredFiftyTwo} & Gaps13521360 & PROVED & \checkmark & verified \\
\texttt{singular\_series\_pos\_evenPair\_oneThousandThreeHundredSixty} & Gaps13521360 & PROVED & \checkmark & verified \\
\texttt{evenPair\_card\_oneThousandThreeHundredSeventy} & Gaps13621370 & PROVED & \checkmark & verified \\
\texttt{evenPair\_card\_oneThousandThreeHundredSixtyEight} & Gaps13621370 & PROVED & \checkmark & verified \\
\texttt{evenPair\_card\_oneThousandThreeHundredSixtyFour} & Gaps13621370 & PROVED & \checkmark & verified \\
\texttt{evenPair\_card\_oneThousandThreeHundredSixtySix} & Gaps13621370 & PROVED & \checkmark & verified \\
\texttt{evenPair\_card\_oneThousandThreeHundredSixtyTwo} & Gaps13621370 & PROVED & \checkmark & verified \\
\texttt{isAdmissible\_evenPair\_oneThousandThreeHundredSeventy} & Gaps13621370 & PROVED & \checkmark & verified \\
\texttt{isAdmissible\_evenPair\_oneThousandThreeHundredSixtyEight} & Gaps13621370 & PROVED & \checkmark & verified \\
\texttt{isAdmissible\_evenPair\_oneThousandThreeHundredSixtyFour} & Gaps13621370 & PROVED & \checkmark & verified \\
\texttt{isAdmissible\_evenPair\_oneThousandThreeHundredSixtySix} & Gaps13621370 & PROVED & \checkmark & verified \\
\texttt{isAdmissible\_evenPair\_oneThousandThreeHundredSixtyTwo} & Gaps13621370 & PROVED & \checkmark & verified \\
\texttt{localFactor\_oneThousandThreeHundredSeventy\_two} & Gaps13621370 & PROVED & \checkmark & verified \\
\texttt{localFactor\_oneThousandThreeHundredSixtyTwo\_two} & Gaps13621370 & PROVED & \checkmark & verified \\
\texttt{nu\_p\_oneThousandThreeHundredSeventy} & Gaps13621370 & PROVED & \checkmark & verified \\
\texttt{nu\_p\_oneThousandThreeHundredSeventy\_two} & Gaps13621370 & PROVED & \checkmark & verified \\
\texttt{nu\_p\_oneThousandThreeHundredSixtyEight} & Gaps13621370 & PROVED & \checkmark & verified \\
\texttt{nu\_p\_oneThousandThreeHundredSixtyFour} & Gaps13621370 & PROVED & \checkmark & verified \\
\texttt{nu\_p\_oneThousandThreeHundredSixtySix} & Gaps13621370 & PROVED & \checkmark & verified \\
\texttt{nu\_p\_oneThousandThreeHundredSixtyTwo} & Gaps13621370 & PROVED & \checkmark & verified \\
\texttt{nu\_p\_oneThousandThreeHundredSixtyTwo\_two} & Gaps13621370 & PROVED & \checkmark & verified \\
\texttt{singular\_series\_finite\_pos\_evenPair\_oneThousandThreeHundredSeventy} & Gaps13621370 & PROVED & \checkmark & verified \\
\texttt{singular\_series\_finite\_pos\_evenPair\_oneThousandThreeHundredSixtyEight} & Gaps13621370 & PROVED & \checkmark & verified \\
\texttt{singular\_series\_finite\_pos\_evenPair\_oneThousandThreeHundredSixtyFour} & Gaps13621370 & PROVED & \checkmark & verified \\
\texttt{singular\_series\_finite\_pos\_evenPair\_oneThousandThreeHundredSixtySix} & Gaps13621370 & PROVED & \checkmark & verified \\
\texttt{singular\_series\_finite\_pos\_evenPair\_oneThousandThreeHundredSixtyTwo} & Gaps13621370 & PROVED & \checkmark & verified \\
\texttt{singular\_series\_pos\_evenPair\_oneThousandThreeHundredSeventy} & Gaps13621370 & PROVED & \checkmark & verified \\
\texttt{singular\_series\_pos\_evenPair\_oneThousandThreeHundredSixtyEight} & Gaps13621370 & PROVED & \checkmark & verified \\
\texttt{singular\_series\_pos\_evenPair\_oneThousandThreeHundredSixtyFour} & Gaps13621370 & PROVED & \checkmark & verified \\
\texttt{singular\_series\_pos\_evenPair\_oneThousandThreeHundredSixtySix} & Gaps13621370 & PROVED & \checkmark & verified \\
\texttt{singular\_series\_pos\_evenPair\_oneThousandThreeHundredSixtyTwo} & Gaps13621370 & PROVED & \checkmark & verified \\
\texttt{evenPair\_card\_oneThousandThreeHundredEighty} & Gaps13721380 & PROVED & \checkmark & verified \\
\texttt{evenPair\_card\_oneThousandThreeHundredSeventyEight} & Gaps13721380 & PROVED & \checkmark & verified \\
\texttt{evenPair\_card\_oneThousandThreeHundredSeventyFour} & Gaps13721380 & PROVED & \checkmark & verified \\
\texttt{evenPair\_card\_oneThousandThreeHundredSeventySix} & Gaps13721380 & PROVED & \checkmark & verified \\
\texttt{evenPair\_card\_oneThousandThreeHundredSeventyTwo} & Gaps13721380 & PROVED & \checkmark & verified \\
\texttt{isAdmissible\_evenPair\_oneThousandThreeHundredEighty} & Gaps13721380 & PROVED & \checkmark & verified \\
\texttt{isAdmissible\_evenPair\_oneThousandThreeHundredSeventyEight} & Gaps13721380 & PROVED & \checkmark & verified \\
\texttt{isAdmissible\_evenPair\_oneThousandThreeHundredSeventyFour} & Gaps13721380 & PROVED & \checkmark & verified \\
\texttt{isAdmissible\_evenPair\_oneThousandThreeHundredSeventySix} & Gaps13721380 & PROVED & \checkmark & verified \\
\texttt{isAdmissible\_evenPair\_oneThousandThreeHundredSeventyTwo} & Gaps13721380 & PROVED & \checkmark & verified \\
\texttt{localFactor\_oneThousandThreeHundredEighty\_two} & Gaps13721380 & PROVED & \checkmark & verified \\
\texttt{localFactor\_oneThousandThreeHundredSeventyTwo\_two} & Gaps13721380 & PROVED & \checkmark & verified \\
\texttt{nu\_p\_oneThousandThreeHundredEighty} & Gaps13721380 & PROVED & \checkmark & verified \\
\texttt{nu\_p\_oneThousandThreeHundredEighty\_two} & Gaps13721380 & PROVED & \checkmark & verified \\
\texttt{nu\_p\_oneThousandThreeHundredSeventyEight} & Gaps13721380 & PROVED & \checkmark & verified \\
\texttt{nu\_p\_oneThousandThreeHundredSeventyFour} & Gaps13721380 & PROVED & \checkmark & verified \\
\texttt{nu\_p\_oneThousandThreeHundredSeventySix} & Gaps13721380 & PROVED & \checkmark & verified \\
\texttt{nu\_p\_oneThousandThreeHundredSeventyTwo} & Gaps13721380 & PROVED & \checkmark & verified \\
\texttt{nu\_p\_oneThousandThreeHundredSeventyTwo\_two} & Gaps13721380 & PROVED & \checkmark & verified \\
\texttt{singular\_series\_finite\_pos\_evenPair\_oneThousandThreeHundredEighty} & Gaps13721380 & PROVED & \checkmark & verified \\
\texttt{singular\_series\_finite\_pos\_evenPair\_oneThousandThreeHundredSeventyEight} & Gaps13721380 & PROVED & \checkmark & verified \\
\texttt{singular\_series\_finite\_pos\_evenPair\_oneThousandThreeHundredSeventyFour} & Gaps13721380 & PROVED & \checkmark & verified \\
\texttt{singular\_series\_finite\_pos\_evenPair\_oneThousandThreeHundredSeventySix} & Gaps13721380 & PROVED & \checkmark & verified \\
\texttt{singular\_series\_finite\_pos\_evenPair\_oneThousandThreeHundredSeventyTwo} & Gaps13721380 & PROVED & \checkmark & verified \\
\texttt{singular\_series\_pos\_evenPair\_oneThousandThreeHundredEighty} & Gaps13721380 & PROVED & \checkmark & verified \\
\texttt{singular\_series\_pos\_evenPair\_oneThousandThreeHundredSeventyEight} & Gaps13721380 & PROVED & \checkmark & verified \\
\texttt{singular\_series\_pos\_evenPair\_oneThousandThreeHundredSeventyFour} & Gaps13721380 & PROVED & \checkmark & verified \\
\texttt{singular\_series\_pos\_evenPair\_oneThousandThreeHundredSeventySix} & Gaps13721380 & PROVED & \checkmark & verified \\
\texttt{singular\_series\_pos\_evenPair\_oneThousandThreeHundredSeventyTwo} & Gaps13721380 & PROVED & \checkmark & verified \\
\texttt{evenPair\_card\_oneThousandThreeHundredEightyEight} & Gaps13821390 & PROVED & \checkmark & verified \\
\texttt{evenPair\_card\_oneThousandThreeHundredEightyFour} & Gaps13821390 & PROVED & \checkmark & verified \\
\texttt{evenPair\_card\_oneThousandThreeHundredEightySix} & Gaps13821390 & PROVED & \checkmark & verified \\
\texttt{evenPair\_card\_oneThousandThreeHundredEightyTwo} & Gaps13821390 & PROVED & \checkmark & verified \\
\texttt{evenPair\_card\_oneThousandThreeHundredNinety} & Gaps13821390 & PROVED & \checkmark & verified \\
\texttt{isAdmissible\_evenPair\_oneThousandThreeHundredEightyEight} & Gaps13821390 & PROVED & \checkmark & verified \\
\texttt{isAdmissible\_evenPair\_oneThousandThreeHundredEightyFour} & Gaps13821390 & PROVED & \checkmark & verified \\
\texttt{isAdmissible\_evenPair\_oneThousandThreeHundredEightySix} & Gaps13821390 & PROVED & \checkmark & verified \\
\texttt{isAdmissible\_evenPair\_oneThousandThreeHundredEightyTwo} & Gaps13821390 & PROVED & \checkmark & verified \\
\texttt{isAdmissible\_evenPair\_oneThousandThreeHundredNinety} & Gaps13821390 & PROVED & \checkmark & verified \\
\texttt{localFactor\_oneThousandThreeHundredEightyTwo\_two} & Gaps13821390 & PROVED & \checkmark & verified \\
\texttt{localFactor\_oneThousandThreeHundredNinety\_two} & Gaps13821390 & PROVED & \checkmark & verified \\
\texttt{nu\_p\_oneThousandThreeHundredEightyEight} & Gaps13821390 & PROVED & \checkmark & verified \\
\texttt{nu\_p\_oneThousandThreeHundredEightyFour} & Gaps13821390 & PROVED & \checkmark & verified \\
\texttt{nu\_p\_oneThousandThreeHundredEightySix} & Gaps13821390 & PROVED & \checkmark & verified \\
\texttt{nu\_p\_oneThousandThreeHundredEightyTwo} & Gaps13821390 & PROVED & \checkmark & verified \\
\texttt{nu\_p\_oneThousandThreeHundredEightyTwo\_two} & Gaps13821390 & PROVED & \checkmark & verified \\
\texttt{nu\_p\_oneThousandThreeHundredNinety} & Gaps13821390 & PROVED & \checkmark & verified \\
\texttt{nu\_p\_oneThousandThreeHundredNinety\_two} & Gaps13821390 & PROVED & \checkmark & verified \\
\texttt{singular\_series\_finite\_pos\_evenPair\_oneThousandThreeHundredEightyEight} & Gaps13821390 & PROVED & \checkmark & verified \\
\texttt{singular\_series\_finite\_pos\_evenPair\_oneThousandThreeHundredEightyFour} & Gaps13821390 & PROVED & \checkmark & verified \\
\texttt{singular\_series\_finite\_pos\_evenPair\_oneThousandThreeHundredEightySix} & Gaps13821390 & PROVED & \checkmark & verified \\
\texttt{singular\_series\_finite\_pos\_evenPair\_oneThousandThreeHundredEightyTwo} & Gaps13821390 & PROVED & \checkmark & verified \\
\texttt{singular\_series\_finite\_pos\_evenPair\_oneThousandThreeHundredNinety} & Gaps13821390 & PROVED & \checkmark & verified \\
\texttt{singular\_series\_pos\_evenPair\_oneThousandThreeHundredEightyEight} & Gaps13821390 & PROVED & \checkmark & verified \\
\texttt{singular\_series\_pos\_evenPair\_oneThousandThreeHundredEightyFour} & Gaps13821390 & PROVED & \checkmark & verified \\
\texttt{singular\_series\_pos\_evenPair\_oneThousandThreeHundredEightySix} & Gaps13821390 & PROVED & \checkmark & verified \\
\texttt{singular\_series\_pos\_evenPair\_oneThousandThreeHundredEightyTwo} & Gaps13821390 & PROVED & \checkmark & verified \\
\texttt{singular\_series\_pos\_evenPair\_oneThousandThreeHundredNinety} & Gaps13821390 & PROVED & \checkmark & verified \\
\texttt{evenPair\_card\_oneThousandFourHundred} & Gaps13921400 & PROVED & \checkmark & verified \\
\texttt{evenPair\_card\_oneThousandThreeHundredNinetyEight} & Gaps13921400 & PROVED & \checkmark & verified \\
\texttt{evenPair\_card\_oneThousandThreeHundredNinetyFour} & Gaps13921400 & PROVED & \checkmark & verified \\
\texttt{evenPair\_card\_oneThousandThreeHundredNinetySix} & Gaps13921400 & PROVED & \checkmark & verified \\
\texttt{evenPair\_card\_oneThousandThreeHundredNinetyTwo} & Gaps13921400 & PROVED & \checkmark & verified \\
\texttt{isAdmissible\_evenPair\_oneThousandFourHundred} & Gaps13921400 & PROVED & \checkmark & verified \\
\texttt{isAdmissible\_evenPair\_oneThousandThreeHundredNinetyEight} & Gaps13921400 & PROVED & \checkmark & verified \\
\texttt{isAdmissible\_evenPair\_oneThousandThreeHundredNinetyFour} & Gaps13921400 & PROVED & \checkmark & verified \\
\texttt{isAdmissible\_evenPair\_oneThousandThreeHundredNinetySix} & Gaps13921400 & PROVED & \checkmark & verified \\
\texttt{isAdmissible\_evenPair\_oneThousandThreeHundredNinetyTwo} & Gaps13921400 & PROVED & \checkmark & verified \\
\texttt{localFactor\_oneThousandFourHundred\_two} & Gaps13921400 & PROVED & \checkmark & verified \\
\texttt{localFactor\_oneThousandThreeHundredNinetyTwo\_two} & Gaps13921400 & PROVED & \checkmark & verified \\
\texttt{nu\_p\_oneThousandFourHundred} & Gaps13921400 & PROVED & \checkmark & verified \\
\texttt{nu\_p\_oneThousandFourHundred\_two} & Gaps13921400 & PROVED & \checkmark & verified \\
\texttt{nu\_p\_oneThousandThreeHundredNinetyEight} & Gaps13921400 & PROVED & \checkmark & verified \\
\texttt{nu\_p\_oneThousandThreeHundredNinetyFour} & Gaps13921400 & PROVED & \checkmark & verified \\
\texttt{nu\_p\_oneThousandThreeHundredNinetySix} & Gaps13921400 & PROVED & \checkmark & verified \\
\texttt{nu\_p\_oneThousandThreeHundredNinetyTwo} & Gaps13921400 & PROVED & \checkmark & verified \\
\texttt{nu\_p\_oneThousandThreeHundredNinetyTwo\_two} & Gaps13921400 & PROVED & \checkmark & verified \\
\texttt{singular\_series\_finite\_pos\_evenPair\_oneThousandFourHundred} & Gaps13921400 & PROVED & \checkmark & verified \\
\texttt{singular\_series\_finite\_pos\_evenPair\_oneThousandThreeHundredNinetyEight} & Gaps13921400 & PROVED & \checkmark & verified \\
\texttt{singular\_series\_finite\_pos\_evenPair\_oneThousandThreeHundredNinetyFour} & Gaps13921400 & PROVED & \checkmark & verified \\
\texttt{singular\_series\_finite\_pos\_evenPair\_oneThousandThreeHundredNinetySix} & Gaps13921400 & PROVED & \checkmark & verified \\
\texttt{singular\_series\_finite\_pos\_evenPair\_oneThousandThreeHundredNinetyTwo} & Gaps13921400 & PROVED & \checkmark & verified \\
\texttt{singular\_series\_pos\_evenPair\_oneThousandFourHundred} & Gaps13921400 & PROVED & \checkmark & verified \\
\texttt{singular\_series\_pos\_evenPair\_oneThousandThreeHundredNinetyEight} & Gaps13921400 & PROVED & \checkmark & verified \\
\texttt{singular\_series\_pos\_evenPair\_oneThousandThreeHundredNinetyFour} & Gaps13921400 & PROVED & \checkmark & verified \\
\texttt{singular\_series\_pos\_evenPair\_oneThousandThreeHundredNinetySix} & Gaps13921400 & PROVED & \checkmark & verified \\
\texttt{singular\_series\_pos\_evenPair\_oneThousandThreeHundredNinetyTwo} & Gaps13921400 & PROVED & \checkmark & verified \\
\texttt{evenPair\_card\_oneThousandFourHundredEight} & Gaps14021410 & PROVED & \checkmark & verified \\
\texttt{evenPair\_card\_oneThousandFourHundredFour} & Gaps14021410 & PROVED & \checkmark & verified \\
\texttt{evenPair\_card\_oneThousandFourHundredSix} & Gaps14021410 & PROVED & \checkmark & verified \\
\texttt{evenPair\_card\_oneThousandFourHundredTen} & Gaps14021410 & PROVED & \checkmark & verified \\
\texttt{evenPair\_card\_oneThousandFourHundredTwo} & Gaps14021410 & PROVED & \checkmark & verified \\
\texttt{isAdmissible\_evenPair\_oneThousandFourHundredEight} & Gaps14021410 & PROVED & \checkmark & verified \\
\texttt{isAdmissible\_evenPair\_oneThousandFourHundredFour} & Gaps14021410 & PROVED & \checkmark & verified \\
\texttt{isAdmissible\_evenPair\_oneThousandFourHundredSix} & Gaps14021410 & PROVED & \checkmark & verified \\
\texttt{isAdmissible\_evenPair\_oneThousandFourHundredTen} & Gaps14021410 & PROVED & \checkmark & verified \\
\texttt{isAdmissible\_evenPair\_oneThousandFourHundredTwo} & Gaps14021410 & PROVED & \checkmark & verified \\
\texttt{localFactor\_oneThousandFourHundredTen\_two} & Gaps14021410 & PROVED & \checkmark & verified \\
\texttt{localFactor\_oneThousandFourHundredTwo\_two} & Gaps14021410 & PROVED & \checkmark & verified \\
\texttt{nu\_p\_oneThousandFourHundredEight} & Gaps14021410 & PROVED & \checkmark & verified \\
\texttt{nu\_p\_oneThousandFourHundredFour} & Gaps14021410 & PROVED & \checkmark & verified \\
\texttt{nu\_p\_oneThousandFourHundredSix} & Gaps14021410 & PROVED & \checkmark & verified \\
\texttt{nu\_p\_oneThousandFourHundredTen} & Gaps14021410 & PROVED & \checkmark & verified \\
\texttt{nu\_p\_oneThousandFourHundredTen\_two} & Gaps14021410 & PROVED & \checkmark & verified \\
\texttt{nu\_p\_oneThousandFourHundredTwo} & Gaps14021410 & PROVED & \checkmark & verified \\
\texttt{nu\_p\_oneThousandFourHundredTwo\_two} & Gaps14021410 & PROVED & \checkmark & verified \\
\texttt{singular\_series\_finite\_pos\_evenPair\_oneThousandFourHundredEight} & Gaps14021410 & PROVED & \checkmark & verified \\
\texttt{singular\_series\_finite\_pos\_evenPair\_oneThousandFourHundredFour} & Gaps14021410 & PROVED & \checkmark & verified \\
\texttt{singular\_series\_finite\_pos\_evenPair\_oneThousandFourHundredSix} & Gaps14021410 & PROVED & \checkmark & verified \\
\texttt{singular\_series\_finite\_pos\_evenPair\_oneThousandFourHundredTen} & Gaps14021410 & PROVED & \checkmark & verified \\
\texttt{singular\_series\_finite\_pos\_evenPair\_oneThousandFourHundredTwo} & Gaps14021410 & PROVED & \checkmark & verified \\
\texttt{singular\_series\_pos\_evenPair\_oneThousandFourHundredEight} & Gaps14021410 & PROVED & \checkmark & verified \\
\texttt{singular\_series\_pos\_evenPair\_oneThousandFourHundredFour} & Gaps14021410 & PROVED & \checkmark & verified \\
\texttt{singular\_series\_pos\_evenPair\_oneThousandFourHundredSix} & Gaps14021410 & PROVED & \checkmark & verified \\
\texttt{singular\_series\_pos\_evenPair\_oneThousandFourHundredTen} & Gaps14021410 & PROVED & \checkmark & verified \\
\texttt{singular\_series\_pos\_evenPair\_oneThousandFourHundredTwo} & Gaps14021410 & PROVED & \checkmark & verified \\
\texttt{evenPair\_card\_oneThousandFourHundredEighteen} & Gaps14121420 & PROVED & \checkmark & verified \\
\texttt{evenPair\_card\_oneThousandFourHundredFourteen} & Gaps14121420 & PROVED & \checkmark & verified \\
\texttt{evenPair\_card\_oneThousandFourHundredSixteen} & Gaps14121420 & PROVED & \checkmark & verified \\
\texttt{evenPair\_card\_oneThousandFourHundredTwelve} & Gaps14121420 & PROVED & \checkmark & verified \\
\texttt{evenPair\_card\_oneThousandFourHundredTwenty} & Gaps14121420 & PROVED & \checkmark & verified \\
\texttt{isAdmissible\_evenPair\_oneThousandFourHundredEighteen} & Gaps14121420 & PROVED & \checkmark & verified \\
\texttt{isAdmissible\_evenPair\_oneThousandFourHundredFourteen} & Gaps14121420 & PROVED & \checkmark & verified \\
\texttt{isAdmissible\_evenPair\_oneThousandFourHundredSixteen} & Gaps14121420 & PROVED & \checkmark & verified \\
\texttt{isAdmissible\_evenPair\_oneThousandFourHundredTwelve} & Gaps14121420 & PROVED & \checkmark & verified \\
\texttt{isAdmissible\_evenPair\_oneThousandFourHundredTwenty} & Gaps14121420 & PROVED & \checkmark & verified \\
\texttt{localFactor\_oneThousandFourHundredTwelve\_two} & Gaps14121420 & PROVED & \checkmark & verified \\
\texttt{localFactor\_oneThousandFourHundredTwenty\_two} & Gaps14121420 & PROVED & \checkmark & verified \\
\texttt{nu\_p\_oneThousandFourHundredEighteen} & Gaps14121420 & PROVED & \checkmark & verified \\
\texttt{nu\_p\_oneThousandFourHundredFourteen} & Gaps14121420 & PROVED & \checkmark & verified \\
\texttt{nu\_p\_oneThousandFourHundredSixteen} & Gaps14121420 & PROVED & \checkmark & verified \\
\texttt{nu\_p\_oneThousandFourHundredTwelve} & Gaps14121420 & PROVED & \checkmark & verified \\
\texttt{nu\_p\_oneThousandFourHundredTwelve\_two} & Gaps14121420 & PROVED & \checkmark & verified \\
\texttt{nu\_p\_oneThousandFourHundredTwenty} & Gaps14121420 & PROVED & \checkmark & verified \\
\texttt{nu\_p\_oneThousandFourHundredTwenty\_two} & Gaps14121420 & PROVED & \checkmark & verified \\
\texttt{singular\_series\_finite\_pos\_evenPair\_oneThousandFourHundredEighteen} & Gaps14121420 & PROVED & \checkmark & verified \\
\texttt{singular\_series\_finite\_pos\_evenPair\_oneThousandFourHundredFourteen} & Gaps14121420 & PROVED & \checkmark & verified \\
\texttt{singular\_series\_finite\_pos\_evenPair\_oneThousandFourHundredSixteen} & Gaps14121420 & PROVED & \checkmark & verified \\
\texttt{singular\_series\_finite\_pos\_evenPair\_oneThousandFourHundredTwelve} & Gaps14121420 & PROVED & \checkmark & verified \\
\texttt{singular\_series\_finite\_pos\_evenPair\_oneThousandFourHundredTwenty} & Gaps14121420 & PROVED & \checkmark & verified \\
\texttt{singular\_series\_pos\_evenPair\_oneThousandFourHundredEighteen} & Gaps14121420 & PROVED & \checkmark & verified \\
\texttt{singular\_series\_pos\_evenPair\_oneThousandFourHundredFourteen} & Gaps14121420 & PROVED & \checkmark & verified \\
\texttt{singular\_series\_pos\_evenPair\_oneThousandFourHundredSixteen} & Gaps14121420 & PROVED & \checkmark & verified \\
\texttt{singular\_series\_pos\_evenPair\_oneThousandFourHundredTwelve} & Gaps14121420 & PROVED & \checkmark & verified \\
\texttt{singular\_series\_pos\_evenPair\_oneThousandFourHundredTwenty} & Gaps14121420 & PROVED & \checkmark & verified \\
\texttt{evenPair\_card\_oneHundredFifty} & Gaps142150 & PROVED & \checkmark & verified \\
\texttt{evenPair\_card\_oneHundredFortyEight} & Gaps142150 & PROVED & \checkmark & verified \\
\texttt{evenPair\_card\_oneHundredFortyFour} & Gaps142150 & PROVED & \checkmark & verified \\
\texttt{evenPair\_card\_oneHundredFortySix} & Gaps142150 & PROVED & \checkmark & verified \\
\texttt{evenPair\_card\_oneHundredFortyTwo} & Gaps142150 & PROVED & \checkmark & verified \\
\texttt{isAdmissible\_evenPair\_oneHundredFifty} & Gaps142150 & PROVED & \checkmark & verified \\
\texttt{isAdmissible\_evenPair\_oneHundredFortyEight} & Gaps142150 & PROVED & \checkmark & verified \\
\texttt{isAdmissible\_evenPair\_oneHundredFortyFour} & Gaps142150 & PROVED & \checkmark & verified \\
\texttt{isAdmissible\_evenPair\_oneHundredFortySix} & Gaps142150 & PROVED & \checkmark & verified \\
\texttt{isAdmissible\_evenPair\_oneHundredFortyTwo} & Gaps142150 & PROVED & \checkmark & verified \\
\texttt{localFactor\_oneHundredFifty\_two} & Gaps142150 & PROVED & \checkmark & verified \\
\texttt{localFactor\_oneHundredFortyTwo\_two} & Gaps142150 & PROVED & \checkmark & verified \\
\texttt{nu\_p\_oneHundredFifty} & Gaps142150 & PROVED & \checkmark & verified \\
\texttt{nu\_p\_oneHundredFifty\_two} & Gaps142150 & PROVED & \checkmark & verified \\
\texttt{nu\_p\_oneHundredFortyEight} & Gaps142150 & PROVED & \checkmark & verified \\
\texttt{nu\_p\_oneHundredFortyFour} & Gaps142150 & PROVED & \checkmark & verified \\
\texttt{nu\_p\_oneHundredFortySix} & Gaps142150 & PROVED & \checkmark & verified \\
\texttt{nu\_p\_oneHundredFortyTwo} & Gaps142150 & PROVED & \checkmark & verified \\
\texttt{nu\_p\_oneHundredFortyTwo\_two} & Gaps142150 & PROVED & \checkmark & verified \\
\texttt{singular\_series\_finite\_pos\_evenPair\_oneHundredFifty} & Gaps142150 & PROVED & \checkmark & verified \\
\texttt{singular\_series\_finite\_pos\_evenPair\_oneHundredFortyEight} & Gaps142150 & PROVED & \checkmark & verified \\
\texttt{singular\_series\_finite\_pos\_evenPair\_oneHundredFortyFour} & Gaps142150 & PROVED & \checkmark & verified \\
\texttt{singular\_series\_finite\_pos\_evenPair\_oneHundredFortySix} & Gaps142150 & PROVED & \checkmark & verified \\
\texttt{singular\_series\_finite\_pos\_evenPair\_oneHundredFortyTwo} & Gaps142150 & PROVED & \checkmark & verified \\
\texttt{singular\_series\_pos\_evenPair\_oneHundredFifty} & Gaps142150 & PROVED & \checkmark & verified \\
\texttt{singular\_series\_pos\_evenPair\_oneHundredFortyEight} & Gaps142150 & PROVED & \checkmark & verified \\
\texttt{singular\_series\_pos\_evenPair\_oneHundredFortyFour} & Gaps142150 & PROVED & \checkmark & verified \\
\texttt{singular\_series\_pos\_evenPair\_oneHundredFortySix} & Gaps142150 & PROVED & \checkmark & verified \\
\texttt{singular\_series\_pos\_evenPair\_oneHundredFortyTwo} & Gaps142150 & PROVED & \checkmark & verified \\
\texttt{evenPair\_card\_oneThousandFourHundredThirty} & Gaps14221430 & PROVED & \checkmark & verified \\
\texttt{evenPair\_card\_oneThousandFourHundredTwentyEight} & Gaps14221430 & PROVED & \checkmark & verified \\
\texttt{evenPair\_card\_oneThousandFourHundredTwentyFour} & Gaps14221430 & PROVED & \checkmark & verified \\
\texttt{evenPair\_card\_oneThousandFourHundredTwentySix} & Gaps14221430 & PROVED & \checkmark & verified \\
\texttt{evenPair\_card\_oneThousandFourHundredTwentyTwo} & Gaps14221430 & PROVED & \checkmark & verified \\
\texttt{isAdmissible\_evenPair\_oneThousandFourHundredThirty} & Gaps14221430 & PROVED & \checkmark & verified \\
\texttt{isAdmissible\_evenPair\_oneThousandFourHundredTwentyEight} & Gaps14221430 & PROVED & \checkmark & verified \\
\texttt{isAdmissible\_evenPair\_oneThousandFourHundredTwentyFour} & Gaps14221430 & PROVED & \checkmark & verified \\
\texttt{isAdmissible\_evenPair\_oneThousandFourHundredTwentySix} & Gaps14221430 & PROVED & \checkmark & verified \\
\texttt{isAdmissible\_evenPair\_oneThousandFourHundredTwentyTwo} & Gaps14221430 & PROVED & \checkmark & verified \\
\texttt{localFactor\_oneThousandFourHundredThirty\_two} & Gaps14221430 & PROVED & \checkmark & verified \\
\texttt{localFactor\_oneThousandFourHundredTwentyTwo\_two} & Gaps14221430 & PROVED & \checkmark & verified \\
\texttt{nu\_p\_oneThousandFourHundredThirty} & Gaps14221430 & PROVED & \checkmark & verified \\
\texttt{nu\_p\_oneThousandFourHundredThirty\_two} & Gaps14221430 & PROVED & \checkmark & verified \\
\texttt{nu\_p\_oneThousandFourHundredTwentyEight} & Gaps14221430 & PROVED & \checkmark & verified \\
\texttt{nu\_p\_oneThousandFourHundredTwentyFour} & Gaps14221430 & PROVED & \checkmark & verified \\
\texttt{nu\_p\_oneThousandFourHundredTwentySix} & Gaps14221430 & PROVED & \checkmark & verified \\
\texttt{nu\_p\_oneThousandFourHundredTwentyTwo} & Gaps14221430 & PROVED & \checkmark & verified \\
\texttt{nu\_p\_oneThousandFourHundredTwentyTwo\_two} & Gaps14221430 & PROVED & \checkmark & verified \\
\texttt{singular\_series\_finite\_pos\_evenPair\_oneThousandFourHundredThirty} & Gaps14221430 & PROVED & \checkmark & verified \\
\texttt{singular\_series\_finite\_pos\_evenPair\_oneThousandFourHundredTwentyEight} & Gaps14221430 & PROVED & \checkmark & verified \\
\texttt{singular\_series\_finite\_pos\_evenPair\_oneThousandFourHundredTwentyFour} & Gaps14221430 & PROVED & \checkmark & verified \\
\texttt{singular\_series\_finite\_pos\_evenPair\_oneThousandFourHundredTwentySix} & Gaps14221430 & PROVED & \checkmark & verified \\
\texttt{singular\_series\_finite\_pos\_evenPair\_oneThousandFourHundredTwentyTwo} & Gaps14221430 & PROVED & \checkmark & verified \\
\texttt{singular\_series\_pos\_evenPair\_oneThousandFourHundredThirty} & Gaps14221430 & PROVED & \checkmark & verified \\
\texttt{singular\_series\_pos\_evenPair\_oneThousandFourHundredTwentyEight} & Gaps14221430 & PROVED & \checkmark & verified \\
\texttt{singular\_series\_pos\_evenPair\_oneThousandFourHundredTwentyFour} & Gaps14221430 & PROVED & \checkmark & verified \\
\texttt{singular\_series\_pos\_evenPair\_oneThousandFourHundredTwentySix} & Gaps14221430 & PROVED & \checkmark & verified \\
\texttt{singular\_series\_pos\_evenPair\_oneThousandFourHundredTwentyTwo} & Gaps14221430 & PROVED & \checkmark & verified \\
\texttt{evenPair\_card\_oneThousandFourHundredForty} & Gaps14321440 & PROVED & \checkmark & verified \\
\texttt{evenPair\_card\_oneThousandFourHundredThirtyEight} & Gaps14321440 & PROVED & \checkmark & verified \\
\texttt{evenPair\_card\_oneThousandFourHundredThirtyFour} & Gaps14321440 & PROVED & \checkmark & verified \\
\texttt{evenPair\_card\_oneThousandFourHundredThirtySix} & Gaps14321440 & PROVED & \checkmark & verified \\
\texttt{evenPair\_card\_oneThousandFourHundredThirtyTwo} & Gaps14321440 & PROVED & \checkmark & verified \\
\texttt{isAdmissible\_evenPair\_oneThousandFourHundredForty} & Gaps14321440 & PROVED & \checkmark & verified \\
\texttt{isAdmissible\_evenPair\_oneThousandFourHundredThirtyEight} & Gaps14321440 & PROVED & \checkmark & verified \\
\texttt{isAdmissible\_evenPair\_oneThousandFourHundredThirtyFour} & Gaps14321440 & PROVED & \checkmark & verified \\
\texttt{isAdmissible\_evenPair\_oneThousandFourHundredThirtySix} & Gaps14321440 & PROVED & \checkmark & verified \\
\texttt{isAdmissible\_evenPair\_oneThousandFourHundredThirtyTwo} & Gaps14321440 & PROVED & \checkmark & verified \\
\texttt{localFactor\_oneThousandFourHundredForty\_two} & Gaps14321440 & PROVED & \checkmark & verified \\
\texttt{localFactor\_oneThousandFourHundredThirtyTwo\_two} & Gaps14321440 & PROVED & \checkmark & verified \\
\texttt{nu\_p\_oneThousandFourHundredForty} & Gaps14321440 & PROVED & \checkmark & verified \\
\texttt{nu\_p\_oneThousandFourHundredForty\_two} & Gaps14321440 & PROVED & \checkmark & verified \\
\texttt{nu\_p\_oneThousandFourHundredThirtyEight} & Gaps14321440 & PROVED & \checkmark & verified \\
\texttt{nu\_p\_oneThousandFourHundredThirtyFour} & Gaps14321440 & PROVED & \checkmark & verified \\
\texttt{nu\_p\_oneThousandFourHundredThirtySix} & Gaps14321440 & PROVED & \checkmark & verified \\
\texttt{nu\_p\_oneThousandFourHundredThirtyTwo} & Gaps14321440 & PROVED & \checkmark & verified \\
\texttt{nu\_p\_oneThousandFourHundredThirtyTwo\_two} & Gaps14321440 & PROVED & \checkmark & verified \\
\texttt{singular\_series\_finite\_pos\_evenPair\_oneThousandFourHundredForty} & Gaps14321440 & PROVED & \checkmark & verified \\
\texttt{singular\_series\_finite\_pos\_evenPair\_oneThousandFourHundredThirtyEight} & Gaps14321440 & PROVED & \checkmark & verified \\
\texttt{singular\_series\_finite\_pos\_evenPair\_oneThousandFourHundredThirtyFour} & Gaps14321440 & PROVED & \checkmark & verified \\
\texttt{singular\_series\_finite\_pos\_evenPair\_oneThousandFourHundredThirtySix} & Gaps14321440 & PROVED & \checkmark & verified \\
\texttt{singular\_series\_finite\_pos\_evenPair\_oneThousandFourHundredThirtyTwo} & Gaps14321440 & PROVED & \checkmark & verified \\
\texttt{singular\_series\_pos\_evenPair\_oneThousandFourHundredForty} & Gaps14321440 & PROVED & \checkmark & verified \\
\texttt{singular\_series\_pos\_evenPair\_oneThousandFourHundredThirtyEight} & Gaps14321440 & PROVED & \checkmark & verified \\
\texttt{singular\_series\_pos\_evenPair\_oneThousandFourHundredThirtyFour} & Gaps14321440 & PROVED & \checkmark & verified \\
\texttt{singular\_series\_pos\_evenPair\_oneThousandFourHundredThirtySix} & Gaps14321440 & PROVED & \checkmark & verified \\
\texttt{singular\_series\_pos\_evenPair\_oneThousandFourHundredThirtyTwo} & Gaps14321440 & PROVED & \checkmark & verified \\
\texttt{evenPair\_card\_oneThousandFourHundredFifty} & Gaps14421450 & PROVED & \checkmark & verified \\
\texttt{evenPair\_card\_oneThousandFourHundredFortyEight} & Gaps14421450 & PROVED & \checkmark & verified \\
\texttt{evenPair\_card\_oneThousandFourHundredFortyFour} & Gaps14421450 & PROVED & \checkmark & verified \\
\texttt{evenPair\_card\_oneThousandFourHundredFortySix} & Gaps14421450 & PROVED & \checkmark & verified \\
\texttt{evenPair\_card\_oneThousandFourHundredFortyTwo} & Gaps14421450 & PROVED & \checkmark & verified \\
\texttt{isAdmissible\_evenPair\_oneThousandFourHundredFifty} & Gaps14421450 & PROVED & \checkmark & verified \\
\texttt{isAdmissible\_evenPair\_oneThousandFourHundredFortyEight} & Gaps14421450 & PROVED & \checkmark & verified \\
\texttt{isAdmissible\_evenPair\_oneThousandFourHundredFortyFour} & Gaps14421450 & PROVED & \checkmark & verified \\
\texttt{isAdmissible\_evenPair\_oneThousandFourHundredFortySix} & Gaps14421450 & PROVED & \checkmark & verified \\
\texttt{isAdmissible\_evenPair\_oneThousandFourHundredFortyTwo} & Gaps14421450 & PROVED & \checkmark & verified \\
\texttt{localFactor\_oneThousandFourHundredFifty\_two} & Gaps14421450 & PROVED & \checkmark & verified \\
\texttt{localFactor\_oneThousandFourHundredFortyTwo\_two} & Gaps14421450 & PROVED & \checkmark & verified \\
\texttt{nu\_p\_oneThousandFourHundredFifty} & Gaps14421450 & PROVED & \checkmark & verified \\
\texttt{nu\_p\_oneThousandFourHundredFifty\_two} & Gaps14421450 & PROVED & \checkmark & verified \\
\texttt{nu\_p\_oneThousandFourHundredFortyEight} & Gaps14421450 & PROVED & \checkmark & verified \\
\texttt{nu\_p\_oneThousandFourHundredFortyFour} & Gaps14421450 & PROVED & \checkmark & verified \\
\texttt{nu\_p\_oneThousandFourHundredFortySix} & Gaps14421450 & PROVED & \checkmark & verified \\
\texttt{nu\_p\_oneThousandFourHundredFortyTwo} & Gaps14421450 & PROVED & \checkmark & verified \\
\texttt{nu\_p\_oneThousandFourHundredFortyTwo\_two} & Gaps14421450 & PROVED & \checkmark & verified \\
\texttt{singular\_series\_finite\_pos\_evenPair\_oneThousandFourHundredFifty} & Gaps14421450 & PROVED & \checkmark & verified \\
\texttt{singular\_series\_finite\_pos\_evenPair\_oneThousandFourHundredFortyEight} & Gaps14421450 & PROVED & \checkmark & verified \\
\texttt{singular\_series\_finite\_pos\_evenPair\_oneThousandFourHundredFortyFour} & Gaps14421450 & PROVED & \checkmark & verified \\
\texttt{singular\_series\_finite\_pos\_evenPair\_oneThousandFourHundredFortySix} & Gaps14421450 & PROVED & \checkmark & verified \\
\texttt{singular\_series\_finite\_pos\_evenPair\_oneThousandFourHundredFortyTwo} & Gaps14421450 & PROVED & \checkmark & verified \\
\texttt{singular\_series\_pos\_evenPair\_oneThousandFourHundredFifty} & Gaps14421450 & PROVED & \checkmark & verified \\
\texttt{singular\_series\_pos\_evenPair\_oneThousandFourHundredFortyEight} & Gaps14421450 & PROVED & \checkmark & verified \\
\texttt{singular\_series\_pos\_evenPair\_oneThousandFourHundredFortyFour} & Gaps14421450 & PROVED & \checkmark & verified \\
\texttt{singular\_series\_pos\_evenPair\_oneThousandFourHundredFortySix} & Gaps14421450 & PROVED & \checkmark & verified \\
\texttt{singular\_series\_pos\_evenPair\_oneThousandFourHundredFortyTwo} & Gaps14421450 & PROVED & \checkmark & verified \\
\texttt{evenPair\_card\_oneThousandFourHundredFiftyEight} & Gaps14521460 & PROVED & \checkmark & verified \\
\texttt{evenPair\_card\_oneThousandFourHundredFiftyFour} & Gaps14521460 & PROVED & \checkmark & verified \\
\texttt{evenPair\_card\_oneThousandFourHundredFiftySix} & Gaps14521460 & PROVED & \checkmark & verified \\
\texttt{evenPair\_card\_oneThousandFourHundredFiftyTwo} & Gaps14521460 & PROVED & \checkmark & verified \\
\texttt{evenPair\_card\_oneThousandFourHundredSixty} & Gaps14521460 & PROVED & \checkmark & verified \\
\texttt{isAdmissible\_evenPair\_oneThousandFourHundredFiftyEight} & Gaps14521460 & PROVED & \checkmark & verified \\
\texttt{isAdmissible\_evenPair\_oneThousandFourHundredFiftyFour} & Gaps14521460 & PROVED & \checkmark & verified \\
\texttt{isAdmissible\_evenPair\_oneThousandFourHundredFiftySix} & Gaps14521460 & PROVED & \checkmark & verified \\
\texttt{isAdmissible\_evenPair\_oneThousandFourHundredFiftyTwo} & Gaps14521460 & PROVED & \checkmark & verified \\
\texttt{isAdmissible\_evenPair\_oneThousandFourHundredSixty} & Gaps14521460 & PROVED & \checkmark & verified \\
\texttt{localFactor\_oneThousandFourHundredFiftyTwo\_two} & Gaps14521460 & PROVED & \checkmark & verified \\
\texttt{localFactor\_oneThousandFourHundredSixty\_two} & Gaps14521460 & PROVED & \checkmark & verified \\
\texttt{nu\_p\_oneThousandFourHundredFiftyEight} & Gaps14521460 & PROVED & \checkmark & verified \\
\texttt{nu\_p\_oneThousandFourHundredFiftyFour} & Gaps14521460 & PROVED & \checkmark & verified \\
\texttt{nu\_p\_oneThousandFourHundredFiftySix} & Gaps14521460 & PROVED & \checkmark & verified \\
\texttt{nu\_p\_oneThousandFourHundredFiftyTwo} & Gaps14521460 & PROVED & \checkmark & verified \\
\texttt{nu\_p\_oneThousandFourHundredFiftyTwo\_two} & Gaps14521460 & PROVED & \checkmark & verified \\
\texttt{nu\_p\_oneThousandFourHundredSixty} & Gaps14521460 & PROVED & \checkmark & verified \\
\texttt{nu\_p\_oneThousandFourHundredSixty\_two} & Gaps14521460 & PROVED & \checkmark & verified \\
\texttt{singular\_series\_finite\_pos\_evenPair\_oneThousandFourHundredFiftyEight} & Gaps14521460 & PROVED & \checkmark & verified \\
\texttt{singular\_series\_finite\_pos\_evenPair\_oneThousandFourHundredFiftyFour} & Gaps14521460 & PROVED & \checkmark & verified \\
\texttt{singular\_series\_finite\_pos\_evenPair\_oneThousandFourHundredFiftySix} & Gaps14521460 & PROVED & \checkmark & verified \\
\texttt{singular\_series\_finite\_pos\_evenPair\_oneThousandFourHundredFiftyTwo} & Gaps14521460 & PROVED & \checkmark & verified \\
\texttt{singular\_series\_finite\_pos\_evenPair\_oneThousandFourHundredSixty} & Gaps14521460 & PROVED & \checkmark & verified \\
\texttt{singular\_series\_pos\_evenPair\_oneThousandFourHundredFiftyEight} & Gaps14521460 & PROVED & \checkmark & verified \\
\texttt{singular\_series\_pos\_evenPair\_oneThousandFourHundredFiftyFour} & Gaps14521460 & PROVED & \checkmark & verified \\
\texttt{singular\_series\_pos\_evenPair\_oneThousandFourHundredFiftySix} & Gaps14521460 & PROVED & \checkmark & verified \\
\texttt{singular\_series\_pos\_evenPair\_oneThousandFourHundredFiftyTwo} & Gaps14521460 & PROVED & \checkmark & verified \\
\texttt{singular\_series\_pos\_evenPair\_oneThousandFourHundredSixty} & Gaps14521460 & PROVED & \checkmark & verified \\
\texttt{evenPair\_card\_oneThousandFourHundredSeventy} & Gaps14621470 & PROVED & \checkmark & verified \\
\texttt{evenPair\_card\_oneThousandFourHundredSixtyEight} & Gaps14621470 & PROVED & \checkmark & verified \\
\texttt{evenPair\_card\_oneThousandFourHundredSixtyFour} & Gaps14621470 & PROVED & \checkmark & verified \\
\texttt{evenPair\_card\_oneThousandFourHundredSixtySix} & Gaps14621470 & PROVED & \checkmark & verified \\
\texttt{evenPair\_card\_oneThousandFourHundredSixtyTwo} & Gaps14621470 & PROVED & \checkmark & verified \\
\texttt{isAdmissible\_evenPair\_oneThousandFourHundredSeventy} & Gaps14621470 & PROVED & \checkmark & verified \\
\texttt{isAdmissible\_evenPair\_oneThousandFourHundredSixtyEight} & Gaps14621470 & PROVED & \checkmark & verified \\
\texttt{isAdmissible\_evenPair\_oneThousandFourHundredSixtyFour} & Gaps14621470 & PROVED & \checkmark & verified \\
\texttt{isAdmissible\_evenPair\_oneThousandFourHundredSixtySix} & Gaps14621470 & PROVED & \checkmark & verified \\
\texttt{isAdmissible\_evenPair\_oneThousandFourHundredSixtyTwo} & Gaps14621470 & PROVED & \checkmark & verified \\
\texttt{localFactor\_oneThousandFourHundredSeventy\_two} & Gaps14621470 & PROVED & \checkmark & verified \\
\texttt{localFactor\_oneThousandFourHundredSixtyTwo\_two} & Gaps14621470 & PROVED & \checkmark & verified \\
\texttt{nu\_p\_oneThousandFourHundredSeventy} & Gaps14621470 & PROVED & \checkmark & verified \\
\texttt{nu\_p\_oneThousandFourHundredSeventy\_two} & Gaps14621470 & PROVED & \checkmark & verified \\
\texttt{nu\_p\_oneThousandFourHundredSixtyEight} & Gaps14621470 & PROVED & \checkmark & verified \\
\texttt{nu\_p\_oneThousandFourHundredSixtyFour} & Gaps14621470 & PROVED & \checkmark & verified \\
\texttt{nu\_p\_oneThousandFourHundredSixtySix} & Gaps14621470 & PROVED & \checkmark & verified \\
\texttt{nu\_p\_oneThousandFourHundredSixtyTwo} & Gaps14621470 & PROVED & \checkmark & verified \\
\texttt{nu\_p\_oneThousandFourHundredSixtyTwo\_two} & Gaps14621470 & PROVED & \checkmark & verified \\
\texttt{singular\_series\_finite\_pos\_evenPair\_oneThousandFourHundredSeventy} & Gaps14621470 & PROVED & \checkmark & verified \\
\texttt{singular\_series\_finite\_pos\_evenPair\_oneThousandFourHundredSixtyEight} & Gaps14621470 & PROVED & \checkmark & verified \\
\texttt{singular\_series\_finite\_pos\_evenPair\_oneThousandFourHundredSixtyFour} & Gaps14621470 & PROVED & \checkmark & verified \\
\texttt{singular\_series\_finite\_pos\_evenPair\_oneThousandFourHundredSixtySix} & Gaps14621470 & PROVED & \checkmark & verified \\
\texttt{singular\_series\_finite\_pos\_evenPair\_oneThousandFourHundredSixtyTwo} & Gaps14621470 & PROVED & \checkmark & verified \\
\texttt{singular\_series\_pos\_evenPair\_oneThousandFourHundredSeventy} & Gaps14621470 & PROVED & \checkmark & verified \\
\texttt{singular\_series\_pos\_evenPair\_oneThousandFourHundredSixtyEight} & Gaps14621470 & PROVED & \checkmark & verified \\
\texttt{singular\_series\_pos\_evenPair\_oneThousandFourHundredSixtyFour} & Gaps14621470 & PROVED & \checkmark & verified \\
\texttt{singular\_series\_pos\_evenPair\_oneThousandFourHundredSixtySix} & Gaps14621470 & PROVED & \checkmark & verified \\
\texttt{singular\_series\_pos\_evenPair\_oneThousandFourHundredSixtyTwo} & Gaps14621470 & PROVED & \checkmark & verified \\
\texttt{evenPair\_card\_oneThousandFourHundredEighty} & Gaps14721480 & PROVED & \checkmark & verified \\
\texttt{evenPair\_card\_oneThousandFourHundredSeventyEight} & Gaps14721480 & PROVED & \checkmark & verified \\
\texttt{evenPair\_card\_oneThousandFourHundredSeventyFour} & Gaps14721480 & PROVED & \checkmark & verified \\
\texttt{evenPair\_card\_oneThousandFourHundredSeventySix} & Gaps14721480 & PROVED & \checkmark & verified \\
\texttt{evenPair\_card\_oneThousandFourHundredSeventyTwo} & Gaps14721480 & PROVED & \checkmark & verified \\
\texttt{isAdmissible\_evenPair\_oneThousandFourHundredEighty} & Gaps14721480 & PROVED & \checkmark & verified \\
\texttt{isAdmissible\_evenPair\_oneThousandFourHundredSeventyEight} & Gaps14721480 & PROVED & \checkmark & verified \\
\texttt{isAdmissible\_evenPair\_oneThousandFourHundredSeventyFour} & Gaps14721480 & PROVED & \checkmark & verified \\
\texttt{isAdmissible\_evenPair\_oneThousandFourHundredSeventySix} & Gaps14721480 & PROVED & \checkmark & verified \\
\texttt{isAdmissible\_evenPair\_oneThousandFourHundredSeventyTwo} & Gaps14721480 & PROVED & \checkmark & verified \\
\texttt{localFactor\_oneThousandFourHundredEighty\_two} & Gaps14721480 & PROVED & \checkmark & verified \\
\texttt{localFactor\_oneThousandFourHundredSeventyTwo\_two} & Gaps14721480 & PROVED & \checkmark & verified \\
\texttt{nu\_p\_oneThousandFourHundredEighty} & Gaps14721480 & PROVED & \checkmark & verified \\
\texttt{nu\_p\_oneThousandFourHundredEighty\_two} & Gaps14721480 & PROVED & \checkmark & verified \\
\texttt{nu\_p\_oneThousandFourHundredSeventyEight} & Gaps14721480 & PROVED & \checkmark & verified \\
\texttt{nu\_p\_oneThousandFourHundredSeventyFour} & Gaps14721480 & PROVED & \checkmark & verified \\
\texttt{nu\_p\_oneThousandFourHundredSeventySix} & Gaps14721480 & PROVED & \checkmark & verified \\
\texttt{nu\_p\_oneThousandFourHundredSeventyTwo} & Gaps14721480 & PROVED & \checkmark & verified \\
\texttt{nu\_p\_oneThousandFourHundredSeventyTwo\_two} & Gaps14721480 & PROVED & \checkmark & verified \\
\texttt{singular\_series\_finite\_pos\_evenPair\_oneThousandFourHundredEighty} & Gaps14721480 & PROVED & \checkmark & verified \\
\texttt{singular\_series\_finite\_pos\_evenPair\_oneThousandFourHundredSeventyEight} & Gaps14721480 & PROVED & \checkmark & verified \\
\texttt{singular\_series\_finite\_pos\_evenPair\_oneThousandFourHundredSeventyFour} & Gaps14721480 & PROVED & \checkmark & verified \\
\texttt{singular\_series\_finite\_pos\_evenPair\_oneThousandFourHundredSeventySix} & Gaps14721480 & PROVED & \checkmark & verified \\
\texttt{singular\_series\_finite\_pos\_evenPair\_oneThousandFourHundredSeventyTwo} & Gaps14721480 & PROVED & \checkmark & verified \\
\texttt{singular\_series\_pos\_evenPair\_oneThousandFourHundredEighty} & Gaps14721480 & PROVED & \checkmark & verified \\
\texttt{singular\_series\_pos\_evenPair\_oneThousandFourHundredSeventyEight} & Gaps14721480 & PROVED & \checkmark & verified \\
\texttt{singular\_series\_pos\_evenPair\_oneThousandFourHundredSeventyFour} & Gaps14721480 & PROVED & \checkmark & verified \\
\texttt{singular\_series\_pos\_evenPair\_oneThousandFourHundredSeventySix} & Gaps14721480 & PROVED & \checkmark & verified \\
\texttt{singular\_series\_pos\_evenPair\_oneThousandFourHundredSeventyTwo} & Gaps14721480 & PROVED & \checkmark & verified \\
\texttt{evenPair\_card\_oneThousandFourHundredEightyEight} & Gaps14821490 & PROVED & \checkmark & verified \\
\texttt{evenPair\_card\_oneThousandFourHundredEightyFour} & Gaps14821490 & PROVED & \checkmark & verified \\
\texttt{evenPair\_card\_oneThousandFourHundredEightySix} & Gaps14821490 & PROVED & \checkmark & verified \\
\texttt{evenPair\_card\_oneThousandFourHundredEightyTwo} & Gaps14821490 & PROVED & \checkmark & verified \\
\texttt{evenPair\_card\_oneThousandFourHundredNinety} & Gaps14821490 & PROVED & \checkmark & verified \\
\texttt{isAdmissible\_evenPair\_oneThousandFourHundredEightyEight} & Gaps14821490 & PROVED & \checkmark & verified \\
\texttt{isAdmissible\_evenPair\_oneThousandFourHundredEightyFour} & Gaps14821490 & PROVED & \checkmark & verified \\
\texttt{isAdmissible\_evenPair\_oneThousandFourHundredEightySix} & Gaps14821490 & PROVED & \checkmark & verified \\
\texttt{isAdmissible\_evenPair\_oneThousandFourHundredEightyTwo} & Gaps14821490 & PROVED & \checkmark & verified \\
\texttt{isAdmissible\_evenPair\_oneThousandFourHundredNinety} & Gaps14821490 & PROVED & \checkmark & verified \\
\texttt{localFactor\_oneThousandFourHundredEightyTwo\_two} & Gaps14821490 & PROVED & \checkmark & verified \\
\texttt{localFactor\_oneThousandFourHundredNinety\_two} & Gaps14821490 & PROVED & \checkmark & verified \\
\texttt{nu\_p\_oneThousandFourHundredEightyEight} & Gaps14821490 & PROVED & \checkmark & verified \\
\texttt{nu\_p\_oneThousandFourHundredEightyFour} & Gaps14821490 & PROVED & \checkmark & verified \\
\texttt{nu\_p\_oneThousandFourHundredEightySix} & Gaps14821490 & PROVED & \checkmark & verified \\
\texttt{nu\_p\_oneThousandFourHundredEightyTwo} & Gaps14821490 & PROVED & \checkmark & verified \\
\texttt{nu\_p\_oneThousandFourHundredEightyTwo\_two} & Gaps14821490 & PROVED & \checkmark & verified \\
\texttt{nu\_p\_oneThousandFourHundredNinety} & Gaps14821490 & PROVED & \checkmark & verified \\
\texttt{nu\_p\_oneThousandFourHundredNinety\_two} & Gaps14821490 & PROVED & \checkmark & verified \\
\texttt{singular\_series\_finite\_pos\_evenPair\_oneThousandFourHundredEightyEight} & Gaps14821490 & PROVED & \checkmark & verified \\
\texttt{singular\_series\_finite\_pos\_evenPair\_oneThousandFourHundredEightyFour} & Gaps14821490 & PROVED & \checkmark & verified \\
\texttt{singular\_series\_finite\_pos\_evenPair\_oneThousandFourHundredEightySix} & Gaps14821490 & PROVED & \checkmark & verified \\
\texttt{singular\_series\_finite\_pos\_evenPair\_oneThousandFourHundredEightyTwo} & Gaps14821490 & PROVED & \checkmark & verified \\
\texttt{singular\_series\_finite\_pos\_evenPair\_oneThousandFourHundredNinety} & Gaps14821490 & PROVED & \checkmark & verified \\
\texttt{singular\_series\_pos\_evenPair\_oneThousandFourHundredEightyEight} & Gaps14821490 & PROVED & \checkmark & verified \\
\texttt{singular\_series\_pos\_evenPair\_oneThousandFourHundredEightyFour} & Gaps14821490 & PROVED & \checkmark & verified \\
\texttt{singular\_series\_pos\_evenPair\_oneThousandFourHundredEightySix} & Gaps14821490 & PROVED & \checkmark & verified \\
\texttt{singular\_series\_pos\_evenPair\_oneThousandFourHundredEightyTwo} & Gaps14821490 & PROVED & \checkmark & verified \\
\texttt{singular\_series\_pos\_evenPair\_oneThousandFourHundredNinety} & Gaps14821490 & PROVED & \checkmark & verified \\
\texttt{evenPair\_card\_oneThousandFiveHundred} & Gaps14921500 & PROVED & \checkmark & verified \\
\texttt{evenPair\_card\_oneThousandFourHundredNinetyEight} & Gaps14921500 & PROVED & \checkmark & verified \\
\texttt{evenPair\_card\_oneThousandFourHundredNinetyFour} & Gaps14921500 & PROVED & \checkmark & verified \\
\texttt{evenPair\_card\_oneThousandFourHundredNinetySix} & Gaps14921500 & PROVED & \checkmark & verified \\
\texttt{evenPair\_card\_oneThousandFourHundredNinetyTwo} & Gaps14921500 & PROVED & \checkmark & verified \\
\texttt{isAdmissible\_evenPair\_oneThousandFiveHundred} & Gaps14921500 & PROVED & \checkmark & verified \\
\texttt{isAdmissible\_evenPair\_oneThousandFourHundredNinetyEight} & Gaps14921500 & PROVED & \checkmark & verified \\
\texttt{isAdmissible\_evenPair\_oneThousandFourHundredNinetyFour} & Gaps14921500 & PROVED & \checkmark & verified \\
\texttt{isAdmissible\_evenPair\_oneThousandFourHundredNinetySix} & Gaps14921500 & PROVED & \checkmark & verified \\
\texttt{isAdmissible\_evenPair\_oneThousandFourHundredNinetyTwo} & Gaps14921500 & PROVED & \checkmark & verified \\
\texttt{localFactor\_oneThousandFiveHundred\_two} & Gaps14921500 & PROVED & \checkmark & verified \\
\texttt{localFactor\_oneThousandFourHundredNinetyTwo\_two} & Gaps14921500 & PROVED & \checkmark & verified \\
\texttt{nu\_p\_oneThousandFiveHundred} & Gaps14921500 & PROVED & \checkmark & verified \\
\texttt{nu\_p\_oneThousandFiveHundred\_two} & Gaps14921500 & PROVED & \checkmark & verified \\
\texttt{nu\_p\_oneThousandFourHundredNinetyEight} & Gaps14921500 & PROVED & \checkmark & verified \\
\texttt{nu\_p\_oneThousandFourHundredNinetyFour} & Gaps14921500 & PROVED & \checkmark & verified \\
\texttt{nu\_p\_oneThousandFourHundredNinetySix} & Gaps14921500 & PROVED & \checkmark & verified \\
\texttt{nu\_p\_oneThousandFourHundredNinetyTwo} & Gaps14921500 & PROVED & \checkmark & verified \\
\texttt{nu\_p\_oneThousandFourHundredNinetyTwo\_two} & Gaps14921500 & PROVED & \checkmark & verified \\
\texttt{singular\_series\_finite\_pos\_evenPair\_oneThousandFiveHundred} & Gaps14921500 & PROVED & \checkmark & verified \\
\texttt{singular\_series\_finite\_pos\_evenPair\_oneThousandFourHundredNinetyEight} & Gaps14921500 & PROVED & \checkmark & verified \\
\texttt{singular\_series\_finite\_pos\_evenPair\_oneThousandFourHundredNinetyFour} & Gaps14921500 & PROVED & \checkmark & verified \\
\texttt{singular\_series\_finite\_pos\_evenPair\_oneThousandFourHundredNinetySix} & Gaps14921500 & PROVED & \checkmark & verified \\
\texttt{singular\_series\_finite\_pos\_evenPair\_oneThousandFourHundredNinetyTwo} & Gaps14921500 & PROVED & \checkmark & verified \\
\texttt{singular\_series\_pos\_evenPair\_oneThousandFiveHundred} & Gaps14921500 & PROVED & \checkmark & verified \\
\texttt{singular\_series\_pos\_evenPair\_oneThousandFourHundredNinetyEight} & Gaps14921500 & PROVED & \checkmark & verified \\
\texttt{singular\_series\_pos\_evenPair\_oneThousandFourHundredNinetyFour} & Gaps14921500 & PROVED & \checkmark & verified \\
\texttt{singular\_series\_pos\_evenPair\_oneThousandFourHundredNinetySix} & Gaps14921500 & PROVED & \checkmark & verified \\
\texttt{singular\_series\_pos\_evenPair\_oneThousandFourHundredNinetyTwo} & Gaps14921500 & PROVED & \checkmark & verified \\
\texttt{evenPair\_card\_oneThousandFiveHundredEight} & Gaps15021510 & PROVED & \checkmark & verified \\
\texttt{evenPair\_card\_oneThousandFiveHundredFour} & Gaps15021510 & PROVED & \checkmark & verified \\
\texttt{evenPair\_card\_oneThousandFiveHundredSix} & Gaps15021510 & PROVED & \checkmark & verified \\
\texttt{evenPair\_card\_oneThousandFiveHundredTen} & Gaps15021510 & PROVED & \checkmark & verified \\
\texttt{evenPair\_card\_oneThousandFiveHundredTwo} & Gaps15021510 & PROVED & \checkmark & verified \\
\texttt{isAdmissible\_evenPair\_oneThousandFiveHundredEight} & Gaps15021510 & PROVED & \checkmark & verified \\
\texttt{isAdmissible\_evenPair\_oneThousandFiveHundredFour} & Gaps15021510 & PROVED & \checkmark & verified \\
\texttt{isAdmissible\_evenPair\_oneThousandFiveHundredSix} & Gaps15021510 & PROVED & \checkmark & verified \\
\texttt{isAdmissible\_evenPair\_oneThousandFiveHundredTen} & Gaps15021510 & PROVED & \checkmark & verified \\
\texttt{isAdmissible\_evenPair\_oneThousandFiveHundredTwo} & Gaps15021510 & PROVED & \checkmark & verified \\
\texttt{localFactor\_oneThousandFiveHundredTen\_two} & Gaps15021510 & PROVED & \checkmark & verified \\
\texttt{localFactor\_oneThousandFiveHundredTwo\_two} & Gaps15021510 & PROVED & \checkmark & verified \\
\texttt{nu\_p\_oneThousandFiveHundredEight} & Gaps15021510 & PROVED & \checkmark & verified \\
\texttt{nu\_p\_oneThousandFiveHundredFour} & Gaps15021510 & PROVED & \checkmark & verified \\
\texttt{nu\_p\_oneThousandFiveHundredSix} & Gaps15021510 & PROVED & \checkmark & verified \\
\texttt{nu\_p\_oneThousandFiveHundredTen} & Gaps15021510 & PROVED & \checkmark & verified \\
\texttt{nu\_p\_oneThousandFiveHundredTen\_two} & Gaps15021510 & PROVED & \checkmark & verified \\
\texttt{nu\_p\_oneThousandFiveHundredTwo} & Gaps15021510 & PROVED & \checkmark & verified \\
\texttt{nu\_p\_oneThousandFiveHundredTwo\_two} & Gaps15021510 & PROVED & \checkmark & verified \\
\texttt{singular\_series\_finite\_pos\_evenPair\_oneThousandFiveHundredEight} & Gaps15021510 & PROVED & \checkmark & verified \\
\texttt{singular\_series\_finite\_pos\_evenPair\_oneThousandFiveHundredFour} & Gaps15021510 & PROVED & \checkmark & verified \\
\texttt{singular\_series\_finite\_pos\_evenPair\_oneThousandFiveHundredSix} & Gaps15021510 & PROVED & \checkmark & verified \\
\texttt{singular\_series\_finite\_pos\_evenPair\_oneThousandFiveHundredTen} & Gaps15021510 & PROVED & \checkmark & verified \\
\texttt{singular\_series\_finite\_pos\_evenPair\_oneThousandFiveHundredTwo} & Gaps15021510 & PROVED & \checkmark & verified \\
\texttt{singular\_series\_pos\_evenPair\_oneThousandFiveHundredEight} & Gaps15021510 & PROVED & \checkmark & verified \\
\texttt{singular\_series\_pos\_evenPair\_oneThousandFiveHundredFour} & Gaps15021510 & PROVED & \checkmark & verified \\
\texttt{singular\_series\_pos\_evenPair\_oneThousandFiveHundredSix} & Gaps15021510 & PROVED & \checkmark & verified \\
\texttt{singular\_series\_pos\_evenPair\_oneThousandFiveHundredTen} & Gaps15021510 & PROVED & \checkmark & verified \\
\texttt{singular\_series\_pos\_evenPair\_oneThousandFiveHundredTwo} & Gaps15021510 & PROVED & \checkmark & verified \\
\texttt{evenPair\_card\_oneThousandFiveHundredEighteen} & Gaps15121520 & PROVED & \checkmark & verified \\
\texttt{evenPair\_card\_oneThousandFiveHundredFourteen} & Gaps15121520 & PROVED & \checkmark & verified \\
\texttt{evenPair\_card\_oneThousandFiveHundredSixteen} & Gaps15121520 & PROVED & \checkmark & verified \\
\texttt{evenPair\_card\_oneThousandFiveHundredTwelve} & Gaps15121520 & PROVED & \checkmark & verified \\
\texttt{evenPair\_card\_oneThousandFiveHundredTwenty} & Gaps15121520 & PROVED & \checkmark & verified \\
\texttt{isAdmissible\_evenPair\_oneThousandFiveHundredEighteen} & Gaps15121520 & PROVED & \checkmark & verified \\
\texttt{isAdmissible\_evenPair\_oneThousandFiveHundredFourteen} & Gaps15121520 & PROVED & \checkmark & verified \\
\texttt{isAdmissible\_evenPair\_oneThousandFiveHundredSixteen} & Gaps15121520 & PROVED & \checkmark & verified \\
\texttt{isAdmissible\_evenPair\_oneThousandFiveHundredTwelve} & Gaps15121520 & PROVED & \checkmark & verified \\
\texttt{isAdmissible\_evenPair\_oneThousandFiveHundredTwenty} & Gaps15121520 & PROVED & \checkmark & verified \\
\texttt{localFactor\_oneThousandFiveHundredTwelve\_two} & Gaps15121520 & PROVED & \checkmark & verified \\
\texttt{localFactor\_oneThousandFiveHundredTwenty\_two} & Gaps15121520 & PROVED & \checkmark & verified \\
\texttt{nu\_p\_oneThousandFiveHundredEighteen} & Gaps15121520 & PROVED & \checkmark & verified \\
\texttt{nu\_p\_oneThousandFiveHundredFourteen} & Gaps15121520 & PROVED & \checkmark & verified \\
\texttt{nu\_p\_oneThousandFiveHundredSixteen} & Gaps15121520 & PROVED & \checkmark & verified \\
\texttt{nu\_p\_oneThousandFiveHundredTwelve} & Gaps15121520 & PROVED & \checkmark & verified \\
\texttt{nu\_p\_oneThousandFiveHundredTwelve\_two} & Gaps15121520 & PROVED & \checkmark & verified \\
\texttt{nu\_p\_oneThousandFiveHundredTwenty} & Gaps15121520 & PROVED & \checkmark & verified \\
\texttt{nu\_p\_oneThousandFiveHundredTwenty\_two} & Gaps15121520 & PROVED & \checkmark & verified \\
\texttt{singular\_series\_finite\_pos\_evenPair\_oneThousandFiveHundredEighteen} & Gaps15121520 & PROVED & \checkmark & verified \\
\texttt{singular\_series\_finite\_pos\_evenPair\_oneThousandFiveHundredFourteen} & Gaps15121520 & PROVED & \checkmark & verified \\
\texttt{singular\_series\_finite\_pos\_evenPair\_oneThousandFiveHundredSixteen} & Gaps15121520 & PROVED & \checkmark & verified \\
\texttt{singular\_series\_finite\_pos\_evenPair\_oneThousandFiveHundredTwelve} & Gaps15121520 & PROVED & \checkmark & verified \\
\texttt{singular\_series\_finite\_pos\_evenPair\_oneThousandFiveHundredTwenty} & Gaps15121520 & PROVED & \checkmark & verified \\
\texttt{singular\_series\_pos\_evenPair\_oneThousandFiveHundredEighteen} & Gaps15121520 & PROVED & \checkmark & verified \\
\texttt{singular\_series\_pos\_evenPair\_oneThousandFiveHundredFourteen} & Gaps15121520 & PROVED & \checkmark & verified \\
\texttt{singular\_series\_pos\_evenPair\_oneThousandFiveHundredSixteen} & Gaps15121520 & PROVED & \checkmark & verified \\
\texttt{singular\_series\_pos\_evenPair\_oneThousandFiveHundredTwelve} & Gaps15121520 & PROVED & \checkmark & verified \\
\texttt{singular\_series\_pos\_evenPair\_oneThousandFiveHundredTwenty} & Gaps15121520 & PROVED & \checkmark & verified \\
\texttt{evenPair\_card\_oneHundredFiftyEight} & Gaps152160 & PROVED & \checkmark & verified \\
\texttt{evenPair\_card\_oneHundredFiftyFour} & Gaps152160 & PROVED & \checkmark & verified \\
\texttt{evenPair\_card\_oneHundredFiftySix} & Gaps152160 & PROVED & \checkmark & verified \\
\texttt{evenPair\_card\_oneHundredFiftyTwo} & Gaps152160 & PROVED & \checkmark & verified \\
\texttt{evenPair\_card\_oneHundredSixty} & Gaps152160 & PROVED & \checkmark & verified \\
\texttt{isAdmissible\_evenPair\_oneHundredFiftyEight} & Gaps152160 & PROVED & \checkmark & verified \\
\texttt{isAdmissible\_evenPair\_oneHundredFiftyFour} & Gaps152160 & PROVED & \checkmark & verified \\
\texttt{isAdmissible\_evenPair\_oneHundredFiftySix} & Gaps152160 & PROVED & \checkmark & verified \\
\texttt{isAdmissible\_evenPair\_oneHundredFiftyTwo} & Gaps152160 & PROVED & \checkmark & verified \\
\texttt{isAdmissible\_evenPair\_oneHundredSixty} & Gaps152160 & PROVED & \checkmark & verified \\
\texttt{localFactor\_oneHundredFiftyTwo\_two} & Gaps152160 & PROVED & \checkmark & verified \\
\texttt{localFactor\_oneHundredSixty\_two} & Gaps152160 & PROVED & \checkmark & verified \\
\texttt{nu\_p\_oneHundredFiftyEight} & Gaps152160 & PROVED & \checkmark & verified \\
\texttt{nu\_p\_oneHundredFiftyFour} & Gaps152160 & PROVED & \checkmark & verified \\
\texttt{nu\_p\_oneHundredFiftySix} & Gaps152160 & PROVED & \checkmark & verified \\
\texttt{nu\_p\_oneHundredFiftyTwo} & Gaps152160 & PROVED & \checkmark & verified \\
\texttt{nu\_p\_oneHundredFiftyTwo\_two} & Gaps152160 & PROVED & \checkmark & verified \\
\texttt{nu\_p\_oneHundredSixty} & Gaps152160 & PROVED & \checkmark & verified \\
\texttt{nu\_p\_oneHundredSixty\_two} & Gaps152160 & PROVED & \checkmark & verified \\
\texttt{singular\_series\_finite\_pos\_evenPair\_oneHundredFiftyEight} & Gaps152160 & PROVED & \checkmark & verified \\
\texttt{singular\_series\_finite\_pos\_evenPair\_oneHundredFiftyFour} & Gaps152160 & PROVED & \checkmark & verified \\
\texttt{singular\_series\_finite\_pos\_evenPair\_oneHundredFiftySix} & Gaps152160 & PROVED & \checkmark & verified \\
\texttt{singular\_series\_finite\_pos\_evenPair\_oneHundredFiftyTwo} & Gaps152160 & PROVED & \checkmark & verified \\
\texttt{singular\_series\_finite\_pos\_evenPair\_oneHundredSixty} & Gaps152160 & PROVED & \checkmark & verified \\
\texttt{singular\_series\_pos\_evenPair\_oneHundredFiftyEight} & Gaps152160 & PROVED & \checkmark & verified \\
\texttt{singular\_series\_pos\_evenPair\_oneHundredFiftyFour} & Gaps152160 & PROVED & \checkmark & verified \\
\texttt{singular\_series\_pos\_evenPair\_oneHundredFiftySix} & Gaps152160 & PROVED & \checkmark & verified \\
\texttt{singular\_series\_pos\_evenPair\_oneHundredFiftyTwo} & Gaps152160 & PROVED & \checkmark & verified \\
\texttt{singular\_series\_pos\_evenPair\_oneHundredSixty} & Gaps152160 & PROVED & \checkmark & verified \\
\texttt{evenPair\_card\_oneThousandFiveHundredThirty} & Gaps15221530 & PROVED & \checkmark & verified \\
\texttt{evenPair\_card\_oneThousandFiveHundredTwentyEight} & Gaps15221530 & PROVED & \checkmark & verified \\
\texttt{evenPair\_card\_oneThousandFiveHundredTwentyFour} & Gaps15221530 & PROVED & \checkmark & verified \\
\texttt{evenPair\_card\_oneThousandFiveHundredTwentySix} & Gaps15221530 & PROVED & \checkmark & verified \\
\texttt{evenPair\_card\_oneThousandFiveHundredTwentyTwo} & Gaps15221530 & PROVED & \checkmark & verified \\
\texttt{isAdmissible\_evenPair\_oneThousandFiveHundredThirty} & Gaps15221530 & PROVED & \checkmark & verified \\
\texttt{isAdmissible\_evenPair\_oneThousandFiveHundredTwentyEight} & Gaps15221530 & PROVED & \checkmark & verified \\
\texttt{isAdmissible\_evenPair\_oneThousandFiveHundredTwentyFour} & Gaps15221530 & PROVED & \checkmark & verified \\
\texttt{isAdmissible\_evenPair\_oneThousandFiveHundredTwentySix} & Gaps15221530 & PROVED & \checkmark & verified \\
\texttt{isAdmissible\_evenPair\_oneThousandFiveHundredTwentyTwo} & Gaps15221530 & PROVED & \checkmark & verified \\
\texttt{localFactor\_oneThousandFiveHundredThirty\_two} & Gaps15221530 & PROVED & \checkmark & verified \\
\texttt{localFactor\_oneThousandFiveHundredTwentyTwo\_two} & Gaps15221530 & PROVED & \checkmark & verified \\
\texttt{nu\_p\_oneThousandFiveHundredThirty} & Gaps15221530 & PROVED & \checkmark & verified \\
\texttt{nu\_p\_oneThousandFiveHundredThirty\_two} & Gaps15221530 & PROVED & \checkmark & verified \\
\texttt{nu\_p\_oneThousandFiveHundredTwentyEight} & Gaps15221530 & PROVED & \checkmark & verified \\
\texttt{nu\_p\_oneThousandFiveHundredTwentyFour} & Gaps15221530 & PROVED & \checkmark & verified \\
\texttt{nu\_p\_oneThousandFiveHundredTwentySix} & Gaps15221530 & PROVED & \checkmark & verified \\
\texttt{nu\_p\_oneThousandFiveHundredTwentyTwo} & Gaps15221530 & PROVED & \checkmark & verified \\
\texttt{nu\_p\_oneThousandFiveHundredTwentyTwo\_two} & Gaps15221530 & PROVED & \checkmark & verified \\
\texttt{singular\_series\_finite\_pos\_evenPair\_oneThousandFiveHundredThirty} & Gaps15221530 & PROVED & \checkmark & verified \\
\texttt{singular\_series\_finite\_pos\_evenPair\_oneThousandFiveHundredTwentyEight} & Gaps15221530 & PROVED & \checkmark & verified \\
\texttt{singular\_series\_finite\_pos\_evenPair\_oneThousandFiveHundredTwentyFour} & Gaps15221530 & PROVED & \checkmark & verified \\
\texttt{singular\_series\_finite\_pos\_evenPair\_oneThousandFiveHundredTwentySix} & Gaps15221530 & PROVED & \checkmark & verified \\
\texttt{singular\_series\_finite\_pos\_evenPair\_oneThousandFiveHundredTwentyTwo} & Gaps15221530 & PROVED & \checkmark & verified \\
\texttt{singular\_series\_pos\_evenPair\_oneThousandFiveHundredThirty} & Gaps15221530 & PROVED & \checkmark & verified \\
\texttt{singular\_series\_pos\_evenPair\_oneThousandFiveHundredTwentyEight} & Gaps15221530 & PROVED & \checkmark & verified \\
\texttt{singular\_series\_pos\_evenPair\_oneThousandFiveHundredTwentyFour} & Gaps15221530 & PROVED & \checkmark & verified \\
\texttt{singular\_series\_pos\_evenPair\_oneThousandFiveHundredTwentySix} & Gaps15221530 & PROVED & \checkmark & verified \\
\texttt{singular\_series\_pos\_evenPair\_oneThousandFiveHundredTwentyTwo} & Gaps15221530 & PROVED & \checkmark & verified \\
\texttt{evenPair\_card\_oneThousandFiveHundredForty} & Gaps15321540 & PROVED & \checkmark & verified \\
\texttt{evenPair\_card\_oneThousandFiveHundredThirtyEight} & Gaps15321540 & PROVED & \checkmark & verified \\
\texttt{evenPair\_card\_oneThousandFiveHundredThirtyFour} & Gaps15321540 & PROVED & \checkmark & verified \\
\texttt{evenPair\_card\_oneThousandFiveHundredThirtySix} & Gaps15321540 & PROVED & \checkmark & verified \\
\texttt{evenPair\_card\_oneThousandFiveHundredThirtyTwo} & Gaps15321540 & PROVED & \checkmark & verified \\
\texttt{isAdmissible\_evenPair\_oneThousandFiveHundredForty} & Gaps15321540 & PROVED & \checkmark & verified \\
\texttt{isAdmissible\_evenPair\_oneThousandFiveHundredThirtyEight} & Gaps15321540 & PROVED & \checkmark & verified \\
\texttt{isAdmissible\_evenPair\_oneThousandFiveHundredThirtyFour} & Gaps15321540 & PROVED & \checkmark & verified \\
\texttt{isAdmissible\_evenPair\_oneThousandFiveHundredThirtySix} & Gaps15321540 & PROVED & \checkmark & verified \\
\texttt{isAdmissible\_evenPair\_oneThousandFiveHundredThirtyTwo} & Gaps15321540 & PROVED & \checkmark & verified \\
\texttt{localFactor\_oneThousandFiveHundredForty\_two} & Gaps15321540 & PROVED & \checkmark & verified \\
\texttt{localFactor\_oneThousandFiveHundredThirtyTwo\_two} & Gaps15321540 & PROVED & \checkmark & verified \\
\texttt{nu\_p\_oneThousandFiveHundredForty} & Gaps15321540 & PROVED & \checkmark & verified \\
\texttt{nu\_p\_oneThousandFiveHundredForty\_two} & Gaps15321540 & PROVED & \checkmark & verified \\
\texttt{nu\_p\_oneThousandFiveHundredThirtyEight} & Gaps15321540 & PROVED & \checkmark & verified \\
\texttt{nu\_p\_oneThousandFiveHundredThirtyFour} & Gaps15321540 & PROVED & \checkmark & verified \\
\texttt{nu\_p\_oneThousandFiveHundredThirtySix} & Gaps15321540 & PROVED & \checkmark & verified \\
\texttt{nu\_p\_oneThousandFiveHundredThirtyTwo} & Gaps15321540 & PROVED & \checkmark & verified \\
\texttt{nu\_p\_oneThousandFiveHundredThirtyTwo\_two} & Gaps15321540 & PROVED & \checkmark & verified \\
\texttt{singular\_series\_finite\_pos\_evenPair\_oneThousandFiveHundredForty} & Gaps15321540 & PROVED & \checkmark & verified \\
\texttt{singular\_series\_finite\_pos\_evenPair\_oneThousandFiveHundredThirtyEight} & Gaps15321540 & PROVED & \checkmark & verified \\
\texttt{singular\_series\_finite\_pos\_evenPair\_oneThousandFiveHundredThirtyFour} & Gaps15321540 & PROVED & \checkmark & verified \\
\texttt{singular\_series\_finite\_pos\_evenPair\_oneThousandFiveHundredThirtySix} & Gaps15321540 & PROVED & \checkmark & verified \\
\texttt{singular\_series\_finite\_pos\_evenPair\_oneThousandFiveHundredThirtyTwo} & Gaps15321540 & PROVED & \checkmark & verified \\
\texttt{singular\_series\_pos\_evenPair\_oneThousandFiveHundredForty} & Gaps15321540 & PROVED & \checkmark & verified \\
\texttt{singular\_series\_pos\_evenPair\_oneThousandFiveHundredThirtyEight} & Gaps15321540 & PROVED & \checkmark & verified \\
\texttt{singular\_series\_pos\_evenPair\_oneThousandFiveHundredThirtyFour} & Gaps15321540 & PROVED & \checkmark & verified \\
\texttt{singular\_series\_pos\_evenPair\_oneThousandFiveHundredThirtySix} & Gaps15321540 & PROVED & \checkmark & verified \\
\texttt{singular\_series\_pos\_evenPair\_oneThousandFiveHundredThirtyTwo} & Gaps15321540 & PROVED & \checkmark & verified \\
\texttt{evenPair\_card\_oneThousandFiveHundredFifty} & Gaps15421550 & PROVED & \checkmark & verified \\
\texttt{evenPair\_card\_oneThousandFiveHundredFortyEight} & Gaps15421550 & PROVED & \checkmark & verified \\
\texttt{evenPair\_card\_oneThousandFiveHundredFortyFour} & Gaps15421550 & PROVED & \checkmark & verified \\
\texttt{evenPair\_card\_oneThousandFiveHundredFortySix} & Gaps15421550 & PROVED & \checkmark & verified \\
\texttt{evenPair\_card\_oneThousandFiveHundredFortyTwo} & Gaps15421550 & PROVED & \checkmark & verified \\
\texttt{isAdmissible\_evenPair\_oneThousandFiveHundredFifty} & Gaps15421550 & PROVED & \checkmark & verified \\
\texttt{isAdmissible\_evenPair\_oneThousandFiveHundredFortyEight} & Gaps15421550 & PROVED & \checkmark & verified \\
\texttt{isAdmissible\_evenPair\_oneThousandFiveHundredFortyFour} & Gaps15421550 & PROVED & \checkmark & verified \\
\texttt{isAdmissible\_evenPair\_oneThousandFiveHundredFortySix} & Gaps15421550 & PROVED & \checkmark & verified \\
\texttt{isAdmissible\_evenPair\_oneThousandFiveHundredFortyTwo} & Gaps15421550 & PROVED & \checkmark & verified \\
\texttt{localFactor\_oneThousandFiveHundredFifty\_two} & Gaps15421550 & PROVED & \checkmark & verified \\
\texttt{localFactor\_oneThousandFiveHundredFortyTwo\_two} & Gaps15421550 & PROVED & \checkmark & verified \\
\texttt{nu\_p\_oneThousandFiveHundredFifty} & Gaps15421550 & PROVED & \checkmark & verified \\
\texttt{nu\_p\_oneThousandFiveHundredFifty\_two} & Gaps15421550 & PROVED & \checkmark & verified \\
\texttt{nu\_p\_oneThousandFiveHundredFortyEight} & Gaps15421550 & PROVED & \checkmark & verified \\
\texttt{nu\_p\_oneThousandFiveHundredFortyFour} & Gaps15421550 & PROVED & \checkmark & verified \\
\texttt{nu\_p\_oneThousandFiveHundredFortySix} & Gaps15421550 & PROVED & \checkmark & verified \\
\texttt{nu\_p\_oneThousandFiveHundredFortyTwo} & Gaps15421550 & PROVED & \checkmark & verified \\
\texttt{nu\_p\_oneThousandFiveHundredFortyTwo\_two} & Gaps15421550 & PROVED & \checkmark & verified \\
\texttt{singular\_series\_finite\_pos\_evenPair\_oneThousandFiveHundredFifty} & Gaps15421550 & PROVED & \checkmark & verified \\
\texttt{singular\_series\_finite\_pos\_evenPair\_oneThousandFiveHundredFortyEight} & Gaps15421550 & PROVED & \checkmark & verified \\
\texttt{singular\_series\_finite\_pos\_evenPair\_oneThousandFiveHundredFortyFour} & Gaps15421550 & PROVED & \checkmark & verified \\
\texttt{singular\_series\_finite\_pos\_evenPair\_oneThousandFiveHundredFortySix} & Gaps15421550 & PROVED & \checkmark & verified \\
\texttt{singular\_series\_finite\_pos\_evenPair\_oneThousandFiveHundredFortyTwo} & Gaps15421550 & PROVED & \checkmark & verified \\
\texttt{singular\_series\_pos\_evenPair\_oneThousandFiveHundredFifty} & Gaps15421550 & PROVED & \checkmark & verified \\
\texttt{singular\_series\_pos\_evenPair\_oneThousandFiveHundredFortyEight} & Gaps15421550 & PROVED & \checkmark & verified \\
\texttt{singular\_series\_pos\_evenPair\_oneThousandFiveHundredFortyFour} & Gaps15421550 & PROVED & \checkmark & verified \\
\texttt{singular\_series\_pos\_evenPair\_oneThousandFiveHundredFortySix} & Gaps15421550 & PROVED & \checkmark & verified \\
\texttt{singular\_series\_pos\_evenPair\_oneThousandFiveHundredFortyTwo} & Gaps15421550 & PROVED & \checkmark & verified \\
\texttt{evenPair\_card\_oneThousandFiveHundredFiftyEight} & Gaps15521560 & PROVED & \checkmark & verified \\
\texttt{evenPair\_card\_oneThousandFiveHundredFiftyFour} & Gaps15521560 & PROVED & \checkmark & verified \\
\texttt{evenPair\_card\_oneThousandFiveHundredFiftySix} & Gaps15521560 & PROVED & \checkmark & verified \\
\texttt{evenPair\_card\_oneThousandFiveHundredFiftyTwo} & Gaps15521560 & PROVED & \checkmark & verified \\
\texttt{evenPair\_card\_oneThousandFiveHundredSixty} & Gaps15521560 & PROVED & \checkmark & verified \\
\texttt{isAdmissible\_evenPair\_oneThousandFiveHundredFiftyEight} & Gaps15521560 & PROVED & \checkmark & verified \\
\texttt{isAdmissible\_evenPair\_oneThousandFiveHundredFiftyFour} & Gaps15521560 & PROVED & \checkmark & verified \\
\texttt{isAdmissible\_evenPair\_oneThousandFiveHundredFiftySix} & Gaps15521560 & PROVED & \checkmark & verified \\
\texttt{isAdmissible\_evenPair\_oneThousandFiveHundredFiftyTwo} & Gaps15521560 & PROVED & \checkmark & verified \\
\texttt{isAdmissible\_evenPair\_oneThousandFiveHundredSixty} & Gaps15521560 & PROVED & \checkmark & verified \\
\texttt{localFactor\_oneThousandFiveHundredFiftyTwo\_two} & Gaps15521560 & PROVED & \checkmark & verified \\
\texttt{localFactor\_oneThousandFiveHundredSixty\_two} & Gaps15521560 & PROVED & \checkmark & verified \\
\texttt{nu\_p\_oneThousandFiveHundredFiftyEight} & Gaps15521560 & PROVED & \checkmark & verified \\
\texttt{nu\_p\_oneThousandFiveHundredFiftyFour} & Gaps15521560 & PROVED & \checkmark & verified \\
\texttt{nu\_p\_oneThousandFiveHundredFiftySix} & Gaps15521560 & PROVED & \checkmark & verified \\
\texttt{nu\_p\_oneThousandFiveHundredFiftyTwo} & Gaps15521560 & PROVED & \checkmark & verified \\
\texttt{nu\_p\_oneThousandFiveHundredFiftyTwo\_two} & Gaps15521560 & PROVED & \checkmark & verified \\
\texttt{nu\_p\_oneThousandFiveHundredSixty} & Gaps15521560 & PROVED & \checkmark & verified \\
\texttt{nu\_p\_oneThousandFiveHundredSixty\_two} & Gaps15521560 & PROVED & \checkmark & verified \\
\texttt{singular\_series\_finite\_pos\_evenPair\_oneThousandFiveHundredFiftyEight} & Gaps15521560 & PROVED & \checkmark & verified \\
\texttt{singular\_series\_finite\_pos\_evenPair\_oneThousandFiveHundredFiftyFour} & Gaps15521560 & PROVED & \checkmark & verified \\
\texttt{singular\_series\_finite\_pos\_evenPair\_oneThousandFiveHundredFiftySix} & Gaps15521560 & PROVED & \checkmark & verified \\
\texttt{singular\_series\_finite\_pos\_evenPair\_oneThousandFiveHundredFiftyTwo} & Gaps15521560 & PROVED & \checkmark & verified \\
\texttt{singular\_series\_finite\_pos\_evenPair\_oneThousandFiveHundredSixty} & Gaps15521560 & PROVED & \checkmark & verified \\
\texttt{singular\_series\_pos\_evenPair\_oneThousandFiveHundredFiftyEight} & Gaps15521560 & PROVED & \checkmark & verified \\
\texttt{singular\_series\_pos\_evenPair\_oneThousandFiveHundredFiftyFour} & Gaps15521560 & PROVED & \checkmark & verified \\
\texttt{singular\_series\_pos\_evenPair\_oneThousandFiveHundredFiftySix} & Gaps15521560 & PROVED & \checkmark & verified \\
\texttt{singular\_series\_pos\_evenPair\_oneThousandFiveHundredFiftyTwo} & Gaps15521560 & PROVED & \checkmark & verified \\
\texttt{singular\_series\_pos\_evenPair\_oneThousandFiveHundredSixty} & Gaps15521560 & PROVED & \checkmark & verified \\
\texttt{evenPair\_card\_oneThousandFiveHundredSeventy} & Gaps15621570 & PROVED & \checkmark & verified \\
\texttt{evenPair\_card\_oneThousandFiveHundredSixtyEight} & Gaps15621570 & PROVED & \checkmark & verified \\
\texttt{evenPair\_card\_oneThousandFiveHundredSixtyFour} & Gaps15621570 & PROVED & \checkmark & verified \\
\texttt{evenPair\_card\_oneThousandFiveHundredSixtySix} & Gaps15621570 & PROVED & \checkmark & verified \\
\texttt{evenPair\_card\_oneThousandFiveHundredSixtyTwo} & Gaps15621570 & PROVED & \checkmark & verified \\
\texttt{isAdmissible\_evenPair\_oneThousandFiveHundredSeventy} & Gaps15621570 & PROVED & \checkmark & verified \\
\texttt{isAdmissible\_evenPair\_oneThousandFiveHundredSixtyEight} & Gaps15621570 & PROVED & \checkmark & verified \\
\texttt{isAdmissible\_evenPair\_oneThousandFiveHundredSixtyFour} & Gaps15621570 & PROVED & \checkmark & verified \\
\texttt{isAdmissible\_evenPair\_oneThousandFiveHundredSixtySix} & Gaps15621570 & PROVED & \checkmark & verified \\
\texttt{isAdmissible\_evenPair\_oneThousandFiveHundredSixtyTwo} & Gaps15621570 & PROVED & \checkmark & verified \\
\texttt{localFactor\_oneThousandFiveHundredSeventy\_two} & Gaps15621570 & PROVED & \checkmark & verified \\
\texttt{localFactor\_oneThousandFiveHundredSixtyTwo\_two} & Gaps15621570 & PROVED & \checkmark & verified \\
\texttt{nu\_p\_oneThousandFiveHundredSeventy} & Gaps15621570 & PROVED & \checkmark & verified \\
\texttt{nu\_p\_oneThousandFiveHundredSeventy\_two} & Gaps15621570 & PROVED & \checkmark & verified \\
\texttt{nu\_p\_oneThousandFiveHundredSixtyEight} & Gaps15621570 & PROVED & \checkmark & verified \\
\texttt{nu\_p\_oneThousandFiveHundredSixtyFour} & Gaps15621570 & PROVED & \checkmark & verified \\
\texttt{nu\_p\_oneThousandFiveHundredSixtySix} & Gaps15621570 & PROVED & \checkmark & verified \\
\texttt{nu\_p\_oneThousandFiveHundredSixtyTwo} & Gaps15621570 & PROVED & \checkmark & verified \\
\texttt{nu\_p\_oneThousandFiveHundredSixtyTwo\_two} & Gaps15621570 & PROVED & \checkmark & verified \\
\texttt{singular\_series\_finite\_pos\_evenPair\_oneThousandFiveHundredSeventy} & Gaps15621570 & PROVED & \checkmark & verified \\
\texttt{singular\_series\_finite\_pos\_evenPair\_oneThousandFiveHundredSixtyEight} & Gaps15621570 & PROVED & \checkmark & verified \\
\texttt{singular\_series\_finite\_pos\_evenPair\_oneThousandFiveHundredSixtyFour} & Gaps15621570 & PROVED & \checkmark & verified \\
\texttt{singular\_series\_finite\_pos\_evenPair\_oneThousandFiveHundredSixtySix} & Gaps15621570 & PROVED & \checkmark & verified \\
\texttt{singular\_series\_finite\_pos\_evenPair\_oneThousandFiveHundredSixtyTwo} & Gaps15621570 & PROVED & \checkmark & verified \\
\texttt{singular\_series\_pos\_evenPair\_oneThousandFiveHundredSeventy} & Gaps15621570 & PROVED & \checkmark & verified \\
\texttt{singular\_series\_pos\_evenPair\_oneThousandFiveHundredSixtyEight} & Gaps15621570 & PROVED & \checkmark & verified \\
\texttt{singular\_series\_pos\_evenPair\_oneThousandFiveHundredSixtyFour} & Gaps15621570 & PROVED & \checkmark & verified \\
\texttt{singular\_series\_pos\_evenPair\_oneThousandFiveHundredSixtySix} & Gaps15621570 & PROVED & \checkmark & verified \\
\texttt{singular\_series\_pos\_evenPair\_oneThousandFiveHundredSixtyTwo} & Gaps15621570 & PROVED & \checkmark & verified \\
\texttt{evenPair\_card\_oneThousandFiveHundredEighty} & Gaps15721580 & PROVED & \checkmark & verified \\
\texttt{evenPair\_card\_oneThousandFiveHundredSeventyEight} & Gaps15721580 & PROVED & \checkmark & verified \\
\texttt{evenPair\_card\_oneThousandFiveHundredSeventyFour} & Gaps15721580 & PROVED & \checkmark & verified \\
\texttt{evenPair\_card\_oneThousandFiveHundredSeventySix} & Gaps15721580 & PROVED & \checkmark & verified \\
\texttt{evenPair\_card\_oneThousandFiveHundredSeventyTwo} & Gaps15721580 & PROVED & \checkmark & verified \\
\texttt{isAdmissible\_evenPair\_oneThousandFiveHundredEighty} & Gaps15721580 & PROVED & \checkmark & verified \\
\texttt{isAdmissible\_evenPair\_oneThousandFiveHundredSeventyEight} & Gaps15721580 & PROVED & \checkmark & verified \\
\texttt{isAdmissible\_evenPair\_oneThousandFiveHundredSeventyFour} & Gaps15721580 & PROVED & \checkmark & verified \\
\texttt{isAdmissible\_evenPair\_oneThousandFiveHundredSeventySix} & Gaps15721580 & PROVED & \checkmark & verified \\
\texttt{isAdmissible\_evenPair\_oneThousandFiveHundredSeventyTwo} & Gaps15721580 & PROVED & \checkmark & verified \\
\texttt{localFactor\_oneThousandFiveHundredEighty\_two} & Gaps15721580 & PROVED & \checkmark & verified \\
\texttt{localFactor\_oneThousandFiveHundredSeventyTwo\_two} & Gaps15721580 & PROVED & \checkmark & verified \\
\texttt{nu\_p\_oneThousandFiveHundredEighty} & Gaps15721580 & PROVED & \checkmark & verified \\
\texttt{nu\_p\_oneThousandFiveHundredEighty\_two} & Gaps15721580 & PROVED & \checkmark & verified \\
\texttt{nu\_p\_oneThousandFiveHundredSeventyEight} & Gaps15721580 & PROVED & \checkmark & verified \\
\texttt{nu\_p\_oneThousandFiveHundredSeventyFour} & Gaps15721580 & PROVED & \checkmark & verified \\
\texttt{nu\_p\_oneThousandFiveHundredSeventySix} & Gaps15721580 & PROVED & \checkmark & verified \\
\texttt{nu\_p\_oneThousandFiveHundredSeventyTwo} & Gaps15721580 & PROVED & \checkmark & verified \\
\texttt{nu\_p\_oneThousandFiveHundredSeventyTwo\_two} & Gaps15721580 & PROVED & \checkmark & verified \\
\texttt{singular\_series\_finite\_pos\_evenPair\_oneThousandFiveHundredEighty} & Gaps15721580 & PROVED & \checkmark & verified \\
\texttt{singular\_series\_finite\_pos\_evenPair\_oneThousandFiveHundredSeventyEight} & Gaps15721580 & PROVED & \checkmark & verified \\
\texttt{singular\_series\_finite\_pos\_evenPair\_oneThousandFiveHundredSeventyFour} & Gaps15721580 & PROVED & \checkmark & verified \\
\texttt{singular\_series\_finite\_pos\_evenPair\_oneThousandFiveHundredSeventySix} & Gaps15721580 & PROVED & \checkmark & verified \\
\texttt{singular\_series\_finite\_pos\_evenPair\_oneThousandFiveHundredSeventyTwo} & Gaps15721580 & PROVED & \checkmark & verified \\
\texttt{singular\_series\_pos\_evenPair\_oneThousandFiveHundredEighty} & Gaps15721580 & PROVED & \checkmark & verified \\
\texttt{singular\_series\_pos\_evenPair\_oneThousandFiveHundredSeventyEight} & Gaps15721580 & PROVED & \checkmark & verified \\
\texttt{singular\_series\_pos\_evenPair\_oneThousandFiveHundredSeventyFour} & Gaps15721580 & PROVED & \checkmark & verified \\
\texttt{singular\_series\_pos\_evenPair\_oneThousandFiveHundredSeventySix} & Gaps15721580 & PROVED & \checkmark & verified \\
\texttt{singular\_series\_pos\_evenPair\_oneThousandFiveHundredSeventyTwo} & Gaps15721580 & PROVED & \checkmark & verified \\
\texttt{evenPair\_card\_oneThousandFiveHundredEightyEight} & Gaps15821590 & PROVED & \checkmark & verified \\
\texttt{evenPair\_card\_oneThousandFiveHundredEightyFour} & Gaps15821590 & PROVED & \checkmark & verified \\
\texttt{evenPair\_card\_oneThousandFiveHundredEightySix} & Gaps15821590 & PROVED & \checkmark & verified \\
\texttt{evenPair\_card\_oneThousandFiveHundredEightyTwo} & Gaps15821590 & PROVED & \checkmark & verified \\
\texttt{evenPair\_card\_oneThousandFiveHundredNinety} & Gaps15821590 & PROVED & \checkmark & verified \\
\texttt{isAdmissible\_evenPair\_oneThousandFiveHundredEightyEight} & Gaps15821590 & PROVED & \checkmark & verified \\
\texttt{isAdmissible\_evenPair\_oneThousandFiveHundredEightyFour} & Gaps15821590 & PROVED & \checkmark & verified \\
\texttt{isAdmissible\_evenPair\_oneThousandFiveHundredEightySix} & Gaps15821590 & PROVED & \checkmark & verified \\
\texttt{isAdmissible\_evenPair\_oneThousandFiveHundredEightyTwo} & Gaps15821590 & PROVED & \checkmark & verified \\
\texttt{isAdmissible\_evenPair\_oneThousandFiveHundredNinety} & Gaps15821590 & PROVED & \checkmark & verified \\
\texttt{localFactor\_oneThousandFiveHundredEightyTwo\_two} & Gaps15821590 & PROVED & \checkmark & verified \\
\texttt{localFactor\_oneThousandFiveHundredNinety\_two} & Gaps15821590 & PROVED & \checkmark & verified \\
\texttt{nu\_p\_oneThousandFiveHundredEightyEight} & Gaps15821590 & PROVED & \checkmark & verified \\
\texttt{nu\_p\_oneThousandFiveHundredEightyFour} & Gaps15821590 & PROVED & \checkmark & verified \\
\texttt{nu\_p\_oneThousandFiveHundredEightySix} & Gaps15821590 & PROVED & \checkmark & verified \\
\texttt{nu\_p\_oneThousandFiveHundredEightyTwo} & Gaps15821590 & PROVED & \checkmark & verified \\
\texttt{nu\_p\_oneThousandFiveHundredEightyTwo\_two} & Gaps15821590 & PROVED & \checkmark & verified \\
\texttt{nu\_p\_oneThousandFiveHundredNinety} & Gaps15821590 & PROVED & \checkmark & verified \\
\texttt{nu\_p\_oneThousandFiveHundredNinety\_two} & Gaps15821590 & PROVED & \checkmark & verified \\
\texttt{singular\_series\_finite\_pos\_evenPair\_oneThousandFiveHundredEightyEight} & Gaps15821590 & PROVED & \checkmark & verified \\
\texttt{singular\_series\_finite\_pos\_evenPair\_oneThousandFiveHundredEightyFour} & Gaps15821590 & PROVED & \checkmark & verified \\
\texttt{singular\_series\_finite\_pos\_evenPair\_oneThousandFiveHundredEightySix} & Gaps15821590 & PROVED & \checkmark & verified \\
\texttt{singular\_series\_finite\_pos\_evenPair\_oneThousandFiveHundredEightyTwo} & Gaps15821590 & PROVED & \checkmark & verified \\
\texttt{singular\_series\_finite\_pos\_evenPair\_oneThousandFiveHundredNinety} & Gaps15821590 & PROVED & \checkmark & verified \\
\texttt{singular\_series\_pos\_evenPair\_oneThousandFiveHundredEightyEight} & Gaps15821590 & PROVED & \checkmark & verified \\
\texttt{singular\_series\_pos\_evenPair\_oneThousandFiveHundredEightyFour} & Gaps15821590 & PROVED & \checkmark & verified \\
\texttt{singular\_series\_pos\_evenPair\_oneThousandFiveHundredEightySix} & Gaps15821590 & PROVED & \checkmark & verified \\
\texttt{singular\_series\_pos\_evenPair\_oneThousandFiveHundredEightyTwo} & Gaps15821590 & PROVED & \checkmark & verified \\
\texttt{singular\_series\_pos\_evenPair\_oneThousandFiveHundredNinety} & Gaps15821590 & PROVED & \checkmark & verified \\
\texttt{evenPair\_card\_oneThousandFiveHundredNinetyEight} & Gaps15921600 & PROVED & \checkmark & verified \\
\texttt{evenPair\_card\_oneThousandFiveHundredNinetyFour} & Gaps15921600 & PROVED & \checkmark & verified \\
\texttt{evenPair\_card\_oneThousandFiveHundredNinetySix} & Gaps15921600 & PROVED & \checkmark & verified \\
\texttt{evenPair\_card\_oneThousandFiveHundredNinetyTwo} & Gaps15921600 & PROVED & \checkmark & verified \\
\texttt{evenPair\_card\_oneThousandSixHundred} & Gaps15921600 & PROVED & \checkmark & verified \\
\texttt{isAdmissible\_evenPair\_oneThousandFiveHundredNinetyEight} & Gaps15921600 & PROVED & \checkmark & verified \\
\texttt{isAdmissible\_evenPair\_oneThousandFiveHundredNinetyFour} & Gaps15921600 & PROVED & \checkmark & verified \\
\texttt{isAdmissible\_evenPair\_oneThousandFiveHundredNinetySix} & Gaps15921600 & PROVED & \checkmark & verified \\
\texttt{isAdmissible\_evenPair\_oneThousandFiveHundredNinetyTwo} & Gaps15921600 & PROVED & \checkmark & verified \\
\texttt{isAdmissible\_evenPair\_oneThousandSixHundred} & Gaps15921600 & PROVED & \checkmark & verified \\
\texttt{localFactor\_oneThousandFiveHundredNinetyTwo\_two} & Gaps15921600 & PROVED & \checkmark & verified \\
\texttt{localFactor\_oneThousandSixHundred\_two} & Gaps15921600 & PROVED & \checkmark & verified \\
\texttt{nu\_p\_oneThousandFiveHundredNinetyEight} & Gaps15921600 & PROVED & \checkmark & verified \\
\texttt{nu\_p\_oneThousandFiveHundredNinetyFour} & Gaps15921600 & PROVED & \checkmark & verified \\
\texttt{nu\_p\_oneThousandFiveHundredNinetySix} & Gaps15921600 & PROVED & \checkmark & verified \\
\texttt{nu\_p\_oneThousandFiveHundredNinetyTwo} & Gaps15921600 & PROVED & \checkmark & verified \\
\texttt{nu\_p\_oneThousandFiveHundredNinetyTwo\_two} & Gaps15921600 & PROVED & \checkmark & verified \\
\texttt{nu\_p\_oneThousandSixHundred} & Gaps15921600 & PROVED & \checkmark & verified \\
\texttt{nu\_p\_oneThousandSixHundred\_two} & Gaps15921600 & PROVED & \checkmark & verified \\
\texttt{singular\_series\_finite\_pos\_evenPair\_oneThousandFiveHundredNinetyEight} & Gaps15921600 & PROVED & \checkmark & verified \\
\texttt{singular\_series\_finite\_pos\_evenPair\_oneThousandFiveHundredNinetyFour} & Gaps15921600 & PROVED & \checkmark & verified \\
\texttt{singular\_series\_finite\_pos\_evenPair\_oneThousandFiveHundredNinetySix} & Gaps15921600 & PROVED & \checkmark & verified \\
\texttt{singular\_series\_finite\_pos\_evenPair\_oneThousandFiveHundredNinetyTwo} & Gaps15921600 & PROVED & \checkmark & verified \\
\texttt{singular\_series\_finite\_pos\_evenPair\_oneThousandSixHundred} & Gaps15921600 & PROVED & \checkmark & verified \\
\texttt{singular\_series\_pos\_evenPair\_oneThousandFiveHundredNinetyEight} & Gaps15921600 & PROVED & \checkmark & verified \\
\texttt{singular\_series\_pos\_evenPair\_oneThousandFiveHundredNinetyFour} & Gaps15921600 & PROVED & \checkmark & verified \\
\texttt{singular\_series\_pos\_evenPair\_oneThousandFiveHundredNinetySix} & Gaps15921600 & PROVED & \checkmark & verified \\
\texttt{singular\_series\_pos\_evenPair\_oneThousandFiveHundredNinetyTwo} & Gaps15921600 & PROVED & \checkmark & verified \\
\texttt{singular\_series\_pos\_evenPair\_oneThousandSixHundred} & Gaps15921600 & PROVED & \checkmark & verified \\
\texttt{evenPair\_card\_oneThousandSixHundredEight} & Gaps16021610 & PROVED & \checkmark & verified \\
\texttt{evenPair\_card\_oneThousandSixHundredFour} & Gaps16021610 & PROVED & \checkmark & verified \\
\texttt{evenPair\_card\_oneThousandSixHundredSix} & Gaps16021610 & PROVED & \checkmark & verified \\
\texttt{evenPair\_card\_oneThousandSixHundredTen} & Gaps16021610 & PROVED & \checkmark & verified \\
\texttt{evenPair\_card\_oneThousandSixHundredTwo} & Gaps16021610 & PROVED & \checkmark & verified \\
\texttt{isAdmissible\_evenPair\_oneThousandSixHundredEight} & Gaps16021610 & PROVED & \checkmark & verified \\
\texttt{isAdmissible\_evenPair\_oneThousandSixHundredFour} & Gaps16021610 & PROVED & \checkmark & verified \\
\texttt{isAdmissible\_evenPair\_oneThousandSixHundredSix} & Gaps16021610 & PROVED & \checkmark & verified \\
\texttt{isAdmissible\_evenPair\_oneThousandSixHundredTen} & Gaps16021610 & PROVED & \checkmark & verified \\
\texttt{isAdmissible\_evenPair\_oneThousandSixHundredTwo} & Gaps16021610 & PROVED & \checkmark & verified \\
\texttt{localFactor\_oneThousandSixHundredTen\_two} & Gaps16021610 & PROVED & \checkmark & verified \\
\texttt{localFactor\_oneThousandSixHundredTwo\_two} & Gaps16021610 & PROVED & \checkmark & verified \\
\texttt{nu\_p\_oneThousandSixHundredEight} & Gaps16021610 & PROVED & \checkmark & verified \\
\texttt{nu\_p\_oneThousandSixHundredFour} & Gaps16021610 & PROVED & \checkmark & verified \\
\texttt{nu\_p\_oneThousandSixHundredSix} & Gaps16021610 & PROVED & \checkmark & verified \\
\texttt{nu\_p\_oneThousandSixHundredTen} & Gaps16021610 & PROVED & \checkmark & verified \\
\texttt{nu\_p\_oneThousandSixHundredTen\_two} & Gaps16021610 & PROVED & \checkmark & verified \\
\texttt{nu\_p\_oneThousandSixHundredTwo} & Gaps16021610 & PROVED & \checkmark & verified \\
\texttt{nu\_p\_oneThousandSixHundredTwo\_two} & Gaps16021610 & PROVED & \checkmark & verified \\
\texttt{singular\_series\_finite\_pos\_evenPair\_oneThousandSixHundredEight} & Gaps16021610 & PROVED & \checkmark & verified \\
\texttt{singular\_series\_finite\_pos\_evenPair\_oneThousandSixHundredFour} & Gaps16021610 & PROVED & \checkmark & verified \\
\texttt{singular\_series\_finite\_pos\_evenPair\_oneThousandSixHundredSix} & Gaps16021610 & PROVED & \checkmark & verified \\
\texttt{singular\_series\_finite\_pos\_evenPair\_oneThousandSixHundredTen} & Gaps16021610 & PROVED & \checkmark & verified \\
\texttt{singular\_series\_finite\_pos\_evenPair\_oneThousandSixHundredTwo} & Gaps16021610 & PROVED & \checkmark & verified \\
\texttt{singular\_series\_pos\_evenPair\_oneThousandSixHundredEight} & Gaps16021610 & PROVED & \checkmark & verified \\
\texttt{singular\_series\_pos\_evenPair\_oneThousandSixHundredFour} & Gaps16021610 & PROVED & \checkmark & verified \\
\texttt{singular\_series\_pos\_evenPair\_oneThousandSixHundredSix} & Gaps16021610 & PROVED & \checkmark & verified \\
\texttt{singular\_series\_pos\_evenPair\_oneThousandSixHundredTen} & Gaps16021610 & PROVED & \checkmark & verified \\
\texttt{singular\_series\_pos\_evenPair\_oneThousandSixHundredTwo} & Gaps16021610 & PROVED & \checkmark & verified \\
\texttt{evenPair\_card\_oneThousandSixHundredEighteen} & Gaps16121620 & PROVED & \checkmark & verified \\
\texttt{evenPair\_card\_oneThousandSixHundredFourteen} & Gaps16121620 & PROVED & \checkmark & verified \\
\texttt{evenPair\_card\_oneThousandSixHundredSixteen} & Gaps16121620 & PROVED & \checkmark & verified \\
\texttt{evenPair\_card\_oneThousandSixHundredTwelve} & Gaps16121620 & PROVED & \checkmark & verified \\
\texttt{evenPair\_card\_oneThousandSixHundredTwenty} & Gaps16121620 & PROVED & \checkmark & verified \\
\texttt{isAdmissible\_evenPair\_oneThousandSixHundredEighteen} & Gaps16121620 & PROVED & \checkmark & verified \\
\texttt{isAdmissible\_evenPair\_oneThousandSixHundredFourteen} & Gaps16121620 & PROVED & \checkmark & verified \\
\texttt{isAdmissible\_evenPair\_oneThousandSixHundredSixteen} & Gaps16121620 & PROVED & \checkmark & verified \\
\texttt{isAdmissible\_evenPair\_oneThousandSixHundredTwelve} & Gaps16121620 & PROVED & \checkmark & verified \\
\texttt{isAdmissible\_evenPair\_oneThousandSixHundredTwenty} & Gaps16121620 & PROVED & \checkmark & verified \\
\texttt{localFactor\_oneThousandSixHundredTwelve\_two} & Gaps16121620 & PROVED & \checkmark & verified \\
\texttt{localFactor\_oneThousandSixHundredTwenty\_two} & Gaps16121620 & PROVED & \checkmark & verified \\
\texttt{nu\_p\_oneThousandSixHundredEighteen} & Gaps16121620 & PROVED & \checkmark & verified \\
\texttt{nu\_p\_oneThousandSixHundredFourteen} & Gaps16121620 & PROVED & \checkmark & verified \\
\texttt{nu\_p\_oneThousandSixHundredSixteen} & Gaps16121620 & PROVED & \checkmark & verified \\
\texttt{nu\_p\_oneThousandSixHundredTwelve} & Gaps16121620 & PROVED & \checkmark & verified \\
\texttt{nu\_p\_oneThousandSixHundredTwelve\_two} & Gaps16121620 & PROVED & \checkmark & verified \\
\texttt{nu\_p\_oneThousandSixHundredTwenty} & Gaps16121620 & PROVED & \checkmark & verified \\
\texttt{nu\_p\_oneThousandSixHundredTwenty\_two} & Gaps16121620 & PROVED & \checkmark & verified \\
\texttt{singular\_series\_finite\_pos\_evenPair\_oneThousandSixHundredEighteen} & Gaps16121620 & PROVED & \checkmark & verified \\
\texttt{singular\_series\_finite\_pos\_evenPair\_oneThousandSixHundredFourteen} & Gaps16121620 & PROVED & \checkmark & verified \\
\texttt{singular\_series\_finite\_pos\_evenPair\_oneThousandSixHundredSixteen} & Gaps16121620 & PROVED & \checkmark & verified \\
\texttt{singular\_series\_finite\_pos\_evenPair\_oneThousandSixHundredTwelve} & Gaps16121620 & PROVED & \checkmark & verified \\
\texttt{singular\_series\_finite\_pos\_evenPair\_oneThousandSixHundredTwenty} & Gaps16121620 & PROVED & \checkmark & verified \\
\texttt{singular\_series\_pos\_evenPair\_oneThousandSixHundredEighteen} & Gaps16121620 & PROVED & \checkmark & verified \\
\texttt{singular\_series\_pos\_evenPair\_oneThousandSixHundredFourteen} & Gaps16121620 & PROVED & \checkmark & verified \\
\texttt{singular\_series\_pos\_evenPair\_oneThousandSixHundredSixteen} & Gaps16121620 & PROVED & \checkmark & verified \\
\texttt{singular\_series\_pos\_evenPair\_oneThousandSixHundredTwelve} & Gaps16121620 & PROVED & \checkmark & verified \\
\texttt{singular\_series\_pos\_evenPair\_oneThousandSixHundredTwenty} & Gaps16121620 & PROVED & \checkmark & verified \\
\texttt{evenPair\_card\_oneHundredSeventy} & Gaps162170 & PROVED & \checkmark & verified \\
\texttt{evenPair\_card\_oneHundredSixtyEight} & Gaps162170 & PROVED & \checkmark & verified \\
\texttt{evenPair\_card\_oneHundredSixtyFour} & Gaps162170 & PROVED & \checkmark & verified \\
\texttt{evenPair\_card\_oneHundredSixtySix} & Gaps162170 & PROVED & \checkmark & verified \\
\texttt{evenPair\_card\_oneHundredSixtyTwo} & Gaps162170 & PROVED & \checkmark & verified \\
\texttt{isAdmissible\_evenPair\_oneHundredSeventy} & Gaps162170 & PROVED & \checkmark & verified \\
\texttt{isAdmissible\_evenPair\_oneHundredSixtyEight} & Gaps162170 & PROVED & \checkmark & verified \\
\texttt{isAdmissible\_evenPair\_oneHundredSixtyFour} & Gaps162170 & PROVED & \checkmark & verified \\
\texttt{isAdmissible\_evenPair\_oneHundredSixtySix} & Gaps162170 & PROVED & \checkmark & verified \\
\texttt{isAdmissible\_evenPair\_oneHundredSixtyTwo} & Gaps162170 & PROVED & \checkmark & verified \\
\texttt{localFactor\_oneHundredSeventy\_two} & Gaps162170 & PROVED & \checkmark & verified \\
\texttt{localFactor\_oneHundredSixtyTwo\_two} & Gaps162170 & PROVED & \checkmark & verified \\
\texttt{nu\_p\_oneHundredSeventy} & Gaps162170 & PROVED & \checkmark & verified \\
\texttt{nu\_p\_oneHundredSeventy\_two} & Gaps162170 & PROVED & \checkmark & verified \\
\texttt{nu\_p\_oneHundredSixtyEight} & Gaps162170 & PROVED & \checkmark & verified \\
\texttt{nu\_p\_oneHundredSixtyFour} & Gaps162170 & PROVED & \checkmark & verified \\
\texttt{nu\_p\_oneHundredSixtySix} & Gaps162170 & PROVED & \checkmark & verified \\
\texttt{nu\_p\_oneHundredSixtyTwo} & Gaps162170 & PROVED & \checkmark & verified \\
\texttt{nu\_p\_oneHundredSixtyTwo\_two} & Gaps162170 & PROVED & \checkmark & verified \\
\texttt{singular\_series\_finite\_pos\_evenPair\_oneHundredSeventy} & Gaps162170 & PROVED & \checkmark & verified \\
\texttt{singular\_series\_finite\_pos\_evenPair\_oneHundredSixtyEight} & Gaps162170 & PROVED & \checkmark & verified \\
\texttt{singular\_series\_finite\_pos\_evenPair\_oneHundredSixtyFour} & Gaps162170 & PROVED & \checkmark & verified \\
\texttt{singular\_series\_finite\_pos\_evenPair\_oneHundredSixtySix} & Gaps162170 & PROVED & \checkmark & verified \\
\texttt{singular\_series\_finite\_pos\_evenPair\_oneHundredSixtyTwo} & Gaps162170 & PROVED & \checkmark & verified \\
\texttt{singular\_series\_pos\_evenPair\_oneHundredSeventy} & Gaps162170 & PROVED & \checkmark & verified \\
\texttt{singular\_series\_pos\_evenPair\_oneHundredSixtyEight} & Gaps162170 & PROVED & \checkmark & verified \\
\texttt{singular\_series\_pos\_evenPair\_oneHundredSixtyFour} & Gaps162170 & PROVED & \checkmark & verified \\
\texttt{singular\_series\_pos\_evenPair\_oneHundredSixtySix} & Gaps162170 & PROVED & \checkmark & verified \\
\texttt{singular\_series\_pos\_evenPair\_oneHundredSixtyTwo} & Gaps162170 & PROVED & \checkmark & verified \\
\texttt{evenPair\_card\_oneThousandSixHundredThirty} & Gaps16221630 & PROVED & \checkmark & verified \\
\texttt{evenPair\_card\_oneThousandSixHundredTwentyEight} & Gaps16221630 & PROVED & \checkmark & verified \\
\texttt{evenPair\_card\_oneThousandSixHundredTwentyFour} & Gaps16221630 & PROVED & \checkmark & verified \\
\texttt{evenPair\_card\_oneThousandSixHundredTwentySix} & Gaps16221630 & PROVED & \checkmark & verified \\
\texttt{evenPair\_card\_oneThousandSixHundredTwentyTwo} & Gaps16221630 & PROVED & \checkmark & verified \\
\texttt{isAdmissible\_evenPair\_oneThousandSixHundredThirty} & Gaps16221630 & PROVED & \checkmark & verified \\
\texttt{isAdmissible\_evenPair\_oneThousandSixHundredTwentyEight} & Gaps16221630 & PROVED & \checkmark & verified \\
\texttt{isAdmissible\_evenPair\_oneThousandSixHundredTwentyFour} & Gaps16221630 & PROVED & \checkmark & verified \\
\texttt{isAdmissible\_evenPair\_oneThousandSixHundredTwentySix} & Gaps16221630 & PROVED & \checkmark & verified \\
\texttt{isAdmissible\_evenPair\_oneThousandSixHundredTwentyTwo} & Gaps16221630 & PROVED & \checkmark & verified \\
\texttt{localFactor\_oneThousandSixHundredThirty\_two} & Gaps16221630 & PROVED & \checkmark & verified \\
\texttt{localFactor\_oneThousandSixHundredTwentyTwo\_two} & Gaps16221630 & PROVED & \checkmark & verified \\
\texttt{nu\_p\_oneThousandSixHundredThirty} & Gaps16221630 & PROVED & \checkmark & verified \\
\texttt{nu\_p\_oneThousandSixHundredThirty\_two} & Gaps16221630 & PROVED & \checkmark & verified \\
\texttt{nu\_p\_oneThousandSixHundredTwentyEight} & Gaps16221630 & PROVED & \checkmark & verified \\
\texttt{nu\_p\_oneThousandSixHundredTwentyFour} & Gaps16221630 & PROVED & \checkmark & verified \\
\texttt{nu\_p\_oneThousandSixHundredTwentySix} & Gaps16221630 & PROVED & \checkmark & verified \\
\texttt{nu\_p\_oneThousandSixHundredTwentyTwo} & Gaps16221630 & PROVED & \checkmark & verified \\
\texttt{nu\_p\_oneThousandSixHundredTwentyTwo\_two} & Gaps16221630 & PROVED & \checkmark & verified \\
\texttt{singular\_series\_finite\_pos\_evenPair\_oneThousandSixHundredThirty} & Gaps16221630 & PROVED & \checkmark & verified \\
\texttt{singular\_series\_finite\_pos\_evenPair\_oneThousandSixHundredTwentyEight} & Gaps16221630 & PROVED & \checkmark & verified \\
\texttt{singular\_series\_finite\_pos\_evenPair\_oneThousandSixHundredTwentyFour} & Gaps16221630 & PROVED & \checkmark & verified \\
\texttt{singular\_series\_finite\_pos\_evenPair\_oneThousandSixHundredTwentySix} & Gaps16221630 & PROVED & \checkmark & verified \\
\texttt{singular\_series\_finite\_pos\_evenPair\_oneThousandSixHundredTwentyTwo} & Gaps16221630 & PROVED & \checkmark & verified \\
\texttt{singular\_series\_pos\_evenPair\_oneThousandSixHundredThirty} & Gaps16221630 & PROVED & \checkmark & verified \\
\texttt{singular\_series\_pos\_evenPair\_oneThousandSixHundredTwentyEight} & Gaps16221630 & PROVED & \checkmark & verified \\
\texttt{singular\_series\_pos\_evenPair\_oneThousandSixHundredTwentyFour} & Gaps16221630 & PROVED & \checkmark & verified \\
\texttt{singular\_series\_pos\_evenPair\_oneThousandSixHundredTwentySix} & Gaps16221630 & PROVED & \checkmark & verified \\
\texttt{singular\_series\_pos\_evenPair\_oneThousandSixHundredTwentyTwo} & Gaps16221630 & PROVED & \checkmark & verified \\
\texttt{evenPair\_card\_oneThousandSixHundredForty} & Gaps16321640 & PROVED & \checkmark & verified \\
\texttt{evenPair\_card\_oneThousandSixHundredThirtyEight} & Gaps16321640 & PROVED & \checkmark & verified \\
\texttt{evenPair\_card\_oneThousandSixHundredThirtyFour} & Gaps16321640 & PROVED & \checkmark & verified \\
\texttt{evenPair\_card\_oneThousandSixHundredThirtySix} & Gaps16321640 & PROVED & \checkmark & verified \\
\texttt{evenPair\_card\_oneThousandSixHundredThirtyTwo} & Gaps16321640 & PROVED & \checkmark & verified \\
\texttt{isAdmissible\_evenPair\_oneThousandSixHundredForty} & Gaps16321640 & PROVED & \checkmark & verified \\
\texttt{isAdmissible\_evenPair\_oneThousandSixHundredThirtyEight} & Gaps16321640 & PROVED & \checkmark & verified \\
\texttt{isAdmissible\_evenPair\_oneThousandSixHundredThirtyFour} & Gaps16321640 & PROVED & \checkmark & verified \\
\texttt{isAdmissible\_evenPair\_oneThousandSixHundredThirtySix} & Gaps16321640 & PROVED & \checkmark & verified \\
\texttt{isAdmissible\_evenPair\_oneThousandSixHundredThirtyTwo} & Gaps16321640 & PROVED & \checkmark & verified \\
\texttt{localFactor\_oneThousandSixHundredForty\_two} & Gaps16321640 & PROVED & \checkmark & verified \\
\texttt{localFactor\_oneThousandSixHundredThirtyTwo\_two} & Gaps16321640 & PROVED & \checkmark & verified \\
\texttt{nu\_p\_oneThousandSixHundredForty} & Gaps16321640 & PROVED & \checkmark & verified \\
\texttt{nu\_p\_oneThousandSixHundredForty\_two} & Gaps16321640 & PROVED & \checkmark & verified \\
\texttt{nu\_p\_oneThousandSixHundredThirtyEight} & Gaps16321640 & PROVED & \checkmark & verified \\
\texttt{nu\_p\_oneThousandSixHundredThirtyFour} & Gaps16321640 & PROVED & \checkmark & verified \\
\texttt{nu\_p\_oneThousandSixHundredThirtySix} & Gaps16321640 & PROVED & \checkmark & verified \\
\texttt{nu\_p\_oneThousandSixHundredThirtyTwo} & Gaps16321640 & PROVED & \checkmark & verified \\
\texttt{nu\_p\_oneThousandSixHundredThirtyTwo\_two} & Gaps16321640 & PROVED & \checkmark & verified \\
\texttt{singular\_series\_finite\_pos\_evenPair\_oneThousandSixHundredForty} & Gaps16321640 & PROVED & \checkmark & verified \\
\texttt{singular\_series\_finite\_pos\_evenPair\_oneThousandSixHundredThirtyEight} & Gaps16321640 & PROVED & \checkmark & verified \\
\texttt{singular\_series\_finite\_pos\_evenPair\_oneThousandSixHundredThirtyFour} & Gaps16321640 & PROVED & \checkmark & verified \\
\texttt{singular\_series\_finite\_pos\_evenPair\_oneThousandSixHundredThirtySix} & Gaps16321640 & PROVED & \checkmark & verified \\
\texttt{singular\_series\_finite\_pos\_evenPair\_oneThousandSixHundredThirtyTwo} & Gaps16321640 & PROVED & \checkmark & verified \\
\texttt{singular\_series\_pos\_evenPair\_oneThousandSixHundredForty} & Gaps16321640 & PROVED & \checkmark & verified \\
\texttt{singular\_series\_pos\_evenPair\_oneThousandSixHundredThirtyEight} & Gaps16321640 & PROVED & \checkmark & verified \\
\texttt{singular\_series\_pos\_evenPair\_oneThousandSixHundredThirtyFour} & Gaps16321640 & PROVED & \checkmark & verified \\
\texttt{singular\_series\_pos\_evenPair\_oneThousandSixHundredThirtySix} & Gaps16321640 & PROVED & \checkmark & verified \\
\texttt{singular\_series\_pos\_evenPair\_oneThousandSixHundredThirtyTwo} & Gaps16321640 & PROVED & \checkmark & verified \\
\texttt{evenPair\_card\_oneThousandSixHundredFifty} & Gaps16421650 & PROVED & \checkmark & verified \\
\texttt{evenPair\_card\_oneThousandSixHundredFortyEight} & Gaps16421650 & PROVED & \checkmark & verified \\
\texttt{evenPair\_card\_oneThousandSixHundredFortyFour} & Gaps16421650 & PROVED & \checkmark & verified \\
\texttt{evenPair\_card\_oneThousandSixHundredFortySix} & Gaps16421650 & PROVED & \checkmark & verified \\
\texttt{evenPair\_card\_oneThousandSixHundredFortyTwo} & Gaps16421650 & PROVED & \checkmark & verified \\
\texttt{isAdmissible\_evenPair\_oneThousandSixHundredFifty} & Gaps16421650 & PROVED & \checkmark & verified \\
\texttt{isAdmissible\_evenPair\_oneThousandSixHundredFortyEight} & Gaps16421650 & PROVED & \checkmark & verified \\
\texttt{isAdmissible\_evenPair\_oneThousandSixHundredFortyFour} & Gaps16421650 & PROVED & \checkmark & verified \\
\texttt{isAdmissible\_evenPair\_oneThousandSixHundredFortySix} & Gaps16421650 & PROVED & \checkmark & verified \\
\texttt{isAdmissible\_evenPair\_oneThousandSixHundredFortyTwo} & Gaps16421650 & PROVED & \checkmark & verified \\
\texttt{localFactor\_oneThousandSixHundredFifty\_two} & Gaps16421650 & PROVED & \checkmark & verified \\
\texttt{localFactor\_oneThousandSixHundredFortyTwo\_two} & Gaps16421650 & PROVED & \checkmark & verified \\
\texttt{nu\_p\_oneThousandSixHundredFifty} & Gaps16421650 & PROVED & \checkmark & verified \\
\texttt{nu\_p\_oneThousandSixHundredFifty\_two} & Gaps16421650 & PROVED & \checkmark & verified \\
\texttt{nu\_p\_oneThousandSixHundredFortyEight} & Gaps16421650 & PROVED & \checkmark & verified \\
\texttt{nu\_p\_oneThousandSixHundredFortyFour} & Gaps16421650 & PROVED & \checkmark & verified \\
\texttt{nu\_p\_oneThousandSixHundredFortySix} & Gaps16421650 & PROVED & \checkmark & verified \\
\texttt{nu\_p\_oneThousandSixHundredFortyTwo} & Gaps16421650 & PROVED & \checkmark & verified \\
\texttt{nu\_p\_oneThousandSixHundredFortyTwo\_two} & Gaps16421650 & PROVED & \checkmark & verified \\
\texttt{singular\_series\_finite\_pos\_evenPair\_oneThousandSixHundredFifty} & Gaps16421650 & PROVED & \checkmark & verified \\
\texttt{singular\_series\_finite\_pos\_evenPair\_oneThousandSixHundredFortyEight} & Gaps16421650 & PROVED & \checkmark & verified \\
\texttt{singular\_series\_finite\_pos\_evenPair\_oneThousandSixHundredFortyFour} & Gaps16421650 & PROVED & \checkmark & verified \\
\texttt{singular\_series\_finite\_pos\_evenPair\_oneThousandSixHundredFortySix} & Gaps16421650 & PROVED & \checkmark & verified \\
\texttt{singular\_series\_finite\_pos\_evenPair\_oneThousandSixHundredFortyTwo} & Gaps16421650 & PROVED & \checkmark & verified \\
\texttt{singular\_series\_pos\_evenPair\_oneThousandSixHundredFifty} & Gaps16421650 & PROVED & \checkmark & verified \\
\texttt{singular\_series\_pos\_evenPair\_oneThousandSixHundredFortyEight} & Gaps16421650 & PROVED & \checkmark & verified \\
\texttt{singular\_series\_pos\_evenPair\_oneThousandSixHundredFortyFour} & Gaps16421650 & PROVED & \checkmark & verified \\
\texttt{singular\_series\_pos\_evenPair\_oneThousandSixHundredFortySix} & Gaps16421650 & PROVED & \checkmark & verified \\
\texttt{singular\_series\_pos\_evenPair\_oneThousandSixHundredFortyTwo} & Gaps16421650 & PROVED & \checkmark & verified \\
\texttt{evenPair\_card\_oneThousandSixHundredFiftyEight} & Gaps16521660 & PROVED & \checkmark & verified \\
\texttt{evenPair\_card\_oneThousandSixHundredFiftyFour} & Gaps16521660 & PROVED & \checkmark & verified \\
\texttt{evenPair\_card\_oneThousandSixHundredFiftySix} & Gaps16521660 & PROVED & \checkmark & verified \\
\texttt{evenPair\_card\_oneThousandSixHundredFiftyTwo} & Gaps16521660 & PROVED & \checkmark & verified \\
\texttt{evenPair\_card\_oneThousandSixHundredSixty} & Gaps16521660 & PROVED & \checkmark & verified \\
\texttt{isAdmissible\_evenPair\_oneThousandSixHundredFiftyEight} & Gaps16521660 & PROVED & \checkmark & verified \\
\texttt{isAdmissible\_evenPair\_oneThousandSixHundredFiftyFour} & Gaps16521660 & PROVED & \checkmark & verified \\
\texttt{isAdmissible\_evenPair\_oneThousandSixHundredFiftySix} & Gaps16521660 & PROVED & \checkmark & verified \\
\texttt{isAdmissible\_evenPair\_oneThousandSixHundredFiftyTwo} & Gaps16521660 & PROVED & \checkmark & verified \\
\texttt{isAdmissible\_evenPair\_oneThousandSixHundredSixty} & Gaps16521660 & PROVED & \checkmark & verified \\
\texttt{localFactor\_oneThousandSixHundredFiftyTwo\_two} & Gaps16521660 & PROVED & \checkmark & verified \\
\texttt{localFactor\_oneThousandSixHundredSixty\_two} & Gaps16521660 & PROVED & \checkmark & verified \\
\texttt{nu\_p\_oneThousandSixHundredFiftyEight} & Gaps16521660 & PROVED & \checkmark & verified \\
\texttt{nu\_p\_oneThousandSixHundredFiftyFour} & Gaps16521660 & PROVED & \checkmark & verified \\
\texttt{nu\_p\_oneThousandSixHundredFiftySix} & Gaps16521660 & PROVED & \checkmark & verified \\
\texttt{nu\_p\_oneThousandSixHundredFiftyTwo} & Gaps16521660 & PROVED & \checkmark & verified \\
\texttt{nu\_p\_oneThousandSixHundredFiftyTwo\_two} & Gaps16521660 & PROVED & \checkmark & verified \\
\texttt{nu\_p\_oneThousandSixHundredSixty} & Gaps16521660 & PROVED & \checkmark & verified \\
\texttt{nu\_p\_oneThousandSixHundredSixty\_two} & Gaps16521660 & PROVED & \checkmark & verified \\
\texttt{singular\_series\_finite\_pos\_evenPair\_oneThousandSixHundredFiftyEight} & Gaps16521660 & PROVED & \checkmark & verified \\
\texttt{singular\_series\_finite\_pos\_evenPair\_oneThousandSixHundredFiftyFour} & Gaps16521660 & PROVED & \checkmark & verified \\
\texttt{singular\_series\_finite\_pos\_evenPair\_oneThousandSixHundredFiftySix} & Gaps16521660 & PROVED & \checkmark & verified \\
\texttt{singular\_series\_finite\_pos\_evenPair\_oneThousandSixHundredFiftyTwo} & Gaps16521660 & PROVED & \checkmark & verified \\
\texttt{singular\_series\_finite\_pos\_evenPair\_oneThousandSixHundredSixty} & Gaps16521660 & PROVED & \checkmark & verified \\
\texttt{singular\_series\_pos\_evenPair\_oneThousandSixHundredFiftyEight} & Gaps16521660 & PROVED & \checkmark & verified \\
\texttt{singular\_series\_pos\_evenPair\_oneThousandSixHundredFiftyFour} & Gaps16521660 & PROVED & \checkmark & verified \\
\texttt{singular\_series\_pos\_evenPair\_oneThousandSixHundredFiftySix} & Gaps16521660 & PROVED & \checkmark & verified \\
\texttt{singular\_series\_pos\_evenPair\_oneThousandSixHundredFiftyTwo} & Gaps16521660 & PROVED & \checkmark & verified \\
\texttt{singular\_series\_pos\_evenPair\_oneThousandSixHundredSixty} & Gaps16521660 & PROVED & \checkmark & verified \\
\texttt{evenPair\_card\_oneThousandSixHundredSeventy} & Gaps16621670 & PROVED & \checkmark & verified \\
\texttt{evenPair\_card\_oneThousandSixHundredSixtyEight} & Gaps16621670 & PROVED & \checkmark & verified \\
\texttt{evenPair\_card\_oneThousandSixHundredSixtyFour} & Gaps16621670 & PROVED & \checkmark & verified \\
\texttt{evenPair\_card\_oneThousandSixHundredSixtySix} & Gaps16621670 & PROVED & \checkmark & verified \\
\texttt{evenPair\_card\_oneThousandSixHundredSixtyTwo} & Gaps16621670 & PROVED & \checkmark & verified \\
\texttt{isAdmissible\_evenPair\_oneThousandSixHundredSeventy} & Gaps16621670 & PROVED & \checkmark & verified \\
\texttt{isAdmissible\_evenPair\_oneThousandSixHundredSixtyEight} & Gaps16621670 & PROVED & \checkmark & verified \\
\texttt{isAdmissible\_evenPair\_oneThousandSixHundredSixtyFour} & Gaps16621670 & PROVED & \checkmark & verified \\
\texttt{isAdmissible\_evenPair\_oneThousandSixHundredSixtySix} & Gaps16621670 & PROVED & \checkmark & verified \\
\texttt{isAdmissible\_evenPair\_oneThousandSixHundredSixtyTwo} & Gaps16621670 & PROVED & \checkmark & verified \\
\texttt{localFactor\_oneThousandSixHundredSeventy\_two} & Gaps16621670 & PROVED & \checkmark & verified \\
\texttt{localFactor\_oneThousandSixHundredSixtyTwo\_two} & Gaps16621670 & PROVED & \checkmark & verified \\
\texttt{nu\_p\_oneThousandSixHundredSeventy} & Gaps16621670 & PROVED & \checkmark & verified \\
\texttt{nu\_p\_oneThousandSixHundredSeventy\_two} & Gaps16621670 & PROVED & \checkmark & verified \\
\texttt{nu\_p\_oneThousandSixHundredSixtyEight} & Gaps16621670 & PROVED & \checkmark & verified \\
\texttt{nu\_p\_oneThousandSixHundredSixtyFour} & Gaps16621670 & PROVED & \checkmark & verified \\
\texttt{nu\_p\_oneThousandSixHundredSixtySix} & Gaps16621670 & PROVED & \checkmark & verified \\
\texttt{nu\_p\_oneThousandSixHundredSixtyTwo} & Gaps16621670 & PROVED & \checkmark & verified \\
\texttt{nu\_p\_oneThousandSixHundredSixtyTwo\_two} & Gaps16621670 & PROVED & \checkmark & verified \\
\texttt{singular\_series\_finite\_pos\_evenPair\_oneThousandSixHundredSeventy} & Gaps16621670 & PROVED & \checkmark & verified \\
\texttt{singular\_series\_finite\_pos\_evenPair\_oneThousandSixHundredSixtyEight} & Gaps16621670 & PROVED & \checkmark & verified \\
\texttt{singular\_series\_finite\_pos\_evenPair\_oneThousandSixHundredSixtyFour} & Gaps16621670 & PROVED & \checkmark & verified \\
\texttt{singular\_series\_finite\_pos\_evenPair\_oneThousandSixHundredSixtySix} & Gaps16621670 & PROVED & \checkmark & verified \\
\texttt{singular\_series\_finite\_pos\_evenPair\_oneThousandSixHundredSixtyTwo} & Gaps16621670 & PROVED & \checkmark & verified \\
\texttt{singular\_series\_pos\_evenPair\_oneThousandSixHundredSeventy} & Gaps16621670 & PROVED & \checkmark & verified \\
\texttt{singular\_series\_pos\_evenPair\_oneThousandSixHundredSixtyEight} & Gaps16621670 & PROVED & \checkmark & verified \\
\texttt{singular\_series\_pos\_evenPair\_oneThousandSixHundredSixtyFour} & Gaps16621670 & PROVED & \checkmark & verified \\
\texttt{singular\_series\_pos\_evenPair\_oneThousandSixHundredSixtySix} & Gaps16621670 & PROVED & \checkmark & verified \\
\texttt{singular\_series\_pos\_evenPair\_oneThousandSixHundredSixtyTwo} & Gaps16621670 & PROVED & \checkmark & verified \\
\texttt{evenPair\_card\_oneThousandSixHundredEighty} & Gaps16721680 & PROVED & \checkmark & verified \\
\texttt{evenPair\_card\_oneThousandSixHundredSeventyEight} & Gaps16721680 & PROVED & \checkmark & verified \\
\texttt{evenPair\_card\_oneThousandSixHundredSeventyFour} & Gaps16721680 & PROVED & \checkmark & verified \\
\texttt{evenPair\_card\_oneThousandSixHundredSeventySix} & Gaps16721680 & PROVED & \checkmark & verified \\
\texttt{evenPair\_card\_oneThousandSixHundredSeventyTwo} & Gaps16721680 & PROVED & \checkmark & verified \\
\texttt{isAdmissible\_evenPair\_oneThousandSixHundredEighty} & Gaps16721680 & PROVED & \checkmark & verified \\
\texttt{isAdmissible\_evenPair\_oneThousandSixHundredSeventyEight} & Gaps16721680 & PROVED & \checkmark & verified \\
\texttt{isAdmissible\_evenPair\_oneThousandSixHundredSeventyFour} & Gaps16721680 & PROVED & \checkmark & verified \\
\texttt{isAdmissible\_evenPair\_oneThousandSixHundredSeventySix} & Gaps16721680 & PROVED & \checkmark & verified \\
\texttt{isAdmissible\_evenPair\_oneThousandSixHundredSeventyTwo} & Gaps16721680 & PROVED & \checkmark & verified \\
\texttt{localFactor\_oneThousandSixHundredEighty\_two} & Gaps16721680 & PROVED & \checkmark & verified \\
\texttt{localFactor\_oneThousandSixHundredSeventyTwo\_two} & Gaps16721680 & PROVED & \checkmark & verified \\
\texttt{nu\_p\_oneThousandSixHundredEighty} & Gaps16721680 & PROVED & \checkmark & verified \\
\texttt{nu\_p\_oneThousandSixHundredEighty\_two} & Gaps16721680 & PROVED & \checkmark & verified \\
\texttt{nu\_p\_oneThousandSixHundredSeventyEight} & Gaps16721680 & PROVED & \checkmark & verified \\
\texttt{nu\_p\_oneThousandSixHundredSeventyFour} & Gaps16721680 & PROVED & \checkmark & verified \\
\texttt{nu\_p\_oneThousandSixHundredSeventySix} & Gaps16721680 & PROVED & \checkmark & verified \\
\texttt{nu\_p\_oneThousandSixHundredSeventyTwo} & Gaps16721680 & PROVED & \checkmark & verified \\
\texttt{nu\_p\_oneThousandSixHundredSeventyTwo\_two} & Gaps16721680 & PROVED & \checkmark & verified \\
\texttt{singular\_series\_finite\_pos\_evenPair\_oneThousandSixHundredEighty} & Gaps16721680 & PROVED & \checkmark & verified \\
\texttt{singular\_series\_finite\_pos\_evenPair\_oneThousandSixHundredSeventyEight} & Gaps16721680 & PROVED & \checkmark & verified \\
\texttt{singular\_series\_finite\_pos\_evenPair\_oneThousandSixHundredSeventyFour} & Gaps16721680 & PROVED & \checkmark & verified \\
\texttt{singular\_series\_finite\_pos\_evenPair\_oneThousandSixHundredSeventySix} & Gaps16721680 & PROVED & \checkmark & verified \\
\texttt{singular\_series\_finite\_pos\_evenPair\_oneThousandSixHundredSeventyTwo} & Gaps16721680 & PROVED & \checkmark & verified \\
\texttt{singular\_series\_pos\_evenPair\_oneThousandSixHundredEighty} & Gaps16721680 & PROVED & \checkmark & verified \\
\texttt{singular\_series\_pos\_evenPair\_oneThousandSixHundredSeventyEight} & Gaps16721680 & PROVED & \checkmark & verified \\
\texttt{singular\_series\_pos\_evenPair\_oneThousandSixHundredSeventyFour} & Gaps16721680 & PROVED & \checkmark & verified \\
\texttt{singular\_series\_pos\_evenPair\_oneThousandSixHundredSeventySix} & Gaps16721680 & PROVED & \checkmark & verified \\
\texttt{singular\_series\_pos\_evenPair\_oneThousandSixHundredSeventyTwo} & Gaps16721680 & PROVED & \checkmark & verified \\
\texttt{evenPair\_card\_oneThousandSixHundredEightyEight} & Gaps16821690 & PROVED & \checkmark & verified \\
\texttt{evenPair\_card\_oneThousandSixHundredEightyFour} & Gaps16821690 & PROVED & \checkmark & verified \\
\texttt{evenPair\_card\_oneThousandSixHundredEightySix} & Gaps16821690 & PROVED & \checkmark & verified \\
\texttt{evenPair\_card\_oneThousandSixHundredEightyTwo} & Gaps16821690 & PROVED & \checkmark & verified \\
\texttt{evenPair\_card\_oneThousandSixHundredNinety} & Gaps16821690 & PROVED & \checkmark & verified \\
\texttt{isAdmissible\_evenPair\_oneThousandSixHundredEightyEight} & Gaps16821690 & PROVED & \checkmark & verified \\
\texttt{isAdmissible\_evenPair\_oneThousandSixHundredEightyFour} & Gaps16821690 & PROVED & \checkmark & verified \\
\texttt{isAdmissible\_evenPair\_oneThousandSixHundredEightySix} & Gaps16821690 & PROVED & \checkmark & verified \\
\texttt{isAdmissible\_evenPair\_oneThousandSixHundredEightyTwo} & Gaps16821690 & PROVED & \checkmark & verified \\
\texttt{isAdmissible\_evenPair\_oneThousandSixHundredNinety} & Gaps16821690 & PROVED & \checkmark & verified \\
\texttt{localFactor\_oneThousandSixHundredEightyTwo\_two} & Gaps16821690 & PROVED & \checkmark & verified \\
\texttt{localFactor\_oneThousandSixHundredNinety\_two} & Gaps16821690 & PROVED & \checkmark & verified \\
\texttt{nu\_p\_oneThousandSixHundredEightyEight} & Gaps16821690 & PROVED & \checkmark & verified \\
\texttt{nu\_p\_oneThousandSixHundredEightyFour} & Gaps16821690 & PROVED & \checkmark & verified \\
\texttt{nu\_p\_oneThousandSixHundredEightySix} & Gaps16821690 & PROVED & \checkmark & verified \\
\texttt{nu\_p\_oneThousandSixHundredEightyTwo} & Gaps16821690 & PROVED & \checkmark & verified \\
\texttt{nu\_p\_oneThousandSixHundredEightyTwo\_two} & Gaps16821690 & PROVED & \checkmark & verified \\
\texttt{nu\_p\_oneThousandSixHundredNinety} & Gaps16821690 & PROVED & \checkmark & verified \\
\texttt{nu\_p\_oneThousandSixHundredNinety\_two} & Gaps16821690 & PROVED & \checkmark & verified \\
\texttt{singular\_series\_finite\_pos\_evenPair\_oneThousandSixHundredEightyEight} & Gaps16821690 & PROVED & \checkmark & verified \\
\texttt{singular\_series\_finite\_pos\_evenPair\_oneThousandSixHundredEightyFour} & Gaps16821690 & PROVED & \checkmark & verified \\
\texttt{singular\_series\_finite\_pos\_evenPair\_oneThousandSixHundredEightySix} & Gaps16821690 & PROVED & \checkmark & verified \\
\texttt{singular\_series\_finite\_pos\_evenPair\_oneThousandSixHundredEightyTwo} & Gaps16821690 & PROVED & \checkmark & verified \\
\texttt{singular\_series\_finite\_pos\_evenPair\_oneThousandSixHundredNinety} & Gaps16821690 & PROVED & \checkmark & verified \\
\texttt{singular\_series\_pos\_evenPair\_oneThousandSixHundredEightyEight} & Gaps16821690 & PROVED & \checkmark & verified \\
\texttt{singular\_series\_pos\_evenPair\_oneThousandSixHundredEightyFour} & Gaps16821690 & PROVED & \checkmark & verified \\
\texttt{singular\_series\_pos\_evenPair\_oneThousandSixHundredEightySix} & Gaps16821690 & PROVED & \checkmark & verified \\
\texttt{singular\_series\_pos\_evenPair\_oneThousandSixHundredEightyTwo} & Gaps16821690 & PROVED & \checkmark & verified \\
\texttt{singular\_series\_pos\_evenPair\_oneThousandSixHundredNinety} & Gaps16821690 & PROVED & \checkmark & verified \\
\texttt{evenPair\_card\_oneThousandSevenHundred} & Gaps16921700 & PROVED & \checkmark & verified \\
\texttt{evenPair\_card\_oneThousandSixHundredNinetyEight} & Gaps16921700 & PROVED & \checkmark & verified \\
\texttt{evenPair\_card\_oneThousandSixHundredNinetyFour} & Gaps16921700 & PROVED & \checkmark & verified \\
\texttt{evenPair\_card\_oneThousandSixHundredNinetySix} & Gaps16921700 & PROVED & \checkmark & verified \\
\texttt{evenPair\_card\_oneThousandSixHundredNinetyTwo} & Gaps16921700 & PROVED & \checkmark & verified \\
\texttt{isAdmissible\_evenPair\_oneThousandSevenHundred} & Gaps16921700 & PROVED & \checkmark & verified \\
\texttt{isAdmissible\_evenPair\_oneThousandSixHundredNinetyEight} & Gaps16921700 & PROVED & \checkmark & verified \\
\texttt{isAdmissible\_evenPair\_oneThousandSixHundredNinetyFour} & Gaps16921700 & PROVED & \checkmark & verified \\
\texttt{isAdmissible\_evenPair\_oneThousandSixHundredNinetySix} & Gaps16921700 & PROVED & \checkmark & verified \\
\texttt{isAdmissible\_evenPair\_oneThousandSixHundredNinetyTwo} & Gaps16921700 & PROVED & \checkmark & verified \\
\texttt{localFactor\_oneThousandSevenHundred\_two} & Gaps16921700 & PROVED & \checkmark & verified \\
\texttt{localFactor\_oneThousandSixHundredNinetyTwo\_two} & Gaps16921700 & PROVED & \checkmark & verified \\
\texttt{nu\_p\_oneThousandSevenHundred} & Gaps16921700 & PROVED & \checkmark & verified \\
\texttt{nu\_p\_oneThousandSevenHundred\_two} & Gaps16921700 & PROVED & \checkmark & verified \\
\texttt{nu\_p\_oneThousandSixHundredNinetyEight} & Gaps16921700 & PROVED & \checkmark & verified \\
\texttt{nu\_p\_oneThousandSixHundredNinetyFour} & Gaps16921700 & PROVED & \checkmark & verified \\
\texttt{nu\_p\_oneThousandSixHundredNinetySix} & Gaps16921700 & PROVED & \checkmark & verified \\
\texttt{nu\_p\_oneThousandSixHundredNinetyTwo} & Gaps16921700 & PROVED & \checkmark & verified \\
\texttt{nu\_p\_oneThousandSixHundredNinetyTwo\_two} & Gaps16921700 & PROVED & \checkmark & verified \\
\texttt{singular\_series\_finite\_pos\_evenPair\_oneThousandSevenHundred} & Gaps16921700 & PROVED & \checkmark & verified \\
\texttt{singular\_series\_finite\_pos\_evenPair\_oneThousandSixHundredNinetyEight} & Gaps16921700 & PROVED & \checkmark & verified \\
\texttt{singular\_series\_finite\_pos\_evenPair\_oneThousandSixHundredNinetyFour} & Gaps16921700 & PROVED & \checkmark & verified \\
\texttt{singular\_series\_finite\_pos\_evenPair\_oneThousandSixHundredNinetySix} & Gaps16921700 & PROVED & \checkmark & verified \\
\texttt{singular\_series\_finite\_pos\_evenPair\_oneThousandSixHundredNinetyTwo} & Gaps16921700 & PROVED & \checkmark & verified \\
\texttt{singular\_series\_pos\_evenPair\_oneThousandSevenHundred} & Gaps16921700 & PROVED & \checkmark & verified \\
\texttt{singular\_series\_pos\_evenPair\_oneThousandSixHundredNinetyEight} & Gaps16921700 & PROVED & \checkmark & verified \\
\texttt{singular\_series\_pos\_evenPair\_oneThousandSixHundredNinetyFour} & Gaps16921700 & PROVED & \checkmark & verified \\
\texttt{singular\_series\_pos\_evenPair\_oneThousandSixHundredNinetySix} & Gaps16921700 & PROVED & \checkmark & verified \\
\texttt{singular\_series\_pos\_evenPair\_oneThousandSixHundredNinetyTwo} & Gaps16921700 & PROVED & \checkmark & verified \\
\texttt{evenPair\_card\_oneThousandSevenHundredEight} & Gaps17021710 & PROVED & \checkmark & verified \\
\texttt{evenPair\_card\_oneThousandSevenHundredFour} & Gaps17021710 & PROVED & \checkmark & verified \\
\texttt{evenPair\_card\_oneThousandSevenHundredSix} & Gaps17021710 & PROVED & \checkmark & verified \\
\texttt{evenPair\_card\_oneThousandSevenHundredTen} & Gaps17021710 & PROVED & \checkmark & verified \\
\texttt{evenPair\_card\_oneThousandSevenHundredTwo} & Gaps17021710 & PROVED & \checkmark & verified \\
\texttt{isAdmissible\_evenPair\_oneThousandSevenHundredEight} & Gaps17021710 & PROVED & \checkmark & verified \\
\texttt{isAdmissible\_evenPair\_oneThousandSevenHundredFour} & Gaps17021710 & PROVED & \checkmark & verified \\
\texttt{isAdmissible\_evenPair\_oneThousandSevenHundredSix} & Gaps17021710 & PROVED & \checkmark & verified \\
\texttt{isAdmissible\_evenPair\_oneThousandSevenHundredTen} & Gaps17021710 & PROVED & \checkmark & verified \\
\texttt{isAdmissible\_evenPair\_oneThousandSevenHundredTwo} & Gaps17021710 & PROVED & \checkmark & verified \\
\texttt{localFactor\_oneThousandSevenHundredTen\_two} & Gaps17021710 & PROVED & \checkmark & verified \\
\texttt{localFactor\_oneThousandSevenHundredTwo\_two} & Gaps17021710 & PROVED & \checkmark & verified \\
\texttt{nu\_p\_oneThousandSevenHundredEight} & Gaps17021710 & PROVED & \checkmark & verified \\
\texttt{nu\_p\_oneThousandSevenHundredFour} & Gaps17021710 & PROVED & \checkmark & verified \\
\texttt{nu\_p\_oneThousandSevenHundredSix} & Gaps17021710 & PROVED & \checkmark & verified \\
\texttt{nu\_p\_oneThousandSevenHundredTen} & Gaps17021710 & PROVED & \checkmark & verified \\
\texttt{nu\_p\_oneThousandSevenHundredTen\_two} & Gaps17021710 & PROVED & \checkmark & verified \\
\texttt{nu\_p\_oneThousandSevenHundredTwo} & Gaps17021710 & PROVED & \checkmark & verified \\
\texttt{nu\_p\_oneThousandSevenHundredTwo\_two} & Gaps17021710 & PROVED & \checkmark & verified \\
\texttt{singular\_series\_finite\_pos\_evenPair\_oneThousandSevenHundredEight} & Gaps17021710 & PROVED & \checkmark & verified \\
\texttt{singular\_series\_finite\_pos\_evenPair\_oneThousandSevenHundredFour} & Gaps17021710 & PROVED & \checkmark & verified \\
\texttt{singular\_series\_finite\_pos\_evenPair\_oneThousandSevenHundredSix} & Gaps17021710 & PROVED & \checkmark & verified \\
\texttt{singular\_series\_finite\_pos\_evenPair\_oneThousandSevenHundredTen} & Gaps17021710 & PROVED & \checkmark & verified \\
\texttt{singular\_series\_finite\_pos\_evenPair\_oneThousandSevenHundredTwo} & Gaps17021710 & PROVED & \checkmark & verified \\
\texttt{singular\_series\_pos\_evenPair\_oneThousandSevenHundredEight} & Gaps17021710 & PROVED & \checkmark & verified \\
\texttt{singular\_series\_pos\_evenPair\_oneThousandSevenHundredFour} & Gaps17021710 & PROVED & \checkmark & verified \\
\texttt{singular\_series\_pos\_evenPair\_oneThousandSevenHundredSix} & Gaps17021710 & PROVED & \checkmark & verified \\
\texttt{singular\_series\_pos\_evenPair\_oneThousandSevenHundredTen} & Gaps17021710 & PROVED & \checkmark & verified \\
\texttt{singular\_series\_pos\_evenPair\_oneThousandSevenHundredTwo} & Gaps17021710 & PROVED & \checkmark & verified \\
\texttt{evenPair\_card\_oneThousandSevenHundredEighteen} & Gaps17121720 & PROVED & \checkmark & verified \\
\texttt{evenPair\_card\_oneThousandSevenHundredFourteen} & Gaps17121720 & PROVED & \checkmark & verified \\
\texttt{evenPair\_card\_oneThousandSevenHundredSixteen} & Gaps17121720 & PROVED & \checkmark & verified \\
\texttt{evenPair\_card\_oneThousandSevenHundredTwelve} & Gaps17121720 & PROVED & \checkmark & verified \\
\texttt{evenPair\_card\_oneThousandSevenHundredTwenty} & Gaps17121720 & PROVED & \checkmark & verified \\
\texttt{isAdmissible\_evenPair\_oneThousandSevenHundredEighteen} & Gaps17121720 & PROVED & \checkmark & verified \\
\texttt{isAdmissible\_evenPair\_oneThousandSevenHundredFourteen} & Gaps17121720 & PROVED & \checkmark & verified \\
\texttt{isAdmissible\_evenPair\_oneThousandSevenHundredSixteen} & Gaps17121720 & PROVED & \checkmark & verified \\
\texttt{isAdmissible\_evenPair\_oneThousandSevenHundredTwelve} & Gaps17121720 & PROVED & \checkmark & verified \\
\texttt{isAdmissible\_evenPair\_oneThousandSevenHundredTwenty} & Gaps17121720 & PROVED & \checkmark & verified \\
\texttt{localFactor\_oneThousandSevenHundredTwelve\_two} & Gaps17121720 & PROVED & \checkmark & verified \\
\texttt{localFactor\_oneThousandSevenHundredTwenty\_two} & Gaps17121720 & PROVED & \checkmark & verified \\
\texttt{nu\_p\_oneThousandSevenHundredEighteen} & Gaps17121720 & PROVED & \checkmark & verified \\
\texttt{nu\_p\_oneThousandSevenHundredFourteen} & Gaps17121720 & PROVED & \checkmark & verified \\
\texttt{nu\_p\_oneThousandSevenHundredSixteen} & Gaps17121720 & PROVED & \checkmark & verified \\
\texttt{nu\_p\_oneThousandSevenHundredTwelve} & Gaps17121720 & PROVED & \checkmark & verified \\
\texttt{nu\_p\_oneThousandSevenHundredTwelve\_two} & Gaps17121720 & PROVED & \checkmark & verified \\
\texttt{nu\_p\_oneThousandSevenHundredTwenty} & Gaps17121720 & PROVED & \checkmark & verified \\
\texttt{nu\_p\_oneThousandSevenHundredTwenty\_two} & Gaps17121720 & PROVED & \checkmark & verified \\
\texttt{singular\_series\_finite\_pos\_evenPair\_oneThousandSevenHundredEighteen} & Gaps17121720 & PROVED & \checkmark & verified \\
\texttt{singular\_series\_finite\_pos\_evenPair\_oneThousandSevenHundredFourteen} & Gaps17121720 & PROVED & \checkmark & verified \\
\texttt{singular\_series\_finite\_pos\_evenPair\_oneThousandSevenHundredSixteen} & Gaps17121720 & PROVED & \checkmark & verified \\
\texttt{singular\_series\_finite\_pos\_evenPair\_oneThousandSevenHundredTwelve} & Gaps17121720 & PROVED & \checkmark & verified \\
\texttt{singular\_series\_finite\_pos\_evenPair\_oneThousandSevenHundredTwenty} & Gaps17121720 & PROVED & \checkmark & verified \\
\texttt{singular\_series\_pos\_evenPair\_oneThousandSevenHundredEighteen} & Gaps17121720 & PROVED & \checkmark & verified \\
\texttt{singular\_series\_pos\_evenPair\_oneThousandSevenHundredFourteen} & Gaps17121720 & PROVED & \checkmark & verified \\
\texttt{singular\_series\_pos\_evenPair\_oneThousandSevenHundredSixteen} & Gaps17121720 & PROVED & \checkmark & verified \\
\texttt{singular\_series\_pos\_evenPair\_oneThousandSevenHundredTwelve} & Gaps17121720 & PROVED & \checkmark & verified \\
\texttt{singular\_series\_pos\_evenPair\_oneThousandSevenHundredTwenty} & Gaps17121720 & PROVED & \checkmark & verified \\
\texttt{evenPair\_card\_oneHundredEighty} & Gaps172180 & PROVED & \checkmark & verified \\
\texttt{evenPair\_card\_oneHundredSeventyEight} & Gaps172180 & PROVED & \checkmark & verified \\
\texttt{evenPair\_card\_oneHundredSeventyFour} & Gaps172180 & PROVED & \checkmark & verified \\
\texttt{evenPair\_card\_oneHundredSeventySix} & Gaps172180 & PROVED & \checkmark & verified \\
\texttt{evenPair\_card\_oneHundredSeventyTwo} & Gaps172180 & PROVED & \checkmark & verified \\
\texttt{isAdmissible\_evenPair\_oneHundredEighty} & Gaps172180 & PROVED & \checkmark & verified \\
\texttt{isAdmissible\_evenPair\_oneHundredSeventyEight} & Gaps172180 & PROVED & \checkmark & verified \\
\texttt{isAdmissible\_evenPair\_oneHundredSeventyFour} & Gaps172180 & PROVED & \checkmark & verified \\
\texttt{isAdmissible\_evenPair\_oneHundredSeventySix} & Gaps172180 & PROVED & \checkmark & verified \\
\texttt{isAdmissible\_evenPair\_oneHundredSeventyTwo} & Gaps172180 & PROVED & \checkmark & verified \\
\texttt{localFactor\_oneHundredEighty\_two} & Gaps172180 & PROVED & \checkmark & verified \\
\texttt{localFactor\_oneHundredSeventyTwo\_two} & Gaps172180 & PROVED & \checkmark & verified \\
\texttt{nu\_p\_oneHundredEighty} & Gaps172180 & PROVED & \checkmark & verified \\
\texttt{nu\_p\_oneHundredEighty\_two} & Gaps172180 & PROVED & \checkmark & verified \\
\texttt{nu\_p\_oneHundredSeventyEight} & Gaps172180 & PROVED & \checkmark & verified \\
\texttt{nu\_p\_oneHundredSeventyFour} & Gaps172180 & PROVED & \checkmark & verified \\
\texttt{nu\_p\_oneHundredSeventySix} & Gaps172180 & PROVED & \checkmark & verified \\
\texttt{nu\_p\_oneHundredSeventyTwo} & Gaps172180 & PROVED & \checkmark & verified \\
\texttt{nu\_p\_oneHundredSeventyTwo\_two} & Gaps172180 & PROVED & \checkmark & verified \\
\texttt{singular\_series\_finite\_pos\_evenPair\_oneHundredEighty} & Gaps172180 & PROVED & \checkmark & verified \\
\texttt{singular\_series\_finite\_pos\_evenPair\_oneHundredSeventyEight} & Gaps172180 & PROVED & \checkmark & verified \\
\texttt{singular\_series\_finite\_pos\_evenPair\_oneHundredSeventyFour} & Gaps172180 & PROVED & \checkmark & verified \\
\texttt{singular\_series\_finite\_pos\_evenPair\_oneHundredSeventySix} & Gaps172180 & PROVED & \checkmark & verified \\
\texttt{singular\_series\_finite\_pos\_evenPair\_oneHundredSeventyTwo} & Gaps172180 & PROVED & \checkmark & verified \\
\texttt{singular\_series\_pos\_evenPair\_oneHundredEighty} & Gaps172180 & PROVED & \checkmark & verified \\
\texttt{singular\_series\_pos\_evenPair\_oneHundredSeventyEight} & Gaps172180 & PROVED & \checkmark & verified \\
\texttt{singular\_series\_pos\_evenPair\_oneHundredSeventyFour} & Gaps172180 & PROVED & \checkmark & verified \\
\texttt{singular\_series\_pos\_evenPair\_oneHundredSeventySix} & Gaps172180 & PROVED & \checkmark & verified \\
\texttt{singular\_series\_pos\_evenPair\_oneHundredSeventyTwo} & Gaps172180 & PROVED & \checkmark & verified \\
\texttt{evenPair\_card\_oneThousandSevenHundredThirty} & Gaps17221730 & PROVED & \checkmark & verified \\
\texttt{evenPair\_card\_oneThousandSevenHundredTwentyEight} & Gaps17221730 & PROVED & \checkmark & verified \\
\texttt{evenPair\_card\_oneThousandSevenHundredTwentyFour} & Gaps17221730 & PROVED & \checkmark & verified \\
\texttt{evenPair\_card\_oneThousandSevenHundredTwentySix} & Gaps17221730 & PROVED & \checkmark & verified \\
\texttt{evenPair\_card\_oneThousandSevenHundredTwentyTwo} & Gaps17221730 & PROVED & \checkmark & verified \\
\texttt{isAdmissible\_evenPair\_oneThousandSevenHundredThirty} & Gaps17221730 & PROVED & \checkmark & verified \\
\texttt{isAdmissible\_evenPair\_oneThousandSevenHundredTwentyEight} & Gaps17221730 & PROVED & \checkmark & verified \\
\texttt{isAdmissible\_evenPair\_oneThousandSevenHundredTwentyFour} & Gaps17221730 & PROVED & \checkmark & verified \\
\texttt{isAdmissible\_evenPair\_oneThousandSevenHundredTwentySix} & Gaps17221730 & PROVED & \checkmark & verified \\
\texttt{isAdmissible\_evenPair\_oneThousandSevenHundredTwentyTwo} & Gaps17221730 & PROVED & \checkmark & verified \\
\texttt{localFactor\_oneThousandSevenHundredThirty\_two} & Gaps17221730 & PROVED & \checkmark & verified \\
\texttt{localFactor\_oneThousandSevenHundredTwentyTwo\_two} & Gaps17221730 & PROVED & \checkmark & verified \\
\texttt{nu\_p\_oneThousandSevenHundredThirty} & Gaps17221730 & PROVED & \checkmark & verified \\
\texttt{nu\_p\_oneThousandSevenHundredThirty\_two} & Gaps17221730 & PROVED & \checkmark & verified \\
\texttt{nu\_p\_oneThousandSevenHundredTwentyEight} & Gaps17221730 & PROVED & \checkmark & verified \\
\texttt{nu\_p\_oneThousandSevenHundredTwentyFour} & Gaps17221730 & PROVED & \checkmark & verified \\
\texttt{nu\_p\_oneThousandSevenHundredTwentySix} & Gaps17221730 & PROVED & \checkmark & verified \\
\texttt{nu\_p\_oneThousandSevenHundredTwentyTwo} & Gaps17221730 & PROVED & \checkmark & verified \\
\texttt{nu\_p\_oneThousandSevenHundredTwentyTwo\_two} & Gaps17221730 & PROVED & \checkmark & verified \\
\texttt{singular\_series\_finite\_pos\_evenPair\_oneThousandSevenHundredThirty} & Gaps17221730 & PROVED & \checkmark & verified \\
\texttt{singular\_series\_finite\_pos\_evenPair\_oneThousandSevenHundredTwentyEight} & Gaps17221730 & PROVED & \checkmark & verified \\
\texttt{singular\_series\_finite\_pos\_evenPair\_oneThousandSevenHundredTwentyFour} & Gaps17221730 & PROVED & \checkmark & verified \\
\texttt{singular\_series\_finite\_pos\_evenPair\_oneThousandSevenHundredTwentySix} & Gaps17221730 & PROVED & \checkmark & verified \\
\texttt{singular\_series\_finite\_pos\_evenPair\_oneThousandSevenHundredTwentyTwo} & Gaps17221730 & PROVED & \checkmark & verified \\
\texttt{singular\_series\_pos\_evenPair\_oneThousandSevenHundredThirty} & Gaps17221730 & PROVED & \checkmark & verified \\
\texttt{singular\_series\_pos\_evenPair\_oneThousandSevenHundredTwentyEight} & Gaps17221730 & PROVED & \checkmark & verified \\
\texttt{singular\_series\_pos\_evenPair\_oneThousandSevenHundredTwentyFour} & Gaps17221730 & PROVED & \checkmark & verified \\
\texttt{singular\_series\_pos\_evenPair\_oneThousandSevenHundredTwentySix} & Gaps17221730 & PROVED & \checkmark & verified \\
\texttt{singular\_series\_pos\_evenPair\_oneThousandSevenHundredTwentyTwo} & Gaps17221730 & PROVED & \checkmark & verified \\
\texttt{evenPair\_card\_oneThousandSevenHundredForty} & Gaps17321740 & PROVED & \checkmark & verified \\
\texttt{evenPair\_card\_oneThousandSevenHundredThirtyEight} & Gaps17321740 & PROVED & \checkmark & verified \\
\texttt{evenPair\_card\_oneThousandSevenHundredThirtyFour} & Gaps17321740 & PROVED & \checkmark & verified \\
\texttt{evenPair\_card\_oneThousandSevenHundredThirtySix} & Gaps17321740 & PROVED & \checkmark & verified \\
\texttt{evenPair\_card\_oneThousandSevenHundredThirtyTwo} & Gaps17321740 & PROVED & \checkmark & verified \\
\texttt{isAdmissible\_evenPair\_oneThousandSevenHundredForty} & Gaps17321740 & PROVED & \checkmark & verified \\
\texttt{isAdmissible\_evenPair\_oneThousandSevenHundredThirtyEight} & Gaps17321740 & PROVED & \checkmark & verified \\
\texttt{isAdmissible\_evenPair\_oneThousandSevenHundredThirtyFour} & Gaps17321740 & PROVED & \checkmark & verified \\
\texttt{isAdmissible\_evenPair\_oneThousandSevenHundredThirtySix} & Gaps17321740 & PROVED & \checkmark & verified \\
\texttt{isAdmissible\_evenPair\_oneThousandSevenHundredThirtyTwo} & Gaps17321740 & PROVED & \checkmark & verified \\
\texttt{localFactor\_oneThousandSevenHundredForty\_two} & Gaps17321740 & PROVED & \checkmark & verified \\
\texttt{localFactor\_oneThousandSevenHundredThirtyTwo\_two} & Gaps17321740 & PROVED & \checkmark & verified \\
\texttt{nu\_p\_oneThousandSevenHundredForty} & Gaps17321740 & PROVED & \checkmark & verified \\
\texttt{nu\_p\_oneThousandSevenHundredForty\_two} & Gaps17321740 & PROVED & \checkmark & verified \\
\texttt{nu\_p\_oneThousandSevenHundredThirtyEight} & Gaps17321740 & PROVED & \checkmark & verified \\
\texttt{nu\_p\_oneThousandSevenHundredThirtyFour} & Gaps17321740 & PROVED & \checkmark & verified \\
\texttt{nu\_p\_oneThousandSevenHundredThirtySix} & Gaps17321740 & PROVED & \checkmark & verified \\
\texttt{nu\_p\_oneThousandSevenHundredThirtyTwo} & Gaps17321740 & PROVED & \checkmark & verified \\
\texttt{nu\_p\_oneThousandSevenHundredThirtyTwo\_two} & Gaps17321740 & PROVED & \checkmark & verified \\
\texttt{singular\_series\_finite\_pos\_evenPair\_oneThousandSevenHundredForty} & Gaps17321740 & PROVED & \checkmark & verified \\
\texttt{singular\_series\_finite\_pos\_evenPair\_oneThousandSevenHundredThirtyEight} & Gaps17321740 & PROVED & \checkmark & verified \\
\texttt{singular\_series\_finite\_pos\_evenPair\_oneThousandSevenHundredThirtyFour} & Gaps17321740 & PROVED & \checkmark & verified \\
\texttt{singular\_series\_finite\_pos\_evenPair\_oneThousandSevenHundredThirtySix} & Gaps17321740 & PROVED & \checkmark & verified \\
\texttt{singular\_series\_finite\_pos\_evenPair\_oneThousandSevenHundredThirtyTwo} & Gaps17321740 & PROVED & \checkmark & verified \\
\texttt{singular\_series\_pos\_evenPair\_oneThousandSevenHundredForty} & Gaps17321740 & PROVED & \checkmark & verified \\
\texttt{singular\_series\_pos\_evenPair\_oneThousandSevenHundredThirtyEight} & Gaps17321740 & PROVED & \checkmark & verified \\
\texttt{singular\_series\_pos\_evenPair\_oneThousandSevenHundredThirtyFour} & Gaps17321740 & PROVED & \checkmark & verified \\
\texttt{singular\_series\_pos\_evenPair\_oneThousandSevenHundredThirtySix} & Gaps17321740 & PROVED & \checkmark & verified \\
\texttt{singular\_series\_pos\_evenPair\_oneThousandSevenHundredThirtyTwo} & Gaps17321740 & PROVED & \checkmark & verified \\
\texttt{evenPair\_card\_oneThousandSevenHundredFifty} & Gaps17421750 & PROVED & \checkmark & verified \\
\texttt{evenPair\_card\_oneThousandSevenHundredFortyEight} & Gaps17421750 & PROVED & \checkmark & verified \\
\texttt{evenPair\_card\_oneThousandSevenHundredFortyFour} & Gaps17421750 & PROVED & \checkmark & verified \\
\texttt{evenPair\_card\_oneThousandSevenHundredFortySix} & Gaps17421750 & PROVED & \checkmark & verified \\
\texttt{evenPair\_card\_oneThousandSevenHundredFortyTwo} & Gaps17421750 & PROVED & \checkmark & verified \\
\texttt{isAdmissible\_evenPair\_oneThousandSevenHundredFifty} & Gaps17421750 & PROVED & \checkmark & verified \\
\texttt{isAdmissible\_evenPair\_oneThousandSevenHundredFortyEight} & Gaps17421750 & PROVED & \checkmark & verified \\
\texttt{isAdmissible\_evenPair\_oneThousandSevenHundredFortyFour} & Gaps17421750 & PROVED & \checkmark & verified \\
\texttt{isAdmissible\_evenPair\_oneThousandSevenHundredFortySix} & Gaps17421750 & PROVED & \checkmark & verified \\
\texttt{isAdmissible\_evenPair\_oneThousandSevenHundredFortyTwo} & Gaps17421750 & PROVED & \checkmark & verified \\
\texttt{localFactor\_oneThousandSevenHundredFifty\_two} & Gaps17421750 & PROVED & \checkmark & verified \\
\texttt{localFactor\_oneThousandSevenHundredFortyTwo\_two} & Gaps17421750 & PROVED & \checkmark & verified \\
\texttt{nu\_p\_oneThousandSevenHundredFifty} & Gaps17421750 & PROVED & \checkmark & verified \\
\texttt{nu\_p\_oneThousandSevenHundredFifty\_two} & Gaps17421750 & PROVED & \checkmark & verified \\
\texttt{nu\_p\_oneThousandSevenHundredFortyEight} & Gaps17421750 & PROVED & \checkmark & verified \\
\texttt{nu\_p\_oneThousandSevenHundredFortyFour} & Gaps17421750 & PROVED & \checkmark & verified \\
\texttt{nu\_p\_oneThousandSevenHundredFortySix} & Gaps17421750 & PROVED & \checkmark & verified \\
\texttt{nu\_p\_oneThousandSevenHundredFortyTwo} & Gaps17421750 & PROVED & \checkmark & verified \\
\texttt{nu\_p\_oneThousandSevenHundredFortyTwo\_two} & Gaps17421750 & PROVED & \checkmark & verified \\
\texttt{singular\_series\_finite\_pos\_evenPair\_oneThousandSevenHundredFifty} & Gaps17421750 & PROVED & \checkmark & verified \\
\texttt{singular\_series\_finite\_pos\_evenPair\_oneThousandSevenHundredFortyEight} & Gaps17421750 & PROVED & \checkmark & verified \\
\texttt{singular\_series\_finite\_pos\_evenPair\_oneThousandSevenHundredFortyFour} & Gaps17421750 & PROVED & \checkmark & verified \\
\texttt{singular\_series\_finite\_pos\_evenPair\_oneThousandSevenHundredFortySix} & Gaps17421750 & PROVED & \checkmark & verified \\
\texttt{singular\_series\_finite\_pos\_evenPair\_oneThousandSevenHundredFortyTwo} & Gaps17421750 & PROVED & \checkmark & verified \\
\texttt{singular\_series\_pos\_evenPair\_oneThousandSevenHundredFifty} & Gaps17421750 & PROVED & \checkmark & verified \\
\texttt{singular\_series\_pos\_evenPair\_oneThousandSevenHundredFortyEight} & Gaps17421750 & PROVED & \checkmark & verified \\
\texttt{singular\_series\_pos\_evenPair\_oneThousandSevenHundredFortyFour} & Gaps17421750 & PROVED & \checkmark & verified \\
\texttt{singular\_series\_pos\_evenPair\_oneThousandSevenHundredFortySix} & Gaps17421750 & PROVED & \checkmark & verified \\
\texttt{singular\_series\_pos\_evenPair\_oneThousandSevenHundredFortyTwo} & Gaps17421750 & PROVED & \checkmark & verified \\
\texttt{evenPair\_card\_oneThousandSevenHundredFiftyEight} & Gaps17521760 & PROVED & \checkmark & verified \\
\texttt{evenPair\_card\_oneThousandSevenHundredFiftyFour} & Gaps17521760 & PROVED & \checkmark & verified \\
\texttt{evenPair\_card\_oneThousandSevenHundredFiftySix} & Gaps17521760 & PROVED & \checkmark & verified \\
\texttt{evenPair\_card\_oneThousandSevenHundredFiftyTwo} & Gaps17521760 & PROVED & \checkmark & verified \\
\texttt{evenPair\_card\_oneThousandSevenHundredSixty} & Gaps17521760 & PROVED & \checkmark & verified \\
\texttt{isAdmissible\_evenPair\_oneThousandSevenHundredFiftyEight} & Gaps17521760 & PROVED & \checkmark & verified \\
\texttt{isAdmissible\_evenPair\_oneThousandSevenHundredFiftyFour} & Gaps17521760 & PROVED & \checkmark & verified \\
\texttt{isAdmissible\_evenPair\_oneThousandSevenHundredFiftySix} & Gaps17521760 & PROVED & \checkmark & verified \\
\texttt{isAdmissible\_evenPair\_oneThousandSevenHundredFiftyTwo} & Gaps17521760 & PROVED & \checkmark & verified \\
\texttt{isAdmissible\_evenPair\_oneThousandSevenHundredSixty} & Gaps17521760 & PROVED & \checkmark & verified \\
\texttt{localFactor\_oneThousandSevenHundredFiftyTwo\_two} & Gaps17521760 & PROVED & \checkmark & verified \\
\texttt{localFactor\_oneThousandSevenHundredSixty\_two} & Gaps17521760 & PROVED & \checkmark & verified \\
\texttt{nu\_p\_oneThousandSevenHundredFiftyEight} & Gaps17521760 & PROVED & \checkmark & verified \\
\texttt{nu\_p\_oneThousandSevenHundredFiftyFour} & Gaps17521760 & PROVED & \checkmark & verified \\
\texttt{nu\_p\_oneThousandSevenHundredFiftySix} & Gaps17521760 & PROVED & \checkmark & verified \\
\texttt{nu\_p\_oneThousandSevenHundredFiftyTwo} & Gaps17521760 & PROVED & \checkmark & verified \\
\texttt{nu\_p\_oneThousandSevenHundredFiftyTwo\_two} & Gaps17521760 & PROVED & \checkmark & verified \\
\texttt{nu\_p\_oneThousandSevenHundredSixty} & Gaps17521760 & PROVED & \checkmark & verified \\
\texttt{nu\_p\_oneThousandSevenHundredSixty\_two} & Gaps17521760 & PROVED & \checkmark & verified \\
\texttt{singular\_series\_finite\_pos\_evenPair\_oneThousandSevenHundredFiftyEight} & Gaps17521760 & PROVED & \checkmark & verified \\
\texttt{singular\_series\_finite\_pos\_evenPair\_oneThousandSevenHundredFiftyFour} & Gaps17521760 & PROVED & \checkmark & verified \\
\texttt{singular\_series\_finite\_pos\_evenPair\_oneThousandSevenHundredFiftySix} & Gaps17521760 & PROVED & \checkmark & verified \\
\texttt{singular\_series\_finite\_pos\_evenPair\_oneThousandSevenHundredFiftyTwo} & Gaps17521760 & PROVED & \checkmark & verified \\
\texttt{singular\_series\_finite\_pos\_evenPair\_oneThousandSevenHundredSixty} & Gaps17521760 & PROVED & \checkmark & verified \\
\texttt{singular\_series\_pos\_evenPair\_oneThousandSevenHundredFiftyEight} & Gaps17521760 & PROVED & \checkmark & verified \\
\texttt{singular\_series\_pos\_evenPair\_oneThousandSevenHundredFiftyFour} & Gaps17521760 & PROVED & \checkmark & verified \\
\texttt{singular\_series\_pos\_evenPair\_oneThousandSevenHundredFiftySix} & Gaps17521760 & PROVED & \checkmark & verified \\
\texttt{singular\_series\_pos\_evenPair\_oneThousandSevenHundredFiftyTwo} & Gaps17521760 & PROVED & \checkmark & verified \\
\texttt{singular\_series\_pos\_evenPair\_oneThousandSevenHundredSixty} & Gaps17521760 & PROVED & \checkmark & verified \\
\texttt{evenPair\_card\_oneThousandSevenHundredSeventy} & Gaps17621770 & PROVED & \checkmark & verified \\
\texttt{evenPair\_card\_oneThousandSevenHundredSixtyEight} & Gaps17621770 & PROVED & \checkmark & verified \\
\texttt{evenPair\_card\_oneThousandSevenHundredSixtyFour} & Gaps17621770 & PROVED & \checkmark & verified \\
\texttt{evenPair\_card\_oneThousandSevenHundredSixtySix} & Gaps17621770 & PROVED & \checkmark & verified \\
\texttt{evenPair\_card\_oneThousandSevenHundredSixtyTwo} & Gaps17621770 & PROVED & \checkmark & verified \\
\texttt{isAdmissible\_evenPair\_oneThousandSevenHundredSeventy} & Gaps17621770 & PROVED & \checkmark & verified \\
\texttt{isAdmissible\_evenPair\_oneThousandSevenHundredSixtyEight} & Gaps17621770 & PROVED & \checkmark & verified \\
\texttt{isAdmissible\_evenPair\_oneThousandSevenHundredSixtyFour} & Gaps17621770 & PROVED & \checkmark & verified \\
\texttt{isAdmissible\_evenPair\_oneThousandSevenHundredSixtySix} & Gaps17621770 & PROVED & \checkmark & verified \\
\texttt{isAdmissible\_evenPair\_oneThousandSevenHundredSixtyTwo} & Gaps17621770 & PROVED & \checkmark & verified \\
\texttt{localFactor\_oneThousandSevenHundredSeventy\_two} & Gaps17621770 & PROVED & \checkmark & verified \\
\texttt{localFactor\_oneThousandSevenHundredSixtyTwo\_two} & Gaps17621770 & PROVED & \checkmark & verified \\
\texttt{nu\_p\_oneThousandSevenHundredSeventy} & Gaps17621770 & PROVED & \checkmark & verified \\
\texttt{nu\_p\_oneThousandSevenHundredSeventy\_two} & Gaps17621770 & PROVED & \checkmark & verified \\
\texttt{nu\_p\_oneThousandSevenHundredSixtyEight} & Gaps17621770 & PROVED & \checkmark & verified \\
\texttt{nu\_p\_oneThousandSevenHundredSixtyFour} & Gaps17621770 & PROVED & \checkmark & verified \\
\texttt{nu\_p\_oneThousandSevenHundredSixtySix} & Gaps17621770 & PROVED & \checkmark & verified \\
\texttt{nu\_p\_oneThousandSevenHundredSixtyTwo} & Gaps17621770 & PROVED & \checkmark & verified \\
\texttt{nu\_p\_oneThousandSevenHundredSixtyTwo\_two} & Gaps17621770 & PROVED & \checkmark & verified \\
\texttt{singular\_series\_finite\_pos\_evenPair\_oneThousandSevenHundredSeventy} & Gaps17621770 & PROVED & \checkmark & verified \\
\texttt{singular\_series\_finite\_pos\_evenPair\_oneThousandSevenHundredSixtyEight} & Gaps17621770 & PROVED & \checkmark & verified \\
\texttt{singular\_series\_finite\_pos\_evenPair\_oneThousandSevenHundredSixtyFour} & Gaps17621770 & PROVED & \checkmark & verified \\
\texttt{singular\_series\_finite\_pos\_evenPair\_oneThousandSevenHundredSixtySix} & Gaps17621770 & PROVED & \checkmark & verified \\
\texttt{singular\_series\_finite\_pos\_evenPair\_oneThousandSevenHundredSixtyTwo} & Gaps17621770 & PROVED & \checkmark & verified \\
\texttt{singular\_series\_pos\_evenPair\_oneThousandSevenHundredSeventy} & Gaps17621770 & PROVED & \checkmark & verified \\
\texttt{singular\_series\_pos\_evenPair\_oneThousandSevenHundredSixtyEight} & Gaps17621770 & PROVED & \checkmark & verified \\
\texttt{singular\_series\_pos\_evenPair\_oneThousandSevenHundredSixtyFour} & Gaps17621770 & PROVED & \checkmark & verified \\
\texttt{singular\_series\_pos\_evenPair\_oneThousandSevenHundredSixtySix} & Gaps17621770 & PROVED & \checkmark & verified \\
\texttt{singular\_series\_pos\_evenPair\_oneThousandSevenHundredSixtyTwo} & Gaps17621770 & PROVED & \checkmark & verified \\
\texttt{evenPair\_card\_oneThousandSevenHundredEighty} & Gaps17721780 & PROVED & \checkmark & verified \\
\texttt{evenPair\_card\_oneThousandSevenHundredSeventyEight} & Gaps17721780 & PROVED & \checkmark & verified \\
\texttt{evenPair\_card\_oneThousandSevenHundredSeventyFour} & Gaps17721780 & PROVED & \checkmark & verified \\
\texttt{evenPair\_card\_oneThousandSevenHundredSeventySix} & Gaps17721780 & PROVED & \checkmark & verified \\
\texttt{evenPair\_card\_oneThousandSevenHundredSeventyTwo} & Gaps17721780 & PROVED & \checkmark & verified \\
\texttt{isAdmissible\_evenPair\_oneThousandSevenHundredEighty} & Gaps17721780 & PROVED & \checkmark & verified \\
\texttt{isAdmissible\_evenPair\_oneThousandSevenHundredSeventyEight} & Gaps17721780 & PROVED & \checkmark & verified \\
\texttt{isAdmissible\_evenPair\_oneThousandSevenHundredSeventyFour} & Gaps17721780 & PROVED & \checkmark & verified \\
\texttt{isAdmissible\_evenPair\_oneThousandSevenHundredSeventySix} & Gaps17721780 & PROVED & \checkmark & verified \\
\texttt{isAdmissible\_evenPair\_oneThousandSevenHundredSeventyTwo} & Gaps17721780 & PROVED & \checkmark & verified \\
\texttt{localFactor\_oneThousandSevenHundredEighty\_two} & Gaps17721780 & PROVED & \checkmark & verified \\
\texttt{localFactor\_oneThousandSevenHundredSeventyTwo\_two} & Gaps17721780 & PROVED & \checkmark & verified \\
\texttt{nu\_p\_oneThousandSevenHundredEighty} & Gaps17721780 & PROVED & \checkmark & verified \\
\texttt{nu\_p\_oneThousandSevenHundredEighty\_two} & Gaps17721780 & PROVED & \checkmark & verified \\
\texttt{nu\_p\_oneThousandSevenHundredSeventyEight} & Gaps17721780 & PROVED & \checkmark & verified \\
\texttt{nu\_p\_oneThousandSevenHundredSeventyFour} & Gaps17721780 & PROVED & \checkmark & verified \\
\texttt{nu\_p\_oneThousandSevenHundredSeventySix} & Gaps17721780 & PROVED & \checkmark & verified \\
\texttt{nu\_p\_oneThousandSevenHundredSeventyTwo} & Gaps17721780 & PROVED & \checkmark & verified \\
\texttt{nu\_p\_oneThousandSevenHundredSeventyTwo\_two} & Gaps17721780 & PROVED & \checkmark & verified \\
\texttt{singular\_series\_finite\_pos\_evenPair\_oneThousandSevenHundredEighty} & Gaps17721780 & PROVED & \checkmark & verified \\
\texttt{singular\_series\_finite\_pos\_evenPair\_oneThousandSevenHundredSeventyEight} & Gaps17721780 & PROVED & \checkmark & verified \\
\texttt{singular\_series\_finite\_pos\_evenPair\_oneThousandSevenHundredSeventyFour} & Gaps17721780 & PROVED & \checkmark & verified \\
\texttt{singular\_series\_finite\_pos\_evenPair\_oneThousandSevenHundredSeventySix} & Gaps17721780 & PROVED & \checkmark & verified \\
\texttt{singular\_series\_finite\_pos\_evenPair\_oneThousandSevenHundredSeventyTwo} & Gaps17721780 & PROVED & \checkmark & verified \\
\texttt{singular\_series\_pos\_evenPair\_oneThousandSevenHundredEighty} & Gaps17721780 & PROVED & \checkmark & verified \\
\texttt{singular\_series\_pos\_evenPair\_oneThousandSevenHundredSeventyEight} & Gaps17721780 & PROVED & \checkmark & verified \\
\texttt{singular\_series\_pos\_evenPair\_oneThousandSevenHundredSeventyFour} & Gaps17721780 & PROVED & \checkmark & verified \\
\texttt{singular\_series\_pos\_evenPair\_oneThousandSevenHundredSeventySix} & Gaps17721780 & PROVED & \checkmark & verified \\
\texttt{singular\_series\_pos\_evenPair\_oneThousandSevenHundredSeventyTwo} & Gaps17721780 & PROVED & \checkmark & verified \\
\texttt{evenPair\_card\_oneThousandSevenHundredEightyEight} & Gaps17821790 & PROVED & \checkmark & verified \\
\texttt{evenPair\_card\_oneThousandSevenHundredEightyFour} & Gaps17821790 & PROVED & \checkmark & verified \\
\texttt{evenPair\_card\_oneThousandSevenHundredEightySix} & Gaps17821790 & PROVED & \checkmark & verified \\
\texttt{evenPair\_card\_oneThousandSevenHundredEightyTwo} & Gaps17821790 & PROVED & \checkmark & verified \\
\texttt{evenPair\_card\_oneThousandSevenHundredNinety} & Gaps17821790 & PROVED & \checkmark & verified \\
\texttt{isAdmissible\_evenPair\_oneThousandSevenHundredEightyEight} & Gaps17821790 & PROVED & \checkmark & verified \\
\texttt{isAdmissible\_evenPair\_oneThousandSevenHundredEightyFour} & Gaps17821790 & PROVED & \checkmark & verified \\
\texttt{isAdmissible\_evenPair\_oneThousandSevenHundredEightySix} & Gaps17821790 & PROVED & \checkmark & verified \\
\texttt{isAdmissible\_evenPair\_oneThousandSevenHundredEightyTwo} & Gaps17821790 & PROVED & \checkmark & verified \\
\texttt{isAdmissible\_evenPair\_oneThousandSevenHundredNinety} & Gaps17821790 & PROVED & \checkmark & verified \\
\texttt{localFactor\_oneThousandSevenHundredEightyTwo\_two} & Gaps17821790 & PROVED & \checkmark & verified \\
\texttt{localFactor\_oneThousandSevenHundredNinety\_two} & Gaps17821790 & PROVED & \checkmark & verified \\
\texttt{nu\_p\_oneThousandSevenHundredEightyEight} & Gaps17821790 & PROVED & \checkmark & verified \\
\texttt{nu\_p\_oneThousandSevenHundredEightyFour} & Gaps17821790 & PROVED & \checkmark & verified \\
\texttt{nu\_p\_oneThousandSevenHundredEightySix} & Gaps17821790 & PROVED & \checkmark & verified \\
\texttt{nu\_p\_oneThousandSevenHundredEightyTwo} & Gaps17821790 & PROVED & \checkmark & verified \\
\texttt{nu\_p\_oneThousandSevenHundredEightyTwo\_two} & Gaps17821790 & PROVED & \checkmark & verified \\
\texttt{nu\_p\_oneThousandSevenHundredNinety} & Gaps17821790 & PROVED & \checkmark & verified \\
\texttt{nu\_p\_oneThousandSevenHundredNinety\_two} & Gaps17821790 & PROVED & \checkmark & verified \\
\texttt{singular\_series\_finite\_pos\_evenPair\_oneThousandSevenHundredEightyEight} & Gaps17821790 & PROVED & \checkmark & verified \\
\texttt{singular\_series\_finite\_pos\_evenPair\_oneThousandSevenHundredEightyFour} & Gaps17821790 & PROVED & \checkmark & verified \\
\texttt{singular\_series\_finite\_pos\_evenPair\_oneThousandSevenHundredEightySix} & Gaps17821790 & PROVED & \checkmark & verified \\
\texttt{singular\_series\_finite\_pos\_evenPair\_oneThousandSevenHundredEightyTwo} & Gaps17821790 & PROVED & \checkmark & verified \\
\texttt{singular\_series\_finite\_pos\_evenPair\_oneThousandSevenHundredNinety} & Gaps17821790 & PROVED & \checkmark & verified \\
\texttt{singular\_series\_pos\_evenPair\_oneThousandSevenHundredEightyEight} & Gaps17821790 & PROVED & \checkmark & verified \\
\texttt{singular\_series\_pos\_evenPair\_oneThousandSevenHundredEightyFour} & Gaps17821790 & PROVED & \checkmark & verified \\
\texttt{singular\_series\_pos\_evenPair\_oneThousandSevenHundredEightySix} & Gaps17821790 & PROVED & \checkmark & verified \\
\texttt{singular\_series\_pos\_evenPair\_oneThousandSevenHundredEightyTwo} & Gaps17821790 & PROVED & \checkmark & verified \\
\texttt{singular\_series\_pos\_evenPair\_oneThousandSevenHundredNinety} & Gaps17821790 & PROVED & \checkmark & verified \\
\texttt{evenPair\_card\_oneThousandEightHundred} & Gaps17921800 & PROVED & \checkmark & verified \\
\texttt{evenPair\_card\_oneThousandSevenHundredNinetyEight} & Gaps17921800 & PROVED & \checkmark & verified \\
\texttt{evenPair\_card\_oneThousandSevenHundredNinetyFour} & Gaps17921800 & PROVED & \checkmark & verified \\
\texttt{evenPair\_card\_oneThousandSevenHundredNinetySix} & Gaps17921800 & PROVED & \checkmark & verified \\
\texttt{evenPair\_card\_oneThousandSevenHundredNinetyTwo} & Gaps17921800 & PROVED & \checkmark & verified \\
\texttt{isAdmissible\_evenPair\_oneThousandEightHundred} & Gaps17921800 & PROVED & \checkmark & verified \\
\texttt{isAdmissible\_evenPair\_oneThousandSevenHundredNinetyEight} & Gaps17921800 & PROVED & \checkmark & verified \\
\texttt{isAdmissible\_evenPair\_oneThousandSevenHundredNinetyFour} & Gaps17921800 & PROVED & \checkmark & verified \\
\texttt{isAdmissible\_evenPair\_oneThousandSevenHundredNinetySix} & Gaps17921800 & PROVED & \checkmark & verified \\
\texttt{isAdmissible\_evenPair\_oneThousandSevenHundredNinetyTwo} & Gaps17921800 & PROVED & \checkmark & verified \\
\texttt{localFactor\_oneThousandEightHundred\_two} & Gaps17921800 & PROVED & \checkmark & verified \\
\texttt{localFactor\_oneThousandSevenHundredNinetyTwo\_two} & Gaps17921800 & PROVED & \checkmark & verified \\
\texttt{nu\_p\_oneThousandEightHundred} & Gaps17921800 & PROVED & \checkmark & verified \\
\texttt{nu\_p\_oneThousandEightHundred\_two} & Gaps17921800 & PROVED & \checkmark & verified \\
\texttt{nu\_p\_oneThousandSevenHundredNinetyEight} & Gaps17921800 & PROVED & \checkmark & verified \\
\texttt{nu\_p\_oneThousandSevenHundredNinetyFour} & Gaps17921800 & PROVED & \checkmark & verified \\
\texttt{nu\_p\_oneThousandSevenHundredNinetySix} & Gaps17921800 & PROVED & \checkmark & verified \\
\texttt{nu\_p\_oneThousandSevenHundredNinetyTwo} & Gaps17921800 & PROVED & \checkmark & verified \\
\texttt{nu\_p\_oneThousandSevenHundredNinetyTwo\_two} & Gaps17921800 & PROVED & \checkmark & verified \\
\texttt{singular\_series\_finite\_pos\_evenPair\_oneThousandEightHundred} & Gaps17921800 & PROVED & \checkmark & verified \\
\texttt{singular\_series\_finite\_pos\_evenPair\_oneThousandSevenHundredNinetyEight} & Gaps17921800 & PROVED & \checkmark & verified \\
\texttt{singular\_series\_finite\_pos\_evenPair\_oneThousandSevenHundredNinetyFour} & Gaps17921800 & PROVED & \checkmark & verified \\
\texttt{singular\_series\_finite\_pos\_evenPair\_oneThousandSevenHundredNinetySix} & Gaps17921800 & PROVED & \checkmark & verified \\
\texttt{singular\_series\_finite\_pos\_evenPair\_oneThousandSevenHundredNinetyTwo} & Gaps17921800 & PROVED & \checkmark & verified \\
\texttt{singular\_series\_pos\_evenPair\_oneThousandEightHundred} & Gaps17921800 & PROVED & \checkmark & verified \\
\texttt{singular\_series\_pos\_evenPair\_oneThousandSevenHundredNinetyEight} & Gaps17921800 & PROVED & \checkmark & verified \\
\texttt{singular\_series\_pos\_evenPair\_oneThousandSevenHundredNinetyFour} & Gaps17921800 & PROVED & \checkmark & verified \\
\texttt{singular\_series\_pos\_evenPair\_oneThousandSevenHundredNinetySix} & Gaps17921800 & PROVED & \checkmark & verified \\
\texttt{singular\_series\_pos\_evenPair\_oneThousandSevenHundredNinetyTwo} & Gaps17921800 & PROVED & \checkmark & verified \\
\texttt{evenPair\_card\_oneThousandEightHundredEight} & Gaps18021810 & PROVED & \checkmark & verified \\
\texttt{evenPair\_card\_oneThousandEightHundredFour} & Gaps18021810 & PROVED & \checkmark & verified \\
\texttt{evenPair\_card\_oneThousandEightHundredSix} & Gaps18021810 & PROVED & \checkmark & verified \\
\texttt{evenPair\_card\_oneThousandEightHundredTen} & Gaps18021810 & PROVED & \checkmark & verified \\
\texttt{evenPair\_card\_oneThousandEightHundredTwo} & Gaps18021810 & PROVED & \checkmark & verified \\
\texttt{isAdmissible\_evenPair\_oneThousandEightHundredEight} & Gaps18021810 & PROVED & \checkmark & verified \\
\texttt{isAdmissible\_evenPair\_oneThousandEightHundredFour} & Gaps18021810 & PROVED & \checkmark & verified \\
\texttt{isAdmissible\_evenPair\_oneThousandEightHundredSix} & Gaps18021810 & PROVED & \checkmark & verified \\
\texttt{isAdmissible\_evenPair\_oneThousandEightHundredTen} & Gaps18021810 & PROVED & \checkmark & verified \\
\texttt{isAdmissible\_evenPair\_oneThousandEightHundredTwo} & Gaps18021810 & PROVED & \checkmark & verified \\
\texttt{localFactor\_oneThousandEightHundredTen\_two} & Gaps18021810 & PROVED & \checkmark & verified \\
\texttt{localFactor\_oneThousandEightHundredTwo\_two} & Gaps18021810 & PROVED & \checkmark & verified \\
\texttt{nu\_p\_oneThousandEightHundredEight} & Gaps18021810 & PROVED & \checkmark & verified \\
\texttt{nu\_p\_oneThousandEightHundredFour} & Gaps18021810 & PROVED & \checkmark & verified \\
\texttt{nu\_p\_oneThousandEightHundredSix} & Gaps18021810 & PROVED & \checkmark & verified \\
\texttt{nu\_p\_oneThousandEightHundredTen} & Gaps18021810 & PROVED & \checkmark & verified \\
\texttt{nu\_p\_oneThousandEightHundredTen\_two} & Gaps18021810 & PROVED & \checkmark & verified \\
\texttt{nu\_p\_oneThousandEightHundredTwo} & Gaps18021810 & PROVED & \checkmark & verified \\
\texttt{nu\_p\_oneThousandEightHundredTwo\_two} & Gaps18021810 & PROVED & \checkmark & verified \\
\texttt{singular\_series\_finite\_pos\_evenPair\_oneThousandEightHundredEight} & Gaps18021810 & PROVED & \checkmark & verified \\
\texttt{singular\_series\_finite\_pos\_evenPair\_oneThousandEightHundredFour} & Gaps18021810 & PROVED & \checkmark & verified \\
\texttt{singular\_series\_finite\_pos\_evenPair\_oneThousandEightHundredSix} & Gaps18021810 & PROVED & \checkmark & verified \\
\texttt{singular\_series\_finite\_pos\_evenPair\_oneThousandEightHundredTen} & Gaps18021810 & PROVED & \checkmark & verified \\
\texttt{singular\_series\_finite\_pos\_evenPair\_oneThousandEightHundredTwo} & Gaps18021810 & PROVED & \checkmark & verified \\
\texttt{singular\_series\_pos\_evenPair\_oneThousandEightHundredEight} & Gaps18021810 & PROVED & \checkmark & verified \\
\texttt{singular\_series\_pos\_evenPair\_oneThousandEightHundredFour} & Gaps18021810 & PROVED & \checkmark & verified \\
\texttt{singular\_series\_pos\_evenPair\_oneThousandEightHundredSix} & Gaps18021810 & PROVED & \checkmark & verified \\
\texttt{singular\_series\_pos\_evenPair\_oneThousandEightHundredTen} & Gaps18021810 & PROVED & \checkmark & verified \\
\texttt{singular\_series\_pos\_evenPair\_oneThousandEightHundredTwo} & Gaps18021810 & PROVED & \checkmark & verified \\
\texttt{evenPair\_card\_oneThousandEightHundredEighteen} & Gaps18121820 & PROVED & \checkmark & verified \\
\texttt{evenPair\_card\_oneThousandEightHundredFourteen} & Gaps18121820 & PROVED & \checkmark & verified \\
\texttt{evenPair\_card\_oneThousandEightHundredSixteen} & Gaps18121820 & PROVED & \checkmark & verified \\
\texttt{evenPair\_card\_oneThousandEightHundredTwelve} & Gaps18121820 & PROVED & \checkmark & verified \\
\texttt{evenPair\_card\_oneThousandEightHundredTwenty} & Gaps18121820 & PROVED & \checkmark & verified \\
\texttt{isAdmissible\_evenPair\_oneThousandEightHundredEighteen} & Gaps18121820 & PROVED & \checkmark & verified \\
\texttt{isAdmissible\_evenPair\_oneThousandEightHundredFourteen} & Gaps18121820 & PROVED & \checkmark & verified \\
\texttt{isAdmissible\_evenPair\_oneThousandEightHundredSixteen} & Gaps18121820 & PROVED & \checkmark & verified \\
\texttt{isAdmissible\_evenPair\_oneThousandEightHundredTwelve} & Gaps18121820 & PROVED & \checkmark & verified \\
\texttt{isAdmissible\_evenPair\_oneThousandEightHundredTwenty} & Gaps18121820 & PROVED & \checkmark & verified \\
\texttt{localFactor\_oneThousandEightHundredTwelve\_two} & Gaps18121820 & PROVED & \checkmark & verified \\
\texttt{localFactor\_oneThousandEightHundredTwenty\_two} & Gaps18121820 & PROVED & \checkmark & verified \\
\texttt{nu\_p\_oneThousandEightHundredEighteen} & Gaps18121820 & PROVED & \checkmark & verified \\
\texttt{nu\_p\_oneThousandEightHundredFourteen} & Gaps18121820 & PROVED & \checkmark & verified \\
\texttt{nu\_p\_oneThousandEightHundredSixteen} & Gaps18121820 & PROVED & \checkmark & verified \\
\texttt{nu\_p\_oneThousandEightHundredTwelve} & Gaps18121820 & PROVED & \checkmark & verified \\
\texttt{nu\_p\_oneThousandEightHundredTwelve\_two} & Gaps18121820 & PROVED & \checkmark & verified \\
\texttt{nu\_p\_oneThousandEightHundredTwenty} & Gaps18121820 & PROVED & \checkmark & verified \\
\texttt{nu\_p\_oneThousandEightHundredTwenty\_two} & Gaps18121820 & PROVED & \checkmark & verified \\
\texttt{singular\_series\_finite\_pos\_evenPair\_oneThousandEightHundredEighteen} & Gaps18121820 & PROVED & \checkmark & verified \\
\texttt{singular\_series\_finite\_pos\_evenPair\_oneThousandEightHundredFourteen} & Gaps18121820 & PROVED & \checkmark & verified \\
\texttt{singular\_series\_finite\_pos\_evenPair\_oneThousandEightHundredSixteen} & Gaps18121820 & PROVED & \checkmark & verified \\
\texttt{singular\_series\_finite\_pos\_evenPair\_oneThousandEightHundredTwelve} & Gaps18121820 & PROVED & \checkmark & verified \\
\texttt{singular\_series\_finite\_pos\_evenPair\_oneThousandEightHundredTwenty} & Gaps18121820 & PROVED & \checkmark & verified \\
\texttt{singular\_series\_pos\_evenPair\_oneThousandEightHundredEighteen} & Gaps18121820 & PROVED & \checkmark & verified \\
\texttt{singular\_series\_pos\_evenPair\_oneThousandEightHundredFourteen} & Gaps18121820 & PROVED & \checkmark & verified \\
\texttt{singular\_series\_pos\_evenPair\_oneThousandEightHundredSixteen} & Gaps18121820 & PROVED & \checkmark & verified \\
\texttt{singular\_series\_pos\_evenPair\_oneThousandEightHundredTwelve} & Gaps18121820 & PROVED & \checkmark & verified \\
\texttt{singular\_series\_pos\_evenPair\_oneThousandEightHundredTwenty} & Gaps18121820 & PROVED & \checkmark & verified \\
\texttt{evenPair\_card\_oneHundredEightyEight} & Gaps182190 & PROVED & \checkmark & verified \\
\texttt{evenPair\_card\_oneHundredEightyFour} & Gaps182190 & PROVED & \checkmark & verified \\
\texttt{evenPair\_card\_oneHundredEightySix} & Gaps182190 & PROVED & \checkmark & verified \\
\texttt{evenPair\_card\_oneHundredEightyTwo} & Gaps182190 & PROVED & \checkmark & verified \\
\texttt{evenPair\_card\_oneHundredNinety} & Gaps182190 & PROVED & \checkmark & verified \\
\texttt{isAdmissible\_evenPair\_oneHundredEightyEight} & Gaps182190 & PROVED & \checkmark & verified \\
\texttt{isAdmissible\_evenPair\_oneHundredEightyFour} & Gaps182190 & PROVED & \checkmark & verified \\
\texttt{isAdmissible\_evenPair\_oneHundredEightySix} & Gaps182190 & PROVED & \checkmark & verified \\
\texttt{isAdmissible\_evenPair\_oneHundredEightyTwo} & Gaps182190 & PROVED & \checkmark & verified \\
\texttt{isAdmissible\_evenPair\_oneHundredNinety} & Gaps182190 & PROVED & \checkmark & verified \\
\texttt{localFactor\_oneHundredEightyTwo\_two} & Gaps182190 & PROVED & \checkmark & verified \\
\texttt{localFactor\_oneHundredNinety\_two} & Gaps182190 & PROVED & \checkmark & verified \\
\texttt{nu\_p\_oneHundredEightyEight} & Gaps182190 & PROVED & \checkmark & verified \\
\texttt{nu\_p\_oneHundredEightyFour} & Gaps182190 & PROVED & \checkmark & verified \\
\texttt{nu\_p\_oneHundredEightySix} & Gaps182190 & PROVED & \checkmark & verified \\
\texttt{nu\_p\_oneHundredEightyTwo} & Gaps182190 & PROVED & \checkmark & verified \\
\texttt{nu\_p\_oneHundredEightyTwo\_two} & Gaps182190 & PROVED & \checkmark & verified \\
\texttt{nu\_p\_oneHundredNinety} & Gaps182190 & PROVED & \checkmark & verified \\
\texttt{nu\_p\_oneHundredNinety\_two} & Gaps182190 & PROVED & \checkmark & verified \\
\texttt{singular\_series\_finite\_pos\_evenPair\_oneHundredEightyEight} & Gaps182190 & PROVED & \checkmark & verified \\
\texttt{singular\_series\_finite\_pos\_evenPair\_oneHundredEightyFour} & Gaps182190 & PROVED & \checkmark & verified \\
\texttt{singular\_series\_finite\_pos\_evenPair\_oneHundredEightySix} & Gaps182190 & PROVED & \checkmark & verified \\
\texttt{singular\_series\_finite\_pos\_evenPair\_oneHundredEightyTwo} & Gaps182190 & PROVED & \checkmark & verified \\
\texttt{singular\_series\_finite\_pos\_evenPair\_oneHundredNinety} & Gaps182190 & PROVED & \checkmark & verified \\
\texttt{singular\_series\_pos\_evenPair\_oneHundredEightyEight} & Gaps182190 & PROVED & \checkmark & verified \\
\texttt{singular\_series\_pos\_evenPair\_oneHundredEightyFour} & Gaps182190 & PROVED & \checkmark & verified \\
\texttt{singular\_series\_pos\_evenPair\_oneHundredEightySix} & Gaps182190 & PROVED & \checkmark & verified \\
\texttt{singular\_series\_pos\_evenPair\_oneHundredEightyTwo} & Gaps182190 & PROVED & \checkmark & verified \\
\texttt{singular\_series\_pos\_evenPair\_oneHundredNinety} & Gaps182190 & PROVED & \checkmark & verified \\
\texttt{evenPair\_card\_oneThousandEightHundredThirty} & Gaps18221830 & PROVED & \checkmark & verified \\
\texttt{evenPair\_card\_oneThousandEightHundredTwentyEight} & Gaps18221830 & PROVED & \checkmark & verified \\
\texttt{evenPair\_card\_oneThousandEightHundredTwentyFour} & Gaps18221830 & PROVED & \checkmark & verified \\
\texttt{evenPair\_card\_oneThousandEightHundredTwentySix} & Gaps18221830 & PROVED & \checkmark & verified \\
\texttt{evenPair\_card\_oneThousandEightHundredTwentyTwo} & Gaps18221830 & PROVED & \checkmark & verified \\
\texttt{isAdmissible\_evenPair\_oneThousandEightHundredThirty} & Gaps18221830 & PROVED & \checkmark & verified \\
\texttt{isAdmissible\_evenPair\_oneThousandEightHundredTwentyEight} & Gaps18221830 & PROVED & \checkmark & verified \\
\texttt{isAdmissible\_evenPair\_oneThousandEightHundredTwentyFour} & Gaps18221830 & PROVED & \checkmark & verified \\
\texttt{isAdmissible\_evenPair\_oneThousandEightHundredTwentySix} & Gaps18221830 & PROVED & \checkmark & verified \\
\texttt{isAdmissible\_evenPair\_oneThousandEightHundredTwentyTwo} & Gaps18221830 & PROVED & \checkmark & verified \\
\texttt{localFactor\_oneThousandEightHundredThirty\_two} & Gaps18221830 & PROVED & \checkmark & verified \\
\texttt{localFactor\_oneThousandEightHundredTwentyTwo\_two} & Gaps18221830 & PROVED & \checkmark & verified \\
\texttt{nu\_p\_oneThousandEightHundredThirty} & Gaps18221830 & PROVED & \checkmark & verified \\
\texttt{nu\_p\_oneThousandEightHundredThirty\_two} & Gaps18221830 & PROVED & \checkmark & verified \\
\texttt{nu\_p\_oneThousandEightHundredTwentyEight} & Gaps18221830 & PROVED & \checkmark & verified \\
\texttt{nu\_p\_oneThousandEightHundredTwentyFour} & Gaps18221830 & PROVED & \checkmark & verified \\
\texttt{nu\_p\_oneThousandEightHundredTwentySix} & Gaps18221830 & PROVED & \checkmark & verified \\
\texttt{nu\_p\_oneThousandEightHundredTwentyTwo} & Gaps18221830 & PROVED & \checkmark & verified \\
\texttt{nu\_p\_oneThousandEightHundredTwentyTwo\_two} & Gaps18221830 & PROVED & \checkmark & verified \\
\texttt{singular\_series\_finite\_pos\_evenPair\_oneThousandEightHundredThirty} & Gaps18221830 & PROVED & \checkmark & verified \\
\texttt{singular\_series\_finite\_pos\_evenPair\_oneThousandEightHundredTwentyEight} & Gaps18221830 & PROVED & \checkmark & verified \\
\texttt{singular\_series\_finite\_pos\_evenPair\_oneThousandEightHundredTwentyFour} & Gaps18221830 & PROVED & \checkmark & verified \\
\texttt{singular\_series\_finite\_pos\_evenPair\_oneThousandEightHundredTwentySix} & Gaps18221830 & PROVED & \checkmark & verified \\
\texttt{singular\_series\_finite\_pos\_evenPair\_oneThousandEightHundredTwentyTwo} & Gaps18221830 & PROVED & \checkmark & verified \\
\texttt{singular\_series\_pos\_evenPair\_oneThousandEightHundredThirty} & Gaps18221830 & PROVED & \checkmark & verified \\
\texttt{singular\_series\_pos\_evenPair\_oneThousandEightHundredTwentyEight} & Gaps18221830 & PROVED & \checkmark & verified \\
\texttt{singular\_series\_pos\_evenPair\_oneThousandEightHundredTwentyFour} & Gaps18221830 & PROVED & \checkmark & verified \\
\texttt{singular\_series\_pos\_evenPair\_oneThousandEightHundredTwentySix} & Gaps18221830 & PROVED & \checkmark & verified \\
\texttt{singular\_series\_pos\_evenPair\_oneThousandEightHundredTwentyTwo} & Gaps18221830 & PROVED & \checkmark & verified \\
\texttt{evenPair\_card\_oneThousandEightHundredForty} & Gaps18321840 & PROVED & \checkmark & verified \\
\texttt{evenPair\_card\_oneThousandEightHundredThirtyEight} & Gaps18321840 & PROVED & \checkmark & verified \\
\texttt{evenPair\_card\_oneThousandEightHundredThirtyFour} & Gaps18321840 & PROVED & \checkmark & verified \\
\texttt{evenPair\_card\_oneThousandEightHundredThirtySix} & Gaps18321840 & PROVED & \checkmark & verified \\
\texttt{evenPair\_card\_oneThousandEightHundredThirtyTwo} & Gaps18321840 & PROVED & \checkmark & verified \\
\texttt{isAdmissible\_evenPair\_oneThousandEightHundredForty} & Gaps18321840 & PROVED & \checkmark & verified \\
\texttt{isAdmissible\_evenPair\_oneThousandEightHundredThirtyEight} & Gaps18321840 & PROVED & \checkmark & verified \\
\texttt{isAdmissible\_evenPair\_oneThousandEightHundredThirtyFour} & Gaps18321840 & PROVED & \checkmark & verified \\
\texttt{isAdmissible\_evenPair\_oneThousandEightHundredThirtySix} & Gaps18321840 & PROVED & \checkmark & verified \\
\texttt{isAdmissible\_evenPair\_oneThousandEightHundredThirtyTwo} & Gaps18321840 & PROVED & \checkmark & verified \\
\texttt{localFactor\_oneThousandEightHundredForty\_two} & Gaps18321840 & PROVED & \checkmark & verified \\
\texttt{localFactor\_oneThousandEightHundredThirtyTwo\_two} & Gaps18321840 & PROVED & \checkmark & verified \\
\texttt{nu\_p\_oneThousandEightHundredForty} & Gaps18321840 & PROVED & \checkmark & verified \\
\texttt{nu\_p\_oneThousandEightHundredForty\_two} & Gaps18321840 & PROVED & \checkmark & verified \\
\texttt{nu\_p\_oneThousandEightHundredThirtyEight} & Gaps18321840 & PROVED & \checkmark & verified \\
\texttt{nu\_p\_oneThousandEightHundredThirtyFour} & Gaps18321840 & PROVED & \checkmark & verified \\
\texttt{nu\_p\_oneThousandEightHundredThirtySix} & Gaps18321840 & PROVED & \checkmark & verified \\
\texttt{nu\_p\_oneThousandEightHundredThirtyTwo} & Gaps18321840 & PROVED & \checkmark & verified \\
\texttt{nu\_p\_oneThousandEightHundredThirtyTwo\_two} & Gaps18321840 & PROVED & \checkmark & verified \\
\texttt{singular\_series\_finite\_pos\_evenPair\_oneThousandEightHundredForty} & Gaps18321840 & PROVED & \checkmark & verified \\
\texttt{singular\_series\_finite\_pos\_evenPair\_oneThousandEightHundredThirtyEight} & Gaps18321840 & PROVED & \checkmark & verified \\
\texttt{singular\_series\_finite\_pos\_evenPair\_oneThousandEightHundredThirtyFour} & Gaps18321840 & PROVED & \checkmark & verified \\
\texttt{singular\_series\_finite\_pos\_evenPair\_oneThousandEightHundredThirtySix} & Gaps18321840 & PROVED & \checkmark & verified \\
\texttt{singular\_series\_finite\_pos\_evenPair\_oneThousandEightHundredThirtyTwo} & Gaps18321840 & PROVED & \checkmark & verified \\
\texttt{singular\_series\_pos\_evenPair\_oneThousandEightHundredForty} & Gaps18321840 & PROVED & \checkmark & verified \\
\texttt{singular\_series\_pos\_evenPair\_oneThousandEightHundredThirtyEight} & Gaps18321840 & PROVED & \checkmark & verified \\
\texttt{singular\_series\_pos\_evenPair\_oneThousandEightHundredThirtyFour} & Gaps18321840 & PROVED & \checkmark & verified \\
\texttt{singular\_series\_pos\_evenPair\_oneThousandEightHundredThirtySix} & Gaps18321840 & PROVED & \checkmark & verified \\
\texttt{singular\_series\_pos\_evenPair\_oneThousandEightHundredThirtyTwo} & Gaps18321840 & PROVED & \checkmark & verified \\
\texttt{evenPair\_card\_oneThousandEightHundredFifty} & Gaps18421850 & PROVED & \checkmark & verified \\
\texttt{evenPair\_card\_oneThousandEightHundredFortyEight} & Gaps18421850 & PROVED & \checkmark & verified \\
\texttt{evenPair\_card\_oneThousandEightHundredFortyFour} & Gaps18421850 & PROVED & \checkmark & verified \\
\texttt{evenPair\_card\_oneThousandEightHundredFortySix} & Gaps18421850 & PROVED & \checkmark & verified \\
\texttt{evenPair\_card\_oneThousandEightHundredFortyTwo} & Gaps18421850 & PROVED & \checkmark & verified \\
\texttt{isAdmissible\_evenPair\_oneThousandEightHundredFifty} & Gaps18421850 & PROVED & \checkmark & verified \\
\texttt{isAdmissible\_evenPair\_oneThousandEightHundredFortyEight} & Gaps18421850 & PROVED & \checkmark & verified \\
\texttt{isAdmissible\_evenPair\_oneThousandEightHundredFortyFour} & Gaps18421850 & PROVED & \checkmark & verified \\
\texttt{isAdmissible\_evenPair\_oneThousandEightHundredFortySix} & Gaps18421850 & PROVED & \checkmark & verified \\
\texttt{isAdmissible\_evenPair\_oneThousandEightHundredFortyTwo} & Gaps18421850 & PROVED & \checkmark & verified \\
\texttt{localFactor\_oneThousandEightHundredFifty\_two} & Gaps18421850 & PROVED & \checkmark & verified \\
\texttt{localFactor\_oneThousandEightHundredFortyTwo\_two} & Gaps18421850 & PROVED & \checkmark & verified \\
\texttt{nu\_p\_oneThousandEightHundredFifty} & Gaps18421850 & PROVED & \checkmark & verified \\
\texttt{nu\_p\_oneThousandEightHundredFifty\_two} & Gaps18421850 & PROVED & \checkmark & verified \\
\texttt{nu\_p\_oneThousandEightHundredFortyEight} & Gaps18421850 & PROVED & \checkmark & verified \\
\texttt{nu\_p\_oneThousandEightHundredFortyFour} & Gaps18421850 & PROVED & \checkmark & verified \\
\texttt{nu\_p\_oneThousandEightHundredFortySix} & Gaps18421850 & PROVED & \checkmark & verified \\
\texttt{nu\_p\_oneThousandEightHundredFortyTwo} & Gaps18421850 & PROVED & \checkmark & verified \\
\texttt{nu\_p\_oneThousandEightHundredFortyTwo\_two} & Gaps18421850 & PROVED & \checkmark & verified \\
\texttt{singular\_series\_finite\_pos\_evenPair\_oneThousandEightHundredFifty} & Gaps18421850 & PROVED & \checkmark & verified \\
\texttt{singular\_series\_finite\_pos\_evenPair\_oneThousandEightHundredFortyEight} & Gaps18421850 & PROVED & \checkmark & verified \\
\texttt{singular\_series\_finite\_pos\_evenPair\_oneThousandEightHundredFortyFour} & Gaps18421850 & PROVED & \checkmark & verified \\
\texttt{singular\_series\_finite\_pos\_evenPair\_oneThousandEightHundredFortySix} & Gaps18421850 & PROVED & \checkmark & verified \\
\texttt{singular\_series\_finite\_pos\_evenPair\_oneThousandEightHundredFortyTwo} & Gaps18421850 & PROVED & \checkmark & verified \\
\texttt{singular\_series\_pos\_evenPair\_oneThousandEightHundredFifty} & Gaps18421850 & PROVED & \checkmark & verified \\
\texttt{singular\_series\_pos\_evenPair\_oneThousandEightHundredFortyEight} & Gaps18421850 & PROVED & \checkmark & verified \\
\texttt{singular\_series\_pos\_evenPair\_oneThousandEightHundredFortyFour} & Gaps18421850 & PROVED & \checkmark & verified \\
\texttt{singular\_series\_pos\_evenPair\_oneThousandEightHundredFortySix} & Gaps18421850 & PROVED & \checkmark & verified \\
\texttt{singular\_series\_pos\_evenPair\_oneThousandEightHundredFortyTwo} & Gaps18421850 & PROVED & \checkmark & verified \\
\texttt{evenPair\_card\_oneThousandEightHundredFiftyEight} & Gaps18521860 & PROVED & \checkmark & verified \\
\texttt{evenPair\_card\_oneThousandEightHundredFiftyFour} & Gaps18521860 & PROVED & \checkmark & verified \\
\texttt{evenPair\_card\_oneThousandEightHundredFiftySix} & Gaps18521860 & PROVED & \checkmark & verified \\
\texttt{evenPair\_card\_oneThousandEightHundredFiftyTwo} & Gaps18521860 & PROVED & \checkmark & verified \\
\texttt{evenPair\_card\_oneThousandEightHundredSixty} & Gaps18521860 & PROVED & \checkmark & verified \\
\texttt{isAdmissible\_evenPair\_oneThousandEightHundredFiftyEight} & Gaps18521860 & PROVED & \checkmark & verified \\
\texttt{isAdmissible\_evenPair\_oneThousandEightHundredFiftyFour} & Gaps18521860 & PROVED & \checkmark & verified \\
\texttt{isAdmissible\_evenPair\_oneThousandEightHundredFiftySix} & Gaps18521860 & PROVED & \checkmark & verified \\
\texttt{isAdmissible\_evenPair\_oneThousandEightHundredFiftyTwo} & Gaps18521860 & PROVED & \checkmark & verified \\
\texttt{isAdmissible\_evenPair\_oneThousandEightHundredSixty} & Gaps18521860 & PROVED & \checkmark & verified \\
\texttt{localFactor\_oneThousandEightHundredFiftyTwo\_two} & Gaps18521860 & PROVED & \checkmark & verified \\
\texttt{localFactor\_oneThousandEightHundredSixty\_two} & Gaps18521860 & PROVED & \checkmark & verified \\
\texttt{nu\_p\_oneThousandEightHundredFiftyEight} & Gaps18521860 & PROVED & \checkmark & verified \\
\texttt{nu\_p\_oneThousandEightHundredFiftyFour} & Gaps18521860 & PROVED & \checkmark & verified \\
\texttt{nu\_p\_oneThousandEightHundredFiftySix} & Gaps18521860 & PROVED & \checkmark & verified \\
\texttt{nu\_p\_oneThousandEightHundredFiftyTwo} & Gaps18521860 & PROVED & \checkmark & verified \\
\texttt{nu\_p\_oneThousandEightHundredFiftyTwo\_two} & Gaps18521860 & PROVED & \checkmark & verified \\
\texttt{nu\_p\_oneThousandEightHundredSixty} & Gaps18521860 & PROVED & \checkmark & verified \\
\texttt{nu\_p\_oneThousandEightHundredSixty\_two} & Gaps18521860 & PROVED & \checkmark & verified \\
\texttt{singular\_series\_finite\_pos\_evenPair\_oneThousandEightHundredFiftyEight} & Gaps18521860 & PROVED & \checkmark & verified \\
\texttt{singular\_series\_finite\_pos\_evenPair\_oneThousandEightHundredFiftyFour} & Gaps18521860 & PROVED & \checkmark & verified \\
\texttt{singular\_series\_finite\_pos\_evenPair\_oneThousandEightHundredFiftySix} & Gaps18521860 & PROVED & \checkmark & verified \\
\texttt{singular\_series\_finite\_pos\_evenPair\_oneThousandEightHundredFiftyTwo} & Gaps18521860 & PROVED & \checkmark & verified \\
\texttt{singular\_series\_finite\_pos\_evenPair\_oneThousandEightHundredSixty} & Gaps18521860 & PROVED & \checkmark & verified \\
\texttt{singular\_series\_pos\_evenPair\_oneThousandEightHundredFiftyEight} & Gaps18521860 & PROVED & \checkmark & verified \\
\texttt{singular\_series\_pos\_evenPair\_oneThousandEightHundredFiftyFour} & Gaps18521860 & PROVED & \checkmark & verified \\
\texttt{singular\_series\_pos\_evenPair\_oneThousandEightHundredFiftySix} & Gaps18521860 & PROVED & \checkmark & verified \\
\texttt{singular\_series\_pos\_evenPair\_oneThousandEightHundredFiftyTwo} & Gaps18521860 & PROVED & \checkmark & verified \\
\texttt{singular\_series\_pos\_evenPair\_oneThousandEightHundredSixty} & Gaps18521860 & PROVED & \checkmark & verified \\
\texttt{evenPair\_card\_oneThousandEightHundredSeventy} & Gaps18621870 & PROVED & \checkmark & verified \\
\texttt{evenPair\_card\_oneThousandEightHundredSixtyEight} & Gaps18621870 & PROVED & \checkmark & verified \\
\texttt{evenPair\_card\_oneThousandEightHundredSixtyFour} & Gaps18621870 & PROVED & \checkmark & verified \\
\texttt{evenPair\_card\_oneThousandEightHundredSixtySix} & Gaps18621870 & PROVED & \checkmark & verified \\
\texttt{evenPair\_card\_oneThousandEightHundredSixtyTwo} & Gaps18621870 & PROVED & \checkmark & verified \\
\texttt{isAdmissible\_evenPair\_oneThousandEightHundredSeventy} & Gaps18621870 & PROVED & \checkmark & verified \\
\texttt{isAdmissible\_evenPair\_oneThousandEightHundredSixtyEight} & Gaps18621870 & PROVED & \checkmark & verified \\
\texttt{isAdmissible\_evenPair\_oneThousandEightHundredSixtyFour} & Gaps18621870 & PROVED & \checkmark & verified \\
\texttt{isAdmissible\_evenPair\_oneThousandEightHundredSixtySix} & Gaps18621870 & PROVED & \checkmark & verified \\
\texttt{isAdmissible\_evenPair\_oneThousandEightHundredSixtyTwo} & Gaps18621870 & PROVED & \checkmark & verified \\
\texttt{localFactor\_oneThousandEightHundredSeventy\_two} & Gaps18621870 & PROVED & \checkmark & verified \\
\texttt{localFactor\_oneThousandEightHundredSixtyTwo\_two} & Gaps18621870 & PROVED & \checkmark & verified \\
\texttt{nu\_p\_oneThousandEightHundredSeventy} & Gaps18621870 & PROVED & \checkmark & verified \\
\texttt{nu\_p\_oneThousandEightHundredSeventy\_two} & Gaps18621870 & PROVED & \checkmark & verified \\
\texttt{nu\_p\_oneThousandEightHundredSixtyEight} & Gaps18621870 & PROVED & \checkmark & verified \\
\texttt{nu\_p\_oneThousandEightHundredSixtyFour} & Gaps18621870 & PROVED & \checkmark & verified \\
\texttt{nu\_p\_oneThousandEightHundredSixtySix} & Gaps18621870 & PROVED & \checkmark & verified \\
\texttt{nu\_p\_oneThousandEightHundredSixtyTwo} & Gaps18621870 & PROVED & \checkmark & verified \\
\texttt{nu\_p\_oneThousandEightHundredSixtyTwo\_two} & Gaps18621870 & PROVED & \checkmark & verified \\
\texttt{singular\_series\_finite\_pos\_evenPair\_oneThousandEightHundredSeventy} & Gaps18621870 & PROVED & \checkmark & verified \\
\texttt{singular\_series\_finite\_pos\_evenPair\_oneThousandEightHundredSixtyEight} & Gaps18621870 & PROVED & \checkmark & verified \\
\texttt{singular\_series\_finite\_pos\_evenPair\_oneThousandEightHundredSixtyFour} & Gaps18621870 & PROVED & \checkmark & verified \\
\texttt{singular\_series\_finite\_pos\_evenPair\_oneThousandEightHundredSixtySix} & Gaps18621870 & PROVED & \checkmark & verified \\
\texttt{singular\_series\_finite\_pos\_evenPair\_oneThousandEightHundredSixtyTwo} & Gaps18621870 & PROVED & \checkmark & verified \\
\texttt{singular\_series\_pos\_evenPair\_oneThousandEightHundredSeventy} & Gaps18621870 & PROVED & \checkmark & verified \\
\texttt{singular\_series\_pos\_evenPair\_oneThousandEightHundredSixtyEight} & Gaps18621870 & PROVED & \checkmark & verified \\
\texttt{singular\_series\_pos\_evenPair\_oneThousandEightHundredSixtyFour} & Gaps18621870 & PROVED & \checkmark & verified \\
\texttt{singular\_series\_pos\_evenPair\_oneThousandEightHundredSixtySix} & Gaps18621870 & PROVED & \checkmark & verified \\
\texttt{singular\_series\_pos\_evenPair\_oneThousandEightHundredSixtyTwo} & Gaps18621870 & PROVED & \checkmark & verified \\
\texttt{evenPair\_card\_oneThousandEightHundredEighty} & Gaps18721880 & PROVED & \checkmark & verified \\
\texttt{evenPair\_card\_oneThousandEightHundredSeventyEight} & Gaps18721880 & PROVED & \checkmark & verified \\
\texttt{evenPair\_card\_oneThousandEightHundredSeventyFour} & Gaps18721880 & PROVED & \checkmark & verified \\
\texttt{evenPair\_card\_oneThousandEightHundredSeventySix} & Gaps18721880 & PROVED & \checkmark & verified \\
\texttt{evenPair\_card\_oneThousandEightHundredSeventyTwo} & Gaps18721880 & PROVED & \checkmark & verified \\
\texttt{isAdmissible\_evenPair\_oneThousandEightHundredEighty} & Gaps18721880 & PROVED & \checkmark & verified \\
\texttt{isAdmissible\_evenPair\_oneThousandEightHundredSeventyEight} & Gaps18721880 & PROVED & \checkmark & verified \\
\texttt{isAdmissible\_evenPair\_oneThousandEightHundredSeventyFour} & Gaps18721880 & PROVED & \checkmark & verified \\
\texttt{isAdmissible\_evenPair\_oneThousandEightHundredSeventySix} & Gaps18721880 & PROVED & \checkmark & verified \\
\texttt{isAdmissible\_evenPair\_oneThousandEightHundredSeventyTwo} & Gaps18721880 & PROVED & \checkmark & verified \\
\texttt{localFactor\_oneThousandEightHundredEighty\_two} & Gaps18721880 & PROVED & \checkmark & verified \\
\texttt{localFactor\_oneThousandEightHundredSeventyTwo\_two} & Gaps18721880 & PROVED & \checkmark & verified \\
\texttt{nu\_p\_oneThousandEightHundredEighty} & Gaps18721880 & PROVED & \checkmark & verified \\
\texttt{nu\_p\_oneThousandEightHundredEighty\_two} & Gaps18721880 & PROVED & \checkmark & verified \\
\texttt{nu\_p\_oneThousandEightHundredSeventyEight} & Gaps18721880 & PROVED & \checkmark & verified \\
\texttt{nu\_p\_oneThousandEightHundredSeventyFour} & Gaps18721880 & PROVED & \checkmark & verified \\
\texttt{nu\_p\_oneThousandEightHundredSeventySix} & Gaps18721880 & PROVED & \checkmark & verified \\
\texttt{nu\_p\_oneThousandEightHundredSeventyTwo} & Gaps18721880 & PROVED & \checkmark & verified \\
\texttt{nu\_p\_oneThousandEightHundredSeventyTwo\_two} & Gaps18721880 & PROVED & \checkmark & verified \\
\texttt{singular\_series\_finite\_pos\_evenPair\_oneThousandEightHundredEighty} & Gaps18721880 & PROVED & \checkmark & verified \\
\texttt{singular\_series\_finite\_pos\_evenPair\_oneThousandEightHundredSeventyEight} & Gaps18721880 & PROVED & \checkmark & verified \\
\texttt{singular\_series\_finite\_pos\_evenPair\_oneThousandEightHundredSeventyFour} & Gaps18721880 & PROVED & \checkmark & verified \\
\texttt{singular\_series\_finite\_pos\_evenPair\_oneThousandEightHundredSeventySix} & Gaps18721880 & PROVED & \checkmark & verified \\
\texttt{singular\_series\_finite\_pos\_evenPair\_oneThousandEightHundredSeventyTwo} & Gaps18721880 & PROVED & \checkmark & verified \\
\texttt{singular\_series\_pos\_evenPair\_oneThousandEightHundredEighty} & Gaps18721880 & PROVED & \checkmark & verified \\
\texttt{singular\_series\_pos\_evenPair\_oneThousandEightHundredSeventyEight} & Gaps18721880 & PROVED & \checkmark & verified \\
\texttt{singular\_series\_pos\_evenPair\_oneThousandEightHundredSeventyFour} & Gaps18721880 & PROVED & \checkmark & verified \\
\texttt{singular\_series\_pos\_evenPair\_oneThousandEightHundredSeventySix} & Gaps18721880 & PROVED & \checkmark & verified \\
\texttt{singular\_series\_pos\_evenPair\_oneThousandEightHundredSeventyTwo} & Gaps18721880 & PROVED & \checkmark & verified \\
\texttt{evenPair\_card\_oneThousandEightHundredEightyEight} & Gaps18821890 & PROVED & \checkmark & verified \\
\texttt{evenPair\_card\_oneThousandEightHundredEightyFour} & Gaps18821890 & PROVED & \checkmark & verified \\
\texttt{evenPair\_card\_oneThousandEightHundredEightySix} & Gaps18821890 & PROVED & \checkmark & verified \\
\texttt{evenPair\_card\_oneThousandEightHundredEightyTwo} & Gaps18821890 & PROVED & \checkmark & verified \\
\texttt{evenPair\_card\_oneThousandEightHundredNinety} & Gaps18821890 & PROVED & \checkmark & verified \\
\texttt{isAdmissible\_evenPair\_oneThousandEightHundredEightyEight} & Gaps18821890 & PROVED & \checkmark & verified \\
\texttt{isAdmissible\_evenPair\_oneThousandEightHundredEightyFour} & Gaps18821890 & PROVED & \checkmark & verified \\
\texttt{isAdmissible\_evenPair\_oneThousandEightHundredEightySix} & Gaps18821890 & PROVED & \checkmark & verified \\
\texttt{isAdmissible\_evenPair\_oneThousandEightHundredEightyTwo} & Gaps18821890 & PROVED & \checkmark & verified \\
\texttt{isAdmissible\_evenPair\_oneThousandEightHundredNinety} & Gaps18821890 & PROVED & \checkmark & verified \\
\texttt{localFactor\_oneThousandEightHundredEightyTwo\_two} & Gaps18821890 & PROVED & \checkmark & verified \\
\texttt{localFactor\_oneThousandEightHundredNinety\_two} & Gaps18821890 & PROVED & \checkmark & verified \\
\texttt{nu\_p\_oneThousandEightHundredEightyEight} & Gaps18821890 & PROVED & \checkmark & verified \\
\texttt{nu\_p\_oneThousandEightHundredEightyFour} & Gaps18821890 & PROVED & \checkmark & verified \\
\texttt{nu\_p\_oneThousandEightHundredEightySix} & Gaps18821890 & PROVED & \checkmark & verified \\
\texttt{nu\_p\_oneThousandEightHundredEightyTwo} & Gaps18821890 & PROVED & \checkmark & verified \\
\texttt{nu\_p\_oneThousandEightHundredEightyTwo\_two} & Gaps18821890 & PROVED & \checkmark & verified \\
\texttt{nu\_p\_oneThousandEightHundredNinety} & Gaps18821890 & PROVED & \checkmark & verified \\
\texttt{nu\_p\_oneThousandEightHundredNinety\_two} & Gaps18821890 & PROVED & \checkmark & verified \\
\texttt{singular\_series\_finite\_pos\_evenPair\_oneThousandEightHundredEightyEight} & Gaps18821890 & PROVED & \checkmark & verified \\
\texttt{singular\_series\_finite\_pos\_evenPair\_oneThousandEightHundredEightyFour} & Gaps18821890 & PROVED & \checkmark & verified \\
\texttt{singular\_series\_finite\_pos\_evenPair\_oneThousandEightHundredEightySix} & Gaps18821890 & PROVED & \checkmark & verified \\
\texttt{singular\_series\_finite\_pos\_evenPair\_oneThousandEightHundredEightyTwo} & Gaps18821890 & PROVED & \checkmark & verified \\
\texttt{singular\_series\_finite\_pos\_evenPair\_oneThousandEightHundredNinety} & Gaps18821890 & PROVED & \checkmark & verified \\
\texttt{singular\_series\_pos\_evenPair\_oneThousandEightHundredEightyEight} & Gaps18821890 & PROVED & \checkmark & verified \\
\texttt{singular\_series\_pos\_evenPair\_oneThousandEightHundredEightyFour} & Gaps18821890 & PROVED & \checkmark & verified \\
\texttt{singular\_series\_pos\_evenPair\_oneThousandEightHundredEightySix} & Gaps18821890 & PROVED & \checkmark & verified \\
\texttt{singular\_series\_pos\_evenPair\_oneThousandEightHundredEightyTwo} & Gaps18821890 & PROVED & \checkmark & verified \\
\texttt{singular\_series\_pos\_evenPair\_oneThousandEightHundredNinety} & Gaps18821890 & PROVED & \checkmark & verified \\
\texttt{evenPair\_card\_oneThousandEightHundredNinetyEight} & Gaps18921900 & PROVED & \checkmark & verified \\
\texttt{evenPair\_card\_oneThousandEightHundredNinetyFour} & Gaps18921900 & PROVED & \checkmark & verified \\
\texttt{evenPair\_card\_oneThousandEightHundredNinetySix} & Gaps18921900 & PROVED & \checkmark & verified \\
\texttt{evenPair\_card\_oneThousandEightHundredNinetyTwo} & Gaps18921900 & PROVED & \checkmark & verified \\
\texttt{evenPair\_card\_oneThousandNineHundred} & Gaps18921900 & PROVED & \checkmark & verified \\
\texttt{isAdmissible\_evenPair\_oneThousandEightHundredNinetyEight} & Gaps18921900 & PROVED & \checkmark & verified \\
\texttt{isAdmissible\_evenPair\_oneThousandEightHundredNinetyFour} & Gaps18921900 & PROVED & \checkmark & verified \\
\texttt{isAdmissible\_evenPair\_oneThousandEightHundredNinetySix} & Gaps18921900 & PROVED & \checkmark & verified \\
\texttt{isAdmissible\_evenPair\_oneThousandEightHundredNinetyTwo} & Gaps18921900 & PROVED & \checkmark & verified \\
\texttt{isAdmissible\_evenPair\_oneThousandNineHundred} & Gaps18921900 & PROVED & \checkmark & verified \\
\texttt{localFactor\_oneThousandEightHundredNinetyTwo\_two} & Gaps18921900 & PROVED & \checkmark & verified \\
\texttt{localFactor\_oneThousandNineHundred\_two} & Gaps18921900 & PROVED & \checkmark & verified \\
\texttt{nu\_p\_oneThousandEightHundredNinetyEight} & Gaps18921900 & PROVED & \checkmark & verified \\
\texttt{nu\_p\_oneThousandEightHundredNinetyFour} & Gaps18921900 & PROVED & \checkmark & verified \\
\texttt{nu\_p\_oneThousandEightHundredNinetySix} & Gaps18921900 & PROVED & \checkmark & verified \\
\texttt{nu\_p\_oneThousandEightHundredNinetyTwo} & Gaps18921900 & PROVED & \checkmark & verified \\
\texttt{nu\_p\_oneThousandEightHundredNinetyTwo\_two} & Gaps18921900 & PROVED & \checkmark & verified \\
\texttt{nu\_p\_oneThousandNineHundred} & Gaps18921900 & PROVED & \checkmark & verified \\
\texttt{nu\_p\_oneThousandNineHundred\_two} & Gaps18921900 & PROVED & \checkmark & verified \\
\texttt{singular\_series\_finite\_pos\_evenPair\_oneThousandEightHundredNinetyEight} & Gaps18921900 & PROVED & \checkmark & verified \\
\texttt{singular\_series\_finite\_pos\_evenPair\_oneThousandEightHundredNinetyFour} & Gaps18921900 & PROVED & \checkmark & verified \\
\texttt{singular\_series\_finite\_pos\_evenPair\_oneThousandEightHundredNinetySix} & Gaps18921900 & PROVED & \checkmark & verified \\
\texttt{singular\_series\_finite\_pos\_evenPair\_oneThousandEightHundredNinetyTwo} & Gaps18921900 & PROVED & \checkmark & verified \\
\texttt{singular\_series\_finite\_pos\_evenPair\_oneThousandNineHundred} & Gaps18921900 & PROVED & \checkmark & verified \\
\texttt{singular\_series\_pos\_evenPair\_oneThousandEightHundredNinetyEight} & Gaps18921900 & PROVED & \checkmark & verified \\
\texttt{singular\_series\_pos\_evenPair\_oneThousandEightHundredNinetyFour} & Gaps18921900 & PROVED & \checkmark & verified \\
\texttt{singular\_series\_pos\_evenPair\_oneThousandEightHundredNinetySix} & Gaps18921900 & PROVED & \checkmark & verified \\
\texttt{singular\_series\_pos\_evenPair\_oneThousandEightHundredNinetyTwo} & Gaps18921900 & PROVED & \checkmark & verified \\
\texttt{singular\_series\_pos\_evenPair\_oneThousandNineHundred} & Gaps18921900 & PROVED & \checkmark & verified \\
\texttt{evenPair\_card\_oneThousandNineHundredEight} & Gaps19021910 & PROVED & \checkmark & verified \\
\texttt{evenPair\_card\_oneThousandNineHundredFour} & Gaps19021910 & PROVED & \checkmark & verified \\
\texttt{evenPair\_card\_oneThousandNineHundredSix} & Gaps19021910 & PROVED & \checkmark & verified \\
\texttt{evenPair\_card\_oneThousandNineHundredTen} & Gaps19021910 & PROVED & \checkmark & verified \\
\texttt{evenPair\_card\_oneThousandNineHundredTwo} & Gaps19021910 & PROVED & \checkmark & verified \\
\texttt{isAdmissible\_evenPair\_oneThousandNineHundredEight} & Gaps19021910 & PROVED & \checkmark & verified \\
\texttt{isAdmissible\_evenPair\_oneThousandNineHundredFour} & Gaps19021910 & PROVED & \checkmark & verified \\
\texttt{isAdmissible\_evenPair\_oneThousandNineHundredSix} & Gaps19021910 & PROVED & \checkmark & verified \\
\texttt{isAdmissible\_evenPair\_oneThousandNineHundredTen} & Gaps19021910 & PROVED & \checkmark & verified \\
\texttt{isAdmissible\_evenPair\_oneThousandNineHundredTwo} & Gaps19021910 & PROVED & \checkmark & verified \\
\texttt{localFactor\_oneThousandNineHundredTen\_two} & Gaps19021910 & PROVED & \checkmark & verified \\
\texttt{localFactor\_oneThousandNineHundredTwo\_two} & Gaps19021910 & PROVED & \checkmark & verified \\
\texttt{nu\_p\_oneThousandNineHundredEight} & Gaps19021910 & PROVED & \checkmark & verified \\
\texttt{nu\_p\_oneThousandNineHundredFour} & Gaps19021910 & PROVED & \checkmark & verified \\
\texttt{nu\_p\_oneThousandNineHundredSix} & Gaps19021910 & PROVED & \checkmark & verified \\
\texttt{nu\_p\_oneThousandNineHundredTen} & Gaps19021910 & PROVED & \checkmark & verified \\
\texttt{nu\_p\_oneThousandNineHundredTen\_two} & Gaps19021910 & PROVED & \checkmark & verified \\
\texttt{nu\_p\_oneThousandNineHundredTwo} & Gaps19021910 & PROVED & \checkmark & verified \\
\texttt{nu\_p\_oneThousandNineHundredTwo\_two} & Gaps19021910 & PROVED & \checkmark & verified \\
\texttt{singular\_series\_finite\_pos\_evenPair\_oneThousandNineHundredEight} & Gaps19021910 & PROVED & \checkmark & verified \\
\texttt{singular\_series\_finite\_pos\_evenPair\_oneThousandNineHundredFour} & Gaps19021910 & PROVED & \checkmark & verified \\
\texttt{singular\_series\_finite\_pos\_evenPair\_oneThousandNineHundredSix} & Gaps19021910 & PROVED & \checkmark & verified \\
\texttt{singular\_series\_finite\_pos\_evenPair\_oneThousandNineHundredTen} & Gaps19021910 & PROVED & \checkmark & verified \\
\texttt{singular\_series\_finite\_pos\_evenPair\_oneThousandNineHundredTwo} & Gaps19021910 & PROVED & \checkmark & verified \\
\texttt{singular\_series\_pos\_evenPair\_oneThousandNineHundredEight} & Gaps19021910 & PROVED & \checkmark & verified \\
\texttt{singular\_series\_pos\_evenPair\_oneThousandNineHundredFour} & Gaps19021910 & PROVED & \checkmark & verified \\
\texttt{singular\_series\_pos\_evenPair\_oneThousandNineHundredSix} & Gaps19021910 & PROVED & \checkmark & verified \\
\texttt{singular\_series\_pos\_evenPair\_oneThousandNineHundredTen} & Gaps19021910 & PROVED & \checkmark & verified \\
\texttt{singular\_series\_pos\_evenPair\_oneThousandNineHundredTwo} & Gaps19021910 & PROVED & \checkmark & verified \\
\texttt{evenPair\_card\_oneThousandNineHundredEighteen} & Gaps19121920 & PROVED & \checkmark & verified \\
\texttt{evenPair\_card\_oneThousandNineHundredFourteen} & Gaps19121920 & PROVED & \checkmark & verified \\
\texttt{evenPair\_card\_oneThousandNineHundredSixteen} & Gaps19121920 & PROVED & \checkmark & verified \\
\texttt{evenPair\_card\_oneThousandNineHundredTwelve} & Gaps19121920 & PROVED & \checkmark & verified \\
\texttt{evenPair\_card\_oneThousandNineHundredTwenty} & Gaps19121920 & PROVED & \checkmark & verified \\
\texttt{isAdmissible\_evenPair\_oneThousandNineHundredEighteen} & Gaps19121920 & PROVED & \checkmark & verified \\
\texttt{isAdmissible\_evenPair\_oneThousandNineHundredFourteen} & Gaps19121920 & PROVED & \checkmark & verified \\
\texttt{isAdmissible\_evenPair\_oneThousandNineHundredSixteen} & Gaps19121920 & PROVED & \checkmark & verified \\
\texttt{isAdmissible\_evenPair\_oneThousandNineHundredTwelve} & Gaps19121920 & PROVED & \checkmark & verified \\
\texttt{isAdmissible\_evenPair\_oneThousandNineHundredTwenty} & Gaps19121920 & PROVED & \checkmark & verified \\
\texttt{localFactor\_oneThousandNineHundredTwelve\_two} & Gaps19121920 & PROVED & \checkmark & verified \\
\texttt{localFactor\_oneThousandNineHundredTwenty\_two} & Gaps19121920 & PROVED & \checkmark & verified \\
\texttt{nu\_p\_oneThousandNineHundredEighteen} & Gaps19121920 & PROVED & \checkmark & verified \\
\texttt{nu\_p\_oneThousandNineHundredFourteen} & Gaps19121920 & PROVED & \checkmark & verified \\
\texttt{nu\_p\_oneThousandNineHundredSixteen} & Gaps19121920 & PROVED & \checkmark & verified \\
\texttt{nu\_p\_oneThousandNineHundredTwelve} & Gaps19121920 & PROVED & \checkmark & verified \\
\texttt{nu\_p\_oneThousandNineHundredTwelve\_two} & Gaps19121920 & PROVED & \checkmark & verified \\
\texttt{nu\_p\_oneThousandNineHundredTwenty} & Gaps19121920 & PROVED & \checkmark & verified \\
\texttt{nu\_p\_oneThousandNineHundredTwenty\_two} & Gaps19121920 & PROVED & \checkmark & verified \\
\texttt{singular\_series\_finite\_pos\_evenPair\_oneThousandNineHundredEighteen} & Gaps19121920 & PROVED & \checkmark & verified \\
\texttt{singular\_series\_finite\_pos\_evenPair\_oneThousandNineHundredFourteen} & Gaps19121920 & PROVED & \checkmark & verified \\
\texttt{singular\_series\_finite\_pos\_evenPair\_oneThousandNineHundredSixteen} & Gaps19121920 & PROVED & \checkmark & verified \\
\texttt{singular\_series\_finite\_pos\_evenPair\_oneThousandNineHundredTwelve} & Gaps19121920 & PROVED & \checkmark & verified \\
\texttt{singular\_series\_finite\_pos\_evenPair\_oneThousandNineHundredTwenty} & Gaps19121920 & PROVED & \checkmark & verified \\
\texttt{singular\_series\_pos\_evenPair\_oneThousandNineHundredEighteen} & Gaps19121920 & PROVED & \checkmark & verified \\
\texttt{singular\_series\_pos\_evenPair\_oneThousandNineHundredFourteen} & Gaps19121920 & PROVED & \checkmark & verified \\
\texttt{singular\_series\_pos\_evenPair\_oneThousandNineHundredSixteen} & Gaps19121920 & PROVED & \checkmark & verified \\
\texttt{singular\_series\_pos\_evenPair\_oneThousandNineHundredTwelve} & Gaps19121920 & PROVED & \checkmark & verified \\
\texttt{singular\_series\_pos\_evenPair\_oneThousandNineHundredTwenty} & Gaps19121920 & PROVED & \checkmark & verified \\
\texttt{evenPair\_card\_oneHundredNinetyEight} & Gaps192200 & PROVED & \checkmark & verified \\
\texttt{evenPair\_card\_oneHundredNinetyFour} & Gaps192200 & PROVED & \checkmark & verified \\
\texttt{evenPair\_card\_oneHundredNinetySix} & Gaps192200 & PROVED & \checkmark & verified \\
\texttt{evenPair\_card\_oneHundredNinetyTwo} & Gaps192200 & PROVED & \checkmark & verified \\
\texttt{evenPair\_card\_twoHundred} & Gaps192200 & PROVED & \checkmark & verified \\
\texttt{isAdmissible\_evenPair\_oneHundredNinetyEight} & Gaps192200 & PROVED & \checkmark & verified \\
\texttt{isAdmissible\_evenPair\_oneHundredNinetyFour} & Gaps192200 & PROVED & \checkmark & verified \\
\texttt{isAdmissible\_evenPair\_oneHundredNinetySix} & Gaps192200 & PROVED & \checkmark & verified \\
\texttt{isAdmissible\_evenPair\_oneHundredNinetyTwo} & Gaps192200 & PROVED & \checkmark & verified \\
\texttt{isAdmissible\_evenPair\_twoHundred} & Gaps192200 & PROVED & \checkmark & verified \\
\texttt{localFactor\_oneHundredNinetyTwo\_two} & Gaps192200 & PROVED & \checkmark & verified \\
\texttt{localFactor\_twoHundred\_two} & Gaps192200 & PROVED & \checkmark & verified \\
\texttt{nu\_p\_oneHundredNinetyEight} & Gaps192200 & PROVED & \checkmark & verified \\
\texttt{nu\_p\_oneHundredNinetyFour} & Gaps192200 & PROVED & \checkmark & verified \\
\texttt{nu\_p\_oneHundredNinetySix} & Gaps192200 & PROVED & \checkmark & verified \\
\texttt{nu\_p\_oneHundredNinetyTwo} & Gaps192200 & PROVED & \checkmark & verified \\
\texttt{nu\_p\_oneHundredNinetyTwo\_two} & Gaps192200 & PROVED & \checkmark & verified \\
\texttt{nu\_p\_twoHundred} & Gaps192200 & PROVED & \checkmark & verified \\
\texttt{nu\_p\_twoHundred\_two} & Gaps192200 & PROVED & \checkmark & verified \\
\texttt{singular\_series\_finite\_pos\_evenPair\_oneHundredNinetyEight} & Gaps192200 & PROVED & \checkmark & verified \\
\texttt{singular\_series\_finite\_pos\_evenPair\_oneHundredNinetyFour} & Gaps192200 & PROVED & \checkmark & verified \\
\texttt{singular\_series\_finite\_pos\_evenPair\_oneHundredNinetySix} & Gaps192200 & PROVED & \checkmark & verified \\
\texttt{singular\_series\_finite\_pos\_evenPair\_oneHundredNinetyTwo} & Gaps192200 & PROVED & \checkmark & verified \\
\texttt{singular\_series\_finite\_pos\_evenPair\_twoHundred} & Gaps192200 & PROVED & \checkmark & verified \\
\texttt{singular\_series\_pos\_evenPair\_oneHundredNinetyEight} & Gaps192200 & PROVED & \checkmark & verified \\
\texttt{singular\_series\_pos\_evenPair\_oneHundredNinetyFour} & Gaps192200 & PROVED & \checkmark & verified \\
\texttt{singular\_series\_pos\_evenPair\_oneHundredNinetySix} & Gaps192200 & PROVED & \checkmark & verified \\
\texttt{singular\_series\_pos\_evenPair\_oneHundredNinetyTwo} & Gaps192200 & PROVED & \checkmark & verified \\
\texttt{singular\_series\_pos\_evenPair\_twoHundred} & Gaps192200 & PROVED & \checkmark & verified \\
\texttt{evenPair\_card\_oneThousandNineHundredThirty} & Gaps19221930 & PROVED & \checkmark & verified \\
\texttt{evenPair\_card\_oneThousandNineHundredTwentyEight} & Gaps19221930 & PROVED & \checkmark & verified \\
\texttt{evenPair\_card\_oneThousandNineHundredTwentyFour} & Gaps19221930 & PROVED & \checkmark & verified \\
\texttt{evenPair\_card\_oneThousandNineHundredTwentySix} & Gaps19221930 & PROVED & \checkmark & verified \\
\texttt{evenPair\_card\_oneThousandNineHundredTwentyTwo} & Gaps19221930 & PROVED & \checkmark & verified \\
\texttt{isAdmissible\_evenPair\_oneThousandNineHundredThirty} & Gaps19221930 & PROVED & \checkmark & verified \\
\texttt{isAdmissible\_evenPair\_oneThousandNineHundredTwentyEight} & Gaps19221930 & PROVED & \checkmark & verified \\
\texttt{isAdmissible\_evenPair\_oneThousandNineHundredTwentyFour} & Gaps19221930 & PROVED & \checkmark & verified \\
\texttt{isAdmissible\_evenPair\_oneThousandNineHundredTwentySix} & Gaps19221930 & PROVED & \checkmark & verified \\
\texttt{isAdmissible\_evenPair\_oneThousandNineHundredTwentyTwo} & Gaps19221930 & PROVED & \checkmark & verified \\
\texttt{localFactor\_oneThousandNineHundredThirty\_two} & Gaps19221930 & PROVED & \checkmark & verified \\
\texttt{localFactor\_oneThousandNineHundredTwentyTwo\_two} & Gaps19221930 & PROVED & \checkmark & verified \\
\texttt{nu\_p\_oneThousandNineHundredThirty} & Gaps19221930 & PROVED & \checkmark & verified \\
\texttt{nu\_p\_oneThousandNineHundredThirty\_two} & Gaps19221930 & PROVED & \checkmark & verified \\
\texttt{nu\_p\_oneThousandNineHundredTwentyEight} & Gaps19221930 & PROVED & \checkmark & verified \\
\texttt{nu\_p\_oneThousandNineHundredTwentyFour} & Gaps19221930 & PROVED & \checkmark & verified \\
\texttt{nu\_p\_oneThousandNineHundredTwentySix} & Gaps19221930 & PROVED & \checkmark & verified \\
\texttt{nu\_p\_oneThousandNineHundredTwentyTwo} & Gaps19221930 & PROVED & \checkmark & verified \\
\texttt{nu\_p\_oneThousandNineHundredTwentyTwo\_two} & Gaps19221930 & PROVED & \checkmark & verified \\
\texttt{singular\_series\_finite\_pos\_evenPair\_oneThousandNineHundredThirty} & Gaps19221930 & PROVED & \checkmark & verified \\
\texttt{singular\_series\_finite\_pos\_evenPair\_oneThousandNineHundredTwentyEight} & Gaps19221930 & PROVED & \checkmark & verified \\
\texttt{singular\_series\_finite\_pos\_evenPair\_oneThousandNineHundredTwentyFour} & Gaps19221930 & PROVED & \checkmark & verified \\
\texttt{singular\_series\_finite\_pos\_evenPair\_oneThousandNineHundredTwentySix} & Gaps19221930 & PROVED & \checkmark & verified \\
\texttt{singular\_series\_finite\_pos\_evenPair\_oneThousandNineHundredTwentyTwo} & Gaps19221930 & PROVED & \checkmark & verified \\
\texttt{singular\_series\_pos\_evenPair\_oneThousandNineHundredThirty} & Gaps19221930 & PROVED & \checkmark & verified \\
\texttt{singular\_series\_pos\_evenPair\_oneThousandNineHundredTwentyEight} & Gaps19221930 & PROVED & \checkmark & verified \\
\texttt{singular\_series\_pos\_evenPair\_oneThousandNineHundredTwentyFour} & Gaps19221930 & PROVED & \checkmark & verified \\
\texttt{singular\_series\_pos\_evenPair\_oneThousandNineHundredTwentySix} & Gaps19221930 & PROVED & \checkmark & verified \\
\texttt{singular\_series\_pos\_evenPair\_oneThousandNineHundredTwentyTwo} & Gaps19221930 & PROVED & \checkmark & verified \\
\texttt{evenPair\_card\_oneThousandNineHundredForty} & Gaps19321940 & PROVED & \checkmark & verified \\
\texttt{evenPair\_card\_oneThousandNineHundredThirtyEight} & Gaps19321940 & PROVED & \checkmark & verified \\
\texttt{evenPair\_card\_oneThousandNineHundredThirtyFour} & Gaps19321940 & PROVED & \checkmark & verified \\
\texttt{evenPair\_card\_oneThousandNineHundredThirtySix} & Gaps19321940 & PROVED & \checkmark & verified \\
\texttt{evenPair\_card\_oneThousandNineHundredThirtyTwo} & Gaps19321940 & PROVED & \checkmark & verified \\
\texttt{isAdmissible\_evenPair\_oneThousandNineHundredForty} & Gaps19321940 & PROVED & \checkmark & verified \\
\texttt{isAdmissible\_evenPair\_oneThousandNineHundredThirtyEight} & Gaps19321940 & PROVED & \checkmark & verified \\
\texttt{isAdmissible\_evenPair\_oneThousandNineHundredThirtyFour} & Gaps19321940 & PROVED & \checkmark & verified \\
\texttt{isAdmissible\_evenPair\_oneThousandNineHundredThirtySix} & Gaps19321940 & PROVED & \checkmark & verified \\
\texttt{isAdmissible\_evenPair\_oneThousandNineHundredThirtyTwo} & Gaps19321940 & PROVED & \checkmark & verified \\
\texttt{localFactor\_oneThousandNineHundredForty\_two} & Gaps19321940 & PROVED & \checkmark & verified \\
\texttt{localFactor\_oneThousandNineHundredThirtyTwo\_two} & Gaps19321940 & PROVED & \checkmark & verified \\
\texttt{nu\_p\_oneThousandNineHundredForty} & Gaps19321940 & PROVED & \checkmark & verified \\
\texttt{nu\_p\_oneThousandNineHundredForty\_two} & Gaps19321940 & PROVED & \checkmark & verified \\
\texttt{nu\_p\_oneThousandNineHundredThirtyEight} & Gaps19321940 & PROVED & \checkmark & verified \\
\texttt{nu\_p\_oneThousandNineHundredThirtyFour} & Gaps19321940 & PROVED & \checkmark & verified \\
\texttt{nu\_p\_oneThousandNineHundredThirtySix} & Gaps19321940 & PROVED & \checkmark & verified \\
\texttt{nu\_p\_oneThousandNineHundredThirtyTwo} & Gaps19321940 & PROVED & \checkmark & verified \\
\texttt{nu\_p\_oneThousandNineHundredThirtyTwo\_two} & Gaps19321940 & PROVED & \checkmark & verified \\
\texttt{singular\_series\_finite\_pos\_evenPair\_oneThousandNineHundredForty} & Gaps19321940 & PROVED & \checkmark & verified \\
\texttt{singular\_series\_finite\_pos\_evenPair\_oneThousandNineHundredThirtyEight} & Gaps19321940 & PROVED & \checkmark & verified \\
\texttt{singular\_series\_finite\_pos\_evenPair\_oneThousandNineHundredThirtyFour} & Gaps19321940 & PROVED & \checkmark & verified \\
\texttt{singular\_series\_finite\_pos\_evenPair\_oneThousandNineHundredThirtySix} & Gaps19321940 & PROVED & \checkmark & verified \\
\texttt{singular\_series\_finite\_pos\_evenPair\_oneThousandNineHundredThirtyTwo} & Gaps19321940 & PROVED & \checkmark & verified \\
\texttt{singular\_series\_pos\_evenPair\_oneThousandNineHundredForty} & Gaps19321940 & PROVED & \checkmark & verified \\
\texttt{singular\_series\_pos\_evenPair\_oneThousandNineHundredThirtyEight} & Gaps19321940 & PROVED & \checkmark & verified \\
\texttt{singular\_series\_pos\_evenPair\_oneThousandNineHundredThirtyFour} & Gaps19321940 & PROVED & \checkmark & verified \\
\texttt{singular\_series\_pos\_evenPair\_oneThousandNineHundredThirtySix} & Gaps19321940 & PROVED & \checkmark & verified \\
\texttt{singular\_series\_pos\_evenPair\_oneThousandNineHundredThirtyTwo} & Gaps19321940 & PROVED & \checkmark & verified \\
\texttt{evenPair\_card\_oneThousandNineHundredFifty} & Gaps19421950 & PROVED & \checkmark & verified \\
\texttt{evenPair\_card\_oneThousandNineHundredFortyEight} & Gaps19421950 & PROVED & \checkmark & verified \\
\texttt{evenPair\_card\_oneThousandNineHundredFortyFour} & Gaps19421950 & PROVED & \checkmark & verified \\
\texttt{evenPair\_card\_oneThousandNineHundredFortySix} & Gaps19421950 & PROVED & \checkmark & verified \\
\texttt{evenPair\_card\_oneThousandNineHundredFortyTwo} & Gaps19421950 & PROVED & \checkmark & verified \\
\texttt{isAdmissible\_evenPair\_oneThousandNineHundredFifty} & Gaps19421950 & PROVED & \checkmark & verified \\
\texttt{isAdmissible\_evenPair\_oneThousandNineHundredFortyEight} & Gaps19421950 & PROVED & \checkmark & verified \\
\texttt{isAdmissible\_evenPair\_oneThousandNineHundredFortyFour} & Gaps19421950 & PROVED & \checkmark & verified \\
\texttt{isAdmissible\_evenPair\_oneThousandNineHundredFortySix} & Gaps19421950 & PROVED & \checkmark & verified \\
\texttt{isAdmissible\_evenPair\_oneThousandNineHundredFortyTwo} & Gaps19421950 & PROVED & \checkmark & verified \\
\texttt{localFactor\_oneThousandNineHundredFifty\_two} & Gaps19421950 & PROVED & \checkmark & verified \\
\texttt{localFactor\_oneThousandNineHundredFortyTwo\_two} & Gaps19421950 & PROVED & \checkmark & verified \\
\texttt{nu\_p\_oneThousandNineHundredFifty} & Gaps19421950 & PROVED & \checkmark & verified \\
\texttt{nu\_p\_oneThousandNineHundredFifty\_two} & Gaps19421950 & PROVED & \checkmark & verified \\
\texttt{nu\_p\_oneThousandNineHundredFortyEight} & Gaps19421950 & PROVED & \checkmark & verified \\
\texttt{nu\_p\_oneThousandNineHundredFortyFour} & Gaps19421950 & PROVED & \checkmark & verified \\
\texttt{nu\_p\_oneThousandNineHundredFortySix} & Gaps19421950 & PROVED & \checkmark & verified \\
\texttt{nu\_p\_oneThousandNineHundredFortyTwo} & Gaps19421950 & PROVED & \checkmark & verified \\
\texttt{nu\_p\_oneThousandNineHundredFortyTwo\_two} & Gaps19421950 & PROVED & \checkmark & verified \\
\texttt{singular\_series\_finite\_pos\_evenPair\_oneThousandNineHundredFifty} & Gaps19421950 & PROVED & \checkmark & verified \\
\texttt{singular\_series\_finite\_pos\_evenPair\_oneThousandNineHundredFortyEight} & Gaps19421950 & PROVED & \checkmark & verified \\
\texttt{singular\_series\_finite\_pos\_evenPair\_oneThousandNineHundredFortyFour} & Gaps19421950 & PROVED & \checkmark & verified \\
\texttt{singular\_series\_finite\_pos\_evenPair\_oneThousandNineHundredFortySix} & Gaps19421950 & PROVED & \checkmark & verified \\
\texttt{singular\_series\_finite\_pos\_evenPair\_oneThousandNineHundredFortyTwo} & Gaps19421950 & PROVED & \checkmark & verified \\
\texttt{singular\_series\_pos\_evenPair\_oneThousandNineHundredFifty} & Gaps19421950 & PROVED & \checkmark & verified \\
\texttt{singular\_series\_pos\_evenPair\_oneThousandNineHundredFortyEight} & Gaps19421950 & PROVED & \checkmark & verified \\
\texttt{singular\_series\_pos\_evenPair\_oneThousandNineHundredFortyFour} & Gaps19421950 & PROVED & \checkmark & verified \\
\texttt{singular\_series\_pos\_evenPair\_oneThousandNineHundredFortySix} & Gaps19421950 & PROVED & \checkmark & verified \\
\texttt{singular\_series\_pos\_evenPair\_oneThousandNineHundredFortyTwo} & Gaps19421950 & PROVED & \checkmark & verified \\
\texttt{evenPair\_card\_oneThousandNineHundredFiftyEight} & Gaps19521960 & PROVED & \checkmark & verified \\
\texttt{evenPair\_card\_oneThousandNineHundredFiftyFour} & Gaps19521960 & PROVED & \checkmark & verified \\
\texttt{evenPair\_card\_oneThousandNineHundredFiftySix} & Gaps19521960 & PROVED & \checkmark & verified \\
\texttt{evenPair\_card\_oneThousandNineHundredFiftyTwo} & Gaps19521960 & PROVED & \checkmark & verified \\
\texttt{evenPair\_card\_oneThousandNineHundredSixty} & Gaps19521960 & PROVED & \checkmark & verified \\
\texttt{isAdmissible\_evenPair\_oneThousandNineHundredFiftyEight} & Gaps19521960 & PROVED & \checkmark & verified \\
\texttt{isAdmissible\_evenPair\_oneThousandNineHundredFiftyFour} & Gaps19521960 & PROVED & \checkmark & verified \\
\texttt{isAdmissible\_evenPair\_oneThousandNineHundredFiftySix} & Gaps19521960 & PROVED & \checkmark & verified \\
\texttt{isAdmissible\_evenPair\_oneThousandNineHundredFiftyTwo} & Gaps19521960 & PROVED & \checkmark & verified \\
\texttt{isAdmissible\_evenPair\_oneThousandNineHundredSixty} & Gaps19521960 & PROVED & \checkmark & verified \\
\texttt{localFactor\_oneThousandNineHundredFiftyTwo\_two} & Gaps19521960 & PROVED & \checkmark & verified \\
\texttt{localFactor\_oneThousandNineHundredSixty\_two} & Gaps19521960 & PROVED & \checkmark & verified \\
\texttt{nu\_p\_oneThousandNineHundredFiftyEight} & Gaps19521960 & PROVED & \checkmark & verified \\
\texttt{nu\_p\_oneThousandNineHundredFiftyFour} & Gaps19521960 & PROVED & \checkmark & verified \\
\texttt{nu\_p\_oneThousandNineHundredFiftySix} & Gaps19521960 & PROVED & \checkmark & verified \\
\texttt{nu\_p\_oneThousandNineHundredFiftyTwo} & Gaps19521960 & PROVED & \checkmark & verified \\
\texttt{nu\_p\_oneThousandNineHundredFiftyTwo\_two} & Gaps19521960 & PROVED & \checkmark & verified \\
\texttt{nu\_p\_oneThousandNineHundredSixty} & Gaps19521960 & PROVED & \checkmark & verified \\
\texttt{nu\_p\_oneThousandNineHundredSixty\_two} & Gaps19521960 & PROVED & \checkmark & verified \\
\texttt{singular\_series\_finite\_pos\_evenPair\_oneThousandNineHundredFiftyEight} & Gaps19521960 & PROVED & \checkmark & verified \\
\texttt{singular\_series\_finite\_pos\_evenPair\_oneThousandNineHundredFiftyFour} & Gaps19521960 & PROVED & \checkmark & verified \\
\texttt{singular\_series\_finite\_pos\_evenPair\_oneThousandNineHundredFiftySix} & Gaps19521960 & PROVED & \checkmark & verified \\
\texttt{singular\_series\_finite\_pos\_evenPair\_oneThousandNineHundredFiftyTwo} & Gaps19521960 & PROVED & \checkmark & verified \\
\texttt{singular\_series\_finite\_pos\_evenPair\_oneThousandNineHundredSixty} & Gaps19521960 & PROVED & \checkmark & verified \\
\texttt{singular\_series\_pos\_evenPair\_oneThousandNineHundredFiftyEight} & Gaps19521960 & PROVED & \checkmark & verified \\
\texttt{singular\_series\_pos\_evenPair\_oneThousandNineHundredFiftyFour} & Gaps19521960 & PROVED & \checkmark & verified \\
\texttt{singular\_series\_pos\_evenPair\_oneThousandNineHundredFiftySix} & Gaps19521960 & PROVED & \checkmark & verified \\
\texttt{singular\_series\_pos\_evenPair\_oneThousandNineHundredFiftyTwo} & Gaps19521960 & PROVED & \checkmark & verified \\
\texttt{singular\_series\_pos\_evenPair\_oneThousandNineHundredSixty} & Gaps19521960 & PROVED & \checkmark & verified \\
\texttt{evenPair\_card\_oneThousandNineHundredSeventy} & Gaps19621970 & PROVED & \checkmark & verified \\
\texttt{evenPair\_card\_oneThousandNineHundredSixtyEight} & Gaps19621970 & PROVED & \checkmark & verified \\
\texttt{evenPair\_card\_oneThousandNineHundredSixtyFour} & Gaps19621970 & PROVED & \checkmark & verified \\
\texttt{evenPair\_card\_oneThousandNineHundredSixtySix} & Gaps19621970 & PROVED & \checkmark & verified \\
\texttt{evenPair\_card\_oneThousandNineHundredSixtyTwo} & Gaps19621970 & PROVED & \checkmark & verified \\
\texttt{isAdmissible\_evenPair\_oneThousandNineHundredSeventy} & Gaps19621970 & PROVED & \checkmark & verified \\
\texttt{isAdmissible\_evenPair\_oneThousandNineHundredSixtyEight} & Gaps19621970 & PROVED & \checkmark & verified \\
\texttt{isAdmissible\_evenPair\_oneThousandNineHundredSixtyFour} & Gaps19621970 & PROVED & \checkmark & verified \\
\texttt{isAdmissible\_evenPair\_oneThousandNineHundredSixtySix} & Gaps19621970 & PROVED & \checkmark & verified \\
\texttt{isAdmissible\_evenPair\_oneThousandNineHundredSixtyTwo} & Gaps19621970 & PROVED & \checkmark & verified \\
\texttt{localFactor\_oneThousandNineHundredSeventy\_two} & Gaps19621970 & PROVED & \checkmark & verified \\
\texttt{localFactor\_oneThousandNineHundredSixtyTwo\_two} & Gaps19621970 & PROVED & \checkmark & verified \\
\texttt{nu\_p\_oneThousandNineHundredSeventy} & Gaps19621970 & PROVED & \checkmark & verified \\
\texttt{nu\_p\_oneThousandNineHundredSeventy\_two} & Gaps19621970 & PROVED & \checkmark & verified \\
\texttt{nu\_p\_oneThousandNineHundredSixtyEight} & Gaps19621970 & PROVED & \checkmark & verified \\
\texttt{nu\_p\_oneThousandNineHundredSixtyFour} & Gaps19621970 & PROVED & \checkmark & verified \\
\texttt{nu\_p\_oneThousandNineHundredSixtySix} & Gaps19621970 & PROVED & \checkmark & verified \\
\texttt{nu\_p\_oneThousandNineHundredSixtyTwo} & Gaps19621970 & PROVED & \checkmark & verified \\
\texttt{nu\_p\_oneThousandNineHundredSixtyTwo\_two} & Gaps19621970 & PROVED & \checkmark & verified \\
\texttt{singular\_series\_finite\_pos\_evenPair\_oneThousandNineHundredSeventy} & Gaps19621970 & PROVED & \checkmark & verified \\
\texttt{singular\_series\_finite\_pos\_evenPair\_oneThousandNineHundredSixtyEight} & Gaps19621970 & PROVED & \checkmark & verified \\
\texttt{singular\_series\_finite\_pos\_evenPair\_oneThousandNineHundredSixtyFour} & Gaps19621970 & PROVED & \checkmark & verified \\
\texttt{singular\_series\_finite\_pos\_evenPair\_oneThousandNineHundredSixtySix} & Gaps19621970 & PROVED & \checkmark & verified \\
\texttt{singular\_series\_finite\_pos\_evenPair\_oneThousandNineHundredSixtyTwo} & Gaps19621970 & PROVED & \checkmark & verified \\
\texttt{singular\_series\_pos\_evenPair\_oneThousandNineHundredSeventy} & Gaps19621970 & PROVED & \checkmark & verified \\
\texttt{singular\_series\_pos\_evenPair\_oneThousandNineHundredSixtyEight} & Gaps19621970 & PROVED & \checkmark & verified \\
\texttt{singular\_series\_pos\_evenPair\_oneThousandNineHundredSixtyFour} & Gaps19621970 & PROVED & \checkmark & verified \\
\texttt{singular\_series\_pos\_evenPair\_oneThousandNineHundredSixtySix} & Gaps19621970 & PROVED & \checkmark & verified \\
\texttt{singular\_series\_pos\_evenPair\_oneThousandNineHundredSixtyTwo} & Gaps19621970 & PROVED & \checkmark & verified \\
\texttt{evenPair\_card\_oneThousandNineHundredEighty} & Gaps19721980 & PROVED & \checkmark & verified \\
\texttt{evenPair\_card\_oneThousandNineHundredSeventyEight} & Gaps19721980 & PROVED & \checkmark & verified \\
\texttt{evenPair\_card\_oneThousandNineHundredSeventyFour} & Gaps19721980 & PROVED & \checkmark & verified \\
\texttt{evenPair\_card\_oneThousandNineHundredSeventySix} & Gaps19721980 & PROVED & \checkmark & verified \\
\texttt{evenPair\_card\_oneThousandNineHundredSeventyTwo} & Gaps19721980 & PROVED & \checkmark & verified \\
\texttt{isAdmissible\_evenPair\_oneThousandNineHundredEighty} & Gaps19721980 & PROVED & \checkmark & verified \\
\texttt{isAdmissible\_evenPair\_oneThousandNineHundredSeventyEight} & Gaps19721980 & PROVED & \checkmark & verified \\
\texttt{isAdmissible\_evenPair\_oneThousandNineHundredSeventyFour} & Gaps19721980 & PROVED & \checkmark & verified \\
\texttt{isAdmissible\_evenPair\_oneThousandNineHundredSeventySix} & Gaps19721980 & PROVED & \checkmark & verified \\
\texttt{isAdmissible\_evenPair\_oneThousandNineHundredSeventyTwo} & Gaps19721980 & PROVED & \checkmark & verified \\
\texttt{localFactor\_oneThousandNineHundredEighty\_two} & Gaps19721980 & PROVED & \checkmark & verified \\
\texttt{localFactor\_oneThousandNineHundredSeventyTwo\_two} & Gaps19721980 & PROVED & \checkmark & verified \\
\texttt{nu\_p\_oneThousandNineHundredEighty} & Gaps19721980 & PROVED & \checkmark & verified \\
\texttt{nu\_p\_oneThousandNineHundredEighty\_two} & Gaps19721980 & PROVED & \checkmark & verified \\
\texttt{nu\_p\_oneThousandNineHundredSeventyEight} & Gaps19721980 & PROVED & \checkmark & verified \\
\texttt{nu\_p\_oneThousandNineHundredSeventyFour} & Gaps19721980 & PROVED & \checkmark & verified \\
\texttt{nu\_p\_oneThousandNineHundredSeventySix} & Gaps19721980 & PROVED & \checkmark & verified \\
\texttt{nu\_p\_oneThousandNineHundredSeventyTwo} & Gaps19721980 & PROVED & \checkmark & verified \\
\texttt{nu\_p\_oneThousandNineHundredSeventyTwo\_two} & Gaps19721980 & PROVED & \checkmark & verified \\
\texttt{singular\_series\_finite\_pos\_evenPair\_oneThousandNineHundredEighty} & Gaps19721980 & PROVED & \checkmark & verified \\
\texttt{singular\_series\_finite\_pos\_evenPair\_oneThousandNineHundredSeventyEight} & Gaps19721980 & PROVED & \checkmark & verified \\
\texttt{singular\_series\_finite\_pos\_evenPair\_oneThousandNineHundredSeventyFour} & Gaps19721980 & PROVED & \checkmark & verified \\
\texttt{singular\_series\_finite\_pos\_evenPair\_oneThousandNineHundredSeventySix} & Gaps19721980 & PROVED & \checkmark & verified \\
\texttt{singular\_series\_finite\_pos\_evenPair\_oneThousandNineHundredSeventyTwo} & Gaps19721980 & PROVED & \checkmark & verified \\
\texttt{singular\_series\_pos\_evenPair\_oneThousandNineHundredEighty} & Gaps19721980 & PROVED & \checkmark & verified \\
\texttt{singular\_series\_pos\_evenPair\_oneThousandNineHundredSeventyEight} & Gaps19721980 & PROVED & \checkmark & verified \\
\texttt{singular\_series\_pos\_evenPair\_oneThousandNineHundredSeventyFour} & Gaps19721980 & PROVED & \checkmark & verified \\
\texttt{singular\_series\_pos\_evenPair\_oneThousandNineHundredSeventySix} & Gaps19721980 & PROVED & \checkmark & verified \\
\texttt{singular\_series\_pos\_evenPair\_oneThousandNineHundredSeventyTwo} & Gaps19721980 & PROVED & \checkmark & verified \\
\texttt{evenPair\_card\_oneThousandNineHundredEightyEight} & Gaps19821990 & PROVED & \checkmark & verified \\
\texttt{evenPair\_card\_oneThousandNineHundredEightyFour} & Gaps19821990 & PROVED & \checkmark & verified \\
\texttt{evenPair\_card\_oneThousandNineHundredEightySix} & Gaps19821990 & PROVED & \checkmark & verified \\
\texttt{evenPair\_card\_oneThousandNineHundredEightyTwo} & Gaps19821990 & PROVED & \checkmark & verified \\
\texttt{evenPair\_card\_oneThousandNineHundredNinety} & Gaps19821990 & PROVED & \checkmark & verified \\
\texttt{isAdmissible\_evenPair\_oneThousandNineHundredEightyEight} & Gaps19821990 & PROVED & \checkmark & verified \\
\texttt{isAdmissible\_evenPair\_oneThousandNineHundredEightyFour} & Gaps19821990 & PROVED & \checkmark & verified \\
\texttt{isAdmissible\_evenPair\_oneThousandNineHundredEightySix} & Gaps19821990 & PROVED & \checkmark & verified \\
\texttt{isAdmissible\_evenPair\_oneThousandNineHundredEightyTwo} & Gaps19821990 & PROVED & \checkmark & verified \\
\texttt{isAdmissible\_evenPair\_oneThousandNineHundredNinety} & Gaps19821990 & PROVED & \checkmark & verified \\
\texttt{localFactor\_oneThousandNineHundredEightyTwo\_two} & Gaps19821990 & PROVED & \checkmark & verified \\
\texttt{localFactor\_oneThousandNineHundredNinety\_two} & Gaps19821990 & PROVED & \checkmark & verified \\
\texttt{nu\_p\_oneThousandNineHundredEightyEight} & Gaps19821990 & PROVED & \checkmark & verified \\
\texttt{nu\_p\_oneThousandNineHundredEightyFour} & Gaps19821990 & PROVED & \checkmark & verified \\
\texttt{nu\_p\_oneThousandNineHundredEightySix} & Gaps19821990 & PROVED & \checkmark & verified \\
\texttt{nu\_p\_oneThousandNineHundredEightyTwo} & Gaps19821990 & PROVED & \checkmark & verified \\
\texttt{nu\_p\_oneThousandNineHundredEightyTwo\_two} & Gaps19821990 & PROVED & \checkmark & verified \\
\texttt{nu\_p\_oneThousandNineHundredNinety} & Gaps19821990 & PROVED & \checkmark & verified \\
\texttt{nu\_p\_oneThousandNineHundredNinety\_two} & Gaps19821990 & PROVED & \checkmark & verified \\
\texttt{singular\_series\_finite\_pos\_evenPair\_oneThousandNineHundredEightyEight} & Gaps19821990 & PROVED & \checkmark & verified \\
\texttt{singular\_series\_finite\_pos\_evenPair\_oneThousandNineHundredEightyFour} & Gaps19821990 & PROVED & \checkmark & verified \\
\texttt{singular\_series\_finite\_pos\_evenPair\_oneThousandNineHundredEightySix} & Gaps19821990 & PROVED & \checkmark & verified \\
\texttt{singular\_series\_finite\_pos\_evenPair\_oneThousandNineHundredEightyTwo} & Gaps19821990 & PROVED & \checkmark & verified \\
\texttt{singular\_series\_finite\_pos\_evenPair\_oneThousandNineHundredNinety} & Gaps19821990 & PROVED & \checkmark & verified \\
\texttt{singular\_series\_pos\_evenPair\_oneThousandNineHundredEightyEight} & Gaps19821990 & PROVED & \checkmark & verified \\
\texttt{singular\_series\_pos\_evenPair\_oneThousandNineHundredEightyFour} & Gaps19821990 & PROVED & \checkmark & verified \\
\texttt{singular\_series\_pos\_evenPair\_oneThousandNineHundredEightySix} & Gaps19821990 & PROVED & \checkmark & verified \\
\texttt{singular\_series\_pos\_evenPair\_oneThousandNineHundredEightyTwo} & Gaps19821990 & PROVED & \checkmark & verified \\
\texttt{singular\_series\_pos\_evenPair\_oneThousandNineHundredNinety} & Gaps19821990 & PROVED & \checkmark & verified \\
\texttt{evenPair\_card\_oneThousandNineHundredNinetyEight} & Gaps19922000 & PROVED & \checkmark & verified \\
\texttt{evenPair\_card\_oneThousandNineHundredNinetyFour} & Gaps19922000 & PROVED & \checkmark & verified \\
\texttt{evenPair\_card\_oneThousandNineHundredNinetySix} & Gaps19922000 & PROVED & \checkmark & verified \\
\texttt{evenPair\_card\_oneThousandNineHundredNinetyTwo} & Gaps19922000 & PROVED & \checkmark & verified \\
\texttt{evenPair\_card\_twoThousand} & Gaps19922000 & PROVED & \checkmark & verified \\
\texttt{isAdmissible\_evenPair\_oneThousandNineHundredNinetyEight} & Gaps19922000 & PROVED & \checkmark & verified \\
\texttt{isAdmissible\_evenPair\_oneThousandNineHundredNinetyFour} & Gaps19922000 & PROVED & \checkmark & verified \\
\texttt{isAdmissible\_evenPair\_oneThousandNineHundredNinetySix} & Gaps19922000 & PROVED & \checkmark & verified \\
\texttt{isAdmissible\_evenPair\_oneThousandNineHundredNinetyTwo} & Gaps19922000 & PROVED & \checkmark & verified \\
\texttt{isAdmissible\_evenPair\_twoThousand} & Gaps19922000 & PROVED & \checkmark & verified \\
\texttt{localFactor\_oneThousandNineHundredNinetyTwo\_two} & Gaps19922000 & PROVED & \checkmark & verified \\
\texttt{localFactor\_twoThousand\_two} & Gaps19922000 & PROVED & \checkmark & verified \\
\texttt{nu\_p\_oneThousandNineHundredNinetyEight} & Gaps19922000 & PROVED & \checkmark & verified \\
\texttt{nu\_p\_oneThousandNineHundredNinetyFour} & Gaps19922000 & PROVED & \checkmark & verified \\
\texttt{nu\_p\_oneThousandNineHundredNinetySix} & Gaps19922000 & PROVED & \checkmark & verified \\
\texttt{nu\_p\_oneThousandNineHundredNinetyTwo} & Gaps19922000 & PROVED & \checkmark & verified \\
\texttt{nu\_p\_oneThousandNineHundredNinetyTwo\_two} & Gaps19922000 & PROVED & \checkmark & verified \\
\texttt{nu\_p\_twoThousand} & Gaps19922000 & PROVED & \checkmark & verified \\
\texttt{nu\_p\_twoThousand\_two} & Gaps19922000 & PROVED & \checkmark & verified \\
\texttt{singular\_series\_finite\_pos\_evenPair\_oneThousandNineHundredNinetyEight} & Gaps19922000 & PROVED & \checkmark & verified \\
\texttt{singular\_series\_finite\_pos\_evenPair\_oneThousandNineHundredNinetyFour} & Gaps19922000 & PROVED & \checkmark & verified \\
\texttt{singular\_series\_finite\_pos\_evenPair\_oneThousandNineHundredNinetySix} & Gaps19922000 & PROVED & \checkmark & verified \\
\texttt{singular\_series\_finite\_pos\_evenPair\_oneThousandNineHundredNinetyTwo} & Gaps19922000 & PROVED & \checkmark & verified \\
\texttt{singular\_series\_finite\_pos\_evenPair\_twoThousand} & Gaps19922000 & PROVED & \checkmark & verified \\
\texttt{singular\_series\_pos\_evenPair\_oneThousandNineHundredNinetyEight} & Gaps19922000 & PROVED & \checkmark & verified \\
\texttt{singular\_series\_pos\_evenPair\_oneThousandNineHundredNinetyFour} & Gaps19922000 & PROVED & \checkmark & verified \\
\texttt{singular\_series\_pos\_evenPair\_oneThousandNineHundredNinetySix} & Gaps19922000 & PROVED & \checkmark & verified \\
\texttt{singular\_series\_pos\_evenPair\_oneThousandNineHundredNinetyTwo} & Gaps19922000 & PROVED & \checkmark & verified \\
\texttt{singular\_series\_pos\_evenPair\_twoThousand} & Gaps19922000 & PROVED & \checkmark & verified \\
\texttt{evenPair\_card\_twoThousandEight} & Gaps20022010 & PROVED & \checkmark & verified \\
\texttt{evenPair\_card\_twoThousandFour} & Gaps20022010 & PROVED & \checkmark & verified \\
\texttt{evenPair\_card\_twoThousandSix} & Gaps20022010 & PROVED & \checkmark & verified \\
\texttt{evenPair\_card\_twoThousandTen} & Gaps20022010 & PROVED & \checkmark & verified \\
\texttt{evenPair\_card\_twoThousandTwo} & Gaps20022010 & PROVED & \checkmark & verified \\
\texttt{isAdmissible\_evenPair\_twoThousandEight} & Gaps20022010 & PROVED & \checkmark & verified \\
\texttt{isAdmissible\_evenPair\_twoThousandFour} & Gaps20022010 & PROVED & \checkmark & verified \\
\texttt{isAdmissible\_evenPair\_twoThousandSix} & Gaps20022010 & PROVED & \checkmark & verified \\
\texttt{isAdmissible\_evenPair\_twoThousandTen} & Gaps20022010 & PROVED & \checkmark & verified \\
\texttt{isAdmissible\_evenPair\_twoThousandTwo} & Gaps20022010 & PROVED & \checkmark & verified \\
\texttt{localFactor\_twoThousandTen\_two} & Gaps20022010 & PROVED & \checkmark & verified \\
\texttt{localFactor\_twoThousandTwo\_two} & Gaps20022010 & PROVED & \checkmark & verified \\
\texttt{nu\_p\_twoThousandEight} & Gaps20022010 & PROVED & \checkmark & verified \\
\texttt{nu\_p\_twoThousandFour} & Gaps20022010 & PROVED & \checkmark & verified \\
\texttt{nu\_p\_twoThousandSix} & Gaps20022010 & PROVED & \checkmark & verified \\
\texttt{nu\_p\_twoThousandTen} & Gaps20022010 & PROVED & \checkmark & verified \\
\texttt{nu\_p\_twoThousandTen\_two} & Gaps20022010 & PROVED & \checkmark & verified \\
\texttt{nu\_p\_twoThousandTwo} & Gaps20022010 & PROVED & \checkmark & verified \\
\texttt{nu\_p\_twoThousandTwo\_two} & Gaps20022010 & PROVED & \checkmark & verified \\
\texttt{singular\_series\_finite\_pos\_evenPair\_twoThousandEight} & Gaps20022010 & PROVED & \checkmark & verified \\
\texttt{singular\_series\_finite\_pos\_evenPair\_twoThousandFour} & Gaps20022010 & PROVED & \checkmark & verified \\
\texttt{singular\_series\_finite\_pos\_evenPair\_twoThousandSix} & Gaps20022010 & PROVED & \checkmark & verified \\
\texttt{singular\_series\_finite\_pos\_evenPair\_twoThousandTen} & Gaps20022010 & PROVED & \checkmark & verified \\
\texttt{singular\_series\_finite\_pos\_evenPair\_twoThousandTwo} & Gaps20022010 & PROVED & \checkmark & verified \\
\texttt{singular\_series\_pos\_evenPair\_twoThousandEight} & Gaps20022010 & PROVED & \checkmark & verified \\
\texttt{singular\_series\_pos\_evenPair\_twoThousandFour} & Gaps20022010 & PROVED & \checkmark & verified \\
\texttt{singular\_series\_pos\_evenPair\_twoThousandSix} & Gaps20022010 & PROVED & \checkmark & verified \\
\texttt{singular\_series\_pos\_evenPair\_twoThousandTen} & Gaps20022010 & PROVED & \checkmark & verified \\
\texttt{singular\_series\_pos\_evenPair\_twoThousandTwo} & Gaps20022010 & PROVED & \checkmark & verified \\
\texttt{evenPair\_card\_twoThousandEighteen} & Gaps20122020 & PROVED & \checkmark & verified \\
\texttt{evenPair\_card\_twoThousandFourteen} & Gaps20122020 & PROVED & \checkmark & verified \\
\texttt{evenPair\_card\_twoThousandSixteen} & Gaps20122020 & PROVED & \checkmark & verified \\
\texttt{evenPair\_card\_twoThousandTwelve} & Gaps20122020 & PROVED & \checkmark & verified \\
\texttt{evenPair\_card\_twoThousandTwenty} & Gaps20122020 & PROVED & \checkmark & verified \\
\texttt{isAdmissible\_evenPair\_twoThousandEighteen} & Gaps20122020 & PROVED & \checkmark & verified \\
\texttt{isAdmissible\_evenPair\_twoThousandFourteen} & Gaps20122020 & PROVED & \checkmark & verified \\
\texttt{isAdmissible\_evenPair\_twoThousandSixteen} & Gaps20122020 & PROVED & \checkmark & verified \\
\texttt{isAdmissible\_evenPair\_twoThousandTwelve} & Gaps20122020 & PROVED & \checkmark & verified \\
\texttt{isAdmissible\_evenPair\_twoThousandTwenty} & Gaps20122020 & PROVED & \checkmark & verified \\
\texttt{localFactor\_twoThousandTwelve\_two} & Gaps20122020 & PROVED & \checkmark & verified \\
\texttt{localFactor\_twoThousandTwenty\_two} & Gaps20122020 & PROVED & \checkmark & verified \\
\texttt{nu\_p\_twoThousandEighteen} & Gaps20122020 & PROVED & \checkmark & verified \\
\texttt{nu\_p\_twoThousandFourteen} & Gaps20122020 & PROVED & \checkmark & verified \\
\texttt{nu\_p\_twoThousandSixteen} & Gaps20122020 & PROVED & \checkmark & verified \\
\texttt{nu\_p\_twoThousandTwelve} & Gaps20122020 & PROVED & \checkmark & verified \\
\texttt{nu\_p\_twoThousandTwelve\_two} & Gaps20122020 & PROVED & \checkmark & verified \\
\texttt{nu\_p\_twoThousandTwenty} & Gaps20122020 & PROVED & \checkmark & verified \\
\texttt{nu\_p\_twoThousandTwenty\_two} & Gaps20122020 & PROVED & \checkmark & verified \\
\texttt{singular\_series\_finite\_pos\_evenPair\_twoThousandEighteen} & Gaps20122020 & PROVED & \checkmark & verified \\
\texttt{singular\_series\_finite\_pos\_evenPair\_twoThousandFourteen} & Gaps20122020 & PROVED & \checkmark & verified \\
\texttt{singular\_series\_finite\_pos\_evenPair\_twoThousandSixteen} & Gaps20122020 & PROVED & \checkmark & verified \\
\texttt{singular\_series\_finite\_pos\_evenPair\_twoThousandTwelve} & Gaps20122020 & PROVED & \checkmark & verified \\
\texttt{singular\_series\_finite\_pos\_evenPair\_twoThousandTwenty} & Gaps20122020 & PROVED & \checkmark & verified \\
\texttt{singular\_series\_pos\_evenPair\_twoThousandEighteen} & Gaps20122020 & PROVED & \checkmark & verified \\
\texttt{singular\_series\_pos\_evenPair\_twoThousandFourteen} & Gaps20122020 & PROVED & \checkmark & verified \\
\texttt{singular\_series\_pos\_evenPair\_twoThousandSixteen} & Gaps20122020 & PROVED & \checkmark & verified \\
\texttt{singular\_series\_pos\_evenPair\_twoThousandTwelve} & Gaps20122020 & PROVED & \checkmark & verified \\
\texttt{singular\_series\_pos\_evenPair\_twoThousandTwenty} & Gaps20122020 & PROVED & \checkmark & verified \\
\texttt{evenPair\_card\_twoHundredEight} & Gaps202210 & PROVED & \checkmark & verified \\
\texttt{evenPair\_card\_twoHundredFour} & Gaps202210 & PROVED & \checkmark & verified \\
\texttt{evenPair\_card\_twoHundredSix} & Gaps202210 & PROVED & \checkmark & verified \\
\texttt{evenPair\_card\_twoHundredTen} & Gaps202210 & PROVED & \checkmark & verified \\
\texttt{evenPair\_card\_twoHundredTwo} & Gaps202210 & PROVED & \checkmark & verified \\
\texttt{isAdmissible\_evenPair\_twoHundredEight} & Gaps202210 & PROVED & \checkmark & verified \\
\texttt{isAdmissible\_evenPair\_twoHundredFour} & Gaps202210 & PROVED & \checkmark & verified \\
\texttt{isAdmissible\_evenPair\_twoHundredSix} & Gaps202210 & PROVED & \checkmark & verified \\
\texttt{isAdmissible\_evenPair\_twoHundredTen} & Gaps202210 & PROVED & \checkmark & verified \\
\texttt{isAdmissible\_evenPair\_twoHundredTwo} & Gaps202210 & PROVED & \checkmark & verified \\
\texttt{localFactor\_twoHundredTen\_two} & Gaps202210 & PROVED & \checkmark & verified \\
\texttt{localFactor\_twoHundredTwo\_two} & Gaps202210 & PROVED & \checkmark & verified \\
\texttt{nu\_p\_twoHundredEight} & Gaps202210 & PROVED & \checkmark & verified \\
\texttt{nu\_p\_twoHundredFour} & Gaps202210 & PROVED & \checkmark & verified \\
\texttt{nu\_p\_twoHundredSix} & Gaps202210 & PROVED & \checkmark & verified \\
\texttt{nu\_p\_twoHundredTen} & Gaps202210 & PROVED & \checkmark & verified \\
\texttt{nu\_p\_twoHundredTen\_two} & Gaps202210 & PROVED & \checkmark & verified \\
\texttt{nu\_p\_twoHundredTwo} & Gaps202210 & PROVED & \checkmark & verified \\
\texttt{nu\_p\_twoHundredTwo\_two} & Gaps202210 & PROVED & \checkmark & verified \\
\texttt{singular\_series\_finite\_pos\_evenPair\_twoHundredEight} & Gaps202210 & PROVED & \checkmark & verified \\
\texttt{singular\_series\_finite\_pos\_evenPair\_twoHundredFour} & Gaps202210 & PROVED & \checkmark & verified \\
\texttt{singular\_series\_finite\_pos\_evenPair\_twoHundredSix} & Gaps202210 & PROVED & \checkmark & verified \\
\texttt{singular\_series\_finite\_pos\_evenPair\_twoHundredTen} & Gaps202210 & PROVED & \checkmark & verified \\
\texttt{singular\_series\_finite\_pos\_evenPair\_twoHundredTwo} & Gaps202210 & PROVED & \checkmark & verified \\
\texttt{singular\_series\_pos\_evenPair\_twoHundredEight} & Gaps202210 & PROVED & \checkmark & verified \\
\texttt{singular\_series\_pos\_evenPair\_twoHundredFour} & Gaps202210 & PROVED & \checkmark & verified \\
\texttt{singular\_series\_pos\_evenPair\_twoHundredSix} & Gaps202210 & PROVED & \checkmark & verified \\
\texttt{singular\_series\_pos\_evenPair\_twoHundredTen} & Gaps202210 & PROVED & \checkmark & verified \\
\texttt{singular\_series\_pos\_evenPair\_twoHundredTwo} & Gaps202210 & PROVED & \checkmark & verified \\
\texttt{evenPair\_card\_twoThousandThirty} & Gaps20222030 & PROVED & \checkmark & verified \\
\texttt{evenPair\_card\_twoThousandTwentyEight} & Gaps20222030 & PROVED & \checkmark & verified \\
\texttt{evenPair\_card\_twoThousandTwentyFour} & Gaps20222030 & PROVED & \checkmark & verified \\
\texttt{evenPair\_card\_twoThousandTwentySix} & Gaps20222030 & PROVED & \checkmark & verified \\
\texttt{evenPair\_card\_twoThousandTwentyTwo} & Gaps20222030 & PROVED & \checkmark & verified \\
\texttt{isAdmissible\_evenPair\_twoThousandThirty} & Gaps20222030 & PROVED & \checkmark & verified \\
\texttt{isAdmissible\_evenPair\_twoThousandTwentyEight} & Gaps20222030 & PROVED & \checkmark & verified \\
\texttt{isAdmissible\_evenPair\_twoThousandTwentyFour} & Gaps20222030 & PROVED & \checkmark & verified \\
\texttt{isAdmissible\_evenPair\_twoThousandTwentySix} & Gaps20222030 & PROVED & \checkmark & verified \\
\texttt{isAdmissible\_evenPair\_twoThousandTwentyTwo} & Gaps20222030 & PROVED & \checkmark & verified \\
\texttt{localFactor\_twoThousandThirty\_two} & Gaps20222030 & PROVED & \checkmark & verified \\
\texttt{localFactor\_twoThousandTwentyTwo\_two} & Gaps20222030 & PROVED & \checkmark & verified \\
\texttt{nu\_p\_twoThousandThirty} & Gaps20222030 & PROVED & \checkmark & verified \\
\texttt{nu\_p\_twoThousandThirty\_two} & Gaps20222030 & PROVED & \checkmark & verified \\
\texttt{nu\_p\_twoThousandTwentyEight} & Gaps20222030 & PROVED & \checkmark & verified \\
\texttt{nu\_p\_twoThousandTwentyFour} & Gaps20222030 & PROVED & \checkmark & verified \\
\texttt{nu\_p\_twoThousandTwentySix} & Gaps20222030 & PROVED & \checkmark & verified \\
\texttt{nu\_p\_twoThousandTwentyTwo} & Gaps20222030 & PROVED & \checkmark & verified \\
\texttt{nu\_p\_twoThousandTwentyTwo\_two} & Gaps20222030 & PROVED & \checkmark & verified \\
\texttt{singular\_series\_finite\_pos\_evenPair\_twoThousandThirty} & Gaps20222030 & PROVED & \checkmark & verified \\
\texttt{singular\_series\_finite\_pos\_evenPair\_twoThousandTwentyEight} & Gaps20222030 & PROVED & \checkmark & verified \\
\texttt{singular\_series\_finite\_pos\_evenPair\_twoThousandTwentyFour} & Gaps20222030 & PROVED & \checkmark & verified \\
\texttt{singular\_series\_finite\_pos\_evenPair\_twoThousandTwentySix} & Gaps20222030 & PROVED & \checkmark & verified \\
\texttt{singular\_series\_finite\_pos\_evenPair\_twoThousandTwentyTwo} & Gaps20222030 & PROVED & \checkmark & verified \\
\texttt{singular\_series\_pos\_evenPair\_twoThousandThirty} & Gaps20222030 & PROVED & \checkmark & verified \\
\texttt{singular\_series\_pos\_evenPair\_twoThousandTwentyEight} & Gaps20222030 & PROVED & \checkmark & verified \\
\texttt{singular\_series\_pos\_evenPair\_twoThousandTwentyFour} & Gaps20222030 & PROVED & \checkmark & verified \\
\texttt{singular\_series\_pos\_evenPair\_twoThousandTwentySix} & Gaps20222030 & PROVED & \checkmark & verified \\
\texttt{singular\_series\_pos\_evenPair\_twoThousandTwentyTwo} & Gaps20222030 & PROVED & \checkmark & verified \\
\texttt{evenPair\_card\_twoThousandForty} & Gaps20322040 & PROVED & \checkmark & verified \\
\texttt{evenPair\_card\_twoThousandThirtyEight} & Gaps20322040 & PROVED & \checkmark & verified \\
\texttt{evenPair\_card\_twoThousandThirtyFour} & Gaps20322040 & PROVED & \checkmark & verified \\
\texttt{evenPair\_card\_twoThousandThirtySix} & Gaps20322040 & PROVED & \checkmark & verified \\
\texttt{evenPair\_card\_twoThousandThirtyTwo} & Gaps20322040 & PROVED & \checkmark & verified \\
\texttt{isAdmissible\_evenPair\_twoThousandForty} & Gaps20322040 & PROVED & \checkmark & verified \\
\texttt{isAdmissible\_evenPair\_twoThousandThirtyEight} & Gaps20322040 & PROVED & \checkmark & verified \\
\texttt{isAdmissible\_evenPair\_twoThousandThirtyFour} & Gaps20322040 & PROVED & \checkmark & verified \\
\texttt{isAdmissible\_evenPair\_twoThousandThirtySix} & Gaps20322040 & PROVED & \checkmark & verified \\
\texttt{isAdmissible\_evenPair\_twoThousandThirtyTwo} & Gaps20322040 & PROVED & \checkmark & verified \\
\texttt{localFactor\_twoThousandForty\_two} & Gaps20322040 & PROVED & \checkmark & verified \\
\texttt{localFactor\_twoThousandThirtyTwo\_two} & Gaps20322040 & PROVED & \checkmark & verified \\
\texttt{nu\_p\_twoThousandForty} & Gaps20322040 & PROVED & \checkmark & verified \\
\texttt{nu\_p\_twoThousandForty\_two} & Gaps20322040 & PROVED & \checkmark & verified \\
\texttt{nu\_p\_twoThousandThirtyEight} & Gaps20322040 & PROVED & \checkmark & verified \\
\texttt{nu\_p\_twoThousandThirtyFour} & Gaps20322040 & PROVED & \checkmark & verified \\
\texttt{nu\_p\_twoThousandThirtySix} & Gaps20322040 & PROVED & \checkmark & verified \\
\texttt{nu\_p\_twoThousandThirtyTwo} & Gaps20322040 & PROVED & \checkmark & verified \\
\texttt{nu\_p\_twoThousandThirtyTwo\_two} & Gaps20322040 & PROVED & \checkmark & verified \\
\texttt{singular\_series\_finite\_pos\_evenPair\_twoThousandForty} & Gaps20322040 & PROVED & \checkmark & verified \\
\texttt{singular\_series\_finite\_pos\_evenPair\_twoThousandThirtyEight} & Gaps20322040 & PROVED & \checkmark & verified \\
\texttt{singular\_series\_finite\_pos\_evenPair\_twoThousandThirtyFour} & Gaps20322040 & PROVED & \checkmark & verified \\
\texttt{singular\_series\_finite\_pos\_evenPair\_twoThousandThirtySix} & Gaps20322040 & PROVED & \checkmark & verified \\
\texttt{singular\_series\_finite\_pos\_evenPair\_twoThousandThirtyTwo} & Gaps20322040 & PROVED & \checkmark & verified \\
\texttt{singular\_series\_pos\_evenPair\_twoThousandForty} & Gaps20322040 & PROVED & \checkmark & verified \\
\texttt{singular\_series\_pos\_evenPair\_twoThousandThirtyEight} & Gaps20322040 & PROVED & \checkmark & verified \\
\texttt{singular\_series\_pos\_evenPair\_twoThousandThirtyFour} & Gaps20322040 & PROVED & \checkmark & verified \\
\texttt{singular\_series\_pos\_evenPair\_twoThousandThirtySix} & Gaps20322040 & PROVED & \checkmark & verified \\
\texttt{singular\_series\_pos\_evenPair\_twoThousandThirtyTwo} & Gaps20322040 & PROVED & \checkmark & verified \\
\texttt{evenPair\_card\_twoThousandFifty} & Gaps20422050 & PROVED & \checkmark & verified \\
\texttt{evenPair\_card\_twoThousandFortyEight} & Gaps20422050 & PROVED & \checkmark & verified \\
\texttt{evenPair\_card\_twoThousandFortyFour} & Gaps20422050 & PROVED & \checkmark & verified \\
\texttt{evenPair\_card\_twoThousandFortySix} & Gaps20422050 & PROVED & \checkmark & verified \\
\texttt{evenPair\_card\_twoThousandFortyTwo} & Gaps20422050 & PROVED & \checkmark & verified \\
\texttt{isAdmissible\_evenPair\_twoThousandFifty} & Gaps20422050 & PROVED & \checkmark & verified \\
\texttt{isAdmissible\_evenPair\_twoThousandFortyEight} & Gaps20422050 & PROVED & \checkmark & verified \\
\texttt{isAdmissible\_evenPair\_twoThousandFortyFour} & Gaps20422050 & PROVED & \checkmark & verified \\
\texttt{isAdmissible\_evenPair\_twoThousandFortySix} & Gaps20422050 & PROVED & \checkmark & verified \\
\texttt{isAdmissible\_evenPair\_twoThousandFortyTwo} & Gaps20422050 & PROVED & \checkmark & verified \\
\texttt{localFactor\_twoThousandFifty\_two} & Gaps20422050 & PROVED & \checkmark & verified \\
\texttt{localFactor\_twoThousandFortyTwo\_two} & Gaps20422050 & PROVED & \checkmark & verified \\
\texttt{nu\_p\_twoThousandFifty} & Gaps20422050 & PROVED & \checkmark & verified \\
\texttt{nu\_p\_twoThousandFifty\_two} & Gaps20422050 & PROVED & \checkmark & verified \\
\texttt{nu\_p\_twoThousandFortyEight} & Gaps20422050 & PROVED & \checkmark & verified \\
\texttt{nu\_p\_twoThousandFortyFour} & Gaps20422050 & PROVED & \checkmark & verified \\
\texttt{nu\_p\_twoThousandFortySix} & Gaps20422050 & PROVED & \checkmark & verified \\
\texttt{nu\_p\_twoThousandFortyTwo} & Gaps20422050 & PROVED & \checkmark & verified \\
\texttt{nu\_p\_twoThousandFortyTwo\_two} & Gaps20422050 & PROVED & \checkmark & verified \\
\texttt{singular\_series\_finite\_pos\_evenPair\_twoThousandFifty} & Gaps20422050 & PROVED & \checkmark & verified \\
\texttt{singular\_series\_finite\_pos\_evenPair\_twoThousandFortyEight} & Gaps20422050 & PROVED & \checkmark & verified \\
\texttt{singular\_series\_finite\_pos\_evenPair\_twoThousandFortyFour} & Gaps20422050 & PROVED & \checkmark & verified \\
\texttt{singular\_series\_finite\_pos\_evenPair\_twoThousandFortySix} & Gaps20422050 & PROVED & \checkmark & verified \\
\texttt{singular\_series\_finite\_pos\_evenPair\_twoThousandFortyTwo} & Gaps20422050 & PROVED & \checkmark & verified \\
\texttt{singular\_series\_pos\_evenPair\_twoThousandFifty} & Gaps20422050 & PROVED & \checkmark & verified \\
\texttt{singular\_series\_pos\_evenPair\_twoThousandFortyEight} & Gaps20422050 & PROVED & \checkmark & verified \\
\texttt{singular\_series\_pos\_evenPair\_twoThousandFortyFour} & Gaps20422050 & PROVED & \checkmark & verified \\
\texttt{singular\_series\_pos\_evenPair\_twoThousandFortySix} & Gaps20422050 & PROVED & \checkmark & verified \\
\texttt{singular\_series\_pos\_evenPair\_twoThousandFortyTwo} & Gaps20422050 & PROVED & \checkmark & verified \\
\texttt{evenPair\_card\_twoThousandFiftyEight} & Gaps20522060 & PROVED & \checkmark & verified \\
\texttt{evenPair\_card\_twoThousandFiftyFour} & Gaps20522060 & PROVED & \checkmark & verified \\
\texttt{evenPair\_card\_twoThousandFiftySix} & Gaps20522060 & PROVED & \checkmark & verified \\
\texttt{evenPair\_card\_twoThousandFiftyTwo} & Gaps20522060 & PROVED & \checkmark & verified \\
\texttt{evenPair\_card\_twoThousandSixty} & Gaps20522060 & PROVED & \checkmark & verified \\
\texttt{isAdmissible\_evenPair\_twoThousandFiftyEight} & Gaps20522060 & PROVED & \checkmark & verified \\
\texttt{isAdmissible\_evenPair\_twoThousandFiftyFour} & Gaps20522060 & PROVED & \checkmark & verified \\
\texttt{isAdmissible\_evenPair\_twoThousandFiftySix} & Gaps20522060 & PROVED & \checkmark & verified \\
\texttt{isAdmissible\_evenPair\_twoThousandFiftyTwo} & Gaps20522060 & PROVED & \checkmark & verified \\
\texttt{isAdmissible\_evenPair\_twoThousandSixty} & Gaps20522060 & PROVED & \checkmark & verified \\
\texttt{localFactor\_twoThousandFiftyTwo\_two} & Gaps20522060 & PROVED & \checkmark & verified \\
\texttt{localFactor\_twoThousandSixty\_two} & Gaps20522060 & PROVED & \checkmark & verified \\
\texttt{nu\_p\_twoThousandFiftyEight} & Gaps20522060 & PROVED & \checkmark & verified \\
\texttt{nu\_p\_twoThousandFiftyFour} & Gaps20522060 & PROVED & \checkmark & verified \\
\texttt{nu\_p\_twoThousandFiftySix} & Gaps20522060 & PROVED & \checkmark & verified \\
\texttt{nu\_p\_twoThousandFiftyTwo} & Gaps20522060 & PROVED & \checkmark & verified \\
\texttt{nu\_p\_twoThousandFiftyTwo\_two} & Gaps20522060 & PROVED & \checkmark & verified \\
\texttt{nu\_p\_twoThousandSixty} & Gaps20522060 & PROVED & \checkmark & verified \\
\texttt{nu\_p\_twoThousandSixty\_two} & Gaps20522060 & PROVED & \checkmark & verified \\
\texttt{singular\_series\_finite\_pos\_evenPair\_twoThousandFiftyEight} & Gaps20522060 & PROVED & \checkmark & verified \\
\texttt{singular\_series\_finite\_pos\_evenPair\_twoThousandFiftyFour} & Gaps20522060 & PROVED & \checkmark & verified \\
\texttt{singular\_series\_finite\_pos\_evenPair\_twoThousandFiftySix} & Gaps20522060 & PROVED & \checkmark & verified \\
\texttt{singular\_series\_finite\_pos\_evenPair\_twoThousandFiftyTwo} & Gaps20522060 & PROVED & \checkmark & verified \\
\texttt{singular\_series\_finite\_pos\_evenPair\_twoThousandSixty} & Gaps20522060 & PROVED & \checkmark & verified \\
\texttt{singular\_series\_pos\_evenPair\_twoThousandFiftyEight} & Gaps20522060 & PROVED & \checkmark & verified \\
\texttt{singular\_series\_pos\_evenPair\_twoThousandFiftyFour} & Gaps20522060 & PROVED & \checkmark & verified \\
\texttt{singular\_series\_pos\_evenPair\_twoThousandFiftySix} & Gaps20522060 & PROVED & \checkmark & verified \\
\texttt{singular\_series\_pos\_evenPair\_twoThousandFiftyTwo} & Gaps20522060 & PROVED & \checkmark & verified \\
\texttt{singular\_series\_pos\_evenPair\_twoThousandSixty} & Gaps20522060 & PROVED & \checkmark & verified \\
\texttt{evenPair\_card\_twoThousandSeventy} & Gaps20622070 & PROVED & \checkmark & verified \\
\texttt{evenPair\_card\_twoThousandSixtyEight} & Gaps20622070 & PROVED & \checkmark & verified \\
\texttt{evenPair\_card\_twoThousandSixtyFour} & Gaps20622070 & PROVED & \checkmark & verified \\
\texttt{evenPair\_card\_twoThousandSixtySix} & Gaps20622070 & PROVED & \checkmark & verified \\
\texttt{evenPair\_card\_twoThousandSixtyTwo} & Gaps20622070 & PROVED & \checkmark & verified \\
\texttt{isAdmissible\_evenPair\_twoThousandSeventy} & Gaps20622070 & PROVED & \checkmark & verified \\
\texttt{isAdmissible\_evenPair\_twoThousandSixtyEight} & Gaps20622070 & PROVED & \checkmark & verified \\
\texttt{isAdmissible\_evenPair\_twoThousandSixtyFour} & Gaps20622070 & PROVED & \checkmark & verified \\
\texttt{isAdmissible\_evenPair\_twoThousandSixtySix} & Gaps20622070 & PROVED & \checkmark & verified \\
\texttt{isAdmissible\_evenPair\_twoThousandSixtyTwo} & Gaps20622070 & PROVED & \checkmark & verified \\
\texttt{localFactor\_twoThousandSeventy\_two} & Gaps20622070 & PROVED & \checkmark & verified \\
\texttt{localFactor\_twoThousandSixtyTwo\_two} & Gaps20622070 & PROVED & \checkmark & verified \\
\texttt{nu\_p\_twoThousandSeventy} & Gaps20622070 & PROVED & \checkmark & verified \\
\texttt{nu\_p\_twoThousandSeventy\_two} & Gaps20622070 & PROVED & \checkmark & verified \\
\texttt{nu\_p\_twoThousandSixtyEight} & Gaps20622070 & PROVED & \checkmark & verified \\
\texttt{nu\_p\_twoThousandSixtyFour} & Gaps20622070 & PROVED & \checkmark & verified \\
\texttt{nu\_p\_twoThousandSixtySix} & Gaps20622070 & PROVED & \checkmark & verified \\
\texttt{nu\_p\_twoThousandSixtyTwo} & Gaps20622070 & PROVED & \checkmark & verified \\
\texttt{nu\_p\_twoThousandSixtyTwo\_two} & Gaps20622070 & PROVED & \checkmark & verified \\
\texttt{singular\_series\_finite\_pos\_evenPair\_twoThousandSeventy} & Gaps20622070 & PROVED & \checkmark & verified \\
\texttt{singular\_series\_finite\_pos\_evenPair\_twoThousandSixtyEight} & Gaps20622070 & PROVED & \checkmark & verified \\
\texttt{singular\_series\_finite\_pos\_evenPair\_twoThousandSixtyFour} & Gaps20622070 & PROVED & \checkmark & verified \\
\texttt{singular\_series\_finite\_pos\_evenPair\_twoThousandSixtySix} & Gaps20622070 & PROVED & \checkmark & verified \\
\texttt{singular\_series\_finite\_pos\_evenPair\_twoThousandSixtyTwo} & Gaps20622070 & PROVED & \checkmark & verified \\
\texttt{singular\_series\_pos\_evenPair\_twoThousandSeventy} & Gaps20622070 & PROVED & \checkmark & verified \\
\texttt{singular\_series\_pos\_evenPair\_twoThousandSixtyEight} & Gaps20622070 & PROVED & \checkmark & verified \\
\texttt{singular\_series\_pos\_evenPair\_twoThousandSixtyFour} & Gaps20622070 & PROVED & \checkmark & verified \\
\texttt{singular\_series\_pos\_evenPair\_twoThousandSixtySix} & Gaps20622070 & PROVED & \checkmark & verified \\
\texttt{singular\_series\_pos\_evenPair\_twoThousandSixtyTwo} & Gaps20622070 & PROVED & \checkmark & verified \\
\texttt{evenPair\_card\_twoThousandEighty} & Gaps20722080 & PROVED & \checkmark & verified \\
\texttt{evenPair\_card\_twoThousandSeventyEight} & Gaps20722080 & PROVED & \checkmark & verified \\
\texttt{evenPair\_card\_twoThousandSeventyFour} & Gaps20722080 & PROVED & \checkmark & verified \\
\texttt{evenPair\_card\_twoThousandSeventySix} & Gaps20722080 & PROVED & \checkmark & verified \\
\texttt{evenPair\_card\_twoThousandSeventyTwo} & Gaps20722080 & PROVED & \checkmark & verified \\
\texttt{isAdmissible\_evenPair\_twoThousandEighty} & Gaps20722080 & PROVED & \checkmark & verified \\
\texttt{isAdmissible\_evenPair\_twoThousandSeventyEight} & Gaps20722080 & PROVED & \checkmark & verified \\
\texttt{isAdmissible\_evenPair\_twoThousandSeventyFour} & Gaps20722080 & PROVED & \checkmark & verified \\
\texttt{isAdmissible\_evenPair\_twoThousandSeventySix} & Gaps20722080 & PROVED & \checkmark & verified \\
\texttt{isAdmissible\_evenPair\_twoThousandSeventyTwo} & Gaps20722080 & PROVED & \checkmark & verified \\
\texttt{localFactor\_twoThousandEighty\_two} & Gaps20722080 & PROVED & \checkmark & verified \\
\texttt{localFactor\_twoThousandSeventyTwo\_two} & Gaps20722080 & PROVED & \checkmark & verified \\
\texttt{nu\_p\_twoThousandEighty} & Gaps20722080 & PROVED & \checkmark & verified \\
\texttt{nu\_p\_twoThousandEighty\_two} & Gaps20722080 & PROVED & \checkmark & verified \\
\texttt{nu\_p\_twoThousandSeventyEight} & Gaps20722080 & PROVED & \checkmark & verified \\
\texttt{nu\_p\_twoThousandSeventyFour} & Gaps20722080 & PROVED & \checkmark & verified \\
\texttt{nu\_p\_twoThousandSeventySix} & Gaps20722080 & PROVED & \checkmark & verified \\
\texttt{nu\_p\_twoThousandSeventyTwo} & Gaps20722080 & PROVED & \checkmark & verified \\
\texttt{nu\_p\_twoThousandSeventyTwo\_two} & Gaps20722080 & PROVED & \checkmark & verified \\
\texttt{singular\_series\_finite\_pos\_evenPair\_twoThousandEighty} & Gaps20722080 & PROVED & \checkmark & verified \\
\texttt{singular\_series\_finite\_pos\_evenPair\_twoThousandSeventyEight} & Gaps20722080 & PROVED & \checkmark & verified \\
\texttt{singular\_series\_finite\_pos\_evenPair\_twoThousandSeventyFour} & Gaps20722080 & PROVED & \checkmark & verified \\
\texttt{singular\_series\_finite\_pos\_evenPair\_twoThousandSeventySix} & Gaps20722080 & PROVED & \checkmark & verified \\
\texttt{singular\_series\_finite\_pos\_evenPair\_twoThousandSeventyTwo} & Gaps20722080 & PROVED & \checkmark & verified \\
\texttt{singular\_series\_pos\_evenPair\_twoThousandEighty} & Gaps20722080 & PROVED & \checkmark & verified \\
\texttt{singular\_series\_pos\_evenPair\_twoThousandSeventyEight} & Gaps20722080 & PROVED & \checkmark & verified \\
\texttt{singular\_series\_pos\_evenPair\_twoThousandSeventyFour} & Gaps20722080 & PROVED & \checkmark & verified \\
\texttt{singular\_series\_pos\_evenPair\_twoThousandSeventySix} & Gaps20722080 & PROVED & \checkmark & verified \\
\texttt{singular\_series\_pos\_evenPair\_twoThousandSeventyTwo} & Gaps20722080 & PROVED & \checkmark & verified \\
\texttt{evenPair\_card\_twoThousandEightyEight} & Gaps20822090 & PROVED & \checkmark & verified \\
\texttt{evenPair\_card\_twoThousandEightyFour} & Gaps20822090 & PROVED & \checkmark & verified \\
\texttt{evenPair\_card\_twoThousandEightySix} & Gaps20822090 & PROVED & \checkmark & verified \\
\texttt{evenPair\_card\_twoThousandEightyTwo} & Gaps20822090 & PROVED & \checkmark & verified \\
\texttt{evenPair\_card\_twoThousandNinety} & Gaps20822090 & PROVED & \checkmark & verified \\
\texttt{isAdmissible\_evenPair\_twoThousandEightyEight} & Gaps20822090 & PROVED & \checkmark & verified \\
\texttt{isAdmissible\_evenPair\_twoThousandEightyFour} & Gaps20822090 & PROVED & \checkmark & verified \\
\texttt{isAdmissible\_evenPair\_twoThousandEightySix} & Gaps20822090 & PROVED & \checkmark & verified \\
\texttt{isAdmissible\_evenPair\_twoThousandEightyTwo} & Gaps20822090 & PROVED & \checkmark & verified \\
\texttt{isAdmissible\_evenPair\_twoThousandNinety} & Gaps20822090 & PROVED & \checkmark & verified \\
\texttt{localFactor\_twoThousandEightyTwo\_two} & Gaps20822090 & PROVED & \checkmark & verified \\
\texttt{localFactor\_twoThousandNinety\_two} & Gaps20822090 & PROVED & \checkmark & verified \\
\texttt{nu\_p\_twoThousandEightyEight} & Gaps20822090 & PROVED & \checkmark & verified \\
\texttt{nu\_p\_twoThousandEightyFour} & Gaps20822090 & PROVED & \checkmark & verified \\
\texttt{nu\_p\_twoThousandEightySix} & Gaps20822090 & PROVED & \checkmark & verified \\
\texttt{nu\_p\_twoThousandEightyTwo} & Gaps20822090 & PROVED & \checkmark & verified \\
\texttt{nu\_p\_twoThousandEightyTwo\_two} & Gaps20822090 & PROVED & \checkmark & verified \\
\texttt{nu\_p\_twoThousandNinety} & Gaps20822090 & PROVED & \checkmark & verified \\
\texttt{nu\_p\_twoThousandNinety\_two} & Gaps20822090 & PROVED & \checkmark & verified \\
\texttt{singular\_series\_finite\_pos\_evenPair\_twoThousandEightyEight} & Gaps20822090 & PROVED & \checkmark & verified \\
\texttt{singular\_series\_finite\_pos\_evenPair\_twoThousandEightyFour} & Gaps20822090 & PROVED & \checkmark & verified \\
\texttt{singular\_series\_finite\_pos\_evenPair\_twoThousandEightySix} & Gaps20822090 & PROVED & \checkmark & verified \\
\texttt{singular\_series\_finite\_pos\_evenPair\_twoThousandEightyTwo} & Gaps20822090 & PROVED & \checkmark & verified \\
\texttt{singular\_series\_finite\_pos\_evenPair\_twoThousandNinety} & Gaps20822090 & PROVED & \checkmark & verified \\
\texttt{singular\_series\_pos\_evenPair\_twoThousandEightyEight} & Gaps20822090 & PROVED & \checkmark & verified \\
\texttt{singular\_series\_pos\_evenPair\_twoThousandEightyFour} & Gaps20822090 & PROVED & \checkmark & verified \\
\texttt{singular\_series\_pos\_evenPair\_twoThousandEightySix} & Gaps20822090 & PROVED & \checkmark & verified \\
\texttt{singular\_series\_pos\_evenPair\_twoThousandEightyTwo} & Gaps20822090 & PROVED & \checkmark & verified \\
\texttt{singular\_series\_pos\_evenPair\_twoThousandNinety} & Gaps20822090 & PROVED & \checkmark & verified \\
\texttt{evenPair\_card\_twoThousandNinetyEight} & Gaps20922100 & PROVED & \checkmark & verified \\
\texttt{evenPair\_card\_twoThousandNinetyFour} & Gaps20922100 & PROVED & \checkmark & verified \\
\texttt{evenPair\_card\_twoThousandNinetySix} & Gaps20922100 & PROVED & \checkmark & verified \\
\texttt{evenPair\_card\_twoThousandNinetyTwo} & Gaps20922100 & PROVED & \checkmark & verified \\
\texttt{evenPair\_card\_twoThousandOneHundred} & Gaps20922100 & PROVED & \checkmark & verified \\
\texttt{isAdmissible\_evenPair\_twoThousandNinetyEight} & Gaps20922100 & PROVED & \checkmark & verified \\
\texttt{isAdmissible\_evenPair\_twoThousandNinetyFour} & Gaps20922100 & PROVED & \checkmark & verified \\
\texttt{isAdmissible\_evenPair\_twoThousandNinetySix} & Gaps20922100 & PROVED & \checkmark & verified \\
\texttt{isAdmissible\_evenPair\_twoThousandNinetyTwo} & Gaps20922100 & PROVED & \checkmark & verified \\
\texttt{isAdmissible\_evenPair\_twoThousandOneHundred} & Gaps20922100 & PROVED & \checkmark & verified \\
\texttt{localFactor\_twoThousandNinetyTwo\_two} & Gaps20922100 & PROVED & \checkmark & verified \\
\texttt{localFactor\_twoThousandOneHundred\_two} & Gaps20922100 & PROVED & \checkmark & verified \\
\texttt{nu\_p\_twoThousandNinetyEight} & Gaps20922100 & PROVED & \checkmark & verified \\
\texttt{nu\_p\_twoThousandNinetyFour} & Gaps20922100 & PROVED & \checkmark & verified \\
\texttt{nu\_p\_twoThousandNinetySix} & Gaps20922100 & PROVED & \checkmark & verified \\
\texttt{nu\_p\_twoThousandNinetyTwo} & Gaps20922100 & PROVED & \checkmark & verified \\
\texttt{nu\_p\_twoThousandNinetyTwo\_two} & Gaps20922100 & PROVED & \checkmark & verified \\
\texttt{nu\_p\_twoThousandOneHundred} & Gaps20922100 & PROVED & \checkmark & verified \\
\texttt{nu\_p\_twoThousandOneHundred\_two} & Gaps20922100 & PROVED & \checkmark & verified \\
\texttt{singular\_series\_finite\_pos\_evenPair\_twoThousandNinetyEight} & Gaps20922100 & PROVED & \checkmark & verified \\
\texttt{singular\_series\_finite\_pos\_evenPair\_twoThousandNinetyFour} & Gaps20922100 & PROVED & \checkmark & verified \\
\texttt{singular\_series\_finite\_pos\_evenPair\_twoThousandNinetySix} & Gaps20922100 & PROVED & \checkmark & verified \\
\texttt{singular\_series\_finite\_pos\_evenPair\_twoThousandNinetyTwo} & Gaps20922100 & PROVED & \checkmark & verified \\
\texttt{singular\_series\_finite\_pos\_evenPair\_twoThousandOneHundred} & Gaps20922100 & PROVED & \checkmark & verified \\
\texttt{singular\_series\_pos\_evenPair\_twoThousandNinetyEight} & Gaps20922100 & PROVED & \checkmark & verified \\
\texttt{singular\_series\_pos\_evenPair\_twoThousandNinetyFour} & Gaps20922100 & PROVED & \checkmark & verified \\
\texttt{singular\_series\_pos\_evenPair\_twoThousandNinetySix} & Gaps20922100 & PROVED & \checkmark & verified \\
\texttt{singular\_series\_pos\_evenPair\_twoThousandNinetyTwo} & Gaps20922100 & PROVED & \checkmark & verified \\
\texttt{singular\_series\_pos\_evenPair\_twoThousandOneHundred} & Gaps20922100 & PROVED & \checkmark & verified \\
\texttt{evenPair\_card\_twoThousandOneHundredEight} & Gaps21022110 & PROVED & \checkmark & verified \\
\texttt{evenPair\_card\_twoThousandOneHundredFour} & Gaps21022110 & PROVED & \checkmark & verified \\
\texttt{evenPair\_card\_twoThousandOneHundredSix} & Gaps21022110 & PROVED & \checkmark & verified \\
\texttt{evenPair\_card\_twoThousandOneHundredTen} & Gaps21022110 & PROVED & \checkmark & verified \\
\texttt{evenPair\_card\_twoThousandOneHundredTwo} & Gaps21022110 & PROVED & \checkmark & verified \\
\texttt{isAdmissible\_evenPair\_twoThousandOneHundredEight} & Gaps21022110 & PROVED & \checkmark & verified \\
\texttt{isAdmissible\_evenPair\_twoThousandOneHundredFour} & Gaps21022110 & PROVED & \checkmark & verified \\
\texttt{isAdmissible\_evenPair\_twoThousandOneHundredSix} & Gaps21022110 & PROVED & \checkmark & verified \\
\texttt{isAdmissible\_evenPair\_twoThousandOneHundredTen} & Gaps21022110 & PROVED & \checkmark & verified \\
\texttt{isAdmissible\_evenPair\_twoThousandOneHundredTwo} & Gaps21022110 & PROVED & \checkmark & verified \\
\texttt{localFactor\_twoThousandOneHundredTen\_two} & Gaps21022110 & PROVED & \checkmark & verified \\
\texttt{localFactor\_twoThousandOneHundredTwo\_two} & Gaps21022110 & PROVED & \checkmark & verified \\
\texttt{nu\_p\_twoThousandOneHundredEight} & Gaps21022110 & PROVED & \checkmark & verified \\
\texttt{nu\_p\_twoThousandOneHundredFour} & Gaps21022110 & PROVED & \checkmark & verified \\
\texttt{nu\_p\_twoThousandOneHundredSix} & Gaps21022110 & PROVED & \checkmark & verified \\
\texttt{nu\_p\_twoThousandOneHundredTen} & Gaps21022110 & PROVED & \checkmark & verified \\
\texttt{nu\_p\_twoThousandOneHundredTen\_two} & Gaps21022110 & PROVED & \checkmark & verified \\
\texttt{nu\_p\_twoThousandOneHundredTwo} & Gaps21022110 & PROVED & \checkmark & verified \\
\texttt{nu\_p\_twoThousandOneHundredTwo\_two} & Gaps21022110 & PROVED & \checkmark & verified \\
\texttt{singular\_series\_finite\_pos\_evenPair\_twoThousandOneHundredEight} & Gaps21022110 & PROVED & \checkmark & verified \\
\texttt{singular\_series\_finite\_pos\_evenPair\_twoThousandOneHundredFour} & Gaps21022110 & PROVED & \checkmark & verified \\
\texttt{singular\_series\_finite\_pos\_evenPair\_twoThousandOneHundredSix} & Gaps21022110 & PROVED & \checkmark & verified \\
\texttt{singular\_series\_finite\_pos\_evenPair\_twoThousandOneHundredTen} & Gaps21022110 & PROVED & \checkmark & verified \\
\texttt{singular\_series\_finite\_pos\_evenPair\_twoThousandOneHundredTwo} & Gaps21022110 & PROVED & \checkmark & verified \\
\texttt{singular\_series\_pos\_evenPair\_twoThousandOneHundredEight} & Gaps21022110 & PROVED & \checkmark & verified \\
\texttt{singular\_series\_pos\_evenPair\_twoThousandOneHundredFour} & Gaps21022110 & PROVED & \checkmark & verified \\
\texttt{singular\_series\_pos\_evenPair\_twoThousandOneHundredSix} & Gaps21022110 & PROVED & \checkmark & verified \\
\texttt{singular\_series\_pos\_evenPair\_twoThousandOneHundredTen} & Gaps21022110 & PROVED & \checkmark & verified \\
\texttt{singular\_series\_pos\_evenPair\_twoThousandOneHundredTwo} & Gaps21022110 & PROVED & \checkmark & verified \\
\texttt{evenPair\_card\_twoThousandOneHundredEighteen} & Gaps21122120 & PROVED & \checkmark & verified \\
\texttt{evenPair\_card\_twoThousandOneHundredFourteen} & Gaps21122120 & PROVED & \checkmark & verified \\
\texttt{evenPair\_card\_twoThousandOneHundredSixteen} & Gaps21122120 & PROVED & \checkmark & verified \\
\texttt{evenPair\_card\_twoThousandOneHundredTwelve} & Gaps21122120 & PROVED & \checkmark & verified \\
\texttt{evenPair\_card\_twoThousandOneHundredTwenty} & Gaps21122120 & PROVED & \checkmark & verified \\
\texttt{isAdmissible\_evenPair\_twoThousandOneHundredEighteen} & Gaps21122120 & PROVED & \checkmark & verified \\
\texttt{isAdmissible\_evenPair\_twoThousandOneHundredFourteen} & Gaps21122120 & PROVED & \checkmark & verified \\
\texttt{isAdmissible\_evenPair\_twoThousandOneHundredSixteen} & Gaps21122120 & PROVED & \checkmark & verified \\
\texttt{isAdmissible\_evenPair\_twoThousandOneHundredTwelve} & Gaps21122120 & PROVED & \checkmark & verified \\
\texttt{isAdmissible\_evenPair\_twoThousandOneHundredTwenty} & Gaps21122120 & PROVED & \checkmark & verified \\
\texttt{localFactor\_twoThousandOneHundredTwelve\_two} & Gaps21122120 & PROVED & \checkmark & verified \\
\texttt{localFactor\_twoThousandOneHundredTwenty\_two} & Gaps21122120 & PROVED & \checkmark & verified \\
\texttt{nu\_p\_twoThousandOneHundredEighteen} & Gaps21122120 & PROVED & \checkmark & verified \\
\texttt{nu\_p\_twoThousandOneHundredFourteen} & Gaps21122120 & PROVED & \checkmark & verified \\
\texttt{nu\_p\_twoThousandOneHundredSixteen} & Gaps21122120 & PROVED & \checkmark & verified \\
\texttt{nu\_p\_twoThousandOneHundredTwelve} & Gaps21122120 & PROVED & \checkmark & verified \\
\texttt{nu\_p\_twoThousandOneHundredTwelve\_two} & Gaps21122120 & PROVED & \checkmark & verified \\
\texttt{nu\_p\_twoThousandOneHundredTwenty} & Gaps21122120 & PROVED & \checkmark & verified \\
\texttt{nu\_p\_twoThousandOneHundredTwenty\_two} & Gaps21122120 & PROVED & \checkmark & verified \\
\texttt{singular\_series\_finite\_pos\_evenPair\_twoThousandOneHundredEighteen} & Gaps21122120 & PROVED & \checkmark & verified \\
\texttt{singular\_series\_finite\_pos\_evenPair\_twoThousandOneHundredFourteen} & Gaps21122120 & PROVED & \checkmark & verified \\
\texttt{singular\_series\_finite\_pos\_evenPair\_twoThousandOneHundredSixteen} & Gaps21122120 & PROVED & \checkmark & verified \\
\texttt{singular\_series\_finite\_pos\_evenPair\_twoThousandOneHundredTwelve} & Gaps21122120 & PROVED & \checkmark & verified \\
\texttt{singular\_series\_finite\_pos\_evenPair\_twoThousandOneHundredTwenty} & Gaps21122120 & PROVED & \checkmark & verified \\
\texttt{singular\_series\_pos\_evenPair\_twoThousandOneHundredEighteen} & Gaps21122120 & PROVED & \checkmark & verified \\
\texttt{singular\_series\_pos\_evenPair\_twoThousandOneHundredFourteen} & Gaps21122120 & PROVED & \checkmark & verified \\
\texttt{singular\_series\_pos\_evenPair\_twoThousandOneHundredSixteen} & Gaps21122120 & PROVED & \checkmark & verified \\
\texttt{singular\_series\_pos\_evenPair\_twoThousandOneHundredTwelve} & Gaps21122120 & PROVED & \checkmark & verified \\
\texttt{singular\_series\_pos\_evenPair\_twoThousandOneHundredTwenty} & Gaps21122120 & PROVED & \checkmark & verified \\
\texttt{evenPair\_card\_twoHundredEighteen} & Gaps212220 & PROVED & \checkmark & verified \\
\texttt{evenPair\_card\_twoHundredFourteen} & Gaps212220 & PROVED & \checkmark & verified \\
\texttt{evenPair\_card\_twoHundredSixteen} & Gaps212220 & PROVED & \checkmark & verified \\
\texttt{evenPair\_card\_twoHundredTwelve} & Gaps212220 & PROVED & \checkmark & verified \\
\texttt{evenPair\_card\_twoHundredTwenty} & Gaps212220 & PROVED & \checkmark & verified \\
\texttt{isAdmissible\_evenPair\_twoHundredEighteen} & Gaps212220 & PROVED & \checkmark & verified \\
\texttt{isAdmissible\_evenPair\_twoHundredFourteen} & Gaps212220 & PROVED & \checkmark & verified \\
\texttt{isAdmissible\_evenPair\_twoHundredSixteen} & Gaps212220 & PROVED & \checkmark & verified \\
\texttt{isAdmissible\_evenPair\_twoHundredTwelve} & Gaps212220 & PROVED & \checkmark & verified \\
\texttt{isAdmissible\_evenPair\_twoHundredTwenty} & Gaps212220 & PROVED & \checkmark & verified \\
\texttt{localFactor\_twoHundredTwelve\_two} & Gaps212220 & PROVED & \checkmark & verified \\
\texttt{localFactor\_twoHundredTwenty\_two} & Gaps212220 & PROVED & \checkmark & verified \\
\texttt{nu\_p\_twoHundredEighteen} & Gaps212220 & PROVED & \checkmark & verified \\
\texttt{nu\_p\_twoHundredFourteen} & Gaps212220 & PROVED & \checkmark & verified \\
\texttt{nu\_p\_twoHundredSixteen} & Gaps212220 & PROVED & \checkmark & verified \\
\texttt{nu\_p\_twoHundredTwelve} & Gaps212220 & PROVED & \checkmark & verified \\
\texttt{nu\_p\_twoHundredTwelve\_two} & Gaps212220 & PROVED & \checkmark & verified \\
\texttt{nu\_p\_twoHundredTwenty} & Gaps212220 & PROVED & \checkmark & verified \\
\texttt{nu\_p\_twoHundredTwenty\_two} & Gaps212220 & PROVED & \checkmark & verified \\
\texttt{singular\_series\_finite\_pos\_evenPair\_twoHundredEighteen} & Gaps212220 & PROVED & \checkmark & verified \\
\texttt{singular\_series\_finite\_pos\_evenPair\_twoHundredFourteen} & Gaps212220 & PROVED & \checkmark & verified \\
\texttt{singular\_series\_finite\_pos\_evenPair\_twoHundredSixteen} & Gaps212220 & PROVED & \checkmark & verified \\
\texttt{singular\_series\_finite\_pos\_evenPair\_twoHundredTwelve} & Gaps212220 & PROVED & \checkmark & verified \\
\texttt{singular\_series\_finite\_pos\_evenPair\_twoHundredTwenty} & Gaps212220 & PROVED & \checkmark & verified \\
\texttt{singular\_series\_pos\_evenPair\_twoHundredEighteen} & Gaps212220 & PROVED & \checkmark & verified \\
\texttt{singular\_series\_pos\_evenPair\_twoHundredFourteen} & Gaps212220 & PROVED & \checkmark & verified \\
\texttt{singular\_series\_pos\_evenPair\_twoHundredSixteen} & Gaps212220 & PROVED & \checkmark & verified \\
\texttt{singular\_series\_pos\_evenPair\_twoHundredTwelve} & Gaps212220 & PROVED & \checkmark & verified \\
\texttt{singular\_series\_pos\_evenPair\_twoHundredTwenty} & Gaps212220 & PROVED & \checkmark & verified \\
\texttt{evenPair\_card\_twoThousandOneHundredThirty} & Gaps21222130 & PROVED & \checkmark & verified \\
\texttt{evenPair\_card\_twoThousandOneHundredTwentyEight} & Gaps21222130 & PROVED & \checkmark & verified \\
\texttt{evenPair\_card\_twoThousandOneHundredTwentyFour} & Gaps21222130 & PROVED & \checkmark & verified \\
\texttt{evenPair\_card\_twoThousandOneHundredTwentySix} & Gaps21222130 & PROVED & \checkmark & verified \\
\texttt{evenPair\_card\_twoThousandOneHundredTwentyTwo} & Gaps21222130 & PROVED & \checkmark & verified \\
\texttt{isAdmissible\_evenPair\_twoThousandOneHundredThirty} & Gaps21222130 & PROVED & \checkmark & verified \\
\texttt{isAdmissible\_evenPair\_twoThousandOneHundredTwentyEight} & Gaps21222130 & PROVED & \checkmark & verified \\
\texttt{isAdmissible\_evenPair\_twoThousandOneHundredTwentyFour} & Gaps21222130 & PROVED & \checkmark & verified \\
\texttt{isAdmissible\_evenPair\_twoThousandOneHundredTwentySix} & Gaps21222130 & PROVED & \checkmark & verified \\
\texttt{isAdmissible\_evenPair\_twoThousandOneHundredTwentyTwo} & Gaps21222130 & PROVED & \checkmark & verified \\
\texttt{localFactor\_twoThousandOneHundredThirty\_two} & Gaps21222130 & PROVED & \checkmark & verified \\
\texttt{localFactor\_twoThousandOneHundredTwentyTwo\_two} & Gaps21222130 & PROVED & \checkmark & verified \\
\texttt{nu\_p\_twoThousandOneHundredThirty} & Gaps21222130 & PROVED & \checkmark & verified \\
\texttt{nu\_p\_twoThousandOneHundredThirty\_two} & Gaps21222130 & PROVED & \checkmark & verified \\
\texttt{nu\_p\_twoThousandOneHundredTwentyEight} & Gaps21222130 & PROVED & \checkmark & verified \\
\texttt{nu\_p\_twoThousandOneHundredTwentyFour} & Gaps21222130 & PROVED & \checkmark & verified \\
\texttt{nu\_p\_twoThousandOneHundredTwentySix} & Gaps21222130 & PROVED & \checkmark & verified \\
\texttt{nu\_p\_twoThousandOneHundredTwentyTwo} & Gaps21222130 & PROVED & \checkmark & verified \\
\texttt{nu\_p\_twoThousandOneHundredTwentyTwo\_two} & Gaps21222130 & PROVED & \checkmark & verified \\
\texttt{singular\_series\_finite\_pos\_evenPair\_twoThousandOneHundredThirty} & Gaps21222130 & PROVED & \checkmark & verified \\
\texttt{singular\_series\_finite\_pos\_evenPair\_twoThousandOneHundredTwentyEight} & Gaps21222130 & PROVED & \checkmark & verified \\
\texttt{singular\_series\_finite\_pos\_evenPair\_twoThousandOneHundredTwentyFour} & Gaps21222130 & PROVED & \checkmark & verified \\
\texttt{singular\_series\_finite\_pos\_evenPair\_twoThousandOneHundredTwentySix} & Gaps21222130 & PROVED & \checkmark & verified \\
\texttt{singular\_series\_finite\_pos\_evenPair\_twoThousandOneHundredTwentyTwo} & Gaps21222130 & PROVED & \checkmark & verified \\
\texttt{singular\_series\_pos\_evenPair\_twoThousandOneHundredThirty} & Gaps21222130 & PROVED & \checkmark & verified \\
\texttt{singular\_series\_pos\_evenPair\_twoThousandOneHundredTwentyEight} & Gaps21222130 & PROVED & \checkmark & verified \\
\texttt{singular\_series\_pos\_evenPair\_twoThousandOneHundredTwentyFour} & Gaps21222130 & PROVED & \checkmark & verified \\
\texttt{singular\_series\_pos\_evenPair\_twoThousandOneHundredTwentySix} & Gaps21222130 & PROVED & \checkmark & verified \\
\texttt{singular\_series\_pos\_evenPair\_twoThousandOneHundredTwentyTwo} & Gaps21222130 & PROVED & \checkmark & verified \\
\texttt{evenPair\_card\_twoThousandOneHundredForty} & Gaps21322140 & PROVED & \checkmark & verified \\
\texttt{evenPair\_card\_twoThousandOneHundredThirtyEight} & Gaps21322140 & PROVED & \checkmark & verified \\
\texttt{evenPair\_card\_twoThousandOneHundredThirtyFour} & Gaps21322140 & PROVED & \checkmark & verified \\
\texttt{evenPair\_card\_twoThousandOneHundredThirtySix} & Gaps21322140 & PROVED & \checkmark & verified \\
\texttt{evenPair\_card\_twoThousandOneHundredThirtyTwo} & Gaps21322140 & PROVED & \checkmark & verified \\
\texttt{isAdmissible\_evenPair\_twoThousandOneHundredForty} & Gaps21322140 & PROVED & \checkmark & verified \\
\texttt{isAdmissible\_evenPair\_twoThousandOneHundredThirtyEight} & Gaps21322140 & PROVED & \checkmark & verified \\
\texttt{isAdmissible\_evenPair\_twoThousandOneHundredThirtyFour} & Gaps21322140 & PROVED & \checkmark & verified \\
\texttt{isAdmissible\_evenPair\_twoThousandOneHundredThirtySix} & Gaps21322140 & PROVED & \checkmark & verified \\
\texttt{isAdmissible\_evenPair\_twoThousandOneHundredThirtyTwo} & Gaps21322140 & PROVED & \checkmark & verified \\
\texttt{localFactor\_twoThousandOneHundredForty\_two} & Gaps21322140 & PROVED & \checkmark & verified \\
\texttt{localFactor\_twoThousandOneHundredThirtyTwo\_two} & Gaps21322140 & PROVED & \checkmark & verified \\
\texttt{nu\_p\_twoThousandOneHundredForty} & Gaps21322140 & PROVED & \checkmark & verified \\
\texttt{nu\_p\_twoThousandOneHundredForty\_two} & Gaps21322140 & PROVED & \checkmark & verified \\
\texttt{nu\_p\_twoThousandOneHundredThirtyEight} & Gaps21322140 & PROVED & \checkmark & verified \\
\texttt{nu\_p\_twoThousandOneHundredThirtyFour} & Gaps21322140 & PROVED & \checkmark & verified \\
\texttt{nu\_p\_twoThousandOneHundredThirtySix} & Gaps21322140 & PROVED & \checkmark & verified \\
\texttt{nu\_p\_twoThousandOneHundredThirtyTwo} & Gaps21322140 & PROVED & \checkmark & verified \\
\texttt{nu\_p\_twoThousandOneHundredThirtyTwo\_two} & Gaps21322140 & PROVED & \checkmark & verified \\
\texttt{singular\_series\_finite\_pos\_evenPair\_twoThousandOneHundredForty} & Gaps21322140 & PROVED & \checkmark & verified \\
\texttt{singular\_series\_finite\_pos\_evenPair\_twoThousandOneHundredThirtyEight} & Gaps21322140 & PROVED & \checkmark & verified \\
\texttt{singular\_series\_finite\_pos\_evenPair\_twoThousandOneHundredThirtyFour} & Gaps21322140 & PROVED & \checkmark & verified \\
\texttt{singular\_series\_finite\_pos\_evenPair\_twoThousandOneHundredThirtySix} & Gaps21322140 & PROVED & \checkmark & verified \\
\texttt{singular\_series\_finite\_pos\_evenPair\_twoThousandOneHundredThirtyTwo} & Gaps21322140 & PROVED & \checkmark & verified \\
\texttt{singular\_series\_pos\_evenPair\_twoThousandOneHundredForty} & Gaps21322140 & PROVED & \checkmark & verified \\
\texttt{singular\_series\_pos\_evenPair\_twoThousandOneHundredThirtyEight} & Gaps21322140 & PROVED & \checkmark & verified \\
\texttt{singular\_series\_pos\_evenPair\_twoThousandOneHundredThirtyFour} & Gaps21322140 & PROVED & \checkmark & verified \\
\texttt{singular\_series\_pos\_evenPair\_twoThousandOneHundredThirtySix} & Gaps21322140 & PROVED & \checkmark & verified \\
\texttt{singular\_series\_pos\_evenPair\_twoThousandOneHundredThirtyTwo} & Gaps21322140 & PROVED & \checkmark & verified \\
\texttt{evenPair\_card\_twoThousandOneHundredFifty} & Gaps21422150 & PROVED & \checkmark & verified \\
\texttt{evenPair\_card\_twoThousandOneHundredFortyEight} & Gaps21422150 & PROVED & \checkmark & verified \\
\texttt{evenPair\_card\_twoThousandOneHundredFortyFour} & Gaps21422150 & PROVED & \checkmark & verified \\
\texttt{evenPair\_card\_twoThousandOneHundredFortySix} & Gaps21422150 & PROVED & \checkmark & verified \\
\texttt{evenPair\_card\_twoThousandOneHundredFortyTwo} & Gaps21422150 & PROVED & \checkmark & verified \\
\texttt{isAdmissible\_evenPair\_twoThousandOneHundredFifty} & Gaps21422150 & PROVED & \checkmark & verified \\
\texttt{isAdmissible\_evenPair\_twoThousandOneHundredFortyEight} & Gaps21422150 & PROVED & \checkmark & verified \\
\texttt{isAdmissible\_evenPair\_twoThousandOneHundredFortyFour} & Gaps21422150 & PROVED & \checkmark & verified \\
\texttt{isAdmissible\_evenPair\_twoThousandOneHundredFortySix} & Gaps21422150 & PROVED & \checkmark & verified \\
\texttt{isAdmissible\_evenPair\_twoThousandOneHundredFortyTwo} & Gaps21422150 & PROVED & \checkmark & verified \\
\texttt{localFactor\_twoThousandOneHundredFifty\_two} & Gaps21422150 & PROVED & \checkmark & verified \\
\texttt{localFactor\_twoThousandOneHundredFortyTwo\_two} & Gaps21422150 & PROVED & \checkmark & verified \\
\texttt{nu\_p\_twoThousandOneHundredFifty} & Gaps21422150 & PROVED & \checkmark & verified \\
\texttt{nu\_p\_twoThousandOneHundredFifty\_two} & Gaps21422150 & PROVED & \checkmark & verified \\
\texttt{nu\_p\_twoThousandOneHundredFortyEight} & Gaps21422150 & PROVED & \checkmark & verified \\
\texttt{nu\_p\_twoThousandOneHundredFortyFour} & Gaps21422150 & PROVED & \checkmark & verified \\
\texttt{nu\_p\_twoThousandOneHundredFortySix} & Gaps21422150 & PROVED & \checkmark & verified \\
\texttt{nu\_p\_twoThousandOneHundredFortyTwo} & Gaps21422150 & PROVED & \checkmark & verified \\
\texttt{nu\_p\_twoThousandOneHundredFortyTwo\_two} & Gaps21422150 & PROVED & \checkmark & verified \\
\texttt{singular\_series\_finite\_pos\_evenPair\_twoThousandOneHundredFifty} & Gaps21422150 & PROVED & \checkmark & verified \\
\texttt{singular\_series\_finite\_pos\_evenPair\_twoThousandOneHundredFortyEight} & Gaps21422150 & PROVED & \checkmark & verified \\
\texttt{singular\_series\_finite\_pos\_evenPair\_twoThousandOneHundredFortyFour} & Gaps21422150 & PROVED & \checkmark & verified \\
\texttt{singular\_series\_finite\_pos\_evenPair\_twoThousandOneHundredFortySix} & Gaps21422150 & PROVED & \checkmark & verified \\
\texttt{singular\_series\_finite\_pos\_evenPair\_twoThousandOneHundredFortyTwo} & Gaps21422150 & PROVED & \checkmark & verified \\
\texttt{singular\_series\_pos\_evenPair\_twoThousandOneHundredFifty} & Gaps21422150 & PROVED & \checkmark & verified \\
\texttt{singular\_series\_pos\_evenPair\_twoThousandOneHundredFortyEight} & Gaps21422150 & PROVED & \checkmark & verified \\
\texttt{singular\_series\_pos\_evenPair\_twoThousandOneHundredFortyFour} & Gaps21422150 & PROVED & \checkmark & verified \\
\texttt{singular\_series\_pos\_evenPair\_twoThousandOneHundredFortySix} & Gaps21422150 & PROVED & \checkmark & verified \\
\texttt{singular\_series\_pos\_evenPair\_twoThousandOneHundredFortyTwo} & Gaps21422150 & PROVED & \checkmark & verified \\
\texttt{evenPair\_card\_twoThousandOneHundredFiftyEight} & Gaps21522160 & PROVED & \checkmark & verified \\
\texttt{evenPair\_card\_twoThousandOneHundredFiftyFour} & Gaps21522160 & PROVED & \checkmark & verified \\
\texttt{evenPair\_card\_twoThousandOneHundredFiftySix} & Gaps21522160 & PROVED & \checkmark & verified \\
\texttt{evenPair\_card\_twoThousandOneHundredFiftyTwo} & Gaps21522160 & PROVED & \checkmark & verified \\
\texttt{evenPair\_card\_twoThousandOneHundredSixty} & Gaps21522160 & PROVED & \checkmark & verified \\
\texttt{isAdmissible\_evenPair\_twoThousandOneHundredFiftyEight} & Gaps21522160 & PROVED & \checkmark & verified \\
\texttt{isAdmissible\_evenPair\_twoThousandOneHundredFiftyFour} & Gaps21522160 & PROVED & \checkmark & verified \\
\texttt{isAdmissible\_evenPair\_twoThousandOneHundredFiftySix} & Gaps21522160 & PROVED & \checkmark & verified \\
\texttt{isAdmissible\_evenPair\_twoThousandOneHundredFiftyTwo} & Gaps21522160 & PROVED & \checkmark & verified \\
\texttt{isAdmissible\_evenPair\_twoThousandOneHundredSixty} & Gaps21522160 & PROVED & \checkmark & verified \\
\texttt{localFactor\_twoThousandOneHundredFiftyTwo\_two} & Gaps21522160 & PROVED & \checkmark & verified \\
\texttt{localFactor\_twoThousandOneHundredSixty\_two} & Gaps21522160 & PROVED & \checkmark & verified \\
\texttt{nu\_p\_twoThousandOneHundredFiftyEight} & Gaps21522160 & PROVED & \checkmark & verified \\
\texttt{nu\_p\_twoThousandOneHundredFiftyFour} & Gaps21522160 & PROVED & \checkmark & verified \\
\texttt{nu\_p\_twoThousandOneHundredFiftySix} & Gaps21522160 & PROVED & \checkmark & verified \\
\texttt{nu\_p\_twoThousandOneHundredFiftyTwo} & Gaps21522160 & PROVED & \checkmark & verified \\
\texttt{nu\_p\_twoThousandOneHundredFiftyTwo\_two} & Gaps21522160 & PROVED & \checkmark & verified \\
\texttt{nu\_p\_twoThousandOneHundredSixty} & Gaps21522160 & PROVED & \checkmark & verified \\
\texttt{nu\_p\_twoThousandOneHundredSixty\_two} & Gaps21522160 & PROVED & \checkmark & verified \\
\texttt{singular\_series\_finite\_pos\_evenPair\_twoThousandOneHundredFiftyEight} & Gaps21522160 & PROVED & \checkmark & verified \\
\texttt{singular\_series\_finite\_pos\_evenPair\_twoThousandOneHundredFiftyFour} & Gaps21522160 & PROVED & \checkmark & verified \\
\texttt{singular\_series\_finite\_pos\_evenPair\_twoThousandOneHundredFiftySix} & Gaps21522160 & PROVED & \checkmark & verified \\
\texttt{singular\_series\_finite\_pos\_evenPair\_twoThousandOneHundredFiftyTwo} & Gaps21522160 & PROVED & \checkmark & verified \\
\texttt{singular\_series\_finite\_pos\_evenPair\_twoThousandOneHundredSixty} & Gaps21522160 & PROVED & \checkmark & verified \\
\texttt{singular\_series\_pos\_evenPair\_twoThousandOneHundredFiftyEight} & Gaps21522160 & PROVED & \checkmark & verified \\
\texttt{singular\_series\_pos\_evenPair\_twoThousandOneHundredFiftyFour} & Gaps21522160 & PROVED & \checkmark & verified \\
\texttt{singular\_series\_pos\_evenPair\_twoThousandOneHundredFiftySix} & Gaps21522160 & PROVED & \checkmark & verified \\
\texttt{singular\_series\_pos\_evenPair\_twoThousandOneHundredFiftyTwo} & Gaps21522160 & PROVED & \checkmark & verified \\
\texttt{singular\_series\_pos\_evenPair\_twoThousandOneHundredSixty} & Gaps21522160 & PROVED & \checkmark & verified \\
\texttt{evenPair\_card\_twoThousandOneHundredSeventy} & Gaps21622170 & PROVED & \checkmark & verified \\
\texttt{evenPair\_card\_twoThousandOneHundredSixtyEight} & Gaps21622170 & PROVED & \checkmark & verified \\
\texttt{evenPair\_card\_twoThousandOneHundredSixtyFour} & Gaps21622170 & PROVED & \checkmark & verified \\
\texttt{evenPair\_card\_twoThousandOneHundredSixtySix} & Gaps21622170 & PROVED & \checkmark & verified \\
\texttt{evenPair\_card\_twoThousandOneHundredSixtyTwo} & Gaps21622170 & PROVED & \checkmark & verified \\
\texttt{isAdmissible\_evenPair\_twoThousandOneHundredSeventy} & Gaps21622170 & PROVED & \checkmark & verified \\
\texttt{isAdmissible\_evenPair\_twoThousandOneHundredSixtyEight} & Gaps21622170 & PROVED & \checkmark & verified \\
\texttt{isAdmissible\_evenPair\_twoThousandOneHundredSixtyFour} & Gaps21622170 & PROVED & \checkmark & verified \\
\texttt{isAdmissible\_evenPair\_twoThousandOneHundredSixtySix} & Gaps21622170 & PROVED & \checkmark & verified \\
\texttt{isAdmissible\_evenPair\_twoThousandOneHundredSixtyTwo} & Gaps21622170 & PROVED & \checkmark & verified \\
\texttt{localFactor\_twoThousandOneHundredSeventy\_two} & Gaps21622170 & PROVED & \checkmark & verified \\
\texttt{localFactor\_twoThousandOneHundredSixtyTwo\_two} & Gaps21622170 & PROVED & \checkmark & verified \\
\texttt{nu\_p\_twoThousandOneHundredSeventy} & Gaps21622170 & PROVED & \checkmark & verified \\
\texttt{nu\_p\_twoThousandOneHundredSeventy\_two} & Gaps21622170 & PROVED & \checkmark & verified \\
\texttt{nu\_p\_twoThousandOneHundredSixtyEight} & Gaps21622170 & PROVED & \checkmark & verified \\
\texttt{nu\_p\_twoThousandOneHundredSixtyFour} & Gaps21622170 & PROVED & \checkmark & verified \\
\texttt{nu\_p\_twoThousandOneHundredSixtySix} & Gaps21622170 & PROVED & \checkmark & verified \\
\texttt{nu\_p\_twoThousandOneHundredSixtyTwo} & Gaps21622170 & PROVED & \checkmark & verified \\
\texttt{nu\_p\_twoThousandOneHundredSixtyTwo\_two} & Gaps21622170 & PROVED & \checkmark & verified \\
\texttt{singular\_series\_finite\_pos\_evenPair\_twoThousandOneHundredSeventy} & Gaps21622170 & PROVED & \checkmark & verified \\
\texttt{singular\_series\_finite\_pos\_evenPair\_twoThousandOneHundredSixtyEight} & Gaps21622170 & PROVED & \checkmark & verified \\
\texttt{singular\_series\_finite\_pos\_evenPair\_twoThousandOneHundredSixtyFour} & Gaps21622170 & PROVED & \checkmark & verified \\
\texttt{singular\_series\_finite\_pos\_evenPair\_twoThousandOneHundredSixtySix} & Gaps21622170 & PROVED & \checkmark & verified \\
\texttt{singular\_series\_finite\_pos\_evenPair\_twoThousandOneHundredSixtyTwo} & Gaps21622170 & PROVED & \checkmark & verified \\
\texttt{singular\_series\_pos\_evenPair\_twoThousandOneHundredSeventy} & Gaps21622170 & PROVED & \checkmark & verified \\
\texttt{singular\_series\_pos\_evenPair\_twoThousandOneHundredSixtyEight} & Gaps21622170 & PROVED & \checkmark & verified \\
\texttt{singular\_series\_pos\_evenPair\_twoThousandOneHundredSixtyFour} & Gaps21622170 & PROVED & \checkmark & verified \\
\texttt{singular\_series\_pos\_evenPair\_twoThousandOneHundredSixtySix} & Gaps21622170 & PROVED & \checkmark & verified \\
\texttt{singular\_series\_pos\_evenPair\_twoThousandOneHundredSixtyTwo} & Gaps21622170 & PROVED & \checkmark & verified \\
\texttt{evenPair\_card\_twoThousandOneHundredEighty} & Gaps21722180 & PROVED & \checkmark & verified \\
\texttt{evenPair\_card\_twoThousandOneHundredSeventyEight} & Gaps21722180 & PROVED & \checkmark & verified \\
\texttt{evenPair\_card\_twoThousandOneHundredSeventyFour} & Gaps21722180 & PROVED & \checkmark & verified \\
\texttt{evenPair\_card\_twoThousandOneHundredSeventySix} & Gaps21722180 & PROVED & \checkmark & verified \\
\texttt{evenPair\_card\_twoThousandOneHundredSeventyTwo} & Gaps21722180 & PROVED & \checkmark & verified \\
\texttt{isAdmissible\_evenPair\_twoThousandOneHundredEighty} & Gaps21722180 & PROVED & \checkmark & verified \\
\texttt{isAdmissible\_evenPair\_twoThousandOneHundredSeventyEight} & Gaps21722180 & PROVED & \checkmark & verified \\
\texttt{isAdmissible\_evenPair\_twoThousandOneHundredSeventyFour} & Gaps21722180 & PROVED & \checkmark & verified \\
\texttt{isAdmissible\_evenPair\_twoThousandOneHundredSeventySix} & Gaps21722180 & PROVED & \checkmark & verified \\
\texttt{isAdmissible\_evenPair\_twoThousandOneHundredSeventyTwo} & Gaps21722180 & PROVED & \checkmark & verified \\
\texttt{localFactor\_twoThousandOneHundredEighty\_two} & Gaps21722180 & PROVED & \checkmark & verified \\
\texttt{localFactor\_twoThousandOneHundredSeventyTwo\_two} & Gaps21722180 & PROVED & \checkmark & verified \\
\texttt{nu\_p\_twoThousandOneHundredEighty} & Gaps21722180 & PROVED & \checkmark & verified \\
\texttt{nu\_p\_twoThousandOneHundredEighty\_two} & Gaps21722180 & PROVED & \checkmark & verified \\
\texttt{nu\_p\_twoThousandOneHundredSeventyEight} & Gaps21722180 & PROVED & \checkmark & verified \\
\texttt{nu\_p\_twoThousandOneHundredSeventyFour} & Gaps21722180 & PROVED & \checkmark & verified \\
\texttt{nu\_p\_twoThousandOneHundredSeventySix} & Gaps21722180 & PROVED & \checkmark & verified \\
\texttt{nu\_p\_twoThousandOneHundredSeventyTwo} & Gaps21722180 & PROVED & \checkmark & verified \\
\texttt{nu\_p\_twoThousandOneHundredSeventyTwo\_two} & Gaps21722180 & PROVED & \checkmark & verified \\
\texttt{singular\_series\_finite\_pos\_evenPair\_twoThousandOneHundredEighty} & Gaps21722180 & PROVED & \checkmark & verified \\
\texttt{singular\_series\_finite\_pos\_evenPair\_twoThousandOneHundredSeventyEight} & Gaps21722180 & PROVED & \checkmark & verified \\
\texttt{singular\_series\_finite\_pos\_evenPair\_twoThousandOneHundredSeventyFour} & Gaps21722180 & PROVED & \checkmark & verified \\
\texttt{singular\_series\_finite\_pos\_evenPair\_twoThousandOneHundredSeventySix} & Gaps21722180 & PROVED & \checkmark & verified \\
\texttt{singular\_series\_finite\_pos\_evenPair\_twoThousandOneHundredSeventyTwo} & Gaps21722180 & PROVED & \checkmark & verified \\
\texttt{singular\_series\_pos\_evenPair\_twoThousandOneHundredEighty} & Gaps21722180 & PROVED & \checkmark & verified \\
\texttt{singular\_series\_pos\_evenPair\_twoThousandOneHundredSeventyEight} & Gaps21722180 & PROVED & \checkmark & verified \\
\texttt{singular\_series\_pos\_evenPair\_twoThousandOneHundredSeventyFour} & Gaps21722180 & PROVED & \checkmark & verified \\
\texttt{singular\_series\_pos\_evenPair\_twoThousandOneHundredSeventySix} & Gaps21722180 & PROVED & \checkmark & verified \\
\texttt{singular\_series\_pos\_evenPair\_twoThousandOneHundredSeventyTwo} & Gaps21722180 & PROVED & \checkmark & verified \\
\texttt{evenPair\_card\_twoThousandOneHundredEightyEight} & Gaps21822190 & PROVED & \checkmark & verified \\
\texttt{evenPair\_card\_twoThousandOneHundredEightyFour} & Gaps21822190 & PROVED & \checkmark & verified \\
\texttt{evenPair\_card\_twoThousandOneHundredEightySix} & Gaps21822190 & PROVED & \checkmark & verified \\
\texttt{evenPair\_card\_twoThousandOneHundredEightyTwo} & Gaps21822190 & PROVED & \checkmark & verified \\
\texttt{evenPair\_card\_twoThousandOneHundredNinety} & Gaps21822190 & PROVED & \checkmark & verified \\
\texttt{isAdmissible\_evenPair\_twoThousandOneHundredEightyEight} & Gaps21822190 & PROVED & \checkmark & verified \\
\texttt{isAdmissible\_evenPair\_twoThousandOneHundredEightyFour} & Gaps21822190 & PROVED & \checkmark & verified \\
\texttt{isAdmissible\_evenPair\_twoThousandOneHundredEightySix} & Gaps21822190 & PROVED & \checkmark & verified \\
\texttt{isAdmissible\_evenPair\_twoThousandOneHundredEightyTwo} & Gaps21822190 & PROVED & \checkmark & verified \\
\texttt{isAdmissible\_evenPair\_twoThousandOneHundredNinety} & Gaps21822190 & PROVED & \checkmark & verified \\
\texttt{localFactor\_twoThousandOneHundredEightyTwo\_two} & Gaps21822190 & PROVED & \checkmark & verified \\
\texttt{localFactor\_twoThousandOneHundredNinety\_two} & Gaps21822190 & PROVED & \checkmark & verified \\
\texttt{nu\_p\_twoThousandOneHundredEightyEight} & Gaps21822190 & PROVED & \checkmark & verified \\
\texttt{nu\_p\_twoThousandOneHundredEightyFour} & Gaps21822190 & PROVED & \checkmark & verified \\
\texttt{nu\_p\_twoThousandOneHundredEightySix} & Gaps21822190 & PROVED & \checkmark & verified \\
\texttt{nu\_p\_twoThousandOneHundredEightyTwo} & Gaps21822190 & PROVED & \checkmark & verified \\
\texttt{nu\_p\_twoThousandOneHundredEightyTwo\_two} & Gaps21822190 & PROVED & \checkmark & verified \\
\texttt{nu\_p\_twoThousandOneHundredNinety} & Gaps21822190 & PROVED & \checkmark & verified \\
\texttt{nu\_p\_twoThousandOneHundredNinety\_two} & Gaps21822190 & PROVED & \checkmark & verified \\
\texttt{singular\_series\_finite\_pos\_evenPair\_twoThousandOneHundredEightyEight} & Gaps21822190 & PROVED & \checkmark & verified \\
\texttt{singular\_series\_finite\_pos\_evenPair\_twoThousandOneHundredEightyFour} & Gaps21822190 & PROVED & \checkmark & verified \\
\texttt{singular\_series\_finite\_pos\_evenPair\_twoThousandOneHundredEightySix} & Gaps21822190 & PROVED & \checkmark & verified \\
\texttt{singular\_series\_finite\_pos\_evenPair\_twoThousandOneHundredEightyTwo} & Gaps21822190 & PROVED & \checkmark & verified \\
\texttt{singular\_series\_finite\_pos\_evenPair\_twoThousandOneHundredNinety} & Gaps21822190 & PROVED & \checkmark & verified \\
\texttt{singular\_series\_pos\_evenPair\_twoThousandOneHundredEightyEight} & Gaps21822190 & PROVED & \checkmark & verified \\
\texttt{singular\_series\_pos\_evenPair\_twoThousandOneHundredEightyFour} & Gaps21822190 & PROVED & \checkmark & verified \\
\texttt{singular\_series\_pos\_evenPair\_twoThousandOneHundredEightySix} & Gaps21822190 & PROVED & \checkmark & verified \\
\texttt{singular\_series\_pos\_evenPair\_twoThousandOneHundredEightyTwo} & Gaps21822190 & PROVED & \checkmark & verified \\
\texttt{singular\_series\_pos\_evenPair\_twoThousandOneHundredNinety} & Gaps21822190 & PROVED & \checkmark & verified \\
\texttt{evenPair\_card\_twoThousandOneHundredNinetyEight} & Gaps21922200 & PROVED & \checkmark & verified \\
\texttt{evenPair\_card\_twoThousandOneHundredNinetyFour} & Gaps21922200 & PROVED & \checkmark & verified \\
\texttt{evenPair\_card\_twoThousandOneHundredNinetySix} & Gaps21922200 & PROVED & \checkmark & verified \\
\texttt{evenPair\_card\_twoThousandOneHundredNinetyTwo} & Gaps21922200 & PROVED & \checkmark & verified \\
\texttt{evenPair\_card\_twoThousandTwoHundred} & Gaps21922200 & PROVED & \checkmark & verified \\
\texttt{isAdmissible\_evenPair\_twoThousandOneHundredNinetyEight} & Gaps21922200 & PROVED & \checkmark & verified \\
\texttt{isAdmissible\_evenPair\_twoThousandOneHundredNinetyFour} & Gaps21922200 & PROVED & \checkmark & verified \\
\texttt{isAdmissible\_evenPair\_twoThousandOneHundredNinetySix} & Gaps21922200 & PROVED & \checkmark & verified \\
\texttt{isAdmissible\_evenPair\_twoThousandOneHundredNinetyTwo} & Gaps21922200 & PROVED & \checkmark & verified \\
\texttt{isAdmissible\_evenPair\_twoThousandTwoHundred} & Gaps21922200 & PROVED & \checkmark & verified \\
\texttt{localFactor\_twoThousandOneHundredNinetyTwo\_two} & Gaps21922200 & PROVED & \checkmark & verified \\
\texttt{localFactor\_twoThousandTwoHundred\_two} & Gaps21922200 & PROVED & \checkmark & verified \\
\texttt{nu\_p\_twoThousandOneHundredNinetyEight} & Gaps21922200 & PROVED & \checkmark & verified \\
\texttt{nu\_p\_twoThousandOneHundredNinetyFour} & Gaps21922200 & PROVED & \checkmark & verified \\
\texttt{nu\_p\_twoThousandOneHundredNinetySix} & Gaps21922200 & PROVED & \checkmark & verified \\
\texttt{nu\_p\_twoThousandOneHundredNinetyTwo} & Gaps21922200 & PROVED & \checkmark & verified \\
\texttt{nu\_p\_twoThousandOneHundredNinetyTwo\_two} & Gaps21922200 & PROVED & \checkmark & verified \\
\texttt{nu\_p\_twoThousandTwoHundred} & Gaps21922200 & PROVED & \checkmark & verified \\
\texttt{nu\_p\_twoThousandTwoHundred\_two} & Gaps21922200 & PROVED & \checkmark & verified \\
\texttt{singular\_series\_finite\_pos\_evenPair\_twoThousandOneHundredNinetyEight} & Gaps21922200 & PROVED & \checkmark & verified \\
\texttt{singular\_series\_finite\_pos\_evenPair\_twoThousandOneHundredNinetyFour} & Gaps21922200 & PROVED & \checkmark & verified \\
\texttt{singular\_series\_finite\_pos\_evenPair\_twoThousandOneHundredNinetySix} & Gaps21922200 & PROVED & \checkmark & verified \\
\texttt{singular\_series\_finite\_pos\_evenPair\_twoThousandOneHundredNinetyTwo} & Gaps21922200 & PROVED & \checkmark & verified \\
\texttt{singular\_series\_finite\_pos\_evenPair\_twoThousandTwoHundred} & Gaps21922200 & PROVED & \checkmark & verified \\
\texttt{singular\_series\_pos\_evenPair\_twoThousandOneHundredNinetyEight} & Gaps21922200 & PROVED & \checkmark & verified \\
\texttt{singular\_series\_pos\_evenPair\_twoThousandOneHundredNinetyFour} & Gaps21922200 & PROVED & \checkmark & verified \\
\texttt{singular\_series\_pos\_evenPair\_twoThousandOneHundredNinetySix} & Gaps21922200 & PROVED & \checkmark & verified \\
\texttt{singular\_series\_pos\_evenPair\_twoThousandOneHundredNinetyTwo} & Gaps21922200 & PROVED & \checkmark & verified \\
\texttt{singular\_series\_pos\_evenPair\_twoThousandTwoHundred} & Gaps21922200 & PROVED & \checkmark & verified \\
\texttt{evenPair\_card\_twoHundredThirty} & Gaps222230 & PROVED & \checkmark & verified \\
\texttt{evenPair\_card\_twoHundredTwentyEight} & Gaps222230 & PROVED & \checkmark & verified \\
\texttt{evenPair\_card\_twoHundredTwentyFour} & Gaps222230 & PROVED & \checkmark & verified \\
\texttt{evenPair\_card\_twoHundredTwentySix} & Gaps222230 & PROVED & \checkmark & verified \\
\texttt{evenPair\_card\_twoHundredTwentyTwo} & Gaps222230 & PROVED & \checkmark & verified \\
\texttt{isAdmissible\_evenPair\_twoHundredThirty} & Gaps222230 & PROVED & \checkmark & verified \\
\texttt{isAdmissible\_evenPair\_twoHundredTwentyEight} & Gaps222230 & PROVED & \checkmark & verified \\
\texttt{isAdmissible\_evenPair\_twoHundredTwentyFour} & Gaps222230 & PROVED & \checkmark & verified \\
\texttt{isAdmissible\_evenPair\_twoHundredTwentySix} & Gaps222230 & PROVED & \checkmark & verified \\
\texttt{isAdmissible\_evenPair\_twoHundredTwentyTwo} & Gaps222230 & PROVED & \checkmark & verified \\
\texttt{localFactor\_twoHundredThirty\_two} & Gaps222230 & PROVED & \checkmark & verified \\
\texttt{localFactor\_twoHundredTwentyTwo\_two} & Gaps222230 & PROVED & \checkmark & verified \\
\texttt{nu\_p\_twoHundredThirty} & Gaps222230 & PROVED & \checkmark & verified \\
\texttt{nu\_p\_twoHundredThirty\_two} & Gaps222230 & PROVED & \checkmark & verified \\
\texttt{nu\_p\_twoHundredTwentyEight} & Gaps222230 & PROVED & \checkmark & verified \\
\texttt{nu\_p\_twoHundredTwentyFour} & Gaps222230 & PROVED & \checkmark & verified \\
\texttt{nu\_p\_twoHundredTwentySix} & Gaps222230 & PROVED & \checkmark & verified \\
\texttt{nu\_p\_twoHundredTwentyTwo} & Gaps222230 & PROVED & \checkmark & verified \\
\texttt{nu\_p\_twoHundredTwentyTwo\_two} & Gaps222230 & PROVED & \checkmark & verified \\
\texttt{singular\_series\_finite\_pos\_evenPair\_twoHundredThirty} & Gaps222230 & PROVED & \checkmark & verified \\
\texttt{singular\_series\_finite\_pos\_evenPair\_twoHundredTwentyEight} & Gaps222230 & PROVED & \checkmark & verified \\
\texttt{singular\_series\_finite\_pos\_evenPair\_twoHundredTwentyFour} & Gaps222230 & PROVED & \checkmark & verified \\
\texttt{singular\_series\_finite\_pos\_evenPair\_twoHundredTwentySix} & Gaps222230 & PROVED & \checkmark & verified \\
\texttt{singular\_series\_finite\_pos\_evenPair\_twoHundredTwentyTwo} & Gaps222230 & PROVED & \checkmark & verified \\
\texttt{singular\_series\_pos\_evenPair\_twoHundredThirty} & Gaps222230 & PROVED & \checkmark & verified \\
\texttt{singular\_series\_pos\_evenPair\_twoHundredTwentyEight} & Gaps222230 & PROVED & \checkmark & verified \\
\texttt{singular\_series\_pos\_evenPair\_twoHundredTwentyFour} & Gaps222230 & PROVED & \checkmark & verified \\
\texttt{singular\_series\_pos\_evenPair\_twoHundredTwentySix} & Gaps222230 & PROVED & \checkmark & verified \\
\texttt{singular\_series\_pos\_evenPair\_twoHundredTwentyTwo} & Gaps222230 & PROVED & \checkmark & verified \\
\texttt{evenPair\_card\_thirty} & Gaps2230 & PROVED & \checkmark & verified \\
\texttt{evenPair\_card\_twentyEight} & Gaps2230 & PROVED & \checkmark & verified \\
\texttt{evenPair\_card\_twentyFour} & Gaps2230 & PROVED & \checkmark & verified \\
\texttt{evenPair\_card\_twentySix} & Gaps2230 & PROVED & \checkmark & verified \\
\texttt{evenPair\_card\_twentyTwo} & Gaps2230 & PROVED & \checkmark & verified \\
\texttt{isAdmissible\_evenPair\_thirty} & Gaps2230 & PROVED & \checkmark & verified \\
\texttt{isAdmissible\_evenPair\_twentyEight} & Gaps2230 & PROVED & \checkmark & verified \\
\texttt{isAdmissible\_evenPair\_twentyFour} & Gaps2230 & PROVED & \checkmark & verified \\
\texttt{isAdmissible\_evenPair\_twentySix} & Gaps2230 & PROVED & \checkmark & verified \\
\texttt{isAdmissible\_evenPair\_twentyTwo} & Gaps2230 & PROVED & \checkmark & verified \\
\texttt{localFactor\_thirty\_two} & Gaps2230 & PROVED & \checkmark & verified \\
\texttt{localFactor\_twentyTwo\_two} & Gaps2230 & PROVED & \checkmark & verified \\
\texttt{nu\_p\_thirty} & Gaps2230 & PROVED & \checkmark & verified \\
\texttt{nu\_p\_thirty\_two} & Gaps2230 & PROVED & \checkmark & verified \\
\texttt{nu\_p\_twentyEight} & Gaps2230 & PROVED & \checkmark & verified \\
\texttt{nu\_p\_twentyFour} & Gaps2230 & PROVED & \checkmark & verified \\
\texttt{nu\_p\_twentySix} & Gaps2230 & PROVED & \checkmark & verified \\
\texttt{nu\_p\_twentyTwo} & Gaps2230 & PROVED & \checkmark & verified \\
\texttt{nu\_p\_twentyTwo\_two} & Gaps2230 & PROVED & \checkmark & verified \\
\texttt{singular\_series\_finite\_pos\_evenPair\_thirty} & Gaps2230 & PROVED & \checkmark & verified \\
\texttt{singular\_series\_finite\_pos\_evenPair\_twentyEight} & Gaps2230 & PROVED & \checkmark & verified \\
\texttt{singular\_series\_finite\_pos\_evenPair\_twentyFour} & Gaps2230 & PROVED & \checkmark & verified \\
\texttt{singular\_series\_finite\_pos\_evenPair\_twentySix} & Gaps2230 & PROVED & \checkmark & verified \\
\texttt{singular\_series\_finite\_pos\_evenPair\_twentyTwo} & Gaps2230 & PROVED & \checkmark & verified \\
\texttt{singular\_series\_pos\_evenPair\_thirty} & Gaps2230 & PROVED & \checkmark & verified \\
\texttt{singular\_series\_pos\_evenPair\_twentyEight} & Gaps2230 & PROVED & \checkmark & verified \\
\texttt{singular\_series\_pos\_evenPair\_twentyFour} & Gaps2230 & PROVED & \checkmark & verified \\
\texttt{singular\_series\_pos\_evenPair\_twentySix} & Gaps2230 & PROVED & \checkmark & verified \\
\texttt{singular\_series\_pos\_evenPair\_twentyTwo} & Gaps2230 & PROVED & \checkmark & verified \\
\texttt{evenPair\_card\_twoHundredForty} & Gaps232240 & PROVED & \checkmark & verified \\
\texttt{evenPair\_card\_twoHundredThirtyEight} & Gaps232240 & PROVED & \checkmark & verified \\
\texttt{evenPair\_card\_twoHundredThirtyFour} & Gaps232240 & PROVED & \checkmark & verified \\
\texttt{evenPair\_card\_twoHundredThirtySix} & Gaps232240 & PROVED & \checkmark & verified \\
\texttt{evenPair\_card\_twoHundredThirtyTwo} & Gaps232240 & PROVED & \checkmark & verified \\
\texttt{isAdmissible\_evenPair\_twoHundredForty} & Gaps232240 & PROVED & \checkmark & verified \\
\texttt{isAdmissible\_evenPair\_twoHundredThirtyEight} & Gaps232240 & PROVED & \checkmark & verified \\
\texttt{isAdmissible\_evenPair\_twoHundredThirtyFour} & Gaps232240 & PROVED & \checkmark & verified \\
\texttt{isAdmissible\_evenPair\_twoHundredThirtySix} & Gaps232240 & PROVED & \checkmark & verified \\
\texttt{isAdmissible\_evenPair\_twoHundredThirtyTwo} & Gaps232240 & PROVED & \checkmark & verified \\
\texttt{localFactor\_twoHundredForty\_two} & Gaps232240 & PROVED & \checkmark & verified \\
\texttt{localFactor\_twoHundredThirtyTwo\_two} & Gaps232240 & PROVED & \checkmark & verified \\
\texttt{nu\_p\_twoHundredForty} & Gaps232240 & PROVED & \checkmark & verified \\
\texttt{nu\_p\_twoHundredForty\_two} & Gaps232240 & PROVED & \checkmark & verified \\
\texttt{nu\_p\_twoHundredThirtyEight} & Gaps232240 & PROVED & \checkmark & verified \\
\texttt{nu\_p\_twoHundredThirtyFour} & Gaps232240 & PROVED & \checkmark & verified \\
\texttt{nu\_p\_twoHundredThirtySix} & Gaps232240 & PROVED & \checkmark & verified \\
\texttt{nu\_p\_twoHundredThirtyTwo} & Gaps232240 & PROVED & \checkmark & verified \\
\texttt{nu\_p\_twoHundredThirtyTwo\_two} & Gaps232240 & PROVED & \checkmark & verified \\
\texttt{singular\_series\_finite\_pos\_evenPair\_twoHundredForty} & Gaps232240 & PROVED & \checkmark & verified \\
\texttt{singular\_series\_finite\_pos\_evenPair\_twoHundredThirtyEight} & Gaps232240 & PROVED & \checkmark & verified \\
\texttt{singular\_series\_finite\_pos\_evenPair\_twoHundredThirtyFour} & Gaps232240 & PROVED & \checkmark & verified \\
\texttt{singular\_series\_finite\_pos\_evenPair\_twoHundredThirtySix} & Gaps232240 & PROVED & \checkmark & verified \\
\texttt{singular\_series\_finite\_pos\_evenPair\_twoHundredThirtyTwo} & Gaps232240 & PROVED & \checkmark & verified \\
\texttt{singular\_series\_pos\_evenPair\_twoHundredForty} & Gaps232240 & PROVED & \checkmark & verified \\
\texttt{singular\_series\_pos\_evenPair\_twoHundredThirtyEight} & Gaps232240 & PROVED & \checkmark & verified \\
\texttt{singular\_series\_pos\_evenPair\_twoHundredThirtyFour} & Gaps232240 & PROVED & \checkmark & verified \\
\texttt{singular\_series\_pos\_evenPair\_twoHundredThirtySix} & Gaps232240 & PROVED & \checkmark & verified \\
\texttt{singular\_series\_pos\_evenPair\_twoHundredThirtyTwo} & Gaps232240 & PROVED & \checkmark & verified \\
\texttt{evenPair\_card\_twoHundredFifty} & Gaps242250 & PROVED & \checkmark & verified \\
\texttt{evenPair\_card\_twoHundredFortyEight} & Gaps242250 & PROVED & \checkmark & verified \\
\texttt{evenPair\_card\_twoHundredFortyFour} & Gaps242250 & PROVED & \checkmark & verified \\
\texttt{evenPair\_card\_twoHundredFortySix} & Gaps242250 & PROVED & \checkmark & verified \\
\texttt{evenPair\_card\_twoHundredFortyTwo} & Gaps242250 & PROVED & \checkmark & verified \\
\texttt{isAdmissible\_evenPair\_twoHundredFifty} & Gaps242250 & PROVED & \checkmark & verified \\
\texttt{isAdmissible\_evenPair\_twoHundredFortyEight} & Gaps242250 & PROVED & \checkmark & verified \\
\texttt{isAdmissible\_evenPair\_twoHundredFortyFour} & Gaps242250 & PROVED & \checkmark & verified \\
\texttt{isAdmissible\_evenPair\_twoHundredFortySix} & Gaps242250 & PROVED & \checkmark & verified \\
\texttt{isAdmissible\_evenPair\_twoHundredFortyTwo} & Gaps242250 & PROVED & \checkmark & verified \\
\texttt{localFactor\_twoHundredFifty\_two} & Gaps242250 & PROVED & \checkmark & verified \\
\texttt{localFactor\_twoHundredFortyTwo\_two} & Gaps242250 & PROVED & \checkmark & verified \\
\texttt{nu\_p\_twoHundredFifty} & Gaps242250 & PROVED & \checkmark & verified \\
\texttt{nu\_p\_twoHundredFifty\_two} & Gaps242250 & PROVED & \checkmark & verified \\
\texttt{nu\_p\_twoHundredFortyEight} & Gaps242250 & PROVED & \checkmark & verified \\
\texttt{nu\_p\_twoHundredFortyFour} & Gaps242250 & PROVED & \checkmark & verified \\
\texttt{nu\_p\_twoHundredFortySix} & Gaps242250 & PROVED & \checkmark & verified \\
\texttt{nu\_p\_twoHundredFortyTwo} & Gaps242250 & PROVED & \checkmark & verified \\
\texttt{nu\_p\_twoHundredFortyTwo\_two} & Gaps242250 & PROVED & \checkmark & verified \\
\texttt{singular\_series\_finite\_pos\_evenPair\_twoHundredFifty} & Gaps242250 & PROVED & \checkmark & verified \\
\texttt{singular\_series\_finite\_pos\_evenPair\_twoHundredFortyEight} & Gaps242250 & PROVED & \checkmark & verified \\
\texttt{singular\_series\_finite\_pos\_evenPair\_twoHundredFortyFour} & Gaps242250 & PROVED & \checkmark & verified \\
\texttt{singular\_series\_finite\_pos\_evenPair\_twoHundredFortySix} & Gaps242250 & PROVED & \checkmark & verified \\
\texttt{singular\_series\_finite\_pos\_evenPair\_twoHundredFortyTwo} & Gaps242250 & PROVED & \checkmark & verified \\
\texttt{singular\_series\_pos\_evenPair\_twoHundredFifty} & Gaps242250 & PROVED & \checkmark & verified \\
\texttt{singular\_series\_pos\_evenPair\_twoHundredFortyEight} & Gaps242250 & PROVED & \checkmark & verified \\
\texttt{singular\_series\_pos\_evenPair\_twoHundredFortyFour} & Gaps242250 & PROVED & \checkmark & verified \\
\texttt{singular\_series\_pos\_evenPair\_twoHundredFortySix} & Gaps242250 & PROVED & \checkmark & verified \\
\texttt{singular\_series\_pos\_evenPair\_twoHundredFortyTwo} & Gaps242250 & PROVED & \checkmark & verified \\
\texttt{evenPair\_card\_twoHundredFiftyEight} & Gaps252260 & PROVED & \checkmark & verified \\
\texttt{evenPair\_card\_twoHundredFiftyFour} & Gaps252260 & PROVED & \checkmark & verified \\
\texttt{evenPair\_card\_twoHundredFiftySix} & Gaps252260 & PROVED & \checkmark & verified \\
\texttt{evenPair\_card\_twoHundredFiftyTwo} & Gaps252260 & PROVED & \checkmark & verified \\
\texttt{evenPair\_card\_twoHundredSixty} & Gaps252260 & PROVED & \checkmark & verified \\
\texttt{isAdmissible\_evenPair\_twoHundredFiftyEight} & Gaps252260 & PROVED & \checkmark & verified \\
\texttt{isAdmissible\_evenPair\_twoHundredFiftyFour} & Gaps252260 & PROVED & \checkmark & verified \\
\texttt{isAdmissible\_evenPair\_twoHundredFiftySix} & Gaps252260 & PROVED & \checkmark & verified \\
\texttt{isAdmissible\_evenPair\_twoHundredFiftyTwo} & Gaps252260 & PROVED & \checkmark & verified \\
\texttt{isAdmissible\_evenPair\_twoHundredSixty} & Gaps252260 & PROVED & \checkmark & verified \\
\texttt{localFactor\_twoHundredFiftyTwo\_two} & Gaps252260 & PROVED & \checkmark & verified \\
\texttt{localFactor\_twoHundredSixty\_two} & Gaps252260 & PROVED & \checkmark & verified \\
\texttt{nu\_p\_twoHundredFiftyEight} & Gaps252260 & PROVED & \checkmark & verified \\
\texttt{nu\_p\_twoHundredFiftyFour} & Gaps252260 & PROVED & \checkmark & verified \\
\texttt{nu\_p\_twoHundredFiftySix} & Gaps252260 & PROVED & \checkmark & verified \\
\texttt{nu\_p\_twoHundredFiftyTwo} & Gaps252260 & PROVED & \checkmark & verified \\
\texttt{nu\_p\_twoHundredFiftyTwo\_two} & Gaps252260 & PROVED & \checkmark & verified \\
\texttt{nu\_p\_twoHundredSixty} & Gaps252260 & PROVED & \checkmark & verified \\
\texttt{nu\_p\_twoHundredSixty\_two} & Gaps252260 & PROVED & \checkmark & verified \\
\texttt{singular\_series\_finite\_pos\_evenPair\_twoHundredFiftyEight} & Gaps252260 & PROVED & \checkmark & verified \\
\texttt{singular\_series\_finite\_pos\_evenPair\_twoHundredFiftyFour} & Gaps252260 & PROVED & \checkmark & verified \\
\texttt{singular\_series\_finite\_pos\_evenPair\_twoHundredFiftySix} & Gaps252260 & PROVED & \checkmark & verified \\
\texttt{singular\_series\_finite\_pos\_evenPair\_twoHundredFiftyTwo} & Gaps252260 & PROVED & \checkmark & verified \\
\texttt{singular\_series\_finite\_pos\_evenPair\_twoHundredSixty} & Gaps252260 & PROVED & \checkmark & verified \\
\texttt{singular\_series\_pos\_evenPair\_twoHundredFiftyEight} & Gaps252260 & PROVED & \checkmark & verified \\
\texttt{singular\_series\_pos\_evenPair\_twoHundredFiftyFour} & Gaps252260 & PROVED & \checkmark & verified \\
\texttt{singular\_series\_pos\_evenPair\_twoHundredFiftySix} & Gaps252260 & PROVED & \checkmark & verified \\
\texttt{singular\_series\_pos\_evenPair\_twoHundredFiftyTwo} & Gaps252260 & PROVED & \checkmark & verified \\
\texttt{singular\_series\_pos\_evenPair\_twoHundredSixty} & Gaps252260 & PROVED & \checkmark & verified \\
\texttt{evenPair\_card\_twoHundredSeventy} & Gaps262270 & PROVED & \checkmark & verified \\
\texttt{evenPair\_card\_twoHundredSixtyEight} & Gaps262270 & PROVED & \checkmark & verified \\
\texttt{evenPair\_card\_twoHundredSixtyFour} & Gaps262270 & PROVED & \checkmark & verified \\
\texttt{evenPair\_card\_twoHundredSixtySix} & Gaps262270 & PROVED & \checkmark & verified \\
\texttt{evenPair\_card\_twoHundredSixtyTwo} & Gaps262270 & PROVED & \checkmark & verified \\
\texttt{isAdmissible\_evenPair\_twoHundredSeventy} & Gaps262270 & PROVED & \checkmark & verified \\
\texttt{isAdmissible\_evenPair\_twoHundredSixtyEight} & Gaps262270 & PROVED & \checkmark & verified \\
\texttt{isAdmissible\_evenPair\_twoHundredSixtyFour} & Gaps262270 & PROVED & \checkmark & verified \\
\texttt{isAdmissible\_evenPair\_twoHundredSixtySix} & Gaps262270 & PROVED & \checkmark & verified \\
\texttt{isAdmissible\_evenPair\_twoHundredSixtyTwo} & Gaps262270 & PROVED & \checkmark & verified \\
\texttt{localFactor\_twoHundredSeventy\_two} & Gaps262270 & PROVED & \checkmark & verified \\
\texttt{localFactor\_twoHundredSixtyTwo\_two} & Gaps262270 & PROVED & \checkmark & verified \\
\texttt{nu\_p\_twoHundredSeventy} & Gaps262270 & PROVED & \checkmark & verified \\
\texttt{nu\_p\_twoHundredSeventy\_two} & Gaps262270 & PROVED & \checkmark & verified \\
\texttt{nu\_p\_twoHundredSixtyEight} & Gaps262270 & PROVED & \checkmark & verified \\
\texttt{nu\_p\_twoHundredSixtyFour} & Gaps262270 & PROVED & \checkmark & verified \\
\texttt{nu\_p\_twoHundredSixtySix} & Gaps262270 & PROVED & \checkmark & verified \\
\texttt{nu\_p\_twoHundredSixtyTwo} & Gaps262270 & PROVED & \checkmark & verified \\
\texttt{nu\_p\_twoHundredSixtyTwo\_two} & Gaps262270 & PROVED & \checkmark & verified \\
\texttt{singular\_series\_finite\_pos\_evenPair\_twoHundredSeventy} & Gaps262270 & PROVED & \checkmark & verified \\
\texttt{singular\_series\_finite\_pos\_evenPair\_twoHundredSixtyEight} & Gaps262270 & PROVED & \checkmark & verified \\
\texttt{singular\_series\_finite\_pos\_evenPair\_twoHundredSixtyFour} & Gaps262270 & PROVED & \checkmark & verified \\
\texttt{singular\_series\_finite\_pos\_evenPair\_twoHundredSixtySix} & Gaps262270 & PROVED & \checkmark & verified \\
\texttt{singular\_series\_finite\_pos\_evenPair\_twoHundredSixtyTwo} & Gaps262270 & PROVED & \checkmark & verified \\
\texttt{singular\_series\_pos\_evenPair\_twoHundredSeventy} & Gaps262270 & PROVED & \checkmark & verified \\
\texttt{singular\_series\_pos\_evenPair\_twoHundredSixtyEight} & Gaps262270 & PROVED & \checkmark & verified \\
\texttt{singular\_series\_pos\_evenPair\_twoHundredSixtyFour} & Gaps262270 & PROVED & \checkmark & verified \\
\texttt{singular\_series\_pos\_evenPair\_twoHundredSixtySix} & Gaps262270 & PROVED & \checkmark & verified \\
\texttt{singular\_series\_pos\_evenPair\_twoHundredSixtyTwo} & Gaps262270 & PROVED & \checkmark & verified \\
\texttt{evenPair\_card\_twoHundredEighty} & Gaps272280 & PROVED & \checkmark & verified \\
\texttt{evenPair\_card\_twoHundredSeventyEight} & Gaps272280 & PROVED & \checkmark & verified \\
\texttt{evenPair\_card\_twoHundredSeventyFour} & Gaps272280 & PROVED & \checkmark & verified \\
\texttt{evenPair\_card\_twoHundredSeventySix} & Gaps272280 & PROVED & \checkmark & verified \\
\texttt{evenPair\_card\_twoHundredSeventyTwo} & Gaps272280 & PROVED & \checkmark & verified \\
\texttt{isAdmissible\_evenPair\_twoHundredEighty} & Gaps272280 & PROVED & \checkmark & verified \\
\texttt{isAdmissible\_evenPair\_twoHundredSeventyEight} & Gaps272280 & PROVED & \checkmark & verified \\
\texttt{isAdmissible\_evenPair\_twoHundredSeventyFour} & Gaps272280 & PROVED & \checkmark & verified \\
\texttt{isAdmissible\_evenPair\_twoHundredSeventySix} & Gaps272280 & PROVED & \checkmark & verified \\
\texttt{isAdmissible\_evenPair\_twoHundredSeventyTwo} & Gaps272280 & PROVED & \checkmark & verified \\
\texttt{localFactor\_twoHundredEighty\_two} & Gaps272280 & PROVED & \checkmark & verified \\
\texttt{localFactor\_twoHundredSeventyTwo\_two} & Gaps272280 & PROVED & \checkmark & verified \\
\texttt{nu\_p\_twoHundredEighty} & Gaps272280 & PROVED & \checkmark & verified \\
\texttt{nu\_p\_twoHundredEighty\_two} & Gaps272280 & PROVED & \checkmark & verified \\
\texttt{nu\_p\_twoHundredSeventyEight} & Gaps272280 & PROVED & \checkmark & verified \\
\texttt{nu\_p\_twoHundredSeventyFour} & Gaps272280 & PROVED & \checkmark & verified \\
\texttt{nu\_p\_twoHundredSeventySix} & Gaps272280 & PROVED & \checkmark & verified \\
\texttt{nu\_p\_twoHundredSeventyTwo} & Gaps272280 & PROVED & \checkmark & verified \\
\texttt{nu\_p\_twoHundredSeventyTwo\_two} & Gaps272280 & PROVED & \checkmark & verified \\
\texttt{singular\_series\_finite\_pos\_evenPair\_twoHundredEighty} & Gaps272280 & PROVED & \checkmark & verified \\
\texttt{singular\_series\_finite\_pos\_evenPair\_twoHundredSeventyEight} & Gaps272280 & PROVED & \checkmark & verified \\
\texttt{singular\_series\_finite\_pos\_evenPair\_twoHundredSeventyFour} & Gaps272280 & PROVED & \checkmark & verified \\
\texttt{singular\_series\_finite\_pos\_evenPair\_twoHundredSeventySix} & Gaps272280 & PROVED & \checkmark & verified \\
\texttt{singular\_series\_finite\_pos\_evenPair\_twoHundredSeventyTwo} & Gaps272280 & PROVED & \checkmark & verified \\
\texttt{singular\_series\_pos\_evenPair\_twoHundredEighty} & Gaps272280 & PROVED & \checkmark & verified \\
\texttt{singular\_series\_pos\_evenPair\_twoHundredSeventyEight} & Gaps272280 & PROVED & \checkmark & verified \\
\texttt{singular\_series\_pos\_evenPair\_twoHundredSeventyFour} & Gaps272280 & PROVED & \checkmark & verified \\
\texttt{singular\_series\_pos\_evenPair\_twoHundredSeventySix} & Gaps272280 & PROVED & \checkmark & verified \\
\texttt{singular\_series\_pos\_evenPair\_twoHundredSeventyTwo} & Gaps272280 & PROVED & \checkmark & verified \\
\texttt{evenPair\_card\_twoHundredEightyEight} & Gaps282290 & PROVED & \checkmark & verified \\
\texttt{evenPair\_card\_twoHundredEightyFour} & Gaps282290 & PROVED & \checkmark & verified \\
\texttt{evenPair\_card\_twoHundredEightySix} & Gaps282290 & PROVED & \checkmark & verified \\
\texttt{evenPair\_card\_twoHundredEightyTwo} & Gaps282290 & PROVED & \checkmark & verified \\
\texttt{evenPair\_card\_twoHundredNinety} & Gaps282290 & PROVED & \checkmark & verified \\
\texttt{isAdmissible\_evenPair\_twoHundredEightyEight} & Gaps282290 & PROVED & \checkmark & verified \\
\texttt{isAdmissible\_evenPair\_twoHundredEightyFour} & Gaps282290 & PROVED & \checkmark & verified \\
\texttt{isAdmissible\_evenPair\_twoHundredEightySix} & Gaps282290 & PROVED & \checkmark & verified \\
\texttt{isAdmissible\_evenPair\_twoHundredEightyTwo} & Gaps282290 & PROVED & \checkmark & verified \\
\texttt{isAdmissible\_evenPair\_twoHundredNinety} & Gaps282290 & PROVED & \checkmark & verified \\
\texttt{localFactor\_twoHundredEightyTwo\_two} & Gaps282290 & PROVED & \checkmark & verified \\
\texttt{localFactor\_twoHundredNinety\_two} & Gaps282290 & PROVED & \checkmark & verified \\
\texttt{nu\_p\_twoHundredEightyEight} & Gaps282290 & PROVED & \checkmark & verified \\
\texttt{nu\_p\_twoHundredEightyFour} & Gaps282290 & PROVED & \checkmark & verified \\
\texttt{nu\_p\_twoHundredEightySix} & Gaps282290 & PROVED & \checkmark & verified \\
\texttt{nu\_p\_twoHundredEightyTwo} & Gaps282290 & PROVED & \checkmark & verified \\
\texttt{nu\_p\_twoHundredEightyTwo\_two} & Gaps282290 & PROVED & \checkmark & verified \\
\texttt{nu\_p\_twoHundredNinety} & Gaps282290 & PROVED & \checkmark & verified \\
\texttt{nu\_p\_twoHundredNinety\_two} & Gaps282290 & PROVED & \checkmark & verified \\
\texttt{singular\_series\_finite\_pos\_evenPair\_twoHundredEightyEight} & Gaps282290 & PROVED & \checkmark & verified \\
\texttt{singular\_series\_finite\_pos\_evenPair\_twoHundredEightyFour} & Gaps282290 & PROVED & \checkmark & verified \\
\texttt{singular\_series\_finite\_pos\_evenPair\_twoHundredEightySix} & Gaps282290 & PROVED & \checkmark & verified \\
\texttt{singular\_series\_finite\_pos\_evenPair\_twoHundredEightyTwo} & Gaps282290 & PROVED & \checkmark & verified \\
\texttt{singular\_series\_finite\_pos\_evenPair\_twoHundredNinety} & Gaps282290 & PROVED & \checkmark & verified \\
\texttt{singular\_series\_pos\_evenPair\_twoHundredEightyEight} & Gaps282290 & PROVED & \checkmark & verified \\
\texttt{singular\_series\_pos\_evenPair\_twoHundredEightyFour} & Gaps282290 & PROVED & \checkmark & verified \\
\texttt{singular\_series\_pos\_evenPair\_twoHundredEightySix} & Gaps282290 & PROVED & \checkmark & verified \\
\texttt{singular\_series\_pos\_evenPair\_twoHundredEightyTwo} & Gaps282290 & PROVED & \checkmark & verified \\
\texttt{singular\_series\_pos\_evenPair\_twoHundredNinety} & Gaps282290 & PROVED & \checkmark & verified \\
\texttt{evenPair\_card\_threeHundred} & Gaps292300 & PROVED & \checkmark & verified \\
\texttt{evenPair\_card\_twoHundredNinetyEight} & Gaps292300 & PROVED & \checkmark & verified \\
\texttt{evenPair\_card\_twoHundredNinetyFour} & Gaps292300 & PROVED & \checkmark & verified \\
\texttt{evenPair\_card\_twoHundredNinetySix} & Gaps292300 & PROVED & \checkmark & verified \\
\texttt{evenPair\_card\_twoHundredNinetyTwo} & Gaps292300 & PROVED & \checkmark & verified \\
\texttt{isAdmissible\_evenPair\_threeHundred} & Gaps292300 & PROVED & \checkmark & verified \\
\texttt{isAdmissible\_evenPair\_twoHundredNinetyEight} & Gaps292300 & PROVED & \checkmark & verified \\
\texttt{isAdmissible\_evenPair\_twoHundredNinetyFour} & Gaps292300 & PROVED & \checkmark & verified \\
\texttt{isAdmissible\_evenPair\_twoHundredNinetySix} & Gaps292300 & PROVED & \checkmark & verified \\
\texttt{isAdmissible\_evenPair\_twoHundredNinetyTwo} & Gaps292300 & PROVED & \checkmark & verified \\
\texttt{localFactor\_threeHundred\_two} & Gaps292300 & PROVED & \checkmark & verified \\
\texttt{localFactor\_twoHundredNinetyTwo\_two} & Gaps292300 & PROVED & \checkmark & verified \\
\texttt{nu\_p\_threeHundred} & Gaps292300 & PROVED & \checkmark & verified \\
\texttt{nu\_p\_threeHundred\_two} & Gaps292300 & PROVED & \checkmark & verified \\
\texttt{nu\_p\_twoHundredNinetyEight} & Gaps292300 & PROVED & \checkmark & verified \\
\texttt{nu\_p\_twoHundredNinetyFour} & Gaps292300 & PROVED & \checkmark & verified \\
\texttt{nu\_p\_twoHundredNinetySix} & Gaps292300 & PROVED & \checkmark & verified \\
\texttt{nu\_p\_twoHundredNinetyTwo} & Gaps292300 & PROVED & \checkmark & verified \\
\texttt{nu\_p\_twoHundredNinetyTwo\_two} & Gaps292300 & PROVED & \checkmark & verified \\
\texttt{singular\_series\_finite\_pos\_evenPair\_threeHundred} & Gaps292300 & PROVED & \checkmark & verified \\
\texttt{singular\_series\_finite\_pos\_evenPair\_twoHundredNinetyEight} & Gaps292300 & PROVED & \checkmark & verified \\
\texttt{singular\_series\_finite\_pos\_evenPair\_twoHundredNinetyFour} & Gaps292300 & PROVED & \checkmark & verified \\
\texttt{singular\_series\_finite\_pos\_evenPair\_twoHundredNinetySix} & Gaps292300 & PROVED & \checkmark & verified \\
\texttt{singular\_series\_finite\_pos\_evenPair\_twoHundredNinetyTwo} & Gaps292300 & PROVED & \checkmark & verified \\
\texttt{singular\_series\_pos\_evenPair\_threeHundred} & Gaps292300 & PROVED & \checkmark & verified \\
\texttt{singular\_series\_pos\_evenPair\_twoHundredNinetyEight} & Gaps292300 & PROVED & \checkmark & verified \\
\texttt{singular\_series\_pos\_evenPair\_twoHundredNinetyFour} & Gaps292300 & PROVED & \checkmark & verified \\
\texttt{singular\_series\_pos\_evenPair\_twoHundredNinetySix} & Gaps292300 & PROVED & \checkmark & verified \\
\texttt{singular\_series\_pos\_evenPair\_twoHundredNinetyTwo} & Gaps292300 & PROVED & \checkmark & verified \\
\texttt{evenPair\_card\_threeHundredEight} & Gaps302310 & PROVED & \checkmark & verified \\
\texttt{evenPair\_card\_threeHundredFour} & Gaps302310 & PROVED & \checkmark & verified \\
\texttt{evenPair\_card\_threeHundredSix} & Gaps302310 & PROVED & \checkmark & verified \\
\texttt{evenPair\_card\_threeHundredTen} & Gaps302310 & PROVED & \checkmark & verified \\
\texttt{evenPair\_card\_threeHundredTwo} & Gaps302310 & PROVED & \checkmark & verified \\
\texttt{isAdmissible\_evenPair\_threeHundredEight} & Gaps302310 & PROVED & \checkmark & verified \\
\texttt{isAdmissible\_evenPair\_threeHundredFour} & Gaps302310 & PROVED & \checkmark & verified \\
\texttt{isAdmissible\_evenPair\_threeHundredSix} & Gaps302310 & PROVED & \checkmark & verified \\
\texttt{isAdmissible\_evenPair\_threeHundredTen} & Gaps302310 & PROVED & \checkmark & verified \\
\texttt{isAdmissible\_evenPair\_threeHundredTwo} & Gaps302310 & PROVED & \checkmark & verified \\
\texttt{localFactor\_threeHundredTen\_two} & Gaps302310 & PROVED & \checkmark & verified \\
\texttt{localFactor\_threeHundredTwo\_two} & Gaps302310 & PROVED & \checkmark & verified \\
\texttt{nu\_p\_threeHundredEight} & Gaps302310 & PROVED & \checkmark & verified \\
\texttt{nu\_p\_threeHundredFour} & Gaps302310 & PROVED & \checkmark & verified \\
\texttt{nu\_p\_threeHundredSix} & Gaps302310 & PROVED & \checkmark & verified \\
\texttt{nu\_p\_threeHundredTen} & Gaps302310 & PROVED & \checkmark & verified \\
\texttt{nu\_p\_threeHundredTen\_two} & Gaps302310 & PROVED & \checkmark & verified \\
\texttt{nu\_p\_threeHundredTwo} & Gaps302310 & PROVED & \checkmark & verified \\
\texttt{nu\_p\_threeHundredTwo\_two} & Gaps302310 & PROVED & \checkmark & verified \\
\texttt{singular\_series\_finite\_pos\_evenPair\_threeHundredEight} & Gaps302310 & PROVED & \checkmark & verified \\
\texttt{singular\_series\_finite\_pos\_evenPair\_threeHundredFour} & Gaps302310 & PROVED & \checkmark & verified \\
\texttt{singular\_series\_finite\_pos\_evenPair\_threeHundredSix} & Gaps302310 & PROVED & \checkmark & verified \\
\texttt{singular\_series\_finite\_pos\_evenPair\_threeHundredTen} & Gaps302310 & PROVED & \checkmark & verified \\
\texttt{singular\_series\_finite\_pos\_evenPair\_threeHundredTwo} & Gaps302310 & PROVED & \checkmark & verified \\
\texttt{singular\_series\_pos\_evenPair\_threeHundredEight} & Gaps302310 & PROVED & \checkmark & verified \\
\texttt{singular\_series\_pos\_evenPair\_threeHundredFour} & Gaps302310 & PROVED & \checkmark & verified \\
\texttt{singular\_series\_pos\_evenPair\_threeHundredSix} & Gaps302310 & PROVED & \checkmark & verified \\
\texttt{singular\_series\_pos\_evenPair\_threeHundredTen} & Gaps302310 & PROVED & \checkmark & verified \\
\texttt{singular\_series\_pos\_evenPair\_threeHundredTwo} & Gaps302310 & PROVED & \checkmark & verified \\
\texttt{evenPair\_card\_threeHundredEighteen} & Gaps312320 & PROVED & \checkmark & verified \\
\texttt{evenPair\_card\_threeHundredFourteen} & Gaps312320 & PROVED & \checkmark & verified \\
\texttt{evenPair\_card\_threeHundredSixteen} & Gaps312320 & PROVED & \checkmark & verified \\
\texttt{evenPair\_card\_threeHundredTwelve} & Gaps312320 & PROVED & \checkmark & verified \\
\texttt{evenPair\_card\_threeHundredTwenty} & Gaps312320 & PROVED & \checkmark & verified \\
\texttt{isAdmissible\_evenPair\_threeHundredEighteen} & Gaps312320 & PROVED & \checkmark & verified \\
\texttt{isAdmissible\_evenPair\_threeHundredFourteen} & Gaps312320 & PROVED & \checkmark & verified \\
\texttt{isAdmissible\_evenPair\_threeHundredSixteen} & Gaps312320 & PROVED & \checkmark & verified \\
\texttt{isAdmissible\_evenPair\_threeHundredTwelve} & Gaps312320 & PROVED & \checkmark & verified \\
\texttt{isAdmissible\_evenPair\_threeHundredTwenty} & Gaps312320 & PROVED & \checkmark & verified \\
\texttt{localFactor\_threeHundredTwelve\_two} & Gaps312320 & PROVED & \checkmark & verified \\
\texttt{localFactor\_threeHundredTwenty\_two} & Gaps312320 & PROVED & \checkmark & verified \\
\texttt{nu\_p\_threeHundredEighteen} & Gaps312320 & PROVED & \checkmark & verified \\
\texttt{nu\_p\_threeHundredFourteen} & Gaps312320 & PROVED & \checkmark & verified \\
\texttt{nu\_p\_threeHundredSixteen} & Gaps312320 & PROVED & \checkmark & verified \\
\texttt{nu\_p\_threeHundredTwelve} & Gaps312320 & PROVED & \checkmark & verified \\
\texttt{nu\_p\_threeHundredTwelve\_two} & Gaps312320 & PROVED & \checkmark & verified \\
\texttt{nu\_p\_threeHundredTwenty} & Gaps312320 & PROVED & \checkmark & verified \\
\texttt{nu\_p\_threeHundredTwenty\_two} & Gaps312320 & PROVED & \checkmark & verified \\
\texttt{singular\_series\_finite\_pos\_evenPair\_threeHundredEighteen} & Gaps312320 & PROVED & \checkmark & verified \\
\texttt{singular\_series\_finite\_pos\_evenPair\_threeHundredFourteen} & Gaps312320 & PROVED & \checkmark & verified \\
\texttt{singular\_series\_finite\_pos\_evenPair\_threeHundredSixteen} & Gaps312320 & PROVED & \checkmark & verified \\
\texttt{singular\_series\_finite\_pos\_evenPair\_threeHundredTwelve} & Gaps312320 & PROVED & \checkmark & verified \\
\texttt{singular\_series\_finite\_pos\_evenPair\_threeHundredTwenty} & Gaps312320 & PROVED & \checkmark & verified \\
\texttt{singular\_series\_pos\_evenPair\_threeHundredEighteen} & Gaps312320 & PROVED & \checkmark & verified \\
\texttt{singular\_series\_pos\_evenPair\_threeHundredFourteen} & Gaps312320 & PROVED & \checkmark & verified \\
\texttt{singular\_series\_pos\_evenPair\_threeHundredSixteen} & Gaps312320 & PROVED & \checkmark & verified \\
\texttt{singular\_series\_pos\_evenPair\_threeHundredTwelve} & Gaps312320 & PROVED & \checkmark & verified \\
\texttt{singular\_series\_pos\_evenPair\_threeHundredTwenty} & Gaps312320 & PROVED & \checkmark & verified \\
\texttt{evenPair\_card\_threeHundredThirty} & Gaps322330 & PROVED & \checkmark & verified \\
\texttt{evenPair\_card\_threeHundredTwentyEight} & Gaps322330 & PROVED & \checkmark & verified \\
\texttt{evenPair\_card\_threeHundredTwentyFour} & Gaps322330 & PROVED & \checkmark & verified \\
\texttt{evenPair\_card\_threeHundredTwentySix} & Gaps322330 & PROVED & \checkmark & verified \\
\texttt{evenPair\_card\_threeHundredTwentyTwo} & Gaps322330 & PROVED & \checkmark & verified \\
\texttt{isAdmissible\_evenPair\_threeHundredThirty} & Gaps322330 & PROVED & \checkmark & verified \\
\texttt{isAdmissible\_evenPair\_threeHundredTwentyEight} & Gaps322330 & PROVED & \checkmark & verified \\
\texttt{isAdmissible\_evenPair\_threeHundredTwentyFour} & Gaps322330 & PROVED & \checkmark & verified \\
\texttt{isAdmissible\_evenPair\_threeHundredTwentySix} & Gaps322330 & PROVED & \checkmark & verified \\
\texttt{isAdmissible\_evenPair\_threeHundredTwentyTwo} & Gaps322330 & PROVED & \checkmark & verified \\
\texttt{localFactor\_threeHundredThirty\_two} & Gaps322330 & PROVED & \checkmark & verified \\
\texttt{localFactor\_threeHundredTwentyTwo\_two} & Gaps322330 & PROVED & \checkmark & verified \\
\texttt{nu\_p\_threeHundredThirty} & Gaps322330 & PROVED & \checkmark & verified \\
\texttt{nu\_p\_threeHundredThirty\_two} & Gaps322330 & PROVED & \checkmark & verified \\
\texttt{nu\_p\_threeHundredTwentyEight} & Gaps322330 & PROVED & \checkmark & verified \\
\texttt{nu\_p\_threeHundredTwentyFour} & Gaps322330 & PROVED & \checkmark & verified \\
\texttt{nu\_p\_threeHundredTwentySix} & Gaps322330 & PROVED & \checkmark & verified \\
\texttt{nu\_p\_threeHundredTwentyTwo} & Gaps322330 & PROVED & \checkmark & verified \\
\texttt{nu\_p\_threeHundredTwentyTwo\_two} & Gaps322330 & PROVED & \checkmark & verified \\
\texttt{singular\_series\_finite\_pos\_evenPair\_threeHundredThirty} & Gaps322330 & PROVED & \checkmark & verified \\
\texttt{singular\_series\_finite\_pos\_evenPair\_threeHundredTwentyEight} & Gaps322330 & PROVED & \checkmark & verified \\
\texttt{singular\_series\_finite\_pos\_evenPair\_threeHundredTwentyFour} & Gaps322330 & PROVED & \checkmark & verified \\
\texttt{singular\_series\_finite\_pos\_evenPair\_threeHundredTwentySix} & Gaps322330 & PROVED & \checkmark & verified \\
\texttt{singular\_series\_finite\_pos\_evenPair\_threeHundredTwentyTwo} & Gaps322330 & PROVED & \checkmark & verified \\
\texttt{singular\_series\_pos\_evenPair\_threeHundredThirty} & Gaps322330 & PROVED & \checkmark & verified \\
\texttt{singular\_series\_pos\_evenPair\_threeHundredTwentyEight} & Gaps322330 & PROVED & \checkmark & verified \\
\texttt{singular\_series\_pos\_evenPair\_threeHundredTwentyFour} & Gaps322330 & PROVED & \checkmark & verified \\
\texttt{singular\_series\_pos\_evenPair\_threeHundredTwentySix} & Gaps322330 & PROVED & \checkmark & verified \\
\texttt{singular\_series\_pos\_evenPair\_threeHundredTwentyTwo} & Gaps322330 & PROVED & \checkmark & verified \\
\texttt{evenPair\_card\_forty} & Gaps3240 & PROVED & \checkmark & verified \\
\texttt{evenPair\_card\_thirtyEight} & Gaps3240 & PROVED & \checkmark & verified \\
\texttt{evenPair\_card\_thirtyFour} & Gaps3240 & PROVED & \checkmark & verified \\
\texttt{evenPair\_card\_thirtySix} & Gaps3240 & PROVED & \checkmark & verified \\
\texttt{evenPair\_card\_thirtyTwo} & Gaps3240 & PROVED & \checkmark & verified \\
\texttt{isAdmissible\_evenPair\_forty} & Gaps3240 & PROVED & \checkmark & verified \\
\texttt{isAdmissible\_evenPair\_thirtyEight} & Gaps3240 & PROVED & \checkmark & verified \\
\texttt{isAdmissible\_evenPair\_thirtyFour} & Gaps3240 & PROVED & \checkmark & verified \\
\texttt{isAdmissible\_evenPair\_thirtySix} & Gaps3240 & PROVED & \checkmark & verified \\
\texttt{isAdmissible\_evenPair\_thirtyTwo} & Gaps3240 & PROVED & \checkmark & verified \\
\texttt{localFactor\_forty\_two} & Gaps3240 & PROVED & \checkmark & verified \\
\texttt{localFactor\_thirtyTwo\_two} & Gaps3240 & PROVED & \checkmark & verified \\
\texttt{nu\_p\_forty} & Gaps3240 & PROVED & \checkmark & verified \\
\texttt{nu\_p\_forty\_two} & Gaps3240 & PROVED & \checkmark & verified \\
\texttt{nu\_p\_thirtyEight} & Gaps3240 & PROVED & \checkmark & verified \\
\texttt{nu\_p\_thirtyFour} & Gaps3240 & PROVED & \checkmark & verified \\
\texttt{nu\_p\_thirtySix} & Gaps3240 & PROVED & \checkmark & verified \\
\texttt{nu\_p\_thirtyTwo} & Gaps3240 & PROVED & \checkmark & verified \\
\texttt{nu\_p\_thirtyTwo\_two} & Gaps3240 & PROVED & \checkmark & verified \\
\texttt{singular\_series\_finite\_pos\_evenPair\_forty} & Gaps3240 & PROVED & \checkmark & verified \\
\texttt{singular\_series\_finite\_pos\_evenPair\_thirtyEight} & Gaps3240 & PROVED & \checkmark & verified \\
\texttt{singular\_series\_finite\_pos\_evenPair\_thirtyFour} & Gaps3240 & PROVED & \checkmark & verified \\
\texttt{singular\_series\_finite\_pos\_evenPair\_thirtySix} & Gaps3240 & PROVED & \checkmark & verified \\
\texttt{singular\_series\_finite\_pos\_evenPair\_thirtyTwo} & Gaps3240 & PROVED & \checkmark & verified \\
\texttt{singular\_series\_pos\_evenPair\_forty} & Gaps3240 & PROVED & \checkmark & verified \\
\texttt{singular\_series\_pos\_evenPair\_thirtyEight} & Gaps3240 & PROVED & \checkmark & verified \\
\texttt{singular\_series\_pos\_evenPair\_thirtyFour} & Gaps3240 & PROVED & \checkmark & verified \\
\texttt{singular\_series\_pos\_evenPair\_thirtySix} & Gaps3240 & PROVED & \checkmark & verified \\
\texttt{singular\_series\_pos\_evenPair\_thirtyTwo} & Gaps3240 & PROVED & \checkmark & verified \\
\texttt{evenPair\_card\_threeHundredForty} & Gaps332340 & PROVED & \checkmark & verified \\
\texttt{evenPair\_card\_threeHundredThirtyEight} & Gaps332340 & PROVED & \checkmark & verified \\
\texttt{evenPair\_card\_threeHundredThirtyFour} & Gaps332340 & PROVED & \checkmark & verified \\
\texttt{evenPair\_card\_threeHundredThirtySix} & Gaps332340 & PROVED & \checkmark & verified \\
\texttt{evenPair\_card\_threeHundredThirtyTwo} & Gaps332340 & PROVED & \checkmark & verified \\
\texttt{isAdmissible\_evenPair\_threeHundredForty} & Gaps332340 & PROVED & \checkmark & verified \\
\texttt{isAdmissible\_evenPair\_threeHundredThirtyEight} & Gaps332340 & PROVED & \checkmark & verified \\
\texttt{isAdmissible\_evenPair\_threeHundredThirtyFour} & Gaps332340 & PROVED & \checkmark & verified \\
\texttt{isAdmissible\_evenPair\_threeHundredThirtySix} & Gaps332340 & PROVED & \checkmark & verified \\
\texttt{isAdmissible\_evenPair\_threeHundredThirtyTwo} & Gaps332340 & PROVED & \checkmark & verified \\
\texttt{localFactor\_threeHundredForty\_two} & Gaps332340 & PROVED & \checkmark & verified \\
\texttt{localFactor\_threeHundredThirtyTwo\_two} & Gaps332340 & PROVED & \checkmark & verified \\
\texttt{nu\_p\_threeHundredForty} & Gaps332340 & PROVED & \checkmark & verified \\
\texttt{nu\_p\_threeHundredForty\_two} & Gaps332340 & PROVED & \checkmark & verified \\
\texttt{nu\_p\_threeHundredThirtyEight} & Gaps332340 & PROVED & \checkmark & verified \\
\texttt{nu\_p\_threeHundredThirtyFour} & Gaps332340 & PROVED & \checkmark & verified \\
\texttt{nu\_p\_threeHundredThirtySix} & Gaps332340 & PROVED & \checkmark & verified \\
\texttt{nu\_p\_threeHundredThirtyTwo} & Gaps332340 & PROVED & \checkmark & verified \\
\texttt{nu\_p\_threeHundredThirtyTwo\_two} & Gaps332340 & PROVED & \checkmark & verified \\
\texttt{singular\_series\_finite\_pos\_evenPair\_threeHundredForty} & Gaps332340 & PROVED & \checkmark & verified \\
\texttt{singular\_series\_finite\_pos\_evenPair\_threeHundredThirtyEight} & Gaps332340 & PROVED & \checkmark & verified \\
\texttt{singular\_series\_finite\_pos\_evenPair\_threeHundredThirtyFour} & Gaps332340 & PROVED & \checkmark & verified \\
\texttt{singular\_series\_finite\_pos\_evenPair\_threeHundredThirtySix} & Gaps332340 & PROVED & \checkmark & verified \\
\texttt{singular\_series\_finite\_pos\_evenPair\_threeHundredThirtyTwo} & Gaps332340 & PROVED & \checkmark & verified \\
\texttt{singular\_series\_pos\_evenPair\_threeHundredForty} & Gaps332340 & PROVED & \checkmark & verified \\
\texttt{singular\_series\_pos\_evenPair\_threeHundredThirtyEight} & Gaps332340 & PROVED & \checkmark & verified \\
\texttt{singular\_series\_pos\_evenPair\_threeHundredThirtyFour} & Gaps332340 & PROVED & \checkmark & verified \\
\texttt{singular\_series\_pos\_evenPair\_threeHundredThirtySix} & Gaps332340 & PROVED & \checkmark & verified \\
\texttt{singular\_series\_pos\_evenPair\_threeHundredThirtyTwo} & Gaps332340 & PROVED & \checkmark & verified \\
\texttt{evenPair\_card\_threeHundredFifty} & Gaps342350 & PROVED & \checkmark & verified \\
\texttt{evenPair\_card\_threeHundredFortyEight} & Gaps342350 & PROVED & \checkmark & verified \\
\texttt{evenPair\_card\_threeHundredFortyFour} & Gaps342350 & PROVED & \checkmark & verified \\
\texttt{evenPair\_card\_threeHundredFortySix} & Gaps342350 & PROVED & \checkmark & verified \\
\texttt{evenPair\_card\_threeHundredFortyTwo} & Gaps342350 & PROVED & \checkmark & verified \\
\texttt{isAdmissible\_evenPair\_threeHundredFifty} & Gaps342350 & PROVED & \checkmark & verified \\
\texttt{isAdmissible\_evenPair\_threeHundredFortyEight} & Gaps342350 & PROVED & \checkmark & verified \\
\texttt{isAdmissible\_evenPair\_threeHundredFortyFour} & Gaps342350 & PROVED & \checkmark & verified \\
\texttt{isAdmissible\_evenPair\_threeHundredFortySix} & Gaps342350 & PROVED & \checkmark & verified \\
\texttt{isAdmissible\_evenPair\_threeHundredFortyTwo} & Gaps342350 & PROVED & \checkmark & verified \\
\texttt{localFactor\_threeHundredFifty\_two} & Gaps342350 & PROVED & \checkmark & verified \\
\texttt{localFactor\_threeHundredFortyTwo\_two} & Gaps342350 & PROVED & \checkmark & verified \\
\texttt{nu\_p\_threeHundredFifty} & Gaps342350 & PROVED & \checkmark & verified \\
\texttt{nu\_p\_threeHundredFifty\_two} & Gaps342350 & PROVED & \checkmark & verified \\
\texttt{nu\_p\_threeHundredFortyEight} & Gaps342350 & PROVED & \checkmark & verified \\
\texttt{nu\_p\_threeHundredFortyFour} & Gaps342350 & PROVED & \checkmark & verified \\
\texttt{nu\_p\_threeHundredFortySix} & Gaps342350 & PROVED & \checkmark & verified \\
\texttt{nu\_p\_threeHundredFortyTwo} & Gaps342350 & PROVED & \checkmark & verified \\
\texttt{nu\_p\_threeHundredFortyTwo\_two} & Gaps342350 & PROVED & \checkmark & verified \\
\texttt{singular\_series\_finite\_pos\_evenPair\_threeHundredFifty} & Gaps342350 & PROVED & \checkmark & verified \\
\texttt{singular\_series\_finite\_pos\_evenPair\_threeHundredFortyEight} & Gaps342350 & PROVED & \checkmark & verified \\
\texttt{singular\_series\_finite\_pos\_evenPair\_threeHundredFortyFour} & Gaps342350 & PROVED & \checkmark & verified \\
\texttt{singular\_series\_finite\_pos\_evenPair\_threeHundredFortySix} & Gaps342350 & PROVED & \checkmark & verified \\
\texttt{singular\_series\_finite\_pos\_evenPair\_threeHundredFortyTwo} & Gaps342350 & PROVED & \checkmark & verified \\
\texttt{singular\_series\_pos\_evenPair\_threeHundredFifty} & Gaps342350 & PROVED & \checkmark & verified \\
\texttt{singular\_series\_pos\_evenPair\_threeHundredFortyEight} & Gaps342350 & PROVED & \checkmark & verified \\
\texttt{singular\_series\_pos\_evenPair\_threeHundredFortyFour} & Gaps342350 & PROVED & \checkmark & verified \\
\texttt{singular\_series\_pos\_evenPair\_threeHundredFortySix} & Gaps342350 & PROVED & \checkmark & verified \\
\texttt{singular\_series\_pos\_evenPair\_threeHundredFortyTwo} & Gaps342350 & PROVED & \checkmark & verified \\
\texttt{evenPair\_card\_threeHundredFiftyEight} & Gaps352360 & PROVED & \checkmark & verified \\
\texttt{evenPair\_card\_threeHundredFiftyFour} & Gaps352360 & PROVED & \checkmark & verified \\
\texttt{evenPair\_card\_threeHundredFiftySix} & Gaps352360 & PROVED & \checkmark & verified \\
\texttt{evenPair\_card\_threeHundredFiftyTwo} & Gaps352360 & PROVED & \checkmark & verified \\
\texttt{evenPair\_card\_threeHundredSixty} & Gaps352360 & PROVED & \checkmark & verified \\
\texttt{isAdmissible\_evenPair\_threeHundredFiftyEight} & Gaps352360 & PROVED & \checkmark & verified \\
\texttt{isAdmissible\_evenPair\_threeHundredFiftyFour} & Gaps352360 & PROVED & \checkmark & verified \\
\texttt{isAdmissible\_evenPair\_threeHundredFiftySix} & Gaps352360 & PROVED & \checkmark & verified \\
\texttt{isAdmissible\_evenPair\_threeHundredFiftyTwo} & Gaps352360 & PROVED & \checkmark & verified \\
\texttt{isAdmissible\_evenPair\_threeHundredSixty} & Gaps352360 & PROVED & \checkmark & verified \\
\texttt{localFactor\_threeHundredFiftyTwo\_two} & Gaps352360 & PROVED & \checkmark & verified \\
\texttt{localFactor\_threeHundredSixty\_two} & Gaps352360 & PROVED & \checkmark & verified \\
\texttt{nu\_p\_threeHundredFiftyEight} & Gaps352360 & PROVED & \checkmark & verified \\
\texttt{nu\_p\_threeHundredFiftyFour} & Gaps352360 & PROVED & \checkmark & verified \\
\texttt{nu\_p\_threeHundredFiftySix} & Gaps352360 & PROVED & \checkmark & verified \\
\texttt{nu\_p\_threeHundredFiftyTwo} & Gaps352360 & PROVED & \checkmark & verified \\
\texttt{nu\_p\_threeHundredFiftyTwo\_two} & Gaps352360 & PROVED & \checkmark & verified \\
\texttt{nu\_p\_threeHundredSixty} & Gaps352360 & PROVED & \checkmark & verified \\
\texttt{nu\_p\_threeHundredSixty\_two} & Gaps352360 & PROVED & \checkmark & verified \\
\texttt{singular\_series\_finite\_pos\_evenPair\_threeHundredFiftyEight} & Gaps352360 & PROVED & \checkmark & verified \\
\texttt{singular\_series\_finite\_pos\_evenPair\_threeHundredFiftyFour} & Gaps352360 & PROVED & \checkmark & verified \\
\texttt{singular\_series\_finite\_pos\_evenPair\_threeHundredFiftySix} & Gaps352360 & PROVED & \checkmark & verified \\
\texttt{singular\_series\_finite\_pos\_evenPair\_threeHundredFiftyTwo} & Gaps352360 & PROVED & \checkmark & verified \\
\texttt{singular\_series\_finite\_pos\_evenPair\_threeHundredSixty} & Gaps352360 & PROVED & \checkmark & verified \\
\texttt{singular\_series\_pos\_evenPair\_threeHundredFiftyEight} & Gaps352360 & PROVED & \checkmark & verified \\
\texttt{singular\_series\_pos\_evenPair\_threeHundredFiftyFour} & Gaps352360 & PROVED & \checkmark & verified \\
\texttt{singular\_series\_pos\_evenPair\_threeHundredFiftySix} & Gaps352360 & PROVED & \checkmark & verified \\
\texttt{singular\_series\_pos\_evenPair\_threeHundredFiftyTwo} & Gaps352360 & PROVED & \checkmark & verified \\
\texttt{singular\_series\_pos\_evenPair\_threeHundredSixty} & Gaps352360 & PROVED & \checkmark & verified \\
\texttt{evenPair\_card\_threeHundredSeventy} & Gaps362370 & PROVED & \checkmark & verified \\
\texttt{evenPair\_card\_threeHundredSixtyEight} & Gaps362370 & PROVED & \checkmark & verified \\
\texttt{evenPair\_card\_threeHundredSixtyFour} & Gaps362370 & PROVED & \checkmark & verified \\
\texttt{evenPair\_card\_threeHundredSixtySix} & Gaps362370 & PROVED & \checkmark & verified \\
\texttt{evenPair\_card\_threeHundredSixtyTwo} & Gaps362370 & PROVED & \checkmark & verified \\
\texttt{isAdmissible\_evenPair\_threeHundredSeventy} & Gaps362370 & PROVED & \checkmark & verified \\
\texttt{isAdmissible\_evenPair\_threeHundredSixtyEight} & Gaps362370 & PROVED & \checkmark & verified \\
\texttt{isAdmissible\_evenPair\_threeHundredSixtyFour} & Gaps362370 & PROVED & \checkmark & verified \\
\texttt{isAdmissible\_evenPair\_threeHundredSixtySix} & Gaps362370 & PROVED & \checkmark & verified \\
\texttt{isAdmissible\_evenPair\_threeHundredSixtyTwo} & Gaps362370 & PROVED & \checkmark & verified \\
\texttt{localFactor\_threeHundredSeventy\_two} & Gaps362370 & PROVED & \checkmark & verified \\
\texttt{localFactor\_threeHundredSixtyTwo\_two} & Gaps362370 & PROVED & \checkmark & verified \\
\texttt{nu\_p\_threeHundredSeventy} & Gaps362370 & PROVED & \checkmark & verified \\
\texttt{nu\_p\_threeHundredSeventy\_two} & Gaps362370 & PROVED & \checkmark & verified \\
\texttt{nu\_p\_threeHundredSixtyEight} & Gaps362370 & PROVED & \checkmark & verified \\
\texttt{nu\_p\_threeHundredSixtyFour} & Gaps362370 & PROVED & \checkmark & verified \\
\texttt{nu\_p\_threeHundredSixtySix} & Gaps362370 & PROVED & \checkmark & verified \\
\texttt{nu\_p\_threeHundredSixtyTwo} & Gaps362370 & PROVED & \checkmark & verified \\
\texttt{nu\_p\_threeHundredSixtyTwo\_two} & Gaps362370 & PROVED & \checkmark & verified \\
\texttt{singular\_series\_finite\_pos\_evenPair\_threeHundredSeventy} & Gaps362370 & PROVED & \checkmark & verified \\
\texttt{singular\_series\_finite\_pos\_evenPair\_threeHundredSixtyEight} & Gaps362370 & PROVED & \checkmark & verified \\
\texttt{singular\_series\_finite\_pos\_evenPair\_threeHundredSixtyFour} & Gaps362370 & PROVED & \checkmark & verified \\
\texttt{singular\_series\_finite\_pos\_evenPair\_threeHundredSixtySix} & Gaps362370 & PROVED & \checkmark & verified \\
\texttt{singular\_series\_finite\_pos\_evenPair\_threeHundredSixtyTwo} & Gaps362370 & PROVED & \checkmark & verified \\
\texttt{singular\_series\_pos\_evenPair\_threeHundredSeventy} & Gaps362370 & PROVED & \checkmark & verified \\
\texttt{singular\_series\_pos\_evenPair\_threeHundredSixtyEight} & Gaps362370 & PROVED & \checkmark & verified \\
\texttt{singular\_series\_pos\_evenPair\_threeHundredSixtyFour} & Gaps362370 & PROVED & \checkmark & verified \\
\texttt{singular\_series\_pos\_evenPair\_threeHundredSixtySix} & Gaps362370 & PROVED & \checkmark & verified \\
\texttt{singular\_series\_pos\_evenPair\_threeHundredSixtyTwo} & Gaps362370 & PROVED & \checkmark & verified \\
\texttt{evenPair\_card\_threeHundredEighty} & Gaps372380 & PROVED & \checkmark & verified \\
\texttt{evenPair\_card\_threeHundredSeventyEight} & Gaps372380 & PROVED & \checkmark & verified \\
\texttt{evenPair\_card\_threeHundredSeventyFour} & Gaps372380 & PROVED & \checkmark & verified \\
\texttt{evenPair\_card\_threeHundredSeventySix} & Gaps372380 & PROVED & \checkmark & verified \\
\texttt{evenPair\_card\_threeHundredSeventyTwo} & Gaps372380 & PROVED & \checkmark & verified \\
\texttt{isAdmissible\_evenPair\_threeHundredEighty} & Gaps372380 & PROVED & \checkmark & verified \\
\texttt{isAdmissible\_evenPair\_threeHundredSeventyEight} & Gaps372380 & PROVED & \checkmark & verified \\
\texttt{isAdmissible\_evenPair\_threeHundredSeventyFour} & Gaps372380 & PROVED & \checkmark & verified \\
\texttt{isAdmissible\_evenPair\_threeHundredSeventySix} & Gaps372380 & PROVED & \checkmark & verified \\
\texttt{isAdmissible\_evenPair\_threeHundredSeventyTwo} & Gaps372380 & PROVED & \checkmark & verified \\
\texttt{localFactor\_threeHundredEighty\_two} & Gaps372380 & PROVED & \checkmark & verified \\
\texttt{localFactor\_threeHundredSeventyTwo\_two} & Gaps372380 & PROVED & \checkmark & verified \\
\texttt{nu\_p\_threeHundredEighty} & Gaps372380 & PROVED & \checkmark & verified \\
\texttt{nu\_p\_threeHundredEighty\_two} & Gaps372380 & PROVED & \checkmark & verified \\
\texttt{nu\_p\_threeHundredSeventyEight} & Gaps372380 & PROVED & \checkmark & verified \\
\texttt{nu\_p\_threeHundredSeventyFour} & Gaps372380 & PROVED & \checkmark & verified \\
\texttt{nu\_p\_threeHundredSeventySix} & Gaps372380 & PROVED & \checkmark & verified \\
\texttt{nu\_p\_threeHundredSeventyTwo} & Gaps372380 & PROVED & \checkmark & verified \\
\texttt{nu\_p\_threeHundredSeventyTwo\_two} & Gaps372380 & PROVED & \checkmark & verified \\
\texttt{singular\_series\_finite\_pos\_evenPair\_threeHundredEighty} & Gaps372380 & PROVED & \checkmark & verified \\
\texttt{singular\_series\_finite\_pos\_evenPair\_threeHundredSeventyEight} & Gaps372380 & PROVED & \checkmark & verified \\
\texttt{singular\_series\_finite\_pos\_evenPair\_threeHundredSeventyFour} & Gaps372380 & PROVED & \checkmark & verified \\
\texttt{singular\_series\_finite\_pos\_evenPair\_threeHundredSeventySix} & Gaps372380 & PROVED & \checkmark & verified \\
\texttt{singular\_series\_finite\_pos\_evenPair\_threeHundredSeventyTwo} & Gaps372380 & PROVED & \checkmark & verified \\
\texttt{singular\_series\_pos\_evenPair\_threeHundredEighty} & Gaps372380 & PROVED & \checkmark & verified \\
\texttt{singular\_series\_pos\_evenPair\_threeHundredSeventyEight} & Gaps372380 & PROVED & \checkmark & verified \\
\texttt{singular\_series\_pos\_evenPair\_threeHundredSeventyFour} & Gaps372380 & PROVED & \checkmark & verified \\
\texttt{singular\_series\_pos\_evenPair\_threeHundredSeventySix} & Gaps372380 & PROVED & \checkmark & verified \\
\texttt{singular\_series\_pos\_evenPair\_threeHundredSeventyTwo} & Gaps372380 & PROVED & \checkmark & verified \\
\texttt{evenPair\_card\_threeHundredEightyEight} & Gaps382390 & PROVED & \checkmark & verified \\
\texttt{evenPair\_card\_threeHundredEightyFour} & Gaps382390 & PROVED & \checkmark & verified \\
\texttt{evenPair\_card\_threeHundredEightySix} & Gaps382390 & PROVED & \checkmark & verified \\
\texttt{evenPair\_card\_threeHundredEightyTwo} & Gaps382390 & PROVED & \checkmark & verified \\
\texttt{evenPair\_card\_threeHundredNinety} & Gaps382390 & PROVED & \checkmark & verified \\
\texttt{isAdmissible\_evenPair\_threeHundredEightyEight} & Gaps382390 & PROVED & \checkmark & verified \\
\texttt{isAdmissible\_evenPair\_threeHundredEightyFour} & Gaps382390 & PROVED & \checkmark & verified \\
\texttt{isAdmissible\_evenPair\_threeHundredEightySix} & Gaps382390 & PROVED & \checkmark & verified \\
\texttt{isAdmissible\_evenPair\_threeHundredEightyTwo} & Gaps382390 & PROVED & \checkmark & verified \\
\texttt{isAdmissible\_evenPair\_threeHundredNinety} & Gaps382390 & PROVED & \checkmark & verified \\
\texttt{localFactor\_threeHundredEightyTwo\_two} & Gaps382390 & PROVED & \checkmark & verified \\
\texttt{localFactor\_threeHundredNinety\_two} & Gaps382390 & PROVED & \checkmark & verified \\
\texttt{nu\_p\_threeHundredEightyEight} & Gaps382390 & PROVED & \checkmark & verified \\
\texttt{nu\_p\_threeHundredEightyFour} & Gaps382390 & PROVED & \checkmark & verified \\
\texttt{nu\_p\_threeHundredEightySix} & Gaps382390 & PROVED & \checkmark & verified \\
\texttt{nu\_p\_threeHundredEightyTwo} & Gaps382390 & PROVED & \checkmark & verified \\
\texttt{nu\_p\_threeHundredEightyTwo\_two} & Gaps382390 & PROVED & \checkmark & verified \\
\texttt{nu\_p\_threeHundredNinety} & Gaps382390 & PROVED & \checkmark & verified \\
\texttt{nu\_p\_threeHundredNinety\_two} & Gaps382390 & PROVED & \checkmark & verified \\
\texttt{singular\_series\_finite\_pos\_evenPair\_threeHundredEightyEight} & Gaps382390 & PROVED & \checkmark & verified \\
\texttt{singular\_series\_finite\_pos\_evenPair\_threeHundredEightyFour} & Gaps382390 & PROVED & \checkmark & verified \\
\texttt{singular\_series\_finite\_pos\_evenPair\_threeHundredEightySix} & Gaps382390 & PROVED & \checkmark & verified \\
\texttt{singular\_series\_finite\_pos\_evenPair\_threeHundredEightyTwo} & Gaps382390 & PROVED & \checkmark & verified \\
\texttt{singular\_series\_finite\_pos\_evenPair\_threeHundredNinety} & Gaps382390 & PROVED & \checkmark & verified \\
\texttt{singular\_series\_pos\_evenPair\_threeHundredEightyEight} & Gaps382390 & PROVED & \checkmark & verified \\
\texttt{singular\_series\_pos\_evenPair\_threeHundredEightyFour} & Gaps382390 & PROVED & \checkmark & verified \\
\texttt{singular\_series\_pos\_evenPair\_threeHundredEightySix} & Gaps382390 & PROVED & \checkmark & verified \\
\texttt{singular\_series\_pos\_evenPair\_threeHundredEightyTwo} & Gaps382390 & PROVED & \checkmark & verified \\
\texttt{singular\_series\_pos\_evenPair\_threeHundredNinety} & Gaps382390 & PROVED & \checkmark & verified \\
\texttt{evenPair\_card\_fourHundred} & Gaps392400 & PROVED & \checkmark & verified \\
\texttt{evenPair\_card\_threeHundredNinetyEight} & Gaps392400 & PROVED & \checkmark & verified \\
\texttt{evenPair\_card\_threeHundredNinetyFour} & Gaps392400 & PROVED & \checkmark & verified \\
\texttt{evenPair\_card\_threeHundredNinetySix} & Gaps392400 & PROVED & \checkmark & verified \\
\texttt{evenPair\_card\_threeHundredNinetyTwo} & Gaps392400 & PROVED & \checkmark & verified \\
\texttt{isAdmissible\_evenPair\_fourHundred} & Gaps392400 & PROVED & \checkmark & verified \\
\texttt{isAdmissible\_evenPair\_threeHundredNinetyEight} & Gaps392400 & PROVED & \checkmark & verified \\
\texttt{isAdmissible\_evenPair\_threeHundredNinetyFour} & Gaps392400 & PROVED & \checkmark & verified \\
\texttt{isAdmissible\_evenPair\_threeHundredNinetySix} & Gaps392400 & PROVED & \checkmark & verified \\
\texttt{isAdmissible\_evenPair\_threeHundredNinetyTwo} & Gaps392400 & PROVED & \checkmark & verified \\
\texttt{localFactor\_fourHundred\_two} & Gaps392400 & PROVED & \checkmark & verified \\
\texttt{localFactor\_threeHundredNinetyTwo\_two} & Gaps392400 & PROVED & \checkmark & verified \\
\texttt{nu\_p\_fourHundred} & Gaps392400 & PROVED & \checkmark & verified \\
\texttt{nu\_p\_fourHundred\_two} & Gaps392400 & PROVED & \checkmark & verified \\
\texttt{nu\_p\_threeHundredNinetyEight} & Gaps392400 & PROVED & \checkmark & verified \\
\texttt{nu\_p\_threeHundredNinetyFour} & Gaps392400 & PROVED & \checkmark & verified \\
\texttt{nu\_p\_threeHundredNinetySix} & Gaps392400 & PROVED & \checkmark & verified \\
\texttt{nu\_p\_threeHundredNinetyTwo} & Gaps392400 & PROVED & \checkmark & verified \\
\texttt{nu\_p\_threeHundredNinetyTwo\_two} & Gaps392400 & PROVED & \checkmark & verified \\
\texttt{singular\_series\_finite\_pos\_evenPair\_fourHundred} & Gaps392400 & PROVED & \checkmark & verified \\
\texttt{singular\_series\_finite\_pos\_evenPair\_threeHundredNinetyEight} & Gaps392400 & PROVED & \checkmark & verified \\
\texttt{singular\_series\_finite\_pos\_evenPair\_threeHundredNinetyFour} & Gaps392400 & PROVED & \checkmark & verified \\
\texttt{singular\_series\_finite\_pos\_evenPair\_threeHundredNinetySix} & Gaps392400 & PROVED & \checkmark & verified \\
\texttt{singular\_series\_finite\_pos\_evenPair\_threeHundredNinetyTwo} & Gaps392400 & PROVED & \checkmark & verified \\
\texttt{singular\_series\_pos\_evenPair\_fourHundred} & Gaps392400 & PROVED & \checkmark & verified \\
\texttt{singular\_series\_pos\_evenPair\_threeHundredNinetyEight} & Gaps392400 & PROVED & \checkmark & verified \\
\texttt{singular\_series\_pos\_evenPair\_threeHundredNinetyFour} & Gaps392400 & PROVED & \checkmark & verified \\
\texttt{singular\_series\_pos\_evenPair\_threeHundredNinetySix} & Gaps392400 & PROVED & \checkmark & verified \\
\texttt{singular\_series\_pos\_evenPair\_threeHundredNinetyTwo} & Gaps392400 & PROVED & \checkmark & verified \\
\texttt{evenPair\_card\_fourHundredEight} & Gaps402410 & PROVED & \checkmark & verified \\
\texttt{evenPair\_card\_fourHundredFour} & Gaps402410 & PROVED & \checkmark & verified \\
\texttt{evenPair\_card\_fourHundredSix} & Gaps402410 & PROVED & \checkmark & verified \\
\texttt{evenPair\_card\_fourHundredTen} & Gaps402410 & PROVED & \checkmark & verified \\
\texttt{evenPair\_card\_fourHundredTwo} & Gaps402410 & PROVED & \checkmark & verified \\
\texttt{isAdmissible\_evenPair\_fourHundredEight} & Gaps402410 & PROVED & \checkmark & verified \\
\texttt{isAdmissible\_evenPair\_fourHundredFour} & Gaps402410 & PROVED & \checkmark & verified \\
\texttt{isAdmissible\_evenPair\_fourHundredSix} & Gaps402410 & PROVED & \checkmark & verified \\
\texttt{isAdmissible\_evenPair\_fourHundredTen} & Gaps402410 & PROVED & \checkmark & verified \\
\texttt{isAdmissible\_evenPair\_fourHundredTwo} & Gaps402410 & PROVED & \checkmark & verified \\
\texttt{localFactor\_fourHundredTen\_two} & Gaps402410 & PROVED & \checkmark & verified \\
\texttt{localFactor\_fourHundredTwo\_two} & Gaps402410 & PROVED & \checkmark & verified \\
\texttt{nu\_p\_fourHundredEight} & Gaps402410 & PROVED & \checkmark & verified \\
\texttt{nu\_p\_fourHundredFour} & Gaps402410 & PROVED & \checkmark & verified \\
\texttt{nu\_p\_fourHundredSix} & Gaps402410 & PROVED & \checkmark & verified \\
\texttt{nu\_p\_fourHundredTen} & Gaps402410 & PROVED & \checkmark & verified \\
\texttt{nu\_p\_fourHundredTen\_two} & Gaps402410 & PROVED & \checkmark & verified \\
\texttt{nu\_p\_fourHundredTwo} & Gaps402410 & PROVED & \checkmark & verified \\
\texttt{nu\_p\_fourHundredTwo\_two} & Gaps402410 & PROVED & \checkmark & verified \\
\texttt{singular\_series\_finite\_pos\_evenPair\_fourHundredEight} & Gaps402410 & PROVED & \checkmark & verified \\
\texttt{singular\_series\_finite\_pos\_evenPair\_fourHundredFour} & Gaps402410 & PROVED & \checkmark & verified \\
\texttt{singular\_series\_finite\_pos\_evenPair\_fourHundredSix} & Gaps402410 & PROVED & \checkmark & verified \\
\texttt{singular\_series\_finite\_pos\_evenPair\_fourHundredTen} & Gaps402410 & PROVED & \checkmark & verified \\
\texttt{singular\_series\_finite\_pos\_evenPair\_fourHundredTwo} & Gaps402410 & PROVED & \checkmark & verified \\
\texttt{singular\_series\_pos\_evenPair\_fourHundredEight} & Gaps402410 & PROVED & \checkmark & verified \\
\texttt{singular\_series\_pos\_evenPair\_fourHundredFour} & Gaps402410 & PROVED & \checkmark & verified \\
\texttt{singular\_series\_pos\_evenPair\_fourHundredSix} & Gaps402410 & PROVED & \checkmark & verified \\
\texttt{singular\_series\_pos\_evenPair\_fourHundredTen} & Gaps402410 & PROVED & \checkmark & verified \\
\texttt{singular\_series\_pos\_evenPair\_fourHundredTwo} & Gaps402410 & PROVED & \checkmark & verified \\
\texttt{evenPair\_card\_fourHundredEighteen} & Gaps412420 & PROVED & \checkmark & verified \\
\texttt{evenPair\_card\_fourHundredFourteen} & Gaps412420 & PROVED & \checkmark & verified \\
\texttt{evenPair\_card\_fourHundredSixteen} & Gaps412420 & PROVED & \checkmark & verified \\
\texttt{evenPair\_card\_fourHundredTwelve} & Gaps412420 & PROVED & \checkmark & verified \\
\texttt{evenPair\_card\_fourHundredTwenty} & Gaps412420 & PROVED & \checkmark & verified \\
\texttt{isAdmissible\_evenPair\_fourHundredEighteen} & Gaps412420 & PROVED & \checkmark & verified \\
\texttt{isAdmissible\_evenPair\_fourHundredFourteen} & Gaps412420 & PROVED & \checkmark & verified \\
\texttt{isAdmissible\_evenPair\_fourHundredSixteen} & Gaps412420 & PROVED & \checkmark & verified \\
\texttt{isAdmissible\_evenPair\_fourHundredTwelve} & Gaps412420 & PROVED & \checkmark & verified \\
\texttt{isAdmissible\_evenPair\_fourHundredTwenty} & Gaps412420 & PROVED & \checkmark & verified \\
\texttt{localFactor\_fourHundredTwelve\_two} & Gaps412420 & PROVED & \checkmark & verified \\
\texttt{localFactor\_fourHundredTwenty\_two} & Gaps412420 & PROVED & \checkmark & verified \\
\texttt{nu\_p\_fourHundredEighteen} & Gaps412420 & PROVED & \checkmark & verified \\
\texttt{nu\_p\_fourHundredFourteen} & Gaps412420 & PROVED & \checkmark & verified \\
\texttt{nu\_p\_fourHundredSixteen} & Gaps412420 & PROVED & \checkmark & verified \\
\texttt{nu\_p\_fourHundredTwelve} & Gaps412420 & PROVED & \checkmark & verified \\
\texttt{nu\_p\_fourHundredTwelve\_two} & Gaps412420 & PROVED & \checkmark & verified \\
\texttt{nu\_p\_fourHundredTwenty} & Gaps412420 & PROVED & \checkmark & verified \\
\texttt{nu\_p\_fourHundredTwenty\_two} & Gaps412420 & PROVED & \checkmark & verified \\
\texttt{singular\_series\_finite\_pos\_evenPair\_fourHundredEighteen} & Gaps412420 & PROVED & \checkmark & verified \\
\texttt{singular\_series\_finite\_pos\_evenPair\_fourHundredFourteen} & Gaps412420 & PROVED & \checkmark & verified \\
\texttt{singular\_series\_finite\_pos\_evenPair\_fourHundredSixteen} & Gaps412420 & PROVED & \checkmark & verified \\
\texttt{singular\_series\_finite\_pos\_evenPair\_fourHundredTwelve} & Gaps412420 & PROVED & \checkmark & verified \\
\texttt{singular\_series\_finite\_pos\_evenPair\_fourHundredTwenty} & Gaps412420 & PROVED & \checkmark & verified \\
\texttt{singular\_series\_pos\_evenPair\_fourHundredEighteen} & Gaps412420 & PROVED & \checkmark & verified \\
\texttt{singular\_series\_pos\_evenPair\_fourHundredFourteen} & Gaps412420 & PROVED & \checkmark & verified \\
\texttt{singular\_series\_pos\_evenPair\_fourHundredSixteen} & Gaps412420 & PROVED & \checkmark & verified \\
\texttt{singular\_series\_pos\_evenPair\_fourHundredTwelve} & Gaps412420 & PROVED & \checkmark & verified \\
\texttt{singular\_series\_pos\_evenPair\_fourHundredTwenty} & Gaps412420 & PROVED & \checkmark & verified \\
\texttt{evenPair\_card\_fourHundredThirty} & Gaps422430 & PROVED & \checkmark & verified \\
\texttt{evenPair\_card\_fourHundredTwentyEight} & Gaps422430 & PROVED & \checkmark & verified \\
\texttt{evenPair\_card\_fourHundredTwentyFour} & Gaps422430 & PROVED & \checkmark & verified \\
\texttt{evenPair\_card\_fourHundredTwentySix} & Gaps422430 & PROVED & \checkmark & verified \\
\texttt{evenPair\_card\_fourHundredTwentyTwo} & Gaps422430 & PROVED & \checkmark & verified \\
\texttt{isAdmissible\_evenPair\_fourHundredThirty} & Gaps422430 & PROVED & \checkmark & verified \\
\texttt{isAdmissible\_evenPair\_fourHundredTwentyEight} & Gaps422430 & PROVED & \checkmark & verified \\
\texttt{isAdmissible\_evenPair\_fourHundredTwentyFour} & Gaps422430 & PROVED & \checkmark & verified \\
\texttt{isAdmissible\_evenPair\_fourHundredTwentySix} & Gaps422430 & PROVED & \checkmark & verified \\
\texttt{isAdmissible\_evenPair\_fourHundredTwentyTwo} & Gaps422430 & PROVED & \checkmark & verified \\
\texttt{localFactor\_fourHundredThirty\_two} & Gaps422430 & PROVED & \checkmark & verified \\
\texttt{localFactor\_fourHundredTwentyTwo\_two} & Gaps422430 & PROVED & \checkmark & verified \\
\texttt{nu\_p\_fourHundredThirty} & Gaps422430 & PROVED & \checkmark & verified \\
\texttt{nu\_p\_fourHundredThirty\_two} & Gaps422430 & PROVED & \checkmark & verified \\
\texttt{nu\_p\_fourHundredTwentyEight} & Gaps422430 & PROVED & \checkmark & verified \\
\texttt{nu\_p\_fourHundredTwentyFour} & Gaps422430 & PROVED & \checkmark & verified \\
\texttt{nu\_p\_fourHundredTwentySix} & Gaps422430 & PROVED & \checkmark & verified \\
\texttt{nu\_p\_fourHundredTwentyTwo} & Gaps422430 & PROVED & \checkmark & verified \\
\texttt{nu\_p\_fourHundredTwentyTwo\_two} & Gaps422430 & PROVED & \checkmark & verified \\
\texttt{singular\_series\_finite\_pos\_evenPair\_fourHundredThirty} & Gaps422430 & PROVED & \checkmark & verified \\
\texttt{singular\_series\_finite\_pos\_evenPair\_fourHundredTwentyEight} & Gaps422430 & PROVED & \checkmark & verified \\
\texttt{singular\_series\_finite\_pos\_evenPair\_fourHundredTwentyFour} & Gaps422430 & PROVED & \checkmark & verified \\
\texttt{singular\_series\_finite\_pos\_evenPair\_fourHundredTwentySix} & Gaps422430 & PROVED & \checkmark & verified \\
\texttt{singular\_series\_finite\_pos\_evenPair\_fourHundredTwentyTwo} & Gaps422430 & PROVED & \checkmark & verified \\
\texttt{singular\_series\_pos\_evenPair\_fourHundredThirty} & Gaps422430 & PROVED & \checkmark & verified \\
\texttt{singular\_series\_pos\_evenPair\_fourHundredTwentyEight} & Gaps422430 & PROVED & \checkmark & verified \\
\texttt{singular\_series\_pos\_evenPair\_fourHundredTwentyFour} & Gaps422430 & PROVED & \checkmark & verified \\
\texttt{singular\_series\_pos\_evenPair\_fourHundredTwentySix} & Gaps422430 & PROVED & \checkmark & verified \\
\texttt{singular\_series\_pos\_evenPair\_fourHundredTwentyTwo} & Gaps422430 & PROVED & \checkmark & verified \\
\texttt{evenPair\_card\_fifty} & Gaps4250 & PROVED & \checkmark & verified \\
\texttt{evenPair\_card\_fortyEight} & Gaps4250 & PROVED & \checkmark & verified \\
\texttt{evenPair\_card\_fortyFour} & Gaps4250 & PROVED & \checkmark & verified \\
\texttt{evenPair\_card\_fortySix} & Gaps4250 & PROVED & \checkmark & verified \\
\texttt{evenPair\_card\_fortyTwo} & Gaps4250 & PROVED & \checkmark & verified \\
\texttt{isAdmissible\_evenPair\_fifty} & Gaps4250 & PROVED & \checkmark & verified \\
\texttt{isAdmissible\_evenPair\_fortyEight} & Gaps4250 & PROVED & \checkmark & verified \\
\texttt{isAdmissible\_evenPair\_fortyFour} & Gaps4250 & PROVED & \checkmark & verified \\
\texttt{isAdmissible\_evenPair\_fortySix} & Gaps4250 & PROVED & \checkmark & verified \\
\texttt{isAdmissible\_evenPair\_fortyTwo} & Gaps4250 & PROVED & \checkmark & verified \\
\texttt{localFactor\_fifty\_two} & Gaps4250 & PROVED & \checkmark & verified \\
\texttt{localFactor\_fortyTwo\_two} & Gaps4250 & PROVED & \checkmark & verified \\
\texttt{nu\_p\_fifty} & Gaps4250 & PROVED & \checkmark & verified \\
\texttt{nu\_p\_fifty\_two} & Gaps4250 & PROVED & \checkmark & verified \\
\texttt{nu\_p\_fortyEight} & Gaps4250 & PROVED & \checkmark & verified \\
\texttt{nu\_p\_fortyFour} & Gaps4250 & PROVED & \checkmark & verified \\
\texttt{nu\_p\_fortySix} & Gaps4250 & PROVED & \checkmark & verified \\
\texttt{nu\_p\_fortyTwo} & Gaps4250 & PROVED & \checkmark & verified \\
\texttt{nu\_p\_fortyTwo\_two} & Gaps4250 & PROVED & \checkmark & verified \\
\texttt{singular\_series\_finite\_pos\_evenPair\_fifty} & Gaps4250 & PROVED & \checkmark & verified \\
\texttt{singular\_series\_finite\_pos\_evenPair\_fortyEight} & Gaps4250 & PROVED & \checkmark & verified \\
\texttt{singular\_series\_finite\_pos\_evenPair\_fortyFour} & Gaps4250 & PROVED & \checkmark & verified \\
\texttt{singular\_series\_finite\_pos\_evenPair\_fortySix} & Gaps4250 & PROVED & \checkmark & verified \\
\texttt{singular\_series\_finite\_pos\_evenPair\_fortyTwo} & Gaps4250 & PROVED & \checkmark & verified \\
\texttt{singular\_series\_pos\_evenPair\_fifty} & Gaps4250 & PROVED & \checkmark & verified \\
\texttt{singular\_series\_pos\_evenPair\_fortyEight} & Gaps4250 & PROVED & \checkmark & verified \\
\texttt{singular\_series\_pos\_evenPair\_fortyFour} & Gaps4250 & PROVED & \checkmark & verified \\
\texttt{singular\_series\_pos\_evenPair\_fortySix} & Gaps4250 & PROVED & \checkmark & verified \\
\texttt{singular\_series\_pos\_evenPair\_fortyTwo} & Gaps4250 & PROVED & \checkmark & verified \\
\texttt{evenPair\_card\_fourHundredForty} & Gaps432440 & PROVED & \checkmark & verified \\
\texttt{evenPair\_card\_fourHundredThirtyEight} & Gaps432440 & PROVED & \checkmark & verified \\
\texttt{evenPair\_card\_fourHundredThirtyFour} & Gaps432440 & PROVED & \checkmark & verified \\
\texttt{evenPair\_card\_fourHundredThirtySix} & Gaps432440 & PROVED & \checkmark & verified \\
\texttt{evenPair\_card\_fourHundredThirtyTwo} & Gaps432440 & PROVED & \checkmark & verified \\
\texttt{isAdmissible\_evenPair\_fourHundredForty} & Gaps432440 & PROVED & \checkmark & verified \\
\texttt{isAdmissible\_evenPair\_fourHundredThirtyEight} & Gaps432440 & PROVED & \checkmark & verified \\
\texttt{isAdmissible\_evenPair\_fourHundredThirtyFour} & Gaps432440 & PROVED & \checkmark & verified \\
\texttt{isAdmissible\_evenPair\_fourHundredThirtySix} & Gaps432440 & PROVED & \checkmark & verified \\
\texttt{isAdmissible\_evenPair\_fourHundredThirtyTwo} & Gaps432440 & PROVED & \checkmark & verified \\
\texttt{localFactor\_fourHundredForty\_two} & Gaps432440 & PROVED & \checkmark & verified \\
\texttt{localFactor\_fourHundredThirtyTwo\_two} & Gaps432440 & PROVED & \checkmark & verified \\
\texttt{nu\_p\_fourHundredForty} & Gaps432440 & PROVED & \checkmark & verified \\
\texttt{nu\_p\_fourHundredForty\_two} & Gaps432440 & PROVED & \checkmark & verified \\
\texttt{nu\_p\_fourHundredThirtyEight} & Gaps432440 & PROVED & \checkmark & verified \\
\texttt{nu\_p\_fourHundredThirtyFour} & Gaps432440 & PROVED & \checkmark & verified \\
\texttt{nu\_p\_fourHundredThirtySix} & Gaps432440 & PROVED & \checkmark & verified \\
\texttt{nu\_p\_fourHundredThirtyTwo} & Gaps432440 & PROVED & \checkmark & verified \\
\texttt{nu\_p\_fourHundredThirtyTwo\_two} & Gaps432440 & PROVED & \checkmark & verified \\
\texttt{singular\_series\_finite\_pos\_evenPair\_fourHundredForty} & Gaps432440 & PROVED & \checkmark & verified \\
\texttt{singular\_series\_finite\_pos\_evenPair\_fourHundredThirtyEight} & Gaps432440 & PROVED & \checkmark & verified \\
\texttt{singular\_series\_finite\_pos\_evenPair\_fourHundredThirtyFour} & Gaps432440 & PROVED & \checkmark & verified \\
\texttt{singular\_series\_finite\_pos\_evenPair\_fourHundredThirtySix} & Gaps432440 & PROVED & \checkmark & verified \\
\texttt{singular\_series\_finite\_pos\_evenPair\_fourHundredThirtyTwo} & Gaps432440 & PROVED & \checkmark & verified \\
\texttt{singular\_series\_pos\_evenPair\_fourHundredForty} & Gaps432440 & PROVED & \checkmark & verified \\
\texttt{singular\_series\_pos\_evenPair\_fourHundredThirtyEight} & Gaps432440 & PROVED & \checkmark & verified \\
\texttt{singular\_series\_pos\_evenPair\_fourHundredThirtyFour} & Gaps432440 & PROVED & \checkmark & verified \\
\texttt{singular\_series\_pos\_evenPair\_fourHundredThirtySix} & Gaps432440 & PROVED & \checkmark & verified \\
\texttt{singular\_series\_pos\_evenPair\_fourHundredThirtyTwo} & Gaps432440 & PROVED & \checkmark & verified \\
\texttt{evenPair\_card\_fourHundredFifty} & Gaps442450 & PROVED & \checkmark & verified \\
\texttt{evenPair\_card\_fourHundredFortyEight} & Gaps442450 & PROVED & \checkmark & verified \\
\texttt{evenPair\_card\_fourHundredFortyFour} & Gaps442450 & PROVED & \checkmark & verified \\
\texttt{evenPair\_card\_fourHundredFortySix} & Gaps442450 & PROVED & \checkmark & verified \\
\texttt{evenPair\_card\_fourHundredFortyTwo} & Gaps442450 & PROVED & \checkmark & verified \\
\texttt{isAdmissible\_evenPair\_fourHundredFifty} & Gaps442450 & PROVED & \checkmark & verified \\
\texttt{isAdmissible\_evenPair\_fourHundredFortyEight} & Gaps442450 & PROVED & \checkmark & verified \\
\texttt{isAdmissible\_evenPair\_fourHundredFortyFour} & Gaps442450 & PROVED & \checkmark & verified \\
\texttt{isAdmissible\_evenPair\_fourHundredFortySix} & Gaps442450 & PROVED & \checkmark & verified \\
\texttt{isAdmissible\_evenPair\_fourHundredFortyTwo} & Gaps442450 & PROVED & \checkmark & verified \\
\texttt{localFactor\_fourHundredFifty\_two} & Gaps442450 & PROVED & \checkmark & verified \\
\texttt{localFactor\_fourHundredFortyTwo\_two} & Gaps442450 & PROVED & \checkmark & verified \\
\texttt{nu\_p\_fourHundredFifty} & Gaps442450 & PROVED & \checkmark & verified \\
\texttt{nu\_p\_fourHundredFifty\_two} & Gaps442450 & PROVED & \checkmark & verified \\
\texttt{nu\_p\_fourHundredFortyEight} & Gaps442450 & PROVED & \checkmark & verified \\
\texttt{nu\_p\_fourHundredFortyFour} & Gaps442450 & PROVED & \checkmark & verified \\
\texttt{nu\_p\_fourHundredFortySix} & Gaps442450 & PROVED & \checkmark & verified \\
\texttt{nu\_p\_fourHundredFortyTwo} & Gaps442450 & PROVED & \checkmark & verified \\
\texttt{nu\_p\_fourHundredFortyTwo\_two} & Gaps442450 & PROVED & \checkmark & verified \\
\texttt{singular\_series\_finite\_pos\_evenPair\_fourHundredFifty} & Gaps442450 & PROVED & \checkmark & verified \\
\texttt{singular\_series\_finite\_pos\_evenPair\_fourHundredFortyEight} & Gaps442450 & PROVED & \checkmark & verified \\
\texttt{singular\_series\_finite\_pos\_evenPair\_fourHundredFortyFour} & Gaps442450 & PROVED & \checkmark & verified \\
\texttt{singular\_series\_finite\_pos\_evenPair\_fourHundredFortySix} & Gaps442450 & PROVED & \checkmark & verified \\
\texttt{singular\_series\_finite\_pos\_evenPair\_fourHundredFortyTwo} & Gaps442450 & PROVED & \checkmark & verified \\
\texttt{singular\_series\_pos\_evenPair\_fourHundredFifty} & Gaps442450 & PROVED & \checkmark & verified \\
\texttt{singular\_series\_pos\_evenPair\_fourHundredFortyEight} & Gaps442450 & PROVED & \checkmark & verified \\
\texttt{singular\_series\_pos\_evenPair\_fourHundredFortyFour} & Gaps442450 & PROVED & \checkmark & verified \\
\texttt{singular\_series\_pos\_evenPair\_fourHundredFortySix} & Gaps442450 & PROVED & \checkmark & verified \\
\texttt{singular\_series\_pos\_evenPair\_fourHundredFortyTwo} & Gaps442450 & PROVED & \checkmark & verified \\
\texttt{evenPair\_card\_fourHundredFiftyEight} & Gaps452460 & PROVED & \checkmark & verified \\
\texttt{evenPair\_card\_fourHundredFiftyFour} & Gaps452460 & PROVED & \checkmark & verified \\
\texttt{evenPair\_card\_fourHundredFiftySix} & Gaps452460 & PROVED & \checkmark & verified \\
\texttt{evenPair\_card\_fourHundredFiftyTwo} & Gaps452460 & PROVED & \checkmark & verified \\
\texttt{evenPair\_card\_fourHundredSixty} & Gaps452460 & PROVED & \checkmark & verified \\
\texttt{isAdmissible\_evenPair\_fourHundredFiftyEight} & Gaps452460 & PROVED & \checkmark & verified \\
\texttt{isAdmissible\_evenPair\_fourHundredFiftyFour} & Gaps452460 & PROVED & \checkmark & verified \\
\texttt{isAdmissible\_evenPair\_fourHundredFiftySix} & Gaps452460 & PROVED & \checkmark & verified \\
\texttt{isAdmissible\_evenPair\_fourHundredFiftyTwo} & Gaps452460 & PROVED & \checkmark & verified \\
\texttt{isAdmissible\_evenPair\_fourHundredSixty} & Gaps452460 & PROVED & \checkmark & verified \\
\texttt{localFactor\_fourHundredFiftyTwo\_two} & Gaps452460 & PROVED & \checkmark & verified \\
\texttt{localFactor\_fourHundredSixty\_two} & Gaps452460 & PROVED & \checkmark & verified \\
\texttt{nu\_p\_fourHundredFiftyEight} & Gaps452460 & PROVED & \checkmark & verified \\
\texttt{nu\_p\_fourHundredFiftyFour} & Gaps452460 & PROVED & \checkmark & verified \\
\texttt{nu\_p\_fourHundredFiftySix} & Gaps452460 & PROVED & \checkmark & verified \\
\texttt{nu\_p\_fourHundredFiftyTwo} & Gaps452460 & PROVED & \checkmark & verified \\
\texttt{nu\_p\_fourHundredFiftyTwo\_two} & Gaps452460 & PROVED & \checkmark & verified \\
\texttt{nu\_p\_fourHundredSixty} & Gaps452460 & PROVED & \checkmark & verified \\
\texttt{nu\_p\_fourHundredSixty\_two} & Gaps452460 & PROVED & \checkmark & verified \\
\texttt{singular\_series\_finite\_pos\_evenPair\_fourHundredFiftyEight} & Gaps452460 & PROVED & \checkmark & verified \\
\texttt{singular\_series\_finite\_pos\_evenPair\_fourHundredFiftyFour} & Gaps452460 & PROVED & \checkmark & verified \\
\texttt{singular\_series\_finite\_pos\_evenPair\_fourHundredFiftySix} & Gaps452460 & PROVED & \checkmark & verified \\
\texttt{singular\_series\_finite\_pos\_evenPair\_fourHundredFiftyTwo} & Gaps452460 & PROVED & \checkmark & verified \\
\texttt{singular\_series\_finite\_pos\_evenPair\_fourHundredSixty} & Gaps452460 & PROVED & \checkmark & verified \\
\texttt{singular\_series\_pos\_evenPair\_fourHundredFiftyEight} & Gaps452460 & PROVED & \checkmark & verified \\
\texttt{singular\_series\_pos\_evenPair\_fourHundredFiftyFour} & Gaps452460 & PROVED & \checkmark & verified \\
\texttt{singular\_series\_pos\_evenPair\_fourHundredFiftySix} & Gaps452460 & PROVED & \checkmark & verified \\
\texttt{singular\_series\_pos\_evenPair\_fourHundredFiftyTwo} & Gaps452460 & PROVED & \checkmark & verified \\
\texttt{singular\_series\_pos\_evenPair\_fourHundredSixty} & Gaps452460 & PROVED & \checkmark & verified \\
\texttt{evenPair\_card\_fourHundredSeventy} & Gaps462470 & PROVED & \checkmark & verified \\
\texttt{evenPair\_card\_fourHundredSixtyEight} & Gaps462470 & PROVED & \checkmark & verified \\
\texttt{evenPair\_card\_fourHundredSixtyFour} & Gaps462470 & PROVED & \checkmark & verified \\
\texttt{evenPair\_card\_fourHundredSixtySix} & Gaps462470 & PROVED & \checkmark & verified \\
\texttt{evenPair\_card\_fourHundredSixtyTwo} & Gaps462470 & PROVED & \checkmark & verified \\
\texttt{isAdmissible\_evenPair\_fourHundredSeventy} & Gaps462470 & PROVED & \checkmark & verified \\
\texttt{isAdmissible\_evenPair\_fourHundredSixtyEight} & Gaps462470 & PROVED & \checkmark & verified \\
\texttt{isAdmissible\_evenPair\_fourHundredSixtyFour} & Gaps462470 & PROVED & \checkmark & verified \\
\texttt{isAdmissible\_evenPair\_fourHundredSixtySix} & Gaps462470 & PROVED & \checkmark & verified \\
\texttt{isAdmissible\_evenPair\_fourHundredSixtyTwo} & Gaps462470 & PROVED & \checkmark & verified \\
\texttt{localFactor\_fourHundredSeventy\_two} & Gaps462470 & PROVED & \checkmark & verified \\
\texttt{localFactor\_fourHundredSixtyTwo\_two} & Gaps462470 & PROVED & \checkmark & verified \\
\texttt{nu\_p\_fourHundredSeventy} & Gaps462470 & PROVED & \checkmark & verified \\
\texttt{nu\_p\_fourHundredSeventy\_two} & Gaps462470 & PROVED & \checkmark & verified \\
\texttt{nu\_p\_fourHundredSixtyEight} & Gaps462470 & PROVED & \checkmark & verified \\
\texttt{nu\_p\_fourHundredSixtyFour} & Gaps462470 & PROVED & \checkmark & verified \\
\texttt{nu\_p\_fourHundredSixtySix} & Gaps462470 & PROVED & \checkmark & verified \\
\texttt{nu\_p\_fourHundredSixtyTwo} & Gaps462470 & PROVED & \checkmark & verified \\
\texttt{nu\_p\_fourHundredSixtyTwo\_two} & Gaps462470 & PROVED & \checkmark & verified \\
\texttt{singular\_series\_finite\_pos\_evenPair\_fourHundredSeventy} & Gaps462470 & PROVED & \checkmark & verified \\
\texttt{singular\_series\_finite\_pos\_evenPair\_fourHundredSixtyEight} & Gaps462470 & PROVED & \checkmark & verified \\
\texttt{singular\_series\_finite\_pos\_evenPair\_fourHundredSixtyFour} & Gaps462470 & PROVED & \checkmark & verified \\
\texttt{singular\_series\_finite\_pos\_evenPair\_fourHundredSixtySix} & Gaps462470 & PROVED & \checkmark & verified \\
\texttt{singular\_series\_finite\_pos\_evenPair\_fourHundredSixtyTwo} & Gaps462470 & PROVED & \checkmark & verified \\
\texttt{singular\_series\_pos\_evenPair\_fourHundredSeventy} & Gaps462470 & PROVED & \checkmark & verified \\
\texttt{singular\_series\_pos\_evenPair\_fourHundredSixtyEight} & Gaps462470 & PROVED & \checkmark & verified \\
\texttt{singular\_series\_pos\_evenPair\_fourHundredSixtyFour} & Gaps462470 & PROVED & \checkmark & verified \\
\texttt{singular\_series\_pos\_evenPair\_fourHundredSixtySix} & Gaps462470 & PROVED & \checkmark & verified \\
\texttt{singular\_series\_pos\_evenPair\_fourHundredSixtyTwo} & Gaps462470 & PROVED & \checkmark & verified \\
\texttt{evenPair\_card\_fourHundredEighty} & Gaps472480 & PROVED & \checkmark & verified \\
\texttt{evenPair\_card\_fourHundredSeventyEight} & Gaps472480 & PROVED & \checkmark & verified \\
\texttt{evenPair\_card\_fourHundredSeventyFour} & Gaps472480 & PROVED & \checkmark & verified \\
\texttt{evenPair\_card\_fourHundredSeventySix} & Gaps472480 & PROVED & \checkmark & verified \\
\texttt{evenPair\_card\_fourHundredSeventyTwo} & Gaps472480 & PROVED & \checkmark & verified \\
\texttt{isAdmissible\_evenPair\_fourHundredEighty} & Gaps472480 & PROVED & \checkmark & verified \\
\texttt{isAdmissible\_evenPair\_fourHundredSeventyEight} & Gaps472480 & PROVED & \checkmark & verified \\
\texttt{isAdmissible\_evenPair\_fourHundredSeventyFour} & Gaps472480 & PROVED & \checkmark & verified \\
\texttt{isAdmissible\_evenPair\_fourHundredSeventySix} & Gaps472480 & PROVED & \checkmark & verified \\
\texttt{isAdmissible\_evenPair\_fourHundredSeventyTwo} & Gaps472480 & PROVED & \checkmark & verified \\
\texttt{localFactor\_fourHundredEighty\_two} & Gaps472480 & PROVED & \checkmark & verified \\
\texttt{localFactor\_fourHundredSeventyTwo\_two} & Gaps472480 & PROVED & \checkmark & verified \\
\texttt{nu\_p\_fourHundredEighty} & Gaps472480 & PROVED & \checkmark & verified \\
\texttt{nu\_p\_fourHundredEighty\_two} & Gaps472480 & PROVED & \checkmark & verified \\
\texttt{nu\_p\_fourHundredSeventyEight} & Gaps472480 & PROVED & \checkmark & verified \\
\texttt{nu\_p\_fourHundredSeventyFour} & Gaps472480 & PROVED & \checkmark & verified \\
\texttt{nu\_p\_fourHundredSeventySix} & Gaps472480 & PROVED & \checkmark & verified \\
\texttt{nu\_p\_fourHundredSeventyTwo} & Gaps472480 & PROVED & \checkmark & verified \\
\texttt{nu\_p\_fourHundredSeventyTwo\_two} & Gaps472480 & PROVED & \checkmark & verified \\
\texttt{singular\_series\_finite\_pos\_evenPair\_fourHundredEighty} & Gaps472480 & PROVED & \checkmark & verified \\
\texttt{singular\_series\_finite\_pos\_evenPair\_fourHundredSeventyEight} & Gaps472480 & PROVED & \checkmark & verified \\
\texttt{singular\_series\_finite\_pos\_evenPair\_fourHundredSeventyFour} & Gaps472480 & PROVED & \checkmark & verified \\
\texttt{singular\_series\_finite\_pos\_evenPair\_fourHundredSeventySix} & Gaps472480 & PROVED & \checkmark & verified \\
\texttt{singular\_series\_finite\_pos\_evenPair\_fourHundredSeventyTwo} & Gaps472480 & PROVED & \checkmark & verified \\
\texttt{singular\_series\_pos\_evenPair\_fourHundredEighty} & Gaps472480 & PROVED & \checkmark & verified \\
\texttt{singular\_series\_pos\_evenPair\_fourHundredSeventyEight} & Gaps472480 & PROVED & \checkmark & verified \\
\texttt{singular\_series\_pos\_evenPair\_fourHundredSeventyFour} & Gaps472480 & PROVED & \checkmark & verified \\
\texttt{singular\_series\_pos\_evenPair\_fourHundredSeventySix} & Gaps472480 & PROVED & \checkmark & verified \\
\texttt{singular\_series\_pos\_evenPair\_fourHundredSeventyTwo} & Gaps472480 & PROVED & \checkmark & verified \\
\texttt{evenPair\_card\_fourHundredEightyEight} & Gaps482490 & PROVED & \checkmark & verified \\
\texttt{evenPair\_card\_fourHundredEightyFour} & Gaps482490 & PROVED & \checkmark & verified \\
\texttt{evenPair\_card\_fourHundredEightySix} & Gaps482490 & PROVED & \checkmark & verified \\
\texttt{evenPair\_card\_fourHundredEightyTwo} & Gaps482490 & PROVED & \checkmark & verified \\
\texttt{evenPair\_card\_fourHundredNinety} & Gaps482490 & PROVED & \checkmark & verified \\
\texttt{isAdmissible\_evenPair\_fourHundredEightyEight} & Gaps482490 & PROVED & \checkmark & verified \\
\texttt{isAdmissible\_evenPair\_fourHundredEightyFour} & Gaps482490 & PROVED & \checkmark & verified \\
\texttt{isAdmissible\_evenPair\_fourHundredEightySix} & Gaps482490 & PROVED & \checkmark & verified \\
\texttt{isAdmissible\_evenPair\_fourHundredEightyTwo} & Gaps482490 & PROVED & \checkmark & verified \\
\texttt{isAdmissible\_evenPair\_fourHundredNinety} & Gaps482490 & PROVED & \checkmark & verified \\
\texttt{localFactor\_fourHundredEightyTwo\_two} & Gaps482490 & PROVED & \checkmark & verified \\
\texttt{localFactor\_fourHundredNinety\_two} & Gaps482490 & PROVED & \checkmark & verified \\
\texttt{nu\_p\_fourHundredEightyEight} & Gaps482490 & PROVED & \checkmark & verified \\
\texttt{nu\_p\_fourHundredEightyFour} & Gaps482490 & PROVED & \checkmark & verified \\
\texttt{nu\_p\_fourHundredEightySix} & Gaps482490 & PROVED & \checkmark & verified \\
\texttt{nu\_p\_fourHundredEightyTwo} & Gaps482490 & PROVED & \checkmark & verified \\
\texttt{nu\_p\_fourHundredEightyTwo\_two} & Gaps482490 & PROVED & \checkmark & verified \\
\texttt{nu\_p\_fourHundredNinety} & Gaps482490 & PROVED & \checkmark & verified \\
\texttt{nu\_p\_fourHundredNinety\_two} & Gaps482490 & PROVED & \checkmark & verified \\
\texttt{singular\_series\_finite\_pos\_evenPair\_fourHundredEightyEight} & Gaps482490 & PROVED & \checkmark & verified \\
\texttt{singular\_series\_finite\_pos\_evenPair\_fourHundredEightyFour} & Gaps482490 & PROVED & \checkmark & verified \\
\texttt{singular\_series\_finite\_pos\_evenPair\_fourHundredEightySix} & Gaps482490 & PROVED & \checkmark & verified \\
\texttt{singular\_series\_finite\_pos\_evenPair\_fourHundredEightyTwo} & Gaps482490 & PROVED & \checkmark & verified \\
\texttt{singular\_series\_finite\_pos\_evenPair\_fourHundredNinety} & Gaps482490 & PROVED & \checkmark & verified \\
\texttt{singular\_series\_pos\_evenPair\_fourHundredEightyEight} & Gaps482490 & PROVED & \checkmark & verified \\
\texttt{singular\_series\_pos\_evenPair\_fourHundredEightyFour} & Gaps482490 & PROVED & \checkmark & verified \\
\texttt{singular\_series\_pos\_evenPair\_fourHundredEightySix} & Gaps482490 & PROVED & \checkmark & verified \\
\texttt{singular\_series\_pos\_evenPair\_fourHundredEightyTwo} & Gaps482490 & PROVED & \checkmark & verified \\
\texttt{singular\_series\_pos\_evenPair\_fourHundredNinety} & Gaps482490 & PROVED & \checkmark & verified \\
\texttt{evenPair\_card\_fiveHundred} & Gaps492500 & PROVED & \checkmark & verified \\
\texttt{evenPair\_card\_fourHundredNinetyEight} & Gaps492500 & PROVED & \checkmark & verified \\
\texttt{evenPair\_card\_fourHundredNinetyFour} & Gaps492500 & PROVED & \checkmark & verified \\
\texttt{evenPair\_card\_fourHundredNinetySix} & Gaps492500 & PROVED & \checkmark & verified \\
\texttt{evenPair\_card\_fourHundredNinetyTwo} & Gaps492500 & PROVED & \checkmark & verified \\
\texttt{isAdmissible\_evenPair\_fiveHundred} & Gaps492500 & PROVED & \checkmark & verified \\
\texttt{isAdmissible\_evenPair\_fourHundredNinetyEight} & Gaps492500 & PROVED & \checkmark & verified \\
\texttt{isAdmissible\_evenPair\_fourHundredNinetyFour} & Gaps492500 & PROVED & \checkmark & verified \\
\texttt{isAdmissible\_evenPair\_fourHundredNinetySix} & Gaps492500 & PROVED & \checkmark & verified \\
\texttt{isAdmissible\_evenPair\_fourHundredNinetyTwo} & Gaps492500 & PROVED & \checkmark & verified \\
\texttt{localFactor\_fiveHundred\_two} & Gaps492500 & PROVED & \checkmark & verified \\
\texttt{localFactor\_fourHundredNinetyTwo\_two} & Gaps492500 & PROVED & \checkmark & verified \\
\texttt{nu\_p\_fiveHundred} & Gaps492500 & PROVED & \checkmark & verified \\
\texttt{nu\_p\_fiveHundred\_two} & Gaps492500 & PROVED & \checkmark & verified \\
\texttt{nu\_p\_fourHundredNinetyEight} & Gaps492500 & PROVED & \checkmark & verified \\
\texttt{nu\_p\_fourHundredNinetyFour} & Gaps492500 & PROVED & \checkmark & verified \\
\texttt{nu\_p\_fourHundredNinetySix} & Gaps492500 & PROVED & \checkmark & verified \\
\texttt{nu\_p\_fourHundredNinetyTwo} & Gaps492500 & PROVED & \checkmark & verified \\
\texttt{nu\_p\_fourHundredNinetyTwo\_two} & Gaps492500 & PROVED & \checkmark & verified \\
\texttt{singular\_series\_finite\_pos\_evenPair\_fiveHundred} & Gaps492500 & PROVED & \checkmark & verified \\
\texttt{singular\_series\_finite\_pos\_evenPair\_fourHundredNinetyEight} & Gaps492500 & PROVED & \checkmark & verified \\
\texttt{singular\_series\_finite\_pos\_evenPair\_fourHundredNinetyFour} & Gaps492500 & PROVED & \checkmark & verified \\
\texttt{singular\_series\_finite\_pos\_evenPair\_fourHundredNinetySix} & Gaps492500 & PROVED & \checkmark & verified \\
\texttt{singular\_series\_finite\_pos\_evenPair\_fourHundredNinetyTwo} & Gaps492500 & PROVED & \checkmark & verified \\
\texttt{singular\_series\_pos\_evenPair\_fiveHundred} & Gaps492500 & PROVED & \checkmark & verified \\
\texttt{singular\_series\_pos\_evenPair\_fourHundredNinetyEight} & Gaps492500 & PROVED & \checkmark & verified \\
\texttt{singular\_series\_pos\_evenPair\_fourHundredNinetyFour} & Gaps492500 & PROVED & \checkmark & verified \\
\texttt{singular\_series\_pos\_evenPair\_fourHundredNinetySix} & Gaps492500 & PROVED & \checkmark & verified \\
\texttt{singular\_series\_pos\_evenPair\_fourHundredNinetyTwo} & Gaps492500 & PROVED & \checkmark & verified \\
\texttt{evenPair\_card\_fiveHundredEight} & Gaps502510 & PROVED & \checkmark & verified \\
\texttt{evenPair\_card\_fiveHundredFour} & Gaps502510 & PROVED & \checkmark & verified \\
\texttt{evenPair\_card\_fiveHundredSix} & Gaps502510 & PROVED & \checkmark & verified \\
\texttt{evenPair\_card\_fiveHundredTen} & Gaps502510 & PROVED & \checkmark & verified \\
\texttt{evenPair\_card\_fiveHundredTwo} & Gaps502510 & PROVED & \checkmark & verified \\
\texttt{isAdmissible\_evenPair\_fiveHundredEight} & Gaps502510 & PROVED & \checkmark & verified \\
\texttt{isAdmissible\_evenPair\_fiveHundredFour} & Gaps502510 & PROVED & \checkmark & verified \\
\texttt{isAdmissible\_evenPair\_fiveHundredSix} & Gaps502510 & PROVED & \checkmark & verified \\
\texttt{isAdmissible\_evenPair\_fiveHundredTen} & Gaps502510 & PROVED & \checkmark & verified \\
\texttt{isAdmissible\_evenPair\_fiveHundredTwo} & Gaps502510 & PROVED & \checkmark & verified \\
\texttt{localFactor\_fiveHundredTen\_two} & Gaps502510 & PROVED & \checkmark & verified \\
\texttt{localFactor\_fiveHundredTwo\_two} & Gaps502510 & PROVED & \checkmark & verified \\
\texttt{nu\_p\_fiveHundredEight} & Gaps502510 & PROVED & \checkmark & verified \\
\texttt{nu\_p\_fiveHundredFour} & Gaps502510 & PROVED & \checkmark & verified \\
\texttt{nu\_p\_fiveHundredSix} & Gaps502510 & PROVED & \checkmark & verified \\
\texttt{nu\_p\_fiveHundredTen} & Gaps502510 & PROVED & \checkmark & verified \\
\texttt{nu\_p\_fiveHundredTen\_two} & Gaps502510 & PROVED & \checkmark & verified \\
\texttt{nu\_p\_fiveHundredTwo} & Gaps502510 & PROVED & \checkmark & verified \\
\texttt{nu\_p\_fiveHundredTwo\_two} & Gaps502510 & PROVED & \checkmark & verified \\
\texttt{singular\_series\_finite\_pos\_evenPair\_fiveHundredEight} & Gaps502510 & PROVED & \checkmark & verified \\
\texttt{singular\_series\_finite\_pos\_evenPair\_fiveHundredFour} & Gaps502510 & PROVED & \checkmark & verified \\
\texttt{singular\_series\_finite\_pos\_evenPair\_fiveHundredSix} & Gaps502510 & PROVED & \checkmark & verified \\
\texttt{singular\_series\_finite\_pos\_evenPair\_fiveHundredTen} & Gaps502510 & PROVED & \checkmark & verified \\
\texttt{singular\_series\_finite\_pos\_evenPair\_fiveHundredTwo} & Gaps502510 & PROVED & \checkmark & verified \\
\texttt{singular\_series\_pos\_evenPair\_fiveHundredEight} & Gaps502510 & PROVED & \checkmark & verified \\
\texttt{singular\_series\_pos\_evenPair\_fiveHundredFour} & Gaps502510 & PROVED & \checkmark & verified \\
\texttt{singular\_series\_pos\_evenPair\_fiveHundredSix} & Gaps502510 & PROVED & \checkmark & verified \\
\texttt{singular\_series\_pos\_evenPair\_fiveHundredTen} & Gaps502510 & PROVED & \checkmark & verified \\
\texttt{singular\_series\_pos\_evenPair\_fiveHundredTwo} & Gaps502510 & PROVED & \checkmark & verified \\
\texttt{evenPair\_card\_fiveHundredEighteen} & Gaps512520 & PROVED & \checkmark & verified \\
\texttt{evenPair\_card\_fiveHundredFourteen} & Gaps512520 & PROVED & \checkmark & verified \\
\texttt{evenPair\_card\_fiveHundredSixteen} & Gaps512520 & PROVED & \checkmark & verified \\
\texttt{evenPair\_card\_fiveHundredTwelve} & Gaps512520 & PROVED & \checkmark & verified \\
\texttt{evenPair\_card\_fiveHundredTwenty} & Gaps512520 & PROVED & \checkmark & verified \\
\texttt{isAdmissible\_evenPair\_fiveHundredEighteen} & Gaps512520 & PROVED & \checkmark & verified \\
\texttt{isAdmissible\_evenPair\_fiveHundredFourteen} & Gaps512520 & PROVED & \checkmark & verified \\
\texttt{isAdmissible\_evenPair\_fiveHundredSixteen} & Gaps512520 & PROVED & \checkmark & verified \\
\texttt{isAdmissible\_evenPair\_fiveHundredTwelve} & Gaps512520 & PROVED & \checkmark & verified \\
\texttt{isAdmissible\_evenPair\_fiveHundredTwenty} & Gaps512520 & PROVED & \checkmark & verified \\
\texttt{localFactor\_fiveHundredTwelve\_two} & Gaps512520 & PROVED & \checkmark & verified \\
\texttt{localFactor\_fiveHundredTwenty\_two} & Gaps512520 & PROVED & \checkmark & verified \\
\texttt{nu\_p\_fiveHundredEighteen} & Gaps512520 & PROVED & \checkmark & verified \\
\texttt{nu\_p\_fiveHundredFourteen} & Gaps512520 & PROVED & \checkmark & verified \\
\texttt{nu\_p\_fiveHundredSixteen} & Gaps512520 & PROVED & \checkmark & verified \\
\texttt{nu\_p\_fiveHundredTwelve} & Gaps512520 & PROVED & \checkmark & verified \\
\texttt{nu\_p\_fiveHundredTwelve\_two} & Gaps512520 & PROVED & \checkmark & verified \\
\texttt{nu\_p\_fiveHundredTwenty} & Gaps512520 & PROVED & \checkmark & verified \\
\texttt{nu\_p\_fiveHundredTwenty\_two} & Gaps512520 & PROVED & \checkmark & verified \\
\texttt{singular\_series\_finite\_pos\_evenPair\_fiveHundredEighteen} & Gaps512520 & PROVED & \checkmark & verified \\
\texttt{singular\_series\_finite\_pos\_evenPair\_fiveHundredFourteen} & Gaps512520 & PROVED & \checkmark & verified \\
\texttt{singular\_series\_finite\_pos\_evenPair\_fiveHundredSixteen} & Gaps512520 & PROVED & \checkmark & verified \\
\texttt{singular\_series\_finite\_pos\_evenPair\_fiveHundredTwelve} & Gaps512520 & PROVED & \checkmark & verified \\
\texttt{singular\_series\_finite\_pos\_evenPair\_fiveHundredTwenty} & Gaps512520 & PROVED & \checkmark & verified \\
\texttt{singular\_series\_pos\_evenPair\_fiveHundredEighteen} & Gaps512520 & PROVED & \checkmark & verified \\
\texttt{singular\_series\_pos\_evenPair\_fiveHundredFourteen} & Gaps512520 & PROVED & \checkmark & verified \\
\texttt{singular\_series\_pos\_evenPair\_fiveHundredSixteen} & Gaps512520 & PROVED & \checkmark & verified \\
\texttt{singular\_series\_pos\_evenPair\_fiveHundredTwelve} & Gaps512520 & PROVED & \checkmark & verified \\
\texttt{singular\_series\_pos\_evenPair\_fiveHundredTwenty} & Gaps512520 & PROVED & \checkmark & verified \\
\texttt{evenPair\_card\_fiveHundredThirty} & Gaps522530 & PROVED & \checkmark & verified \\
\texttt{evenPair\_card\_fiveHundredTwentyEight} & Gaps522530 & PROVED & \checkmark & verified \\
\texttt{evenPair\_card\_fiveHundredTwentyFour} & Gaps522530 & PROVED & \checkmark & verified \\
\texttt{evenPair\_card\_fiveHundredTwentySix} & Gaps522530 & PROVED & \checkmark & verified \\
\texttt{evenPair\_card\_fiveHundredTwentyTwo} & Gaps522530 & PROVED & \checkmark & verified \\
\texttt{isAdmissible\_evenPair\_fiveHundredThirty} & Gaps522530 & PROVED & \checkmark & verified \\
\texttt{isAdmissible\_evenPair\_fiveHundredTwentyEight} & Gaps522530 & PROVED & \checkmark & verified \\
\texttt{isAdmissible\_evenPair\_fiveHundredTwentyFour} & Gaps522530 & PROVED & \checkmark & verified \\
\texttt{isAdmissible\_evenPair\_fiveHundredTwentySix} & Gaps522530 & PROVED & \checkmark & verified \\
\texttt{isAdmissible\_evenPair\_fiveHundredTwentyTwo} & Gaps522530 & PROVED & \checkmark & verified \\
\texttt{localFactor\_fiveHundredThirty\_two} & Gaps522530 & PROVED & \checkmark & verified \\
\texttt{localFactor\_fiveHundredTwentyTwo\_two} & Gaps522530 & PROVED & \checkmark & verified \\
\texttt{nu\_p\_fiveHundredThirty} & Gaps522530 & PROVED & \checkmark & verified \\
\texttt{nu\_p\_fiveHundredThirty\_two} & Gaps522530 & PROVED & \checkmark & verified \\
\texttt{nu\_p\_fiveHundredTwentyEight} & Gaps522530 & PROVED & \checkmark & verified \\
\texttt{nu\_p\_fiveHundredTwentyFour} & Gaps522530 & PROVED & \checkmark & verified \\
\texttt{nu\_p\_fiveHundredTwentySix} & Gaps522530 & PROVED & \checkmark & verified \\
\texttt{nu\_p\_fiveHundredTwentyTwo} & Gaps522530 & PROVED & \checkmark & verified \\
\texttt{nu\_p\_fiveHundredTwentyTwo\_two} & Gaps522530 & PROVED & \checkmark & verified \\
\texttt{singular\_series\_finite\_pos\_evenPair\_fiveHundredThirty} & Gaps522530 & PROVED & \checkmark & verified \\
\texttt{singular\_series\_finite\_pos\_evenPair\_fiveHundredTwentyEight} & Gaps522530 & PROVED & \checkmark & verified \\
\texttt{singular\_series\_finite\_pos\_evenPair\_fiveHundredTwentyFour} & Gaps522530 & PROVED & \checkmark & verified \\
\texttt{singular\_series\_finite\_pos\_evenPair\_fiveHundredTwentySix} & Gaps522530 & PROVED & \checkmark & verified \\
\texttt{singular\_series\_finite\_pos\_evenPair\_fiveHundredTwentyTwo} & Gaps522530 & PROVED & \checkmark & verified \\
\texttt{singular\_series\_pos\_evenPair\_fiveHundredThirty} & Gaps522530 & PROVED & \checkmark & verified \\
\texttt{singular\_series\_pos\_evenPair\_fiveHundredTwentyEight} & Gaps522530 & PROVED & \checkmark & verified \\
\texttt{singular\_series\_pos\_evenPair\_fiveHundredTwentyFour} & Gaps522530 & PROVED & \checkmark & verified \\
\texttt{singular\_series\_pos\_evenPair\_fiveHundredTwentySix} & Gaps522530 & PROVED & \checkmark & verified \\
\texttt{singular\_series\_pos\_evenPair\_fiveHundredTwentyTwo} & Gaps522530 & PROVED & \checkmark & verified \\
\texttt{evenPair\_card\_fiftyEight} & Gaps5260 & PROVED & \checkmark & verified \\
\texttt{evenPair\_card\_fiftyFour} & Gaps5260 & PROVED & \checkmark & verified \\
\texttt{evenPair\_card\_fiftySix} & Gaps5260 & PROVED & \checkmark & verified \\
\texttt{evenPair\_card\_fiftyTwo} & Gaps5260 & PROVED & \checkmark & verified \\
\texttt{evenPair\_card\_sixty} & Gaps5260 & PROVED & \checkmark & verified \\
\texttt{isAdmissible\_evenPair\_fiftyEight} & Gaps5260 & PROVED & \checkmark & verified \\
\texttt{isAdmissible\_evenPair\_fiftyFour} & Gaps5260 & PROVED & \checkmark & verified \\
\texttt{isAdmissible\_evenPair\_fiftySix} & Gaps5260 & PROVED & \checkmark & verified \\
\texttt{isAdmissible\_evenPair\_fiftyTwo} & Gaps5260 & PROVED & \checkmark & verified \\
\texttt{isAdmissible\_evenPair\_sixty} & Gaps5260 & PROVED & \checkmark & verified \\
\texttt{localFactor\_fiftyTwo\_two} & Gaps5260 & PROVED & \checkmark & verified \\
\texttt{localFactor\_sixty\_two} & Gaps5260 & PROVED & \checkmark & verified \\
\texttt{nu\_p\_fiftyEight} & Gaps5260 & PROVED & \checkmark & verified \\
\texttt{nu\_p\_fiftyFour} & Gaps5260 & PROVED & \checkmark & verified \\
\texttt{nu\_p\_fiftySix} & Gaps5260 & PROVED & \checkmark & verified \\
\texttt{nu\_p\_fiftyTwo} & Gaps5260 & PROVED & \checkmark & verified \\
\texttt{nu\_p\_fiftyTwo\_two} & Gaps5260 & PROVED & \checkmark & verified \\
\texttt{nu\_p\_sixty} & Gaps5260 & PROVED & \checkmark & verified \\
\texttt{nu\_p\_sixty\_two} & Gaps5260 & PROVED & \checkmark & verified \\
\texttt{singular\_series\_finite\_pos\_evenPair\_fiftyEight} & Gaps5260 & PROVED & \checkmark & verified \\
\texttt{singular\_series\_finite\_pos\_evenPair\_fiftyFour} & Gaps5260 & PROVED & \checkmark & verified \\
\texttt{singular\_series\_finite\_pos\_evenPair\_fiftySix} & Gaps5260 & PROVED & \checkmark & verified \\
\texttt{singular\_series\_finite\_pos\_evenPair\_fiftyTwo} & Gaps5260 & PROVED & \checkmark & verified \\
\texttt{singular\_series\_finite\_pos\_evenPair\_sixty} & Gaps5260 & PROVED & \checkmark & verified \\
\texttt{singular\_series\_pos\_evenPair\_fiftyEight} & Gaps5260 & PROVED & \checkmark & verified \\
\texttt{singular\_series\_pos\_evenPair\_fiftyFour} & Gaps5260 & PROVED & \checkmark & verified \\
\texttt{singular\_series\_pos\_evenPair\_fiftySix} & Gaps5260 & PROVED & \checkmark & verified \\
\texttt{singular\_series\_pos\_evenPair\_fiftyTwo} & Gaps5260 & PROVED & \checkmark & verified \\
\texttt{singular\_series\_pos\_evenPair\_sixty} & Gaps5260 & PROVED & \checkmark & verified \\
\texttt{evenPair\_card\_fiveHundredForty} & Gaps532540 & PROVED & \checkmark & verified \\
\texttt{evenPair\_card\_fiveHundredThirtyEight} & Gaps532540 & PROVED & \checkmark & verified \\
\texttt{evenPair\_card\_fiveHundredThirtyFour} & Gaps532540 & PROVED & \checkmark & verified \\
\texttt{evenPair\_card\_fiveHundredThirtySix} & Gaps532540 & PROVED & \checkmark & verified \\
\texttt{evenPair\_card\_fiveHundredThirtyTwo} & Gaps532540 & PROVED & \checkmark & verified \\
\texttt{isAdmissible\_evenPair\_fiveHundredForty} & Gaps532540 & PROVED & \checkmark & verified \\
\texttt{isAdmissible\_evenPair\_fiveHundredThirtyEight} & Gaps532540 & PROVED & \checkmark & verified \\
\texttt{isAdmissible\_evenPair\_fiveHundredThirtyFour} & Gaps532540 & PROVED & \checkmark & verified \\
\texttt{isAdmissible\_evenPair\_fiveHundredThirtySix} & Gaps532540 & PROVED & \checkmark & verified \\
\texttt{isAdmissible\_evenPair\_fiveHundredThirtyTwo} & Gaps532540 & PROVED & \checkmark & verified \\
\texttt{localFactor\_fiveHundredForty\_two} & Gaps532540 & PROVED & \checkmark & verified \\
\texttt{localFactor\_fiveHundredThirtyTwo\_two} & Gaps532540 & PROVED & \checkmark & verified \\
\texttt{nu\_p\_fiveHundredForty} & Gaps532540 & PROVED & \checkmark & verified \\
\texttt{nu\_p\_fiveHundredForty\_two} & Gaps532540 & PROVED & \checkmark & verified \\
\texttt{nu\_p\_fiveHundredThirtyEight} & Gaps532540 & PROVED & \checkmark & verified \\
\texttt{nu\_p\_fiveHundredThirtyFour} & Gaps532540 & PROVED & \checkmark & verified \\
\texttt{nu\_p\_fiveHundredThirtySix} & Gaps532540 & PROVED & \checkmark & verified \\
\texttt{nu\_p\_fiveHundredThirtyTwo} & Gaps532540 & PROVED & \checkmark & verified \\
\texttt{nu\_p\_fiveHundredThirtyTwo\_two} & Gaps532540 & PROVED & \checkmark & verified \\
\texttt{singular\_series\_finite\_pos\_evenPair\_fiveHundredForty} & Gaps532540 & PROVED & \checkmark & verified \\
\texttt{singular\_series\_finite\_pos\_evenPair\_fiveHundredThirtyEight} & Gaps532540 & PROVED & \checkmark & verified \\
\texttt{singular\_series\_finite\_pos\_evenPair\_fiveHundredThirtyFour} & Gaps532540 & PROVED & \checkmark & verified \\
\texttt{singular\_series\_finite\_pos\_evenPair\_fiveHundredThirtySix} & Gaps532540 & PROVED & \checkmark & verified \\
\texttt{singular\_series\_finite\_pos\_evenPair\_fiveHundredThirtyTwo} & Gaps532540 & PROVED & \checkmark & verified \\
\texttt{singular\_series\_pos\_evenPair\_fiveHundredForty} & Gaps532540 & PROVED & \checkmark & verified \\
\texttt{singular\_series\_pos\_evenPair\_fiveHundredThirtyEight} & Gaps532540 & PROVED & \checkmark & verified \\
\texttt{singular\_series\_pos\_evenPair\_fiveHundredThirtyFour} & Gaps532540 & PROVED & \checkmark & verified \\
\texttt{singular\_series\_pos\_evenPair\_fiveHundredThirtySix} & Gaps532540 & PROVED & \checkmark & verified \\
\texttt{singular\_series\_pos\_evenPair\_fiveHundredThirtyTwo} & Gaps532540 & PROVED & \checkmark & verified \\
\texttt{evenPair\_card\_fiveHundredFifty} & Gaps542550 & PROVED & \checkmark & verified \\
\texttt{evenPair\_card\_fiveHundredFortyEight} & Gaps542550 & PROVED & \checkmark & verified \\
\texttt{evenPair\_card\_fiveHundredFortyFour} & Gaps542550 & PROVED & \checkmark & verified \\
\texttt{evenPair\_card\_fiveHundredFortySix} & Gaps542550 & PROVED & \checkmark & verified \\
\texttt{evenPair\_card\_fiveHundredFortyTwo} & Gaps542550 & PROVED & \checkmark & verified \\
\texttt{isAdmissible\_evenPair\_fiveHundredFifty} & Gaps542550 & PROVED & \checkmark & verified \\
\texttt{isAdmissible\_evenPair\_fiveHundredFortyEight} & Gaps542550 & PROVED & \checkmark & verified \\
\texttt{isAdmissible\_evenPair\_fiveHundredFortyFour} & Gaps542550 & PROVED & \checkmark & verified \\
\texttt{isAdmissible\_evenPair\_fiveHundredFortySix} & Gaps542550 & PROVED & \checkmark & verified \\
\texttt{isAdmissible\_evenPair\_fiveHundredFortyTwo} & Gaps542550 & PROVED & \checkmark & verified \\
\texttt{localFactor\_fiveHundredFifty\_two} & Gaps542550 & PROVED & \checkmark & verified \\
\texttt{localFactor\_fiveHundredFortyTwo\_two} & Gaps542550 & PROVED & \checkmark & verified \\
\texttt{nu\_p\_fiveHundredFifty} & Gaps542550 & PROVED & \checkmark & verified \\
\texttt{nu\_p\_fiveHundredFifty\_two} & Gaps542550 & PROVED & \checkmark & verified \\
\texttt{nu\_p\_fiveHundredFortyEight} & Gaps542550 & PROVED & \checkmark & verified \\
\texttt{nu\_p\_fiveHundredFortyFour} & Gaps542550 & PROVED & \checkmark & verified \\
\texttt{nu\_p\_fiveHundredFortySix} & Gaps542550 & PROVED & \checkmark & verified \\
\texttt{nu\_p\_fiveHundredFortyTwo} & Gaps542550 & PROVED & \checkmark & verified \\
\texttt{nu\_p\_fiveHundredFortyTwo\_two} & Gaps542550 & PROVED & \checkmark & verified \\
\texttt{singular\_series\_finite\_pos\_evenPair\_fiveHundredFifty} & Gaps542550 & PROVED & \checkmark & verified \\
\texttt{singular\_series\_finite\_pos\_evenPair\_fiveHundredFortyEight} & Gaps542550 & PROVED & \checkmark & verified \\
\texttt{singular\_series\_finite\_pos\_evenPair\_fiveHundredFortyFour} & Gaps542550 & PROVED & \checkmark & verified \\
\texttt{singular\_series\_finite\_pos\_evenPair\_fiveHundredFortySix} & Gaps542550 & PROVED & \checkmark & verified \\
\texttt{singular\_series\_finite\_pos\_evenPair\_fiveHundredFortyTwo} & Gaps542550 & PROVED & \checkmark & verified \\
\texttt{singular\_series\_pos\_evenPair\_fiveHundredFifty} & Gaps542550 & PROVED & \checkmark & verified \\
\texttt{singular\_series\_pos\_evenPair\_fiveHundredFortyEight} & Gaps542550 & PROVED & \checkmark & verified \\
\texttt{singular\_series\_pos\_evenPair\_fiveHundredFortyFour} & Gaps542550 & PROVED & \checkmark & verified \\
\texttt{singular\_series\_pos\_evenPair\_fiveHundredFortySix} & Gaps542550 & PROVED & \checkmark & verified \\
\texttt{singular\_series\_pos\_evenPair\_fiveHundredFortyTwo} & Gaps542550 & PROVED & \checkmark & verified \\
\texttt{evenPair\_card\_fiveHundredFiftyEight} & Gaps552560 & PROVED & \checkmark & verified \\
\texttt{evenPair\_card\_fiveHundredFiftyFour} & Gaps552560 & PROVED & \checkmark & verified \\
\texttt{evenPair\_card\_fiveHundredFiftySix} & Gaps552560 & PROVED & \checkmark & verified \\
\texttt{evenPair\_card\_fiveHundredFiftyTwo} & Gaps552560 & PROVED & \checkmark & verified \\
\texttt{evenPair\_card\_fiveHundredSixty} & Gaps552560 & PROVED & \checkmark & verified \\
\texttt{isAdmissible\_evenPair\_fiveHundredFiftyEight} & Gaps552560 & PROVED & \checkmark & verified \\
\texttt{isAdmissible\_evenPair\_fiveHundredFiftyFour} & Gaps552560 & PROVED & \checkmark & verified \\
\texttt{isAdmissible\_evenPair\_fiveHundredFiftySix} & Gaps552560 & PROVED & \checkmark & verified \\
\texttt{isAdmissible\_evenPair\_fiveHundredFiftyTwo} & Gaps552560 & PROVED & \checkmark & verified \\
\texttt{isAdmissible\_evenPair\_fiveHundredSixty} & Gaps552560 & PROVED & \checkmark & verified \\
\texttt{localFactor\_fiveHundredFiftyTwo\_two} & Gaps552560 & PROVED & \checkmark & verified \\
\texttt{localFactor\_fiveHundredSixty\_two} & Gaps552560 & PROVED & \checkmark & verified \\
\texttt{nu\_p\_fiveHundredFiftyEight} & Gaps552560 & PROVED & \checkmark & verified \\
\texttt{nu\_p\_fiveHundredFiftyFour} & Gaps552560 & PROVED & \checkmark & verified \\
\texttt{nu\_p\_fiveHundredFiftySix} & Gaps552560 & PROVED & \checkmark & verified \\
\texttt{nu\_p\_fiveHundredFiftyTwo} & Gaps552560 & PROVED & \checkmark & verified \\
\texttt{nu\_p\_fiveHundredFiftyTwo\_two} & Gaps552560 & PROVED & \checkmark & verified \\
\texttt{nu\_p\_fiveHundredSixty} & Gaps552560 & PROVED & \checkmark & verified \\
\texttt{nu\_p\_fiveHundredSixty\_two} & Gaps552560 & PROVED & \checkmark & verified \\
\texttt{singular\_series\_finite\_pos\_evenPair\_fiveHundredFiftyEight} & Gaps552560 & PROVED & \checkmark & verified \\
\texttt{singular\_series\_finite\_pos\_evenPair\_fiveHundredFiftyFour} & Gaps552560 & PROVED & \checkmark & verified \\
\texttt{singular\_series\_finite\_pos\_evenPair\_fiveHundredFiftySix} & Gaps552560 & PROVED & \checkmark & verified \\
\texttt{singular\_series\_finite\_pos\_evenPair\_fiveHundredFiftyTwo} & Gaps552560 & PROVED & \checkmark & verified \\
\texttt{singular\_series\_finite\_pos\_evenPair\_fiveHundredSixty} & Gaps552560 & PROVED & \checkmark & verified \\
\texttt{singular\_series\_pos\_evenPair\_fiveHundredFiftyEight} & Gaps552560 & PROVED & \checkmark & verified \\
\texttt{singular\_series\_pos\_evenPair\_fiveHundredFiftyFour} & Gaps552560 & PROVED & \checkmark & verified \\
\texttt{singular\_series\_pos\_evenPair\_fiveHundredFiftySix} & Gaps552560 & PROVED & \checkmark & verified \\
\texttt{singular\_series\_pos\_evenPair\_fiveHundredFiftyTwo} & Gaps552560 & PROVED & \checkmark & verified \\
\texttt{singular\_series\_pos\_evenPair\_fiveHundredSixty} & Gaps552560 & PROVED & \checkmark & verified \\
\texttt{evenPair\_card\_fiveHundredSeventy} & Gaps562570 & PROVED & \checkmark & verified \\
\texttt{evenPair\_card\_fiveHundredSixtyEight} & Gaps562570 & PROVED & \checkmark & verified \\
\texttt{evenPair\_card\_fiveHundredSixtyFour} & Gaps562570 & PROVED & \checkmark & verified \\
\texttt{evenPair\_card\_fiveHundredSixtySix} & Gaps562570 & PROVED & \checkmark & verified \\
\texttt{evenPair\_card\_fiveHundredSixtyTwo} & Gaps562570 & PROVED & \checkmark & verified \\
\texttt{isAdmissible\_evenPair\_fiveHundredSeventy} & Gaps562570 & PROVED & \checkmark & verified \\
\texttt{isAdmissible\_evenPair\_fiveHundredSixtyEight} & Gaps562570 & PROVED & \checkmark & verified \\
\texttt{isAdmissible\_evenPair\_fiveHundredSixtyFour} & Gaps562570 & PROVED & \checkmark & verified \\
\texttt{isAdmissible\_evenPair\_fiveHundredSixtySix} & Gaps562570 & PROVED & \checkmark & verified \\
\texttt{isAdmissible\_evenPair\_fiveHundredSixtyTwo} & Gaps562570 & PROVED & \checkmark & verified \\
\texttt{localFactor\_fiveHundredSeventy\_two} & Gaps562570 & PROVED & \checkmark & verified \\
\texttt{localFactor\_fiveHundredSixtyTwo\_two} & Gaps562570 & PROVED & \checkmark & verified \\
\texttt{nu\_p\_fiveHundredSeventy} & Gaps562570 & PROVED & \checkmark & verified \\
\texttt{nu\_p\_fiveHundredSeventy\_two} & Gaps562570 & PROVED & \checkmark & verified \\
\texttt{nu\_p\_fiveHundredSixtyEight} & Gaps562570 & PROVED & \checkmark & verified \\
\texttt{nu\_p\_fiveHundredSixtyFour} & Gaps562570 & PROVED & \checkmark & verified \\
\texttt{nu\_p\_fiveHundredSixtySix} & Gaps562570 & PROVED & \checkmark & verified \\
\texttt{nu\_p\_fiveHundredSixtyTwo} & Gaps562570 & PROVED & \checkmark & verified \\
\texttt{nu\_p\_fiveHundredSixtyTwo\_two} & Gaps562570 & PROVED & \checkmark & verified \\
\texttt{singular\_series\_finite\_pos\_evenPair\_fiveHundredSeventy} & Gaps562570 & PROVED & \checkmark & verified \\
\texttt{singular\_series\_finite\_pos\_evenPair\_fiveHundredSixtyEight} & Gaps562570 & PROVED & \checkmark & verified \\
\texttt{singular\_series\_finite\_pos\_evenPair\_fiveHundredSixtyFour} & Gaps562570 & PROVED & \checkmark & verified \\
\texttt{singular\_series\_finite\_pos\_evenPair\_fiveHundredSixtySix} & Gaps562570 & PROVED & \checkmark & verified \\
\texttt{singular\_series\_finite\_pos\_evenPair\_fiveHundredSixtyTwo} & Gaps562570 & PROVED & \checkmark & verified \\
\texttt{singular\_series\_pos\_evenPair\_fiveHundredSeventy} & Gaps562570 & PROVED & \checkmark & verified \\
\texttt{singular\_series\_pos\_evenPair\_fiveHundredSixtyEight} & Gaps562570 & PROVED & \checkmark & verified \\
\texttt{singular\_series\_pos\_evenPair\_fiveHundredSixtyFour} & Gaps562570 & PROVED & \checkmark & verified \\
\texttt{singular\_series\_pos\_evenPair\_fiveHundredSixtySix} & Gaps562570 & PROVED & \checkmark & verified \\
\texttt{singular\_series\_pos\_evenPair\_fiveHundredSixtyTwo} & Gaps562570 & PROVED & \checkmark & verified \\
\texttt{evenPair\_card\_fiveHundredEighty} & Gaps572580 & PROVED & \checkmark & verified \\
\texttt{evenPair\_card\_fiveHundredSeventyEight} & Gaps572580 & PROVED & \checkmark & verified \\
\texttt{evenPair\_card\_fiveHundredSeventyFour} & Gaps572580 & PROVED & \checkmark & verified \\
\texttt{evenPair\_card\_fiveHundredSeventySix} & Gaps572580 & PROVED & \checkmark & verified \\
\texttt{evenPair\_card\_fiveHundredSeventyTwo} & Gaps572580 & PROVED & \checkmark & verified \\
\texttt{isAdmissible\_evenPair\_fiveHundredEighty} & Gaps572580 & PROVED & \checkmark & verified \\
\texttt{isAdmissible\_evenPair\_fiveHundredSeventyEight} & Gaps572580 & PROVED & \checkmark & verified \\
\texttt{isAdmissible\_evenPair\_fiveHundredSeventyFour} & Gaps572580 & PROVED & \checkmark & verified \\
\texttt{isAdmissible\_evenPair\_fiveHundredSeventySix} & Gaps572580 & PROVED & \checkmark & verified \\
\texttt{isAdmissible\_evenPair\_fiveHundredSeventyTwo} & Gaps572580 & PROVED & \checkmark & verified \\
\texttt{localFactor\_fiveHundredEighty\_two} & Gaps572580 & PROVED & \checkmark & verified \\
\texttt{localFactor\_fiveHundredSeventyTwo\_two} & Gaps572580 & PROVED & \checkmark & verified \\
\texttt{nu\_p\_fiveHundredEighty} & Gaps572580 & PROVED & \checkmark & verified \\
\texttt{nu\_p\_fiveHundredEighty\_two} & Gaps572580 & PROVED & \checkmark & verified \\
\texttt{nu\_p\_fiveHundredSeventyEight} & Gaps572580 & PROVED & \checkmark & verified \\
\texttt{nu\_p\_fiveHundredSeventyFour} & Gaps572580 & PROVED & \checkmark & verified \\
\texttt{nu\_p\_fiveHundredSeventySix} & Gaps572580 & PROVED & \checkmark & verified \\
\texttt{nu\_p\_fiveHundredSeventyTwo} & Gaps572580 & PROVED & \checkmark & verified \\
\texttt{nu\_p\_fiveHundredSeventyTwo\_two} & Gaps572580 & PROVED & \checkmark & verified \\
\texttt{singular\_series\_finite\_pos\_evenPair\_fiveHundredEighty} & Gaps572580 & PROVED & \checkmark & verified \\
\texttt{singular\_series\_finite\_pos\_evenPair\_fiveHundredSeventyEight} & Gaps572580 & PROVED & \checkmark & verified \\
\texttt{singular\_series\_finite\_pos\_evenPair\_fiveHundredSeventyFour} & Gaps572580 & PROVED & \checkmark & verified \\
\texttt{singular\_series\_finite\_pos\_evenPair\_fiveHundredSeventySix} & Gaps572580 & PROVED & \checkmark & verified \\
\texttt{singular\_series\_finite\_pos\_evenPair\_fiveHundredSeventyTwo} & Gaps572580 & PROVED & \checkmark & verified \\
\texttt{singular\_series\_pos\_evenPair\_fiveHundredEighty} & Gaps572580 & PROVED & \checkmark & verified \\
\texttt{singular\_series\_pos\_evenPair\_fiveHundredSeventyEight} & Gaps572580 & PROVED & \checkmark & verified \\
\texttt{singular\_series\_pos\_evenPair\_fiveHundredSeventyFour} & Gaps572580 & PROVED & \checkmark & verified \\
\texttt{singular\_series\_pos\_evenPair\_fiveHundredSeventySix} & Gaps572580 & PROVED & \checkmark & verified \\
\texttt{singular\_series\_pos\_evenPair\_fiveHundredSeventyTwo} & Gaps572580 & PROVED & \checkmark & verified \\
\texttt{evenPair\_card\_fiveHundredEightyEight} & Gaps582590 & PROVED & \checkmark & verified \\
\texttt{evenPair\_card\_fiveHundredEightyFour} & Gaps582590 & PROVED & \checkmark & verified \\
\texttt{evenPair\_card\_fiveHundredEightySix} & Gaps582590 & PROVED & \checkmark & verified \\
\texttt{evenPair\_card\_fiveHundredEightyTwo} & Gaps582590 & PROVED & \checkmark & verified \\
\texttt{evenPair\_card\_fiveHundredNinety} & Gaps582590 & PROVED & \checkmark & verified \\
\texttt{isAdmissible\_evenPair\_fiveHundredEightyEight} & Gaps582590 & PROVED & \checkmark & verified \\
\texttt{isAdmissible\_evenPair\_fiveHundredEightyFour} & Gaps582590 & PROVED & \checkmark & verified \\
\texttt{isAdmissible\_evenPair\_fiveHundredEightySix} & Gaps582590 & PROVED & \checkmark & verified \\
\texttt{isAdmissible\_evenPair\_fiveHundredEightyTwo} & Gaps582590 & PROVED & \checkmark & verified \\
\texttt{isAdmissible\_evenPair\_fiveHundredNinety} & Gaps582590 & PROVED & \checkmark & verified \\
\texttt{localFactor\_fiveHundredEightyTwo\_two} & Gaps582590 & PROVED & \checkmark & verified \\
\texttt{localFactor\_fiveHundredNinety\_two} & Gaps582590 & PROVED & \checkmark & verified \\
\texttt{nu\_p\_fiveHundredEightyEight} & Gaps582590 & PROVED & \checkmark & verified \\
\texttt{nu\_p\_fiveHundredEightyFour} & Gaps582590 & PROVED & \checkmark & verified \\
\texttt{nu\_p\_fiveHundredEightySix} & Gaps582590 & PROVED & \checkmark & verified \\
\texttt{nu\_p\_fiveHundredEightyTwo} & Gaps582590 & PROVED & \checkmark & verified \\
\texttt{nu\_p\_fiveHundredEightyTwo\_two} & Gaps582590 & PROVED & \checkmark & verified \\
\texttt{nu\_p\_fiveHundredNinety} & Gaps582590 & PROVED & \checkmark & verified \\
\texttt{nu\_p\_fiveHundredNinety\_two} & Gaps582590 & PROVED & \checkmark & verified \\
\texttt{singular\_series\_finite\_pos\_evenPair\_fiveHundredEightyEight} & Gaps582590 & PROVED & \checkmark & verified \\
\texttt{singular\_series\_finite\_pos\_evenPair\_fiveHundredEightyFour} & Gaps582590 & PROVED & \checkmark & verified \\
\texttt{singular\_series\_finite\_pos\_evenPair\_fiveHundredEightySix} & Gaps582590 & PROVED & \checkmark & verified \\
\texttt{singular\_series\_finite\_pos\_evenPair\_fiveHundredEightyTwo} & Gaps582590 & PROVED & \checkmark & verified \\
\texttt{singular\_series\_finite\_pos\_evenPair\_fiveHundredNinety} & Gaps582590 & PROVED & \checkmark & verified \\
\texttt{singular\_series\_pos\_evenPair\_fiveHundredEightyEight} & Gaps582590 & PROVED & \checkmark & verified \\
\texttt{singular\_series\_pos\_evenPair\_fiveHundredEightyFour} & Gaps582590 & PROVED & \checkmark & verified \\
\texttt{singular\_series\_pos\_evenPair\_fiveHundredEightySix} & Gaps582590 & PROVED & \checkmark & verified \\
\texttt{singular\_series\_pos\_evenPair\_fiveHundredEightyTwo} & Gaps582590 & PROVED & \checkmark & verified \\
\texttt{singular\_series\_pos\_evenPair\_fiveHundredNinety} & Gaps582590 & PROVED & \checkmark & verified \\
\texttt{evenPair\_card\_fiveHundredNinetyEight} & Gaps592600 & PROVED & \checkmark & verified \\
\texttt{evenPair\_card\_fiveHundredNinetyFour} & Gaps592600 & PROVED & \checkmark & verified \\
\texttt{evenPair\_card\_fiveHundredNinetySix} & Gaps592600 & PROVED & \checkmark & verified \\
\texttt{evenPair\_card\_fiveHundredNinetyTwo} & Gaps592600 & PROVED & \checkmark & verified \\
\texttt{evenPair\_card\_sixHundred} & Gaps592600 & PROVED & \checkmark & verified \\
\texttt{isAdmissible\_evenPair\_fiveHundredNinetyEight} & Gaps592600 & PROVED & \checkmark & verified \\
\texttt{isAdmissible\_evenPair\_fiveHundredNinetyFour} & Gaps592600 & PROVED & \checkmark & verified \\
\texttt{isAdmissible\_evenPair\_fiveHundredNinetySix} & Gaps592600 & PROVED & \checkmark & verified \\
\texttt{isAdmissible\_evenPair\_fiveHundredNinetyTwo} & Gaps592600 & PROVED & \checkmark & verified \\
\texttt{isAdmissible\_evenPair\_sixHundred} & Gaps592600 & PROVED & \checkmark & verified \\
\texttt{localFactor\_fiveHundredNinetyTwo\_two} & Gaps592600 & PROVED & \checkmark & verified \\
\texttt{localFactor\_sixHundred\_two} & Gaps592600 & PROVED & \checkmark & verified \\
\texttt{nu\_p\_fiveHundredNinetyEight} & Gaps592600 & PROVED & \checkmark & verified \\
\texttt{nu\_p\_fiveHundredNinetyFour} & Gaps592600 & PROVED & \checkmark & verified \\
\texttt{nu\_p\_fiveHundredNinetySix} & Gaps592600 & PROVED & \checkmark & verified \\
\texttt{nu\_p\_fiveHundredNinetyTwo} & Gaps592600 & PROVED & \checkmark & verified \\
\texttt{nu\_p\_fiveHundredNinetyTwo\_two} & Gaps592600 & PROVED & \checkmark & verified \\
\texttt{nu\_p\_sixHundred} & Gaps592600 & PROVED & \checkmark & verified \\
\texttt{nu\_p\_sixHundred\_two} & Gaps592600 & PROVED & \checkmark & verified \\
\texttt{singular\_series\_finite\_pos\_evenPair\_fiveHundredNinetyEight} & Gaps592600 & PROVED & \checkmark & verified \\
\texttt{singular\_series\_finite\_pos\_evenPair\_fiveHundredNinetyFour} & Gaps592600 & PROVED & \checkmark & verified \\
\texttt{singular\_series\_finite\_pos\_evenPair\_fiveHundredNinetySix} & Gaps592600 & PROVED & \checkmark & verified \\
\texttt{singular\_series\_finite\_pos\_evenPair\_fiveHundredNinetyTwo} & Gaps592600 & PROVED & \checkmark & verified \\
\texttt{singular\_series\_finite\_pos\_evenPair\_sixHundred} & Gaps592600 & PROVED & \checkmark & verified \\
\texttt{singular\_series\_pos\_evenPair\_fiveHundredNinetyEight} & Gaps592600 & PROVED & \checkmark & verified \\
\texttt{singular\_series\_pos\_evenPair\_fiveHundredNinetyFour} & Gaps592600 & PROVED & \checkmark & verified \\
\texttt{singular\_series\_pos\_evenPair\_fiveHundredNinetySix} & Gaps592600 & PROVED & \checkmark & verified \\
\texttt{singular\_series\_pos\_evenPair\_fiveHundredNinetyTwo} & Gaps592600 & PROVED & \checkmark & verified \\
\texttt{singular\_series\_pos\_evenPair\_sixHundred} & Gaps592600 & PROVED & \checkmark & verified \\
\texttt{evenPair\_card\_sixHundredEight} & Gaps602610 & PROVED & \checkmark & verified \\
\texttt{evenPair\_card\_sixHundredFour} & Gaps602610 & PROVED & \checkmark & verified \\
\texttt{evenPair\_card\_sixHundredSix} & Gaps602610 & PROVED & \checkmark & verified \\
\texttt{evenPair\_card\_sixHundredTen} & Gaps602610 & PROVED & \checkmark & verified \\
\texttt{evenPair\_card\_sixHundredTwo} & Gaps602610 & PROVED & \checkmark & verified \\
\texttt{isAdmissible\_evenPair\_sixHundredEight} & Gaps602610 & PROVED & \checkmark & verified \\
\texttt{isAdmissible\_evenPair\_sixHundredFour} & Gaps602610 & PROVED & \checkmark & verified \\
\texttt{isAdmissible\_evenPair\_sixHundredSix} & Gaps602610 & PROVED & \checkmark & verified \\
\texttt{isAdmissible\_evenPair\_sixHundredTen} & Gaps602610 & PROVED & \checkmark & verified \\
\texttt{isAdmissible\_evenPair\_sixHundredTwo} & Gaps602610 & PROVED & \checkmark & verified \\
\texttt{localFactor\_sixHundredTen\_two} & Gaps602610 & PROVED & \checkmark & verified \\
\texttt{localFactor\_sixHundredTwo\_two} & Gaps602610 & PROVED & \checkmark & verified \\
\texttt{nu\_p\_sixHundredEight} & Gaps602610 & PROVED & \checkmark & verified \\
\texttt{nu\_p\_sixHundredFour} & Gaps602610 & PROVED & \checkmark & verified \\
\texttt{nu\_p\_sixHundredSix} & Gaps602610 & PROVED & \checkmark & verified \\
\texttt{nu\_p\_sixHundredTen} & Gaps602610 & PROVED & \checkmark & verified \\
\texttt{nu\_p\_sixHundredTen\_two} & Gaps602610 & PROVED & \checkmark & verified \\
\texttt{nu\_p\_sixHundredTwo} & Gaps602610 & PROVED & \checkmark & verified \\
\texttt{nu\_p\_sixHundredTwo\_two} & Gaps602610 & PROVED & \checkmark & verified \\
\texttt{singular\_series\_finite\_pos\_evenPair\_sixHundredEight} & Gaps602610 & PROVED & \checkmark & verified \\
\texttt{singular\_series\_finite\_pos\_evenPair\_sixHundredFour} & Gaps602610 & PROVED & \checkmark & verified \\
\texttt{singular\_series\_finite\_pos\_evenPair\_sixHundredSix} & Gaps602610 & PROVED & \checkmark & verified \\
\texttt{singular\_series\_finite\_pos\_evenPair\_sixHundredTen} & Gaps602610 & PROVED & \checkmark & verified \\
\texttt{singular\_series\_finite\_pos\_evenPair\_sixHundredTwo} & Gaps602610 & PROVED & \checkmark & verified \\
\texttt{singular\_series\_pos\_evenPair\_sixHundredEight} & Gaps602610 & PROVED & \checkmark & verified \\
\texttt{singular\_series\_pos\_evenPair\_sixHundredFour} & Gaps602610 & PROVED & \checkmark & verified \\
\texttt{singular\_series\_pos\_evenPair\_sixHundredSix} & Gaps602610 & PROVED & \checkmark & verified \\
\texttt{singular\_series\_pos\_evenPair\_sixHundredTen} & Gaps602610 & PROVED & \checkmark & verified \\
\texttt{singular\_series\_pos\_evenPair\_sixHundredTwo} & Gaps602610 & PROVED & \checkmark & verified \\
\texttt{evenPair\_card\_sixHundredEighteen} & Gaps612620 & PROVED & \checkmark & verified \\
\texttt{evenPair\_card\_sixHundredFourteen} & Gaps612620 & PROVED & \checkmark & verified \\
\texttt{evenPair\_card\_sixHundredSixteen} & Gaps612620 & PROVED & \checkmark & verified \\
\texttt{evenPair\_card\_sixHundredTwelve} & Gaps612620 & PROVED & \checkmark & verified \\
\texttt{evenPair\_card\_sixHundredTwenty} & Gaps612620 & PROVED & \checkmark & verified \\
\texttt{isAdmissible\_evenPair\_sixHundredEighteen} & Gaps612620 & PROVED & \checkmark & verified \\
\texttt{isAdmissible\_evenPair\_sixHundredFourteen} & Gaps612620 & PROVED & \checkmark & verified \\
\texttt{isAdmissible\_evenPair\_sixHundredSixteen} & Gaps612620 & PROVED & \checkmark & verified \\
\texttt{isAdmissible\_evenPair\_sixHundredTwelve} & Gaps612620 & PROVED & \checkmark & verified \\
\texttt{isAdmissible\_evenPair\_sixHundredTwenty} & Gaps612620 & PROVED & \checkmark & verified \\
\texttt{localFactor\_sixHundredTwelve\_two} & Gaps612620 & PROVED & \checkmark & verified \\
\texttt{localFactor\_sixHundredTwenty\_two} & Gaps612620 & PROVED & \checkmark & verified \\
\texttt{nu\_p\_sixHundredEighteen} & Gaps612620 & PROVED & \checkmark & verified \\
\texttt{nu\_p\_sixHundredFourteen} & Gaps612620 & PROVED & \checkmark & verified \\
\texttt{nu\_p\_sixHundredSixteen} & Gaps612620 & PROVED & \checkmark & verified \\
\texttt{nu\_p\_sixHundredTwelve} & Gaps612620 & PROVED & \checkmark & verified \\
\texttt{nu\_p\_sixHundredTwelve\_two} & Gaps612620 & PROVED & \checkmark & verified \\
\texttt{nu\_p\_sixHundredTwenty} & Gaps612620 & PROVED & \checkmark & verified \\
\texttt{nu\_p\_sixHundredTwenty\_two} & Gaps612620 & PROVED & \checkmark & verified \\
\texttt{singular\_series\_finite\_pos\_evenPair\_sixHundredEighteen} & Gaps612620 & PROVED & \checkmark & verified \\
\texttt{singular\_series\_finite\_pos\_evenPair\_sixHundredFourteen} & Gaps612620 & PROVED & \checkmark & verified \\
\texttt{singular\_series\_finite\_pos\_evenPair\_sixHundredSixteen} & Gaps612620 & PROVED & \checkmark & verified \\
\texttt{singular\_series\_finite\_pos\_evenPair\_sixHundredTwelve} & Gaps612620 & PROVED & \checkmark & verified \\
\texttt{singular\_series\_finite\_pos\_evenPair\_sixHundredTwenty} & Gaps612620 & PROVED & \checkmark & verified \\
\texttt{singular\_series\_pos\_evenPair\_sixHundredEighteen} & Gaps612620 & PROVED & \checkmark & verified \\
\texttt{singular\_series\_pos\_evenPair\_sixHundredFourteen} & Gaps612620 & PROVED & \checkmark & verified \\
\texttt{singular\_series\_pos\_evenPair\_sixHundredSixteen} & Gaps612620 & PROVED & \checkmark & verified \\
\texttt{singular\_series\_pos\_evenPair\_sixHundredTwelve} & Gaps612620 & PROVED & \checkmark & verified \\
\texttt{singular\_series\_pos\_evenPair\_sixHundredTwenty} & Gaps612620 & PROVED & \checkmark & verified \\
\texttt{evenPair\_card\_sixHundredThirty} & Gaps622630 & PROVED & \checkmark & verified \\
\texttt{evenPair\_card\_sixHundredTwentyEight} & Gaps622630 & PROVED & \checkmark & verified \\
\texttt{evenPair\_card\_sixHundredTwentyFour} & Gaps622630 & PROVED & \checkmark & verified \\
\texttt{evenPair\_card\_sixHundredTwentySix} & Gaps622630 & PROVED & \checkmark & verified \\
\texttt{evenPair\_card\_sixHundredTwentyTwo} & Gaps622630 & PROVED & \checkmark & verified \\
\texttt{isAdmissible\_evenPair\_sixHundredThirty} & Gaps622630 & PROVED & \checkmark & verified \\
\texttt{isAdmissible\_evenPair\_sixHundredTwentyEight} & Gaps622630 & PROVED & \checkmark & verified \\
\texttt{isAdmissible\_evenPair\_sixHundredTwentyFour} & Gaps622630 & PROVED & \checkmark & verified \\
\texttt{isAdmissible\_evenPair\_sixHundredTwentySix} & Gaps622630 & PROVED & \checkmark & verified \\
\texttt{isAdmissible\_evenPair\_sixHundredTwentyTwo} & Gaps622630 & PROVED & \checkmark & verified \\
\texttt{localFactor\_sixHundredThirty\_two} & Gaps622630 & PROVED & \checkmark & verified \\
\texttt{localFactor\_sixHundredTwentyTwo\_two} & Gaps622630 & PROVED & \checkmark & verified \\
\texttt{nu\_p\_sixHundredThirty} & Gaps622630 & PROVED & \checkmark & verified \\
\texttt{nu\_p\_sixHundredThirty\_two} & Gaps622630 & PROVED & \checkmark & verified \\
\texttt{nu\_p\_sixHundredTwentyEight} & Gaps622630 & PROVED & \checkmark & verified \\
\texttt{nu\_p\_sixHundredTwentyFour} & Gaps622630 & PROVED & \checkmark & verified \\
\texttt{nu\_p\_sixHundredTwentySix} & Gaps622630 & PROVED & \checkmark & verified \\
\texttt{nu\_p\_sixHundredTwentyTwo} & Gaps622630 & PROVED & \checkmark & verified \\
\texttt{nu\_p\_sixHundredTwentyTwo\_two} & Gaps622630 & PROVED & \checkmark & verified \\
\texttt{singular\_series\_finite\_pos\_evenPair\_sixHundredThirty} & Gaps622630 & PROVED & \checkmark & verified \\
\texttt{singular\_series\_finite\_pos\_evenPair\_sixHundredTwentyEight} & Gaps622630 & PROVED & \checkmark & verified \\
\texttt{singular\_series\_finite\_pos\_evenPair\_sixHundredTwentyFour} & Gaps622630 & PROVED & \checkmark & verified \\
\texttt{singular\_series\_finite\_pos\_evenPair\_sixHundredTwentySix} & Gaps622630 & PROVED & \checkmark & verified \\
\texttt{singular\_series\_finite\_pos\_evenPair\_sixHundredTwentyTwo} & Gaps622630 & PROVED & \checkmark & verified \\
\texttt{singular\_series\_pos\_evenPair\_sixHundredThirty} & Gaps622630 & PROVED & \checkmark & verified \\
\texttt{singular\_series\_pos\_evenPair\_sixHundredTwentyEight} & Gaps622630 & PROVED & \checkmark & verified \\
\texttt{singular\_series\_pos\_evenPair\_sixHundredTwentyFour} & Gaps622630 & PROVED & \checkmark & verified \\
\texttt{singular\_series\_pos\_evenPair\_sixHundredTwentySix} & Gaps622630 & PROVED & \checkmark & verified \\
\texttt{singular\_series\_pos\_evenPair\_sixHundredTwentyTwo} & Gaps622630 & PROVED & \checkmark & verified \\
\texttt{evenPair\_card\_seventy} & Gaps6270 & PROVED & \checkmark & verified \\
\texttt{evenPair\_card\_sixtyEight} & Gaps6270 & PROVED & \checkmark & verified \\
\texttt{evenPair\_card\_sixtyFour} & Gaps6270 & PROVED & \checkmark & verified \\
\texttt{evenPair\_card\_sixtySix} & Gaps6270 & PROVED & \checkmark & verified \\
\texttt{evenPair\_card\_sixtyTwo} & Gaps6270 & PROVED & \checkmark & verified \\
\texttt{isAdmissible\_evenPair\_seventy} & Gaps6270 & PROVED & \checkmark & verified \\
\texttt{isAdmissible\_evenPair\_sixtyEight} & Gaps6270 & PROVED & \checkmark & verified \\
\texttt{isAdmissible\_evenPair\_sixtyFour} & Gaps6270 & PROVED & \checkmark & verified \\
\texttt{isAdmissible\_evenPair\_sixtySix} & Gaps6270 & PROVED & \checkmark & verified \\
\texttt{isAdmissible\_evenPair\_sixtyTwo} & Gaps6270 & PROVED & \checkmark & verified \\
\texttt{localFactor\_seventy\_two} & Gaps6270 & PROVED & \checkmark & verified \\
\texttt{localFactor\_sixtyTwo\_two} & Gaps6270 & PROVED & \checkmark & verified \\
\texttt{nu\_p\_seventy} & Gaps6270 & PROVED & \checkmark & verified \\
\texttt{nu\_p\_seventy\_two} & Gaps6270 & PROVED & \checkmark & verified \\
\texttt{nu\_p\_sixtyEight} & Gaps6270 & PROVED & \checkmark & verified \\
\texttt{nu\_p\_sixtyFour} & Gaps6270 & PROVED & \checkmark & verified \\
\texttt{nu\_p\_sixtySix} & Gaps6270 & PROVED & \checkmark & verified \\
\texttt{nu\_p\_sixtyTwo} & Gaps6270 & PROVED & \checkmark & verified \\
\texttt{nu\_p\_sixtyTwo\_two} & Gaps6270 & PROVED & \checkmark & verified \\
\texttt{singular\_series\_finite\_pos\_evenPair\_seventy} & Gaps6270 & PROVED & \checkmark & verified \\
\texttt{singular\_series\_finite\_pos\_evenPair\_sixtyEight} & Gaps6270 & PROVED & \checkmark & verified \\
\texttt{singular\_series\_finite\_pos\_evenPair\_sixtyFour} & Gaps6270 & PROVED & \checkmark & verified \\
\texttt{singular\_series\_finite\_pos\_evenPair\_sixtySix} & Gaps6270 & PROVED & \checkmark & verified \\
\texttt{singular\_series\_finite\_pos\_evenPair\_sixtyTwo} & Gaps6270 & PROVED & \checkmark & verified \\
\texttt{singular\_series\_pos\_evenPair\_seventy} & Gaps6270 & PROVED & \checkmark & verified \\
\texttt{singular\_series\_pos\_evenPair\_sixtyEight} & Gaps6270 & PROVED & \checkmark & verified \\
\texttt{singular\_series\_pos\_evenPair\_sixtyFour} & Gaps6270 & PROVED & \checkmark & verified \\
\texttt{singular\_series\_pos\_evenPair\_sixtySix} & Gaps6270 & PROVED & \checkmark & verified \\
\texttt{singular\_series\_pos\_evenPair\_sixtyTwo} & Gaps6270 & PROVED & \checkmark & verified \\
\texttt{evenPair\_card\_sixHundredForty} & Gaps632640 & PROVED & \checkmark & verified \\
\texttt{evenPair\_card\_sixHundredThirtyEight} & Gaps632640 & PROVED & \checkmark & verified \\
\texttt{evenPair\_card\_sixHundredThirtyFour} & Gaps632640 & PROVED & \checkmark & verified \\
\texttt{evenPair\_card\_sixHundredThirtySix} & Gaps632640 & PROVED & \checkmark & verified \\
\texttt{evenPair\_card\_sixHundredThirtyTwo} & Gaps632640 & PROVED & \checkmark & verified \\
\texttt{isAdmissible\_evenPair\_sixHundredForty} & Gaps632640 & PROVED & \checkmark & verified \\
\texttt{isAdmissible\_evenPair\_sixHundredThirtyEight} & Gaps632640 & PROVED & \checkmark & verified \\
\texttt{isAdmissible\_evenPair\_sixHundredThirtyFour} & Gaps632640 & PROVED & \checkmark & verified \\
\texttt{isAdmissible\_evenPair\_sixHundredThirtySix} & Gaps632640 & PROVED & \checkmark & verified \\
\texttt{isAdmissible\_evenPair\_sixHundredThirtyTwo} & Gaps632640 & PROVED & \checkmark & verified \\
\texttt{localFactor\_sixHundredForty\_two} & Gaps632640 & PROVED & \checkmark & verified \\
\texttt{localFactor\_sixHundredThirtyTwo\_two} & Gaps632640 & PROVED & \checkmark & verified \\
\texttt{nu\_p\_sixHundredForty} & Gaps632640 & PROVED & \checkmark & verified \\
\texttt{nu\_p\_sixHundredForty\_two} & Gaps632640 & PROVED & \checkmark & verified \\
\texttt{nu\_p\_sixHundredThirtyEight} & Gaps632640 & PROVED & \checkmark & verified \\
\texttt{nu\_p\_sixHundredThirtyFour} & Gaps632640 & PROVED & \checkmark & verified \\
\texttt{nu\_p\_sixHundredThirtySix} & Gaps632640 & PROVED & \checkmark & verified \\
\texttt{nu\_p\_sixHundredThirtyTwo} & Gaps632640 & PROVED & \checkmark & verified \\
\texttt{nu\_p\_sixHundredThirtyTwo\_two} & Gaps632640 & PROVED & \checkmark & verified \\
\texttt{singular\_series\_finite\_pos\_evenPair\_sixHundredForty} & Gaps632640 & PROVED & \checkmark & verified \\
\texttt{singular\_series\_finite\_pos\_evenPair\_sixHundredThirtyEight} & Gaps632640 & PROVED & \checkmark & verified \\
\texttt{singular\_series\_finite\_pos\_evenPair\_sixHundredThirtyFour} & Gaps632640 & PROVED & \checkmark & verified \\
\texttt{singular\_series\_finite\_pos\_evenPair\_sixHundredThirtySix} & Gaps632640 & PROVED & \checkmark & verified \\
\texttt{singular\_series\_finite\_pos\_evenPair\_sixHundredThirtyTwo} & Gaps632640 & PROVED & \checkmark & verified \\
\texttt{singular\_series\_pos\_evenPair\_sixHundredForty} & Gaps632640 & PROVED & \checkmark & verified \\
\texttt{singular\_series\_pos\_evenPair\_sixHundredThirtyEight} & Gaps632640 & PROVED & \checkmark & verified \\
\texttt{singular\_series\_pos\_evenPair\_sixHundredThirtyFour} & Gaps632640 & PROVED & \checkmark & verified \\
\texttt{singular\_series\_pos\_evenPair\_sixHundredThirtySix} & Gaps632640 & PROVED & \checkmark & verified \\
\texttt{singular\_series\_pos\_evenPair\_sixHundredThirtyTwo} & Gaps632640 & PROVED & \checkmark & verified \\
\texttt{evenPair\_card\_sixHundredFifty} & Gaps642650 & PROVED & \checkmark & verified \\
\texttt{evenPair\_card\_sixHundredFortyEight} & Gaps642650 & PROVED & \checkmark & verified \\
\texttt{evenPair\_card\_sixHundredFortyFour} & Gaps642650 & PROVED & \checkmark & verified \\
\texttt{evenPair\_card\_sixHundredFortySix} & Gaps642650 & PROVED & \checkmark & verified \\
\texttt{evenPair\_card\_sixHundredFortyTwo} & Gaps642650 & PROVED & \checkmark & verified \\
\texttt{isAdmissible\_evenPair\_sixHundredFifty} & Gaps642650 & PROVED & \checkmark & verified \\
\texttt{isAdmissible\_evenPair\_sixHundredFortyEight} & Gaps642650 & PROVED & \checkmark & verified \\
\texttt{isAdmissible\_evenPair\_sixHundredFortyFour} & Gaps642650 & PROVED & \checkmark & verified \\
\texttt{isAdmissible\_evenPair\_sixHundredFortySix} & Gaps642650 & PROVED & \checkmark & verified \\
\texttt{isAdmissible\_evenPair\_sixHundredFortyTwo} & Gaps642650 & PROVED & \checkmark & verified \\
\texttt{localFactor\_sixHundredFifty\_two} & Gaps642650 & PROVED & \checkmark & verified \\
\texttt{localFactor\_sixHundredFortyTwo\_two} & Gaps642650 & PROVED & \checkmark & verified \\
\texttt{nu\_p\_sixHundredFifty} & Gaps642650 & PROVED & \checkmark & verified \\
\texttt{nu\_p\_sixHundredFifty\_two} & Gaps642650 & PROVED & \checkmark & verified \\
\texttt{nu\_p\_sixHundredFortyEight} & Gaps642650 & PROVED & \checkmark & verified \\
\texttt{nu\_p\_sixHundredFortyFour} & Gaps642650 & PROVED & \checkmark & verified \\
\texttt{nu\_p\_sixHundredFortySix} & Gaps642650 & PROVED & \checkmark & verified \\
\texttt{nu\_p\_sixHundredFortyTwo} & Gaps642650 & PROVED & \checkmark & verified \\
\texttt{nu\_p\_sixHundredFortyTwo\_two} & Gaps642650 & PROVED & \checkmark & verified \\
\texttt{singular\_series\_finite\_pos\_evenPair\_sixHundredFifty} & Gaps642650 & PROVED & \checkmark & verified \\
\texttt{singular\_series\_finite\_pos\_evenPair\_sixHundredFortyEight} & Gaps642650 & PROVED & \checkmark & verified \\
\texttt{singular\_series\_finite\_pos\_evenPair\_sixHundredFortyFour} & Gaps642650 & PROVED & \checkmark & verified \\
\texttt{singular\_series\_finite\_pos\_evenPair\_sixHundredFortySix} & Gaps642650 & PROVED & \checkmark & verified \\
\texttt{singular\_series\_finite\_pos\_evenPair\_sixHundredFortyTwo} & Gaps642650 & PROVED & \checkmark & verified \\
\texttt{singular\_series\_pos\_evenPair\_sixHundredFifty} & Gaps642650 & PROVED & \checkmark & verified \\
\texttt{singular\_series\_pos\_evenPair\_sixHundredFortyEight} & Gaps642650 & PROVED & \checkmark & verified \\
\texttt{singular\_series\_pos\_evenPair\_sixHundredFortyFour} & Gaps642650 & PROVED & \checkmark & verified \\
\texttt{singular\_series\_pos\_evenPair\_sixHundredFortySix} & Gaps642650 & PROVED & \checkmark & verified \\
\texttt{singular\_series\_pos\_evenPair\_sixHundredFortyTwo} & Gaps642650 & PROVED & \checkmark & verified \\
\texttt{evenPair\_card\_sixHundredFiftyEight} & Gaps652660 & PROVED & \checkmark & verified \\
\texttt{evenPair\_card\_sixHundredFiftyFour} & Gaps652660 & PROVED & \checkmark & verified \\
\texttt{evenPair\_card\_sixHundredFiftySix} & Gaps652660 & PROVED & \checkmark & verified \\
\texttt{evenPair\_card\_sixHundredFiftyTwo} & Gaps652660 & PROVED & \checkmark & verified \\
\texttt{evenPair\_card\_sixHundredSixty} & Gaps652660 & PROVED & \checkmark & verified \\
\texttt{isAdmissible\_evenPair\_sixHundredFiftyEight} & Gaps652660 & PROVED & \checkmark & verified \\
\texttt{isAdmissible\_evenPair\_sixHundredFiftyFour} & Gaps652660 & PROVED & \checkmark & verified \\
\texttt{isAdmissible\_evenPair\_sixHundredFiftySix} & Gaps652660 & PROVED & \checkmark & verified \\
\texttt{isAdmissible\_evenPair\_sixHundredFiftyTwo} & Gaps652660 & PROVED & \checkmark & verified \\
\texttt{isAdmissible\_evenPair\_sixHundredSixty} & Gaps652660 & PROVED & \checkmark & verified \\
\texttt{localFactor\_sixHundredFiftyTwo\_two} & Gaps652660 & PROVED & \checkmark & verified \\
\texttt{localFactor\_sixHundredSixty\_two} & Gaps652660 & PROVED & \checkmark & verified \\
\texttt{nu\_p\_sixHundredFiftyEight} & Gaps652660 & PROVED & \checkmark & verified \\
\texttt{nu\_p\_sixHundredFiftyFour} & Gaps652660 & PROVED & \checkmark & verified \\
\texttt{nu\_p\_sixHundredFiftySix} & Gaps652660 & PROVED & \checkmark & verified \\
\texttt{nu\_p\_sixHundredFiftyTwo} & Gaps652660 & PROVED & \checkmark & verified \\
\texttt{nu\_p\_sixHundredFiftyTwo\_two} & Gaps652660 & PROVED & \checkmark & verified \\
\texttt{nu\_p\_sixHundredSixty} & Gaps652660 & PROVED & \checkmark & verified \\
\texttt{nu\_p\_sixHundredSixty\_two} & Gaps652660 & PROVED & \checkmark & verified \\
\texttt{singular\_series\_finite\_pos\_evenPair\_sixHundredFiftyEight} & Gaps652660 & PROVED & \checkmark & verified \\
\texttt{singular\_series\_finite\_pos\_evenPair\_sixHundredFiftyFour} & Gaps652660 & PROVED & \checkmark & verified \\
\texttt{singular\_series\_finite\_pos\_evenPair\_sixHundredFiftySix} & Gaps652660 & PROVED & \checkmark & verified \\
\texttt{singular\_series\_finite\_pos\_evenPair\_sixHundredFiftyTwo} & Gaps652660 & PROVED & \checkmark & verified \\
\texttt{singular\_series\_finite\_pos\_evenPair\_sixHundredSixty} & Gaps652660 & PROVED & \checkmark & verified \\
\texttt{singular\_series\_pos\_evenPair\_sixHundredFiftyEight} & Gaps652660 & PROVED & \checkmark & verified \\
\texttt{singular\_series\_pos\_evenPair\_sixHundredFiftyFour} & Gaps652660 & PROVED & \checkmark & verified \\
\texttt{singular\_series\_pos\_evenPair\_sixHundredFiftySix} & Gaps652660 & PROVED & \checkmark & verified \\
\texttt{singular\_series\_pos\_evenPair\_sixHundredFiftyTwo} & Gaps652660 & PROVED & \checkmark & verified \\
\texttt{singular\_series\_pos\_evenPair\_sixHundredSixty} & Gaps652660 & PROVED & \checkmark & verified \\
\texttt{evenPair\_card\_sixHundredSeventy} & Gaps662670 & PROVED & \checkmark & verified \\
\texttt{evenPair\_card\_sixHundredSixtyEight} & Gaps662670 & PROVED & \checkmark & verified \\
\texttt{evenPair\_card\_sixHundredSixtyFour} & Gaps662670 & PROVED & \checkmark & verified \\
\texttt{evenPair\_card\_sixHundredSixtySix} & Gaps662670 & PROVED & \checkmark & verified \\
\texttt{evenPair\_card\_sixHundredSixtyTwo} & Gaps662670 & PROVED & \checkmark & verified \\
\texttt{isAdmissible\_evenPair\_sixHundredSeventy} & Gaps662670 & PROVED & \checkmark & verified \\
\texttt{isAdmissible\_evenPair\_sixHundredSixtyEight} & Gaps662670 & PROVED & \checkmark & verified \\
\texttt{isAdmissible\_evenPair\_sixHundredSixtyFour} & Gaps662670 & PROVED & \checkmark & verified \\
\texttt{isAdmissible\_evenPair\_sixHundredSixtySix} & Gaps662670 & PROVED & \checkmark & verified \\
\texttt{isAdmissible\_evenPair\_sixHundredSixtyTwo} & Gaps662670 & PROVED & \checkmark & verified \\
\texttt{localFactor\_sixHundredSeventy\_two} & Gaps662670 & PROVED & \checkmark & verified \\
\texttt{localFactor\_sixHundredSixtyTwo\_two} & Gaps662670 & PROVED & \checkmark & verified \\
\texttt{nu\_p\_sixHundredSeventy} & Gaps662670 & PROVED & \checkmark & verified \\
\texttt{nu\_p\_sixHundredSeventy\_two} & Gaps662670 & PROVED & \checkmark & verified \\
\texttt{nu\_p\_sixHundredSixtyEight} & Gaps662670 & PROVED & \checkmark & verified \\
\texttt{nu\_p\_sixHundredSixtyFour} & Gaps662670 & PROVED & \checkmark & verified \\
\texttt{nu\_p\_sixHundredSixtySix} & Gaps662670 & PROVED & \checkmark & verified \\
\texttt{nu\_p\_sixHundredSixtyTwo} & Gaps662670 & PROVED & \checkmark & verified \\
\texttt{nu\_p\_sixHundredSixtyTwo\_two} & Gaps662670 & PROVED & \checkmark & verified \\
\texttt{singular\_series\_finite\_pos\_evenPair\_sixHundredSeventy} & Gaps662670 & PROVED & \checkmark & verified \\
\texttt{singular\_series\_finite\_pos\_evenPair\_sixHundredSixtyEight} & Gaps662670 & PROVED & \checkmark & verified \\
\texttt{singular\_series\_finite\_pos\_evenPair\_sixHundredSixtyFour} & Gaps662670 & PROVED & \checkmark & verified \\
\texttt{singular\_series\_finite\_pos\_evenPair\_sixHundredSixtySix} & Gaps662670 & PROVED & \checkmark & verified \\
\texttt{singular\_series\_finite\_pos\_evenPair\_sixHundredSixtyTwo} & Gaps662670 & PROVED & \checkmark & verified \\
\texttt{singular\_series\_pos\_evenPair\_sixHundredSeventy} & Gaps662670 & PROVED & \checkmark & verified \\
\texttt{singular\_series\_pos\_evenPair\_sixHundredSixtyEight} & Gaps662670 & PROVED & \checkmark & verified \\
\texttt{singular\_series\_pos\_evenPair\_sixHundredSixtyFour} & Gaps662670 & PROVED & \checkmark & verified \\
\texttt{singular\_series\_pos\_evenPair\_sixHundredSixtySix} & Gaps662670 & PROVED & \checkmark & verified \\
\texttt{singular\_series\_pos\_evenPair\_sixHundredSixtyTwo} & Gaps662670 & PROVED & \checkmark & verified \\
\texttt{evenPair\_card\_sixHundredEighty} & Gaps672680 & PROVED & \checkmark & verified \\
\texttt{evenPair\_card\_sixHundredSeventyEight} & Gaps672680 & PROVED & \checkmark & verified \\
\texttt{evenPair\_card\_sixHundredSeventyFour} & Gaps672680 & PROVED & \checkmark & verified \\
\texttt{evenPair\_card\_sixHundredSeventySix} & Gaps672680 & PROVED & \checkmark & verified \\
\texttt{evenPair\_card\_sixHundredSeventyTwo} & Gaps672680 & PROVED & \checkmark & verified \\
\texttt{isAdmissible\_evenPair\_sixHundredEighty} & Gaps672680 & PROVED & \checkmark & verified \\
\texttt{isAdmissible\_evenPair\_sixHundredSeventyEight} & Gaps672680 & PROVED & \checkmark & verified \\
\texttt{isAdmissible\_evenPair\_sixHundredSeventyFour} & Gaps672680 & PROVED & \checkmark & verified \\
\texttt{isAdmissible\_evenPair\_sixHundredSeventySix} & Gaps672680 & PROVED & \checkmark & verified \\
\texttt{isAdmissible\_evenPair\_sixHundredSeventyTwo} & Gaps672680 & PROVED & \checkmark & verified \\
\texttt{localFactor\_sixHundredEighty\_two} & Gaps672680 & PROVED & \checkmark & verified \\
\texttt{localFactor\_sixHundredSeventyTwo\_two} & Gaps672680 & PROVED & \checkmark & verified \\
\texttt{nu\_p\_sixHundredEighty} & Gaps672680 & PROVED & \checkmark & verified \\
\texttt{nu\_p\_sixHundredEighty\_two} & Gaps672680 & PROVED & \checkmark & verified \\
\texttt{nu\_p\_sixHundredSeventyEight} & Gaps672680 & PROVED & \checkmark & verified \\
\texttt{nu\_p\_sixHundredSeventyFour} & Gaps672680 & PROVED & \checkmark & verified \\
\texttt{nu\_p\_sixHundredSeventySix} & Gaps672680 & PROVED & \checkmark & verified \\
\texttt{nu\_p\_sixHundredSeventyTwo} & Gaps672680 & PROVED & \checkmark & verified \\
\texttt{nu\_p\_sixHundredSeventyTwo\_two} & Gaps672680 & PROVED & \checkmark & verified \\
\texttt{singular\_series\_finite\_pos\_evenPair\_sixHundredEighty} & Gaps672680 & PROVED & \checkmark & verified \\
\texttt{singular\_series\_finite\_pos\_evenPair\_sixHundredSeventyEight} & Gaps672680 & PROVED & \checkmark & verified \\
\texttt{singular\_series\_finite\_pos\_evenPair\_sixHundredSeventyFour} & Gaps672680 & PROVED & \checkmark & verified \\
\texttt{singular\_series\_finite\_pos\_evenPair\_sixHundredSeventySix} & Gaps672680 & PROVED & \checkmark & verified \\
\texttt{singular\_series\_finite\_pos\_evenPair\_sixHundredSeventyTwo} & Gaps672680 & PROVED & \checkmark & verified \\
\texttt{singular\_series\_pos\_evenPair\_sixHundredEighty} & Gaps672680 & PROVED & \checkmark & verified \\
\texttt{singular\_series\_pos\_evenPair\_sixHundredSeventyEight} & Gaps672680 & PROVED & \checkmark & verified \\
\texttt{singular\_series\_pos\_evenPair\_sixHundredSeventyFour} & Gaps672680 & PROVED & \checkmark & verified \\
\texttt{singular\_series\_pos\_evenPair\_sixHundredSeventySix} & Gaps672680 & PROVED & \checkmark & verified \\
\texttt{singular\_series\_pos\_evenPair\_sixHundredSeventyTwo} & Gaps672680 & PROVED & \checkmark & verified \\
\texttt{evenPair\_card\_sixHundredEightyEight} & Gaps682690 & PROVED & \checkmark & verified \\
\texttt{evenPair\_card\_sixHundredEightyFour} & Gaps682690 & PROVED & \checkmark & verified \\
\texttt{evenPair\_card\_sixHundredEightySix} & Gaps682690 & PROVED & \checkmark & verified \\
\texttt{evenPair\_card\_sixHundredEightyTwo} & Gaps682690 & PROVED & \checkmark & verified \\
\texttt{evenPair\_card\_sixHundredNinety} & Gaps682690 & PROVED & \checkmark & verified \\
\texttt{isAdmissible\_evenPair\_sixHundredEightyEight} & Gaps682690 & PROVED & \checkmark & verified \\
\texttt{isAdmissible\_evenPair\_sixHundredEightyFour} & Gaps682690 & PROVED & \checkmark & verified \\
\texttt{isAdmissible\_evenPair\_sixHundredEightySix} & Gaps682690 & PROVED & \checkmark & verified \\
\texttt{isAdmissible\_evenPair\_sixHundredEightyTwo} & Gaps682690 & PROVED & \checkmark & verified \\
\texttt{isAdmissible\_evenPair\_sixHundredNinety} & Gaps682690 & PROVED & \checkmark & verified \\
\texttt{localFactor\_sixHundredEightyTwo\_two} & Gaps682690 & PROVED & \checkmark & verified \\
\texttt{localFactor\_sixHundredNinety\_two} & Gaps682690 & PROVED & \checkmark & verified \\
\texttt{nu\_p\_sixHundredEightyEight} & Gaps682690 & PROVED & \checkmark & verified \\
\texttt{nu\_p\_sixHundredEightyFour} & Gaps682690 & PROVED & \checkmark & verified \\
\texttt{nu\_p\_sixHundredEightySix} & Gaps682690 & PROVED & \checkmark & verified \\
\texttt{nu\_p\_sixHundredEightyTwo} & Gaps682690 & PROVED & \checkmark & verified \\
\texttt{nu\_p\_sixHundredEightyTwo\_two} & Gaps682690 & PROVED & \checkmark & verified \\
\texttt{nu\_p\_sixHundredNinety} & Gaps682690 & PROVED & \checkmark & verified \\
\texttt{nu\_p\_sixHundredNinety\_two} & Gaps682690 & PROVED & \checkmark & verified \\
\texttt{singular\_series\_finite\_pos\_evenPair\_sixHundredEightyEight} & Gaps682690 & PROVED & \checkmark & verified \\
\texttt{singular\_series\_finite\_pos\_evenPair\_sixHundredEightyFour} & Gaps682690 & PROVED & \checkmark & verified \\
\texttt{singular\_series\_finite\_pos\_evenPair\_sixHundredEightySix} & Gaps682690 & PROVED & \checkmark & verified \\
\texttt{singular\_series\_finite\_pos\_evenPair\_sixHundredEightyTwo} & Gaps682690 & PROVED & \checkmark & verified \\
\texttt{singular\_series\_finite\_pos\_evenPair\_sixHundredNinety} & Gaps682690 & PROVED & \checkmark & verified \\
\texttt{singular\_series\_pos\_evenPair\_sixHundredEightyEight} & Gaps682690 & PROVED & \checkmark & verified \\
\texttt{singular\_series\_pos\_evenPair\_sixHundredEightyFour} & Gaps682690 & PROVED & \checkmark & verified \\
\texttt{singular\_series\_pos\_evenPair\_sixHundredEightySix} & Gaps682690 & PROVED & \checkmark & verified \\
\texttt{singular\_series\_pos\_evenPair\_sixHundredEightyTwo} & Gaps682690 & PROVED & \checkmark & verified \\
\texttt{singular\_series\_pos\_evenPair\_sixHundredNinety} & Gaps682690 & PROVED & \checkmark & verified \\
\texttt{evenPair\_card\_sevenHundred} & Gaps692700 & PROVED & \checkmark & verified \\
\texttt{evenPair\_card\_sixHundredNinetyEight} & Gaps692700 & PROVED & \checkmark & verified \\
\texttt{evenPair\_card\_sixHundredNinetyFour} & Gaps692700 & PROVED & \checkmark & verified \\
\texttt{evenPair\_card\_sixHundredNinetySix} & Gaps692700 & PROVED & \checkmark & verified \\
\texttt{evenPair\_card\_sixHundredNinetyTwo} & Gaps692700 & PROVED & \checkmark & verified \\
\texttt{isAdmissible\_evenPair\_sevenHundred} & Gaps692700 & PROVED & \checkmark & verified \\
\texttt{isAdmissible\_evenPair\_sixHundredNinetyEight} & Gaps692700 & PROVED & \checkmark & verified \\
\texttt{isAdmissible\_evenPair\_sixHundredNinetyFour} & Gaps692700 & PROVED & \checkmark & verified \\
\texttt{isAdmissible\_evenPair\_sixHundredNinetySix} & Gaps692700 & PROVED & \checkmark & verified \\
\texttt{isAdmissible\_evenPair\_sixHundredNinetyTwo} & Gaps692700 & PROVED & \checkmark & verified \\
\texttt{localFactor\_sevenHundred\_two} & Gaps692700 & PROVED & \checkmark & verified \\
\texttt{localFactor\_sixHundredNinetyTwo\_two} & Gaps692700 & PROVED & \checkmark & verified \\
\texttt{nu\_p\_sevenHundred} & Gaps692700 & PROVED & \checkmark & verified \\
\texttt{nu\_p\_sevenHundred\_two} & Gaps692700 & PROVED & \checkmark & verified \\
\texttt{nu\_p\_sixHundredNinetyEight} & Gaps692700 & PROVED & \checkmark & verified \\
\texttt{nu\_p\_sixHundredNinetyFour} & Gaps692700 & PROVED & \checkmark & verified \\
\texttt{nu\_p\_sixHundredNinetySix} & Gaps692700 & PROVED & \checkmark & verified \\
\texttt{nu\_p\_sixHundredNinetyTwo} & Gaps692700 & PROVED & \checkmark & verified \\
\texttt{nu\_p\_sixHundredNinetyTwo\_two} & Gaps692700 & PROVED & \checkmark & verified \\
\texttt{singular\_series\_finite\_pos\_evenPair\_sevenHundred} & Gaps692700 & PROVED & \checkmark & verified \\
\texttt{singular\_series\_finite\_pos\_evenPair\_sixHundredNinetyEight} & Gaps692700 & PROVED & \checkmark & verified \\
\texttt{singular\_series\_finite\_pos\_evenPair\_sixHundredNinetyFour} & Gaps692700 & PROVED & \checkmark & verified \\
\texttt{singular\_series\_finite\_pos\_evenPair\_sixHundredNinetySix} & Gaps692700 & PROVED & \checkmark & verified \\
\texttt{singular\_series\_finite\_pos\_evenPair\_sixHundredNinetyTwo} & Gaps692700 & PROVED & \checkmark & verified \\
\texttt{singular\_series\_pos\_evenPair\_sevenHundred} & Gaps692700 & PROVED & \checkmark & verified \\
\texttt{singular\_series\_pos\_evenPair\_sixHundredNinetyEight} & Gaps692700 & PROVED & \checkmark & verified \\
\texttt{singular\_series\_pos\_evenPair\_sixHundredNinetyFour} & Gaps692700 & PROVED & \checkmark & verified \\
\texttt{singular\_series\_pos\_evenPair\_sixHundredNinetySix} & Gaps692700 & PROVED & \checkmark & verified \\
\texttt{singular\_series\_pos\_evenPair\_sixHundredNinetyTwo} & Gaps692700 & PROVED & \checkmark & verified \\
\texttt{evenPair\_card\_sevenHundredEight} & Gaps702710 & PROVED & \checkmark & verified \\
\texttt{evenPair\_card\_sevenHundredFour} & Gaps702710 & PROVED & \checkmark & verified \\
\texttt{evenPair\_card\_sevenHundredSix} & Gaps702710 & PROVED & \checkmark & verified \\
\texttt{evenPair\_card\_sevenHundredTen} & Gaps702710 & PROVED & \checkmark & verified \\
\texttt{evenPair\_card\_sevenHundredTwo} & Gaps702710 & PROVED & \checkmark & verified \\
\texttt{isAdmissible\_evenPair\_sevenHundredEight} & Gaps702710 & PROVED & \checkmark & verified \\
\texttt{isAdmissible\_evenPair\_sevenHundredFour} & Gaps702710 & PROVED & \checkmark & verified \\
\texttt{isAdmissible\_evenPair\_sevenHundredSix} & Gaps702710 & PROVED & \checkmark & verified \\
\texttt{isAdmissible\_evenPair\_sevenHundredTen} & Gaps702710 & PROVED & \checkmark & verified \\
\texttt{isAdmissible\_evenPair\_sevenHundredTwo} & Gaps702710 & PROVED & \checkmark & verified \\
\texttt{localFactor\_sevenHundredTen\_two} & Gaps702710 & PROVED & \checkmark & verified \\
\texttt{localFactor\_sevenHundredTwo\_two} & Gaps702710 & PROVED & \checkmark & verified \\
\texttt{nu\_p\_sevenHundredEight} & Gaps702710 & PROVED & \checkmark & verified \\
\texttt{nu\_p\_sevenHundredFour} & Gaps702710 & PROVED & \checkmark & verified \\
\texttt{nu\_p\_sevenHundredSix} & Gaps702710 & PROVED & \checkmark & verified \\
\texttt{nu\_p\_sevenHundredTen} & Gaps702710 & PROVED & \checkmark & verified \\
\texttt{nu\_p\_sevenHundredTen\_two} & Gaps702710 & PROVED & \checkmark & verified \\
\texttt{nu\_p\_sevenHundredTwo} & Gaps702710 & PROVED & \checkmark & verified \\
\texttt{nu\_p\_sevenHundredTwo\_two} & Gaps702710 & PROVED & \checkmark & verified \\
\texttt{singular\_series\_finite\_pos\_evenPair\_sevenHundredEight} & Gaps702710 & PROVED & \checkmark & verified \\
\texttt{singular\_series\_finite\_pos\_evenPair\_sevenHundredFour} & Gaps702710 & PROVED & \checkmark & verified \\
\texttt{singular\_series\_finite\_pos\_evenPair\_sevenHundredSix} & Gaps702710 & PROVED & \checkmark & verified \\
\texttt{singular\_series\_finite\_pos\_evenPair\_sevenHundredTen} & Gaps702710 & PROVED & \checkmark & verified \\
\texttt{singular\_series\_finite\_pos\_evenPair\_sevenHundredTwo} & Gaps702710 & PROVED & \checkmark & verified \\
\texttt{singular\_series\_pos\_evenPair\_sevenHundredEight} & Gaps702710 & PROVED & \checkmark & verified \\
\texttt{singular\_series\_pos\_evenPair\_sevenHundredFour} & Gaps702710 & PROVED & \checkmark & verified \\
\texttt{singular\_series\_pos\_evenPair\_sevenHundredSix} & Gaps702710 & PROVED & \checkmark & verified \\
\texttt{singular\_series\_pos\_evenPair\_sevenHundredTen} & Gaps702710 & PROVED & \checkmark & verified \\
\texttt{singular\_series\_pos\_evenPair\_sevenHundredTwo} & Gaps702710 & PROVED & \checkmark & verified \\
\texttt{evenPair\_card\_sevenHundredEighteen} & Gaps712720 & PROVED & \checkmark & verified \\
\texttt{evenPair\_card\_sevenHundredFourteen} & Gaps712720 & PROVED & \checkmark & verified \\
\texttt{evenPair\_card\_sevenHundredSixteen} & Gaps712720 & PROVED & \checkmark & verified \\
\texttt{evenPair\_card\_sevenHundredTwelve} & Gaps712720 & PROVED & \checkmark & verified \\
\texttt{evenPair\_card\_sevenHundredTwenty} & Gaps712720 & PROVED & \checkmark & verified \\
\texttt{isAdmissible\_evenPair\_sevenHundredEighteen} & Gaps712720 & PROVED & \checkmark & verified \\
\texttt{isAdmissible\_evenPair\_sevenHundredFourteen} & Gaps712720 & PROVED & \checkmark & verified \\
\texttt{isAdmissible\_evenPair\_sevenHundredSixteen} & Gaps712720 & PROVED & \checkmark & verified \\
\texttt{isAdmissible\_evenPair\_sevenHundredTwelve} & Gaps712720 & PROVED & \checkmark & verified \\
\texttt{isAdmissible\_evenPair\_sevenHundredTwenty} & Gaps712720 & PROVED & \checkmark & verified \\
\texttt{localFactor\_sevenHundredTwelve\_two} & Gaps712720 & PROVED & \checkmark & verified \\
\texttt{localFactor\_sevenHundredTwenty\_two} & Gaps712720 & PROVED & \checkmark & verified \\
\texttt{nu\_p\_sevenHundredEighteen} & Gaps712720 & PROVED & \checkmark & verified \\
\texttt{nu\_p\_sevenHundredFourteen} & Gaps712720 & PROVED & \checkmark & verified \\
\texttt{nu\_p\_sevenHundredSixteen} & Gaps712720 & PROVED & \checkmark & verified \\
\texttt{nu\_p\_sevenHundredTwelve} & Gaps712720 & PROVED & \checkmark & verified \\
\texttt{nu\_p\_sevenHundredTwelve\_two} & Gaps712720 & PROVED & \checkmark & verified \\
\texttt{nu\_p\_sevenHundredTwenty} & Gaps712720 & PROVED & \checkmark & verified \\
\texttt{nu\_p\_sevenHundredTwenty\_two} & Gaps712720 & PROVED & \checkmark & verified \\
\texttt{singular\_series\_finite\_pos\_evenPair\_sevenHundredEighteen} & Gaps712720 & PROVED & \checkmark & verified \\
\texttt{singular\_series\_finite\_pos\_evenPair\_sevenHundredFourteen} & Gaps712720 & PROVED & \checkmark & verified \\
\texttt{singular\_series\_finite\_pos\_evenPair\_sevenHundredSixteen} & Gaps712720 & PROVED & \checkmark & verified \\
\texttt{singular\_series\_finite\_pos\_evenPair\_sevenHundredTwelve} & Gaps712720 & PROVED & \checkmark & verified \\
\texttt{singular\_series\_finite\_pos\_evenPair\_sevenHundredTwenty} & Gaps712720 & PROVED & \checkmark & verified \\
\texttt{singular\_series\_pos\_evenPair\_sevenHundredEighteen} & Gaps712720 & PROVED & \checkmark & verified \\
\texttt{singular\_series\_pos\_evenPair\_sevenHundredFourteen} & Gaps712720 & PROVED & \checkmark & verified \\
\texttt{singular\_series\_pos\_evenPair\_sevenHundredSixteen} & Gaps712720 & PROVED & \checkmark & verified \\
\texttt{singular\_series\_pos\_evenPair\_sevenHundredTwelve} & Gaps712720 & PROVED & \checkmark & verified \\
\texttt{singular\_series\_pos\_evenPair\_sevenHundredTwenty} & Gaps712720 & PROVED & \checkmark & verified \\
\texttt{evenPair\_card\_sevenHundredThirty} & Gaps722730 & PROVED & \checkmark & verified \\
\texttt{evenPair\_card\_sevenHundredTwentyEight} & Gaps722730 & PROVED & \checkmark & verified \\
\texttt{evenPair\_card\_sevenHundredTwentyFour} & Gaps722730 & PROVED & \checkmark & verified \\
\texttt{evenPair\_card\_sevenHundredTwentySix} & Gaps722730 & PROVED & \checkmark & verified \\
\texttt{evenPair\_card\_sevenHundredTwentyTwo} & Gaps722730 & PROVED & \checkmark & verified \\
\texttt{isAdmissible\_evenPair\_sevenHundredThirty} & Gaps722730 & PROVED & \checkmark & verified \\
\texttt{isAdmissible\_evenPair\_sevenHundredTwentyEight} & Gaps722730 & PROVED & \checkmark & verified \\
\texttt{isAdmissible\_evenPair\_sevenHundredTwentyFour} & Gaps722730 & PROVED & \checkmark & verified \\
\texttt{isAdmissible\_evenPair\_sevenHundredTwentySix} & Gaps722730 & PROVED & \checkmark & verified \\
\texttt{isAdmissible\_evenPair\_sevenHundredTwentyTwo} & Gaps722730 & PROVED & \checkmark & verified \\
\texttt{localFactor\_sevenHundredThirty\_two} & Gaps722730 & PROVED & \checkmark & verified \\
\texttt{localFactor\_sevenHundredTwentyTwo\_two} & Gaps722730 & PROVED & \checkmark & verified \\
\texttt{nu\_p\_sevenHundredThirty} & Gaps722730 & PROVED & \checkmark & verified \\
\texttt{nu\_p\_sevenHundredThirty\_two} & Gaps722730 & PROVED & \checkmark & verified \\
\texttt{nu\_p\_sevenHundredTwentyEight} & Gaps722730 & PROVED & \checkmark & verified \\
\texttt{nu\_p\_sevenHundredTwentyFour} & Gaps722730 & PROVED & \checkmark & verified \\
\texttt{nu\_p\_sevenHundredTwentySix} & Gaps722730 & PROVED & \checkmark & verified \\
\texttt{nu\_p\_sevenHundredTwentyTwo} & Gaps722730 & PROVED & \checkmark & verified \\
\texttt{nu\_p\_sevenHundredTwentyTwo\_two} & Gaps722730 & PROVED & \checkmark & verified \\
\texttt{singular\_series\_finite\_pos\_evenPair\_sevenHundredThirty} & Gaps722730 & PROVED & \checkmark & verified \\
\texttt{singular\_series\_finite\_pos\_evenPair\_sevenHundredTwentyEight} & Gaps722730 & PROVED & \checkmark & verified \\
\texttt{singular\_series\_finite\_pos\_evenPair\_sevenHundredTwentyFour} & Gaps722730 & PROVED & \checkmark & verified \\
\texttt{singular\_series\_finite\_pos\_evenPair\_sevenHundredTwentySix} & Gaps722730 & PROVED & \checkmark & verified \\
\texttt{singular\_series\_finite\_pos\_evenPair\_sevenHundredTwentyTwo} & Gaps722730 & PROVED & \checkmark & verified \\
\texttt{singular\_series\_pos\_evenPair\_sevenHundredThirty} & Gaps722730 & PROVED & \checkmark & verified \\
\texttt{singular\_series\_pos\_evenPair\_sevenHundredTwentyEight} & Gaps722730 & PROVED & \checkmark & verified \\
\texttt{singular\_series\_pos\_evenPair\_sevenHundredTwentyFour} & Gaps722730 & PROVED & \checkmark & verified \\
\texttt{singular\_series\_pos\_evenPair\_sevenHundredTwentySix} & Gaps722730 & PROVED & \checkmark & verified \\
\texttt{singular\_series\_pos\_evenPair\_sevenHundredTwentyTwo} & Gaps722730 & PROVED & \checkmark & verified \\
\texttt{evenPair\_card\_eighty} & Gaps7280 & PROVED & \checkmark & verified \\
\texttt{evenPair\_card\_seventyEight} & Gaps7280 & PROVED & \checkmark & verified \\
\texttt{evenPair\_card\_seventyFour} & Gaps7280 & PROVED & \checkmark & verified \\
\texttt{evenPair\_card\_seventySix} & Gaps7280 & PROVED & \checkmark & verified \\
\texttt{evenPair\_card\_seventyTwo} & Gaps7280 & PROVED & \checkmark & verified \\
\texttt{isAdmissible\_evenPair\_eighty} & Gaps7280 & PROVED & \checkmark & verified \\
\texttt{isAdmissible\_evenPair\_seventyEight} & Gaps7280 & PROVED & \checkmark & verified \\
\texttt{isAdmissible\_evenPair\_seventyFour} & Gaps7280 & PROVED & \checkmark & verified \\
\texttt{isAdmissible\_evenPair\_seventySix} & Gaps7280 & PROVED & \checkmark & verified \\
\texttt{isAdmissible\_evenPair\_seventyTwo} & Gaps7280 & PROVED & \checkmark & verified \\
\texttt{localFactor\_eighty\_two} & Gaps7280 & PROVED & \checkmark & verified \\
\texttt{localFactor\_seventyTwo\_two} & Gaps7280 & PROVED & \checkmark & verified \\
\texttt{nu\_p\_eighty} & Gaps7280 & PROVED & \checkmark & verified \\
\texttt{nu\_p\_eighty\_two} & Gaps7280 & PROVED & \checkmark & verified \\
\texttt{nu\_p\_seventyEight} & Gaps7280 & PROVED & \checkmark & verified \\
\texttt{nu\_p\_seventyFour} & Gaps7280 & PROVED & \checkmark & verified \\
\texttt{nu\_p\_seventySix} & Gaps7280 & PROVED & \checkmark & verified \\
\texttt{nu\_p\_seventyTwo} & Gaps7280 & PROVED & \checkmark & verified \\
\texttt{nu\_p\_seventyTwo\_two} & Gaps7280 & PROVED & \checkmark & verified \\
\texttt{singular\_series\_finite\_pos\_evenPair\_eighty} & Gaps7280 & PROVED & \checkmark & verified \\
\texttt{singular\_series\_finite\_pos\_evenPair\_seventyEight} & Gaps7280 & PROVED & \checkmark & verified \\
\texttt{singular\_series\_finite\_pos\_evenPair\_seventyFour} & Gaps7280 & PROVED & \checkmark & verified \\
\texttt{singular\_series\_finite\_pos\_evenPair\_seventySix} & Gaps7280 & PROVED & \checkmark & verified \\
\texttt{singular\_series\_finite\_pos\_evenPair\_seventyTwo} & Gaps7280 & PROVED & \checkmark & verified \\
\texttt{singular\_series\_pos\_evenPair\_eighty} & Gaps7280 & PROVED & \checkmark & verified \\
\texttt{singular\_series\_pos\_evenPair\_seventyEight} & Gaps7280 & PROVED & \checkmark & verified \\
\texttt{singular\_series\_pos\_evenPair\_seventyFour} & Gaps7280 & PROVED & \checkmark & verified \\
\texttt{singular\_series\_pos\_evenPair\_seventySix} & Gaps7280 & PROVED & \checkmark & verified \\
\texttt{singular\_series\_pos\_evenPair\_seventyTwo} & Gaps7280 & PROVED & \checkmark & verified \\
\texttt{evenPair\_card\_sevenHundredForty} & Gaps732740 & PROVED & \checkmark & verified \\
\texttt{evenPair\_card\_sevenHundredThirtyEight} & Gaps732740 & PROVED & \checkmark & verified \\
\texttt{evenPair\_card\_sevenHundredThirtyFour} & Gaps732740 & PROVED & \checkmark & verified \\
\texttt{evenPair\_card\_sevenHundredThirtySix} & Gaps732740 & PROVED & \checkmark & verified \\
\texttt{evenPair\_card\_sevenHundredThirtyTwo} & Gaps732740 & PROVED & \checkmark & verified \\
\texttt{isAdmissible\_evenPair\_sevenHundredForty} & Gaps732740 & PROVED & \checkmark & verified \\
\texttt{isAdmissible\_evenPair\_sevenHundredThirtyEight} & Gaps732740 & PROVED & \checkmark & verified \\
\texttt{isAdmissible\_evenPair\_sevenHundredThirtyFour} & Gaps732740 & PROVED & \checkmark & verified \\
\texttt{isAdmissible\_evenPair\_sevenHundredThirtySix} & Gaps732740 & PROVED & \checkmark & verified \\
\texttt{isAdmissible\_evenPair\_sevenHundredThirtyTwo} & Gaps732740 & PROVED & \checkmark & verified \\
\texttt{localFactor\_sevenHundredForty\_two} & Gaps732740 & PROVED & \checkmark & verified \\
\texttt{localFactor\_sevenHundredThirtyTwo\_two} & Gaps732740 & PROVED & \checkmark & verified \\
\texttt{nu\_p\_sevenHundredForty} & Gaps732740 & PROVED & \checkmark & verified \\
\texttt{nu\_p\_sevenHundredForty\_two} & Gaps732740 & PROVED & \checkmark & verified \\
\texttt{nu\_p\_sevenHundredThirtyEight} & Gaps732740 & PROVED & \checkmark & verified \\
\texttt{nu\_p\_sevenHundredThirtyFour} & Gaps732740 & PROVED & \checkmark & verified \\
\texttt{nu\_p\_sevenHundredThirtySix} & Gaps732740 & PROVED & \checkmark & verified \\
\texttt{nu\_p\_sevenHundredThirtyTwo} & Gaps732740 & PROVED & \checkmark & verified \\
\texttt{nu\_p\_sevenHundredThirtyTwo\_two} & Gaps732740 & PROVED & \checkmark & verified \\
\texttt{singular\_series\_finite\_pos\_evenPair\_sevenHundredForty} & Gaps732740 & PROVED & \checkmark & verified \\
\texttt{singular\_series\_finite\_pos\_evenPair\_sevenHundredThirtyEight} & Gaps732740 & PROVED & \checkmark & verified \\
\texttt{singular\_series\_finite\_pos\_evenPair\_sevenHundredThirtyFour} & Gaps732740 & PROVED & \checkmark & verified \\
\texttt{singular\_series\_finite\_pos\_evenPair\_sevenHundredThirtySix} & Gaps732740 & PROVED & \checkmark & verified \\
\texttt{singular\_series\_finite\_pos\_evenPair\_sevenHundredThirtyTwo} & Gaps732740 & PROVED & \checkmark & verified \\
\texttt{singular\_series\_pos\_evenPair\_sevenHundredForty} & Gaps732740 & PROVED & \checkmark & verified \\
\texttt{singular\_series\_pos\_evenPair\_sevenHundredThirtyEight} & Gaps732740 & PROVED & \checkmark & verified \\
\texttt{singular\_series\_pos\_evenPair\_sevenHundredThirtyFour} & Gaps732740 & PROVED & \checkmark & verified \\
\texttt{singular\_series\_pos\_evenPair\_sevenHundredThirtySix} & Gaps732740 & PROVED & \checkmark & verified \\
\texttt{singular\_series\_pos\_evenPair\_sevenHundredThirtyTwo} & Gaps732740 & PROVED & \checkmark & verified \\
\texttt{evenPair\_card\_sevenHundredFifty} & Gaps742750 & PROVED & \checkmark & verified \\
\texttt{evenPair\_card\_sevenHundredFortyEight} & Gaps742750 & PROVED & \checkmark & verified \\
\texttt{evenPair\_card\_sevenHundredFortyFour} & Gaps742750 & PROVED & \checkmark & verified \\
\texttt{evenPair\_card\_sevenHundredFortySix} & Gaps742750 & PROVED & \checkmark & verified \\
\texttt{evenPair\_card\_sevenHundredFortyTwo} & Gaps742750 & PROVED & \checkmark & verified \\
\texttt{isAdmissible\_evenPair\_sevenHundredFifty} & Gaps742750 & PROVED & \checkmark & verified \\
\texttt{isAdmissible\_evenPair\_sevenHundredFortyEight} & Gaps742750 & PROVED & \checkmark & verified \\
\texttt{isAdmissible\_evenPair\_sevenHundredFortyFour} & Gaps742750 & PROVED & \checkmark & verified \\
\texttt{isAdmissible\_evenPair\_sevenHundredFortySix} & Gaps742750 & PROVED & \checkmark & verified \\
\texttt{isAdmissible\_evenPair\_sevenHundredFortyTwo} & Gaps742750 & PROVED & \checkmark & verified \\
\texttt{localFactor\_sevenHundredFifty\_two} & Gaps742750 & PROVED & \checkmark & verified \\
\texttt{localFactor\_sevenHundredFortyTwo\_two} & Gaps742750 & PROVED & \checkmark & verified \\
\texttt{nu\_p\_sevenHundredFifty} & Gaps742750 & PROVED & \checkmark & verified \\
\texttt{nu\_p\_sevenHundredFifty\_two} & Gaps742750 & PROVED & \checkmark & verified \\
\texttt{nu\_p\_sevenHundredFortyEight} & Gaps742750 & PROVED & \checkmark & verified \\
\texttt{nu\_p\_sevenHundredFortyFour} & Gaps742750 & PROVED & \checkmark & verified \\
\texttt{nu\_p\_sevenHundredFortySix} & Gaps742750 & PROVED & \checkmark & verified \\
\texttt{nu\_p\_sevenHundredFortyTwo} & Gaps742750 & PROVED & \checkmark & verified \\
\texttt{nu\_p\_sevenHundredFortyTwo\_two} & Gaps742750 & PROVED & \checkmark & verified \\
\texttt{singular\_series\_finite\_pos\_evenPair\_sevenHundredFifty} & Gaps742750 & PROVED & \checkmark & verified \\
\texttt{singular\_series\_finite\_pos\_evenPair\_sevenHundredFortyEight} & Gaps742750 & PROVED & \checkmark & verified \\
\texttt{singular\_series\_finite\_pos\_evenPair\_sevenHundredFortyFour} & Gaps742750 & PROVED & \checkmark & verified \\
\texttt{singular\_series\_finite\_pos\_evenPair\_sevenHundredFortySix} & Gaps742750 & PROVED & \checkmark & verified \\
\texttt{singular\_series\_finite\_pos\_evenPair\_sevenHundredFortyTwo} & Gaps742750 & PROVED & \checkmark & verified \\
\texttt{singular\_series\_pos\_evenPair\_sevenHundredFifty} & Gaps742750 & PROVED & \checkmark & verified \\
\texttt{singular\_series\_pos\_evenPair\_sevenHundredFortyEight} & Gaps742750 & PROVED & \checkmark & verified \\
\texttt{singular\_series\_pos\_evenPair\_sevenHundredFortyFour} & Gaps742750 & PROVED & \checkmark & verified \\
\texttt{singular\_series\_pos\_evenPair\_sevenHundredFortySix} & Gaps742750 & PROVED & \checkmark & verified \\
\texttt{singular\_series\_pos\_evenPair\_sevenHundredFortyTwo} & Gaps742750 & PROVED & \checkmark & verified \\
\texttt{evenPair\_card\_sevenHundredFiftyEight} & Gaps752760 & PROVED & \checkmark & verified \\
\texttt{evenPair\_card\_sevenHundredFiftyFour} & Gaps752760 & PROVED & \checkmark & verified \\
\texttt{evenPair\_card\_sevenHundredFiftySix} & Gaps752760 & PROVED & \checkmark & verified \\
\texttt{evenPair\_card\_sevenHundredFiftyTwo} & Gaps752760 & PROVED & \checkmark & verified \\
\texttt{evenPair\_card\_sevenHundredSixty} & Gaps752760 & PROVED & \checkmark & verified \\
\texttt{isAdmissible\_evenPair\_sevenHundredFiftyEight} & Gaps752760 & PROVED & \checkmark & verified \\
\texttt{isAdmissible\_evenPair\_sevenHundredFiftyFour} & Gaps752760 & PROVED & \checkmark & verified \\
\texttt{isAdmissible\_evenPair\_sevenHundredFiftySix} & Gaps752760 & PROVED & \checkmark & verified \\
\texttt{isAdmissible\_evenPair\_sevenHundredFiftyTwo} & Gaps752760 & PROVED & \checkmark & verified \\
\texttt{isAdmissible\_evenPair\_sevenHundredSixty} & Gaps752760 & PROVED & \checkmark & verified \\
\texttt{localFactor\_sevenHundredFiftyTwo\_two} & Gaps752760 & PROVED & \checkmark & verified \\
\texttt{localFactor\_sevenHundredSixty\_two} & Gaps752760 & PROVED & \checkmark & verified \\
\texttt{nu\_p\_sevenHundredFiftyEight} & Gaps752760 & PROVED & \checkmark & verified \\
\texttt{nu\_p\_sevenHundredFiftyFour} & Gaps752760 & PROVED & \checkmark & verified \\
\texttt{nu\_p\_sevenHundredFiftySix} & Gaps752760 & PROVED & \checkmark & verified \\
\texttt{nu\_p\_sevenHundredFiftyTwo} & Gaps752760 & PROVED & \checkmark & verified \\
\texttt{nu\_p\_sevenHundredFiftyTwo\_two} & Gaps752760 & PROVED & \checkmark & verified \\
\texttt{nu\_p\_sevenHundredSixty} & Gaps752760 & PROVED & \checkmark & verified \\
\texttt{nu\_p\_sevenHundredSixty\_two} & Gaps752760 & PROVED & \checkmark & verified \\
\texttt{singular\_series\_finite\_pos\_evenPair\_sevenHundredFiftyEight} & Gaps752760 & PROVED & \checkmark & verified \\
\texttt{singular\_series\_finite\_pos\_evenPair\_sevenHundredFiftyFour} & Gaps752760 & PROVED & \checkmark & verified \\
\texttt{singular\_series\_finite\_pos\_evenPair\_sevenHundredFiftySix} & Gaps752760 & PROVED & \checkmark & verified \\
\texttt{singular\_series\_finite\_pos\_evenPair\_sevenHundredFiftyTwo} & Gaps752760 & PROVED & \checkmark & verified \\
\texttt{singular\_series\_finite\_pos\_evenPair\_sevenHundredSixty} & Gaps752760 & PROVED & \checkmark & verified \\
\texttt{singular\_series\_pos\_evenPair\_sevenHundredFiftyEight} & Gaps752760 & PROVED & \checkmark & verified \\
\texttt{singular\_series\_pos\_evenPair\_sevenHundredFiftyFour} & Gaps752760 & PROVED & \checkmark & verified \\
\texttt{singular\_series\_pos\_evenPair\_sevenHundredFiftySix} & Gaps752760 & PROVED & \checkmark & verified \\
\texttt{singular\_series\_pos\_evenPair\_sevenHundredFiftyTwo} & Gaps752760 & PROVED & \checkmark & verified \\
\texttt{singular\_series\_pos\_evenPair\_sevenHundredSixty} & Gaps752760 & PROVED & \checkmark & verified \\
\texttt{evenPair\_card\_sevenHundredSeventy} & Gaps762770 & PROVED & \checkmark & verified \\
\texttt{evenPair\_card\_sevenHundredSixtyEight} & Gaps762770 & PROVED & \checkmark & verified \\
\texttt{evenPair\_card\_sevenHundredSixtyFour} & Gaps762770 & PROVED & \checkmark & verified \\
\texttt{evenPair\_card\_sevenHundredSixtySix} & Gaps762770 & PROVED & \checkmark & verified \\
\texttt{evenPair\_card\_sevenHundredSixtyTwo} & Gaps762770 & PROVED & \checkmark & verified \\
\texttt{isAdmissible\_evenPair\_sevenHundredSeventy} & Gaps762770 & PROVED & \checkmark & verified \\
\texttt{isAdmissible\_evenPair\_sevenHundredSixtyEight} & Gaps762770 & PROVED & \checkmark & verified \\
\texttt{isAdmissible\_evenPair\_sevenHundredSixtyFour} & Gaps762770 & PROVED & \checkmark & verified \\
\texttt{isAdmissible\_evenPair\_sevenHundredSixtySix} & Gaps762770 & PROVED & \checkmark & verified \\
\texttt{isAdmissible\_evenPair\_sevenHundredSixtyTwo} & Gaps762770 & PROVED & \checkmark & verified \\
\texttt{localFactor\_sevenHundredSeventy\_two} & Gaps762770 & PROVED & \checkmark & verified \\
\texttt{localFactor\_sevenHundredSixtyTwo\_two} & Gaps762770 & PROVED & \checkmark & verified \\
\texttt{nu\_p\_sevenHundredSeventy} & Gaps762770 & PROVED & \checkmark & verified \\
\texttt{nu\_p\_sevenHundredSeventy\_two} & Gaps762770 & PROVED & \checkmark & verified \\
\texttt{nu\_p\_sevenHundredSixtyEight} & Gaps762770 & PROVED & \checkmark & verified \\
\texttt{nu\_p\_sevenHundredSixtyFour} & Gaps762770 & PROVED & \checkmark & verified \\
\texttt{nu\_p\_sevenHundredSixtySix} & Gaps762770 & PROVED & \checkmark & verified \\
\texttt{nu\_p\_sevenHundredSixtyTwo} & Gaps762770 & PROVED & \checkmark & verified \\
\texttt{nu\_p\_sevenHundredSixtyTwo\_two} & Gaps762770 & PROVED & \checkmark & verified \\
\texttt{singular\_series\_finite\_pos\_evenPair\_sevenHundredSeventy} & Gaps762770 & PROVED & \checkmark & verified \\
\texttt{singular\_series\_finite\_pos\_evenPair\_sevenHundredSixtyEight} & Gaps762770 & PROVED & \checkmark & verified \\
\texttt{singular\_series\_finite\_pos\_evenPair\_sevenHundredSixtyFour} & Gaps762770 & PROVED & \checkmark & verified \\
\texttt{singular\_series\_finite\_pos\_evenPair\_sevenHundredSixtySix} & Gaps762770 & PROVED & \checkmark & verified \\
\texttt{singular\_series\_finite\_pos\_evenPair\_sevenHundredSixtyTwo} & Gaps762770 & PROVED & \checkmark & verified \\
\texttt{singular\_series\_pos\_evenPair\_sevenHundredSeventy} & Gaps762770 & PROVED & \checkmark & verified \\
\texttt{singular\_series\_pos\_evenPair\_sevenHundredSixtyEight} & Gaps762770 & PROVED & \checkmark & verified \\
\texttt{singular\_series\_pos\_evenPair\_sevenHundredSixtyFour} & Gaps762770 & PROVED & \checkmark & verified \\
\texttt{singular\_series\_pos\_evenPair\_sevenHundredSixtySix} & Gaps762770 & PROVED & \checkmark & verified \\
\texttt{singular\_series\_pos\_evenPair\_sevenHundredSixtyTwo} & Gaps762770 & PROVED & \checkmark & verified \\
\texttt{evenPair\_card\_sevenHundredEighty} & Gaps772780 & PROVED & \checkmark & verified \\
\texttt{evenPair\_card\_sevenHundredSeventyEight} & Gaps772780 & PROVED & \checkmark & verified \\
\texttt{evenPair\_card\_sevenHundredSeventyFour} & Gaps772780 & PROVED & \checkmark & verified \\
\texttt{evenPair\_card\_sevenHundredSeventySix} & Gaps772780 & PROVED & \checkmark & verified \\
\texttt{evenPair\_card\_sevenHundredSeventyTwo} & Gaps772780 & PROVED & \checkmark & verified \\
\texttt{isAdmissible\_evenPair\_sevenHundredEighty} & Gaps772780 & PROVED & \checkmark & verified \\
\texttt{isAdmissible\_evenPair\_sevenHundredSeventyEight} & Gaps772780 & PROVED & \checkmark & verified \\
\texttt{isAdmissible\_evenPair\_sevenHundredSeventyFour} & Gaps772780 & PROVED & \checkmark & verified \\
\texttt{isAdmissible\_evenPair\_sevenHundredSeventySix} & Gaps772780 & PROVED & \checkmark & verified \\
\texttt{isAdmissible\_evenPair\_sevenHundredSeventyTwo} & Gaps772780 & PROVED & \checkmark & verified \\
\texttt{localFactor\_sevenHundredEighty\_two} & Gaps772780 & PROVED & \checkmark & verified \\
\texttt{localFactor\_sevenHundredSeventyTwo\_two} & Gaps772780 & PROVED & \checkmark & verified \\
\texttt{nu\_p\_sevenHundredEighty} & Gaps772780 & PROVED & \checkmark & verified \\
\texttt{nu\_p\_sevenHundredEighty\_two} & Gaps772780 & PROVED & \checkmark & verified \\
\texttt{nu\_p\_sevenHundredSeventyEight} & Gaps772780 & PROVED & \checkmark & verified \\
\texttt{nu\_p\_sevenHundredSeventyFour} & Gaps772780 & PROVED & \checkmark & verified \\
\texttt{nu\_p\_sevenHundredSeventySix} & Gaps772780 & PROVED & \checkmark & verified \\
\texttt{nu\_p\_sevenHundredSeventyTwo} & Gaps772780 & PROVED & \checkmark & verified \\
\texttt{nu\_p\_sevenHundredSeventyTwo\_two} & Gaps772780 & PROVED & \checkmark & verified \\
\texttt{singular\_series\_finite\_pos\_evenPair\_sevenHundredEighty} & Gaps772780 & PROVED & \checkmark & verified \\
\texttt{singular\_series\_finite\_pos\_evenPair\_sevenHundredSeventyEight} & Gaps772780 & PROVED & \checkmark & verified \\
\texttt{singular\_series\_finite\_pos\_evenPair\_sevenHundredSeventyFour} & Gaps772780 & PROVED & \checkmark & verified \\
\texttt{singular\_series\_finite\_pos\_evenPair\_sevenHundredSeventySix} & Gaps772780 & PROVED & \checkmark & verified \\
\texttt{singular\_series\_finite\_pos\_evenPair\_sevenHundredSeventyTwo} & Gaps772780 & PROVED & \checkmark & verified \\
\texttt{singular\_series\_pos\_evenPair\_sevenHundredEighty} & Gaps772780 & PROVED & \checkmark & verified \\
\texttt{singular\_series\_pos\_evenPair\_sevenHundredSeventyEight} & Gaps772780 & PROVED & \checkmark & verified \\
\texttt{singular\_series\_pos\_evenPair\_sevenHundredSeventyFour} & Gaps772780 & PROVED & \checkmark & verified \\
\texttt{singular\_series\_pos\_evenPair\_sevenHundredSeventySix} & Gaps772780 & PROVED & \checkmark & verified \\
\texttt{singular\_series\_pos\_evenPair\_sevenHundredSeventyTwo} & Gaps772780 & PROVED & \checkmark & verified \\
\texttt{evenPair\_card\_sevenHundredEightyEight} & Gaps782790 & PROVED & \checkmark & verified \\
\texttt{evenPair\_card\_sevenHundredEightyFour} & Gaps782790 & PROVED & \checkmark & verified \\
\texttt{evenPair\_card\_sevenHundredEightySix} & Gaps782790 & PROVED & \checkmark & verified \\
\texttt{evenPair\_card\_sevenHundredEightyTwo} & Gaps782790 & PROVED & \checkmark & verified \\
\texttt{evenPair\_card\_sevenHundredNinety} & Gaps782790 & PROVED & \checkmark & verified \\
\texttt{isAdmissible\_evenPair\_sevenHundredEightyEight} & Gaps782790 & PROVED & \checkmark & verified \\
\texttt{isAdmissible\_evenPair\_sevenHundredEightyFour} & Gaps782790 & PROVED & \checkmark & verified \\
\texttt{isAdmissible\_evenPair\_sevenHundredEightySix} & Gaps782790 & PROVED & \checkmark & verified \\
\texttt{isAdmissible\_evenPair\_sevenHundredEightyTwo} & Gaps782790 & PROVED & \checkmark & verified \\
\texttt{isAdmissible\_evenPair\_sevenHundredNinety} & Gaps782790 & PROVED & \checkmark & verified \\
\texttt{localFactor\_sevenHundredEightyTwo\_two} & Gaps782790 & PROVED & \checkmark & verified \\
\texttt{localFactor\_sevenHundredNinety\_two} & Gaps782790 & PROVED & \checkmark & verified \\
\texttt{nu\_p\_sevenHundredEightyEight} & Gaps782790 & PROVED & \checkmark & verified \\
\texttt{nu\_p\_sevenHundredEightyFour} & Gaps782790 & PROVED & \checkmark & verified \\
\texttt{nu\_p\_sevenHundredEightySix} & Gaps782790 & PROVED & \checkmark & verified \\
\texttt{nu\_p\_sevenHundredEightyTwo} & Gaps782790 & PROVED & \checkmark & verified \\
\texttt{nu\_p\_sevenHundredEightyTwo\_two} & Gaps782790 & PROVED & \checkmark & verified \\
\texttt{nu\_p\_sevenHundredNinety} & Gaps782790 & PROVED & \checkmark & verified \\
\texttt{nu\_p\_sevenHundredNinety\_two} & Gaps782790 & PROVED & \checkmark & verified \\
\texttt{singular\_series\_finite\_pos\_evenPair\_sevenHundredEightyEight} & Gaps782790 & PROVED & \checkmark & verified \\
\texttt{singular\_series\_finite\_pos\_evenPair\_sevenHundredEightyFour} & Gaps782790 & PROVED & \checkmark & verified \\
\texttt{singular\_series\_finite\_pos\_evenPair\_sevenHundredEightySix} & Gaps782790 & PROVED & \checkmark & verified \\
\texttt{singular\_series\_finite\_pos\_evenPair\_sevenHundredEightyTwo} & Gaps782790 & PROVED & \checkmark & verified \\
\texttt{singular\_series\_finite\_pos\_evenPair\_sevenHundredNinety} & Gaps782790 & PROVED & \checkmark & verified \\
\texttt{singular\_series\_pos\_evenPair\_sevenHundredEightyEight} & Gaps782790 & PROVED & \checkmark & verified \\
\texttt{singular\_series\_pos\_evenPair\_sevenHundredEightyFour} & Gaps782790 & PROVED & \checkmark & verified \\
\texttt{singular\_series\_pos\_evenPair\_sevenHundredEightySix} & Gaps782790 & PROVED & \checkmark & verified \\
\texttt{singular\_series\_pos\_evenPair\_sevenHundredEightyTwo} & Gaps782790 & PROVED & \checkmark & verified \\
\texttt{singular\_series\_pos\_evenPair\_sevenHundredNinety} & Gaps782790 & PROVED & \checkmark & verified \\
\texttt{evenPair\_card\_eightHundred} & Gaps792800 & PROVED & \checkmark & verified \\
\texttt{evenPair\_card\_sevenHundredNinetyEight} & Gaps792800 & PROVED & \checkmark & verified \\
\texttt{evenPair\_card\_sevenHundredNinetyFour} & Gaps792800 & PROVED & \checkmark & verified \\
\texttt{evenPair\_card\_sevenHundredNinetySix} & Gaps792800 & PROVED & \checkmark & verified \\
\texttt{evenPair\_card\_sevenHundredNinetyTwo} & Gaps792800 & PROVED & \checkmark & verified \\
\texttt{isAdmissible\_evenPair\_eightHundred} & Gaps792800 & PROVED & \checkmark & verified \\
\texttt{isAdmissible\_evenPair\_sevenHundredNinetyEight} & Gaps792800 & PROVED & \checkmark & verified \\
\texttt{isAdmissible\_evenPair\_sevenHundredNinetyFour} & Gaps792800 & PROVED & \checkmark & verified \\
\texttt{isAdmissible\_evenPair\_sevenHundredNinetySix} & Gaps792800 & PROVED & \checkmark & verified \\
\texttt{isAdmissible\_evenPair\_sevenHundredNinetyTwo} & Gaps792800 & PROVED & \checkmark & verified \\
\texttt{localFactor\_eightHundred\_two} & Gaps792800 & PROVED & \checkmark & verified \\
\texttt{localFactor\_sevenHundredNinetyTwo\_two} & Gaps792800 & PROVED & \checkmark & verified \\
\texttt{nu\_p\_eightHundred} & Gaps792800 & PROVED & \checkmark & verified \\
\texttt{nu\_p\_eightHundred\_two} & Gaps792800 & PROVED & \checkmark & verified \\
\texttt{nu\_p\_sevenHundredNinetyEight} & Gaps792800 & PROVED & \checkmark & verified \\
\texttt{nu\_p\_sevenHundredNinetyFour} & Gaps792800 & PROVED & \checkmark & verified \\
\texttt{nu\_p\_sevenHundredNinetySix} & Gaps792800 & PROVED & \checkmark & verified \\
\texttt{nu\_p\_sevenHundredNinetyTwo} & Gaps792800 & PROVED & \checkmark & verified \\
\texttt{nu\_p\_sevenHundredNinetyTwo\_two} & Gaps792800 & PROVED & \checkmark & verified \\
\texttt{singular\_series\_finite\_pos\_evenPair\_eightHundred} & Gaps792800 & PROVED & \checkmark & verified \\
\texttt{singular\_series\_finite\_pos\_evenPair\_sevenHundredNinetyEight} & Gaps792800 & PROVED & \checkmark & verified \\
\texttt{singular\_series\_finite\_pos\_evenPair\_sevenHundredNinetyFour} & Gaps792800 & PROVED & \checkmark & verified \\
\texttt{singular\_series\_finite\_pos\_evenPair\_sevenHundredNinetySix} & Gaps792800 & PROVED & \checkmark & verified \\
\texttt{singular\_series\_finite\_pos\_evenPair\_sevenHundredNinetyTwo} & Gaps792800 & PROVED & \checkmark & verified \\
\texttt{singular\_series\_pos\_evenPair\_eightHundred} & Gaps792800 & PROVED & \checkmark & verified \\
\texttt{singular\_series\_pos\_evenPair\_sevenHundredNinetyEight} & Gaps792800 & PROVED & \checkmark & verified \\
\texttt{singular\_series\_pos\_evenPair\_sevenHundredNinetyFour} & Gaps792800 & PROVED & \checkmark & verified \\
\texttt{singular\_series\_pos\_evenPair\_sevenHundredNinetySix} & Gaps792800 & PROVED & \checkmark & verified \\
\texttt{singular\_series\_pos\_evenPair\_sevenHundredNinetyTwo} & Gaps792800 & PROVED & \checkmark & verified \\
\texttt{evenPair\_card\_eightHundredEight} & Gaps802810 & PROVED & \checkmark & verified \\
\texttt{evenPair\_card\_eightHundredFour} & Gaps802810 & PROVED & \checkmark & verified \\
\texttt{evenPair\_card\_eightHundredSix} & Gaps802810 & PROVED & \checkmark & verified \\
\texttt{evenPair\_card\_eightHundredTen} & Gaps802810 & PROVED & \checkmark & verified \\
\texttt{evenPair\_card\_eightHundredTwo} & Gaps802810 & PROVED & \checkmark & verified \\
\texttt{isAdmissible\_evenPair\_eightHundredEight} & Gaps802810 & PROVED & \checkmark & verified \\
\texttt{isAdmissible\_evenPair\_eightHundredFour} & Gaps802810 & PROVED & \checkmark & verified \\
\texttt{isAdmissible\_evenPair\_eightHundredSix} & Gaps802810 & PROVED & \checkmark & verified \\
\texttt{isAdmissible\_evenPair\_eightHundredTen} & Gaps802810 & PROVED & \checkmark & verified \\
\texttt{isAdmissible\_evenPair\_eightHundredTwo} & Gaps802810 & PROVED & \checkmark & verified \\
\texttt{localFactor\_eightHundredTen\_two} & Gaps802810 & PROVED & \checkmark & verified \\
\texttt{localFactor\_eightHundredTwo\_two} & Gaps802810 & PROVED & \checkmark & verified \\
\texttt{nu\_p\_eightHundredEight} & Gaps802810 & PROVED & \checkmark & verified \\
\texttt{nu\_p\_eightHundredFour} & Gaps802810 & PROVED & \checkmark & verified \\
\texttt{nu\_p\_eightHundredSix} & Gaps802810 & PROVED & \checkmark & verified \\
\texttt{nu\_p\_eightHundredTen} & Gaps802810 & PROVED & \checkmark & verified \\
\texttt{nu\_p\_eightHundredTen\_two} & Gaps802810 & PROVED & \checkmark & verified \\
\texttt{nu\_p\_eightHundredTwo} & Gaps802810 & PROVED & \checkmark & verified \\
\texttt{nu\_p\_eightHundredTwo\_two} & Gaps802810 & PROVED & \checkmark & verified \\
\texttt{singular\_series\_finite\_pos\_evenPair\_eightHundredEight} & Gaps802810 & PROVED & \checkmark & verified \\
\texttt{singular\_series\_finite\_pos\_evenPair\_eightHundredFour} & Gaps802810 & PROVED & \checkmark & verified \\
\texttt{singular\_series\_finite\_pos\_evenPair\_eightHundredSix} & Gaps802810 & PROVED & \checkmark & verified \\
\texttt{singular\_series\_finite\_pos\_evenPair\_eightHundredTen} & Gaps802810 & PROVED & \checkmark & verified \\
\texttt{singular\_series\_finite\_pos\_evenPair\_eightHundredTwo} & Gaps802810 & PROVED & \checkmark & verified \\
\texttt{singular\_series\_pos\_evenPair\_eightHundredEight} & Gaps802810 & PROVED & \checkmark & verified \\
\texttt{singular\_series\_pos\_evenPair\_eightHundredFour} & Gaps802810 & PROVED & \checkmark & verified \\
\texttt{singular\_series\_pos\_evenPair\_eightHundredSix} & Gaps802810 & PROVED & \checkmark & verified \\
\texttt{singular\_series\_pos\_evenPair\_eightHundredTen} & Gaps802810 & PROVED & \checkmark & verified \\
\texttt{singular\_series\_pos\_evenPair\_eightHundredTwo} & Gaps802810 & PROVED & \checkmark & verified \\
\texttt{evenPair\_card\_eightHundredEighteen} & Gaps812820 & PROVED & \checkmark & verified \\
\texttt{evenPair\_card\_eightHundredFourteen} & Gaps812820 & PROVED & \checkmark & verified \\
\texttt{evenPair\_card\_eightHundredSixteen} & Gaps812820 & PROVED & \checkmark & verified \\
\texttt{evenPair\_card\_eightHundredTwelve} & Gaps812820 & PROVED & \checkmark & verified \\
\texttt{evenPair\_card\_eightHundredTwenty} & Gaps812820 & PROVED & \checkmark & verified \\
\texttt{isAdmissible\_evenPair\_eightHundredEighteen} & Gaps812820 & PROVED & \checkmark & verified \\
\texttt{isAdmissible\_evenPair\_eightHundredFourteen} & Gaps812820 & PROVED & \checkmark & verified \\
\texttt{isAdmissible\_evenPair\_eightHundredSixteen} & Gaps812820 & PROVED & \checkmark & verified \\
\texttt{isAdmissible\_evenPair\_eightHundredTwelve} & Gaps812820 & PROVED & \checkmark & verified \\
\texttt{isAdmissible\_evenPair\_eightHundredTwenty} & Gaps812820 & PROVED & \checkmark & verified \\
\texttt{localFactor\_eightHundredTwelve\_two} & Gaps812820 & PROVED & \checkmark & verified \\
\texttt{localFactor\_eightHundredTwenty\_two} & Gaps812820 & PROVED & \checkmark & verified \\
\texttt{nu\_p\_eightHundredEighteen} & Gaps812820 & PROVED & \checkmark & verified \\
\texttt{nu\_p\_eightHundredFourteen} & Gaps812820 & PROVED & \checkmark & verified \\
\texttt{nu\_p\_eightHundredSixteen} & Gaps812820 & PROVED & \checkmark & verified \\
\texttt{nu\_p\_eightHundredTwelve} & Gaps812820 & PROVED & \checkmark & verified \\
\texttt{nu\_p\_eightHundredTwelve\_two} & Gaps812820 & PROVED & \checkmark & verified \\
\texttt{nu\_p\_eightHundredTwenty} & Gaps812820 & PROVED & \checkmark & verified \\
\texttt{nu\_p\_eightHundredTwenty\_two} & Gaps812820 & PROVED & \checkmark & verified \\
\texttt{singular\_series\_finite\_pos\_evenPair\_eightHundredEighteen} & Gaps812820 & PROVED & \checkmark & verified \\
\texttt{singular\_series\_finite\_pos\_evenPair\_eightHundredFourteen} & Gaps812820 & PROVED & \checkmark & verified \\
\texttt{singular\_series\_finite\_pos\_evenPair\_eightHundredSixteen} & Gaps812820 & PROVED & \checkmark & verified \\
\texttt{singular\_series\_finite\_pos\_evenPair\_eightHundredTwelve} & Gaps812820 & PROVED & \checkmark & verified \\
\texttt{singular\_series\_finite\_pos\_evenPair\_eightHundredTwenty} & Gaps812820 & PROVED & \checkmark & verified \\
\texttt{singular\_series\_pos\_evenPair\_eightHundredEighteen} & Gaps812820 & PROVED & \checkmark & verified \\
\texttt{singular\_series\_pos\_evenPair\_eightHundredFourteen} & Gaps812820 & PROVED & \checkmark & verified \\
\texttt{singular\_series\_pos\_evenPair\_eightHundredSixteen} & Gaps812820 & PROVED & \checkmark & verified \\
\texttt{singular\_series\_pos\_evenPair\_eightHundredTwelve} & Gaps812820 & PROVED & \checkmark & verified \\
\texttt{singular\_series\_pos\_evenPair\_eightHundredTwenty} & Gaps812820 & PROVED & \checkmark & verified \\
\texttt{evenPair\_card\_eightHundredThirty} & Gaps822830 & PROVED & \checkmark & verified \\
\texttt{evenPair\_card\_eightHundredTwentyEight} & Gaps822830 & PROVED & \checkmark & verified \\
\texttt{evenPair\_card\_eightHundredTwentyFour} & Gaps822830 & PROVED & \checkmark & verified \\
\texttt{evenPair\_card\_eightHundredTwentySix} & Gaps822830 & PROVED & \checkmark & verified \\
\texttt{evenPair\_card\_eightHundredTwentyTwo} & Gaps822830 & PROVED & \checkmark & verified \\
\texttt{isAdmissible\_evenPair\_eightHundredThirty} & Gaps822830 & PROVED & \checkmark & verified \\
\texttt{isAdmissible\_evenPair\_eightHundredTwentyEight} & Gaps822830 & PROVED & \checkmark & verified \\
\texttt{isAdmissible\_evenPair\_eightHundredTwentyFour} & Gaps822830 & PROVED & \checkmark & verified \\
\texttt{isAdmissible\_evenPair\_eightHundredTwentySix} & Gaps822830 & PROVED & \checkmark & verified \\
\texttt{isAdmissible\_evenPair\_eightHundredTwentyTwo} & Gaps822830 & PROVED & \checkmark & verified \\
\texttt{localFactor\_eightHundredThirty\_two} & Gaps822830 & PROVED & \checkmark & verified \\
\texttt{localFactor\_eightHundredTwentyTwo\_two} & Gaps822830 & PROVED & \checkmark & verified \\
\texttt{nu\_p\_eightHundredThirty} & Gaps822830 & PROVED & \checkmark & verified \\
\texttt{nu\_p\_eightHundredThirty\_two} & Gaps822830 & PROVED & \checkmark & verified \\
\texttt{nu\_p\_eightHundredTwentyEight} & Gaps822830 & PROVED & \checkmark & verified \\
\texttt{nu\_p\_eightHundredTwentyFour} & Gaps822830 & PROVED & \checkmark & verified \\
\texttt{nu\_p\_eightHundredTwentySix} & Gaps822830 & PROVED & \checkmark & verified \\
\texttt{nu\_p\_eightHundredTwentyTwo} & Gaps822830 & PROVED & \checkmark & verified \\
\texttt{nu\_p\_eightHundredTwentyTwo\_two} & Gaps822830 & PROVED & \checkmark & verified \\
\texttt{singular\_series\_finite\_pos\_evenPair\_eightHundredThirty} & Gaps822830 & PROVED & \checkmark & verified \\
\texttt{singular\_series\_finite\_pos\_evenPair\_eightHundredTwentyEight} & Gaps822830 & PROVED & \checkmark & verified \\
\texttt{singular\_series\_finite\_pos\_evenPair\_eightHundredTwentyFour} & Gaps822830 & PROVED & \checkmark & verified \\
\texttt{singular\_series\_finite\_pos\_evenPair\_eightHundredTwentySix} & Gaps822830 & PROVED & \checkmark & verified \\
\texttt{singular\_series\_finite\_pos\_evenPair\_eightHundredTwentyTwo} & Gaps822830 & PROVED & \checkmark & verified \\
\texttt{singular\_series\_pos\_evenPair\_eightHundredThirty} & Gaps822830 & PROVED & \checkmark & verified \\
\texttt{singular\_series\_pos\_evenPair\_eightHundredTwentyEight} & Gaps822830 & PROVED & \checkmark & verified \\
\texttt{singular\_series\_pos\_evenPair\_eightHundredTwentyFour} & Gaps822830 & PROVED & \checkmark & verified \\
\texttt{singular\_series\_pos\_evenPair\_eightHundredTwentySix} & Gaps822830 & PROVED & \checkmark & verified \\
\texttt{singular\_series\_pos\_evenPair\_eightHundredTwentyTwo} & Gaps822830 & PROVED & \checkmark & verified \\
\texttt{evenPair\_card\_eightyEight} & Gaps8290 & PROVED & \checkmark & verified \\
\texttt{evenPair\_card\_eightyFour} & Gaps8290 & PROVED & \checkmark & verified \\
\texttt{evenPair\_card\_eightySix} & Gaps8290 & PROVED & \checkmark & verified \\
\texttt{evenPair\_card\_eightyTwo} & Gaps8290 & PROVED & \checkmark & verified \\
\texttt{evenPair\_card\_ninety} & Gaps8290 & PROVED & \checkmark & verified \\
\texttt{isAdmissible\_evenPair\_eightyEight} & Gaps8290 & PROVED & \checkmark & verified \\
\texttt{isAdmissible\_evenPair\_eightyFour} & Gaps8290 & PROVED & \checkmark & verified \\
\texttt{isAdmissible\_evenPair\_eightySix} & Gaps8290 & PROVED & \checkmark & verified \\
\texttt{isAdmissible\_evenPair\_eightyTwo} & Gaps8290 & PROVED & \checkmark & verified \\
\texttt{isAdmissible\_evenPair\_ninety} & Gaps8290 & PROVED & \checkmark & verified \\
\texttt{localFactor\_eightyTwo\_two} & Gaps8290 & PROVED & \checkmark & verified \\
\texttt{localFactor\_ninety\_two} & Gaps8290 & PROVED & \checkmark & verified \\
\texttt{nu\_p\_eightyEight} & Gaps8290 & PROVED & \checkmark & verified \\
\texttt{nu\_p\_eightyFour} & Gaps8290 & PROVED & \checkmark & verified \\
\texttt{nu\_p\_eightySix} & Gaps8290 & PROVED & \checkmark & verified \\
\texttt{nu\_p\_eightyTwo} & Gaps8290 & PROVED & \checkmark & verified \\
\texttt{nu\_p\_eightyTwo\_two} & Gaps8290 & PROVED & \checkmark & verified \\
\texttt{nu\_p\_ninety} & Gaps8290 & PROVED & \checkmark & verified \\
\texttt{nu\_p\_ninety\_two} & Gaps8290 & PROVED & \checkmark & verified \\
\texttt{singular\_series\_finite\_pos\_evenPair\_eightyEight} & Gaps8290 & PROVED & \checkmark & verified \\
\texttt{singular\_series\_finite\_pos\_evenPair\_eightyFour} & Gaps8290 & PROVED & \checkmark & verified \\
\texttt{singular\_series\_finite\_pos\_evenPair\_eightySix} & Gaps8290 & PROVED & \checkmark & verified \\
\texttt{singular\_series\_finite\_pos\_evenPair\_eightyTwo} & Gaps8290 & PROVED & \checkmark & verified \\
\texttt{singular\_series\_finite\_pos\_evenPair\_ninety} & Gaps8290 & PROVED & \checkmark & verified \\
\texttt{singular\_series\_pos\_evenPair\_eightyEight} & Gaps8290 & PROVED & \checkmark & verified \\
\texttt{singular\_series\_pos\_evenPair\_eightyFour} & Gaps8290 & PROVED & \checkmark & verified \\
\texttt{singular\_series\_pos\_evenPair\_eightySix} & Gaps8290 & PROVED & \checkmark & verified \\
\texttt{singular\_series\_pos\_evenPair\_eightyTwo} & Gaps8290 & PROVED & \checkmark & verified \\
\texttt{singular\_series\_pos\_evenPair\_ninety} & Gaps8290 & PROVED & \checkmark & verified \\
\texttt{evenPair\_card\_eightHundredForty} & Gaps832840 & PROVED & \checkmark & verified \\
\texttt{evenPair\_card\_eightHundredThirtyEight} & Gaps832840 & PROVED & \checkmark & verified \\
\texttt{evenPair\_card\_eightHundredThirtyFour} & Gaps832840 & PROVED & \checkmark & verified \\
\texttt{evenPair\_card\_eightHundredThirtySix} & Gaps832840 & PROVED & \checkmark & verified \\
\texttt{evenPair\_card\_eightHundredThirtyTwo} & Gaps832840 & PROVED & \checkmark & verified \\
\texttt{isAdmissible\_evenPair\_eightHundredForty} & Gaps832840 & PROVED & \checkmark & verified \\
\texttt{isAdmissible\_evenPair\_eightHundredThirtyEight} & Gaps832840 & PROVED & \checkmark & verified \\
\texttt{isAdmissible\_evenPair\_eightHundredThirtyFour} & Gaps832840 & PROVED & \checkmark & verified \\
\texttt{isAdmissible\_evenPair\_eightHundredThirtySix} & Gaps832840 & PROVED & \checkmark & verified \\
\texttt{isAdmissible\_evenPair\_eightHundredThirtyTwo} & Gaps832840 & PROVED & \checkmark & verified \\
\texttt{localFactor\_eightHundredForty\_two} & Gaps832840 & PROVED & \checkmark & verified \\
\texttt{localFactor\_eightHundredThirtyTwo\_two} & Gaps832840 & PROVED & \checkmark & verified \\
\texttt{nu\_p\_eightHundredForty} & Gaps832840 & PROVED & \checkmark & verified \\
\texttt{nu\_p\_eightHundredForty\_two} & Gaps832840 & PROVED & \checkmark & verified \\
\texttt{nu\_p\_eightHundredThirtyEight} & Gaps832840 & PROVED & \checkmark & verified \\
\texttt{nu\_p\_eightHundredThirtyFour} & Gaps832840 & PROVED & \checkmark & verified \\
\texttt{nu\_p\_eightHundredThirtySix} & Gaps832840 & PROVED & \checkmark & verified \\
\texttt{nu\_p\_eightHundredThirtyTwo} & Gaps832840 & PROVED & \checkmark & verified \\
\texttt{nu\_p\_eightHundredThirtyTwo\_two} & Gaps832840 & PROVED & \checkmark & verified \\
\texttt{singular\_series\_finite\_pos\_evenPair\_eightHundredForty} & Gaps832840 & PROVED & \checkmark & verified \\
\texttt{singular\_series\_finite\_pos\_evenPair\_eightHundredThirtyEight} & Gaps832840 & PROVED & \checkmark & verified \\
\texttt{singular\_series\_finite\_pos\_evenPair\_eightHundredThirtyFour} & Gaps832840 & PROVED & \checkmark & verified \\
\texttt{singular\_series\_finite\_pos\_evenPair\_eightHundredThirtySix} & Gaps832840 & PROVED & \checkmark & verified \\
\texttt{singular\_series\_finite\_pos\_evenPair\_eightHundredThirtyTwo} & Gaps832840 & PROVED & \checkmark & verified \\
\texttt{singular\_series\_pos\_evenPair\_eightHundredForty} & Gaps832840 & PROVED & \checkmark & verified \\
\texttt{singular\_series\_pos\_evenPair\_eightHundredThirtyEight} & Gaps832840 & PROVED & \checkmark & verified \\
\texttt{singular\_series\_pos\_evenPair\_eightHundredThirtyFour} & Gaps832840 & PROVED & \checkmark & verified \\
\texttt{singular\_series\_pos\_evenPair\_eightHundredThirtySix} & Gaps832840 & PROVED & \checkmark & verified \\
\texttt{singular\_series\_pos\_evenPair\_eightHundredThirtyTwo} & Gaps832840 & PROVED & \checkmark & verified \\
\texttt{evenPair\_card\_eightHundredFifty} & Gaps842850 & PROVED & \checkmark & verified \\
\texttt{evenPair\_card\_eightHundredFortyEight} & Gaps842850 & PROVED & \checkmark & verified \\
\texttt{evenPair\_card\_eightHundredFortyFour} & Gaps842850 & PROVED & \checkmark & verified \\
\texttt{evenPair\_card\_eightHundredFortySix} & Gaps842850 & PROVED & \checkmark & verified \\
\texttt{evenPair\_card\_eightHundredFortyTwo} & Gaps842850 & PROVED & \checkmark & verified \\
\texttt{isAdmissible\_evenPair\_eightHundredFifty} & Gaps842850 & PROVED & \checkmark & verified \\
\texttt{isAdmissible\_evenPair\_eightHundredFortyEight} & Gaps842850 & PROVED & \checkmark & verified \\
\texttt{isAdmissible\_evenPair\_eightHundredFortyFour} & Gaps842850 & PROVED & \checkmark & verified \\
\texttt{isAdmissible\_evenPair\_eightHundredFortySix} & Gaps842850 & PROVED & \checkmark & verified \\
\texttt{isAdmissible\_evenPair\_eightHundredFortyTwo} & Gaps842850 & PROVED & \checkmark & verified \\
\texttt{localFactor\_eightHundredFifty\_two} & Gaps842850 & PROVED & \checkmark & verified \\
\texttt{localFactor\_eightHundredFortyTwo\_two} & Gaps842850 & PROVED & \checkmark & verified \\
\texttt{nu\_p\_eightHundredFifty} & Gaps842850 & PROVED & \checkmark & verified \\
\texttt{nu\_p\_eightHundredFifty\_two} & Gaps842850 & PROVED & \checkmark & verified \\
\texttt{nu\_p\_eightHundredFortyEight} & Gaps842850 & PROVED & \checkmark & verified \\
\texttt{nu\_p\_eightHundredFortyFour} & Gaps842850 & PROVED & \checkmark & verified \\
\texttt{nu\_p\_eightHundredFortySix} & Gaps842850 & PROVED & \checkmark & verified \\
\texttt{nu\_p\_eightHundredFortyTwo} & Gaps842850 & PROVED & \checkmark & verified \\
\texttt{nu\_p\_eightHundredFortyTwo\_two} & Gaps842850 & PROVED & \checkmark & verified \\
\texttt{singular\_series\_finite\_pos\_evenPair\_eightHundredFifty} & Gaps842850 & PROVED & \checkmark & verified \\
\texttt{singular\_series\_finite\_pos\_evenPair\_eightHundredFortyEight} & Gaps842850 & PROVED & \checkmark & verified \\
\texttt{singular\_series\_finite\_pos\_evenPair\_eightHundredFortyFour} & Gaps842850 & PROVED & \checkmark & verified \\
\texttt{singular\_series\_finite\_pos\_evenPair\_eightHundredFortySix} & Gaps842850 & PROVED & \checkmark & verified \\
\texttt{singular\_series\_finite\_pos\_evenPair\_eightHundredFortyTwo} & Gaps842850 & PROVED & \checkmark & verified \\
\texttt{singular\_series\_pos\_evenPair\_eightHundredFifty} & Gaps842850 & PROVED & \checkmark & verified \\
\texttt{singular\_series\_pos\_evenPair\_eightHundredFortyEight} & Gaps842850 & PROVED & \checkmark & verified \\
\texttt{singular\_series\_pos\_evenPair\_eightHundredFortyFour} & Gaps842850 & PROVED & \checkmark & verified \\
\texttt{singular\_series\_pos\_evenPair\_eightHundredFortySix} & Gaps842850 & PROVED & \checkmark & verified \\
\texttt{singular\_series\_pos\_evenPair\_eightHundredFortyTwo} & Gaps842850 & PROVED & \checkmark & verified \\
\texttt{evenPair\_card\_eightHundredFiftyEight} & Gaps852860 & PROVED & \checkmark & verified \\
\texttt{evenPair\_card\_eightHundredFiftyFour} & Gaps852860 & PROVED & \checkmark & verified \\
\texttt{evenPair\_card\_eightHundredFiftySix} & Gaps852860 & PROVED & \checkmark & verified \\
\texttt{evenPair\_card\_eightHundredFiftyTwo} & Gaps852860 & PROVED & \checkmark & verified \\
\texttt{evenPair\_card\_eightHundredSixty} & Gaps852860 & PROVED & \checkmark & verified \\
\texttt{isAdmissible\_evenPair\_eightHundredFiftyEight} & Gaps852860 & PROVED & \checkmark & verified \\
\texttt{isAdmissible\_evenPair\_eightHundredFiftyFour} & Gaps852860 & PROVED & \checkmark & verified \\
\texttt{isAdmissible\_evenPair\_eightHundredFiftySix} & Gaps852860 & PROVED & \checkmark & verified \\
\texttt{isAdmissible\_evenPair\_eightHundredFiftyTwo} & Gaps852860 & PROVED & \checkmark & verified \\
\texttt{isAdmissible\_evenPair\_eightHundredSixty} & Gaps852860 & PROVED & \checkmark & verified \\
\texttt{localFactor\_eightHundredFiftyTwo\_two} & Gaps852860 & PROVED & \checkmark & verified \\
\texttt{localFactor\_eightHundredSixty\_two} & Gaps852860 & PROVED & \checkmark & verified \\
\texttt{nu\_p\_eightHundredFiftyEight} & Gaps852860 & PROVED & \checkmark & verified \\
\texttt{nu\_p\_eightHundredFiftyFour} & Gaps852860 & PROVED & \checkmark & verified \\
\texttt{nu\_p\_eightHundredFiftySix} & Gaps852860 & PROVED & \checkmark & verified \\
\texttt{nu\_p\_eightHundredFiftyTwo} & Gaps852860 & PROVED & \checkmark & verified \\
\texttt{nu\_p\_eightHundredFiftyTwo\_two} & Gaps852860 & PROVED & \checkmark & verified \\
\texttt{nu\_p\_eightHundredSixty} & Gaps852860 & PROVED & \checkmark & verified \\
\texttt{nu\_p\_eightHundredSixty\_two} & Gaps852860 & PROVED & \checkmark & verified \\
\texttt{singular\_series\_finite\_pos\_evenPair\_eightHundredFiftyEight} & Gaps852860 & PROVED & \checkmark & verified \\
\texttt{singular\_series\_finite\_pos\_evenPair\_eightHundredFiftyFour} & Gaps852860 & PROVED & \checkmark & verified \\
\texttt{singular\_series\_finite\_pos\_evenPair\_eightHundredFiftySix} & Gaps852860 & PROVED & \checkmark & verified \\
\texttt{singular\_series\_finite\_pos\_evenPair\_eightHundredFiftyTwo} & Gaps852860 & PROVED & \checkmark & verified \\
\texttt{singular\_series\_finite\_pos\_evenPair\_eightHundredSixty} & Gaps852860 & PROVED & \checkmark & verified \\
\texttt{singular\_series\_pos\_evenPair\_eightHundredFiftyEight} & Gaps852860 & PROVED & \checkmark & verified \\
\texttt{singular\_series\_pos\_evenPair\_eightHundredFiftyFour} & Gaps852860 & PROVED & \checkmark & verified \\
\texttt{singular\_series\_pos\_evenPair\_eightHundredFiftySix} & Gaps852860 & PROVED & \checkmark & verified \\
\texttt{singular\_series\_pos\_evenPair\_eightHundredFiftyTwo} & Gaps852860 & PROVED & \checkmark & verified \\
\texttt{singular\_series\_pos\_evenPair\_eightHundredSixty} & Gaps852860 & PROVED & \checkmark & verified \\
\texttt{evenPair\_card\_eightHundredSeventy} & Gaps862870 & PROVED & \checkmark & verified \\
\texttt{evenPair\_card\_eightHundredSixtyEight} & Gaps862870 & PROVED & \checkmark & verified \\
\texttt{evenPair\_card\_eightHundredSixtyFour} & Gaps862870 & PROVED & \checkmark & verified \\
\texttt{evenPair\_card\_eightHundredSixtySix} & Gaps862870 & PROVED & \checkmark & verified \\
\texttt{evenPair\_card\_eightHundredSixtyTwo} & Gaps862870 & PROVED & \checkmark & verified \\
\texttt{isAdmissible\_evenPair\_eightHundredSeventy} & Gaps862870 & PROVED & \checkmark & verified \\
\texttt{isAdmissible\_evenPair\_eightHundredSixtyEight} & Gaps862870 & PROVED & \checkmark & verified \\
\texttt{isAdmissible\_evenPair\_eightHundredSixtyFour} & Gaps862870 & PROVED & \checkmark & verified \\
\texttt{isAdmissible\_evenPair\_eightHundredSixtySix} & Gaps862870 & PROVED & \checkmark & verified \\
\texttt{isAdmissible\_evenPair\_eightHundredSixtyTwo} & Gaps862870 & PROVED & \checkmark & verified \\
\texttt{localFactor\_eightHundredSeventy\_two} & Gaps862870 & PROVED & \checkmark & verified \\
\texttt{localFactor\_eightHundredSixtyTwo\_two} & Gaps862870 & PROVED & \checkmark & verified \\
\texttt{nu\_p\_eightHundredSeventy} & Gaps862870 & PROVED & \checkmark & verified \\
\texttt{nu\_p\_eightHundredSeventy\_two} & Gaps862870 & PROVED & \checkmark & verified \\
\texttt{nu\_p\_eightHundredSixtyEight} & Gaps862870 & PROVED & \checkmark & verified \\
\texttt{nu\_p\_eightHundredSixtyFour} & Gaps862870 & PROVED & \checkmark & verified \\
\texttt{nu\_p\_eightHundredSixtySix} & Gaps862870 & PROVED & \checkmark & verified \\
\texttt{nu\_p\_eightHundredSixtyTwo} & Gaps862870 & PROVED & \checkmark & verified \\
\texttt{nu\_p\_eightHundredSixtyTwo\_two} & Gaps862870 & PROVED & \checkmark & verified \\
\texttt{singular\_series\_finite\_pos\_evenPair\_eightHundredSeventy} & Gaps862870 & PROVED & \checkmark & verified \\
\texttt{singular\_series\_finite\_pos\_evenPair\_eightHundredSixtyEight} & Gaps862870 & PROVED & \checkmark & verified \\
\texttt{singular\_series\_finite\_pos\_evenPair\_eightHundredSixtyFour} & Gaps862870 & PROVED & \checkmark & verified \\
\texttt{singular\_series\_finite\_pos\_evenPair\_eightHundredSixtySix} & Gaps862870 & PROVED & \checkmark & verified \\
\texttt{singular\_series\_finite\_pos\_evenPair\_eightHundredSixtyTwo} & Gaps862870 & PROVED & \checkmark & verified \\
\texttt{singular\_series\_pos\_evenPair\_eightHundredSeventy} & Gaps862870 & PROVED & \checkmark & verified \\
\texttt{singular\_series\_pos\_evenPair\_eightHundredSixtyEight} & Gaps862870 & PROVED & \checkmark & verified \\
\texttt{singular\_series\_pos\_evenPair\_eightHundredSixtyFour} & Gaps862870 & PROVED & \checkmark & verified \\
\texttt{singular\_series\_pos\_evenPair\_eightHundredSixtySix} & Gaps862870 & PROVED & \checkmark & verified \\
\texttt{singular\_series\_pos\_evenPair\_eightHundredSixtyTwo} & Gaps862870 & PROVED & \checkmark & verified \\
\texttt{evenPair\_card\_eightHundredEighty} & Gaps872880 & PROVED & \checkmark & verified \\
\texttt{evenPair\_card\_eightHundredSeventyEight} & Gaps872880 & PROVED & \checkmark & verified \\
\texttt{evenPair\_card\_eightHundredSeventyFour} & Gaps872880 & PROVED & \checkmark & verified \\
\texttt{evenPair\_card\_eightHundredSeventySix} & Gaps872880 & PROVED & \checkmark & verified \\
\texttt{evenPair\_card\_eightHundredSeventyTwo} & Gaps872880 & PROVED & \checkmark & verified \\
\texttt{isAdmissible\_evenPair\_eightHundredEighty} & Gaps872880 & PROVED & \checkmark & verified \\
\texttt{isAdmissible\_evenPair\_eightHundredSeventyEight} & Gaps872880 & PROVED & \checkmark & verified \\
\texttt{isAdmissible\_evenPair\_eightHundredSeventyFour} & Gaps872880 & PROVED & \checkmark & verified \\
\texttt{isAdmissible\_evenPair\_eightHundredSeventySix} & Gaps872880 & PROVED & \checkmark & verified \\
\texttt{isAdmissible\_evenPair\_eightHundredSeventyTwo} & Gaps872880 & PROVED & \checkmark & verified \\
\texttt{localFactor\_eightHundredEighty\_two} & Gaps872880 & PROVED & \checkmark & verified \\
\texttt{localFactor\_eightHundredSeventyTwo\_two} & Gaps872880 & PROVED & \checkmark & verified \\
\texttt{nu\_p\_eightHundredEighty} & Gaps872880 & PROVED & \checkmark & verified \\
\texttt{nu\_p\_eightHundredEighty\_two} & Gaps872880 & PROVED & \checkmark & verified \\
\texttt{nu\_p\_eightHundredSeventyEight} & Gaps872880 & PROVED & \checkmark & verified \\
\texttt{nu\_p\_eightHundredSeventyFour} & Gaps872880 & PROVED & \checkmark & verified \\
\texttt{nu\_p\_eightHundredSeventySix} & Gaps872880 & PROVED & \checkmark & verified \\
\texttt{nu\_p\_eightHundredSeventyTwo} & Gaps872880 & PROVED & \checkmark & verified \\
\texttt{nu\_p\_eightHundredSeventyTwo\_two} & Gaps872880 & PROVED & \checkmark & verified \\
\texttt{singular\_series\_finite\_pos\_evenPair\_eightHundredEighty} & Gaps872880 & PROVED & \checkmark & verified \\
\texttt{singular\_series\_finite\_pos\_evenPair\_eightHundredSeventyEight} & Gaps872880 & PROVED & \checkmark & verified \\
\texttt{singular\_series\_finite\_pos\_evenPair\_eightHundredSeventyFour} & Gaps872880 & PROVED & \checkmark & verified \\
\texttt{singular\_series\_finite\_pos\_evenPair\_eightHundredSeventySix} & Gaps872880 & PROVED & \checkmark & verified \\
\texttt{singular\_series\_finite\_pos\_evenPair\_eightHundredSeventyTwo} & Gaps872880 & PROVED & \checkmark & verified \\
\texttt{singular\_series\_pos\_evenPair\_eightHundredEighty} & Gaps872880 & PROVED & \checkmark & verified \\
\texttt{singular\_series\_pos\_evenPair\_eightHundredSeventyEight} & Gaps872880 & PROVED & \checkmark & verified \\
\texttt{singular\_series\_pos\_evenPair\_eightHundredSeventyFour} & Gaps872880 & PROVED & \checkmark & verified \\
\texttt{singular\_series\_pos\_evenPair\_eightHundredSeventySix} & Gaps872880 & PROVED & \checkmark & verified \\
\texttt{singular\_series\_pos\_evenPair\_eightHundredSeventyTwo} & Gaps872880 & PROVED & \checkmark & verified \\
\texttt{evenPair\_card\_eightHundredEightyEight} & Gaps882890 & PROVED & \checkmark & verified \\
\texttt{evenPair\_card\_eightHundredEightyFour} & Gaps882890 & PROVED & \checkmark & verified \\
\texttt{evenPair\_card\_eightHundredEightySix} & Gaps882890 & PROVED & \checkmark & verified \\
\texttt{evenPair\_card\_eightHundredEightyTwo} & Gaps882890 & PROVED & \checkmark & verified \\
\texttt{evenPair\_card\_eightHundredNinety} & Gaps882890 & PROVED & \checkmark & verified \\
\texttt{isAdmissible\_evenPair\_eightHundredEightyEight} & Gaps882890 & PROVED & \checkmark & verified \\
\texttt{isAdmissible\_evenPair\_eightHundredEightyFour} & Gaps882890 & PROVED & \checkmark & verified \\
\texttt{isAdmissible\_evenPair\_eightHundredEightySix} & Gaps882890 & PROVED & \checkmark & verified \\
\texttt{isAdmissible\_evenPair\_eightHundredEightyTwo} & Gaps882890 & PROVED & \checkmark & verified \\
\texttt{isAdmissible\_evenPair\_eightHundredNinety} & Gaps882890 & PROVED & \checkmark & verified \\
\texttt{localFactor\_eightHundredEightyTwo\_two} & Gaps882890 & PROVED & \checkmark & verified \\
\texttt{localFactor\_eightHundredNinety\_two} & Gaps882890 & PROVED & \checkmark & verified \\
\texttt{nu\_p\_eightHundredEightyEight} & Gaps882890 & PROVED & \checkmark & verified \\
\texttt{nu\_p\_eightHundredEightyFour} & Gaps882890 & PROVED & \checkmark & verified \\
\texttt{nu\_p\_eightHundredEightySix} & Gaps882890 & PROVED & \checkmark & verified \\
\texttt{nu\_p\_eightHundredEightyTwo} & Gaps882890 & PROVED & \checkmark & verified \\
\texttt{nu\_p\_eightHundredEightyTwo\_two} & Gaps882890 & PROVED & \checkmark & verified \\
\texttt{nu\_p\_eightHundredNinety} & Gaps882890 & PROVED & \checkmark & verified \\
\texttt{nu\_p\_eightHundredNinety\_two} & Gaps882890 & PROVED & \checkmark & verified \\
\texttt{singular\_series\_finite\_pos\_evenPair\_eightHundredEightyEight} & Gaps882890 & PROVED & \checkmark & verified \\
\texttt{singular\_series\_finite\_pos\_evenPair\_eightHundredEightyFour} & Gaps882890 & PROVED & \checkmark & verified \\
\texttt{singular\_series\_finite\_pos\_evenPair\_eightHundredEightySix} & Gaps882890 & PROVED & \checkmark & verified \\
\texttt{singular\_series\_finite\_pos\_evenPair\_eightHundredEightyTwo} & Gaps882890 & PROVED & \checkmark & verified \\
\texttt{singular\_series\_finite\_pos\_evenPair\_eightHundredNinety} & Gaps882890 & PROVED & \checkmark & verified \\
\texttt{singular\_series\_pos\_evenPair\_eightHundredEightyEight} & Gaps882890 & PROVED & \checkmark & verified \\
\texttt{singular\_series\_pos\_evenPair\_eightHundredEightyFour} & Gaps882890 & PROVED & \checkmark & verified \\
\texttt{singular\_series\_pos\_evenPair\_eightHundredEightySix} & Gaps882890 & PROVED & \checkmark & verified \\
\texttt{singular\_series\_pos\_evenPair\_eightHundredEightyTwo} & Gaps882890 & PROVED & \checkmark & verified \\
\texttt{singular\_series\_pos\_evenPair\_eightHundredNinety} & Gaps882890 & PROVED & \checkmark & verified \\
\texttt{evenPair\_card\_eightHundredNinetyEight} & Gaps892900 & PROVED & \checkmark & verified \\
\texttt{evenPair\_card\_eightHundredNinetyFour} & Gaps892900 & PROVED & \checkmark & verified \\
\texttt{evenPair\_card\_eightHundredNinetySix} & Gaps892900 & PROVED & \checkmark & verified \\
\texttt{evenPair\_card\_eightHundredNinetyTwo} & Gaps892900 & PROVED & \checkmark & verified \\
\texttt{evenPair\_card\_nineHundred} & Gaps892900 & PROVED & \checkmark & verified \\
\texttt{isAdmissible\_evenPair\_eightHundredNinetyEight} & Gaps892900 & PROVED & \checkmark & verified \\
\texttt{isAdmissible\_evenPair\_eightHundredNinetyFour} & Gaps892900 & PROVED & \checkmark & verified \\
\texttt{isAdmissible\_evenPair\_eightHundredNinetySix} & Gaps892900 & PROVED & \checkmark & verified \\
\texttt{isAdmissible\_evenPair\_eightHundredNinetyTwo} & Gaps892900 & PROVED & \checkmark & verified \\
\texttt{isAdmissible\_evenPair\_nineHundred} & Gaps892900 & PROVED & \checkmark & verified \\
\texttt{localFactor\_eightHundredNinetyTwo\_two} & Gaps892900 & PROVED & \checkmark & verified \\
\texttt{localFactor\_nineHundred\_two} & Gaps892900 & PROVED & \checkmark & verified \\
\texttt{nu\_p\_eightHundredNinetyEight} & Gaps892900 & PROVED & \checkmark & verified \\
\texttt{nu\_p\_eightHundredNinetyFour} & Gaps892900 & PROVED & \checkmark & verified \\
\texttt{nu\_p\_eightHundredNinetySix} & Gaps892900 & PROVED & \checkmark & verified \\
\texttt{nu\_p\_eightHundredNinetyTwo} & Gaps892900 & PROVED & \checkmark & verified \\
\texttt{nu\_p\_eightHundredNinetyTwo\_two} & Gaps892900 & PROVED & \checkmark & verified \\
\texttt{nu\_p\_nineHundred} & Gaps892900 & PROVED & \checkmark & verified \\
\texttt{nu\_p\_nineHundred\_two} & Gaps892900 & PROVED & \checkmark & verified \\
\texttt{singular\_series\_finite\_pos\_evenPair\_eightHundredNinetyEight} & Gaps892900 & PROVED & \checkmark & verified \\
\texttt{singular\_series\_finite\_pos\_evenPair\_eightHundredNinetyFour} & Gaps892900 & PROVED & \checkmark & verified \\
\texttt{singular\_series\_finite\_pos\_evenPair\_eightHundredNinetySix} & Gaps892900 & PROVED & \checkmark & verified \\
\texttt{singular\_series\_finite\_pos\_evenPair\_eightHundredNinetyTwo} & Gaps892900 & PROVED & \checkmark & verified \\
\texttt{singular\_series\_finite\_pos\_evenPair\_nineHundred} & Gaps892900 & PROVED & \checkmark & verified \\
\texttt{singular\_series\_pos\_evenPair\_eightHundredNinetyEight} & Gaps892900 & PROVED & \checkmark & verified \\
\texttt{singular\_series\_pos\_evenPair\_eightHundredNinetyFour} & Gaps892900 & PROVED & \checkmark & verified \\
\texttt{singular\_series\_pos\_evenPair\_eightHundredNinetySix} & Gaps892900 & PROVED & \checkmark & verified \\
\texttt{singular\_series\_pos\_evenPair\_eightHundredNinetyTwo} & Gaps892900 & PROVED & \checkmark & verified \\
\texttt{singular\_series\_pos\_evenPair\_nineHundred} & Gaps892900 & PROVED & \checkmark & verified \\
\texttt{evenPair\_card\_nineHundredEight} & Gaps902910 & PROVED & \checkmark & verified \\
\texttt{evenPair\_card\_nineHundredFour} & Gaps902910 & PROVED & \checkmark & verified \\
\texttt{evenPair\_card\_nineHundredSix} & Gaps902910 & PROVED & \checkmark & verified \\
\texttt{evenPair\_card\_nineHundredTen} & Gaps902910 & PROVED & \checkmark & verified \\
\texttt{evenPair\_card\_nineHundredTwo} & Gaps902910 & PROVED & \checkmark & verified \\
\texttt{isAdmissible\_evenPair\_nineHundredEight} & Gaps902910 & PROVED & \checkmark & verified \\
\texttt{isAdmissible\_evenPair\_nineHundredFour} & Gaps902910 & PROVED & \checkmark & verified \\
\texttt{isAdmissible\_evenPair\_nineHundredSix} & Gaps902910 & PROVED & \checkmark & verified \\
\texttt{isAdmissible\_evenPair\_nineHundredTen} & Gaps902910 & PROVED & \checkmark & verified \\
\texttt{isAdmissible\_evenPair\_nineHundredTwo} & Gaps902910 & PROVED & \checkmark & verified \\
\texttt{localFactor\_nineHundredTen\_two} & Gaps902910 & PROVED & \checkmark & verified \\
\texttt{localFactor\_nineHundredTwo\_two} & Gaps902910 & PROVED & \checkmark & verified \\
\texttt{nu\_p\_nineHundredEight} & Gaps902910 & PROVED & \checkmark & verified \\
\texttt{nu\_p\_nineHundredFour} & Gaps902910 & PROVED & \checkmark & verified \\
\texttt{nu\_p\_nineHundredSix} & Gaps902910 & PROVED & \checkmark & verified \\
\texttt{nu\_p\_nineHundredTen} & Gaps902910 & PROVED & \checkmark & verified \\
\texttt{nu\_p\_nineHundredTen\_two} & Gaps902910 & PROVED & \checkmark & verified \\
\texttt{nu\_p\_nineHundredTwo} & Gaps902910 & PROVED & \checkmark & verified \\
\texttt{nu\_p\_nineHundredTwo\_two} & Gaps902910 & PROVED & \checkmark & verified \\
\texttt{singular\_series\_finite\_pos\_evenPair\_nineHundredEight} & Gaps902910 & PROVED & \checkmark & verified \\
\texttt{singular\_series\_finite\_pos\_evenPair\_nineHundredFour} & Gaps902910 & PROVED & \checkmark & verified \\
\texttt{singular\_series\_finite\_pos\_evenPair\_nineHundredSix} & Gaps902910 & PROVED & \checkmark & verified \\
\texttt{singular\_series\_finite\_pos\_evenPair\_nineHundredTen} & Gaps902910 & PROVED & \checkmark & verified \\
\texttt{singular\_series\_finite\_pos\_evenPair\_nineHundredTwo} & Gaps902910 & PROVED & \checkmark & verified \\
\texttt{singular\_series\_pos\_evenPair\_nineHundredEight} & Gaps902910 & PROVED & \checkmark & verified \\
\texttt{singular\_series\_pos\_evenPair\_nineHundredFour} & Gaps902910 & PROVED & \checkmark & verified \\
\texttt{singular\_series\_pos\_evenPair\_nineHundredSix} & Gaps902910 & PROVED & \checkmark & verified \\
\texttt{singular\_series\_pos\_evenPair\_nineHundredTen} & Gaps902910 & PROVED & \checkmark & verified \\
\texttt{singular\_series\_pos\_evenPair\_nineHundredTwo} & Gaps902910 & PROVED & \checkmark & verified \\
\texttt{evenPair\_card\_nineHundredEighteen} & Gaps912920 & PROVED & \checkmark & verified \\
\texttt{evenPair\_card\_nineHundredFourteen} & Gaps912920 & PROVED & \checkmark & verified \\
\texttt{evenPair\_card\_nineHundredSixteen} & Gaps912920 & PROVED & \checkmark & verified \\
\texttt{evenPair\_card\_nineHundredTwelve} & Gaps912920 & PROVED & \checkmark & verified \\
\texttt{evenPair\_card\_nineHundredTwenty} & Gaps912920 & PROVED & \checkmark & verified \\
\texttt{isAdmissible\_evenPair\_nineHundredEighteen} & Gaps912920 & PROVED & \checkmark & verified \\
\texttt{isAdmissible\_evenPair\_nineHundredFourteen} & Gaps912920 & PROVED & \checkmark & verified \\
\texttt{isAdmissible\_evenPair\_nineHundredSixteen} & Gaps912920 & PROVED & \checkmark & verified \\
\texttt{isAdmissible\_evenPair\_nineHundredTwelve} & Gaps912920 & PROVED & \checkmark & verified \\
\texttt{isAdmissible\_evenPair\_nineHundredTwenty} & Gaps912920 & PROVED & \checkmark & verified \\
\texttt{localFactor\_nineHundredTwelve\_two} & Gaps912920 & PROVED & \checkmark & verified \\
\texttt{localFactor\_nineHundredTwenty\_two} & Gaps912920 & PROVED & \checkmark & verified \\
\texttt{nu\_p\_nineHundredEighteen} & Gaps912920 & PROVED & \checkmark & verified \\
\texttt{nu\_p\_nineHundredFourteen} & Gaps912920 & PROVED & \checkmark & verified \\
\texttt{nu\_p\_nineHundredSixteen} & Gaps912920 & PROVED & \checkmark & verified \\
\texttt{nu\_p\_nineHundredTwelve} & Gaps912920 & PROVED & \checkmark & verified \\
\texttt{nu\_p\_nineHundredTwelve\_two} & Gaps912920 & PROVED & \checkmark & verified \\
\texttt{nu\_p\_nineHundredTwenty} & Gaps912920 & PROVED & \checkmark & verified \\
\texttt{nu\_p\_nineHundredTwenty\_two} & Gaps912920 & PROVED & \checkmark & verified \\
\texttt{singular\_series\_finite\_pos\_evenPair\_nineHundredEighteen} & Gaps912920 & PROVED & \checkmark & verified \\
\texttt{singular\_series\_finite\_pos\_evenPair\_nineHundredFourteen} & Gaps912920 & PROVED & \checkmark & verified \\
\texttt{singular\_series\_finite\_pos\_evenPair\_nineHundredSixteen} & Gaps912920 & PROVED & \checkmark & verified \\
\texttt{singular\_series\_finite\_pos\_evenPair\_nineHundredTwelve} & Gaps912920 & PROVED & \checkmark & verified \\
\texttt{singular\_series\_finite\_pos\_evenPair\_nineHundredTwenty} & Gaps912920 & PROVED & \checkmark & verified \\
\texttt{singular\_series\_pos\_evenPair\_nineHundredEighteen} & Gaps912920 & PROVED & \checkmark & verified \\
\texttt{singular\_series\_pos\_evenPair\_nineHundredFourteen} & Gaps912920 & PROVED & \checkmark & verified \\
\texttt{singular\_series\_pos\_evenPair\_nineHundredSixteen} & Gaps912920 & PROVED & \checkmark & verified \\
\texttt{singular\_series\_pos\_evenPair\_nineHundredTwelve} & Gaps912920 & PROVED & \checkmark & verified \\
\texttt{singular\_series\_pos\_evenPair\_nineHundredTwenty} & Gaps912920 & PROVED & \checkmark & verified \\
\texttt{evenPair\_card\_ninetyEight} & Gaps92100 & PROVED & \checkmark & verified \\
\texttt{evenPair\_card\_ninetyFour} & Gaps92100 & PROVED & \checkmark & verified \\
\texttt{evenPair\_card\_ninetySix} & Gaps92100 & PROVED & \checkmark & verified \\
\texttt{evenPair\_card\_ninetyTwo} & Gaps92100 & PROVED & \checkmark & verified \\
\texttt{evenPair\_card\_oneHundred} & Gaps92100 & PROVED & \checkmark & verified \\
\texttt{isAdmissible\_evenPair\_ninetyEight} & Gaps92100 & PROVED & \checkmark & verified \\
\texttt{isAdmissible\_evenPair\_ninetyFour} & Gaps92100 & PROVED & \checkmark & verified \\
\texttt{isAdmissible\_evenPair\_ninetySix} & Gaps92100 & PROVED & \checkmark & verified \\
\texttt{isAdmissible\_evenPair\_ninetyTwo} & Gaps92100 & PROVED & \checkmark & verified \\
\texttt{isAdmissible\_evenPair\_oneHundred} & Gaps92100 & PROVED & \checkmark & verified \\
\texttt{localFactor\_ninetyTwo\_two} & Gaps92100 & PROVED & \checkmark & verified \\
\texttt{localFactor\_oneHundred\_two} & Gaps92100 & PROVED & \checkmark & verified \\
\texttt{nu\_p\_ninetyEight} & Gaps92100 & PROVED & \checkmark & verified \\
\texttt{nu\_p\_ninetyFour} & Gaps92100 & PROVED & \checkmark & verified \\
\texttt{nu\_p\_ninetySix} & Gaps92100 & PROVED & \checkmark & verified \\
\texttt{nu\_p\_ninetyTwo} & Gaps92100 & PROVED & \checkmark & verified \\
\texttt{nu\_p\_ninetyTwo\_two} & Gaps92100 & PROVED & \checkmark & verified \\
\texttt{nu\_p\_oneHundred} & Gaps92100 & PROVED & \checkmark & verified \\
\texttt{nu\_p\_oneHundred\_two} & Gaps92100 & PROVED & \checkmark & verified \\
\texttt{singular\_series\_finite\_pos\_evenPair\_ninetyEight} & Gaps92100 & PROVED & \checkmark & verified \\
\texttt{singular\_series\_finite\_pos\_evenPair\_ninetyFour} & Gaps92100 & PROVED & \checkmark & verified \\
\texttt{singular\_series\_finite\_pos\_evenPair\_ninetySix} & Gaps92100 & PROVED & \checkmark & verified \\
\texttt{singular\_series\_finite\_pos\_evenPair\_ninetyTwo} & Gaps92100 & PROVED & \checkmark & verified \\
\texttt{singular\_series\_finite\_pos\_evenPair\_oneHundred} & Gaps92100 & PROVED & \checkmark & verified \\
\texttt{singular\_series\_pos\_evenPair\_ninetyEight} & Gaps92100 & PROVED & \checkmark & verified \\
\texttt{singular\_series\_pos\_evenPair\_ninetyFour} & Gaps92100 & PROVED & \checkmark & verified \\
\texttt{singular\_series\_pos\_evenPair\_ninetySix} & Gaps92100 & PROVED & \checkmark & verified \\
\texttt{singular\_series\_pos\_evenPair\_ninetyTwo} & Gaps92100 & PROVED & \checkmark & verified \\
\texttt{singular\_series\_pos\_evenPair\_oneHundred} & Gaps92100 & PROVED & \checkmark & verified \\
\texttt{evenPair\_card\_nineHundredThirty} & Gaps922930 & PROVED & \checkmark & verified \\
\texttt{evenPair\_card\_nineHundredTwentyEight} & Gaps922930 & PROVED & \checkmark & verified \\
\texttt{evenPair\_card\_nineHundredTwentyFour} & Gaps922930 & PROVED & \checkmark & verified \\
\texttt{evenPair\_card\_nineHundredTwentySix} & Gaps922930 & PROVED & \checkmark & verified \\
\texttt{evenPair\_card\_nineHundredTwentyTwo} & Gaps922930 & PROVED & \checkmark & verified \\
\texttt{isAdmissible\_evenPair\_nineHundredThirty} & Gaps922930 & PROVED & \checkmark & verified \\
\texttt{isAdmissible\_evenPair\_nineHundredTwentyEight} & Gaps922930 & PROVED & \checkmark & verified \\
\texttt{isAdmissible\_evenPair\_nineHundredTwentyFour} & Gaps922930 & PROVED & \checkmark & verified \\
\texttt{isAdmissible\_evenPair\_nineHundredTwentySix} & Gaps922930 & PROVED & \checkmark & verified \\
\texttt{isAdmissible\_evenPair\_nineHundredTwentyTwo} & Gaps922930 & PROVED & \checkmark & verified \\
\texttt{localFactor\_nineHundredThirty\_two} & Gaps922930 & PROVED & \checkmark & verified \\
\texttt{localFactor\_nineHundredTwentyTwo\_two} & Gaps922930 & PROVED & \checkmark & verified \\
\texttt{nu\_p\_nineHundredThirty} & Gaps922930 & PROVED & \checkmark & verified \\
\texttt{nu\_p\_nineHundredThirty\_two} & Gaps922930 & PROVED & \checkmark & verified \\
\texttt{nu\_p\_nineHundredTwentyEight} & Gaps922930 & PROVED & \checkmark & verified \\
\texttt{nu\_p\_nineHundredTwentyFour} & Gaps922930 & PROVED & \checkmark & verified \\
\texttt{nu\_p\_nineHundredTwentySix} & Gaps922930 & PROVED & \checkmark & verified \\
\texttt{nu\_p\_nineHundredTwentyTwo} & Gaps922930 & PROVED & \checkmark & verified \\
\texttt{nu\_p\_nineHundredTwentyTwo\_two} & Gaps922930 & PROVED & \checkmark & verified \\
\texttt{singular\_series\_finite\_pos\_evenPair\_nineHundredThirty} & Gaps922930 & PROVED & \checkmark & verified \\
\texttt{singular\_series\_finite\_pos\_evenPair\_nineHundredTwentyEight} & Gaps922930 & PROVED & \checkmark & verified \\
\texttt{singular\_series\_finite\_pos\_evenPair\_nineHundredTwentyFour} & Gaps922930 & PROVED & \checkmark & verified \\
\texttt{singular\_series\_finite\_pos\_evenPair\_nineHundredTwentySix} & Gaps922930 & PROVED & \checkmark & verified \\
\texttt{singular\_series\_finite\_pos\_evenPair\_nineHundredTwentyTwo} & Gaps922930 & PROVED & \checkmark & verified \\
\texttt{singular\_series\_pos\_evenPair\_nineHundredThirty} & Gaps922930 & PROVED & \checkmark & verified \\
\texttt{singular\_series\_pos\_evenPair\_nineHundredTwentyEight} & Gaps922930 & PROVED & \checkmark & verified \\
\texttt{singular\_series\_pos\_evenPair\_nineHundredTwentyFour} & Gaps922930 & PROVED & \checkmark & verified \\
\texttt{singular\_series\_pos\_evenPair\_nineHundredTwentySix} & Gaps922930 & PROVED & \checkmark & verified \\
\texttt{singular\_series\_pos\_evenPair\_nineHundredTwentyTwo} & Gaps922930 & PROVED & \checkmark & verified \\
\texttt{evenPair\_card\_nineHundredForty} & Gaps932940 & PROVED & \checkmark & verified \\
\texttt{evenPair\_card\_nineHundredThirtyEight} & Gaps932940 & PROVED & \checkmark & verified \\
\texttt{evenPair\_card\_nineHundredThirtyFour} & Gaps932940 & PROVED & \checkmark & verified \\
\texttt{evenPair\_card\_nineHundredThirtySix} & Gaps932940 & PROVED & \checkmark & verified \\
\texttt{evenPair\_card\_nineHundredThirtyTwo} & Gaps932940 & PROVED & \checkmark & verified \\
\texttt{isAdmissible\_evenPair\_nineHundredForty} & Gaps932940 & PROVED & \checkmark & verified \\
\texttt{isAdmissible\_evenPair\_nineHundredThirtyEight} & Gaps932940 & PROVED & \checkmark & verified \\
\texttt{isAdmissible\_evenPair\_nineHundredThirtyFour} & Gaps932940 & PROVED & \checkmark & verified \\
\texttt{isAdmissible\_evenPair\_nineHundredThirtySix} & Gaps932940 & PROVED & \checkmark & verified \\
\texttt{isAdmissible\_evenPair\_nineHundredThirtyTwo} & Gaps932940 & PROVED & \checkmark & verified \\
\texttt{localFactor\_nineHundredForty\_two} & Gaps932940 & PROVED & \checkmark & verified \\
\texttt{localFactor\_nineHundredThirtyTwo\_two} & Gaps932940 & PROVED & \checkmark & verified \\
\texttt{nu\_p\_nineHundredForty} & Gaps932940 & PROVED & \checkmark & verified \\
\texttt{nu\_p\_nineHundredForty\_two} & Gaps932940 & PROVED & \checkmark & verified \\
\texttt{nu\_p\_nineHundredThirtyEight} & Gaps932940 & PROVED & \checkmark & verified \\
\texttt{nu\_p\_nineHundredThirtyFour} & Gaps932940 & PROVED & \checkmark & verified \\
\texttt{nu\_p\_nineHundredThirtySix} & Gaps932940 & PROVED & \checkmark & verified \\
\texttt{nu\_p\_nineHundredThirtyTwo} & Gaps932940 & PROVED & \checkmark & verified \\
\texttt{nu\_p\_nineHundredThirtyTwo\_two} & Gaps932940 & PROVED & \checkmark & verified \\
\texttt{singular\_series\_finite\_pos\_evenPair\_nineHundredForty} & Gaps932940 & PROVED & \checkmark & verified \\
\texttt{singular\_series\_finite\_pos\_evenPair\_nineHundredThirtyEight} & Gaps932940 & PROVED & \checkmark & verified \\
\texttt{singular\_series\_finite\_pos\_evenPair\_nineHundredThirtyFour} & Gaps932940 & PROVED & \checkmark & verified \\
\texttt{singular\_series\_finite\_pos\_evenPair\_nineHundredThirtySix} & Gaps932940 & PROVED & \checkmark & verified \\
\texttt{singular\_series\_finite\_pos\_evenPair\_nineHundredThirtyTwo} & Gaps932940 & PROVED & \checkmark & verified \\
\texttt{singular\_series\_pos\_evenPair\_nineHundredForty} & Gaps932940 & PROVED & \checkmark & verified \\
\texttt{singular\_series\_pos\_evenPair\_nineHundredThirtyEight} & Gaps932940 & PROVED & \checkmark & verified \\
\texttt{singular\_series\_pos\_evenPair\_nineHundredThirtyFour} & Gaps932940 & PROVED & \checkmark & verified \\
\texttt{singular\_series\_pos\_evenPair\_nineHundredThirtySix} & Gaps932940 & PROVED & \checkmark & verified \\
\texttt{singular\_series\_pos\_evenPair\_nineHundredThirtyTwo} & Gaps932940 & PROVED & \checkmark & verified \\
\texttt{evenPair\_card\_nineHundredFifty} & Gaps942950 & PROVED & \checkmark & verified \\
\texttt{evenPair\_card\_nineHundredFortyEight} & Gaps942950 & PROVED & \checkmark & verified \\
\texttt{evenPair\_card\_nineHundredFortyFour} & Gaps942950 & PROVED & \checkmark & verified \\
\texttt{evenPair\_card\_nineHundredFortySix} & Gaps942950 & PROVED & \checkmark & verified \\
\texttt{evenPair\_card\_nineHundredFortyTwo} & Gaps942950 & PROVED & \checkmark & verified \\
\texttt{isAdmissible\_evenPair\_nineHundredFifty} & Gaps942950 & PROVED & \checkmark & verified \\
\texttt{isAdmissible\_evenPair\_nineHundredFortyEight} & Gaps942950 & PROVED & \checkmark & verified \\
\texttt{isAdmissible\_evenPair\_nineHundredFortyFour} & Gaps942950 & PROVED & \checkmark & verified \\
\texttt{isAdmissible\_evenPair\_nineHundredFortySix} & Gaps942950 & PROVED & \checkmark & verified \\
\texttt{isAdmissible\_evenPair\_nineHundredFortyTwo} & Gaps942950 & PROVED & \checkmark & verified \\
\texttt{localFactor\_nineHundredFifty\_two} & Gaps942950 & PROVED & \checkmark & verified \\
\texttt{localFactor\_nineHundredFortyTwo\_two} & Gaps942950 & PROVED & \checkmark & verified \\
\texttt{nu\_p\_nineHundredFifty} & Gaps942950 & PROVED & \checkmark & verified \\
\texttt{nu\_p\_nineHundredFifty\_two} & Gaps942950 & PROVED & \checkmark & verified \\
\texttt{nu\_p\_nineHundredFortyEight} & Gaps942950 & PROVED & \checkmark & verified \\
\texttt{nu\_p\_nineHundredFortyFour} & Gaps942950 & PROVED & \checkmark & verified \\
\texttt{nu\_p\_nineHundredFortySix} & Gaps942950 & PROVED & \checkmark & verified \\
\texttt{nu\_p\_nineHundredFortyTwo} & Gaps942950 & PROVED & \checkmark & verified \\
\texttt{nu\_p\_nineHundredFortyTwo\_two} & Gaps942950 & PROVED & \checkmark & verified \\
\texttt{singular\_series\_finite\_pos\_evenPair\_nineHundredFifty} & Gaps942950 & PROVED & \checkmark & verified \\
\texttt{singular\_series\_finite\_pos\_evenPair\_nineHundredFortyEight} & Gaps942950 & PROVED & \checkmark & verified \\
\texttt{singular\_series\_finite\_pos\_evenPair\_nineHundredFortyFour} & Gaps942950 & PROVED & \checkmark & verified \\
\texttt{singular\_series\_finite\_pos\_evenPair\_nineHundredFortySix} & Gaps942950 & PROVED & \checkmark & verified \\
\texttt{singular\_series\_finite\_pos\_evenPair\_nineHundredFortyTwo} & Gaps942950 & PROVED & \checkmark & verified \\
\texttt{singular\_series\_pos\_evenPair\_nineHundredFifty} & Gaps942950 & PROVED & \checkmark & verified \\
\texttt{singular\_series\_pos\_evenPair\_nineHundredFortyEight} & Gaps942950 & PROVED & \checkmark & verified \\
\texttt{singular\_series\_pos\_evenPair\_nineHundredFortyFour} & Gaps942950 & PROVED & \checkmark & verified \\
\texttt{singular\_series\_pos\_evenPair\_nineHundredFortySix} & Gaps942950 & PROVED & \checkmark & verified \\
\texttt{singular\_series\_pos\_evenPair\_nineHundredFortyTwo} & Gaps942950 & PROVED & \checkmark & verified \\
\texttt{evenPair\_card\_nineHundredFiftyEight} & Gaps952960 & PROVED & \checkmark & verified \\
\texttt{evenPair\_card\_nineHundredFiftyFour} & Gaps952960 & PROVED & \checkmark & verified \\
\texttt{evenPair\_card\_nineHundredFiftySix} & Gaps952960 & PROVED & \checkmark & verified \\
\texttt{evenPair\_card\_nineHundredFiftyTwo} & Gaps952960 & PROVED & \checkmark & verified \\
\texttt{evenPair\_card\_nineHundredSixty} & Gaps952960 & PROVED & \checkmark & verified \\
\texttt{isAdmissible\_evenPair\_nineHundredFiftyEight} & Gaps952960 & PROVED & \checkmark & verified \\
\texttt{isAdmissible\_evenPair\_nineHundredFiftyFour} & Gaps952960 & PROVED & \checkmark & verified \\
\texttt{isAdmissible\_evenPair\_nineHundredFiftySix} & Gaps952960 & PROVED & \checkmark & verified \\
\texttt{isAdmissible\_evenPair\_nineHundredFiftyTwo} & Gaps952960 & PROVED & \checkmark & verified \\
\texttt{isAdmissible\_evenPair\_nineHundredSixty} & Gaps952960 & PROVED & \checkmark & verified \\
\texttt{localFactor\_nineHundredFiftyTwo\_two} & Gaps952960 & PROVED & \checkmark & verified \\
\texttt{localFactor\_nineHundredSixty\_two} & Gaps952960 & PROVED & \checkmark & verified \\
\texttt{nu\_p\_nineHundredFiftyEight} & Gaps952960 & PROVED & \checkmark & verified \\
\texttt{nu\_p\_nineHundredFiftyFour} & Gaps952960 & PROVED & \checkmark & verified \\
\texttt{nu\_p\_nineHundredFiftySix} & Gaps952960 & PROVED & \checkmark & verified \\
\texttt{nu\_p\_nineHundredFiftyTwo} & Gaps952960 & PROVED & \checkmark & verified \\
\texttt{nu\_p\_nineHundredFiftyTwo\_two} & Gaps952960 & PROVED & \checkmark & verified \\
\texttt{nu\_p\_nineHundredSixty} & Gaps952960 & PROVED & \checkmark & verified \\
\texttt{nu\_p\_nineHundredSixty\_two} & Gaps952960 & PROVED & \checkmark & verified \\
\texttt{singular\_series\_finite\_pos\_evenPair\_nineHundredFiftyEight} & Gaps952960 & PROVED & \checkmark & verified \\
\texttt{singular\_series\_finite\_pos\_evenPair\_nineHundredFiftyFour} & Gaps952960 & PROVED & \checkmark & verified \\
\texttt{singular\_series\_finite\_pos\_evenPair\_nineHundredFiftySix} & Gaps952960 & PROVED & \checkmark & verified \\
\texttt{singular\_series\_finite\_pos\_evenPair\_nineHundredFiftyTwo} & Gaps952960 & PROVED & \checkmark & verified \\
\texttt{singular\_series\_finite\_pos\_evenPair\_nineHundredSixty} & Gaps952960 & PROVED & \checkmark & verified \\
\texttt{singular\_series\_pos\_evenPair\_nineHundredFiftyEight} & Gaps952960 & PROVED & \checkmark & verified \\
\texttt{singular\_series\_pos\_evenPair\_nineHundredFiftyFour} & Gaps952960 & PROVED & \checkmark & verified \\
\texttt{singular\_series\_pos\_evenPair\_nineHundredFiftySix} & Gaps952960 & PROVED & \checkmark & verified \\
\texttt{singular\_series\_pos\_evenPair\_nineHundredFiftyTwo} & Gaps952960 & PROVED & \checkmark & verified \\
\texttt{singular\_series\_pos\_evenPair\_nineHundredSixty} & Gaps952960 & PROVED & \checkmark & verified \\
\texttt{evenPair\_card\_nineHundredSeventy} & Gaps962970 & PROVED & \checkmark & verified \\
\texttt{evenPair\_card\_nineHundredSixtyEight} & Gaps962970 & PROVED & \checkmark & verified \\
\texttt{evenPair\_card\_nineHundredSixtyFour} & Gaps962970 & PROVED & \checkmark & verified \\
\texttt{evenPair\_card\_nineHundredSixtySix} & Gaps962970 & PROVED & \checkmark & verified \\
\texttt{evenPair\_card\_nineHundredSixtyTwo} & Gaps962970 & PROVED & \checkmark & verified \\
\texttt{isAdmissible\_evenPair\_nineHundredSeventy} & Gaps962970 & PROVED & \checkmark & verified \\
\texttt{isAdmissible\_evenPair\_nineHundredSixtyEight} & Gaps962970 & PROVED & \checkmark & verified \\
\texttt{isAdmissible\_evenPair\_nineHundredSixtyFour} & Gaps962970 & PROVED & \checkmark & verified \\
\texttt{isAdmissible\_evenPair\_nineHundredSixtySix} & Gaps962970 & PROVED & \checkmark & verified \\
\texttt{isAdmissible\_evenPair\_nineHundredSixtyTwo} & Gaps962970 & PROVED & \checkmark & verified \\
\texttt{localFactor\_nineHundredSeventy\_two} & Gaps962970 & PROVED & \checkmark & verified \\
\texttt{localFactor\_nineHundredSixtyTwo\_two} & Gaps962970 & PROVED & \checkmark & verified \\
\texttt{nu\_p\_nineHundredSeventy} & Gaps962970 & PROVED & \checkmark & verified \\
\texttt{nu\_p\_nineHundredSeventy\_two} & Gaps962970 & PROVED & \checkmark & verified \\
\texttt{nu\_p\_nineHundredSixtyEight} & Gaps962970 & PROVED & \checkmark & verified \\
\texttt{nu\_p\_nineHundredSixtyFour} & Gaps962970 & PROVED & \checkmark & verified \\
\texttt{nu\_p\_nineHundredSixtySix} & Gaps962970 & PROVED & \checkmark & verified \\
\texttt{nu\_p\_nineHundredSixtyTwo} & Gaps962970 & PROVED & \checkmark & verified \\
\texttt{nu\_p\_nineHundredSixtyTwo\_two} & Gaps962970 & PROVED & \checkmark & verified \\
\texttt{singular\_series\_finite\_pos\_evenPair\_nineHundredSeventy} & Gaps962970 & PROVED & \checkmark & verified \\
\texttt{singular\_series\_finite\_pos\_evenPair\_nineHundredSixtyEight} & Gaps962970 & PROVED & \checkmark & verified \\
\texttt{singular\_series\_finite\_pos\_evenPair\_nineHundredSixtyFour} & Gaps962970 & PROVED & \checkmark & verified \\
\texttt{singular\_series\_finite\_pos\_evenPair\_nineHundredSixtySix} & Gaps962970 & PROVED & \checkmark & verified \\
\texttt{singular\_series\_finite\_pos\_evenPair\_nineHundredSixtyTwo} & Gaps962970 & PROVED & \checkmark & verified \\
\texttt{singular\_series\_pos\_evenPair\_nineHundredSeventy} & Gaps962970 & PROVED & \checkmark & verified \\
\texttt{singular\_series\_pos\_evenPair\_nineHundredSixtyEight} & Gaps962970 & PROVED & \checkmark & verified \\
\texttt{singular\_series\_pos\_evenPair\_nineHundredSixtyFour} & Gaps962970 & PROVED & \checkmark & verified \\
\texttt{singular\_series\_pos\_evenPair\_nineHundredSixtySix} & Gaps962970 & PROVED & \checkmark & verified \\
\texttt{singular\_series\_pos\_evenPair\_nineHundredSixtyTwo} & Gaps962970 & PROVED & \checkmark & verified \\
\texttt{evenPair\_card\_nineHundredEighty} & Gaps972980 & PROVED & \checkmark & verified \\
\texttt{evenPair\_card\_nineHundredSeventyEight} & Gaps972980 & PROVED & \checkmark & verified \\
\texttt{evenPair\_card\_nineHundredSeventyFour} & Gaps972980 & PROVED & \checkmark & verified \\
\texttt{evenPair\_card\_nineHundredSeventySix} & Gaps972980 & PROVED & \checkmark & verified \\
\texttt{evenPair\_card\_nineHundredSeventyTwo} & Gaps972980 & PROVED & \checkmark & verified \\
\texttt{isAdmissible\_evenPair\_nineHundredEighty} & Gaps972980 & PROVED & \checkmark & verified \\
\texttt{isAdmissible\_evenPair\_nineHundredSeventyEight} & Gaps972980 & PROVED & \checkmark & verified \\
\texttt{isAdmissible\_evenPair\_nineHundredSeventyFour} & Gaps972980 & PROVED & \checkmark & verified \\
\texttt{isAdmissible\_evenPair\_nineHundredSeventySix} & Gaps972980 & PROVED & \checkmark & verified \\
\texttt{isAdmissible\_evenPair\_nineHundredSeventyTwo} & Gaps972980 & PROVED & \checkmark & verified \\
\texttt{localFactor\_nineHundredEighty\_two} & Gaps972980 & PROVED & \checkmark & verified \\
\texttt{localFactor\_nineHundredSeventyTwo\_two} & Gaps972980 & PROVED & \checkmark & verified \\
\texttt{nu\_p\_nineHundredEighty} & Gaps972980 & PROVED & \checkmark & verified \\
\texttt{nu\_p\_nineHundredEighty\_two} & Gaps972980 & PROVED & \checkmark & verified \\
\texttt{nu\_p\_nineHundredSeventyEight} & Gaps972980 & PROVED & \checkmark & verified \\
\texttt{nu\_p\_nineHundredSeventyFour} & Gaps972980 & PROVED & \checkmark & verified \\
\texttt{nu\_p\_nineHundredSeventySix} & Gaps972980 & PROVED & \checkmark & verified \\
\texttt{nu\_p\_nineHundredSeventyTwo} & Gaps972980 & PROVED & \checkmark & verified \\
\texttt{nu\_p\_nineHundredSeventyTwo\_two} & Gaps972980 & PROVED & \checkmark & verified \\
\texttt{singular\_series\_finite\_pos\_evenPair\_nineHundredEighty} & Gaps972980 & PROVED & \checkmark & verified \\
\texttt{singular\_series\_finite\_pos\_evenPair\_nineHundredSeventyEight} & Gaps972980 & PROVED & \checkmark & verified \\
\texttt{singular\_series\_finite\_pos\_evenPair\_nineHundredSeventyFour} & Gaps972980 & PROVED & \checkmark & verified \\
\texttt{singular\_series\_finite\_pos\_evenPair\_nineHundredSeventySix} & Gaps972980 & PROVED & \checkmark & verified \\
\texttt{singular\_series\_finite\_pos\_evenPair\_nineHundredSeventyTwo} & Gaps972980 & PROVED & \checkmark & verified \\
\texttt{singular\_series\_pos\_evenPair\_nineHundredEighty} & Gaps972980 & PROVED & \checkmark & verified \\
\texttt{singular\_series\_pos\_evenPair\_nineHundredSeventyEight} & Gaps972980 & PROVED & \checkmark & verified \\
\texttt{singular\_series\_pos\_evenPair\_nineHundredSeventyFour} & Gaps972980 & PROVED & \checkmark & verified \\
\texttt{singular\_series\_pos\_evenPair\_nineHundredSeventySix} & Gaps972980 & PROVED & \checkmark & verified \\
\texttt{singular\_series\_pos\_evenPair\_nineHundredSeventyTwo} & Gaps972980 & PROVED & \checkmark & verified \\
\texttt{evenPair\_card\_nineHundredEightyEight} & Gaps982990 & PROVED & \checkmark & verified \\
\texttt{evenPair\_card\_nineHundredEightyFour} & Gaps982990 & PROVED & \checkmark & verified \\
\texttt{evenPair\_card\_nineHundredEightySix} & Gaps982990 & PROVED & \checkmark & verified \\
\texttt{evenPair\_card\_nineHundredEightyTwo} & Gaps982990 & PROVED & \checkmark & verified \\
\texttt{evenPair\_card\_nineHundredNinety} & Gaps982990 & PROVED & \checkmark & verified \\
\texttt{isAdmissible\_evenPair\_nineHundredEightyEight} & Gaps982990 & PROVED & \checkmark & verified \\
\texttt{isAdmissible\_evenPair\_nineHundredEightyFour} & Gaps982990 & PROVED & \checkmark & verified \\
\texttt{isAdmissible\_evenPair\_nineHundredEightySix} & Gaps982990 & PROVED & \checkmark & verified \\
\texttt{isAdmissible\_evenPair\_nineHundredEightyTwo} & Gaps982990 & PROVED & \checkmark & verified \\
\texttt{isAdmissible\_evenPair\_nineHundredNinety} & Gaps982990 & PROVED & \checkmark & verified \\
\texttt{localFactor\_nineHundredEightyTwo\_two} & Gaps982990 & PROVED & \checkmark & verified \\
\texttt{localFactor\_nineHundredNinety\_two} & Gaps982990 & PROVED & \checkmark & verified \\
\texttt{nu\_p\_nineHundredEightyEight} & Gaps982990 & PROVED & \checkmark & verified \\
\texttt{nu\_p\_nineHundredEightyFour} & Gaps982990 & PROVED & \checkmark & verified \\
\texttt{nu\_p\_nineHundredEightySix} & Gaps982990 & PROVED & \checkmark & verified \\
\texttt{nu\_p\_nineHundredEightyTwo} & Gaps982990 & PROVED & \checkmark & verified \\
\texttt{nu\_p\_nineHundredEightyTwo\_two} & Gaps982990 & PROVED & \checkmark & verified \\
\texttt{nu\_p\_nineHundredNinety} & Gaps982990 & PROVED & \checkmark & verified \\
\texttt{nu\_p\_nineHundredNinety\_two} & Gaps982990 & PROVED & \checkmark & verified \\
\texttt{singular\_series\_finite\_pos\_evenPair\_nineHundredEightyEight} & Gaps982990 & PROVED & \checkmark & verified \\
\texttt{singular\_series\_finite\_pos\_evenPair\_nineHundredEightyFour} & Gaps982990 & PROVED & \checkmark & verified \\
\texttt{singular\_series\_finite\_pos\_evenPair\_nineHundredEightySix} & Gaps982990 & PROVED & \checkmark & verified \\
\texttt{singular\_series\_finite\_pos\_evenPair\_nineHundredEightyTwo} & Gaps982990 & PROVED & \checkmark & verified \\
\texttt{singular\_series\_finite\_pos\_evenPair\_nineHundredNinety} & Gaps982990 & PROVED & \checkmark & verified \\
\texttt{singular\_series\_pos\_evenPair\_nineHundredEightyEight} & Gaps982990 & PROVED & \checkmark & verified \\
\texttt{singular\_series\_pos\_evenPair\_nineHundredEightyFour} & Gaps982990 & PROVED & \checkmark & verified \\
\texttt{singular\_series\_pos\_evenPair\_nineHundredEightySix} & Gaps982990 & PROVED & \checkmark & verified \\
\texttt{singular\_series\_pos\_evenPair\_nineHundredEightyTwo} & Gaps982990 & PROVED & \checkmark & verified \\
\texttt{singular\_series\_pos\_evenPair\_nineHundredNinety} & Gaps982990 & PROVED & \checkmark & verified \\
\texttt{evenPair\_card\_nineHundredNinetyEight} & Gaps9921000 & PROVED & \checkmark & verified \\
\texttt{evenPair\_card\_nineHundredNinetyFour} & Gaps9921000 & PROVED & \checkmark & verified \\
\texttt{evenPair\_card\_nineHundredNinetySix} & Gaps9921000 & PROVED & \checkmark & verified \\
\texttt{evenPair\_card\_nineHundredNinetyTwo} & Gaps9921000 & PROVED & \checkmark & verified \\
\texttt{evenPair\_card\_oneThousand} & Gaps9921000 & PROVED & \checkmark & verified \\
\texttt{isAdmissible\_evenPair\_nineHundredNinetyEight} & Gaps9921000 & PROVED & \checkmark & verified \\
\texttt{isAdmissible\_evenPair\_nineHundredNinetyFour} & Gaps9921000 & PROVED & \checkmark & verified \\
\texttt{isAdmissible\_evenPair\_nineHundredNinetySix} & Gaps9921000 & PROVED & \checkmark & verified \\
\texttt{isAdmissible\_evenPair\_nineHundredNinetyTwo} & Gaps9921000 & PROVED & \checkmark & verified \\
\texttt{isAdmissible\_evenPair\_oneThousand} & Gaps9921000 & PROVED & \checkmark & verified \\
\texttt{localFactor\_nineHundredNinetyTwo\_two} & Gaps9921000 & PROVED & \checkmark & verified \\
\texttt{localFactor\_oneThousand\_two} & Gaps9921000 & PROVED & \checkmark & verified \\
\texttt{nu\_p\_nineHundredNinetyEight} & Gaps9921000 & PROVED & \checkmark & verified \\
\texttt{nu\_p\_nineHundredNinetyFour} & Gaps9921000 & PROVED & \checkmark & verified \\
\texttt{nu\_p\_nineHundredNinetySix} & Gaps9921000 & PROVED & \checkmark & verified \\
\texttt{nu\_p\_nineHundredNinetyTwo} & Gaps9921000 & PROVED & \checkmark & verified \\
\texttt{nu\_p\_nineHundredNinetyTwo\_two} & Gaps9921000 & PROVED & \checkmark & verified \\
\texttt{nu\_p\_oneThousand} & Gaps9921000 & PROVED & \checkmark & verified \\
\texttt{nu\_p\_oneThousand\_two} & Gaps9921000 & PROVED & \checkmark & verified \\
\texttt{singular\_series\_finite\_pos\_evenPair\_nineHundredNinetyEight} & Gaps9921000 & PROVED & \checkmark & verified \\
\texttt{singular\_series\_finite\_pos\_evenPair\_nineHundredNinetyFour} & Gaps9921000 & PROVED & \checkmark & verified \\
\texttt{singular\_series\_finite\_pos\_evenPair\_nineHundredNinetySix} & Gaps9921000 & PROVED & \checkmark & verified \\
\texttt{singular\_series\_finite\_pos\_evenPair\_nineHundredNinetyTwo} & Gaps9921000 & PROVED & \checkmark & verified \\
\texttt{singular\_series\_finite\_pos\_evenPair\_oneThousand} & Gaps9921000 & PROVED & \checkmark & verified \\
\texttt{singular\_series\_pos\_evenPair\_nineHundredNinetyEight} & Gaps9921000 & PROVED & \checkmark & verified \\
\texttt{singular\_series\_pos\_evenPair\_nineHundredNinetyFour} & Gaps9921000 & PROVED & \checkmark & verified \\
\texttt{singular\_series\_pos\_evenPair\_nineHundredNinetySix} & Gaps9921000 & PROVED & \checkmark & verified \\
\texttt{singular\_series\_pos\_evenPair\_nineHundredNinetyTwo} & Gaps9921000 & PROVED & \checkmark & verified \\
\texttt{singular\_series\_pos\_evenPair\_oneThousand} & Gaps9921000 & PROVED & \checkmark & verified \\
\texttt{evenPair\_card\_eight} & MoreExamples & PROVED & \checkmark & verified \\
\texttt{evenPair\_card\_four} & MoreExamples & PROVED & \checkmark & verified \\
\texttt{evenPair\_card\_of\_ne\_zero} & MoreExamples & PROVED & \checkmark & verified \\
\texttt{evenPair\_card\_six} & MoreExamples & PROVED & \checkmark & verified \\
\texttt{evenPair\_card\_ten} & MoreExamples & PROVED & \checkmark & verified \\
\texttt{isAdmissible\_evenPair\_eight} & MoreExamples & PROVED & \checkmark & verified \\
\texttt{isAdmissible\_evenPair\_four} & MoreExamples & PROVED & \checkmark & verified \\
\texttt{isAdmissible\_evenPair\_six} & MoreExamples & PROVED & \checkmark & verified \\
\texttt{isAdmissible\_evenPair\_ten} & MoreExamples & PROVED & \checkmark & verified \\
\texttt{localFactorAt\_eight\_two} & MoreExamples & PROVED & \checkmark & verified \\
\texttt{localFactorAt\_evenPair\_two} & MoreExamples & PROVED & \checkmark & verified \\
\texttt{localFactorAt\_four\_five} & MoreExamples & PROVED & \checkmark & verified \\
\texttt{localFactorAt\_four\_odd} & MoreExamples & PROVED & \checkmark & verified \\
\texttt{localFactorAt\_four\_three} & MoreExamples & PROVED & \checkmark & verified \\
\texttt{localFactorAt\_four\_two} & MoreExamples & PROVED & \checkmark & verified \\
\texttt{localFactorAt\_six\_five} & MoreExamples & PROVED & \checkmark & verified \\
\texttt{localFactorAt\_six\_odd\_ne\_three} & MoreExamples & PROVED & \checkmark & verified \\
\texttt{localFactorAt\_six\_three} & MoreExamples & PROVED & \checkmark & verified \\
\texttt{localFactorAt\_six\_two} & MoreExamples & PROVED & \checkmark & verified \\
\texttt{localFactorAt\_ten\_two} & MoreExamples & PROVED & \checkmark & verified \\
\texttt{localFactor\_eight\_two} & MoreExamples & PROVED & \checkmark & verified \\
\texttt{localFactor\_evenPair\_two} & MoreExamples & PROVED & \checkmark & verified \\
\texttt{localFactor\_four\_five} & MoreExamples & PROVED & \checkmark & verified \\
\texttt{localFactor\_four\_odd} & MoreExamples & PROVED & \checkmark & verified \\
\texttt{localFactor\_four\_three} & MoreExamples & PROVED & \checkmark & verified \\
\texttt{localFactor\_four\_two} & MoreExamples & PROVED & \checkmark & verified \\
\texttt{localFactor\_six\_five} & MoreExamples & PROVED & \checkmark & verified \\
\texttt{localFactor\_six\_odd\_ne\_three} & MoreExamples & PROVED & \checkmark & verified \\
\texttt{localFactor\_six\_three} & MoreExamples & PROVED & \checkmark & verified \\
\texttt{localFactor\_six\_two} & MoreExamples & PROVED & \checkmark & verified \\
\texttt{localFactor\_ten\_two} & MoreExamples & PROVED & \checkmark & verified \\
\texttt{nu\_p\_evenPair} & MoreExamples & PROVED & \checkmark & verified \\
\texttt{nu\_p\_evenPair\_odd\_of\_dvd} & MoreExamples & PROVED & \checkmark & verified \\
\texttt{nu\_p\_evenPair\_odd\_of\_not\_dvd} & MoreExamples & PROVED & \checkmark & verified \\
\texttt{nu\_p\_evenPair\_two} & MoreExamples & PROVED & \checkmark & verified \\
\texttt{nu\_p\_four} & MoreExamples & PROVED & \checkmark & verified \\
\texttt{nu\_p\_four\_odd} & MoreExamples & PROVED & \checkmark & verified \\
\texttt{nu\_p\_four\_two} & MoreExamples & PROVED & \checkmark & verified \\
\texttt{nu\_p\_six} & MoreExamples & PROVED & \checkmark & verified \\
\texttt{nu\_p\_six\_odd\_ne\_three} & MoreExamples & PROVED & \checkmark & verified \\
\texttt{nu\_p\_six\_three} & MoreExamples & PROVED & \checkmark & verified \\
\texttt{nu\_p\_six\_two} & MoreExamples & PROVED & \checkmark & verified \\
\texttt{singular\_series\_finite\_pos\_evenPair} & MoreExamples & PROVED & \checkmark & verified \\
\texttt{singular\_series\_finite\_pos\_evenPair\_eight} & MoreExamples & PROVED & \checkmark & verified \\
\texttt{singular\_series\_finite\_pos\_evenPair\_four} & MoreExamples & PROVED & \checkmark & verified \\
\texttt{singular\_series\_finite\_pos\_evenPair\_six} & MoreExamples & PROVED & \checkmark & verified \\
\texttt{singular\_series\_finite\_pos\_evenPair\_ten} & MoreExamples & PROVED & \checkmark & verified \\
\texttt{singular\_series\_pos\_evenPair\_eight} & MoreExamples & PROVED & \checkmark & verified \\
\texttt{singular\_series\_pos\_evenPair\_four} & MoreExamples & PROVED & \checkmark & verified \\
\texttt{singular\_series\_pos\_evenPair\_six} & MoreExamples & PROVED & \checkmark & verified \\
\texttt{singular\_series\_pos\_evenPair\_ten} & MoreExamples & PROVED & \checkmark & verified \\
\texttt{h\_conv\_of\_admissible} & Wire & PROVED & \checkmark & verified \\
\texttt{isAdmissible\_iff} & Wire & PROVED & \checkmark & verified \\
\texttt{isAdmissible\_raw} & Wire & PROVED & \checkmark & verified \\
\texttt{singularSeriesFinite\_tendsto\_pos\_of\_admissible} & Wire & PROVED & \checkmark & verified \\
\texttt{singular\_series\_finite\_pos\_of\_admissible} & Wire & PROVED & \checkmark & verified \\
\texttt{singular\_series\_pos\_of\_admissible} & Wire & PROVED & \checkmark & verified \\
\texttt{singular\_series\_pos\_supersedes\_conditional} & Wire & PROVED & \checkmark & verified \\
\texttt{singular\_series\_pos\_unconditional} & Wire & PROVED & \checkmark & verified \\
\texttt{sg\_11} & SophieGermain & PROVED & \checkmark & verified \\
\texttt{sg\_2} & SophieGermain & PROVED & \checkmark & verified \\
\texttt{sg\_23} & SophieGermain & PROVED & \checkmark & verified \\
\texttt{sg\_29} & SophieGermain & PROVED & \checkmark & verified \\
\texttt{sg\_3} & SophieGermain & PROVED & \checkmark & verified \\
\texttt{sg\_41} & SophieGermain & PROVED & \checkmark & verified \\
\texttt{sg\_5} & SophieGermain & PROVED & \checkmark & verified \\
\texttt{sg\_53} & SophieGermain & PROVED & \checkmark & verified \\
\texttt{sg\_mod\_six} & SophieGermain & PROVED & \checkmark & verified \\
\texttt{sg\_safe\_mod} & SophieGermain & PROVED & \checkmark & verified \\
\texttt{cos\_two\_pi\_div\_five} & Spectral & PROVED & \checkmark & verified \\
\texttt{golden\_in\_cycleSpectrum\_five} & Spectral & PROVED & \checkmark & verified \\
\texttt{golden\_sub\_one\_eq\_two\_cos} & Spectral & PROVED & \checkmark & verified \\
\texttt{golden\_unique\_to\_five} & Spectral & PROVED & \checkmark & verified \\
\texttt{neg\_golden\_in\_C5\_spectrum} & Spectral & PROVED & \checkmark & verified \\
\texttt{two\_cos\_four\_pi\_div\_five} & Spectral & PROVED & \checkmark & verified \\
\texttt{abs\_primeGaussian\_le\_two} & SpectralGate1 & PROVED & \checkmark & verified \\
\texttt{coeFn\_mulLpCLM} & SpectralGate1 & PROVED & \checkmark & verified \\
\texttt{coeFn\_mulLpFun} & SpectralGate1 & PROVED & \checkmark & verified \\
\texttt{coeFn\_primeGaussianMulCLM} & SpectralGate1 & PROVED & \checkmark & verified \\
\texttt{continuous\_primeBump} & SpectralGate1 & PROVED & \checkmark & verified \\
\texttt{continuous\_primeGaussian} & SpectralGate1 & PROVED & \checkmark & verified \\
\texttt{eLpNorm\_mulLpFun\_le} & SpectralGate1 & PROVED & \checkmark & verified \\
\texttt{isSelfAdjoint\_mulLpCLM} & SpectralGate1 & PROVED & \checkmark & verified \\
\texttt{isSelfAdjoint\_primeGaussianMulCLM} & SpectralGate1 & PROVED & \checkmark & verified \\
\texttt{primeBump\_le\_geom} & SpectralGate1 & PROVED & \checkmark & verified \\
\texttt{primeBump\_nonneg} & SpectralGate1 & PROVED & \checkmark & verified \\
\texttt{primeGaussian\_le\_two} & SpectralGate1 & PROVED & \checkmark & verified \\
\texttt{primeGaussian\_nonneg} & SpectralGate1 & PROVED & \checkmark & verified \\
\texttt{primeGaussianℂ\_memLp\_top} & SpectralGate1 & PROVED & \checkmark & verified \\
\texttt{primeGaussianℂ\_norm\_le} & SpectralGate1 & PROVED & \checkmark & verified \\
\texttt{summable\_primeBump} & SpectralGate1 & PROVED & \checkmark & verified \\
\texttt{sigma1\_prime} & SuperperfectNumbers & PROVED & \checkmark & verified \\
\texttt{sigma\_two\_pow} & SuperperfectNumbers & PROVED & \checkmark & verified \\
\texttt{six\_not\_superperfect} & SuperperfectNumbers & PROVED & \checkmark & verified \\
\texttt{superperfect\_16} & SuperperfectNumbers & PROVED & \checkmark & verified \\
\texttt{superperfect\_2} & SuperperfectNumbers & PROVED & \checkmark & verified \\
\texttt{superperfect\_4} & SuperperfectNumbers & PROVED & \checkmark & verified \\
\texttt{superperfect\_64} & SuperperfectNumbers & PROVED & \checkmark & verified \\
\texttt{superperfect\_two\_pow\_of\_mersenne\_prime} & SuperperfectNumbers & PROVED & \checkmark & verified \\
\texttt{brockian\_table\_card} & TransitionKernel & PROVED & \checkmark & verified \\
\texttt{cousin\_pins\_mod\_three} & TransitionKernel & PROVED & \checkmark & verified \\
\texttt{cousin\_run\_cap} & TransitionKernel & PROVED & \checkmark & verified \\
\texttt{forbidden\_transition} & TransitionKernel & PROVED & \checkmark & verified \\
\texttt{gap10\_run\_cap} & TransitionKernel & PROVED & \checkmark & verified \\
\texttt{kernel\_row\_sum} & TransitionKernel & PROVED & \checkmark & verified \\
\texttt{quadruplet\_pins\_mod\_five} & TransitionKernel & PROVED & \checkmark & verified \\
\texttt{sexy\_free\_mod\_three} & TransitionKernel & PROVED & \checkmark & verified \\
\texttt{sexy\_run\_cap} & TransitionKernel & PROVED & \checkmark & verified \\
\texttt{totalSum\_count} & TransitionKernel & PROVED & \checkmark & verified \\
\texttt{totalSum\_eq} & TransitionKernel & PROVED & \checkmark & verified \\
\texttt{twin\_admissible\_singleton} & TransitionKernel & PROVED & \checkmark & verified \\
\texttt{twin\_pins\_mod\_three} & TransitionKernel & PROVED & \checkmark & verified \\
\texttt{twin\_table\_card} & TransitionKernel & PROVED & \checkmark & verified \\
\texttt{isAdmissible\_twinOffsets} & TwinPrimeConstant & PROVED & \checkmark & verified \\
\texttt{localFactorAt\_twin} & TwinPrimeConstant & PROVED & \checkmark & verified \\
\texttt{localFactorAt\_twin\_eq\_tFactor} & TwinPrimeConstant & PROVED & \checkmark & verified \\
\texttt{localFactorAt\_twin\_five} & TwinPrimeConstant & PROVED & \checkmark & verified \\
\texttt{localFactorAt\_twin\_odd} & TwinPrimeConstant & PROVED & \checkmark & verified \\
\texttt{localFactorAt\_twin\_seven} & TwinPrimeConstant & PROVED & \checkmark & verified \\
\texttt{localFactorAt\_twin\_three} & TwinPrimeConstant & PROVED & \checkmark & verified \\
\texttt{localFactorAt\_twin\_two} & TwinPrimeConstant & PROVED & \checkmark & verified \\
\texttt{localFactor\_twin\_eq\_tFactor} & TwinPrimeConstant & PROVED & \checkmark & verified \\
\texttt{localFactor\_twin\_five} & TwinPrimeConstant & PROVED & \checkmark & verified \\
\texttt{localFactor\_twin\_odd} & TwinPrimeConstant & PROVED & \checkmark & verified \\
\texttt{localFactor\_twin\_seven} & TwinPrimeConstant & PROVED & \checkmark & verified \\
\texttt{localFactor\_twin\_three} & TwinPrimeConstant & PROVED & \checkmark & verified \\
\texttt{localFactor\_twin\_two} & TwinPrimeConstant & PROVED & \checkmark & verified \\
\texttt{nu\_p\_twin} & TwinPrimeConstant & PROVED & \checkmark & verified \\
\texttt{nu\_p\_twin\_odd} & TwinPrimeConstant & PROVED & \checkmark & verified \\
\texttt{nu\_p\_twin\_two} & TwinPrimeConstant & PROVED & \checkmark & verified \\
\texttt{singularSeriesFinite\_twin\_eq} & TwinPrimeConstant & PROVED & \checkmark & verified \\
\texttt{singularSeriesFinite\_twin\_of\_two\_le} & TwinPrimeConstant & PROVED & \checkmark & verified \\
\texttt{singularSeriesFinite\_twin\_pos} & TwinPrimeConstant & PROVED & \checkmark & verified \\
\texttt{singularSeries\_twin\_eq\_two\_mul\_constant} & TwinPrimeConstant & PROVED & \checkmark & verified \\
\texttt{singularSeries\_twin\_pos} & TwinPrimeConstant & PROVED & \checkmark & verified \\
\texttt{twinOffsets\_card} & TwinPrimeConstant & PROVED & \checkmark & verified \\
\texttt{twinOffsets\_eq} & TwinPrimeConstant & PROVED & \checkmark & verified \\
\texttt{twinPrimeConstantFinite\_pos} & TwinPrimeConstant & PROVED & \checkmark & verified \\
\texttt{twinPrimeConstant\_pos} & TwinPrimeConstant & PROVED & \checkmark & verified \\
\texttt{twin\_11} & TwinPrimes & PROVED & \checkmark & verified \\
\texttt{twin\_17} & TwinPrimes & PROVED & \checkmark & verified \\
\texttt{twin\_29} & TwinPrimes & PROVED & \checkmark & verified \\
\texttt{twin\_3} & TwinPrimes & PROVED & \checkmark & verified \\
\texttt{twin\_41} & TwinPrimes & PROVED & \checkmark & verified \\
\texttt{twin\_5} & TwinPrimes & PROVED & \checkmark & verified \\
\texttt{twin\_59} & TwinPrimes & PROVED & \checkmark & verified \\
\texttt{twin\_71} & TwinPrimes & PROVED & \checkmark & verified \\
\texttt{twin\_form\_6k} & TwinPrimes & PROVED & \checkmark & verified \\
\texttt{twin\_middle\_div\_six} & TwinPrimes & PROVED & \checkmark & verified \\
\texttt{not\_unitaryPerfect\_28} & UnitaryPerfect & PROVED & \checkmark & verified \\
\texttt{unitaryDivisorSum\_eq\_sigma\_6} & UnitaryPerfect & PROVED & \checkmark & verified \\
\texttt{unitaryPerfect\_6} & UnitaryPerfect & PROVED & \checkmark & verified \\
\texttt{unitaryPerfect\_60} & UnitaryPerfect & PROVED & \checkmark & verified \\
\texttt{unitaryPerfect\_90} & UnitaryPerfect & PROVED & \checkmark & verified \\
\texttt{twelve\_abundant} & WeirdNumbers & PROVED & \checkmark & verified \\
\texttt{twelve\_not\_weird} & WeirdNumbers & PROVED & \checkmark & verified \\
\texttt{twelve\_semiperfect} & WeirdNumbers & PROVED & \checkmark & verified \\
\texttt{twenty\_not\_weird} & WeirdNumbers & PROVED & \checkmark & verified \\
\texttt{weird\_70} & WeirdNumbers & PROVED & \checkmark & verified \\
\texttt{weird\_836} & WeirdNumbers & PROVED & \checkmark & verified \\
\texttt{green\_identity\_integral} & Weyl & PROVED & \checkmark & verified \\
\texttt{lagrange\_identity} & Weyl & PROVED & \checkmark & verified \\
\texttt{wronskian\_const\_one\_witness} & Weyl & PROVED & \checkmark & verified \\
\texttt{wronskian\_hasDerivAt} & Weyl & PROVED & \checkmark & verified \\
\texttt{wronskian\_isConst} & Weyl & PROVED & \checkmark & verified \\
\texttt{coeff\_ne\_zero} & Bridge & PROVED & \checkmark & verified \\
\texttt{continuous\_y} & Bridge & PROVED & \checkmark & verified \\
\texttt{continuous\_y'} & Bridge & PROVED & \checkmark & verified \\
\texttt{global\_boundary\_identity} & Bridge & PROVED & \checkmark & verified \\
\texttt{hasDerivAt\_wronskianConj} & Bridge & PROVED & \checkmark & verified \\
\texttt{integrable\_normSq} & Bridge & PROVED & \checkmark & verified \\
\texttt{integrable\_normSq\_add} & Bridge & PROVED & \checkmark & verified \\
\texttt{integrable\_wronskianConj\_deriv} & Bridge & PROVED & \checkmark & verified \\
\texttt{integrable\_y'\_normSq} & Bridge & PROVED & \checkmark & verified \\
\texttt{integrable\_y\_normSq} & Bridge & PROVED & \checkmark & verified \\
\texttt{integral\_wronskianConj\_eq} & Bridge & PROVED & \checkmark & verified \\
\texttt{integral\_wronskianConj\_eq\_mul} & Bridge & PROVED & \checkmark & verified \\
\texttt{limUnder\_wronskianConj\_atBot\_eq\_zero} & Bridge & PROVED & \checkmark & verified \\
\texttt{limUnder\_wronskianConj\_atTop\_eq\_zero} & Bridge & PROVED & \checkmark & verified \\
\texttt{no\_nonzero\_L2\_solution} & Bridge & PROVED & \checkmark & verified \\
\texttt{norm\_wronskianConj\_le} & Bridge & PROVED & \checkmark & verified \\
\texttt{not\_integrableOn\_const\_pos\_Ici} & Bridge & PROVED & \checkmark & verified \\
\texttt{not\_integrableOn\_const\_pos\_Iic} & Bridge & PROVED & \checkmark & verified \\
\texttt{tendsto\_wronskianConj\_atBot} & Bridge & PROVED & \checkmark & verified \\
\texttt{tendsto\_wronskianConj\_atBot\_zero} & Bridge & PROVED & \checkmark & verified \\
\texttt{tendsto\_wronskianConj\_atTop} & Bridge & PROVED & \checkmark & verified \\
\texttt{tendsto\_wronskianConj\_atTop\_zero} & Bridge & PROVED & \checkmark & verified \\
\texttt{wronskianConj\_eq\_two\_I\_im} & Bridge & PROVED & \checkmark & verified \\
\texttt{apply\_ne\_I\_smul} & Cayley & PROVED & \checkmark & verified \\
\texttt{apply\_ne\_neg\_I\_smul} & Cayley & PROVED & \checkmark & verified \\
\texttt{deficiencySpace\_eq\_bot\_iff} & Cayley & PROVED & \checkmark & verified \\
\texttt{essentiallySelfAdjoint\_iff} & Cayley & PROVED & \checkmark & verified \\
\texttt{mem\_orthogonal\_rangeSMulSub\_iff} & Cayley & PROVED & \checkmark & verified \\
\texttt{mem\_rangeSMulSub} & Cayley & PROVED & \checkmark & verified \\
\texttt{norm\_add\_I\_smul\_eq} & Cayley & PROVED & \checkmark & verified \\
\texttt{essSelfAdjoint\_of\_dense\_ranges} & Chain & PROVED & \checkmark & verified \\
\texttt{eq\_univ\_of\_dense\_isClosed} & ClosedRange & PROVED & \checkmark & verified \\
\texttt{range\_eq\_top\_of\_essentiallySelfAdjoint\_of\_isClosed\_ranges} & ClosedRange & PROVED & \checkmark & verified \\
\texttt{smulPMap\_isClosed} & ClosedRange & PROVED & \checkmark & verified \\
\texttt{smulPMap\_isClosed\_and\_symmetric} & ClosedRange & PROVED & \checkmark & verified \\
\texttt{dense\_rangeSMulSub\_of\_le} & ClosedRangeClosure & PROVED & \checkmark & verified \\
\texttt{isClosed\_rangeAddI\_and\_rangeSubI} & ClosedRangeClosure & PROVED & \checkmark & verified \\
\texttt{isClosed\_rangeSMulSub\_of\_isClosed\_of\_isSymmetric} & ClosedRangeClosure & PROVED & \checkmark & verified \\
\texttt{rangeSMulSub\_mono} & ClosedRangeClosure & PROVED & \checkmark & verified \\
\texttt{adjoint\_domain\_le\_closure\_domain\_of\_essentiallySelfAdjoint\_of\_rangeAddI} & ClosedShiftedRanges & PROVED & \checkmark & verified \\
\texttt{closure\_eq\_adjoint\_of\_essentiallySelfAdjoint} & ClosedShiftedRanges & PROVED & \checkmark & verified \\
\texttt{closure\_eq\_adjoint\_of\_essentiallySelfAdjoint\_of\_isClosed\_ranges} & ClosedShiftedRanges & PROVED & \checkmark & verified \\
\texttt{closure\_isSelfAdjoint\_of\_essentiallySelfAdjoint} & ClosedShiftedRanges & PROVED & \checkmark & verified \\
\texttt{closure\_isSelfAdjoint\_of\_essentiallySelfAdjoint\_of\_isClosed\_ranges} & ClosedShiftedRanges & PROVED & \checkmark & verified \\
\texttt{closure\_shifted\_ranges\_eq\_univ\_of\_essentiallySelfAdjoint} & ClosedShiftedRanges & PROVED & \checkmark & verified \\
\texttt{closure\_shifted\_ranges\_eq\_univ\_of\_essentiallySelfAdjoint\_of\_isClosed} & ClosedShiftedRanges & PROVED & \checkmark & verified \\
\texttt{dense\_closure\_shifted\_ranges\_of\_essentiallySelfAdjoint} & ClosedShiftedRanges & PROVED & \checkmark & verified \\
\texttt{isClosed\_closure\_shifted\_ranges} & ClosedShiftedRanges & PROVED & \checkmark & verified \\
\texttt{rangeAddI\_le\_closure} & ClosedShiftedRanges & PROVED & \checkmark & verified \\
\texttt{rangeSubI\_le\_closure} & ClosedShiftedRanges & PROVED & \checkmark & verified \\
\texttt{adjoint\_isClosed'} & Closure & PROVED & \checkmark & verified \\
\texttt{closure\_eq\_self\_of\_isClosed} & Closure & PROVED & \checkmark & verified \\
\texttt{inner\_adjoint\_left} & Closure & PROVED & \checkmark & verified \\
\texttt{isClosed\_deficiencySet} & Closure & PROVED & \checkmark & verified \\
\texttt{mem\_deficiencySet\_iff\_mem\_deficiencySpace} & Closure & PROVED & \checkmark & verified \\
\texttt{smulPMap\_adjoint\_isClosed} & Closure & PROVED & \checkmark & verified \\
\texttt{smulPMap\_dense} & Closure & PROVED & \checkmark & verified \\
\texttt{smulPMap\_isClosable} & Closure & PROVED & \checkmark & verified \\
\texttt{symmetric\_adjoint\_eq} & Closure & PROVED & \checkmark & verified \\
\texttt{symmetric\_closure\_le\_adjoint} & Closure & PROVED & \checkmark & verified \\
\texttt{symmetric\_domain\_le\_adjoint\_domain} & Closure & PROVED & \checkmark & verified \\
\texttt{symmetric\_isClosable} & Closure & PROVED & \checkmark & verified \\
\texttt{symmetric\_le\_adjoint} & Closure & PROVED & \checkmark & verified \\
\texttt{isConfining} & Confining & PROVED & \checkmark & verified \\
\texttt{bound\_C\_blocks\_zeros\_outside\_ball} & Confining & PROVED & \checkmark & verified \\
\texttt{brockian\_realizer\_admits\_no\_finite\_bound} & Confining & PROVED & \checkmark & verified \\
\texttt{confiningCandidate\_not\_bddAbove} & Confining & PROVED & \checkmark & verified \\
\texttt{decaying\_not\_isConfining} & Confining & PROVED & \checkmark & verified \\
\texttt{gate1\_vs\_confining\_shape} & Confining & PROVED & \checkmark & verified \\
\texttt{isConfining\_iff\_tendsto} & Confining & PROVED & \checkmark & verified \\
\texttt{isConfining\_not\_bddAbove} & Confining & PROVED & \checkmark & verified \\
\texttt{isConfining\_unbounded} & Confining & PROVED & \checkmark & verified \\
\texttt{isConfining\_unboundedMultiplierShape} & Confining & PROVED & \checkmark & verified \\
\texttt{no\_brockian\_eigenvector\_outside\_bound} & Confining & PROVED & \checkmark & verified \\
\texttt{not\_both\_decaying\_and\_confining} & Confining & PROVED & \checkmark & verified \\
\texttt{not\_isConfining\_of\_abs\_le} & Confining & PROVED & \checkmark & verified \\
\texttt{primeGaussian\_blocks\_zeros\_outside\_two} & Confining & PROVED & \checkmark & verified \\
\texttt{primeGaussian\_is\_decaying} & Confining & PROVED & \checkmark & verified \\
\texttt{primeGaussian\_not\_isConfining} & Confining & PROVED & \checkmark & verified \\
\texttt{primeGaussian\_not\_unboundedMultiplierShape} & Confining & PROVED & \checkmark & verified \\
\texttt{quadraticCandidate\_V} & Confining & PROVED & \checkmark & verified \\
\texttt{quadratic\_isConfining} & Confining & PROVED & \checkmark & verified \\
\texttt{quadratic\_unboundedMultiplierShape} & Confining & PROVED & \checkmark & verified \\
\texttt{brockian\_eigenvalue\_real\_of\_candidate} & ConfiningShape & PROVED & \checkmark & verified \\
\texttt{clm\_bound\_blocks\_zeros\_outside\_ball} & ConfiningShape & PROVED & \checkmark & verified \\
\texttt{clm\_bound\_no\_brockian\_eigenvector} & ConfiningShape & PROVED & \checkmark & verified \\
\texttt{compactResolventShape\_of\_candidate} & ConfiningShape & PROVED & \checkmark & verified \\
\texttt{compactResolventShape\_of\_isConfining} & ConfiningShape & PROVED & \checkmark & verified \\
\texttt{confining\_needed\_for\_shape\_and\_unbounded\_spectrum} & ConfiningShape & PROVED & \checkmark & verified \\
\texttt{eigenvalueCountingAsymptotic\_of\_matches} & ConfiningShape & PROVED & \checkmark & verified \\
\texttt{eigenvalueCountingMatchesNT\_comm} & ConfiningShape & PROVED & \checkmark & verified \\
\texttt{gate1\_vs\_confining\_shape\_package} & ConfiningShape & PROVED & \checkmark & verified \\
\texttt{not\_compactResolventShape\_of\_abs\_le} & ConfiningShape & PROVED & \checkmark & verified \\
\texttt{not\_discreteSpectrumCandidate\_of\_bound} & ConfiningShape & PROVED & \checkmark & verified \\
\texttt{not\_point\_spectrum\_unbounded\_of\_bound} & ConfiningShape & PROVED & \checkmark & verified \\
\texttt{primeGaussian\_clm\_blocks\_large\_zeros} & ConfiningShape & PROVED & \checkmark & verified \\
\texttt{primeGaussian\_not\_compactResolventShape} & ConfiningShape & PROVED & \checkmark & verified \\
\texttt{quadratic\_compactResolventShape} & ConfiningShape & PROVED & \checkmark & verified \\
\texttt{quadratic\_isConfining\_reexport} & ConfiningShape & PROVED & \checkmark & verified \\
\texttt{realizer\_admits\_no\_finite\_bound} & ConfiningShape & PROVED & \checkmark & verified \\
\texttt{spectrum\_real\_of\_symm} & ConfiningShape & PROVED & \checkmark & verified \\
\texttt{continuous\_growing\_exp\_normSq} & ConstMass & PROVED & \checkmark & verified \\
\texttt{exists\_growing\_mode\_limitPointRadius} & ConstMass & PROVED & \checkmark & verified \\
\texttt{growing\_exp\_IsLimitPointRadius} & ConstMass & PROVED & \checkmark & verified \\
\texttt{growing\_exp\_mass\_monotone} & ConstMass & PROVED & \checkmark & verified \\
\texttt{growing\_exp\_mass\_tendsto\_atTop} & ConstMass & PROVED & \checkmark & verified \\
\texttt{growing\_exp\_radius\_tendsto\_zero} & ConstMass & PROVED & \checkmark & verified \\
\texttt{integral\_exp\_mul\_tendsto\_atTop} & ConstMass & PROVED & \checkmark & verified \\
\texttt{integral\_nonneg\_mono} & ConstMass & PROVED & \checkmark & verified \\
\texttt{normSq\_cexp\_eq} & ConstMass & PROVED & \checkmark & verified \\
\texttt{coeFn\_potentialMul} & DeficiencyODE & PROVED & \checkmark & verified \\
\texttt{deficiencyRepresentsODE\_of\_weakRegularity} & DeficiencyODE & CONDITIONAL & \checkmark & verified \\
\texttt{inner\_g\_potential} & DeficiencyODE & PROVED & \checkmark & verified \\
\texttt{inner\_g\_schwartz} & DeficiencyODE & PROVED & \checkmark & verified \\
\texttt{inner\_g\_schwartz\_D2} & DeficiencyODE & PROVED & \checkmark & verified \\
\texttt{schrodinger\_essentiallySelfAdjoint\_of\_weakRegularity} & DeficiencyODE & CONDITIONAL & \checkmark & verified \\
\texttt{atTop\_of\_radius\_tendsto\_zero} & Dichotomy & PROVED & \checkmark & verified \\
\texttt{limitPointRadius\_radius\_tendsto\_zero} & Dichotomy & PROVED & \checkmark & verified \\
\texttt{limitPoint\_or\_limitCircle\_radius} & Dichotomy & PROVED & \checkmark & verified \\
\texttt{radius\_tendsto\_zero\_iff} & Dichotomy & PROVED & \checkmark & verified \\
\texttt{radius\_tendsto\_zero\_of\_atTop} & Dichotomy & PROVED & \checkmark & verified \\
\texttt{tendsto\_atTop\_of\_monotone\_not\_bddAbove} & Dichotomy & PROVED & \checkmark & verified \\
\texttt{boundary\_L2\_identity} & Disk & PROVED & \checkmark & verified \\
\texttt{circle\_key} & Disk & PROVED & \checkmark & verified \\
\texttt{green\_identity\_integral} & Disk & PROVED & \checkmark & verified \\
\texttt{integral\_normSq\_mono} & Disk & PROVED & \checkmark & verified \\
\texttt{lagrange\_identity} & Disk & PROVED & \checkmark & verified \\
\texttt{radius\_formula} & Disk & PROVED & \checkmark & verified \\
\texttt{weyl\_disk\_circle} & Disk & PROVED & \checkmark & verified \\
\texttt{weyl\_nested\_circle} & Disk & PROVED & \checkmark & verified \\
\texttt{weyl\_radius\_antitone} & Disk & PROVED & \checkmark & verified \\
\texttt{wronskian\_hasDerivAt} & Disk & PROVED & \checkmark & verified \\
\texttt{wronskian\_isConst} & Disk & PROVED & \checkmark & verified \\
\texttt{diskRadius\_eq\_weylRadius} & DiskBridge & PROVED & \checkmark & verified \\
\texttt{diskRadius\_tendsto\_zero\_of\_limitPointRadius} & DiskBridge & PROVED & \checkmark & verified \\
\texttt{diskRadius\_tendsto\_zero\_of\_mass\_atTop} & DiskBridge & PROVED & \checkmark & verified \\
\texttt{mass\_monotone} & DiskBridge & PROVED & \checkmark & verified \\
\texttt{clm\_deficiency\_eq\_bot} & ESA & PROVED & \checkmark & verified \\
\texttt{clm\_dense} & ESA & PROVED & \checkmark & verified \\
\texttt{clm\_domain} & ESA & PROVED & \checkmark & verified \\
\texttt{clm\_essentiallySelfAdjoint} & ESA & PROVED & \checkmark & verified \\
\texttt{clm\_isSymmetric} & ESA & PROVED & \checkmark & verified \\
\texttt{id\_essentiallySelfAdjoint} & ESA & PROVED & \checkmark & verified \\
\texttt{vec\_eq\_zero\_of\_inner} & ESA & PROVED & \checkmark & verified \\
\texttt{adjoint\_closure} & Extension & PROVED & \checkmark & verified \\
\texttt{adjoint\_eigen\_I\_eq\_zero\_of\_essSA} & Extension & PROVED & \checkmark & verified \\
\texttt{adjoint\_eigen\_neg\_I\_eq\_zero\_of\_essSA} & Extension & PROVED & \checkmark & verified \\
\texttt{adjoint\_graph\_topologicalClosure\_eq} & Extension & PROVED & \checkmark & verified \\
\texttt{closure\_eigenvalue\_im\_zero} & Extension & PROVED & \checkmark & verified \\
\texttt{closure\_isSymmetric} & Extension & PROVED & \checkmark & verified \\
\texttt{closure\_le\_of\_isClosed\_extension} & Extension & PROVED & \checkmark & verified \\
\texttt{closure\_le\_of\_isSelfAdjoint\_extension} & Extension & PROVED & \checkmark & verified \\
\texttt{eigenvalue\_im\_zero} & Extension & PROVED & \checkmark & verified \\
\texttt{essentiallySelfAdjoint\_iff'} & Extension & PROVED & \checkmark & verified \\
\texttt{isSelfAdjoint\_closure\_iff\_eq\_adjoint} & Extension & PROVED & \checkmark & verified \\
\texttt{isSymmetric\_of\_le\_adjoint} & Extension & PROVED & \checkmark & verified \\
\texttt{le\_closure\_le\_adjoint} & Extension & PROVED & \checkmark & verified \\
\texttt{smulPMap\_closure\_eigenvalue\_im\_zero} & Extension & PROVED & \checkmark & verified \\
\texttt{smulPMap\_closure\_isSymmetric} & Extension & PROVED & \checkmark & verified \\
\texttt{dense\_domain\_position} & FourierMultiplier & PROVED & \checkmark & verified \\
\texttt{essentiallySelfAdjoint\_position} & FourierMultiplier & PROVED & \checkmark & verified \\
\texttt{dense\_domain\_transfer} & FourierMultiplier & PROVED & \checkmark & verified \\
\texttt{dense\_rangeAddI\_transfer\_iff} & FourierMultiplier & PROVED & \checkmark & verified \\
\texttt{dense\_rangeSMulSub\_transfer\_iff} & FourierMultiplier & PROVED & \checkmark & verified \\
\texttt{dense\_rangeSubI\_transfer\_iff} & FourierMultiplier & PROVED & \checkmark & verified \\
\texttt{essentiallySelfAdjoint\_of\_multiplier\_dense\_ranges} & FourierMultiplier & PROVED & \checkmark & verified \\
\texttt{essentiallySelfAdjoint\_of\_multiplier\_esa} & FourierMultiplier & PROVED & \checkmark & verified \\
\texttt{essentiallySelfAdjoint\_of\_multiplier\_shift\_dense} & FourierMultiplier & PROVED & \checkmark & verified \\
\texttt{FreeLaplacianModel} & FreeLaplacian & PROVED & \checkmark & verified \\
\texttt{dense\_domain\_vadd\_clm} & FreeLaplacian & PROVED & \checkmark & verified \\
\texttt{essentiallySelfAdjoint\_conjCLM} & FreeLaplacian & PROVED & \checkmark & verified \\
\texttt{essentiallySelfAdjoint\_conj\_id} & FreeLaplacian & PROVED & \checkmark & verified \\
\texttt{id\_essentiallySelfAdjoint} & FreeLaplacian & PROVED & \checkmark & verified \\
\texttt{inner\_map\_symm} & FreeLaplacian & PROVED & \checkmark & verified \\
\texttt{inner\_symm\_map} & FreeLaplacian & PROVED & \checkmark & verified \\
\texttt{isSelfAdjoint\_conjCLM} & FreeLaplacian & PROVED & \checkmark & verified \\
\texttt{isSymmetric\_vadd\_clm} & FreeLaplacian & PROVED & \checkmark & verified \\
\texttt{vadd\_clm\_domain} & FreeLaplacian & PROVED & \checkmark & verified \\
\texttt{conjCLM\_essentiallySelfAdjoint} & FreeLaplacian2 & PROVED & \checkmark & verified \\
\texttt{conjCLM\_toPMap\_essentiallySelfAdjoint\_iff} & FreeLaplacian2 & PROVED & \checkmark & verified \\
\texttt{dense\_map\_iff} & FreeLaplacian2 & PROVED & \checkmark & verified \\
\texttt{essentiallySelfAdjoint\_transfer} & FreeLaplacian2 & PROVED & \checkmark & verified \\
\texttt{freeLaplacian\_essentiallySelfAdjoint\_of\_fourier} & FreeLaplacian2 & CONDITIONAL & \checkmark & verified \\
\texttt{isSelfAdjoint\_diagonal} & FreeLaplacian2 & PROVED & \checkmark & verified \\
\texttt{isSelfAdjoint\_multCLM} & FreeLaplacian2 & PROVED & \checkmark & verified \\
\texttt{multCLM\_essentiallySelfAdjoint} & FreeLaplacian2 & PROVED & \checkmark & verified \\
\texttt{rangeSMulSub\_image} & FreeLaplacian2 & PROVED & \checkmark & verified \\
\texttt{sqMult\_conj\_essentiallySelfAdjoint} & FreeLaplacian2 & PROVED & \checkmark & verified \\
\texttt{sqMult\_essentiallySelfAdjoint} & FreeLaplacian2 & PROVED & \checkmark & verified \\
\texttt{add\_primeGaussian\_dense\_range\_sub} & Gate1Bounded & PROVED & \checkmark & verified \\
\texttt{add\_primeGaussian\_essentiallySelfAdjoint} & Gate1Bounded & PROVED & \checkmark & verified \\
\texttt{add\_primeGaussian\_isSelfAdjoint} & Gate1Bounded & PROVED & \checkmark & verified \\
\texttt{primeGaussianMul\_dense\_range\_sub} & Gate1Bounded & PROVED & \checkmark & verified \\
\texttt{primeGaussianMul\_essentiallySelfAdjoint} & Gate1Bounded & PROVED & \checkmark & verified \\
\texttt{harmonicOscillatorPMap\_dense} & HarmonicOscillator & PROVED & \checkmark & verified \\
\texttt{harmonicOscillatorPMap\_isSymmetric} & HarmonicOscillator & PROVED & \checkmark & verified \\
\texttt{harmonicOscillatorPMap\_toFun\_ofInjective} & HarmonicOscillator & PROVED & \checkmark & verified \\
\texttt{harmonicOscillator\_confining\_shape} & HarmonicOscillator & PROVED & \checkmark & verified \\
\texttt{harmonicOscillator\_core\_package} & HarmonicOscillator & PROVED & \checkmark & verified \\
\texttt{oscillatorCoreMap\_expanded} & HarmonicOscillator & PROVED & \checkmark & verified \\
\texttt{oscillatorCoreMap\_symm} & HarmonicOscillator & PROVED & \checkmark & verified \\
\texttt{quadraticMul\_symm} & HarmonicOscillator & PROVED & \checkmark & verified \\
\texttt{quadratic\_hasTemperateGrowth} & HarmonicOscillator & PROVED & \checkmark & verified \\
\texttt{dense\_range\_add\_sub\_of\_selfAdjoint} & Kato & PROVED & \checkmark & verified \\
\texttt{isSelfAdjoint\_add} & Kato & PROVED & \checkmark & verified \\
\texttt{boundedPerturbationTransfer\_of\_resolvent\_norm\_lt\_one} & KatoNeumann & PROVED & \checkmark & verified \\
\texttt{katoFactor\_rightInverse\_of\_norm\_lt\_one} & KatoNeumann & PROVED & \checkmark & verified \\
\texttt{rangeAddI\_perturb\_eq\_univ\_of\_resolvent\_norm\_lt\_one} & KatoNeumann & PROVED & \checkmark & verified \\
\texttt{rangeSubI\_perturb\_eq\_univ\_of\_resolvent\_norm\_lt\_one} & KatoNeumann & PROVED & \checkmark & verified \\
\texttt{rightResolvent\_perturb\_of\_norm\_lt\_one} & KatoNeumann & PROVED & \checkmark & verified \\
\texttt{boundedPerturbationTransfer\_of\_resolvent\_norm\_mul\_lt\_one} & KatoNeumannEstimates & PROVED & \checkmark & verified \\
\texttt{katoFactor\_rightInverse\_of\_norm\_mul\_lt\_one} & KatoNeumannEstimates & PROVED & \checkmark & verified \\
\texttt{norm\_comp\_lt\_one\_of\_norm\_mul\_lt\_one} & KatoNeumannEstimates & PROVED & \checkmark & verified \\
\texttt{rangeAddI\_perturb\_eq\_univ\_of\_resolvent\_norm\_mul\_lt\_one} & KatoNeumannEstimates & PROVED & \checkmark & verified \\
\texttt{rangeSubI\_perturb\_eq\_univ\_of\_resolvent\_norm\_mul\_lt\_one} & KatoNeumannEstimates & PROVED & \checkmark & verified \\
\texttt{rightResolvent\_perturb\_of\_norm\_mul\_lt\_one} & KatoNeumannEstimates & PROVED & \checkmark & verified \\
\texttt{boundedPerturbationTransfer\_iff\_dense\_ranges} & KatoRangeDensity & PROVED & \checkmark & verified \\
\texttt{boundedPerturbationTransfer\_of\_essentiallySelfAdjoint\_perturb} & KatoRangeDensity & PROVED & \checkmark & verified \\
\texttt{boundedPerturbationTransfer\_zero\_of\_essentiallySelfAdjoint} & KatoRangeDensity & PROVED & \checkmark & verified \\
\texttt{boundedSelfAdjoint\_perturb\_dense\_ranges} & KatoRangeDensity & PROVED & \checkmark & verified \\
\texttt{boundedSelfAdjoint\_perturb\_essentiallySelfAdjoint} & KatoRangeDensity & PROVED & \checkmark & verified \\
\texttt{dense\_rangeAddI\_perturb\_of\_transfer} & KatoRangeDensity & PROVED & \checkmark & verified \\
\texttt{dense\_rangeSubI\_perturb\_of\_transfer} & KatoRangeDensity & PROVED & \checkmark & verified \\
\texttt{essentiallySelfAdjoint\_perturb\_iff\_transfer} & KatoRangeDensity & PROVED & \checkmark & verified \\
\texttt{perturb\_zero\_eq} & KatoRangeDensity & PROVED & \checkmark & verified \\
\texttt{perturbed\_ranges\_eq\_univ\_of\_transfer\_of\_isClosed} & KatoRangeDensity & PROVED & \checkmark & verified \\
\texttt{rangeAddI\_perturb\_eq\_univ\_of\_transfer\_of\_isClosed} & KatoRangeDensity & PROVED & \checkmark & verified \\
\texttt{rangeSubI\_perturb\_eq\_univ\_of\_transfer\_of\_isClosed} & KatoRangeDensity & PROVED & \checkmark & verified \\
\texttt{boundedPerturbationTransfer\_of\_resolvent\_factors} & KatoRellichScaffold & PROVED & \checkmark & verified \\
\texttt{boundedPerturbationTransfer\_zero\_of\_rightResolvents} & KatoRellichScaffold & PROVED & \checkmark & verified \\
\texttt{dense\_rangeSMulSub\_of\_rightResolvent} & KatoRellichScaffold & PROVED & \checkmark & verified \\
\texttt{factorRightInverse\_zero} & KatoRellichScaffold & PROVED & \checkmark & verified \\
\texttt{rangeAddI\_perturb\_eq\_univ\_of\_resolvent\_factor} & KatoRellichScaffold & PROVED & \checkmark & verified \\
\texttt{rangeSMulSub\_eq\_univ\_of\_rightResolvent} & KatoRellichScaffold & PROVED & \checkmark & verified \\
\texttt{rangeSubI\_perturb\_eq\_univ\_of\_resolvent\_factor} & KatoRellichScaffold & PROVED & \checkmark & verified \\
\texttt{rightResolvent\_perturb\_of\_factor\_rightInverse} & KatoRellichScaffold & PROVED & \checkmark & verified \\
\texttt{boundedPerturbationTransfer\_of\_ranges\_eq\_univ} & KatoRellichTransfer & PROVED & \checkmark & verified \\
\texttt{boundedPerturbationTransfer\_of\_resolvent\_comp\_norm\_lt\_one} & KatoRellichTransfer & PROVED & \checkmark & verified \\
\texttt{boundedPerturbationTransfer\_of\_resolvent\_product\_norm\_lt\_one} & KatoRellichTransfer & PROVED & \checkmark & verified \\
\texttt{essentiallySelfAdjoint\_perturb\_of\_ranges\_eq\_univ} & KatoRellichTransfer & PROVED & \checkmark & verified \\
\texttt{essentiallySelfAdjoint\_perturb\_of\_resolvent\_norm\_lt\_one} & KatoRellichTransfer & PROVED & \checkmark & verified \\
\texttt{essentiallySelfAdjoint\_perturb\_of\_resolvent\_norm\_mul\_lt\_one} & KatoRellichTransfer & PROVED & \checkmark & verified \\
\texttt{essentiallySelfAdjoint\_perturb\_of\_resolvent\_norm\_mul\_lt\_one\_via\_chain} & KatoRellichTransfer & PROVED & \checkmark & verified \\
\texttt{perturbed\_ranges\_eq\_univ\_of\_resolvent\_norm\_lt\_one} & KatoRellichTransfer & PROVED & \checkmark & verified \\
\texttt{perturbed\_ranges\_eq\_univ\_of\_resolvent\_norm\_mul\_lt\_one} & KatoRellichTransfer & PROVED & \checkmark & verified \\
\texttt{boundedPerturbationTransfer\_of\_unitShiftRightResolvents\_norm\_lt\_one} & KatoResolventConstruction & PROVED & \checkmark & verified \\
\texttt{essentiallySelfAdjoint\_perturb\_of\_unitShiftRightResolvents\_norm\_lt\_one} & KatoResolventConstruction & PROVED & \checkmark & verified \\
\texttt{essentiallySelfAdjoint\_perturb\_of\_unitShiftRightResolvents\_norm\_lt\_one\_via\_chain} & KatoResolventConstruction & PROVED & \checkmark & verified \\
\texttt{norm\_mul\_Radd\_lt\_one\_of\_unitShiftRightResolvents} & KatoResolventConstruction & PROVED & \checkmark & verified \\
\texttt{norm\_mul\_Rsub\_lt\_one\_of\_unitShiftRightResolvents} & KatoResolventConstruction & PROVED & \checkmark & verified \\
\texttt{boundedPerturbationTransfer} & KatoResolventPackage & PROVED & \checkmark & verified \\
\texttt{norm\_mul\_add\_lt\_one} & KatoResolventPackage & PROVED & \checkmark & verified \\
\texttt{norm\_mul\_sub\_lt\_one} & KatoResolventPackage & PROVED & \checkmark & verified \\
\texttt{perturbed\_ranges\_eq\_univ} & KatoResolventPackage & PROVED & \checkmark & verified \\
\texttt{essentiallySelfAdjoint\_perturb\_of\_resolventAtI} & KatoResolventPackage & PROVED & \checkmark & verified \\
\texttt{essentiallySelfAdjoint\_perturb\_of\_resolventAtI\_via\_chain} & KatoResolventPackage & PROVED & \checkmark & verified \\
\texttt{dense\_range\_add\_sub\_of\_selfAdjoint} & KatoTarget & PROVED & \checkmark & verified \\
\texttt{isSelfAdjoint\_add} & KatoTarget & PROVED & \checkmark & verified \\
\texttt{boundedPerturbationTransfer\_clm} & KatoUnbounded & PROVED & \checkmark & verified \\
\texttt{essentiallySelfAdjoint\_perturb} & KatoUnbounded & CONDITIONAL & \checkmark & verified \\
\texttt{essentiallySelfAdjoint\_perturb\_iff} & KatoUnbounded & PROVED & \checkmark & verified \\
\texttt{perturb\_apply} & KatoUnbounded & PROVED & \checkmark & verified \\
\texttt{perturb\_apply\_ne\_I\_smul} & KatoUnbounded & PROVED & \checkmark & verified \\
\texttt{perturb\_apply\_ne\_neg\_I\_smul} & KatoUnbounded & PROVED & \checkmark & verified \\
\texttt{perturb\_clm\_eq} & KatoUnbounded & PROVED & \checkmark & verified \\
\texttt{perturb\_clm\_essentiallySelfAdjoint} & KatoUnbounded & PROVED & \checkmark & verified \\
\texttt{perturb\_dense\_domain} & KatoUnbounded & PROVED & \checkmark & verified \\
\texttt{perturb\_domain} & KatoUnbounded & PROVED & \checkmark & verified \\
\texttt{perturb\_isSymmetric} & KatoUnbounded & PROVED & \checkmark & verified \\
\texttt{perturb\_norm\_add\_I\_smul\_eq} & KatoUnbounded & PROVED & \checkmark & verified \\
\texttt{perturb\_norm\_sub\_smul\_ge} & KatoUnbounded & PROVED & \checkmark & verified \\
\texttt{const\_potential\_isLimitPoint} & LP & PROVED & \checkmark & verified \\
\texttt{exists\_sqrt\_pos\_re} & LP & PROVED & \checkmark & verified \\
\texttt{indep\_of\_wronskian\_ne\_zero} & LP & PROVED & \checkmark & verified \\
\texttt{not\_integrableOn\_exp\_mul\_Ioi} & LP & PROVED & \checkmark & verified \\
\texttt{wronskian\_hasDerivAt} & LP & PROVED & \checkmark & verified \\
\texttt{wronskian\_isConst} & LP & PROVED & \checkmark & verified \\
\texttt{L2NearInfty\_of\_strong} & LimitPointContinuous & PROVED & \checkmark & verified \\
\texttt{abs\_le\_const\_potential} & LimitPointContinuous & PROVED & \checkmark & verified \\
\texttt{const\_continuous\_isLimitPoint} & LimitPointContinuous & PROVED & \checkmark & verified \\
\texttt{contBounded\_const\_isLimitPoint} & LimitPointContinuous & PROVED & \checkmark & verified \\
\texttt{contBounded\_isLimitPoint\_of\_FS\_obligation} & LimitPointContinuous & PROVED & \checkmark & verified \\
\texttt{continuous\_const\_potential} & LimitPointContinuous & PROVED & \checkmark & verified \\
\texttt{continuous\_deriv\_of\_isSolution} & LimitPointContinuous & PROVED & \checkmark & verified \\
\texttt{continuous\_of\_isSolution} & LimitPointContinuous & PROVED & \checkmark & verified \\
\texttt{isLimitPoint\_of\_LP\_const} & LimitPointContinuous & PROVED & \checkmark & verified \\
\texttt{isLimitPoint\_of\_fundSystem\_member\_not\_L2} & LimitPointContinuous & PROVED & \checkmark & verified \\
\texttt{isLimitPoint\_of\_fundSystem\_not\_both\_L2} & LimitPointContinuous & PROVED & \checkmark & verified \\
\texttt{isLimitPoint\_of\_non\_L2\_solution} & LimitPointContinuous & PROVED & \checkmark & verified \\
\texttt{isSolution\_iff\_isSolutionOn} & LimitPointContinuous & PROVED & \checkmark & verified \\
\texttt{no\_nonzero\_global\_L2\_of\_contBounded} & LimitPointContinuous & PROVED & \checkmark & verified \\
\texttt{normSq\_y''\_le\_of\_bounded} & LimitPointContinuous & PROVED & \checkmark & verified \\
\texttt{norm\_y''\_le\_of\_bounded} & LimitPointContinuous & PROVED & \checkmark & verified \\
\texttt{wronskian\_const\_of\_solutions} & LimitPointContinuous & PROVED & \checkmark & verified \\
\texttt{wronskian\_hasDerivAt} & LimitPointContinuous & PROVED & \checkmark & verified \\
\texttt{wronskian\_isConst} & LimitPointContinuous & PROVED & \checkmark & verified \\
\texttt{FreeMulModel} & MulReal & PROVED & \checkmark & verified \\
\texttt{essentiallySelfAdjoint\_mulOp} & MulReal & PROVED & \checkmark & verified \\
\texttt{isSelfAdjoint\_mulOp} & MulReal & PROVED & \checkmark & verified \\
\texttt{add\_clm\_mul\_essentiallySelfAdjoint} & MulReal & PROVED & \checkmark & verified \\
\texttt{add\_constMul\_essentiallySelfAdjoint} & MulReal & PROVED & \checkmark & verified \\
\texttt{add\_mulLpCLM\_essentiallySelfAdjoint} & MulReal & PROVED & \checkmark & verified \\
\texttt{add\_mulLpCLM\_isSelfAdjoint} & MulReal & PROVED & \checkmark & verified \\
\texttt{constFreeMulModel\_essentiallySelfAdjoint} & MulReal & PROVED & \checkmark & verified \\
\texttt{constFun\_memLp\_top} & MulReal & PROVED & \checkmark & verified \\
\texttt{constFun\_norm\_le} & MulReal & PROVED & \checkmark & verified \\
\texttt{constFun\_real} & MulReal & PROVED & \checkmark & verified \\
\texttt{constMul\_essentiallySelfAdjoint} & MulReal & PROVED & \checkmark & verified \\
\texttt{isSelfAdjoint\_constMulCLM} & MulReal & PROVED & \checkmark & verified \\
\texttt{isSelfAdjoint\_oneMulCLM} & MulReal & PROVED & \checkmark & verified \\
\texttt{mulLpCLM\_essentiallySelfAdjoint} & MulReal & PROVED & \checkmark & verified \\
\texttt{oneMul\_essentiallySelfAdjoint} & MulReal & PROVED & \checkmark & verified \\
\texttt{primeGaussianFreeMulModel\_essentiallySelfAdjoint} & MulReal & PROVED & \checkmark & verified \\
\texttt{primeGaussianMul\_essentiallySelfAdjoint} & MulReal & PROVED & \checkmark & verified \\
\texttt{IsSymmetric} & Operator & PROVED & \checkmark & verified \\
\texttt{eq\_zero\_of\_apply\_eq\_smul} & Operator & PROVED & \checkmark & verified \\
\texttt{im\_eq\_zero\_of\_apply\_eq\_smul} & Operator & PROVED & \checkmark & verified \\
\texttt{inner\_apply} & Operator & PROVED & \checkmark & verified \\
\texttt{inner\_self\_im} & Operator & PROVED & \checkmark & verified \\
\texttt{norm\_sub\_smul\_ge} & Operator & PROVED & \checkmark & verified \\
\texttt{mem\_deficiencySpace\_iff} & Operator & PROVED & \checkmark & verified \\
\texttt{smulPMap\_domain} & Operator & PROVED & \checkmark & verified \\
\texttt{smulPMap\_isSymmetric} & Operator & PROVED & \checkmark & verified \\
\texttt{brockian\_eigenvalue\_norm} & OperatorChoice & PROVED & \checkmark & verified \\
\texttt{norm\_eigenvalue\_le\_of\_bound} & OperatorChoice & PROVED & \checkmark & verified \\
\texttt{norm\_eigenvalue\_le\_opNorm} & OperatorChoice & PROVED & \checkmark & verified \\
\texttt{norm\_primeGaussianMulCLM\_le} & OperatorChoice & PROVED & \checkmark & verified \\
\texttt{not\_eigenvalue\_of\_bound\_lt} & OperatorChoice & PROVED & \checkmark & verified \\
\texttt{not\_realize\_zero\_of\_bound\_lt} & OperatorChoice & PROVED & \checkmark & verified \\
\texttt{not\_realize\_zero\_of\_toPMap} & OperatorChoice & PROVED & \checkmark & verified \\
\texttt{primeGaussianMulCLM\_opNorm\_le\_two} & OperatorChoice & PROVED & \checkmark & verified \\
\texttt{primeGaussian\_not\_realize\_large\_zero} & OperatorChoice & PROVED & \checkmark & verified \\
\texttt{rh\_operator\_needs\_unbounded\_spectrum} & OperatorChoice & PROVED & \checkmark & verified \\
\texttt{finiteMultiplicity\_nonzero\_spectrum} & OscillatorDiscrete & PROVED & \checkmark & verified \\
\texttt{nonzero\_spectrum\_is\_point\_spectrum} & OscillatorDiscrete & PROVED & \checkmark & verified \\
\texttt{finiteDimensional\_eigenspace\_of\_isCompact} & OscillatorDiscrete & PROVED & \checkmark & verified \\
\texttt{finiteMultiplicityNonzeroSpectrum\_of\_isCompact} & OscillatorDiscrete & PROVED & \checkmark & verified \\
\texttt{harmonicOscillator\_resolvent\_finiteMultiplicity\_of\_weightedRellich} & OscillatorDiscrete & PROVED & \checkmark & verified \\
\texttt{harmonicOscillator\_resolvent\_nonzero\_spectrum\_of\_weightedRellich} & OscillatorDiscrete & PROVED & \checkmark & verified \\
\texttt{nonzeroSpectrumIsPointSpectrum\_of\_isCompact} & OscillatorDiscrete & PROVED & \checkmark & verified \\
\texttt{radius\_pos\_limit\_of\_mass\_finite} & RadiusDichotomy & PROVED & \checkmark & verified \\
\texttt{radius\_tendsto\_zero\_iff\_counterexample} & RadiusDichotomy & PROVED & \checkmark & verified \\
\texttt{radius\_tendsto\_zero\_iff\_of\_pos} & RadiusDichotomy & PROVED & \checkmark & verified \\
\texttt{radius\_to\_zero\_of\_mass\_infinite} & RadiusDichotomy & PROVED & \checkmark & verified \\
\texttt{essentiallySelfAdjoint\_perturb\_of\_essentiallySelfAdjoint\_of\_isClosed\_ranges} & ResolventFromESA & PROVED & \checkmark & verified \\
\texttt{exists\_rightResolvent\_of\_range\_eq\_univ} & ResolventFromESA & PROVED & \checkmark & verified \\
\texttt{norm\_le\_inv\_im\_mul\_norm\_shifted} & ResolventFromESA & PROVED & \checkmark & verified \\
\texttt{norm\_rightResolventOfRangeEqUniv\_le} & ResolventFromESA & PROVED & \checkmark & verified \\
\texttt{rightResolventOfRangeEqUniv\_maps\_domain} & ResolventFromESA & PROVED & \checkmark & verified \\
\texttt{rightResolventOfRangeEqUniv\_right\_inverse} & ResolventFromESA & PROVED & \checkmark & verified \\
\texttt{rightResolvents\_of\_essentiallySelfAdjoint\_of\_isClosed\_ranges} & ResolventFromESA & PROVED & \checkmark & verified \\
\texttt{rightResolvents\_of\_surjective\_shifted\_ranges} & ResolventFromESA & PROVED & \checkmark & verified \\
\texttt{shiftedMap\_surjective\_of\_range\_eq\_univ} & ResolventFromESA & PROVED & \checkmark & verified \\
\texttt{shiftedResolventLinearMap\_bound} & ResolventFromESA & PROVED & \checkmark & verified \\
\texttt{shiftedResolventLinearMap\_maps\_domain} & ResolventFromESA & PROVED & \checkmark & verified \\
\texttt{shiftedResolventLinearMap\_right\_inverse} & ResolventFromESA & PROVED & \checkmark & verified \\
\texttt{shiftedRightInverseLinearMap\_spec} & ResolventFromESA & PROVED & \checkmark & verified \\
\texttt{deficiencySpace\_eq\_bot\_of\_ode\_bridge} & SchrodingerESA & PROVED & \checkmark & verified \\
\texttt{dense\_ranges\_of\_ode\_bridge} & SchrodingerESA & PROVED & \checkmark & verified \\
\texttt{essSelfAdjoint\_of\_ode\_bridge\_via\_chain} & SchrodingerESA & PROVED & \checkmark & verified \\
\texttt{essentiallySelfAdjoint\_of\_ode\_bridge} & SchrodingerESA & PROVED & \checkmark & verified \\
\texttt{free\_plus\_primeGaussian\_essentiallySelfAdjoint} & SchrodingerESA & PROVED & \checkmark & verified \\
\texttt{primeGaussian\_essentiallySelfAdjoint} & SchrodingerESA & PROVED & \checkmark & verified \\
\texttt{schrodinger\_closure\_eq\_adjoint} & SchrodingerGate1Closed & PROVED & \checkmark & verified \\
\texttt{schrodinger\_closure\_isSelfAdjoint} & SchrodingerGate1Closed & PROVED & \checkmark & verified \\
\texttt{schrodinger\_closure\_shifted\_ranges\_eq\_univ} & SchrodingerGate1Closed & PROVED & \checkmark & verified \\
\texttt{schrodinger\_core\_essentiallySelfAdjoint} & SchrodingerGate1Closed & PROVED & \checkmark & verified \\
\texttt{freeCoreMap\_apply} & SchrodingerGate1Final & PROVED & \checkmark & verified \\
\texttt{freeSchrodingerPMap\_dense} & SchrodingerGate1Final & PROVED & \checkmark & verified \\
\texttt{freeSchrodingerPMap\_domain} & SchrodingerGate1Final & PROVED & \checkmark & verified \\
\texttt{freeSchrodingerPMap\_isSymmetric} & SchrodingerGate1Final & PROVED & \checkmark & verified \\
\texttt{freeSchrodingerPMap\_toFun\_ofInjective} & SchrodingerGate1Final & PROVED & \checkmark & verified \\
\texttt{schrodingerPMap\_eq\_perturb\_free} & SchrodingerGate1Final & PROVED & \checkmark & verified \\
\texttt{schrodinger\_essentiallySelfAdjoint\_iff\_katoTransfer} & SchrodingerGate1Final & PROVED & \checkmark & verified \\
\texttt{schrodinger\_essentiallySelfAdjoint\_of\_distributionalPrimitiveIdentity} & SchrodingerGate1Final & PROVED & \checkmark & verified \\
\texttt{schrodinger\_essentiallySelfAdjoint\_of\_katoTransfer} & SchrodingerGate1Final & PROVED & \checkmark & verified \\
\texttt{schrodinger\_essentiallySelfAdjoint\_of\_weakSolutionVanishing} & SchrodingerGate1Final & PROVED & \checkmark & verified \\
\texttt{D2\_apply} & SchrodingerMinimal & PROVED & \checkmark & verified \\
\texttt{Lconj\_apply} & SchrodingerMinimal & PROVED & \checkmark & verified \\
\texttt{coeFn\_schwartzToL2} & SchrodingerMinimal & PROVED & \checkmark & verified \\
\texttt{coreMap\_apply} & SchrodingerMinimal & PROVED & \checkmark & verified \\
\texttt{coreMap\_symm} & SchrodingerMinimal & PROVED & \checkmark & verified \\
\texttt{inner\_toLp} & SchrodingerMinimal & PROVED & \checkmark & verified \\
\texttt{isSelfAdjoint\_potentialMulCLM} & SchrodingerMinimal & PROVED & \checkmark & verified \\
\texttt{kinetic\_symm} & SchrodingerMinimal & PROVED & \checkmark & verified \\
\texttt{potential\_symm} & SchrodingerMinimal & PROVED & \checkmark & verified \\
\texttt{schrodingerPMap\_dense} & SchrodingerMinimal & PROVED & \checkmark & verified \\
\texttt{schrodingerPMap\_domain} & SchrodingerMinimal & PROVED & \checkmark & verified \\
\texttt{schrodingerPMap\_isSymmetric} & SchrodingerMinimal & PROVED & \checkmark & verified \\
\texttt{schrodingerPMap\_toFun\_ofInjective} & SchrodingerMinimal & PROVED & \checkmark & verified \\
\texttt{schrodinger\_essentiallySelfAdjoint\_of\_ode} & SchrodingerMinimal & CONDITIONAL & \checkmark & verified \\
\texttt{schwartzToL2\_apply} & SchrodingerMinimal & PROVED & \checkmark & verified \\
\texttt{schwartzToL2\_injective} & SchrodingerMinimal & PROVED & \checkmark & verified \\
\texttt{schwartz\_ibp1} & SchrodingerMinimal & PROVED & \checkmark & verified \\
\texttt{schwartz\_ibp2} & SchrodingerMinimal & PROVED & \checkmark & verified \\
\texttt{eigenvalue\_im\_zero} & SymmetryPackage & PROVED & \checkmark & verified \\
\texttt{injective\_of\_im\_ne\_zero} & SymmetryPackage & PROVED & \checkmark & verified \\
\texttt{norm\_sub\_smul\_ge} & SymmetryPackage & PROVED & \checkmark & verified \\
\texttt{not\_eigenvalue\_of\_im\_ne\_zero} & SymmetryPackage & PROVED & \checkmark & verified \\
\texttt{quadratic\_form\_im\_zero} & SymmetryPackage & PROVED & \checkmark & verified \\
\texttt{smulPMap\_isSymmetric} & SymmetryPackage & PROVED & \checkmark & verified \\
\texttt{smulPMap\_not\_eigenvalue\_of\_im\_ne\_zero} & SymmetryPackage & PROVED & \checkmark & verified \\
\texttt{smulPMap\_quadratic\_form\_im\_zero} & SymmetryPackage & PROVED & \checkmark & verified \\
\texttt{deficiencySpace\_eq\_bot\_of\_closure\_isSelfAdjoint} & Upstream & PROVED & \checkmark & verified \\
\texttt{essentiallySelfAdjoint\_iff\_closure\_isSelfAdjoint} & Upstream & PROVED & \checkmark & verified \\
\texttt{isCompactOperator} & WeightedRellich & PROVED & \checkmark & verified \\
\texttt{isCompact\_closure\_image\_closedBall} & WeightedRellich & PROVED & \checkmark & verified \\
\texttt{compact\_resolvents\_of\_factorizations} & WeightedRellich & PROVED & \checkmark & verified \\
\texttt{N\_op\_tendsto\_atTop\_of\_WeylLawMatch} & WeylLawTarget & PROVED & \checkmark & verified \\
\texttt{N\_op\_tendsto\_atTop\_of\_matches\_rvm} & WeylLawTarget & PROVED & \checkmark & verified \\
\texttt{WeylLawMatchDiff\_iff\_eigenvalueCountingMatchesNT} & WeylLawTarget & PROVED & \checkmark & verified \\
\texttt{WeylLawMatch\_iff\_eigenvalueCountingAsymptotic} & WeylLawTarget & PROVED & \checkmark & verified \\
\texttt{WeylLawMatch\_of\_diff} & WeylLawTarget & PROVED & \checkmark & verified \\
\texttt{WeylLawMatch\_of\_eigenvalueCountingMatchesNT} & WeylLawTarget & PROVED & \checkmark & verified \\
\texttt{counting\_diverges\_of\_candidate} & WeylLawTarget & CONDITIONAL & \checkmark & verified \\
\texttt{counting\_diverges\_of\_discrete\_and\_WeylLawMatch} & WeylLawTarget & CONDITIONAL & \checkmark & verified \\
\texttt{counting\_diverges\_of\_discrete\_and\_rvm} & WeylLawTarget & CONDITIONAL & \checkmark & verified \\
\texttt{counting\_diverges\_of\_exists} & WeylLawTarget & CONDITIONAL & \checkmark & verified \\
\texttt{point\_spectrum\_unbounded\_of\_candidate} & WeylLawTarget & PROVED & \checkmark & verified \\
\texttt{riemannVonMangoldtMain\_tendsto\_atTop} & WeylLawTarget & PROVED & \checkmark & verified \\
\texttt{dense\_domain\_position} & WeylPlancherelScaffold & PROVED & \checkmark & verified \\
\texttt{dense\_range\_addI\_position} & WeylPlancherelScaffold & PROVED & \checkmark & verified \\
\texttt{dense\_range\_subI\_position} & WeylPlancherelScaffold & PROVED & \checkmark & verified \\
\texttt{essentiallySelfAdjoint\_position} & WeylPlancherelScaffold & PROVED & \checkmark & verified \\
\texttt{essentiallySelfAdjoint\_position\_of\_multiplier\_esa} & WeylPlancherelScaffold & PROVED & \checkmark & verified \\
\texttt{isPlancherel} & WeylPlancherelScaffold & PROVED & \checkmark & verified \\
\texttt{essentiallySelfAdjoint\_of\_plancherel\_multiplier\_dense\_ranges} & WeylPlancherelScaffold & PROVED & \checkmark & verified \\
\texttt{isPlancherelUnitary} & WeylPlancherelScaffold & PROVED & \checkmark & verified \\
\texttt{conjugateLp\_eq\_star} & WeylWeakEnergy & PROVED & \checkmark & verified \\
\texttt{conjugateRHSLp\_eq\_potential\_sub} & WeylWeakEnergy & PROVED & \checkmark & verified \\
\texttt{continuous\_laplacianSymbol} & WeylWeakEnergy & PROVED & \checkmark & verified \\
\texttt{fourier\_laplacian\_ae} & WeylWeakEnergy & PROVED & \checkmark & verified \\
\texttt{im\_inner\_conjugateRHSLp} & WeylWeakEnergy & PROVED & \checkmark & verified \\
\texttt{im\_inner\_laplacian\_eq\_zero} & WeylWeakEnergy & PROVED & \checkmark & verified \\
\texttt{laplacianSymbol\_mul\_locallyIntegrable} & WeylWeakEnergy & PROVED & \checkmark & verified \\
\texttt{potentialMulCLM\_coeFn} & WeylWeakEnergy & PROVED & \checkmark & verified \\
\texttt{schrodinger\_essentiallySelfAdjoint\_of\_continuous\_bounded} & WeylWeakEnergy & PROVED & \checkmark & verified \\
\texttt{weakSolutionVanishing\_of\_continuous\_bounded} & WeylWeakEnergy & PROVED & \checkmark & verified \\
\texttt{integral\_identity\_of\_hasDerivAt} & WeylWeakPrimitiveClassical & PROVED & \checkmark & verified \\
\texttt{memLp\_of\_l2\_representative} & WeylWeakPrimitiveClassical & PROVED & \checkmark & verified \\
\texttt{primitiveModel\_of\_distributionalPrimitiveData} & WeylWeakPrimitiveClassical & PROVED & \checkmark & verified \\
\texttt{primitiveModel\_zero} & WeylWeakPrimitiveClassical & PROVED & \checkmark & verified \\
\texttt{rhs\_integral\_identity\_of\_hasDerivAt} & WeylWeakPrimitiveClassical & PROVED & \checkmark & verified \\
\texttt{weakToPrimitiveRegularity\_of\_continuous\_bounded\_of\_distributionalPrimitiveIdentity} & WeylWeakPrimitiveClassical & PROVED & \checkmark & verified \\
\texttt{weakToPrimitiveRegularity\_of\_continuous\_bounded\_of\_weakSolutionVanishing} & WeylWeakPrimitiveClassical & PROVED & \checkmark & verified \\
\texttt{weakToPrimitiveRegularity\_of\_distributionalPrimitiveIdentity} & WeylWeakPrimitiveClassical & PROVED & \checkmark & verified \\
\texttt{weakToPrimitiveRegularity\_of\_weakSolutionVanishing} & WeylWeakPrimitiveClassical & PROVED & \checkmark & verified \\
\texttt{primitiveModel\_of\_distributionalPrimitiveIdentity} & WeylWeakPrimitiveLocal & PROVED & \checkmark & verified \\
\texttt{primitiveModel\_y\_eq\_coe\_of\_canonical\_continuous} & WeylWeakPrimitiveLocal & PROVED & \checkmark & verified \\
\texttt{schrodinger\_essentiallySelfAdjoint\_of\_distributional\_primitives} & WeylWeakPrimitiveLocal & PROVED & \checkmark & verified \\
\texttt{weakToIntegralRegularity\_of\_distributional\_primitives} & WeylWeakPrimitiveLocal & PROVED & \checkmark & verified \\
\texttt{weakToPrimitiveRegularity\_of\_distributional\_primitives} & WeylWeakPrimitiveLocal & PROVED & \checkmark & verified \\
\texttt{conjugateLp\_coeFn} & WeylWeakRegularityClosed & PROVED & \checkmark & verified \\
\texttt{conjugateLp\_toTemperedDistribution\_apply} & WeylWeakRegularityClosed & PROVED & \checkmark & verified \\
\texttt{conjugateRHSLp\_coeFn} & WeylWeakRegularityClosed & PROVED & \checkmark & verified \\
\texttt{conjugateRHSLp\_toTemperedDistribution\_apply} & WeylWeakRegularityClosed & PROVED & \checkmark & verified \\
\texttt{conjugateRHS\_memLp} & WeylWeakRegularityClosed & PROVED & \checkmark & verified \\
\texttt{conjugateRepresentative\_memLp} & WeylWeakRegularityClosed & PROVED & \checkmark & verified \\
\texttt{laplacian\_eq\_secondDeriv} & WeylWeakRegularityClosed & PROVED & \checkmark & verified \\
\texttt{lineDeriv\_one\_eq\_deriv} & WeylWeakRegularityClosed & PROVED & \checkmark & verified \\
\texttt{secondDeriv\_conjugateLp\_eq\_conjugateRHSLp} & WeylWeakRegularityClosed & PROVED & \checkmark & verified \\
\texttt{classicalL2Representative\_of\_integralModel} & WeylWeakRegularityCore & PROVED & \checkmark & verified \\
\texttt{continuous\_representatives\_eq\_of\_ae} & WeylWeakRegularityCore & PROVED & \checkmark & verified \\
\texttt{continuous\_schrodingerRHS} & WeylWeakRegularityCore & PROVED & \checkmark & verified \\
\texttt{integralModel\_hasDerivAt\_y} & WeylWeakRegularityCore & PROVED & \checkmark & verified \\
\texttt{integralModel\_hasDerivAt\_yPrime} & WeylWeakRegularityCore & PROVED & \checkmark & verified \\
\texttt{integralModel\_isL2Solution} & WeylWeakRegularityCore & PROVED & \checkmark & verified \\
\texttt{integralModel\_representative\_unique} & WeylWeakRegularityCore & PROVED & \checkmark & verified \\
\texttt{schrodinger\_essentiallySelfAdjoint\_of\_weakToIntegral} & WeylWeakRegularityCore & PROVED & \checkmark & verified \\
\texttt{weakToClassicalRegularity\_of\_weakToIntegral} & WeylWeakRegularityCore & PROVED & \checkmark & verified \\
\texttt{coeFn\_integrableOn\_Icc} & WeylWeakRegularityDischarge & PROVED & \checkmark & verified \\
\texttt{coeFn\_locallyIntegrable} & WeylWeakRegularityDischarge & PROVED & \checkmark & verified \\
\texttt{continuous\_rhs\_primitive\_of\_coeFn} & WeylWeakRegularityDischarge & PROVED & \checkmark & verified \\
\texttt{integralModel\_of\_primitiveModel} & WeylWeakRegularityDischarge & PROVED & \checkmark & verified \\
\texttt{locallyIntegrable\_continuous\_primitive} & WeylWeakRegularityDischarge & PROVED & \checkmark & verified \\
\texttt{locallyIntegrable\_intervalIntegrable} & WeylWeakRegularityDischarge & PROVED & \checkmark & verified \\
\texttt{schrodingerRHS\_coeFn\_locallyIntegrable} & WeylWeakRegularityDischarge & PROVED & \checkmark & verified \\
\texttt{schrodinger\_essentiallySelfAdjoint\_of\_weakToPrimitive} & WeylWeakRegularityDischarge & PROVED & \checkmark & verified \\
\texttt{weakToClassicalRegularity\_of\_weakToPrimitive} & WeylWeakRegularityDischarge & PROVED & \checkmark & verified \\
\texttt{weakToIntegralRegularity\_of\_weakToPrimitive} & WeylWeakRegularityDischarge & PROVED & \checkmark & verified \\
\texttt{classicalL2Representative\_eq\_zero\_of\_bounded\_nonreal} & WeylWeakRegularityScaffold & PROVED & \checkmark & verified \\
\texttt{deficiencyRepresentsODE\_of\_weakToClassical} & WeylWeakRegularityScaffold & PROVED & \checkmark & verified \\
\texttt{deficiencyVector\_weakSchrodingerEquation} & WeylWeakRegularityScaffold & PROVED & \checkmark & verified \\
\texttt{existingWeakRegularity\_of\_weakToClassical} & WeylWeakRegularityScaffold & PROVED & \checkmark & verified \\
\texttt{schrodinger\_essentiallySelfAdjoint\_of\_weakToClassical} & WeylWeakRegularityScaffold & PROVED & \checkmark & verified \\
\texttt{weakToClassicalRegularity\_iff\_existing} & WeylWeakRegularityScaffold & PROVED & \checkmark & verified \\
\texttt{weakToClassicalRegularity\_of\_existing} & WeylWeakRegularityScaffold & PROVED & \checkmark & verified \\
\texttt{prime\_dvd\_factorial\_add\_one} & WilsonPrimes & PROVED & \checkmark & verified \\
\texttt{seven\_not\_wilson} & WilsonPrimes & PROVED & \checkmark & verified \\
\texttt{wilson\_13} & WilsonPrimes & PROVED & \checkmark & verified \\
\texttt{wilson\_5} & WilsonPrimes & PROVED & \checkmark & verified \\
\texttt{completedRiemannZeta\_functional\_equation} & XiFunctionalEquation & PROVED & \checkmark & verified \\
\texttt{riemannXi\_apply} & XiFunctionalEquation & PROVED & \checkmark & verified \\
\texttt{riemannXi\_eq\_zero\_iff\_zeta\_zero\_of\_mem\_critical\_strip} & XiFunctionalEquation & PROVED & \checkmark & verified \\
\texttt{riemannXi\_functional\_equation} & XiFunctionalEquation & PROVED & \checkmark & verified \\
\texttt{zeta\_zero\_of\_riemannXi\_zero} & XiFunctionalEquation & PROVED & \checkmark & verified \\
\texttt{zeta\_zero\_one\_sub\_of\_mem\_critical\_strip} & XiFunctionalEquation & PROVED & \checkmark & verified \\
\bottomrule
\end{longtable}


\begin{thebibliography}{9}
\bibitem{mathlib} The mathlib Community, \emph{The Lean Mathematical Library}, CPP 2020.
  \url{https://github.com/leanprover-community/mathlib4}.
\bibitem{aristotle} Harmonic, \emph{Aristotle: Automated Theorem Proving for Lean}.
  \url{https://aristotle.harmonic.fun}.
\bibitem{axle} Axiom, \emph{AXLE: A Cloud Infrastructure for Lean 4 Theorem Proving Utilities}.
  \url{https://axle.axiommath.ai}.
\end{thebibliography}

\end{document}
